\section*{Appendix}
\label{sec:appendix}

This appendix collects the diagnostic and explanatory figures that support the
main-text results without expanding the confirmatory claim set. Figures
\ref{fig:appendix_implementation_drag}--\ref{fig:appendix_seed_correlation}
describe the frozen primary evaluation and the post-holdout ensemble mechanism;
Figures \ref{fig:appendix_synthetic_marginals}--\ref{fig:appendix_optimizer}
document synthetic-data and training behaviour; and Figures
\ref{fig:appendix_causal_economics}--\ref{fig:appendix_compressed_reconciliation}
provide additional post-holdout causal evidence. The evidence classification is
stated explicitly in every caption.

\subsection*{A. Primary Evaluation and Implementation Diagnostics}

Figure \ref{fig:appendix_implementation_drag} places turnover and realised
implementation drag on a common strategy axis. It shows that the static,
rolling, and dynamic NN-vine optimisers incur substantially greater trading
intensity than the RL ensemble. Equal weight has zero target turnover by
construction and therefore supplies an implementation baseline rather than an
actively rebalanced comparator.

\begin{figure}[!tbp]
\centering
% AUTO-GENERATED. DO NOT EDIT BY HAND.
% Figure: Turnover and implementation drag
% Evidence class: frozen_confirmatory_primary_evaluation
% Regenerate with publication_pipeline_draft/generate_publication_tikz.py.
\begin{tikzpicture}
\begin{groupplot}[
 group style={group size=2 by 1, horizontal sep=0.70cm,
   group name=implementation},
 publication axis, height=0.43\linewidth,
 ytick={1,2,3,4,5,6,7},
 y dir=reverse, ymin=0.5, ymax=7.5,
]
\nextgroupplot[
 width=0.56\linewidth,
 title={Trading intensity}, xlabel={Mean monthly turnover},
 yticklabels={Equal weight,Shrinkage MV,DCC--GARCH,Static vine,Rolling vine,Dynamic NN-vine,NN-vine TD3},
]
\addplot[xbar, bar width=5pt, fill=pubNavy!70, draw=pubNavy]
 coordinates {(0,1) (0.05195489657,2) (0.1534098552,3) (1.050859561,4) (1.236889067,5) (1.087010719,6) (0.3169771455,7)};
\nextgroupplot[
 width=0.30\linewidth,
 title={Realized cost drag}, xlabel={Total-return drag (pp)},
 yticklabels=\empty,
]
\addplot[only marks, mark=diamond*, mark size=2.7pt, pubRose]
 coordinates {(0,1) (0.1701913512,2) (0.5461043281,3) (5.908050628,4) (5.80145347,5) (5.955308169,6) (1.215260924,7)};
\end{groupplot}
\end{tikzpicture}

\caption{Implementation intensity and realised performance drag in the frozen
primary evaluation. Bars report mean monthly full-$L_1$ turnover and diamonds
report the percentage-point reduction in total return caused by the common
transaction and financing cost model. Evidence class: frozen confirmatory
primary evaluation.}
\label{fig:appendix_implementation_drag}
\end{figure}

Figure \ref{fig:appendix_monthly_excess} displays when relative performance was
earned. The variation around zero makes clear that the ranking is not generated
by a uniform monthly advantage and motivates the paired time-series inference
used in the main text.

\begin{figure}[!tbp]
\centering
% AUTO-GENERATED. DO NOT EDIT BY HAND.
% Figure: Monthly excess returns
% Evidence class: frozen_confirmatory_primary_evaluation
% Regenerate with publication_pipeline_draft/generate_publication_tikz.py.
\begin{tikzpicture}
\begin{axis}[
  publication axis,
  width=0.88\linewidth,
  height=0.34\linewidth,
  date coordinates in=x,
  xticklabel={\year--\month},
  xticklabel style={rotate=30, anchor=north east},
  xlabel={Holding-period end}, ylabel={Excess return (pp)},
  title={Month-by-month relative performance},
  xmin=2024-08-01, xmax=2026-06-30,
  enlarge x limits=false,
  legend style={at={(0.98,0.04)}, anchor=south east, legend columns=2},
]
\addplot[pubGray, thin] coordinates {(2024-08-01,0) (2026-06-30,0)};
\addplot[pubNavy, ultra thick, mark=none] coordinates {(2024-08-30,1.447583273) (2024-09-30,-9.040474753) (2024-10-31,2.628782264) (2024-11-29,0.06375246403) (2024-12-31,0.08323199687) (2025-01-27,3.099296534) (2025-02-28,-0.3451189006) (2025-03-31,-0.1872186648) (2025-04-30,2.119410876) (2025-05-30,1.458402913) (2025-06-30,-0.3954057164) (2025-07-31,-0.1950016146) (2025-08-29,0.4866387543) (2025-09-30,5.030954645) (2025-10-31,0.5816517245) (2025-11-28,0.636636359) (2025-12-31,-0.1369367959) (2026-01-30,2.430509221) (2026-03-31,-1.416433606) (2026-04-30,0.1582169603) (2026-05-29,0.5051877477) (2026-06-30,-2.097749078)};
\addlegendentry{NN-vine TD3}
\addplot[pubTeal, dashed, thick, mark=none] coordinates {(2024-08-30,1.050059107) (2024-09-30,-8.282340905) (2024-10-31,2.739634593) (2024-11-29,0.8099742042) (2024-12-31,0.2509552845) (2025-01-27,2.383115428) (2025-02-28,-0.3531530459) (2025-03-31,0.534260806) (2025-04-30,3.844921364) (2025-05-30,-0.5777758672) (2025-06-30,-0.7724152051) (2025-07-31,-0.5372314698) (2025-08-29,-3.982582465) (2025-09-30,5.535769892) (2025-10-31,1.783916425) (2025-11-28,0.5863120891) (2025-12-31,-0.4367742663) (2026-01-30,2.237355101) (2026-03-31,-0.5685456406) (2026-04-30,1.822850428) (2026-05-29,1.800690663) (2026-06-30,-3.442076455)};
\addlegendentry{DCC--GARCH}
\addplot[pubGold, dashdotted, thick, mark=none] coordinates {(2024-08-30,1.388204605) (2024-09-30,-15.05972302) (2024-10-31,0.6460831666) (2024-11-29,1.385547681) (2024-12-31,2.37004188) (2025-01-27,5.371407794) (2025-02-28,-2.461594376) (2025-03-31,-1.396295921) (2025-04-30,4.635373604) (2025-05-30,2.788527129) (2025-06-30,-1.369704615) (2025-07-31,1.497951153) (2025-08-29,2.732777405) (2025-09-30,9.618997211) (2025-10-31,-0.9823727019) (2025-11-28,2.015892735) (2025-12-31,-1.098896154) (2026-01-30,0.8398103223) (2026-03-31,-1.282770534) (2026-04-30,1.441831305) (2026-05-29,-0.1847581267) (2026-06-30,-2.027292439)};
\addlegendentry{Static vine}
\addplot[pubPurple, densely dashed, thick, mark=none] coordinates {(2024-08-30,1.017562027) (2024-09-30,-16.2719203) (2024-10-31,1.47453085) (2024-11-29,0.4585108237) (2024-12-31,2.730763476) (2025-01-27,5.516704516) (2025-02-28,-2.343208578) (2025-03-31,-0.9533431362) (2025-04-30,3.730730279) (2025-05-30,3.258764904) (2025-06-30,-2.05398001) (2025-07-31,1.658805649) (2025-08-29,5.880162842) (2025-09-30,6.923615449) (2025-10-31,-1.001232925) (2025-11-28,2.313256281) (2025-12-31,-1.439051008) (2026-01-30,-1.114528231) (2026-03-31,-1.55626984) (2026-04-30,1.081324527) (2026-05-29,0.3525817464) (2026-06-30,-0.8243262308)};
\addlegendentry{Dynamic NN-vine}
\end{axis}
\end{tikzpicture}

\caption{Monthly net returns in excess of equal weight for the primary ensemble
and selected benchmarks over the 22 complete locked periods. The horizontal
zero line denotes equal-weight performance. Evidence class: frozen confirmatory
primary evaluation.}
\label{fig:appendix_monthly_excess}
\end{figure}

The frozen primary table uses the preregistered target-to-target turnover
convention. Figure \ref{fig:appendix_drift_accounting} reports the subsequent
deterministic reconciliation with realised pre-trade drifted holdings. The
small difference for the ensemble supports the conclusion that its low
implementation burden is not an artefact of the original convention.

\begin{figure}[!tbp]
\centering
% AUTO-GENERATED. DO NOT EDIT BY HAND.
% Figure: Drift-aware turnover and cost accounting
% Evidence class: post_holdout_exploratory
% Regenerate with publication_pipeline_draft/generate_publication_tikz.py.
\begin{tikzpicture}
\begin{groupplot}[
 group style={group size=2 by 1, horizontal sep=2cm},
 publication axis, width=0.42\linewidth, height=0.25\linewidth,
]
\nextgroupplot[title={Turnover}, xlabel={Monthly turnover}, ytick={1,2}, yticklabels={Mean seed,Ensemble}, y dir=reverse, ymin=0.55, ymax=2.45, grid=none, xlabel style={yshift=-1.5mm}]
\draw[pubGray, line width=0.8pt] (axis cs:0.5747084513,1) -- (axis cs:0.584447881,1);
\addplot[only marks, mark=*, mark size=3pt, pubSlate] coordinates {(0.5747084513,1)};
\addplot[only marks, mark=diamond*, mark size=3.3pt, pubNavy] coordinates {(0.584447881,1)};
\node[font=\tiny, anchor=south, inner sep=2pt, text=pubNavy] at (axis cs:0.5795781661,1) {+$\Delta$ 0.010};
\draw[pubGray, line width=0.8pt] (axis cs:0.3169771455,2) -- (axis cs:0.3303314715,2);
\addplot[only marks, mark=*, mark size=3pt, pubSlate] coordinates {(0.3169771455,2)};
\addplot[only marks, mark=diamond*, mark size=3.3pt, pubNavy] coordinates {(0.3303314715,2)};
\node[font=\tiny, anchor=south, inner sep=2pt, text=pubNavy] at (axis cs:0.3236543085,2) {+$\Delta$ 0.013};
\nextgroupplot[title={Trading cost}, xlabel={Transaction cost (bps)}, ytick={1,2}, yticklabels={Mean seed,Ensemble}, y dir=reverse, ymin=0.55, ymax=2.45, grid=none, xlabel style={yshift=-1.5mm}]
\draw[pubGray, line width=0.8pt] (axis cs:5.747084513,1) -- (axis cs:5.84447881,1);
\addplot[only marks, mark=*, mark size=3pt, pubSlate] coordinates {(5.747084513,1)};
\addplot[only marks, mark=diamond*, mark size=3.3pt, pubNavy] coordinates {(5.84447881,1)};
\node[font=\tiny, anchor=south, inner sep=2pt, text=pubNavy] at (axis cs:5.795781661,1) {+$\Delta$ 0.097};
\draw[pubGray, line width=0.8pt] (axis cs:3.169771455,2) -- (axis cs:3.303314715,2);
\addplot[only marks, mark=*, mark size=3pt, pubSlate] coordinates {(3.169771455,2)};
\addplot[only marks, mark=diamond*, mark size=3.3pt, pubNavy] coordinates {(3.303314715,2)};
\node[font=\tiny, anchor=south, inner sep=2pt, text=pubNavy] at (axis cs:3.236543085,2) {+$\Delta$ 0.134};
\end{groupplot}
\matrix[draw=none, font=\scriptsize, column sep=3pt, row sep=2pt,
 anchor=east]
 at ($(current bounding box.east)+(-1mm,0)$) {
 \node[draw=pubSlate, fill=pubSlate, circle, inner sep=1.2pt] {}; &
 \node[anchor=west] {Target-to-target}; \\
 \node[draw=pubNavy, fill=pubNavy, diamond, inner sep=1.2pt] {}; &
 \node[anchor=west] {Drift-aware}; \\
};
\end{tikzpicture}

\caption{Target-to-target and drift-aware turnover and transaction-cost
accounting for the mean individual policy and arithmetic target-weight
ensemble. This deterministic post-holdout audit does not overwrite the frozen
primary result. Evidence class: post-holdout exploratory.}
\label{fig:appendix_drift_accounting}
\end{figure}

Figure \ref{fig:appendix_seed_correlation} examines whether the 20 trained
policies encode genuinely different allocation paths. The comparatively low
off-diagonal correlations contrast with the high common-path return
correlations and explain how target-weight averaging can cancel leverage,
short exposure, and trades.

\begin{figure}[!tbp]
\centering
% AUTO-GENERATED. DO NOT EDIT BY HAND.
% Figure: Pairwise seed weight-path correlation
% Evidence class: post_holdout_explanatory
% Regenerate with publication_pipeline_draft/generate_publication_tikz.py.
\begin{tikzpicture}
\begin{axis}[
  publication axis, width=0.62\linewidth, height=0.50\linewidth,
  title={Pairwise correlation of complete target-weight paths},
  xlabel={Training seed}, ylabel={Training seed},
  xmin=0.5, xmax=20.5, ymin=0.5, ymax=20.5,
  xtick={1,5,10,15,20}, ytick={1,5,10,15,20},
  xticklabels={20260741,20260745,20260750,20260755,20260760}, yticklabels={20260741,20260745,20260750,20260755,20260760},
  y dir=reverse,
  colormap={correlationmap}{rgb255(0cm)=(150,18,42)
    rgb255(0.5cm)=(247,247,248)
    rgb255(1cm)=(10,42,78)},
  point meta min=-1, point meta max=1,
  colorbar, colorbar style={width=1.8mm, ylabel={Correlation},
    ylabel style={font=\scriptsize}, tick label style={font=\scriptsize},
    ytick={-1,-0.5,0,0.5,1}},
]
\addplot[matrix plot*, mesh/cols=20, point meta=explicit]
 table[meta index=2] {
x y meta
1 1 1
2 1 0.03514258873
3 1 0.6779791263
4 1 0.1009734499
5 1 -0.05936015389
6 1 0.2792914979
7 1 0.640927952
8 1 0.7809711263
9 1 0.2621458859
10 1 0.2826881125
11 1 0.510146009
12 1 0.1840585416
13 1 -0.03738278411
14 1 0.2344985604
15 1 0.348889502
16 1 0.1465195878
17 1 0.6445533817
18 1 0.131487965
19 1 0.6772993068
20 1 0.670077678
1 2 0.03514258873
2 2 1
3 2 0.1487271899
4 2 -0.01164303212
5 2 -0.08098176301
6 2 -0.3787515036
7 2 -0.2241786654
8 2 -0.2120457469
9 2 -0.09930568542
10 2 -0.4829313464
11 2 0.01669604583
12 2 0.4586987003
13 2 -0.02443511985
14 2 -0.2700232864
15 2 0.1291750184
16 2 -0.121854151
17 2 -0.2814959234
18 2 0.06860611614
19 2 -0.2153078933
20 2 -0.0868604715
1 3 0.6779791263
2 3 0.1487271899
3 3 1
4 3 -0.02516430804
5 3 0.02212361937
6 3 0.1493420311
7 3 0.6970352291
8 3 0.6376217076
9 3 0.3408326309
10 3 0.1226394649
11 3 0.5344803838
12 3 0.1350135126
13 3 0.1282783468
14 3 0.2466580957
15 3 0.3978960582
16 3 0.1100809189
17 3 0.5932781629
18 3 -0.03218275136
19 3 0.5156294428
20 3 0.6045011875
1 4 0.1009734499
2 4 -0.01164303212
3 4 -0.02516430804
4 4 1
5 4 0.4868820007
6 4 0.4365246749
7 4 0.2105913113
8 4 0.1671278995
9 4 0.1803524503
10 4 0.4280857586
11 4 -0.1326646089
12 4 0.2282007919
13 4 0.4692184156
14 4 0.06236509905
15 4 0.140814013
16 4 0.3551003658
17 4 0.2542910456
18 4 0.4730761816
19 4 0.4040963165
20 4 0.3366693703
1 5 -0.05936015389
2 5 -0.08098176301
3 5 0.02212361937
4 5 0.4868820007
5 5 1
6 5 0.2476265685
7 5 0.3322912128
8 5 0.08364840175
9 5 0.6722005416
10 5 0.1768085446
11 5 -0.1507738624
12 5 0.1213759821
13 5 0.2801135456
14 5 -0.1691893371
15 5 -0.3669093877
16 5 0.6591185552
17 5 0.2023426199
18 5 0.2232944683
19 5 0.2119568884
20 5 0.4281960579
1 6 0.2792914979
2 6 -0.3787515036
3 6 0.1493420311
4 6 0.4365246749
5 6 0.2476265685
6 6 1
7 6 0.5075766346
8 6 0.4478029355
9 6 0.1008627942
10 6 0.5974298692
11 6 -0.04102692117
12 6 0.1555367542
13 6 0.395198882
14 6 0.1862241375
15 6 0.2442701178
16 6 0.2256088852
17 6 0.6290511985
18 6 0.3134346305
19 6 0.6017106909
20 6 0.4689443155
1 7 0.640927952
2 7 -0.2241786654
3 7 0.6970352291
4 7 0.2105913113
5 7 0.3322912128
6 7 0.5075766346
7 7 1
8 7 0.7914871995
9 7 0.5553593952
10 7 0.4629042763
11 7 0.4428024129
12 7 0.04403675623
13 7 0.2506034386
14 7 0.3692584725
15 7 0.2395321919
16 7 0.3800128045
17 7 0.8360614594
18 7 0.09162954661
19 7 0.7672604576
20 7 0.7115941586
1 8 0.7809711263
2 8 -0.2120457469
3 8 0.6376217076
4 8 0.1671278995
5 8 0.08364840175
6 8 0.4478029355
7 8 0.7914871995
8 8 1
9 8 0.3383192408
10 8 0.4056025215
11 8 0.4944415814
12 8 -0.002081436133
13 8 0.09111801985
14 8 0.3089453959
15 8 0.3233138099
16 8 0.2159890492
17 8 0.8471881057
18 8 0.02908906135
19 8 0.7736921237
20 8 0.6604255567
1 9 0.2621458859
2 9 -0.09930568542
3 9 0.3408326309
4 9 0.1803524503
5 9 0.6722005416
6 9 0.1008627942
7 9 0.5553593952
8 9 0.3383192408
9 9 1
10 9 0.2029607988
11 9 0.2936860522
12 9 0.01720509555
13 9 0.06813541873
14 9 0.06227559947
15 9 -0.2922056541
16 9 0.5868805251
17 9 0.3527249067
18 9 0.07812602316
19 9 0.3433763021
20 9 0.5350640444
1 10 0.2826881125
2 10 -0.4829313464
3 10 0.1226394649
4 10 0.4280857586
5 10 0.1768085446
6 10 0.5974298692
7 10 0.4629042763
8 10 0.4056025215
9 10 0.2029607988
10 10 1
11 10 0.2435741893
12 10 -0.13192106
13 10 0.388124747
14 10 0.5112290007
15 10 0.2254409807
16 10 0.2038974208
17 10 0.5782296624
18 10 0.1738769677
19 10 0.5464317341
20 10 0.3044605302
1 11 0.510146009
2 11 0.01669604583
3 11 0.5344803838
4 11 -0.1326646089
5 11 -0.1507738624
6 11 -0.04102692117
7 11 0.4428024129
8 11 0.4944415814
9 11 0.2936860522
10 11 0.2435741893
11 11 1
12 11 -0.1745313877
13 11 -0.1477403967
14 11 0.366544044
15 11 0.2920730282
16 11 -0.01197131548
17 11 0.3985993867
18 11 -0.1138323983
19 11 0.4474695057
20 11 0.2278796797
1 12 0.1840585416
2 12 0.4586987003
3 12 0.1350135126
4 12 0.2282007919
5 12 0.1213759821
6 12 0.1555367542
7 12 0.04403675623
8 12 -0.002081436133
9 12 0.01720509555
10 12 -0.13192106
11 12 -0.1745313877
12 12 1
13 12 0.09915009437
14 12 -0.200402577
15 12 0.05147206248
16 12 0.3223432528
17 12 0.04150037618
18 12 0.354869622
19 12 0.07771174503
20 12 0.2697320312
1 13 -0.03738278411
2 13 -0.02443511985
3 13 0.1282783468
4 13 0.4692184156
5 13 0.2801135456
6 13 0.395198882
7 13 0.2506034386
8 13 0.09111801985
9 13 0.06813541873
10 13 0.388124747
11 13 -0.1477403967
12 13 0.09915009437
13 13 1
14 13 0.05590582153
15 13 0.3684770155
16 13 0.04131722258
17 13 0.2816736109
18 13 0.07806015635
19 13 0.09360462405
20 13 0.1297922883
1 14 0.2344985604
2 14 -0.2700232864
3 14 0.2466580957
4 14 0.06236509905
5 14 -0.1691893371
6 14 0.1862241375
7 14 0.3692584725
8 14 0.3089453959
9 14 0.06227559947
10 14 0.5112290007
11 14 0.366544044
12 14 -0.200402577
13 14 0.05590582153
14 14 1
15 14 0.3405471318
16 14 0.1539214225
17 14 0.4190105886
18 14 0.1194684563
19 14 0.4379165563
20 14 0.05755063947
1 15 0.348889502
2 15 0.1291750184
3 15 0.3978960582
4 15 0.140814013
5 15 -0.3669093877
6 15 0.2442701178
7 15 0.2395321919
8 15 0.3233138099
9 15 -0.2922056541
10 15 0.2254409807
11 15 0.2920730282
12 15 0.05147206248
13 15 0.3684770155
14 15 0.3405471318
15 15 1
16 15 -0.2767923969
17 15 0.402357044
18 15 0.09644175481
19 15 0.377493426
20 15 -0.01087100353
1 16 0.1465195878
2 16 -0.121854151
3 16 0.1100809189
4 16 0.3551003658
5 16 0.6591185552
6 16 0.2256088852
7 16 0.3800128045
8 16 0.2159890492
9 16 0.5868805251
10 16 0.2038974208
11 16 -0.01197131548
12 16 0.3223432528
13 16 0.04131722258
14 16 0.1539214225
15 16 -0.2767923969
16 16 1
17 16 0.2950449133
18 16 0.4617771247
19 16 0.2863052795
20 16 0.5010134516
1 17 0.6445533817
2 17 -0.2814959234
3 17 0.5932781629
4 17 0.2542910456
5 17 0.2023426199
6 17 0.6290511985
7 17 0.8360614594
8 17 0.8471881057
9 17 0.3527249067
10 17 0.5782296624
11 17 0.3985993867
12 17 0.04150037618
13 17 0.2816736109
14 17 0.4190105886
15 17 0.402357044
16 17 0.2950449133
17 17 1
18 17 0.1103086858
19 17 0.7895025014
20 17 0.6635796592
1 18 0.131487965
2 18 0.06860611614
3 18 -0.03218275136
4 18 0.4730761816
5 18 0.2232944683
6 18 0.3134346305
7 18 0.09162954661
8 18 0.02908906135
9 18 0.07812602316
10 18 0.1738769677
11 18 -0.1138323983
12 18 0.354869622
13 18 0.07806015635
14 18 0.1194684563
15 18 0.09644175481
16 18 0.4617771247
17 18 0.1103086858
18 18 1
19 18 0.2275394684
20 18 0.3497873448
1 19 0.6772993068
2 19 -0.2153078933
3 19 0.5156294428
4 19 0.4040963165
5 19 0.2119568884
6 19 0.6017106909
7 19 0.7672604576
8 19 0.7736921237
9 19 0.3433763021
10 19 0.5464317341
11 19 0.4474695057
12 19 0.07771174503
13 19 0.09360462405
14 19 0.4379165563
15 19 0.377493426
16 19 0.2863052795
17 19 0.7895025014
18 19 0.2275394684
19 19 1
20 19 0.614564792
1 20 0.670077678
2 20 -0.0868604715
3 20 0.6045011875
4 20 0.3366693703
5 20 0.4281960579
6 20 0.4689443155
7 20 0.7115941586
8 20 0.6604255567
9 20 0.5350640444
10 20 0.3044605302
11 20 0.2278796797
12 20 0.2697320312
13 20 0.1297922883
14 20 0.05755063947
15 20 -0.01087100353
16 20 0.5010134516
17 20 0.6635796592
18 20 0.3497873448
19 20 0.614564792
20 20 1
};
\end{axis}
\end{tikzpicture}

\caption{Pairwise correlations of the flattened complete target-weight paths
for the 20 frozen primary seeds. The diagonal is one by construction; the
off-diagonal structure measures optimisation-induced allocation diversity.
Evidence class: post-holdout explanatory.}
\label{fig:appendix_seed_correlation}
\end{figure}

\clearpage
\subsection*{B. Synthetic-Data Fidelity Diagnostics}

Figures \ref{fig:appendix_synthetic_marginals}--\ref{fig:appendix_synthetic_temporal}
compare the synthetic pre-training sample with the historical training prefix.
They are diagnostics of the generator's training distribution, not tests on
the reserved evaluation period. The identity lines indicate exact agreement;
deviations should be interpreted jointly with the registered compatibility
gates and the short historical sample.

\begin{figure}[!tbp]
\centering
% AUTO-GENERATED. DO NOT EDIT BY HAND.
% Figure: Synthetic marginal fidelity
% Evidence class: training_prefix_synthetic_diagnostic
% Regenerate with publication_pipeline_draft/generate_publication_tikz.py.
\begin{tikzpicture}
\begin{groupplot}[
  group style={group size=2 by 2, horizontal sep=1.75cm, vertical sep=2cm},
  publication axis, width=0.48\linewidth, height=0.3\linewidth,
  tick label style={font=\tiny, text=pubInk},
]
\nextgroupplot[title={Mean}, xmin=-0.001228863755, xmax=0.01340333895, ymin=-0.001228863755, ymax=0.01340333895, xlabel={Historical}, ylabel={Synthetic}]
\addplot[pubSlate, dashed, thin] coordinates {(-0.001228863755,-0.001228863755) (0.01340333895,0.01340333895)};
\addplot[publication point, pubNavy, fill=pubNavy] coordinates {(0.000187155861233758,0.0005885261008) (0.00194818722101499,0.002022409578) (0.00760404225359474,0.007764369593) (0.00662345620181135,0.00676275009) (0.0117033053729972,0.01198731933) (0.00893118742245167,0.009105712758) (0.000348743759706622,0.0009682863426)};
\node[font=\tiny, anchor=south west, inner sep=2.5pt, text=pubSlate] at (axis cs:0.0001871558612,0.0005885261008) {CHINEXT};
\node[font=\tiny, anchor=south west, inner sep=2.5pt, text=pubSlate] at (axis cs:0.001948187221,0.002022409578) {DIVIDEND};
\node[font=\tiny, anchor=south west, inner sep=2.5pt, text=pubSlate] at (axis cs:0.007604042254,0.007764369593) {DOW};
\node[font=\tiny, anchor=south west, inner sep=2.5pt, text=pubSlate] at (axis cs:0.006623456202,0.00676275009) {GOLD};
\node[font=\tiny, anchor=south west, inner sep=2.5pt, text=pubSlate] at (axis cs:0.01170330537,0.01198731933) {NASDAQ};
\node[font=\tiny, anchor=south west, inner sep=2.5pt, text=pubSlate] at (axis cs:0.008931187422,0.009105712758) {SP500};
\node[font=\tiny, anchor=south west, inner sep=2.5pt, text=pubSlate] at (axis cs:0.0003487437597,0.0009682863426) {SSE50};
\nextgroupplot[title={Standard deviation}, xmin=0.02679572602, xmax=0.0960722822, ymin=0.02679572602, ymax=0.0960722822, xlabel={Historical}, ylabel={Synthetic}]
\addplot[pubSlate, dashed, thin] coordinates {(0.02679572602,0.02679572602) (0.0960722822,0.0960722822)};
\addplot[publication point, pubNavy, fill=pubNavy] coordinates {(0.0893680993446617,0.08892141528) (0.0563590770449414,0.05724981138) (0.0456597504419373,0.04554729386) (0.0337980966455524,0.03349990887) (0.0546257373369223,0.05424454687) (0.0460427439293137,0.04581990519) (0.0574011659548262,0.05737320839)};
\node[font=\tiny, anchor=south west, inner sep=2.5pt, text=pubSlate] at (axis cs:0.08936809934,0.08892141528) {CHINEXT};
\node[font=\tiny, anchor=south west, inner sep=2.5pt, text=pubSlate] at (axis cs:0.05635907704,0.05724981138) {DIVIDEND};
\node[font=\tiny, anchor=south west, inner sep=2.5pt, text=pubSlate] at (axis cs:0.04565975044,0.04554729386) {DOW};
\node[font=\tiny, anchor=south west, inner sep=2.5pt, text=pubSlate] at (axis cs:0.03379809665,0.03349990887) {GOLD};
\node[font=\tiny, anchor=south west, inner sep=2.5pt, text=pubSlate] at (axis cs:0.05462573734,0.05424454687) {NASDAQ};
\node[font=\tiny, anchor=south west, inner sep=2.5pt, text=pubSlate] at (axis cs:0.04604274393,0.04581990519) {SP500};
\node[font=\tiny, anchor=south west, inner sep=2.5pt, text=pubSlate] at (axis cs:0.05740116595,0.05737320839) {SSE50};
\nextgroupplot[title={5\% quantile}, xmin=-0.1470693087, xmax=-0.03023422739, ymin=-0.1470693087, ymax=-0.03023422739, xlabel={Historical}, ylabel={Synthetic}]
\addplot[pubSlate, dashed, thin] coordinates {(-0.1470693087,-0.1470693087) (-0.03023422739,-0.03023422739)};
\addplot[publication point, pubNavy, fill=pubNavy] coordinates {(-0.13406598621456,-0.135762688) (-0.0774944255504085,-0.08274152743) (-0.0683751853692367,-0.0689004967) (-0.0415408481631415,-0.0422416269) (-0.094245996267562,-0.09494924755) (-0.0779512913476148,-0.08420411925) (-0.0812259283694992,-0.08155935532)};
\node[font=\tiny, anchor=south west, inner sep=2.5pt, text=pubSlate] at (axis cs:-0.1340659862,-0.135762688) {CHINEXT};
\node[font=\tiny, anchor=south west, inner sep=2.5pt, text=pubSlate] at (axis cs:-0.07749442555,-0.08274152743) {DIVIDEND};
\node[font=\tiny, anchor=south west, inner sep=2.5pt, text=pubSlate] at (axis cs:-0.06837518537,-0.0689004967) {DOW};
\node[font=\tiny, anchor=south west, inner sep=2.5pt, text=pubSlate] at (axis cs:-0.04154084816,-0.0422416269) {GOLD};
\node[font=\tiny, anchor=south west, inner sep=2.5pt, text=pubSlate] at (axis cs:-0.09424599627,-0.09494924755) {NASDAQ};
\node[font=\tiny, anchor=south west, inner sep=2.5pt, text=pubSlate] at (axis cs:-0.07795129135,-0.08420411925) {SP500};
\node[font=\tiny, anchor=south west, inner sep=2.5pt, text=pubSlate] at (axis cs:-0.08122592837,-0.08155935532) {SSE50};
\nextgroupplot[title={5\% CVaR}, xmin=-0.2275545267, xmax=-0.04396851429, ymin=-0.2275545267, ymax=-0.04396851429, xlabel={Historical}, ylabel={Synthetic}]
\addplot[pubSlate, dashed, thin] coordinates {(-0.2275545267,-0.2275545267) (-0.04396851429,-0.04396851429)};
\addplot[publication point, pubNavy, fill=pubNavy] coordinates {(-0.206099919439797,-0.2097881384) (-0.134576097508507,-0.1412769702) (-0.102535989715556,-0.1042841761) (-0.0618108269798256,-0.06173490259) (-0.115233663171402,-0.1152629598) (-0.105759258482592,-0.1064234603) (-0.128563337118877,-0.1310019934)};
\node[font=\tiny, anchor=south west, inner sep=2.5pt, text=pubSlate] at (axis cs:-0.2060999194,-0.2097881384) {CHINEXT};
\node[font=\tiny, anchor=south west, inner sep=2.5pt, text=pubSlate] at (axis cs:-0.1345760975,-0.1412769702) {DIVIDEND};
\node[font=\tiny, anchor=south west, inner sep=2.5pt, text=pubSlate] at (axis cs:-0.1025359897,-0.1042841761) {DOW};
\node[font=\tiny, anchor=south west, inner sep=2.5pt, text=pubSlate] at (axis cs:-0.06181082698,-0.06173490259) {GOLD};
\node[font=\tiny, anchor=south west, inner sep=2.5pt, text=pubSlate] at (axis cs:-0.1152336632,-0.1152629598) {NASDAQ};
\node[font=\tiny, anchor=south west, inner sep=2.5pt, text=pubSlate] at (axis cs:-0.1057592585,-0.1064234603) {SP500};
\node[font=\tiny, anchor=south west, inner sep=2.5pt, text=pubSlate] at (axis cs:-0.1285633371,-0.1310019934) {SSE50};
\end{groupplot}
\end{tikzpicture}

\caption{Marginal fidelity of synthetic pre-training returns. Each panel
compares an asset-level historical training-prefix statistic with its synthetic
counterpart; the gray diagonal denotes exact agreement. Evidence class:
training-prefix synthetic diagnostic.}
\label{fig:appendix_synthetic_marginals}
\end{figure}

\begin{figure}[!tbp]
\centering
% AUTO-GENERATED. DO NOT EDIT BY HAND.
% Figure: Synthetic dependence fidelity
% Evidence class: training_prefix_synthetic_diagnostic
% Regenerate with publication_pipeline_draft/generate_publication_tikz.py.
\begin{tikzpicture}
\begin{groupplot}[
  group style={group size=2 by 1, horizontal sep=1.75cm}, publication axis,
  width=0.45\linewidth, height=0.34\linewidth,
  tick label style={font=\tiny, text=pubInk},
]
\nextgroupplot[title={Pairwise correlation}, xmin=-0.2860961691, xmax=1.061446298, ymin=-0.2860961691, ymax=1.061446298, xlabel={Historical}, ylabel={Synthetic}]
\addplot[pubSlate, dashed, thin] coordinates {(-0.2860961691,-0.2860961691) (1.061446298,1.061446298)};
\addplot[publication point, pubTeal, fill=pubTeal] coordinates {(0.944630895942701,0.9273822349) (0.949151092190377,0.9437377154) (0.353054445799659,0.3481656788) (0.324686883773623,0.3485842799) (0.299694338201908,0.2573860473) (-0.0816607177252361,-0.09073720555) (0.826845655220616,0.8586617192) (0.363858348135544,0.3764146672) (0.316465554012541,0.3450529802) (0.348157369907034,0.3039032024) (-0.10800686102042,-0.09662994988) (0.350534179460681,0.3823891082) (0.320421710296357,0.368927776) (0.256370224571845,0.2595668966) (-0.10307552335701,-0.1011686246) (0.744966014969528,0.7876120137) (0.650772359726252,0.5737343236) (-0.0122899853781621,-0.08706440116) (-0.136259726650586,-0.1738009635) (-0.023830597600314,0.01459325385)};
\addplot[publication point, pubRose, fill=pubRose, mark=triangle*] coordinates {(0.514562494526225,0.4130728301)};
\nextgroupplot[title={5\% lower-tail co-exceedance}, xmin=-0.08333333333, xmax=0.9166666667, ymin=-0.08333333333, ymax=0.9166666667, xlabel={Historical}, ylabel={Synthetic}]
\addplot[pubSlate, dashed, thin] coordinates {(-0.08333333333,-0.08333333333) (0.9166666667,0.9166666667)};
\addplot[publication point, pubTeal, fill=pubTeal] coordinates {(0.666666666666667,0.6941666667) (0.833333333333333,0.7558333333) (0.0,0.1683333333) (0.0,0.1783333333) (0.166666666666667,0.1558333333) (0.0,0.04833333333) (0.5,0.6075) (0.0,0.1816666667) (0.0,0.1816666667) (0.166666666666667,0.1741666667) (0.0,0.04583333333) (0.0,0.1916666667) (0.0,0.1958333333) (0.0,0.1533333333) (0.0,0.04583333333) (0.5,0.605) (0.5,0.3883333333) (0.166666666666667,0.05916666667) (0.5,0.2783333333) (0.166666666666667,0.03666666667) (0.166666666666667,0.07083333333)};
\addplot[publication point, pubRose, fill=pubRose, mark=triangle*] coordinates {};
\end{groupplot}
\end{tikzpicture}

\caption{Cross-sectional dependence fidelity of the synthetic pre-training
sample. Circles satisfy the registered statistical compatibility gate and
triangles identify pairs outside it. Historical lower-tail rates are supported
by few events and are therefore interpreted with exact finite-sample
uncertainty. Evidence class: training-prefix synthetic diagnostic.}
\label{fig:appendix_synthetic_dependence}
\end{figure}

\begin{figure}[!tbp]
\centering
% AUTO-GENERATED. DO NOT EDIT BY HAND.
% Figure: Synthetic temporal fidelity
% Evidence class: training_prefix_synthetic_diagnostic
% Regenerate with publication_pipeline_draft/generate_publication_tikz.py.
\begin{tikzpicture}
\begin{groupplot}[
 group style={group size=2 by 1, horizontal sep=1.75cm}, publication axis,
 width=0.45\linewidth, height=0.34\linewidth,
 tick label style={font=\tiny, text=pubInk},
]
\nextgroupplot[title={Lag-1 return dependence}, xmin=-0.1929093845, xmax=0.2334903989, ymin=-0.1929093845, ymax=0.2334903989, xlabel={Historical}, ylabel={Synthetic}]
\addplot[pubSlate, dashed, thin] coordinates {(-0.1929093845,-0.1929093845) (0.2334903989,0.2334903989)};
\addplot[publication point, pubPurple, fill=pubPurple] coordinates {(-0.140859575710659,-0.1516448893) (-0.0987760333996144,-0.1096018115) (-0.10495747845599,-0.1454148199) (0.0355214893075414,-0.01358875363) (0.073133561766783,-0.09316314452) (0.192225903775257,0.1504411644) (-0.0688930637166957,-0.009379307192)};
\node[font=\tiny, anchor=south west, inner sep=2.5pt, text=pubSlate] at (axis cs:-0.1408595757,-0.1516448893) {SP500};
\node[font=\tiny, anchor=south west, inner sep=2.5pt, text=pubSlate] at (axis cs:-0.0987760334,-0.1096018115) {NASDAQ};
\node[font=\tiny, anchor=south west, inner sep=2.5pt, text=pubSlate] at (axis cs:-0.1049574785,-0.1454148199) {DOW};
\node[font=\tiny, anchor=south west, inner sep=2.5pt, text=pubSlate] at (axis cs:0.03552148931,-0.01358875363) {SSE50};
\node[font=\tiny, anchor=south west, inner sep=2.5pt, text=pubSlate] at (axis cs:0.07313356177,-0.09316314452) {DIVIDEND};
\node[font=\tiny, anchor=south west, inner sep=2.5pt, text=pubSlate] at (axis cs:0.1922259038,0.1504411644) {CHINEXT};
\node[font=\tiny, anchor=south west, inner sep=2.5pt, text=pubSlate] at (axis cs:-0.06889306372,-0.009379307192) {GOLD};
\nextgroupplot[title={Lag-1 squared-return dependence}, xmin=0.000199464589, xmax=0.47554341, ymin=0.000199464589, ymax=0.47554341, xlabel={Historical}, ylabel={Synthetic}]
\addplot[pubSlate, dashed, thin] coordinates {(0.000199464589,0.000199464589) (0.47554341,0.47554341)};
\addplot[publication point, pubPurple, fill=pubPurple] coordinates {(0.383062365818521,0.3925025308) (0.0801719408253526,0.1670885382) (0.429542382979633,0.2862802852) (0.100600156994286,0.1327820622) (0.0874798179267253,0.1773355398) (0.198197964031468,0.2267387025) (0.0462004915594057,0.1386122327)};
\node[font=\tiny, anchor=south west, inner sep=2.5pt, text=pubSlate] at (axis cs:0.3830623658,0.3925025308) {SP500};
\node[font=\tiny, anchor=south west, inner sep=2.5pt, text=pubSlate] at (axis cs:0.08017194083,0.1670885382) {NASDAQ};
\node[font=\tiny, anchor=south west, inner sep=2.5pt, text=pubSlate] at (axis cs:0.429542383,0.2862802852) {DOW};
\node[font=\tiny, anchor=south west, inner sep=2.5pt, text=pubSlate] at (axis cs:0.100600157,0.1327820622) {SSE50};
\node[font=\tiny, anchor=south west, inner sep=2.5pt, text=pubSlate] at (axis cs:0.08747981793,0.1773355398) {DIVIDEND};
\node[font=\tiny, anchor=south west, inner sep=2.5pt, text=pubSlate] at (axis cs:0.198197964,0.2267387025) {CHINEXT};
\node[font=\tiny, anchor=south west, inner sep=2.5pt, text=pubSlate] at (axis cs:0.04620049156,0.1386122327) {GOLD};
\end{groupplot}
\end{tikzpicture}

\caption{Temporal fidelity for lag-one returns and lag-one squared returns.
Asset labels identify the sources of departures from the identity line.
Evidence class: training-prefix synthetic diagnostic.}
\label{fig:appendix_synthetic_temporal}
\end{figure}

\clearpage
\subsection*{C. Training and Optimisation Diagnostics}

The retained diagnostic logs provide complementary behavioural and numerical
views of learning. Figure \ref{fig:appendix_pretraining_stability} tracks
trailing policy outcomes through pre-training, while Figure
\ref{fig:appendix_optimizer} tracks actor and critic optimisation quantities.
These figures display all 20 frozen optimisation seeds supplied by the
immutable aggregate archive; they do not turn episodes or gradient updates
into independent statistical observations.

\begin{figure}[!tbp]
\centering
% AUTO-GENERATED. DO NOT EDIT BY HAND.
% Figure: Pre-training stability across 20 seeds
% Evidence class: training_prefix_diagnostic
% Regenerate with publication_pipeline_draft/generate_publication_tikz.py.
\begin{tikzpicture}
\begin{groupplot}[
 group style={group size=2 by 2, horizontal sep=1.35cm, vertical sep=2.25cm},
 publication axis, width=0.49\linewidth, height=0.3\linewidth,
 title style={font=\footnotesize\bfseries, text=pubInk, yshift=1.2mm},
]
\nextgroupplot[title={Rolling reward}, xlabel={Pre-training episode}]
\addplot[pubNavy, opacity=0.20, line width=0.35pt] coordinates {(10,-0.01069687658) (14,-0.01396718539) (18,-0.01672859141) (22,-0.01885469712) (26,-0.00403631802) (30,-0.0009516549679) (34,0.01438650499) (38,0.003479770923) (42,-0.001681751945) (46,0.005642343349) (50,0.01559098048) (54,0.0197438646) (58,0.02251468446) (62,0.0113440883) (66,0.02938443927) (70,0.03467506627) (74,0.04907287643) (78,0.05135555488) (82,0.05152708818) (86,0.05915138427) (90,0.07724455548) (94,0.07499500757) (98,0.08709604556) (102,0.08681717011) (106,0.1057935812) (110,0.1051625826) (114,0.1132544886) (118,0.1172275662) (122,0.1249531932) (126,0.1215526155) (130,0.1304436679) (134,0.1170565073) (138,0.1203937015) (142,0.1242937921) (146,0.1230806464) (150,0.1101164887) (154,0.1121163648) (158,0.09440581474) (162,0.09499361709) (166,0.09506638997) (170,0.08282244522) (174,0.08227344644) (178,0.07888248464) (182,0.08173458664) (186,0.08600762618) (190,0.08063425947) (194,0.08453554077) (198,0.09108552902) (202,0.08954921692) (206,0.0946831601) (210,0.0970670791) (214,0.08423273458) (218,0.09449715767) (222,0.09140269347) (226,0.0944586599) (230,0.08946568691) (234,0.08457769093) (238,0.08547670932) (242,0.08212797517) (246,0.08830044804) (250,0.08495999158) (254,0.09909284027) (258,0.09549297702) (262,0.09384925332) (266,0.08705741646) (270,0.09261962908) (274,0.09684801033) (278,0.09788458432) (282,0.1058157321) (286,0.1138673645) (290,0.1265998757) (294,0.1140808254) (298,0.1174004656) (302,0.1221040582) (306,0.1250820313) (310,0.1216240582) (314,0.1279278471) (318,0.1228994067) (322,0.1200491124) (326,0.118981953) (330,0.1154346749) (334,0.1113710257) (338,0.1024667315) (342,0.1137097018) (346,0.1043455572) (350,0.1035473389) (354,0.09308025385) (358,0.1013363351) (362,0.1100808963) (366,0.1098905423) (370,0.1133589622) (374,0.1182333066) (378,0.1224182766) (382,0.117371561) (386,0.12433748) (390,0.1216975705) (394,0.1134955504) (398,0.1211564693) (402,0.1111634405) (406,0.09937166657) (410,0.1097587662) (414,0.1156686152) (418,0.1186205602) (422,0.09890937567) (426,0.1035695395) (430,0.08706655284) (434,0.09205017919) (438,0.09386445852) (442,0.08837697901) (446,0.09968222603) (450,0.1011859173) (454,0.09081910917) (458,0.0973532546) (462,0.09940544222) (466,0.09250689572) (470,0.08823414883) (474,0.1035123964) (478,0.09725234186) (482,0.1064915912) (486,0.09725699714) (490,0.1053031839) (494,0.1023864551) (498,0.09664884263) (502,0.1070623703) (506,0.1201674701) (510,0.1148371747) (514,0.1026925213) (518,0.09706512179) (522,0.1050299451) (526,0.1078524601) (530,0.108519668) (534,0.112838575) (538,0.1288291701) (542,0.1146376574) (546,0.1107797854) (550,0.1123907023) (554,0.1196915726) (558,0.1111737464) (562,0.1081598198) (566,0.1241764759) (570,0.1191880443) (574,0.1109843357) (578,0.1084446504) (582,0.1078631035) (586,0.096922482) (590,0.08741065709) (594,0.08842672923) (598,0.1103225383) (602,0.111326211) (606,0.1022580699) (610,0.1196096602) (614,0.1180089974) (618,0.1181250644) (622,0.1140707449) (626,0.1245472258) (630,0.1241114767) (634,0.1248231926) (638,0.1322918484) (642,0.1451873812) (646,0.1351311995) (650,0.1272739697) (654,0.1294404262) (658,0.1285033602) (662,0.1416486866) (666,0.1292430972) (670,0.1436515729) (674,0.1552048605) (678,0.1395137624) (682,0.1484559041) (686,0.1505025186) (690,0.1475157732) (694,0.1504478693) (698,0.1587436012) (702,0.1587338352) (706,0.1544472577) (710,0.1520024274) (714,0.1364269899) (718,0.1424460199) (722,0.1278554415) (726,0.1315112055) (730,0.121152564) (734,0.1165312346) (738,0.1111516477) (742,0.1158770926) (746,0.0941040566) (750,0.09047834278) (754,0.08224560872) (758,0.08535733374) (762,0.0819820145) (766,0.09534827043) (770,0.09401995369) (774,0.09313243923) (778,0.0799872319) (782,0.09207701435) (786,0.1027839889) (790,0.1014256905) (794,0.1082142558) (798,0.1038822614) (802,0.110916015) (806,0.1109269504) (810,0.1076140929) (814,0.1124712064) (818,0.114970513) (822,0.1143802113) (826,0.1231397792) (830,0.129117564) (834,0.133767567) (838,0.1307258428) (842,0.1292510862) (846,0.1269603009) (850,0.1315932843) (854,0.1412006393) (858,0.1309716464) (862,0.1438421391) (866,0.1369739155) (870,0.136166576) (874,0.1340484157) (878,0.1243349424) (882,0.1048572718) (886,0.09121418134) (890,0.1015888689) (894,0.104824092) (898,0.1074845036) (902,0.1139501333) (906,0.1066755967) (910,0.1047965692) (914,0.1069849564) (918,0.1074042414) (922,0.1007907339) (926,0.1015523987) (930,0.1031921084) (934,0.1062658146) (938,0.09496151258) (942,0.1084407625) (946,0.09997497187) (950,0.09298534434) (954,0.09450886398) (958,0.09575496547) (962,0.08232665183) (966,0.08256747654) (970,0.07962214462) (974,0.0973694809) (978,0.1128867955) (982,0.1131094638) (986,0.1082642258) (990,0.1071059924) (994,0.09989321695) (998,0.1150950959)};
\addplot[pubNavy, opacity=0.20, line width=0.35pt] coordinates {(10,-0.01727179216) (14,-0.0265897913) (18,-0.02299266073) (22,-0.02983951978) (26,-0.01929154646) (30,-0.03409180611) (34,-0.02822168079) (38,-0.02803453919) (42,-0.03120506792) (46,-0.03892580137) (50,-0.0256190156) (54,-0.03484812673) (58,-0.03549827674) (62,-0.05408906678) (66,-0.03957733575) (70,-0.04633143791) (74,-0.03031361639) (78,-0.02550792921) (82,-0.01809729138) (86,-0.008406612048) (90,0.006243744826) (94,0.007907321291) (98,0.02965477646) (102,0.01123195914) (106,0.03504017615) (110,0.03932942491) (114,0.04653752306) (118,0.05421873378) (122,0.05836159051) (126,0.05406121862) (130,0.06882756828) (134,0.04679461306) (138,0.04685623699) (142,0.04207546551) (146,0.04631599383) (150,0.03662890586) (154,0.05571362655) (158,0.03850882505) (162,0.03531132488) (166,0.03959771352) (170,0.02374130208) (174,0.02685940372) (178,0.01450865266) (182,0.0279231534) (186,0.03173477012) (190,0.03097492422) (194,0.03801061987) (198,0.04536989945) (202,0.04830685312) (206,0.04284450923) (210,0.05116696309) (214,0.03948624981) (218,0.04738722766) (222,0.04937207609) (226,0.05941250712) (230,0.05437945539) (234,0.0509914039) (238,0.04973595298) (242,0.03936543369) (246,0.04095940273) (250,0.04311095092) (254,0.05695361977) (258,0.05774808659) (262,0.05697471319) (266,0.04750388216) (270,0.04614628409) (274,0.04985630585) (278,0.0463414995) (282,0.05859873157) (286,0.05261918807) (290,0.0626917261) (294,0.04766018442) (298,0.04283278952) (302,0.03940098841) (306,0.04040915502) (310,0.0368685969) (314,0.04019597457) (318,0.03270728413) (322,0.03231562554) (326,0.02572887667) (330,0.0277541193) (334,0.03042275133) (338,0.02642633758) (342,0.03955292392) (346,0.02636813595) (350,0.02734588671) (354,0.009835939777) (358,0.01832295225) (362,0.02903937291) (366,0.02873343133) (370,0.02375545789) (374,0.03130555512) (378,0.03294979371) (382,0.02927885479) (386,0.02701607072) (390,0.02587707655) (394,0.02527854511) (398,0.03026396104) (402,0.01509511301) (406,0.007712889745) (410,0.008577219826) (414,0.01957986653) (418,0.02646650577) (422,0.01668929824) (426,0.02468408692) (430,0.01655294132) (434,0.02060115829) (438,0.02497873859) (442,0.02463970799) (446,0.04160747312) (450,0.04521282612) (454,0.04171402861) (458,0.05618622195) (462,0.06652374743) (466,0.05640349724) (470,0.04706109081) (474,0.05579931132) (478,0.04688537099) (482,0.04946970331) (486,0.03948953982) (490,0.03988772283) (494,0.03497049368) (498,0.03182192568) (502,0.03637808595) (506,0.03742279421) (510,0.03721860381) (514,0.02252426095) (518,0.02464593396) (522,0.03880666763) (526,0.03196850379) (530,0.02999816054) (534,0.0393337508) (538,0.05489227531) (542,0.03545705742) (546,0.02701348819) (550,0.02593932153) (554,0.04066218388) (558,0.03471293939) (562,0.03823407582) (566,0.05398471091) (570,0.0412654509) (574,0.03913590629) (578,0.03805736764) (582,0.04205823573) (586,0.04236325841) (590,0.04516116373) (594,0.05634053803) (598,0.07456593325) (602,0.07116850033) (606,0.06701447045) (610,0.08024210655) (614,0.06823172429) (618,0.06894727973) (622,0.0728358446) (626,0.07890249756) (630,0.07585374581) (634,0.07233138984) (638,0.06564386854) (642,0.05833145063) (646,0.04906450576) (650,0.04471020014) (654,0.0491150195) (658,0.05338436676) (662,0.06429500904) (666,0.05435892533) (670,0.06766118923) (674,0.07479860477) (678,0.06585028087) (682,0.07067873875) (686,0.07069662771) (690,0.07395805721) (694,0.08192398527) (698,0.08056907661) (702,0.08008478399) (706,0.06535739028) (710,0.06480174367) (714,0.04970544581) (718,0.06114174119) (722,0.0518846936) (726,0.05206615124) (730,0.04936516806) (734,0.05092824099) (738,0.04318836393) (742,0.03770023913) (746,0.03383894762) (750,0.03228984808) (754,0.02464803102) (758,0.02956509421) (762,0.02665308883) (766,0.03902960444) (770,0.0342046149) (774,0.03185147061) (778,0.02723313688) (782,0.031671968) (786,0.04507000242) (790,0.04398637738) (794,0.044535855) (798,0.03664497937) (802,0.05846054227) (806,0.05647686948) (810,0.0606327112) (814,0.06328186677) (818,0.0752558127) (822,0.07693049145) (826,0.0810525425) (830,0.09147996505) (834,0.1022769878) (838,0.1010414911) (842,0.1030843003) (846,0.09613347766) (850,0.09484458269) (854,0.09059568752) (858,0.08759352899) (862,0.09747763252) (866,0.07894693399) (870,0.07295781078) (874,0.07210356164) (878,0.06565750293) (882,0.05563541623) (886,0.03143721445) (890,0.04172581608) (894,0.04893517053) (898,0.058782252) (902,0.06247144251) (906,0.05906360514) (910,0.05015378556) (914,0.05629229436) (918,0.06338186104) (922,0.05890639167) (926,0.05568990472) (930,0.05589368454) (934,0.05325317578) (938,0.0321381072) (942,0.0375977927) (946,0.02825729306) (950,0.03031115961) (954,0.03835816336) (958,0.0438113718) (962,0.04263529136) (966,0.03411899515) (970,0.03213102331) (974,0.04947963197) (978,0.0689759075) (982,0.06429113901) (986,0.08275624862) (990,0.08511637629) (994,0.07324741875) (998,0.07888417076)};
\addplot[pubNavy, opacity=0.20, line width=0.35pt] coordinates {(10,-0.0005041767083) (14,-0.03351304622) (18,-0.03959563888) (22,-0.0350148568) (26,-0.0272529233) (30,-0.03067002215) (34,-0.02105686834) (38,-0.03242132962) (42,-0.03979735007) (46,-0.02852283934) (50,-0.01697375974) (54,-0.01129148803) (58,-0.006815014423) (62,-0.01233407517) (66,0.01040109115) (70,0.01704559884) (74,0.02668480841) (78,0.02915627482) (82,0.03751623788) (86,0.04270023215) (90,0.06197589329) (94,0.06003002882) (98,0.06411871911) (102,0.06593955039) (106,0.0815057516) (110,0.07662443426) (114,0.08005034264) (118,0.08504284714) (122,0.09345633845) (126,0.09428728215) (130,0.1016387806) (134,0.08980515845) (138,0.09680442458) (142,0.1083308687) (146,0.1088271647) (150,0.09389902834) (154,0.09266872701) (158,0.07666822821) (162,0.07391854136) (166,0.07161917313) (170,0.06277011524) (174,0.06514038208) (178,0.06278489763) (182,0.06599438433) (186,0.06609601295) (190,0.06278305029) (194,0.0639809152) (198,0.07054393453) (202,0.08101268739) (206,0.08442743019) (210,0.08411134185) (214,0.07742447066) (218,0.08308836721) (222,0.07737986896) (226,0.08019135098) (230,0.06935785634) (234,0.07295168266) (238,0.06782567111) (242,0.07045526597) (246,0.07462495849) (250,0.06846800864) (254,0.07593158138) (258,0.08040386585) (262,0.08100853133) (266,0.06738503658) (270,0.07028912062) (274,0.07339577852) (278,0.07026007239) (282,0.07582123323) (286,0.08657165847) (290,0.09306511491) (294,0.07436434899) (298,0.07364427531) (302,0.07844643627) (306,0.07634326133) (310,0.06902031014) (314,0.06890045218) (318,0.05445101191) (322,0.05614652556) (326,0.05453831825) (330,0.05723712144) (334,0.05580391287) (338,0.0492290679) (342,0.06622665009) (346,0.05331579877) (350,0.0499394328) (354,0.02672874028) (358,0.01710759638) (362,0.0299527846) (366,0.03841330732) (370,0.04629701067) (374,0.05369175619) (378,0.0518484345) (382,0.04438409133) (386,0.05130312575) (390,0.0419439806) (394,0.03169272042) (398,0.04427252652) (402,0.0328301427) (406,0.03338133595) (410,0.05709435625) (414,0.06906015032) (418,0.07351771028) (422,0.05778963446) (426,0.06600492287) (430,0.05682576202) (434,0.05815962655) (438,0.05900315533) (442,0.05000248398) (446,0.06886494314) (450,0.07671295872) (454,0.07733173859) (458,0.08861341601) (462,0.09212833214) (466,0.0786885559) (470,0.07512213707) (474,0.08463462977) (478,0.07910236011) (482,0.09429193607) (486,0.08999118734) (490,0.09699967747) (494,0.09778941702) (498,0.09479535134) (502,0.1052773071) (506,0.1115043341) (510,0.1138826778) (514,0.1044422244) (518,0.1089469522) (522,0.1242519036) (526,0.124305628) (530,0.1151646763) (534,0.1205066323) (538,0.1288370889) (542,0.1245158063) (546,0.1257864705) (550,0.1318105128) (554,0.1337504882) (558,0.1360356168) (562,0.1303147241) (566,0.131912566) (570,0.132450489) (574,0.1271352813) (578,0.1296123188) (582,0.1275826871) (586,0.1239696305) (590,0.1057316748) (594,0.1069221933) (598,0.1117098758) (602,0.09701151088) (606,0.08888252903) (610,0.09591273396) (614,0.1001133024) (618,0.1026451785) (622,0.09152902495) (626,0.0984987124) (630,0.1027697051) (634,0.1051398124) (638,0.1044403368) (642,0.1150263938) (646,0.1124887933) (650,0.1203944215) (654,0.1217023142) (658,0.1229547894) (662,0.1182645165) (666,0.1175469483) (670,0.1163358999) (674,0.1179293708) (678,0.1165444095) (682,0.1233625572) (686,0.1286350238) (690,0.1278393705) (694,0.1158648811) (698,0.1114745121) (702,0.1160460641) (706,0.1202935729) (710,0.1324319255) (714,0.1199524315) (718,0.1223371089) (722,0.1221130121) (726,0.1193092334) (730,0.1130716199) (734,0.1068765202) (738,0.1081875449) (742,0.1187416518) (746,0.1102250621) (750,0.1084718107) (754,0.09914141544) (758,0.09843074561) (762,0.1006436259) (766,0.1083686402) (770,0.1060950454) (774,0.1063545545) (778,0.1083241171) (782,0.1088297193) (786,0.1120023784) (790,0.1071334231) (794,0.1114582144) (798,0.1127073556) (802,0.1224706747) (806,0.1218504236) (810,0.1247553672) (814,0.1270832675) (818,0.1343081919) (822,0.1315529454) (826,0.1295690083) (830,0.1380252575) (834,0.1362534116) (838,0.1304660669) (842,0.1306714189) (846,0.1273871301) (850,0.1228253426) (854,0.123302425) (858,0.1226503045) (862,0.1225176961) (866,0.1115888873) (870,0.1148917602) (874,0.1067270737) (878,0.09973258939) (882,0.0922610505) (886,0.08616190699) (890,0.09074896004) (894,0.09031232111) (898,0.0965072829) (902,0.1000803672) (906,0.09373399089) (910,0.08772088953) (914,0.08913482651) (918,0.08298973674) (922,0.07548364568) (926,0.08218140968) (930,0.08326995253) (934,0.08525045658) (938,0.07105889121) (942,0.07993845908) (946,0.06928688535) (950,0.07436671938) (954,0.07769964899) (958,0.07446413079) (962,0.06468606545) (966,0.05466454379) (970,0.06247267139) (974,0.07737145179) (978,0.08415783805) (982,0.08202885135) (986,0.08060205195) (990,0.08185686754) (994,0.08533426363) (998,0.0881243977)};
\addplot[pubNavy, opacity=0.20, line width=0.35pt] coordinates {(10,-0.02184085052) (14,-0.0555973604) (18,-0.0518126174) (22,-0.0432105471) (26,-0.02954957447) (30,-0.02666310326) (34,-0.01230575644) (38,-0.02014873216) (42,-0.01868166413) (46,-0.01044155715) (50,0.002253512147) (54,0.004222557466) (58,0.008746388354) (62,-0.006164482532) (66,0.02045142868) (70,0.02017773784) (74,0.03885043411) (78,0.04204213556) (82,0.04066814961) (86,0.04679286218) (90,0.0667157335) (94,0.06109185794) (98,0.07072833295) (102,0.05774184517) (106,0.07705551557) (110,0.08016792514) (114,0.09377588767) (118,0.09577627672) (122,0.1039918124) (126,0.1036382402) (130,0.1073705727) (134,0.08543204501) (138,0.09129849595) (142,0.0831756321) (146,0.08540269613) (150,0.07494752434) (154,0.08721253268) (158,0.06814676261) (162,0.05940584312) (166,0.06016601134) (170,0.04778756199) (174,0.04407535767) (178,0.0300361746) (182,0.05279203631) (186,0.05795235677) (190,0.05625522452) (194,0.06652429212) (198,0.06888651092) (202,0.07655597945) (206,0.06822972974) (210,0.07313527107) (214,0.05917480415) (218,0.06525914393) (222,0.06689453748) (226,0.07746561497) (230,0.07369037349) (234,0.07910997988) (238,0.07603703365) (242,0.06956075178) (246,0.06834180294) (250,0.05875869004) (254,0.06809837109) (258,0.07694813426) (262,0.07960118518) (266,0.07274014702) (270,0.07124157875) (274,0.07456061376) (278,0.07122223887) (282,0.0755332977) (286,0.08246326346) (290,0.08619572445) (294,0.06871948727) (298,0.07332912873) (302,0.08043358653) (306,0.07836851994) (310,0.0667243431) (314,0.06519697939) (318,0.06350927388) (322,0.05373611971) (326,0.05206451901) (330,0.05236063001) (334,0.04073683176) (338,0.03564222283) (342,0.05198446939) (346,0.03847798717) (350,0.04416143708) (354,0.02662960255) (358,0.03219953146) (362,0.04353480988) (366,0.04538910751) (370,0.05061904546) (374,0.05632762238) (378,0.05824734248) (382,0.05551102119) (386,0.06006895744) (390,0.05690026596) (394,0.05198579841) (398,0.05607557543) (402,0.03125551702) (406,0.02600454644) (410,0.04337973602) (414,0.05448554576) (418,0.05999797072) (422,0.03107896361) (426,0.03444185639) (430,0.02672813744) (434,0.0387265198) (438,0.04526366877) (442,0.04076765637) (446,0.05417475875) (450,0.06268154084) (454,0.06497732197) (458,0.07073083229) (462,0.07797358866) (466,0.06994521894) (470,0.06539332829) (474,0.0843436216) (478,0.08303266036) (482,0.08782953311) (486,0.07132035349) (490,0.07880093312) (494,0.07491135016) (498,0.06810623305) (502,0.07114899538) (506,0.07295046109) (510,0.07746944705) (514,0.06414393009) (518,0.05536730551) (522,0.07284359231) (526,0.07512353087) (530,0.06424062794) (534,0.07220996684) (538,0.09203992842) (542,0.07842193803) (546,0.07219180341) (550,0.07701562728) (554,0.08972236871) (558,0.07627427108) (562,0.07447566505) (566,0.0940145363) (570,0.09152576996) (574,0.08600452916) (578,0.09012439337) (582,0.08770689328) (586,0.08140335229) (590,0.0761792674) (594,0.08544929543) (598,0.1078400892) (602,0.1110950351) (606,0.1008377296) (610,0.117661774) (614,0.1152674312) (618,0.1156008188) (622,0.1119907236) (626,0.1165379876) (630,0.1176826658) (634,0.1107055103) (638,0.1106616679) (642,0.1120482207) (646,0.09960069938) (650,0.08192276471) (654,0.08178921964) (658,0.08534658603) (662,0.09120578379) (666,0.07642151332) (670,0.08398469066) (674,0.09970203297) (678,0.08769225293) (682,0.1020261538) (686,0.1070707311) (690,0.1005340432) (694,0.1059734552) (698,0.1064103309) (702,0.1014892462) (706,0.09140771245) (710,0.09584786232) (714,0.07740003623) (718,0.09568859622) (722,0.0822609333) (726,0.08336223791) (730,0.07130786351) (734,0.07445054878) (738,0.06688591582) (742,0.05712124557) (746,0.04482804019) (750,0.04474116292) (754,0.0375928032) (758,0.0297115109) (762,0.01941551269) (766,0.0395532069) (770,0.03771270937) (774,0.03187378649) (778,0.02502043237) (782,0.02523524899) (786,0.03860953783) (790,0.05320594058) (794,0.05365375116) (798,0.06183818792) (802,0.08788515389) (806,0.09767605778) (810,0.1005559102) (814,0.1006467567) (818,0.1041785128) (822,0.1043028677) (826,0.1089061847) (830,0.1189704711) (834,0.128845725) (838,0.1266853227) (842,0.1228001264) (846,0.1197309821) (850,0.1147824364) (854,0.1121198564) (858,0.1011546675) (862,0.1140918997) (866,0.09203529839) (870,0.09421807351) (874,0.08616581538) (878,0.08216808105) (882,0.07239221672) (886,0.0564632481) (890,0.06863363555) (894,0.07598621966) (898,0.08297073875) (902,0.09229257934) (906,0.08809990788) (910,0.08763574457) (914,0.09455830844) (918,0.09777871298) (922,0.0919350187) (926,0.09129656406) (930,0.08110731214) (934,0.07695085692) (938,0.05793229633) (942,0.06699457591) (946,0.05333873284) (950,0.05097659305) (954,0.05229667406) (958,0.05535648765) (962,0.04003190363) (966,0.03299987165) (970,0.03315668864) (974,0.0574375854) (978,0.07182384373) (982,0.07111561468) (986,0.07440188543) (990,0.07442219291) (994,0.05985702928) (998,0.07742477759)};
\addplot[pubNavy, opacity=0.20, line width=0.35pt] coordinates {(10,-0.04031765922) (14,-0.03898422697) (18,-0.03720837615) (22,-0.04013604593) (26,-0.01799457382) (30,-0.02222120032) (34,-0.01364166435) (38,-0.02874091026) (42,-0.02448414464) (46,-0.01953293184) (50,-0.004519322171) (54,-0.003410582636) (58,0.002432926525) (62,-0.01168573096) (66,0.004435210394) (70,0.00539101291) (74,0.01762442472) (78,0.01857440472) (82,0.01710252564) (86,0.0340420305) (90,0.04125567704) (94,0.03509228466) (98,0.03832888154) (102,0.02726484319) (106,0.04297856039) (110,0.04486775771) (114,0.05473178187) (118,0.05464148217) (122,0.05835190272) (126,0.05541561626) (130,0.06565235881) (134,0.04225766577) (138,0.04281135727) (142,0.04425897832) (146,0.0489846643) (150,0.0342035854) (154,0.04492359671) (158,0.02103755117) (162,0.01254639792) (166,0.01505597029) (170,-0.002711463309) (174,0.0002633326203) (178,-0.01295474528) (182,0.0005730775729) (186,0.008292480362) (190,0.005420414665) (194,0.01399330634) (198,0.01925406724) (202,0.02361751951) (206,0.01657064536) (210,0.01881985656) (214,0.003397325557) (218,0.01369360909) (222,0.02192534545) (226,0.03434329839) (230,0.02558356506) (234,0.02787480303) (238,0.02305749735) (242,0.02068474947) (246,0.02140291139) (250,0.02407544283) (254,0.03916938787) (258,0.05127311576) (262,0.0463968217) (266,0.05216843744) (270,0.0494187591) (274,0.04848998239) (278,0.05339952449) (282,0.06368694762) (286,0.06666819191) (290,0.07748494437) (294,0.06369525837) (298,0.06583719246) (302,0.07183429361) (306,0.07310431626) (310,0.06312153469) (314,0.06242164564) (318,0.05877296254) (322,0.05194242582) (326,0.04923411952) (330,0.04451015639) (334,0.04324095618) (338,0.03784708716) (342,0.05436004387) (346,0.04056442381) (350,0.03598461232) (354,0.01126229484) (358,0.02187384778) (362,0.03597855829) (366,0.03504834264) (370,0.04246301007) (374,0.04066985879) (378,0.0481469586) (382,0.03843928442) (386,0.04697393612) (390,0.04661982207) (394,0.04208353247) (398,0.05044945602) (402,0.03463678782) (406,0.02985075615) (410,0.04862865464) (414,0.05648316684) (418,0.05997182364) (422,0.05252648342) (426,0.05736416766) (430,0.05470261557) (434,0.05316185971) (438,0.05439939314) (442,0.05282659729) (446,0.07406062926) (450,0.08137191921) (454,0.08669662976) (458,0.0926522159) (462,0.1008459858) (466,0.09142139249) (470,0.07974327739) (474,0.09209652496) (478,0.08755753245) (482,0.1003533181) (486,0.09754844792) (490,0.09685187264) (494,0.08712370713) (498,0.08687889541) (502,0.08952234807) (506,0.08386392366) (510,0.07476616713) (514,0.06186675427) (518,0.06151590033) (522,0.0747839378) (526,0.07055578954) (530,0.06247775584) (534,0.06171260314) (538,0.07004041633) (542,0.05616276488) (546,0.06301235511) (550,0.06153613933) (554,0.07489912901) (558,0.06407063553) (562,0.06509801118) (566,0.07296065674) (570,0.07013937389) (574,0.05822156458) (578,0.05341136861) (582,0.05556342675) (586,0.05233918987) (590,0.05162495259) (594,0.05359715117) (598,0.06448289008) (602,0.0527215373) (606,0.0450858593) (610,0.06116832028) (614,0.06201040904) (618,0.05975990447) (622,0.06369112747) (626,0.08134102812) (630,0.0826637318) (634,0.08286766402) (638,0.08366810612) (642,0.0842958011) (646,0.07587760003) (650,0.07665380826) (654,0.08087042021) (658,0.08029869795) (662,0.08185372652) (666,0.08129110145) (670,0.08224425107) (674,0.07863325123) (678,0.06588853928) (682,0.07376411687) (686,0.07587497428) (690,0.07727692377) (694,0.07215596145) (698,0.06374746368) (702,0.06598714786) (706,0.06477895989) (710,0.07556715216) (714,0.05347045725) (718,0.06074357836) (722,0.0605169158) (726,0.06057880198) (730,0.05946069553) (734,0.05800169949) (738,0.05468992321) (742,0.05705692508) (746,0.06079920725) (750,0.05573776707) (754,0.04077711904) (758,0.04054814841) (762,0.0403321773) (766,0.05865962472) (770,0.05823254292) (774,0.05794939741) (778,0.05534661936) (782,0.05758654957) (786,0.06147215728) (790,0.05262413716) (794,0.04809876014) (798,0.04271572656) (802,0.07098489866) (806,0.06056405237) (810,0.05427396519) (814,0.04835590661) (818,0.05438925628) (822,0.04545794845) (826,0.04651153598) (830,0.04826784952) (834,0.05047112607) (838,0.05017349843) (842,0.0541427527) (846,0.0534586778) (850,0.05033525629) (854,0.04572297091) (858,0.05035338277) (862,0.06457068472) (866,0.04946576989) (870,0.05299622414) (874,0.04744066914) (878,0.03668962498) (882,0.0295205023) (886,0.02423509296) (890,0.03267292944) (894,0.04069629476) (898,0.03993928701) (902,0.05059290885) (906,0.04582347987) (910,0.05193056364) (914,0.04928530461) (918,0.05850069126) (922,0.05617810946) (926,0.06338364572) (930,0.06493829663) (934,0.05912387347) (938,0.04727141263) (942,0.05069005348) (946,0.04306755482) (950,0.04837227591) (954,0.05032169008) (958,0.04215393963) (962,0.03334334801) (966,0.02327826786) (970,0.01806525004) (974,0.0329358812) (978,0.04480640728) (982,0.04534123339) (986,0.05682964346) (990,0.05664634041) (994,0.05515847732) (998,0.06375146845)};
\addplot[pubNavy, opacity=0.20, line width=0.35pt] coordinates {(10,-0.01258736669) (14,-0.02678758765) (18,-0.03726923421) (22,-0.03129076324) (26,-0.008933419126) (30,-0.01161786062) (34,-0.006822284321) (38,-0.01866845266) (42,-0.01473239647) (46,-0.01976028568) (50,-0.005369212998) (54,-0.02483443064) (58,-0.02474838716) (62,-0.04597825776) (66,-0.03004819025) (70,-0.0361746944) (74,-0.02733154863) (78,-0.03400076787) (82,-0.02818028225) (86,-0.01714247189) (90,0.001412193472) (94,-0.005813950562) (98,0.01111081383) (102,-0.01955629602) (106,0.009871335037) (110,0.0225505525) (114,0.03496614667) (118,0.03942672547) (122,0.04090338152) (126,0.03469616986) (130,0.04451582305) (134,0.01928476152) (138,0.01827361699) (142,0.005168337024) (146,0.0151867762) (150,0.01011522122) (154,0.02353865755) (158,0.01620342663) (162,0.01281265104) (166,0.01512284102) (170,-0.009096536794) (174,-0.001768217902) (178,0.001231541876) (182,0.01741942809) (186,0.0251789627) (190,0.01868550072) (194,0.02263148093) (198,0.0205249788) (202,0.01626789601) (206,0.01379968991) (210,0.01502041032) (214,0.002511117824) (218,0.01388252664) (222,0.01477633316) (226,0.0181921215) (230,0.003015151048) (234,-0.00433547584) (238,-0.00559230645) (242,-0.00308000484) (246,0.002887239239) (250,0.005420244581) (254,0.01593826699) (258,0.03061489824) (262,0.02554925275) (266,0.01636287849) (270,0.01338871952) (274,0.02301726462) (278,0.02674992089) (282,0.03781722737) (286,0.04370166965) (290,0.04307693042) (294,0.02542030392) (298,0.01499825037) (302,0.01846224264) (306,0.008370553765) (310,-0.003443966083) (314,0.002541488075) (318,0.008097945577) (322,0.002815093763) (326,-0.0023483695) (330,-0.008054810379) (334,-0.01088170094) (338,-0.01524229452) (342,0.007335436396) (346,0.00635080026) (350,-0.004172388583) (354,-0.01145963811) (358,-0.008530937245) (362,-0.01166819769) (366,-0.02382741749) (370,-0.03219615251) (374,-0.01778072903) (378,-0.01635470218) (382,-0.01794802813) (386,-0.02064986448) (390,-0.02812375083) (394,-0.02997713685) (398,-0.009080845275) (402,-0.03253707519) (406,-0.01918842719) (410,0.0009069894754) (414,0.009377185637) (418,0.0252456788) (422,0.01467424792) (426,0.02626138003) (430,0.02226708046) (434,0.03365135731) (438,0.03402273928) (442,0.02252211973) (446,0.03251863813) (450,0.05006014054) (454,0.04218409231) (458,0.03883314949) (462,0.05026889278) (466,0.03807238368) (470,0.03738430924) (474,0.04910848989) (478,0.03466520154) (482,0.03149452647) (486,0.01998917979) (490,0.02996347745) (494,0.03752118256) (498,0.03133539805) (502,0.02429131213) (506,0.02221874311) (510,0.02142862543) (514,0.01866736584) (518,-0.006082436414) (522,0.01228956623) (526,0.00582371162) (530,0.01589111208) (534,0.02378971317) (538,0.03348157709) (542,0.02347588175) (546,0.01776571441) (550,0.02728394979) (554,0.05852596477) (558,0.04144824094) (562,0.0327634707) (566,0.04229194611) (570,0.063152205) (574,0.05705650652) (578,0.06768645094) (582,0.06360477732) (586,0.06793436239) (590,0.06022845907) (594,0.06204482302) (598,0.07682702915) (602,0.07635548686) (606,0.05891871457) (610,0.0897491836) (614,0.09023627333) (618,0.09024594621) (622,0.08506440311) (626,0.08638125967) (630,0.07832960768) (634,0.06562404154) (638,0.05389900853) (642,0.06092563756) (646,0.06257125605) (650,0.05408096322) (654,0.06118035378) (658,0.06766880808) (662,0.08154098868) (666,0.06795790319) (670,0.07390029334) (674,0.08334690675) (678,0.07864183949) (682,0.09267530282) (686,0.09317055893) (690,0.09567498183) (694,0.0930932471) (698,0.09245335856) (702,0.08660699399) (706,0.07051759523) (710,0.07789488789) (714,0.04697763335) (718,0.06166191697) (722,0.05183227161) (726,0.04693419805) (730,0.04837355409) (734,0.04799197824) (738,0.05363786178) (742,0.05270848108) (746,0.05201630499) (750,0.05563401976) (754,0.03768061839) (758,0.04189044202) (762,0.03902154268) (766,0.05112693274) (770,0.05729862516) (774,0.05200491938) (778,0.04216904551) (782,0.0579194889) (786,0.06256283093) (790,0.0608312774) (794,0.05635322366) (798,0.0511351478) (802,0.07076473507) (806,0.07046544578) (810,0.0600649217) (814,0.05365535341) (818,0.05991779771) (822,0.06450031628) (826,0.07285457611) (830,0.07436515694) (834,0.07341846823) (838,0.07576273205) (842,0.07012460913) (846,0.06811782187) (850,0.06513070312) (854,0.06470900259) (858,0.06133635619) (862,0.07032166992) (866,0.06807143428) (870,0.06669032461) (874,0.06372827154) (878,0.05808835052) (882,0.0467806741) (886,0.03096258254) (890,0.03423426427) (894,0.04345803076) (898,0.0534108773) (902,0.06443059674) (906,0.05874164959) (910,0.06411997623) (914,0.07349633672) (918,0.06867631365) (922,0.0612268357) (926,0.06412585056) (930,0.0525289084) (934,0.06488266958) (938,0.04684438907) (942,0.06225102133) (946,0.0510527386) (950,0.05107600405) (954,0.04937281518) (958,0.05401817999) (962,0.04699164061) (966,0.04245047478) (970,0.0331672101) (974,0.04333994944) (978,0.06361594086) (982,0.06915322685) (986,0.06718373792) (990,0.05727701405) (994,0.05518011013) (998,0.06019689731)};
\addplot[pubNavy, opacity=0.20, line width=0.35pt] coordinates {(10,-0.05308977987) (14,-0.05914591933) (18,-0.04701855652) (22,-0.04144130705) (26,-0.03521840697) (30,-0.02986140172) (34,-0.02227226771) (38,-0.03423903879) (42,-0.03572555186) (46,-0.02869042505) (50,-0.01940992696) (54,-0.01767054537) (58,-0.001532324625) (62,-0.001383091109) (66,0.01463773948) (70,0.01680502668) (74,0.02954791314) (78,0.03543826956) (82,0.03862477307) (86,0.04968741157) (90,0.07143100709) (94,0.06984668236) (98,0.08256366965) (102,0.0810955338) (106,0.09740146808) (110,0.09333972266) (114,0.09602399447) (118,0.107356353) (122,0.1134368405) (126,0.1171099284) (130,0.1233994435) (134,0.1163512997) (138,0.1202556608) (142,0.1285870502) (146,0.1300037119) (150,0.1245417851) (154,0.1230232469) (158,0.1141342468) (162,0.1151053155) (166,0.1122797577) (170,0.1012322572) (174,0.09836734273) (178,0.09986177433) (182,0.1004974536) (186,0.1010634435) (190,0.09385168018) (194,0.09137437679) (198,0.09009560446) (202,0.1003881798) (206,0.09592320664) (210,0.1020483938) (214,0.1023900927) (218,0.114554319) (222,0.1150588011) (226,0.1214783686) (230,0.1093334004) (234,0.1145413844) (238,0.1103998255) (242,0.1143735948) (246,0.1238507282) (250,0.1217575337) (254,0.1299022116) (258,0.123876121) (262,0.1184453924) (266,0.09913205788) (270,0.09386773183) (274,0.09900379522) (278,0.1014766846) (282,0.1023739292) (286,0.1047912326) (290,0.1176209604) (294,0.1049968006) (298,0.1035849267) (302,0.09604179307) (306,0.1027622184) (310,0.09663173313) (314,0.1051370599) (318,0.11148304) (322,0.111213881) (326,0.1084668882) (330,0.1065913693) (334,0.1069264134) (338,0.09899346488) (342,0.1076732586) (346,0.09610500947) (350,0.09466819072) (354,0.08109323758) (358,0.07627580731) (362,0.07818411824) (366,0.08007735281) (370,0.0846850802) (374,0.0925078658) (378,0.08856772673) (382,0.08603393976) (386,0.08759673677) (390,0.08241862128) (394,0.08010937925) (398,0.09296688906) (402,0.08879642564) (406,0.09506429291) (410,0.1075195726) (414,0.1042878134) (418,0.1164350617) (422,0.1011027922) (426,0.1074376381) (430,0.1002199527) (434,0.1014905806) (438,0.09949313409) (442,0.09212808272) (446,0.09929712925) (450,0.1038627794) (454,0.09361051037) (458,0.08905621062) (462,0.08708147393) (466,0.0753150344) (470,0.06879971276) (474,0.07888725684) (478,0.06806191649) (482,0.07744645233) (486,0.0753087415) (490,0.08460243172) (494,0.0832053402) (498,0.07708923901) (502,0.08349484117) (506,0.09825980185) (510,0.09760652606) (514,0.09713169747) (518,0.0940012417) (522,0.1029215842) (526,0.1089609828) (530,0.1020043178) (534,0.1068019625) (538,0.1178555246) (542,0.1041018241) (546,0.1033984175) (550,0.1021968586) (554,0.1040068726) (558,0.1082214385) (562,0.1058181897) (566,0.1103567584) (570,0.115231799) (574,0.1121390737) (578,0.1084612643) (582,0.1067897073) (586,0.09970860632) (590,0.09248903921) (594,0.09024817208) (598,0.104646346) (602,0.1064522128) (606,0.08949431758) (610,0.1002496328) (614,0.1061711951) (618,0.1068760405) (622,0.1020703488) (626,0.1112876567) (630,0.1150804747) (634,0.1198236452) (638,0.1210664559) (642,0.1313948409) (646,0.1249887629) (650,0.1268435779) (654,0.1310261576) (658,0.1346715738) (662,0.1316302137) (666,0.1249434153) (670,0.1304980072) (674,0.1282941466) (678,0.1231268813) (682,0.1336889458) (686,0.1371535713) (690,0.128829656) (694,0.1245531494) (698,0.1097601182) (702,0.1127147334) (706,0.1115962449) (710,0.1218545371) (714,0.1152276504) (718,0.1161945896) (722,0.1071277255) (726,0.1087542942) (730,0.1016718792) (734,0.1075135438) (738,0.1107006603) (742,0.1196297841) (746,0.109174592) (750,0.1215397358) (754,0.1138862917) (758,0.1127929325) (762,0.113693262) (766,0.1211324774) (770,0.1271254676) (774,0.1216212556) (778,0.1200296035) (782,0.118809628) (786,0.1179346203) (790,0.1232102738) (794,0.1271036184) (798,0.1310060748) (802,0.1377677074) (806,0.1386723214) (810,0.1411813166) (814,0.1443354092) (818,0.1454556864) (822,0.1421633348) (826,0.138879767) (830,0.1436984712) (834,0.151679001) (838,0.1445563789) (842,0.1437434771) (846,0.140215089) (850,0.1328747506) (854,0.1389472977) (858,0.1411925568) (862,0.1395751471) (866,0.1323000368) (870,0.1382609825) (874,0.1293719627) (878,0.1193186759) (882,0.1135064277) (886,0.1033244275) (890,0.1080151128) (894,0.1004431262) (898,0.1059037877) (902,0.1151728093) (906,0.1054496536) (910,0.09762700333) (914,0.09977163108) (918,0.09838994086) (922,0.08916909045) (926,0.102597836) (930,0.1026668515) (934,0.1073606176) (938,0.09835821343) (942,0.1191157223) (946,0.1152752979) (950,0.1247375001) (954,0.1297017473) (958,0.1212017507) (962,0.1054614852) (966,0.1050307933) (970,0.1109877641) (974,0.129280818) (978,0.1305340711) (982,0.1272668986) (986,0.1237705952) (990,0.1112460963) (994,0.1041849642) (998,0.1050684034)};
\addplot[pubNavy, opacity=0.20, line width=0.35pt] coordinates {(10,-0.02913269242) (14,-0.03560664833) (18,-0.02502105679) (22,-0.0196438776) (26,-0.004335982691) (30,0.001472650086) (34,-0.0009531076361) (38,-0.01366450523) (42,-0.01866386037) (46,-0.03132860258) (50,-0.01995473181) (54,-0.05272061694) (58,-0.053485024) (62,-0.07508105592) (66,-0.06759887694) (70,-0.08449263169) (74,-0.07964628748) (78,-0.09549319052) (82,-0.09419272368) (86,-0.0758471295) (90,-0.05766017276) (94,-0.05855727189) (98,-0.02968045836) (102,-0.04834974324) (106,-0.007715241422) (110,0.01045352038) (114,0.0245810554) (118,0.03929553011) (122,0.04819644924) (126,0.05652911925) (130,0.07284726245) (134,0.05309719994) (138,0.05431710339) (142,0.04608286352) (146,0.05420687565) (150,0.05022787766) (154,0.05262177095) (158,0.03785847973) (162,0.03119126144) (166,0.02051145105) (170,0.0037253655) (174,0.0004756462477) (178,-0.01490000318) (182,-0.003714099336) (186,-0.002687252685) (190,-0.00866743956) (194,-0.0008022383292) (198,0.006263984688) (202,0.01097342107) (206,0.009894406091) (210,0.02230709396) (214,0.02866567543) (218,0.03645662541) (222,0.04445828735) (226,0.0676648085) (230,0.06201747266) (234,0.06422505071) (238,0.06144374577) (242,0.06411345778) (246,0.06569182486) (250,0.0589704864) (254,0.07776348038) (258,0.0776904949) (262,0.06961428976) (266,0.06580982022) (270,0.05453809442) (274,0.05195008922) (278,0.05080160289) (282,0.05724375434) (286,0.05605760253) (290,0.06113326837) (294,0.04508148832) (298,0.04250976965) (302,0.04204453465) (306,0.04774249401) (310,0.03347794811) (314,0.04301086288) (318,0.04891148567) (322,0.05206506751) (326,0.05214972766) (330,0.0512284571) (334,0.05501944259) (338,0.04924452298) (342,0.06966555171) (346,0.05697895283) (350,0.05508298775) (354,0.04142270201) (358,0.04233509347) (362,0.05810819511) (366,0.05862925867) (370,0.05029852931) (374,0.05421227285) (378,0.05676808908) (382,0.05504391365) (386,0.05667023854) (390,0.05402033224) (394,0.05265828707) (398,0.06740754158) (402,0.04906357765) (406,0.04680620062) (410,0.05524618753) (414,0.05260613055) (418,0.06894246909) (422,0.06301858026) (426,0.06390340208) (430,0.05258547515) (434,0.05650943776) (438,0.06145282189) (442,0.05684740449) (446,0.06578281074) (450,0.07119104033) (454,0.07136612434) (458,0.0730191268) (462,0.09025603342) (466,0.07033067441) (470,0.05564900728) (474,0.07957456361) (478,0.06867733002) (482,0.08100635111) (486,0.0672135473) (490,0.07085561329) (494,0.08080889452) (498,0.07205066318) (502,0.06939110286) (506,0.07756408251) (510,0.06772120404) (514,0.06023499742) (518,0.04243862154) (522,0.05373114219) (526,0.05650389308) (530,0.04728631928) (534,0.04742923155) (538,0.06124523807) (542,0.04293078523) (546,0.03442113142) (550,0.034216122) (554,0.04937802781) (558,0.03955787706) (562,0.03728306401) (566,0.05351762742) (570,0.05857142242) (574,0.04435299426) (578,0.0440711542) (582,0.03616272449) (586,0.03287898051) (590,0.0351366009) (594,0.03026938526) (598,0.04931473627) (602,0.05719200735) (606,0.04366528903) (610,0.05853147849) (614,0.05759505619) (618,0.05980550935) (622,0.06695506189) (626,0.0823945168) (630,0.0832130281) (634,0.08360333486) (638,0.07703473757) (642,0.08883289347) (646,0.0755625329) (650,0.06650791167) (654,0.07028017373) (658,0.07710650163) (662,0.08677541764) (666,0.0793427906) (670,0.08518819556) (674,0.09914442584) (678,0.08492276646) (682,0.0887756819) (686,0.09300726156) (690,0.09131757875) (694,0.09520062382) (698,0.08895758727) (702,0.08654540203) (706,0.0830814674) (710,0.07822283405) (714,0.06158114962) (718,0.06740307074) (722,0.05711543695) (726,0.05822213407) (730,0.0615878489) (734,0.06323593462) (738,0.06292916005) (742,0.06311138237) (746,0.05289408732) (750,0.05349182246) (754,0.04987097666) (758,0.03974943561) (762,0.0484599986) (766,0.0632703111) (770,0.07361443198) (774,0.06346980825) (778,0.0576789867) (782,0.05126358036) (786,0.06012480245) (790,0.05684025264) (794,0.05706819566) (798,0.0658778385) (802,0.07911218431) (806,0.08111724946) (810,0.07742674859) (814,0.07202915329) (818,0.07380036196) (822,0.07627163759) (826,0.08013688759) (830,0.08558196486) (834,0.1004219582) (838,0.09458715396) (842,0.09950781412) (846,0.1039220618) (850,0.09657972603) (854,0.09626712469) (858,0.08440082551) (862,0.08961035207) (866,0.07047434609) (870,0.06324938865) (874,0.05940251765) (878,0.05285450785) (882,0.03172640754) (886,0.01142304449) (890,0.02204148903) (894,0.02495089043) (898,0.02975728557) (902,0.03664640363) (906,0.04367005369) (910,0.04732048558) (914,0.05279801126) (918,0.05696334792) (922,0.05512889935) (926,0.04990254434) (930,0.05381504789) (934,0.05016343846) (938,0.02973686663) (942,0.03847493345) (946,0.02458865405) (950,0.03281326571) (954,0.03443017851) (958,0.03735031006) (962,0.03135165531) (966,0.02627494528) (970,0.02781723708) (974,0.04407343103) (978,0.0575134605) (982,0.06325985718) (986,0.06762349339) (990,0.07403145239) (994,0.07138775983) (998,0.07531373005)};
\addplot[pubNavy, opacity=0.20, line width=0.35pt] coordinates {(10,-0.03234552148) (14,-0.05166211273) (18,-0.03255280124) (22,-0.03505938935) (26,-0.01870787299) (30,-0.01713282903) (34,-0.009766004006) (38,-0.01651606008) (42,-0.01316088002) (46,-0.009717610961) (50,0.001423942574) (54,0.002045177014) (58,0.004068162028) (62,-0.0003244184448) (66,0.01782332217) (70,0.02085534864) (74,0.0362907636) (78,0.03846101218) (82,0.04174177561) (86,0.05408350243) (90,0.07278557214) (94,0.06756929705) (98,0.0764925104) (102,0.07255971676) (106,0.0931987778) (110,0.08863846557) (114,0.09844973595) (118,0.1025623395) (122,0.1080329768) (126,0.1033321933) (130,0.1100694581) (134,0.09515353724) (138,0.09207949151) (142,0.09092676666) (146,0.09029028655) (150,0.08291721082) (154,0.09082761562) (158,0.0730778415) (162,0.07495989805) (166,0.06402664446) (170,0.0513499806) (174,0.05297272636) (178,0.04523232174) (182,0.0562185789) (186,0.05671374249) (190,0.04849119267) (194,0.05305362671) (198,0.05325921518) (202,0.04948819685) (206,0.03254769681) (210,0.04417100473) (214,0.02685257515) (218,0.03992150161) (222,0.03689646988) (226,0.04211649795) (230,0.03371641502) (234,0.03183180939) (238,0.02853584677) (242,0.02861820672) (246,0.03786037162) (250,0.02906948269) (254,0.05048641176) (258,0.05712317198) (262,0.05335416953) (266,0.04808518589) (270,0.04049627778) (274,0.04926949688) (278,0.04502668936) (282,0.04516457178) (286,0.04650600727) (290,0.04796260009) (294,0.0370589674) (298,0.03162243928) (302,0.03888249574) (306,0.04032238158) (310,0.03995258235) (314,0.0534559578) (318,0.06307607782) (322,0.0599930576) (326,0.06602668397) (330,0.06186365703) (334,0.06604850903) (338,0.07129154528) (342,0.08737558161) (346,0.08612423127) (350,0.08208616501) (354,0.058937759) (358,0.05888877118) (362,0.06432268621) (366,0.05495960307) (370,0.05509318603) (374,0.06675982477) (378,0.06956208294) (382,0.07167664198) (386,0.06577990749) (390,0.06058372003) (394,0.05913888338) (398,0.06975162382) (402,0.05776851006) (406,0.06317804535) (410,0.07663897043) (414,0.07436732479) (418,0.08391164637) (422,0.0752233573) (426,0.07605236445) (430,0.07167091437) (434,0.0788783892) (438,0.07891343029) (442,0.06717115146) (446,0.07916162468) (450,0.09340476748) (454,0.08822985852) (458,0.09540410723) (462,0.1038958365) (466,0.09412274162) (470,0.08705372484) (474,0.09249268591) (478,0.08706834078) (482,0.09895923959) (486,0.09640737704) (490,0.1048744931) (494,0.1040057356) (498,0.09912931479) (502,0.1021781406) (506,0.1030125301) (510,0.09499973081) (514,0.08291102136) (518,0.08235759869) (522,0.09290137694) (526,0.09267753918) (530,0.08649261662) (534,0.08782532872) (538,0.08812459946) (542,0.08537139391) (546,0.08410139442) (550,0.08870869386) (554,0.09066131577) (558,0.09650541437) (562,0.09936587759) (566,0.09817126079) (570,0.100912191) (574,0.08836090462) (578,0.08841930905) (582,0.08687442473) (586,0.08558681861) (590,0.07678778137) (594,0.07515747773) (598,0.08271067211) (602,0.07693060657) (606,0.0749209265) (610,0.07537875291) (614,0.08146854464) (618,0.08430380404) (622,0.08654459557) (626,0.09097928584) (630,0.09461628321) (634,0.08454192131) (638,0.08768364379) (642,0.08972325442) (646,0.08613035955) (650,0.08877741367) (654,0.08481901497) (658,0.08607729803) (662,0.08736291549) (666,0.08266076927) (670,0.07774476162) (674,0.07341722899) (678,0.06699403823) (682,0.08170032521) (686,0.08241301072) (690,0.08612093608) (694,0.08629318) (698,0.08755765243) (702,0.08235527501) (706,0.08211673555) (710,0.08817519164) (714,0.07004022002) (718,0.08046457963) (722,0.08478365099) (726,0.08910728155) (730,0.08685703937) (734,0.08825577266) (738,0.09192396049) (742,0.09661322921) (746,0.09022102012) (750,0.08581542984) (754,0.07568398229) (758,0.07715804198) (762,0.08610341932) (766,0.09567325864) (770,0.09439579752) (774,0.0850957188) (778,0.08390873976) (782,0.08383906361) (786,0.08485683394) (790,0.07461596112) (794,0.07216770209) (798,0.071647196) (802,0.08686514861) (806,0.08338723541) (810,0.0669054471) (814,0.05921139199) (818,0.06348379659) (822,0.06626553705) (826,0.0726484217) (830,0.07872562354) (834,0.08872596303) (838,0.08186832073) (842,0.08628039477) (846,0.08876813125) (850,0.08968780202) (854,0.08748760117) (858,0.08649234798) (862,0.102534091) (866,0.09559614271) (870,0.09691992717) (874,0.09055383956) (878,0.07963147194) (882,0.06961913492) (886,0.05591199692) (890,0.06699512863) (894,0.06650432143) (898,0.06776909991) (902,0.0740908267) (906,0.07552532554) (910,0.0767973351) (914,0.08007220518) (918,0.0794606037) (922,0.07479501055) (926,0.07461221763) (930,0.07829195458) (934,0.0864581288) (938,0.07211828356) (942,0.07787647538) (946,0.06635663405) (950,0.07301029025) (954,0.07650594655) (958,0.07946057233) (962,0.07005340065) (966,0.06655232589) (970,0.06951103659) (974,0.07703181424) (978,0.08464449284) (982,0.07996855483) (986,0.076624843) (990,0.07406993846) (994,0.0792972943) (998,0.08524685816)};
\addplot[pubNavy, opacity=0.20, line width=0.35pt] coordinates {(10,-0.05540238998) (14,-0.03719331944) (18,-0.03079692397) (22,-0.01890687617) (26,-0.003418121674) (30,-0.01335824697) (34,-0.003497556944) (38,-0.01710821373) (42,-0.01536025847) (46,-0.02472924412) (50,-0.01183612691) (54,-0.02851357719) (58,-0.02393592637) (62,-0.03823702958) (66,-0.03208195811) (70,-0.04283518037) (74,-0.03661611408) (78,-0.03920930537) (82,-0.03523972004) (86,-0.01877547828) (90,-0.002760752624) (94,-0.00398341479) (98,0.02593714329) (102,0.006965515854) (106,0.03724568108) (110,0.04903070771) (114,0.06668272169) (118,0.07707987725) (122,0.08811694739) (126,0.09458259903) (130,0.1089311083) (134,0.08377540568) (138,0.09190820743) (142,0.08042868966) (146,0.08784418353) (150,0.08190421648) (154,0.08866677248) (158,0.0716581854) (162,0.06373857056) (166,0.05852277468) (170,0.04043513797) (174,0.02363171734) (178,0.01191101221) (182,0.01680407468) (186,0.01949894821) (190,0.01201223141) (194,0.02429090285) (198,0.02433750038) (202,0.0278468026) (206,0.01560229778) (210,0.02949396458) (214,0.02616959904) (218,0.03492384226) (222,0.04144429602) (226,0.05623969315) (230,0.05775478649) (234,0.0584463117) (238,0.06270391064) (242,0.05131107562) (246,0.05218485366) (250,0.04614257375) (254,0.06116928739) (258,0.06935885183) (262,0.07057107878) (266,0.06780406459) (270,0.06285418925) (274,0.06757551603) (278,0.06257455802) (282,0.0667046645) (286,0.06448352254) (290,0.07302066367) (294,0.0623787328) (298,0.06125273355) (302,0.06765759247) (306,0.06616092617) (310,0.05719267349) (314,0.06316375888) (318,0.05888591992) (322,0.06523358975) (326,0.05738495023) (330,0.06662657294) (334,0.07242172387) (338,0.06755043039) (342,0.07931017135) (346,0.07195313198) (350,0.066228276) (354,0.04230315011) (358,0.04365828726) (362,0.05912012003) (366,0.06143705704) (370,0.05374843325) (374,0.06505821294) (378,0.06746479661) (382,0.06183355714) (386,0.06433724264) (390,0.06186764887) (394,0.05948932396) (398,0.06950022162) (402,0.04593376806) (406,0.05111002275) (410,0.06585521011) (414,0.06676594582) (418,0.07742090957) (422,0.06398034176) (426,0.07120869508) (430,0.05152865493) (434,0.05682544939) (438,0.06586095632) (442,0.05852765701) (446,0.07333774718) (450,0.08059241386) (454,0.07908290884) (458,0.07906383185) (462,0.08411620146) (466,0.07629640383) (470,0.06792459898) (474,0.08158774162) (478,0.07294808268) (482,0.08455847125) (486,0.06480532038) (490,0.07311733189) (494,0.07549271688) (498,0.07140279109) (502,0.06852720287) (506,0.078782961) (510,0.08220814273) (514,0.06986425712) (518,0.06054313485) (522,0.07835484718) (526,0.07768809384) (530,0.0795001762) (534,0.07988490441) (538,0.09441058077) (542,0.07789967535) (546,0.06827103477) (550,0.06125735788) (554,0.07070901705) (558,0.05408129436) (562,0.05124577837) (566,0.06827835006) (570,0.05729204253) (574,0.04658899477) (578,0.04566692435) (582,0.04166264448) (586,0.0418739077) (590,0.04757496382) (594,0.05622200822) (598,0.07649334129) (602,0.09330342494) (606,0.08763809613) (610,0.1071012197) (614,0.1087975854) (618,0.1052574608) (622,0.1077186319) (626,0.1145691058) (630,0.1125090892) (634,0.1002946949) (638,0.09592504144) (642,0.09006523645) (646,0.07939871718) (650,0.07258194295) (654,0.07222731549) (658,0.06933600675) (662,0.07834542387) (666,0.06122155945) (670,0.0755512189) (674,0.08467063153) (678,0.07204489674) (682,0.08371374421) (686,0.08977882089) (690,0.09280202402) (694,0.1018575666) (698,0.09447959732) (702,0.0893562856) (706,0.08251012983) (710,0.08671358428) (714,0.07371996212) (718,0.08269733252) (722,0.07258509797) (726,0.07195444844) (730,0.07284779917) (734,0.07600992509) (738,0.0799847773) (742,0.06870988329) (746,0.05993050738) (750,0.06107663394) (754,0.05642405771) (758,0.05473951363) (762,0.04814021407) (766,0.05760050681) (770,0.06344762456) (774,0.0513577796) (778,0.04240519221) (782,0.04637161947) (786,0.05620438748) (790,0.05100305501) (794,0.0558393704) (798,0.06028842301) (802,0.07696624967) (806,0.07313642116) (810,0.08075347453) (814,0.08467711061) (818,0.09173083125) (822,0.1051819491) (826,0.1092420884) (830,0.1086983241) (834,0.1133512661) (838,0.1183699558) (842,0.1214449655) (846,0.1235575502) (850,0.118758298) (854,0.1227689314) (858,0.1158517816) (862,0.1194198871) (866,0.1088242458) (870,0.09994681274) (874,0.08706251133) (878,0.0839243146) (882,0.07980189658) (886,0.05655516365) (890,0.05850588023) (894,0.055075184) (898,0.05862375463) (902,0.06904863139) (906,0.0574313246) (910,0.05666635817) (914,0.05819298892) (918,0.05468157718) (922,0.05426471922) (926,0.06032290282) (930,0.0545285878) (934,0.06480857839) (938,0.05711093536) (942,0.07539849375) (946,0.06483292045) (950,0.05902500893) (954,0.06180254663) (958,0.06546466916) (962,0.06362009092) (966,0.06099001931) (970,0.0647530704) (974,0.07488399884) (978,0.08722092215) (982,0.08771826787) (986,0.0871751906) (990,0.07330925735) (994,0.0701100458) (998,0.07974274593)};
\addplot[pubNavy, opacity=0.20, line width=0.35pt] coordinates {(10,-0.0250836299) (14,-0.04667050378) (18,-0.05630518092) (22,-0.06332657589) (26,-0.05799305975) (30,-0.04926280549) (34,-0.03270662943) (38,-0.03900885079) (42,-0.03638206952) (46,-0.03310967267) (50,-0.0171267168) (54,-0.02461784224) (58,-0.0175706749) (62,-0.03204069819) (66,-0.008308188209) (70,-0.003387408097) (74,0.01592464973) (78,0.0171573418) (82,0.02078708049) (86,0.03657372166) (90,0.05283296594) (94,0.05046084746) (98,0.06720202646) (102,0.05940577398) (106,0.08404926125) (110,0.08884373115) (114,0.1014266835) (118,0.1065360411) (122,0.1155307737) (126,0.1171630135) (130,0.1237326352) (134,0.1022564648) (138,0.1070694988) (142,0.1080573756) (146,0.1072143678) (150,0.09349795298) (154,0.09963783565) (158,0.08264009286) (162,0.08288115126) (166,0.07168690657) (170,0.06032840597) (174,0.05756583476) (178,0.04983326009) (182,0.05999584984) (186,0.06782918028) (190,0.06526774656) (194,0.06960960309) (198,0.07111455217) (202,0.0798454722) (206,0.07963955081) (210,0.09389777686) (214,0.0894185538) (218,0.1051566271) (222,0.1032028569) (226,0.1105231982) (230,0.1048317177) (234,0.1013390929) (238,0.09439328431) (242,0.09057437978) (246,0.0943587452) (250,0.08901507274) (254,0.09894385219) (258,0.09007606088) (262,0.08499331393) (266,0.07432943601) (270,0.06881339481) (274,0.06976649333) (278,0.06993780275) (282,0.07285228395) (286,0.0729055532) (290,0.08615229231) (294,0.07365490733) (298,0.07385194594) (302,0.07637795576) (306,0.0808384051) (310,0.07595366911) (314,0.08129567265) (318,0.07892225093) (322,0.08048473765) (326,0.07524495852) (330,0.07449389603) (334,0.07541303765) (338,0.06901491442) (342,0.08145450937) (346,0.07126157799) (350,0.05601613189) (354,0.04567702433) (358,0.04273394152) (362,0.04793433095) (366,0.03969344443) (370,0.04118731864) (374,0.0519490137) (378,0.04628018365) (382,0.04053570988) (386,0.04468493256) (390,0.03912062789) (394,0.03409157134) (398,0.05635231845) (402,0.05113311565) (406,0.04719188306) (410,0.06161025876) (414,0.06821557108) (418,0.08001532611) (422,0.06645565175) (426,0.07545026803) (430,0.06388126685) (434,0.06682495041) (438,0.0720478076) (442,0.06704411094) (446,0.07992836121) (450,0.08533190867) (454,0.08163517211) (458,0.08693459968) (462,0.09902144933) (466,0.08733407791) (470,0.08487450004) (474,0.097249094) (478,0.08623932068) (482,0.0942685776) (486,0.08534408164) (490,0.08711810295) (494,0.08465508999) (498,0.08383358542) (502,0.08567598118) (506,0.08545088539) (510,0.08287516432) (514,0.07246081902) (518,0.06260238443) (522,0.07851871584) (526,0.07882555901) (530,0.07359700946) (534,0.08151462235) (538,0.09437745472) (542,0.07530045939) (546,0.07779768495) (550,0.08147891502) (554,0.09680321758) (558,0.07510754394) (562,0.05928117065) (566,0.06644453601) (570,0.06942841306) (574,0.06352852205) (578,0.06391990959) (582,0.05972694714) (586,0.0553125203) (590,0.05821885527) (594,0.05774792018) (598,0.06683777328) (602,0.06606591128) (606,0.05505395341) (610,0.08625113323) (614,0.09174912074) (618,0.09511533789) (622,0.0894076216) (626,0.09928989505) (630,0.09930846705) (634,0.09634564126) (638,0.09310765907) (642,0.1036966392) (646,0.09897506381) (650,0.08812112232) (654,0.08991169172) (658,0.09252116221) (662,0.1032222247) (666,0.09145891284) (670,0.09270706376) (674,0.09893899396) (678,0.0892810617) (682,0.09766141417) (686,0.09772237827) (690,0.103144329) (694,0.1051773131) (698,0.1136812092) (702,0.117961811) (706,0.1168583035) (710,0.1216250092) (714,0.1054680205) (718,0.1087033262) (722,0.1074563846) (726,0.1135785112) (730,0.1109243328) (734,0.1088472472) (738,0.115127755) (742,0.1183318812) (746,0.1027892699) (750,0.09772252588) (754,0.08409564525) (758,0.09223230755) (762,0.08846991394) (766,0.104674549) (770,0.1056589106) (774,0.1068810253) (778,0.09690476494) (782,0.1114160619) (786,0.1118877174) (790,0.09105252899) (794,0.08841069002) (798,0.08323969278) (802,0.09500241455) (806,0.09676999489) (810,0.08013228777) (814,0.08242880178) (818,0.0835453695) (822,0.08018130497) (826,0.09182059213) (830,0.09089757571) (834,0.08888293752) (838,0.09651393994) (842,0.1026987818) (846,0.09789680459) (850,0.1042922486) (854,0.1040612751) (858,0.09776031185) (862,0.11996754) (866,0.1102439437) (870,0.1097776793) (874,0.1135232125) (878,0.1046776307) (882,0.09441219225) (886,0.07914366924) (890,0.0843400419) (894,0.08671940694) (898,0.08086271586) (902,0.08823681009) (906,0.08776771284) (910,0.08374792034) (914,0.0823154816) (918,0.08814172786) (922,0.08061026684) (926,0.07646527024) (930,0.07972369889) (934,0.09454205628) (938,0.08071254188) (942,0.09159421702) (946,0.08379601293) (950,0.09454870001) (954,0.09474085265) (958,0.09220969326) (962,0.08713686197) (966,0.07318740389) (970,0.07180919393) (974,0.08775935168) (978,0.1084250787) (982,0.1018940998) (986,0.1012978118) (990,0.1005838457) (994,0.1044096546) (998,0.1145214879)};
\addplot[pubNavy, opacity=0.20, line width=0.35pt] coordinates {(10,0.01482941027) (14,-0.005419987388) (18,-0.02981738916) (22,-0.03315511178) (26,-0.02184004151) (30,-0.02879139351) (34,-0.01836868075) (38,-0.02197231299) (42,-0.01986022824) (46,-0.01425930805) (50,-0.004694730492) (54,-0.00428533025) (58,-0.002059956615) (62,-0.01139543794) (66,0.01124559438) (70,0.02231104211) (74,0.03948915469) (78,0.0423955167) (82,0.04732037076) (86,0.05164527714) (90,0.06109055506) (94,0.05917850779) (98,0.06571581712) (102,0.06201385988) (106,0.07986096864) (110,0.07727364189) (114,0.08168896447) (118,0.0844594154) (122,0.08860946489) (126,0.08576467416) (130,0.09480373428) (134,0.07805932537) (138,0.08847668854) (142,0.09751231514) (146,0.1022591316) (150,0.08578568401) (154,0.09846841768) (158,0.08090634376) (162,0.08190003785) (166,0.08448727936) (170,0.07426613339) (174,0.07494048766) (178,0.06561819911) (182,0.0731802528) (186,0.07097818714) (190,0.06810882878) (194,0.07246143495) (198,0.0754982858) (202,0.08262243851) (206,0.07195763671) (210,0.07332095842) (214,0.0624971923) (218,0.07187150181) (222,0.07142156199) (226,0.08076549012) (230,0.07297906549) (234,0.07888401432) (238,0.07499535543) (242,0.07468828807) (246,0.07328834836) (250,0.06974996118) (254,0.08174032432) (258,0.08804430444) (262,0.08151428348) (266,0.06843010583) (270,0.06432405196) (274,0.0618738862) (278,0.05751638011) (282,0.06683835135) (286,0.0739205991) (290,0.0778952618) (294,0.06591403418) (298,0.07160310624) (302,0.07762357284) (306,0.07787915439) (310,0.07286260731) (314,0.07602950619) (318,0.08168181383) (322,0.08262904076) (326,0.08354334357) (330,0.08408994653) (334,0.08200208707) (338,0.08060803835) (342,0.09547371199) (346,0.07849914963) (350,0.07816115412) (354,0.05782714465) (358,0.05586433988) (362,0.06665767488) (366,0.06625002154) (370,0.07539654941) (374,0.08665858784) (378,0.08185958788) (382,0.07436733215) (386,0.06904954471) (390,0.05933898566) (394,0.04406213684) (398,0.05415135874) (402,0.0514361659) (406,0.04644509808) (410,0.06185724673) (414,0.06797477136) (418,0.07463089509) (422,0.05749300088) (426,0.06255789103) (430,0.06359136756) (434,0.07834259776) (438,0.0792508944) (442,0.07539986118) (446,0.09981823195) (450,0.1085372756) (454,0.106332768) (458,0.1135880712) (462,0.1149134866) (466,0.1032600091) (470,0.09435195298) (474,0.09927196009) (478,0.08942301519) (482,0.08349158648) (486,0.0772224932) (490,0.08667145013) (494,0.08614130065) (498,0.07344789394) (502,0.07844004932) (506,0.08192966884) (510,0.078180135) (514,0.06720783987) (518,0.06365813257) (522,0.07846059252) (526,0.07732534642) (530,0.07826099507) (534,0.08153115909) (538,0.09313470297) (542,0.08200730161) (546,0.07561696) (550,0.08244190533) (554,0.08895270803) (558,0.0845793758) (562,0.08397358727) (566,0.09061749289) (570,0.09331503982) (574,0.09481685459) (578,0.09367098607) (582,0.09248399633) (586,0.09134202165) (590,0.07967361733) (594,0.08580570583) (598,0.09483347006) (602,0.0832106996) (606,0.0736706927) (610,0.08677768585) (614,0.09041778445) (618,0.09662585125) (622,0.08733341549) (626,0.08852337615) (630,0.09216498742) (634,0.08658158921) (638,0.08440953173) (642,0.0865581904) (646,0.08041728329) (650,0.08530025345) (654,0.09944649104) (658,0.1123703788) (662,0.116135065) (666,0.1006681217) (670,0.1015580911) (674,0.1012847399) (678,0.1037953042) (682,0.120213472) (686,0.122255248) (690,0.1197208742) (694,0.1217124136) (698,0.1093018183) (702,0.1039299064) (706,0.09561589816) (710,0.1029781851) (714,0.08820919339) (718,0.09823489086) (722,0.09688196834) (726,0.0927090382) (730,0.08319605434) (734,0.08108442191) (738,0.08606917714) (742,0.0806766699) (746,0.07466500323) (750,0.08658149598) (754,0.07735774886) (758,0.07623682732) (762,0.07716365655) (766,0.0904493528) (770,0.08832884156) (774,0.09010878364) (778,0.09443207008) (782,0.09071412722) (786,0.09815990467) (790,0.09730440318) (794,0.1005911533) (798,0.1003278496) (802,0.1166005023) (806,0.1207828149) (810,0.1213659899) (814,0.1184379746) (818,0.1223741317) (822,0.122767572) (826,0.1175186631) (830,0.12139099) (834,0.1281277004) (838,0.1335182871) (842,0.1370965492) (846,0.1343034609) (850,0.1239390651) (854,0.1152388214) (858,0.1067351256) (862,0.1128520282) (866,0.09806226298) (870,0.1027325159) (874,0.0951178044) (878,0.08753686127) (882,0.07741049867) (886,0.06313396981) (890,0.06088469629) (894,0.06284324225) (898,0.07245493796) (902,0.08022152249) (906,0.08643421579) (910,0.07453659999) (914,0.07195165258) (918,0.08378391675) (922,0.07512171426) (926,0.08064276257) (930,0.0873257153) (934,0.08680018874) (938,0.07055372073) (942,0.07834126657) (946,0.06914747852) (950,0.0745841685) (954,0.07606370318) (958,0.07820475463) (962,0.06770554051) (966,0.06571618046) (970,0.06391585782) (974,0.08390914055) (978,0.09782645767) (982,0.09210931813) (986,0.0928585556) (990,0.08841082148) (994,0.07792223372) (998,0.0874775808)};
\addplot[pubNavy, opacity=0.20, line width=0.35pt] coordinates {(10,-0.001975990707) (14,-0.03258775822) (18,-0.05051177325) (22,-0.04131407204) (26,-0.02703304783) (30,-0.03400659787) (34,-0.01889828918) (38,-0.01575436365) (42,-0.02141637315) (46,-0.02104839421) (50,-0.01016939919) (54,-0.01881284312) (58,-0.01880449281) (62,-0.03695922788) (66,-0.008769878395) (70,-0.007853546255) (74,0.001627212328) (78,0.00254658354) (82,0.004605179193) (86,0.01066780484) (90,0.02407346094) (94,0.0206201648) (98,0.03137597677) (102,0.02168596151) (106,0.04003880293) (110,0.04165045755) (114,0.05027492786) (118,0.0478901596) (122,0.05375238821) (126,0.05150071809) (130,0.05961757842) (134,0.03364285011) (138,0.03436529601) (142,0.02908727701) (146,0.03501838145) (150,0.02475124479) (154,0.04208128313) (158,0.01871186861) (162,0.008035558263) (166,0.01476387339) (170,0.001116250827) (174,0.001775012734) (178,-0.008946990321) (182,0.005479719044) (186,0.01548227119) (190,0.01272862116) (194,0.0226194053) (198,0.02166835796) (202,0.01842072289) (206,0.002734662135) (210,0.01620539347) (214,0.008443699827) (218,0.01984999406) (222,0.02203254155) (226,0.02761383217) (230,0.02413205341) (234,0.02239120161) (238,0.01843339544) (242,0.01440177184) (246,0.02469328043) (250,0.02602014114) (254,0.0433360835) (258,0.053010842) (262,0.04449959737) (266,0.03010995532) (270,0.02221133229) (274,0.02858199583) (278,0.02636924839) (282,0.03620864347) (286,0.02963552931) (290,0.03193418129) (294,0.01519058809) (298,0.008035528653) (302,0.002040755507) (306,-0.001544764046) (310,-0.005277054768) (314,0.003886933885) (318,0.009677138189) (322,0.009909203804) (326,0.003967459567) (330,0.005705653117) (334,0.009937107703) (338,0.0171185293) (342,0.03432106508) (346,0.03012665668) (350,0.01932279718) (354,0.02151611433) (358,0.02234947097) (362,0.02975241634) (366,0.02742681185) (370,0.03144268714) (374,0.03862200412) (378,0.03975670955) (382,0.03281214599) (386,0.02989714358) (390,0.02335004203) (394,0.01046031178) (398,0.02974358609) (402,0.0219463851) (406,0.01423684291) (410,0.0369395382) (414,0.04447012735) (418,0.05445235775) (422,0.04418440013) (426,0.05352242906) (430,0.04878774467) (434,0.05765243648) (438,0.05842656991) (442,0.05291020314) (446,0.0704690151) (450,0.07171125672) (454,0.06554567718) (458,0.07698800311) (462,0.08028243709) (466,0.07371617738) (470,0.05994550237) (474,0.05981604608) (478,0.04797632218) (482,0.05466691698) (486,0.04077532495) (490,0.03624075375) (494,0.03662146121) (498,0.03148199136) (502,0.04054846389) (506,0.03584451009) (510,0.02711646415) (514,0.01188027356) (518,0.006234718285) (522,0.02821295432) (526,0.0310967295) (530,0.0311036649) (534,0.03566419773) (538,0.06400794365) (542,0.04499479402) (546,0.04019426064) (550,0.04719716183) (554,0.06745725775) (558,0.05714660083) (562,0.05155496685) (566,0.06390513233) (570,0.05304208457) (574,0.05078810839) (578,0.04920282151) (582,0.05180754852) (586,0.04367226394) (590,0.0465911743) (594,0.0601589235) (598,0.07399709183) (602,0.07379675161) (606,0.06113127962) (610,0.0865992594) (614,0.08804834212) (618,0.08848783794) (622,0.09543542915) (626,0.09745487174) (630,0.09273873965) (634,0.09119218464) (638,0.08553820267) (642,0.08457795303) (646,0.0745157288) (650,0.06901566649) (654,0.07565950373) (658,0.08381544864) (662,0.09156927043) (666,0.07575123341) (670,0.08099019884) (674,0.08687256206) (678,0.08095000398) (682,0.08874078291) (686,0.102618247) (690,0.09933300455) (694,0.1001491038) (698,0.1025344984) (702,0.09244470936) (706,0.07584255449) (710,0.07905627479) (714,0.04749132048) (718,0.06432550879) (722,0.06536360976) (726,0.06283135942) (730,0.05829351371) (734,0.05976596666) (738,0.05708636671) (742,0.05671081155) (746,0.04840789711) (750,0.04432766886) (754,0.04168811376) (758,0.0476528837) (762,0.04816935334) (766,0.061673045) (770,0.06207313031) (774,0.05460635797) (778,0.0469155634) (782,0.04480299613) (786,0.04858769381) (790,0.04881147492) (794,0.04599030716) (798,0.05062425851) (802,0.07559206929) (806,0.07996815952) (810,0.079235464) (814,0.08270937382) (818,0.08312667007) (822,0.08653146895) (826,0.08506715765) (830,0.08763656364) (834,0.1000969932) (838,0.1051111555) (842,0.1010721501) (846,0.1025515869) (850,0.09714797162) (854,0.09235136745) (858,0.07529377229) (862,0.09332502071) (866,0.07341136603) (870,0.08091343505) (874,0.07606576082) (878,0.07132007546) (882,0.06204690086) (886,0.04204277724) (890,0.04846318849) (894,0.05593837972) (898,0.06119665693) (902,0.06719762827) (906,0.0659772906) (910,0.06132636199) (914,0.06299552859) (918,0.0683698346) (922,0.05567292673) (926,0.06550124071) (930,0.06656111394) (934,0.06830278519) (938,0.04991464378) (942,0.0620565173) (946,0.04479372338) (950,0.04751297968) (954,0.04815678751) (958,0.05121382475) (962,0.0471296754) (966,0.0392511955) (970,0.03636744816) (974,0.05415874192) (978,0.06592275562) (982,0.05682763745) (986,0.06528233077) (990,0.06598474108) (994,0.05685732241) (998,0.0655786552)};
\addplot[pubNavy, opacity=0.20, line width=0.35pt] coordinates {(10,-0.009943272238) (14,-0.02779827897) (18,-0.03527976794) (22,-0.03917704767) (26,-0.0252302275) (30,-0.02467719735) (34,-0.00967468015) (38,-0.02623649841) (42,-0.02871029524) (46,-0.03372682538) (50,-0.02440580802) (54,-0.03203274541) (58,-0.02779888979) (62,-0.03240667553) (66,-0.01859522532) (70,-0.02408052799) (74,-0.01394764576) (78,-0.01238208766) (82,-0.01146817898) (86,0.004876881054) (90,0.02858265711) (94,0.02627101201) (98,0.04558376694) (102,0.02929788014) (106,0.04375782926) (110,0.04127678317) (114,0.04526794396) (118,0.05199136719) (122,0.04998651525) (126,0.04795374211) (130,0.05475317462) (134,0.02547027287) (138,0.0210454076) (142,0.01706292381) (146,0.02529562334) (150,0.01671494429) (154,0.03012209156) (158,0.02910251404) (162,0.03248699769) (166,0.03234061677) (170,0.02819030658) (174,0.03653882519) (178,0.03531962349) (182,0.05124938195) (186,0.05908285931) (190,0.05538793651) (194,0.06357991701) (198,0.06472095231) (202,0.06688750817) (206,0.04700407411) (210,0.04960056357) (214,0.05072360776) (218,0.05403928423) (222,0.05338976031) (226,0.05578701495) (230,0.05177348421) (234,0.04966896474) (238,0.05247120844) (242,0.04841928111) (246,0.05369936683) (250,0.05029103796) (254,0.06331200899) (258,0.07207180404) (262,0.0614700214) (266,0.05252913259) (270,0.0439986206) (274,0.0432434107) (278,0.04327916488) (282,0.04899768841) (286,0.04437388677) (290,0.04926132863) (294,0.04590846144) (298,0.04134198767) (302,0.04234095555) (306,0.04458987991) (310,0.04335876959) (314,0.05132911237) (318,0.06179645772) (322,0.06154748131) (326,0.06482062804) (330,0.0665436729) (334,0.0632567235) (338,0.05507858621) (342,0.06298128925) (346,0.06115647933) (350,0.05526043262) (354,0.03325757525) (358,0.03511157523) (362,0.03282775443) (366,0.02858774445) (370,0.02332376445) (374,0.03058807816) (378,0.03650768486) (382,0.036188633) (386,0.03081263396) (390,0.02945882722) (394,0.02209162768) (398,0.03446394592) (402,0.029401867) (406,0.02528438865) (410,0.02379204914) (414,0.03657405001) (418,0.04532252167) (422,0.03678145482) (426,0.04074661233) (430,0.03455236463) (434,0.03837183575) (438,0.03945309486) (442,0.03278019853) (446,0.04442724805) (450,0.05155244933) (454,0.04760347216) (458,0.0561851253) (462,0.07249787819) (466,0.06245900693) (470,0.05574764459) (474,0.06313568944) (478,0.05090689343) (482,0.05482227792) (486,0.04904371802) (490,0.05465652976) (494,0.0591346344) (498,0.05347872867) (502,0.05542183555) (506,0.05908524534) (510,0.05617010294) (514,0.03959577745) (518,0.03633058191) (522,0.03705428489) (526,0.03166128916) (530,0.04070021992) (534,0.04101439488) (538,0.05996883018) (542,0.033304745) (546,0.02296306443) (550,0.0220114794) (554,0.02679833013) (558,0.02586434493) (562,0.0414507249) (566,0.05048688679) (570,0.05235555411) (574,0.06215056225) (578,0.04484229658) (582,0.03938843427) (586,0.03473067705) (590,0.03358995414) (594,0.03405692772) (598,0.05721040548) (602,0.04800692725) (606,0.03848416838) (610,0.04226710928) (614,0.03663790644) (618,0.03160013448) (622,0.02825513581) (626,0.03869623435) (630,0.04560521564) (634,0.044012643) (638,0.04309044792) (642,0.04741392107) (646,0.03739131421) (650,0.0368139658) (654,0.04498955343) (658,0.04830100293) (662,0.06149666299) (666,0.05592537471) (670,0.06485527608) (674,0.06680850265) (678,0.06091353853) (682,0.07519667488) (686,0.08800911747) (690,0.08341512613) (694,0.1017103386) (698,0.08759721141) (702,0.08227478998) (706,0.08571046532) (710,0.08306515894) (714,0.06397738217) (718,0.07281773611) (722,0.07204542535) (726,0.0687774522) (730,0.06012566049) (734,0.05225451089) (738,0.04355339911) (742,0.03410340036) (746,0.02611127885) (750,0.0347588907) (754,0.03449302025) (758,0.03751723431) (762,0.04042451272) (766,0.05506780523) (770,0.0537947491) (774,0.05154892607) (778,0.0502930016) (782,0.04931717763) (786,0.07053496348) (790,0.07314914817) (794,0.08306008239) (798,0.08218700766) (802,0.09020109566) (806,0.09045004186) (810,0.08791937432) (814,0.09094179046) (818,0.09273177971) (822,0.09132454319) (826,0.08712672659) (830,0.09779679877) (834,0.1020521778) (838,0.09556550303) (842,0.09501471191) (846,0.09853212323) (850,0.09693307072) (854,0.09624389922) (858,0.09693874957) (862,0.106069432) (866,0.096818501) (870,0.1068300474) (874,0.1031138926) (878,0.1008425669) (882,0.08335955524) (886,0.07292095427) (890,0.08501371235) (894,0.07947935051) (898,0.08035928765) (902,0.08600470177) (906,0.07738969351) (910,0.07507504208) (914,0.07450036049) (918,0.07264238622) (922,0.07372056786) (926,0.08243494796) (930,0.09199138909) (934,0.09440915203) (938,0.0754433078) (942,0.08002171355) (946,0.06775005382) (950,0.06551579964) (954,0.07437572201) (958,0.07565620886) (962,0.06942215388) (966,0.06304291978) (970,0.06797828269) (974,0.07615931858) (978,0.08109624421) (982,0.07532656119) (986,0.07525517282) (990,0.0873725231) (994,0.09277413428) (998,0.1031329415)};
\addplot[pubNavy, opacity=0.20, line width=0.35pt] coordinates {(10,-0.05610656765) (14,-0.05957737966) (18,-0.05367991201) (22,-0.05551812205) (26,-0.04046799759) (30,-0.04523184439) (34,-0.03098095493) (38,-0.03760927815) (42,-0.03821984656) (46,-0.03047549381) (50,-0.01801435475) (54,-0.01169419091) (58,-0.00678968659) (62,-0.01389433795) (66,0.006487276639) (70,0.007743963231) (74,0.02219730394) (78,0.02594806504) (82,0.03036529882) (86,0.04113386713) (90,0.0588463419) (94,0.05239732493) (98,0.06302504914) (102,0.05824193982) (106,0.07523028692) (110,0.07362697811) (114,0.07870716696) (118,0.08478767089) (122,0.09067003618) (126,0.09449337723) (130,0.1020883357) (134,0.09334817721) (138,0.09023902509) (142,0.09624417122) (146,0.09894647341) (150,0.09228590303) (154,0.09423826213) (158,0.08673039945) (162,0.0913653272) (166,0.08802743637) (170,0.08438530552) (174,0.09004124501) (178,0.08521791447) (182,0.08916515352) (186,0.09135509115) (190,0.09317125062) (194,0.1004500228) (198,0.1032017767) (202,0.1144684559) (206,0.09799508151) (210,0.1002190778) (214,0.09183757955) (218,0.1000746881) (222,0.09567497816) (226,0.09962668692) (230,0.09434252312) (234,0.09567156013) (238,0.0861460555) (242,0.08170161742) (246,0.08294513712) (250,0.08037807418) (254,0.0878523689) (258,0.09869020335) (262,0.09522488338) (266,0.08705685168) (270,0.08751370375) (274,0.08833036339) (278,0.08998890446) (282,0.09179884903) (286,0.1019135942) (290,0.1106527152) (294,0.09107541877) (298,0.08797844177) (302,0.09112668182) (306,0.09238444016) (310,0.08461559191) (314,0.07548720604) (318,0.07256792597) (322,0.06708584494) (326,0.06747914041) (330,0.06737925895) (334,0.07297402532) (338,0.07172908302) (342,0.08603504284) (346,0.07778971685) (350,0.06976708963) (354,0.05253619792) (358,0.05850885211) (362,0.06617598327) (366,0.07014319611) (370,0.08197400457) (374,0.08005060465) (378,0.07341523044) (382,0.06139597683) (386,0.05810097072) (390,0.05225340062) (394,0.05129488105) (398,0.07050263955) (402,0.05733282982) (406,0.06441257799) (410,0.07826893987) (414,0.08494650628) (418,0.09113618026) (422,0.07155107132) (426,0.08142920159) (430,0.07264156006) (434,0.07264195014) (438,0.07728519186) (442,0.06887223977) (446,0.08041079367) (450,0.08211271218) (454,0.07281323252) (458,0.06888880732) (462,0.07060200624) (466,0.06617000964) (470,0.06370437782) (474,0.07857143469) (478,0.07188534238) (482,0.08304696659) (486,0.07440886394) (490,0.07414192784) (494,0.07775793772) (498,0.08381915607) (502,0.08569643051) (506,0.1004551437) (510,0.1009029912) (514,0.08354098882) (518,0.07503469055) (522,0.08606695557) (526,0.08924965973) (530,0.08094752421) (534,0.08369833359) (538,0.09855231903) (542,0.08056151266) (546,0.08171333159) (550,0.08493119667) (554,0.08910592922) (558,0.0841223387) (562,0.08249524256) (566,0.09703858934) (570,0.09516073583) (574,0.09543448842) (578,0.1064145182) (582,0.1076961453) (586,0.101263553) (590,0.105280629) (594,0.09942393374) (598,0.1138326679) (602,0.1056601212) (606,0.09890640385) (610,0.1130530849) (614,0.119677733) (618,0.1167960456) (622,0.1187941686) (626,0.1266305054) (630,0.1208412559) (634,0.1192765969) (638,0.1169409946) (642,0.1193004873) (646,0.1185712053) (650,0.1100489074) (654,0.1121343747) (658,0.1154014365) (662,0.1147859162) (666,0.1094023947) (670,0.1089184504) (674,0.102343101) (678,0.09511394126) (682,0.09772216456) (686,0.09649697535) (690,0.1018670957) (694,0.09856781842) (698,0.09461834216) (702,0.0982053849) (706,0.09545145828) (710,0.09865931217) (714,0.09035863046) (718,0.09327683047) (722,0.0898134098) (726,0.09391457438) (730,0.09452260436) (734,0.1023271125) (738,0.1035436404) (742,0.1035925392) (746,0.1009558735) (750,0.1026912457) (754,0.0930530336) (758,0.0904845274) (762,0.09030517266) (766,0.09456009029) (770,0.09911984436) (774,0.09355637268) (778,0.09705898438) (782,0.09344459937) (786,0.091397737) (790,0.09255844398) (794,0.09106272819) (798,0.08376819465) (802,0.09986123865) (806,0.1026309224) (810,0.1121074015) (814,0.1227571299) (818,0.124290609) (822,0.1221234916) (826,0.1291310927) (830,0.13357987) (834,0.1406446236) (838,0.1389504102) (842,0.1499138755) (846,0.1518696363) (850,0.1450931721) (854,0.1486088976) (858,0.1455800432) (862,0.1465429317) (866,0.1315232116) (870,0.1298566058) (874,0.1222104304) (878,0.1127682225) (882,0.1023855844) (886,0.09140477011) (890,0.09592922276) (894,0.09055302151) (898,0.08998087827) (902,0.09874190541) (906,0.08563292794) (910,0.07934625179) (914,0.08102246382) (918,0.09387615612) (922,0.08474107541) (926,0.07861397865) (930,0.07675987437) (934,0.07260210966) (938,0.06771386155) (942,0.06971184664) (946,0.07088848818) (950,0.07938387308) (954,0.08114172154) (958,0.08323690277) (962,0.07505137305) (966,0.06014504583) (970,0.06166666697) (974,0.08514254193) (978,0.1051558579) (982,0.1043018902) (986,0.103491212) (990,0.09604333726) (994,0.09626122933) (998,0.1037265875)};
\addplot[pubNavy, opacity=0.20, line width=0.35pt] coordinates {(10,-0.02215112034) (14,-0.0257060458) (18,-0.01628770949) (22,-0.02569563651) (26,-0.01245031725) (30,-0.01420175341) (34,-0.005163862811) (38,-0.01462995318) (42,-0.01648495471) (46,-0.01044592183) (50,-0.001316307359) (54,-0.002246507634) (58,0.003208447975) (62,-0.01107205121) (66,0.006792913874) (70,0.005579509083) (74,0.01648083694) (78,0.0152251076) (82,0.01225294634) (86,0.02048256898) (90,0.0410084061) (94,0.03549616072) (98,0.03582217701) (102,0.02049413234) (106,0.03422526346) (110,0.0332889952) (114,0.04165632011) (118,0.04325237626) (122,0.04512533988) (126,0.04693677872) (130,0.06092232844) (134,0.03202135306) (138,0.03417632721) (142,0.02998432481) (146,0.04503087837) (150,0.03873593157) (154,0.05970378024) (158,0.04647516442) (162,0.04224998571) (166,0.04056517551) (170,0.03731812865) (174,0.0383418055) (178,0.02960719931) (182,0.04445441281) (186,0.04935778949) (190,0.04415351771) (194,0.05074731801) (198,0.05173741593) (202,0.05503834329) (206,0.04767374386) (210,0.0547243329) (214,0.0497209526) (218,0.05594139192) (222,0.05634947278) (226,0.06665027213) (230,0.06177428906) (234,0.06288957568) (238,0.06058806361) (242,0.06026276533) (246,0.06167717913) (250,0.0610879633) (254,0.07463154678) (258,0.07745540369) (262,0.07802136936) (266,0.06695942864) (270,0.06373535705) (274,0.06800907742) (278,0.06784506193) (282,0.07328192634) (286,0.07953183678) (290,0.09659285966) (294,0.0827342624) (298,0.08091364868) (302,0.08239287532) (306,0.0819811168) (310,0.07120874211) (314,0.07905891211) (318,0.07805552289) (322,0.07258019732) (326,0.06907633354) (330,0.0681157088) (334,0.06483931674) (338,0.0600960591) (342,0.07443452721) (346,0.05815435344) (350,0.05801552205) (354,0.03326955388) (358,0.03911504591) (362,0.04529894013) (366,0.04601301857) (370,0.05554658574) (374,0.05825167655) (378,0.06402308503) (382,0.06925573907) (386,0.0678902349) (390,0.06155130373) (394,0.06399919291) (398,0.07834708421) (402,0.06443916723) (406,0.06145560916) (410,0.07669774146) (414,0.08664721291) (418,0.0943873754) (422,0.08623334854) (426,0.09948775897) (430,0.08868643731) (434,0.08802030235) (438,0.09146941548) (442,0.08176386063) (446,0.09255350332) (450,0.09706845689) (454,0.1011699453) (458,0.1129861332) (462,0.1220944259) (466,0.1151866083) (470,0.09834590347) (474,0.1031291447) (478,0.09967889267) (482,0.1105599397) (486,0.1038979536) (490,0.1091045649) (494,0.1083215698) (498,0.1092924099) (502,0.1114086545) (506,0.1106729799) (510,0.1019643604) (514,0.08764208801) (518,0.08145321309) (522,0.09372740636) (526,0.09759088801) (530,0.09220592954) (534,0.09296325088) (538,0.1133803765) (542,0.07685510392) (546,0.06683448796) (550,0.06787857157) (554,0.07438829048) (558,0.07041197107) (562,0.05480330032) (566,0.07365498095) (570,0.07305725424) (574,0.07098352987) (578,0.06220567839) (582,0.06008075066) (586,0.05057448206) (590,0.05690885987) (594,0.06596750411) (598,0.08872066535) (602,0.09184797823) (606,0.08157920771) (610,0.09311253416) (614,0.09877584975) (618,0.09684127696) (622,0.09603394274) (626,0.1009100222) (630,0.1069630665) (634,0.1078412918) (638,0.1045117359) (642,0.1092975086) (646,0.09999642574) (650,0.08464822059) (654,0.09106484624) (658,0.09898114754) (662,0.1105653225) (666,0.09418016894) (670,0.1009635496) (674,0.1073259321) (678,0.09036405222) (682,0.1000214251) (686,0.1048311221) (690,0.1014305825) (694,0.1063523439) (698,0.1128762889) (702,0.110348949) (706,0.1032940099) (710,0.1075869481) (714,0.09182592922) (718,0.1060175257) (722,0.09463141094) (726,0.09654304356) (730,0.09296189841) (734,0.08846708712) (738,0.08753072303) (742,0.08359562006) (746,0.06908907582) (750,0.05860846552) (754,0.046052397) (758,0.044244123) (762,0.02976860227) (766,0.03951307304) (770,0.03920487924) (774,0.03691797966) (778,0.03192451114) (782,0.03521323924) (786,0.04491272244) (790,0.04355856511) (794,0.04722450608) (798,0.04983328168) (802,0.07228214897) (806,0.07521745847) (810,0.07775153949) (814,0.0832648614) (818,0.09003543645) (822,0.08656057129) (826,0.08499732393) (830,0.09049229202) (834,0.09437674664) (838,0.1038890629) (842,0.105791246) (846,0.1033298832) (850,0.09953077588) (854,0.09883938926) (858,0.1002085027) (862,0.1092904145) (866,0.0995085273) (870,0.09916079249) (874,0.09735047337) (878,0.09263896812) (882,0.08109051883) (886,0.05984263706) (890,0.06565436788) (894,0.06920388115) (898,0.0745601236) (902,0.08053973347) (906,0.07337118475) (910,0.06735860601) (914,0.07184489051) (918,0.07222542378) (922,0.06907181462) (926,0.08113101216) (930,0.09082470384) (934,0.09748337344) (938,0.07917519526) (942,0.08682260922) (946,0.08172096367) (950,0.08402729198) (954,0.0888740935) (958,0.08864036248) (962,0.07709343305) (966,0.07312963999) (970,0.07135326293) (974,0.08945148857) (978,0.09652834952) (982,0.08980176702) (986,0.09929260364) (990,0.1032258686) (994,0.0965686813) (998,0.1048431922)};
\addplot[pubNavy, opacity=0.20, line width=0.35pt] coordinates {(10,-0.05466230349) (14,-0.05463294018) (18,-0.06723012173) (22,-0.05640613607) (26,-0.04873749021) (30,-0.0503554747) (34,-0.03649167141) (38,-0.03995238038) (42,-0.04348089863) (46,-0.03951810082) (50,-0.02666004972) (54,-0.02388080229) (58,-0.01617748178) (62,-0.02393205661) (66,-0.0001812830186) (70,0.004734919703) (74,0.01981539506) (78,0.0276206153) (82,0.03036706084) (86,0.03688163344) (90,0.0535916184) (94,0.05705400977) (98,0.06779407783) (102,0.0546881531) (106,0.07242416943) (110,0.07200684561) (114,0.07684235223) (118,0.08002541217) (122,0.07423252315) (126,0.07090151376) (130,0.08187195247) (134,0.05708053416) (138,0.05513635717) (142,0.0534277201) (146,0.05705108178) (150,0.03984196723) (154,0.0452936366) (158,0.03852041877) (162,0.03756179747) (166,0.04067733768) (170,0.02898691439) (174,0.03490721427) (178,0.02779908676) (182,0.0390209729) (186,0.05308493421) (190,0.04317958671) (194,0.05164333092) (198,0.0526665727) (202,0.06066664619) (206,0.04209434058) (210,0.04128520676) (214,0.03028398799) (218,0.03030611542) (222,0.03438932567) (226,0.04135097448) (230,0.03613423931) (234,0.03311294679) (238,0.03618149144) (242,0.03237379218) (246,0.04043472953) (250,0.03980379808) (254,0.05575004956) (258,0.07380131621) (262,0.06393982506) (266,0.05882558359) (270,0.05276035398) (274,0.05600864719) (278,0.05679445306) (282,0.05864840665) (286,0.05255131519) (290,0.05352757255) (294,0.039674477) (298,0.02859853923) (302,0.03123146681) (306,0.01657386281) (310,0.01167510819) (314,0.01790810229) (318,0.0194876737) (322,0.02190599004) (326,0.01699840632) (330,0.01706240552) (334,0.02616911572) (338,0.02252935302) (342,0.03601801588) (346,0.04180003935) (350,0.03698677975) (354,0.01770709154) (358,0.01270010204) (362,0.01820478057) (366,0.01847741365) (370,0.01728070664) (374,0.02548311413) (378,0.02462574703) (382,0.01838487144) (386,0.01505174106) (390,0.0107359927) (394,0.006914661555) (398,0.02169175721) (402,0.01443635198) (406,0.02063625224) (410,0.03859860145) (414,0.04619127992) (418,0.05543606944) (422,0.04830640711) (426,0.05908380483) (430,0.0582541342) (434,0.06879574596) (438,0.06746590639) (442,0.06615386243) (446,0.08312086615) (450,0.09384483315) (454,0.1000092412) (458,0.1015620731) (462,0.1087985371) (466,0.09386737429) (470,0.09163000824) (474,0.1000162279) (478,0.09139319644) (482,0.09822663932) (486,0.09671572829) (490,0.1003891632) (494,0.1005205701) (498,0.09289488929) (502,0.09814234964) (506,0.1063637079) (510,0.1027215466) (514,0.08887227597) (518,0.08584082693) (522,0.1006973368) (526,0.09368307235) (530,0.09971150136) (534,0.0974577839) (538,0.1084783367) (542,0.1055619411) (546,0.09882030866) (550,0.1010718552) (554,0.1020753613) (558,0.09655313832) (562,0.09588332252) (566,0.1025272609) (570,0.09827853606) (574,0.09214355916) (578,0.08247309029) (582,0.08044514393) (586,0.07594225665) (590,0.06040625893) (594,0.06339284576) (598,0.07989688474) (602,0.07327067741) (606,0.06639975649) (610,0.08034695153) (614,0.08482142884) (618,0.08686403927) (622,0.08732800983) (626,0.1018694343) (630,0.1003007188) (634,0.09669250456) (638,0.08810755235) (642,0.09486816429) (646,0.09514422075) (650,0.08954606531) (654,0.0971876717) (658,0.1027277911) (662,0.1048610296) (666,0.09580209029) (670,0.09888531782) (674,0.09882257173) (678,0.1032283915) (682,0.1175075355) (686,0.1239751275) (690,0.1273851005) (694,0.1210955446) (698,0.1153236097) (702,0.1118485912) (706,0.1059897132) (710,0.1138152394) (714,0.09628543947) (718,0.09889620853) (722,0.1029487951) (726,0.09639774391) (730,0.087378163) (734,0.07903201652) (738,0.07685281245) (742,0.08527996289) (746,0.0756980627) (750,0.07650522815) (754,0.06406418591) (758,0.06588223184) (762,0.0633629564) (766,0.07754156058) (770,0.07725714822) (774,0.07253705519) (778,0.06651780888) (782,0.07816960198) (786,0.08796005082) (790,0.08628541327) (794,0.08404360353) (798,0.0806131439) (802,0.09735656425) (806,0.09706705065) (810,0.1040859836) (814,0.1074471008) (818,0.1130214356) (822,0.1212513082) (826,0.1167683793) (830,0.1204031829) (834,0.1298972309) (838,0.126750242) (842,0.1310957076) (846,0.1338300924) (850,0.1286684567) (854,0.1226835504) (858,0.1189217912) (862,0.1264072578) (866,0.10577115) (870,0.1078903965) (874,0.09802492656) (878,0.09150192339) (882,0.07735097367) (886,0.05796614057) (890,0.05972490746) (894,0.06051115938) (898,0.06302380538) (902,0.07285879889) (906,0.07104216351) (910,0.06510530775) (914,0.06986669467) (918,0.07426867149) (922,0.071909662) (926,0.0733661921) (930,0.07699004414) (934,0.07759905386) (938,0.06556847858) (942,0.07690359607) (946,0.06444070159) (950,0.06782416994) (954,0.07080553711) (958,0.07132282947) (962,0.06834165392) (966,0.0637392044) (970,0.06253023257) (974,0.08123744716) (978,0.09747986294) (982,0.09131248629) (986,0.1022283633) (990,0.1014188116) (994,0.1051527326) (998,0.112310503)};
\addplot[pubNavy, opacity=0.20, line width=0.35pt] coordinates {(10,-0.01972277419) (14,-0.02078775047) (18,-0.02266276483) (22,-0.02437577018) (26,-0.01077566094) (30,-0.01930418332) (34,-0.01021309029) (38,-0.01733368095) (42,-0.02492625654) (46,-0.02614400668) (50,-0.01919855539) (54,-0.01959001268) (58,-0.01517829037) (62,-0.02748608657) (66,-0.009869770078) (70,-0.01473564472) (74,-0.001226155445) (78,0.002508500936) (82,0.003702109728) (86,0.01322603969) (90,0.03333219302) (94,0.03214973401) (98,0.03802088973) (102,0.03248735492) (106,0.0488799055) (110,0.04348433409) (114,0.05467315136) (118,0.06016405403) (122,0.06492166105) (126,0.06421371885) (130,0.0731108207) (134,0.0484089626) (138,0.05120046699) (142,0.05207605114) (146,0.05605668978) (150,0.04705586011) (154,0.05795914062) (158,0.03830128004) (162,0.03562984532) (166,0.03880497624) (170,0.02779477217) (174,0.03112754961) (178,0.02415183387) (182,0.0321749721) (186,0.03844578574) (190,0.03557802441) (194,0.04294664655) (198,0.05185226878) (202,0.05734761655) (206,0.05250000531) (210,0.06028099034) (214,0.05051796181) (218,0.05776200177) (222,0.0578269402) (226,0.06907184031) (230,0.06245670774) (234,0.06950206913) (238,0.06596615288) (242,0.05995582376) (246,0.06073919463) (250,0.05574976197) (254,0.07084705407) (258,0.07411829176) (262,0.0681185334) (266,0.04748717418) (270,0.03919429293) (274,0.03900238208) (278,0.03921483326) (282,0.04502861229) (286,0.04700483767) (290,0.04798135866) (294,0.04006736738) (298,0.03627406878) (302,0.04251819129) (306,0.04021141503) (310,0.03617164606) (314,0.03908736711) (318,0.04645382674) (322,0.04852405784) (326,0.04759509695) (330,0.05307524639) (334,0.05547457341) (338,0.06058926324) (342,0.07063233819) (346,0.0562081134) (350,0.05838719655) (354,0.03690604859) (358,0.04466448869) (362,0.05335015188) (366,0.06217113442) (370,0.06339767292) (374,0.06929116987) (378,0.06931798273) (382,0.06528306378) (386,0.0568759725) (390,0.05068959112) (394,0.04558320439) (398,0.05161161577) (402,0.04724215264) (406,0.03814105811) (410,0.04830195051) (414,0.05234443749) (418,0.05810324131) (422,0.05360111761) (426,0.05511120969) (430,0.0527751098) (434,0.05639352375) (438,0.06053756714) (442,0.05484113581) (446,0.07232830161) (450,0.07200650241) (454,0.07289013208) (458,0.07636946026) (462,0.08524481618) (466,0.06709137225) (470,0.05408451525) (474,0.06460023044) (478,0.05878569661) (482,0.06391658155) (486,0.05787221075) (490,0.06670859545) (494,0.06571598458) (498,0.06094841717) (502,0.06122538149) (506,0.06014447844) (510,0.05601279465) (514,0.04506245229) (518,0.04616015375) (522,0.05698535424) (526,0.05447122209) (530,0.04191933296) (534,0.04806026976) (538,0.05588097454) (542,0.03880276203) (546,0.03244686828) (550,0.04361615879) (554,0.06059845777) (558,0.06418867114) (562,0.06292479077) (566,0.074469788) (570,0.07579007551) (574,0.07815721915) (578,0.07688361479) (582,0.07748453166) (586,0.07156860396) (590,0.07275247286) (594,0.08175439854) (598,0.0938024393) (602,0.08537054087) (606,0.07458031735) (610,0.07820592336) (614,0.08041714601) (618,0.07706338811) (622,0.0747905799) (626,0.08034348833) (630,0.08656950376) (634,0.09028762266) (638,0.08667535077) (642,0.09029816684) (646,0.08648330014) (650,0.08020602918) (654,0.0795912561) (658,0.07936947809) (662,0.08380534333) (666,0.07718341498) (670,0.08372393344) (674,0.0826294455) (678,0.07485353683) (682,0.0771828564) (686,0.07921043532) (690,0.07864355338) (694,0.07733281473) (698,0.08465121488) (702,0.08656678488) (706,0.08607936582) (710,0.1031246234) (714,0.08288533537) (718,0.09665794197) (722,0.09782851271) (726,0.1056485437) (730,0.1024544786) (734,0.1011481485) (738,0.1000006234) (742,0.1017915791) (746,0.09213610085) (750,0.08349594671) (754,0.07245146832) (758,0.07588930176) (762,0.06780075139) (766,0.08075185267) (770,0.07709541266) (774,0.07086316558) (778,0.06737784966) (782,0.06186899367) (786,0.0637022068) (790,0.06946035668) (794,0.06840732106) (798,0.07040763077) (802,0.09146276972) (806,0.09532135843) (810,0.0981545771) (814,0.09924357112) (818,0.1027635261) (822,0.1044882674) (826,0.1055683941) (830,0.1129500374) (834,0.1200007449) (838,0.1213812054) (842,0.1199540672) (846,0.1124676961) (850,0.1093090803) (854,0.1009980143) (858,0.09781968528) (862,0.09661841914) (866,0.07844196267) (870,0.08109145412) (874,0.07153073971) (878,0.06253179676) (882,0.06095109903) (886,0.0435799459) (890,0.04957267792) (894,0.05166986559) (898,0.05661874724) (902,0.06259025651) (906,0.05943554176) (910,0.05737523687) (914,0.06608972515) (918,0.07325618482) (922,0.06886795715) (926,0.0749117235) (930,0.08150637603) (934,0.08770620563) (938,0.05933128368) (942,0.06817750103) (946,0.06394963146) (950,0.0714914496) (954,0.07684607083) (958,0.07847252245) (962,0.06843014635) (966,0.04983328009) (970,0.04883779611) (974,0.06140174038) (978,0.07707852306) (982,0.0641372278) (986,0.07275402419) (990,0.07554157942) (994,0.06771892149) (998,0.07226993334)};
\addplot[pubNavy, opacity=0.20, line width=0.35pt] coordinates {(10,-0.03851967776) (14,-0.04617430163) (18,-0.04108095658) (22,-0.04330555986) (26,-0.03314870366) (30,-0.03272002167) (34,-0.02709239986) (38,-0.02872740255) (42,-0.02722426987) (46,-0.02909427099) (50,-0.01575297663) (54,-0.03107612419) (58,-0.02692847175) (62,-0.04412592486) (66,-0.02782102785) (70,-0.03028065197) (74,-0.01722064964) (78,-0.01292496999) (82,-0.005703940936) (86,0.000438485849) (90,0.0217797981) (94,0.01958220509) (98,0.04423928875) (102,0.03788937347) (106,0.06906894705) (110,0.08079531197) (114,0.09537267627) (118,0.1054112548) (122,0.1183500018) (126,0.1181700674) (130,0.1300654816) (134,0.1209578502) (138,0.1318551514) (142,0.1289268882) (146,0.1320873844) (150,0.123035631) (154,0.1334851137) (158,0.1236636132) (162,0.1197644911) (166,0.1218322167) (170,0.1093839822) (174,0.102841472) (178,0.09575779704) (182,0.09415295912) (186,0.09124101288) (190,0.08852745587) (194,0.09297729175) (198,0.0907717579) (202,0.09893769994) (206,0.0916070794) (210,0.09526036349) (214,0.08786416443) (218,0.09667756732) (222,0.0984920013) (226,0.1119143782) (230,0.109452729) (234,0.115110842) (238,0.1125919941) (242,0.1087088339) (246,0.1142863288) (250,0.1115523924) (254,0.1234830846) (258,0.121177214) (262,0.1190380664) (266,0.109287159) (270,0.1087606411) (274,0.1091019737) (278,0.1068590901) (282,0.1088005544) (286,0.1142333778) (290,0.120128606) (294,0.1053107339) (298,0.1057400995) (302,0.1118015261) (306,0.1126554878) (310,0.103306351) (314,0.1031788933) (318,0.103061552) (322,0.1005562356) (326,0.09730488836) (330,0.1007188353) (334,0.1010626754) (338,0.09810676382) (342,0.1135192898) (346,0.1030241691) (350,0.1032172138) (354,0.0776741029) (358,0.08106274569) (362,0.09482133105) (366,0.09828979752) (370,0.1015371493) (374,0.1099337273) (378,0.106503589) (382,0.09378779218) (386,0.09562461544) (390,0.09224756245) (394,0.08568695952) (398,0.09530692888) (402,0.08724618925) (406,0.07835296638) (410,0.09046609293) (414,0.09464718328) (418,0.1076877564) (422,0.08969888111) (426,0.09440684749) (430,0.08427836108) (434,0.09748443269) (438,0.1007528007) (442,0.09786646495) (446,0.1129789787) (450,0.116302568) (454,0.1161137604) (458,0.1290846521) (462,0.1357697526) (466,0.1284617169) (470,0.1191362149) (474,0.1339612503) (478,0.130533137) (482,0.1411672217) (486,0.1277387839) (490,0.1336174132) (494,0.1313488873) (498,0.1258557184) (502,0.1285629145) (506,0.134752327) (510,0.1331926308) (514,0.1161050491) (518,0.1041869035) (522,0.1205512884) (526,0.1278661766) (530,0.1062075087) (534,0.1114126768) (538,0.126702146) (542,0.1123692209) (546,0.1053098126) (550,0.1079821391) (554,0.1186663769) (558,0.1102534068) (562,0.1096485118) (566,0.1243175438) (570,0.1196335214) (574,0.1102301641) (578,0.1102556323) (582,0.1164174703) (586,0.1099247399) (590,0.1032563311) (594,0.1049606929) (598,0.1236229335) (602,0.1215675873) (606,0.1076998457) (610,0.1188561383) (614,0.1212831891) (618,0.1244405523) (622,0.1203729486) (626,0.1299828889) (630,0.1337601416) (634,0.1294530967) (638,0.1286612403) (642,0.1400685884) (646,0.1260085894) (650,0.1214074233) (654,0.1187557397) (658,0.1235876281) (662,0.1347730354) (666,0.12847167) (670,0.1351800149) (674,0.1438673006) (678,0.1381955158) (682,0.1519880207) (686,0.1583506195) (690,0.1538808161) (694,0.1505709091) (698,0.1530270673) (702,0.1524304412) (706,0.151441157) (710,0.1554786378) (714,0.143046844) (718,0.1494859153) (722,0.1351820481) (726,0.136596007) (730,0.128694343) (734,0.1177649386) (738,0.1168326668) (742,0.1187926722) (746,0.1038794446) (750,0.1045268799) (754,0.09688324381) (758,0.09368935269) (762,0.09679368349) (766,0.09814591663) (770,0.1042590262) (774,0.100502984) (778,0.09495050439) (782,0.1031644918) (786,0.1087699832) (790,0.1139257579) (794,0.1182894325) (798,0.1167257893) (802,0.1250810752) (806,0.1298879835) (810,0.1386680079) (814,0.1429956121) (818,0.1445778619) (822,0.1457712194) (826,0.1538953142) (830,0.1594056779) (834,0.168050442) (838,0.160040987) (842,0.1542409243) (846,0.1559423698) (850,0.1567682487) (854,0.1667600554) (858,0.1595153292) (862,0.1621794327) (866,0.1487186321) (870,0.1452718059) (874,0.133903865) (878,0.118762156) (882,0.109992006) (886,0.09336590506) (890,0.1037954236) (894,0.1033560492) (898,0.1088738833) (902,0.1135376401) (906,0.1038726441) (910,0.09841340108) (914,0.1063968006) (918,0.1126833056) (922,0.1069836153) (926,0.1107298838) (930,0.1132935818) (934,0.1161109751) (938,0.1047146771) (942,0.1203665563) (946,0.114464056) (950,0.1121359605) (954,0.1148568872) (958,0.1184053836) (962,0.1042039114) (966,0.0890935266) (970,0.09283958995) (974,0.111910167) (978,0.1243637159) (982,0.121001444) (986,0.1178596573) (990,0.1110978249) (994,0.1012582104) (998,0.1108451552)};
\addplot[pubNavy, opacity=0.20, line width=0.35pt] coordinates {(10,-0.0736519381) (14,-0.03612871936) (18,-0.03884219528) (22,-0.04428295585) (26,-0.02345090878) (30,-0.02372423583) (34,-0.009481817767) (38,-0.01827215188) (42,-0.02313841383) (46,-0.02014770423) (50,-0.009647317388) (54,-0.004659236733) (58,0.001911613302) (62,-0.009483575611) (66,-0.001781074364) (70,-0.003415149936) (74,0.005668685414) (78,0.006029955045) (82,0.003305106467) (86,0.01254959832) (90,0.03112376069) (94,0.0296543659) (98,0.04122440609) (102,0.03070580906) (106,0.04879285253) (110,0.05169363142) (114,0.06020724551) (118,0.06585050441) (122,0.06912167287) (126,0.06756241494) (130,0.08185438392) (134,0.06148622963) (138,0.05496835076) (142,0.05684714785) (146,0.05959813524) (150,0.03998343109) (154,0.05126062019) (158,0.03381564087) (162,0.0310300492) (166,0.03336046925) (170,0.01358594659) (174,0.01608135461) (178,0.01665852474) (182,0.02626662252) (186,0.03214604672) (190,0.02988220983) (194,0.02901868777) (198,0.02389161713) (202,0.02230445599) (206,0.01800319661) (210,0.00970473642) (214,-0.008507881437) (218,0.003362455105) (222,0.01468818884) (226,0.02074615459) (230,0.007081352617) (234,0.002378904318) (238,-0.002591543328) (242,0.002621602935) (246,0.01561866418) (250,0.02251978895) (254,0.03746783191) (258,0.054857961) (262,0.05015879112) (266,0.03637743745) (270,0.03328160155) (274,0.03529416362) (278,0.02939413846) (282,0.03741643776) (286,0.04329077533) (290,0.04590421924) (294,0.02414192009) (298,0.01485351767) (302,0.01139703816) (306,0.002803841435) (310,-0.004870653376) (314,-0.004898785763) (318,-0.00565572787) (322,-0.005806229835) (326,-0.01365029679) (330,-0.01899094837) (334,-0.01855625115) (338,-0.01374706554) (342,0.004645252752) (346,-0.002546337421) (350,-0.004181078062) (354,-0.01430393586) (358,-0.01382880474) (362,0.001410327307) (366,-0.002984072798) (370,-0.004253320357) (374,0.008571948265) (378,0.008831961723) (382,-0.0005752634366) (386,0.001067874715) (390,-0.006716834954) (394,-0.0109860388) (398,0.01192123336) (402,0.01006990771) (406,0.01130830391) (410,0.02668946548) (414,0.03679069234) (418,0.06268510655) (422,0.04808024396) (426,0.05262941459) (430,0.05523019519) (434,0.07104604234) (438,0.06745237213) (442,0.0626919578) (446,0.0710550685) (450,0.07247245599) (454,0.06913140598) (458,0.08244713917) (462,0.09438646323) (466,0.08247444887) (470,0.06833493614) (474,0.08520534292) (478,0.07888464921) (482,0.09187156506) (486,0.084589645) (490,0.08970675468) (494,0.09191985885) (498,0.09703986154) (502,0.1043118776) (506,0.1006657245) (510,0.1009066516) (514,0.08736587067) (518,0.0808080238) (522,0.09655981241) (526,0.0938317348) (530,0.08837092179) (534,0.09163736728) (538,0.1054314336) (542,0.1025916647) (546,0.09287223956) (550,0.08851713765) (554,0.09892556735) (558,0.0904332329) (562,0.08332005725) (566,0.09594767187) (570,0.0992061817) (574,0.09632779004) (578,0.08674877105) (582,0.08290150455) (586,0.07025717356) (590,0.05641246842) (594,0.05846694348) (598,0.07594960355) (602,0.07490109006) (606,0.05825770396) (610,0.07969660448) (614,0.08139460261) (618,0.08105520988) (622,0.07680838135) (626,0.08206294599) (630,0.08662520313) (634,0.09323097965) (638,0.09178212825) (642,0.0956813723) (646,0.09135777708) (650,0.08431279679) (654,0.09675445671) (658,0.1037173804) (662,0.1088905539) (666,0.09532731855) (670,0.09317068306) (674,0.1051433324) (678,0.102566644) (682,0.117435886) (686,0.1176665005) (690,0.1146642524) (694,0.1135112432) (698,0.1107554163) (702,0.1086019204) (706,0.1065217982) (710,0.1126076366) (714,0.1024544027) (718,0.110695039) (722,0.1070631052) (726,0.09971786409) (730,0.09333799973) (734,0.09140064441) (738,0.09203039761) (742,0.08869294751) (746,0.07537374107) (750,0.07896618017) (754,0.06708575559) (758,0.05953975446) (762,0.04918978381) (766,0.06396291958) (770,0.06406778753) (774,0.05742451156) (778,0.05009123124) (782,0.0575844961) (786,0.06024837498) (790,0.0663455996) (794,0.07250593386) (798,0.07382775523) (802,0.09532060331) (806,0.09469559123) (810,0.09478185225) (814,0.09210252861) (818,0.1038483048) (822,0.1086067373) (826,0.112950641) (830,0.1193625083) (834,0.1303928831) (838,0.1273133177) (842,0.1328461498) (846,0.1261144499) (850,0.1215919119) (854,0.1189117101) (858,0.1201480068) (862,0.124087356) (866,0.10980191) (870,0.1064205211) (874,0.09597378395) (878,0.08698717819) (882,0.08021778788) (886,0.0663073964) (890,0.0658859091) (894,0.06682476923) (898,0.07587541572) (902,0.07833865508) (906,0.07375103886) (910,0.06891725697) (914,0.07943394206) (918,0.07961049499) (922,0.06885772959) (926,0.0732757327) (930,0.07165285935) (934,0.07447621336) (938,0.05978337947) (942,0.07254007843) (946,0.06008177976) (950,0.06699273953) (954,0.07242742677) (958,0.07116874355) (962,0.06345919036) (966,0.05321363804) (970,0.06422130802) (974,0.08105592845) (978,0.09741138489) (982,0.09658749432) (986,0.0987094391) (990,0.08869700373) (994,0.08175984837) (998,0.09452375004)};
\addplot[pubNavy, ultra thick] coordinates {(10,-0.02361737512) (14,-0.03586768385) (18,-0.03723880518) (22,-0.03711821851) (26,-0.02264547515) (30,-0.0256701503) (34,-0.0129737104) (38,-0.02106052257) (42,-0.02381127923) (46,-0.0254366254) (50,-0.01379455177) (54,-0.01824169424) (58,-0.01099665239) (62,-0.01891319728) (66,-0.0009811786913) (70,0.0006737558026) (74,0.01620274334) (78,0.0161912247) (82,0.01467773599) (86,0.02726229974) (90,0.04113204157) (94,0.03529422269) (98,0.04491152784) (102,0.03518836419) (106,0.05897442628) (110,0.06185023851) (114,0.07176253696) (118,0.07855264471) (122,0.08117473527) (126,0.07833309396) (130,0.08833784337) (134,0.0697727775) (138,0.07180652285) (142,0.06863791876) (146,0.07250041569) (150,0.062587701) (154,0.07345815646) (158,0.05731096352) (162,0.05082791441) (166,0.04960005618) (170,0.03887663331) (174,0.03744031535) (178,0.02982168696) (182,0.04785189738) (186,0.05489933835) (190,0.04632235519) (194,0.05234847881) (198,0.05296289394) (202,0.05900713137) (206,0.04733890898) (210,0.05294564799) (214,0.05011945721) (218,0.05499033808) (222,0.05486961654) (226,0.06715754032) (230,0.06189588086) (234,0.0635573132) (238,0.06207382821) (242,0.06010929455) (246,0.06120818688) (250,0.057254226) (254,0.06947271258) (258,0.07553321301) (262,0.07009268427) (266,0.06638462443) (270,0.05869614184) (274,0.0589412667) (278,0.05715541659) (282,0.06519580606) (286,0.06557585723) (290,0.07525280402) (294,0.06303699558) (298,0.06354496301) (302,0.06974594304) (306,0.06963262122) (310,0.06015710409) (314,0.06279270226) (318,0.06034118882) (322,0.05806979158) (326,0.05596163424) (330,0.05955038924) (334,0.05953031818) (338,0.05758732265) (342,0.07014894495) (346,0.05756665314) (350,0.05563828225) (354,0.03508780124) (358,0.04072506969) (362,0.04661663554) (366,0.04570106304) (370,0.05045878739) (374,0.05526994761) (378,0.05750771578) (382,0.05527746742) (386,0.05677310552) (390,0.05147149587) (394,0.04482267061) (398,0.05511346709) (402,0.04658796035) (406,0.0422930781) (410,0.05617027189) (414,0.06162455633) (418,0.07123008969) (422,0.05764131767) (426,0.06323064656) (430,0.05602797861) (434,0.06249228848) (438,0.06665666423) (442,0.06060980741) (446,0.07369918822) (450,0.08098216654) (454,0.07820732371) (458,0.08469086942) (462,0.09119218278) (466,0.07749247987) (470,0.06856732445) (474,0.08448912569) (478,0.07899350466) (482,0.08619400218) (486,0.07626561735) (490,0.08563694093) (494,0.08393021509) (498,0.08045419754) (502,0.08458541118) (506,0.08465740453) (510,0.08254165352) (514,0.07116253807) (518,0.0631302585) (522,0.07848965418) (526,0.07825682643) (530,0.07888058563) (534,0.08152289072) (538,0.09375607885) (542,0.07816080669) (546,0.07390438171) (550,0.07924727115) (554,0.08902931863) (558,0.07569090751) (562,0.06978683812) (566,0.08254364045) (570,0.08365792273) (574,0.08208087416) (578,0.07967835254) (582,0.0789648378) (586,0.07091288876) (590,0.060317359) (594,0.06468017494) (598,0.08130377842) (602,0.08007065308) (606,0.07412550502) (610,0.08668847262) (614,0.09032702889) (618,0.09268064205) (622,0.08837051854) (626,0.09797679207) (630,0.09696237513) (634,0.09478831045) (638,0.0899448403) (642,0.09258316556) (646,0.08892053861) (650,0.08448050869) (654,0.08736535335) (658,0.08929923012) (662,0.09739574757) (666,0.08705984106) (670,0.08894762966) (674,0.0990417099) (678,0.08848665732) (682,0.09769178936) (686,0.1001703126) (690,0.1009823128) (694,0.1035174398) (698,0.1044724146) (702,0.09984731556) (706,0.09342958537) (710,0.1008187486) (714,0.08554726438) (718,0.0961732691) (722,0.09222241037) (726,0.09331180629) (730,0.08711760119) (734,0.08467009728) (738,0.08679995008) (742,0.08443779148) (746,0.07501937215) (750,0.07773570416) (754,0.06557497075) (758,0.06271099315) (762,0.0562763701) (766,0.07075224008) (770,0.07535492232) (774,0.06716648691) (778,0.06209839779) (782,0.05989424128) (786,0.06711858514) (790,0.07130475242) (794,0.07233681797) (798,0.07273747562) (802,0.09083193269) (806,0.09500847483) (810,0.09135061329) (814,0.09152215954) (818,0.09774765289) (822,0.1043955675) (826,0.1072372894) (830,0.1108241808) (834,0.1166760055) (838,0.1198755806) (842,0.1206995164) (846,0.1160993391) (850,0.1120457583) (854,0.1080905657) (858,0.1006815851) (862,0.113471964) (866,0.09878539514) (870,0.1013396643) (874,0.09554579418) (878,0.08726201973) (882,0.07860619762) (886,0.05890438881) (890,0.06577013849) (894,0.06666454533) (898,0.07350753078) (902,0.07928008878) (906,0.0746381822) (910,0.07172692848) (914,0.07399834861) (918,0.07686463759) (922,0.07281511493) (926,0.07568849687) (930,0.07900782673) (934,0.08142475522) (938,0.06664117006) (942,0.07615104491) (946,0.06559477725) (950,0.06965780977) (954,0.07521971259) (958,0.07506016982) (962,0.06802359721) (966,0.06056753257) (970,0.06322304519) (974,0.07720163301) (978,0.08593270749) (982,0.08487355961) (986,0.08496571961) (990,0.08624444969) (994,0.08052857133) (998,0.08780098925)};
\nextgroupplot[title={Rolling terminal wealth}, xlabel={Pre-training episode}]
\addplot[pubNavy, opacity=0.20, line width=0.35pt] coordinates {(10,115045.5536) (14,115998.2215) (18,114339.88) (22,114139.0415) (26,116164.4466) (30,115935.2567) (34,117791.2596) (38,116190.0822) (42,115578.6423) (46,116144.1137) (50,117489.8198) (54,117359.1956) (58,118087.1337) (62,116272.1748) (66,118704.783) (70,119799.678) (74,121930.2672) (78,121821.5223) (82,121709.958) (86,122622.2845) (90,124682.6524) (94,124114.3239) (98,125951.0904) (102,125844.9542) (106,128369.9163) (110,127839.5753) (114,129059.4365) (118,129116.2694) (122,129962.2131) (126,129306.3931) (130,130614.0225) (134,129223.9925) (138,129768.3278) (142,130605.7966) (146,130398.343) (150,128597.4474) (154,129122.9464) (158,126930.0906) (162,127209.7249) (166,127013.0733) (170,124785.2448) (174,124828.9265) (178,124308.5784) (182,124629.1914) (186,124630.8064) (190,123842.3639) (194,124443.6736) (198,125295.5737) (202,125056.0886) (206,125284.1501) (210,125783.3785) (214,123652.852) (218,125462.9687) (222,125179.4682) (226,125472.5595) (230,124763.7428) (234,124061.8011) (238,124615.9148) (242,124142.7512) (246,125431.0061) (250,124963.5673) (254,127107.8775) (258,126501.5974) (262,126219.337) (266,125252.4665) (270,126007.4234) (274,126681.0318) (278,127085.8875) (282,128272.4327) (286,129589.1334) (290,131263.2414) (294,129223.8463) (298,129595.5652) (302,130713.6918) (306,130939.6882) (310,130628.2538) (314,131537.9232) (318,130889.7232) (322,130959.4709) (326,130788.9547) (330,130080.4545) (334,129221.9342) (338,128062.998) (342,130459.4525) (346,128817.1462) (350,128994.9694) (354,127249.3884) (358,129462.9205) (362,130954.5954) (366,130861.2597) (370,132470.1939) (374,133162.3294) (378,134468.4451) (382,133875.3223) (386,134831.9082) (390,134428.5087) (394,134064.3213) (398,135273.5096) (402,133603.2627) (406,131106.2193) (410,132439.0804) (414,133063.6232) (418,133575.0916) (422,129806.9953) (426,130590.8799) (430,128232.4732) (434,128864.3347) (438,129311.5787) (442,128649.3954) (446,129573.0732) (450,130014.3905) (454,128489.9551) (458,130068.7147) (462,130824.8579) (466,129711.0842) (470,128849.8222) (474,131819.8249) (478,130061.0851) (482,131460.2018) (486,130056.3817) (490,132192.518) (494,131647.5204) (498,130974.0075) (502,132342.92) (506,134527.6953) (510,133245.3794) (514,131263.9044) (518,130524.7108) (522,131104.6982) (526,131612.0442) (530,132399.4615) (534,133044.1467) (538,135110.654) (542,132076.1872) (546,131699.9412) (550,131702.3061) (554,133088.6159) (558,131554.3714) (562,132133.8377) (566,135170.6808) (570,134622.2942) (574,134330.4154) (578,133812.3427) (582,133681.7781) (586,132272.692) (590,130846.2705) (594,131053.5835) (598,134773.1271) (602,134896.4335) (606,133221.4963) (610,134681.0192) (614,134275.6315) (618,134742.8836) (622,134098.1022) (626,134466.2602) (630,134242.7072) (634,134115.333) (638,135677.5997) (642,137972.5945) (646,135726.4651) (650,134088.5176) (654,134290.2184) (658,134218.9284) (662,136680.8254) (666,133700.7075) (670,136570.857) (674,138530.0689) (678,136035.5026) (682,137637.2297) (686,137912.6539) (690,137177.8873) (694,137407.6447) (698,138996.8409) (702,138947.6316) (706,138191.3163) (710,137679.1356) (714,135597.9636) (718,136372.6475) (722,133217.1551) (726,134041.1622) (730,132697.7127) (734,131308.1843) (738,130289.8972) (742,130911.6244) (746,126906.9155) (750,126157.8827) (754,125026.9066) (758,125555.9202) (762,125160.5254) (766,126835.2248) (770,126783.6955) (774,126404.9972) (778,124483.7823) (782,126701.0335) (786,128743.7539) (790,128614.4308) (794,129836.9299) (798,130133.3271) (802,130812.3392) (806,130610.3171) (810,131236.4066) (814,131856.9401) (818,132102.9908) (822,131912.7124) (826,132913.4092) (830,133052.9359) (834,134244.3947) (838,133766.1181) (842,133931.553) (846,133446.3133) (850,133487.8248) (854,135481.4926) (858,134185.5923) (862,135026.5296) (866,134265.0171) (870,133812.7473) (874,133423.6288) (878,132209.1291) (882,129255.5249) (886,127083.8743) (890,128662.8178) (894,129111.7178) (898,129372.5483) (902,130583.7942) (906,128524.2084) (910,128149.5748) (914,128494.2698) (918,128587.3695) (922,128418.678) (926,128884.8576) (930,128583.7515) (934,128635.2041) (938,126467.9371) (942,129509.0953) (946,128504.2558) (950,127238.4869) (954,127581.3982) (958,127623.3175) (962,125657.8623) (966,125540.5507) (970,124658.9596) (974,127284.9625) (978,130216.6353) (982,130711.2375) (986,130606.8176) (990,130020.8014) (994,128103.2022) (998,130298.9241)};
\addplot[pubNavy, opacity=0.20, line width=0.35pt] coordinates {(10,115583.8591) (14,114840.2888) (18,113788.9522) (22,113052.6428) (26,114144.229) (30,111900.4746) (34,112133.7294) (38,111932.4784) (42,111381.6764) (46,110707.2056) (50,112986.9514) (54,111942.0432) (58,112582.805) (62,110127.2337) (66,111992.6171) (70,111583.8106) (74,113590.7833) (78,114077.0757) (82,115114.2476) (86,116695.4038) (90,118787.4251) (94,119075.0344) (98,121634.2188) (102,119369.5064) (106,121614.7945) (110,121863.2167) (114,122915.5788) (118,123703.3805) (122,124183.396) (126,123957.2859) (130,126342.2775) (134,123419.0758) (138,123919.4497) (142,123197.8542) (146,123861.776) (150,122498.1065) (154,124735.924) (158,122890.2543) (162,123448.0467) (166,124181.019) (170,121793.2786) (174,122112.5815) (178,120866.345) (182,122603.3149) (186,122960.5551) (190,122923.7526) (194,125042.3943) (198,126083.4707) (202,126777.8222) (206,125534.5265) (210,126476.7119) (214,124117.6024) (218,125571.0697) (222,125831.7173) (226,126676.2762) (230,125838.2848) (234,125414.7463) (238,125590.6224) (242,123405.2354) (246,124203.7472) (250,124524.1398) (254,126069.2835) (258,126326.9998) (262,126017.4141) (266,124492.4546) (270,124401.63) (274,125014.0536) (278,124374.3504) (282,126220.8866) (286,125363.6496) (290,126853.5468) (294,124592.6314) (298,123558.4401) (302,124629.0174) (306,124821.697) (310,124488.5575) (314,125011.9106) (318,124483.3482) (322,124538.1596) (326,123688.7704) (330,124038.3285) (334,124271.9023) (338,124320.936) (342,126189.5609) (346,124412.4083) (350,124915.9106) (354,121296.8868) (358,123349.2409) (362,124376.3231) (366,124164.2724) (370,123391.2372) (374,124719.1235) (378,124956.4881) (382,124457.8974) (386,123815.7298) (390,123442.65) (394,123832.2154) (398,124412.7743) (402,121628.6774) (406,120223.1245) (410,120045.5986) (414,121282.3992) (418,122320.9667) (422,120559.1798) (426,123108.8835) (430,121773.845) (434,122239.4457) (438,122630.1048) (442,122230.2633) (446,123760.9511) (450,124185.2907) (454,123401.8424) (458,125173.7875) (462,126521.8533) (466,125234.7597) (470,123771.7084) (474,125668.0286) (478,122894.9373) (482,123298.8811) (486,122095.4824) (490,123025.756) (494,122347.8879) (498,122105.8914) (502,123117.6756) (506,123275.774) (510,123027.8202) (514,120905.1129) (518,121210.7712) (522,123437.0602) (526,122073.8153) (530,122211.5685) (534,123697.4324) (538,125679.0085) (542,121963.5101) (546,120888.0287) (550,120394.2266) (554,122714.3824) (558,121818.6611) (562,124148.8509) (566,126452.1638) (570,124921.0913) (574,125089.7488) (578,124768.533) (582,125327.0388) (586,125376.563) (590,125558.1391) (594,127946.3837) (598,130551.3379) (602,129644.272) (606,129298.2268) (610,129777.8777) (614,127735.2908) (618,127964.544) (622,128008.495) (626,128096.0669) (630,127542.9425) (634,127241.911) (638,126477.9185) (642,126254.3128) (646,123931.9168) (650,123406.3387) (654,124319.8631) (658,124979.1895) (662,127035.2836) (666,125263.1641) (670,126929.8847) (674,128371.3035) (678,127600.1581) (682,128437.0483) (686,127909.1888) (690,128345.0802) (694,128918.3173) (698,128827.4265) (702,128249.2752) (706,125816.6638) (710,126247.5978) (714,124226.8628) (718,126650.6256) (722,125162.257) (726,124669.522) (730,124066.7385) (734,123905.8752) (738,123370.9315) (742,122906.0105) (746,121843.3229) (750,121617.8307) (754,121341.9086) (758,121586.3079) (762,120871.196) (766,121634.1119) (770,120983.8716) (774,120706.4371) (778,120103.7597) (782,120835.975) (786,122693.2556) (790,122410.0415) (794,122383.0335) (798,122454.462) (802,126015.9353) (806,125288.1152) (810,126800.462) (814,127774.4131) (818,129687.4314) (822,129710.3238) (826,130851.5973) (830,132205.2826) (834,134200.648) (838,134686.6904) (842,134658.4509) (846,133867.2285) (850,132398.344) (854,131204.4541) (858,130243.7141) (862,131016.6544) (866,128055.7462) (870,127401.9475) (874,127054.3647) (878,125782.9193) (882,124594.0308) (886,120067.0832) (890,121828.7848) (894,122946.9932) (898,124379.785) (902,125254.0477) (906,124660.869) (910,123634.5454) (914,124388.5234) (918,124930.5884) (922,125310.1203) (926,125032.746) (930,124253.9405) (934,124059.9779) (938,121426.6858) (942,122174.7644) (946,121059.0498) (950,120811.8247) (954,122006.4041) (958,122449.5754) (962,122163.2573) (966,121074.2733) (970,120180.5276) (974,122517.1987) (978,125564.4091) (982,125417.6771) (986,128175.3232) (990,128213.9391) (994,125914.9928) (998,126436.7215)};
\addplot[pubNavy, opacity=0.20, line width=0.35pt] coordinates {(10,116593.7201) (14,113731.2671) (18,112146.2143) (22,112264.961) (26,112913.4604) (30,112122.6347) (34,112872.8452) (38,111305.9714) (42,110371.9114) (46,111237.8781) (50,112354.4562) (54,112181.1868) (58,112608.4119) (62,111358.0142) (66,113772.7303) (70,114803.999) (74,115764.95) (78,115669.721) (82,116604.4556) (86,117039.713) (90,119167.507) (94,118827.229) (98,119532.5739) (102,120112.5593) (106,122453.7637) (110,121692.7005) (114,122502.2446) (118,122803.6053) (122,123740.3248) (126,123856.4246) (130,124624.9887) (134,123359.2397) (138,124116.5946) (142,125721.9293) (146,125619.8977) (150,123180.9347) (154,123082.1211) (158,120686.0219) (162,120208.0358) (166,119665.538) (170,118278.9674) (174,118684.2001) (178,118608.478) (182,119040.5068) (186,118926.7098) (190,118448.7235) (194,118665.6841) (198,119772.5862) (202,121162.9349) (206,121368.4305) (210,121459.8719) (214,120563.0986) (218,121583.1714) (222,120710.7922) (226,120840.4793) (230,119254.6329) (234,119573.7467) (238,119134.3875) (242,119601.8838) (246,120736.0308) (250,120134.2131) (254,121470.0913) (258,122346.6809) (262,122529.5567) (266,121070.1456) (270,121923.9094) (274,123319.6844) (278,123582.4701) (282,124817.8405) (286,126893.452) (290,128517.4511) (294,125686.5513) (298,125909.8921) (302,127631.5545) (306,127639.2086) (310,127093.7314) (314,127622.352) (318,125903.7024) (322,125946.565) (326,125437.648) (330,125751.649) (334,125610.2137) (338,125223.9188) (342,128810.502) (346,127284.0892) (350,127155.4912) (354,123968.4766) (358,124268.2783) (362,125883.498) (366,126904.4549) (370,128638.3769) (374,130144.9464) (378,130248.6342) (382,129403.23) (386,130044.7678) (390,128123.4963) (394,127407.3668) (398,129182.9345) (402,127307.7302) (406,126587.2044) (410,128487.4724) (414,130204.0387) (418,131123.605) (422,128008.5741) (426,130243.9965) (430,128671.6284) (434,128634.2231) (438,128484.4357) (442,126247.3411) (446,127614.5983) (450,128946.8617) (454,128414.612) (458,129815.0947) (462,130538.8985) (466,127756.7821) (470,126970.8463) (474,128268.4727) (478,125950.215) (482,128186.3817) (486,127459.7809) (490,128631.2874) (494,128529.8656) (498,127534.6132) (502,129328.1699) (506,129894.8495) (510,130044.6703) (514,127825.9649) (518,128326.3517) (522,130604.475) (526,130492.8935) (530,129102.9781) (534,129726.1824) (538,130849.894) (542,129627.4708) (546,129582.1488) (550,130352.3025) (554,130559.6165) (558,130766.8036) (562,130036.188) (566,130197.2015) (570,130011.3172) (574,128628.8021) (578,128495.097) (582,128175.799) (586,127660.0198) (590,124417.8683) (594,124635.3411) (598,125358.1229) (602,122874.8934) (606,121554.0663) (610,122326.997) (614,122784.1018) (618,123409.5136) (622,121538.028) (626,122624.5754) (630,123044.4593) (634,123297.6405) (638,123214.3577) (642,125221.8752) (646,124634.6765) (650,125690.3009) (654,125789.278) (658,126006.2645) (662,125519.7934) (666,125157.4548) (670,124945.9746) (674,125125.3386) (678,125056.7548) (682,125985.1932) (686,126899.4799) (690,126856.321) (694,124625.0025) (698,123921.1229) (702,124682.941) (706,125350.1026) (710,127175.0917) (714,125325.4924) (718,125869.1918) (722,125769.2854) (726,125424.6136) (730,124463.2228) (734,123296.2273) (738,123476.2435) (742,124964.4015) (746,123708.802) (750,123444.3269) (754,122021.3336) (758,121955.0663) (762,122246.5891) (766,123059.5595) (770,122693.1267) (774,122757.1315) (778,123028.5036) (782,123293.6513) (786,123871.2769) (790,123119.817) (794,123681.2943) (798,123812.2343) (802,125130.0992) (806,124876.6789) (810,125234.0968) (814,125682.3399) (818,126891.6697) (822,126399.3998) (826,125856.179) (830,127342.1392) (834,126908.2508) (838,125965.3765) (842,126078.9152) (846,125590.2533) (850,125027.4856) (854,125189.5205) (858,125130.7325) (862,125244.4065) (866,123776.7326) (870,124328.6845) (874,123380.8853) (878,122240.3214) (882,121139.8708) (886,120346.1242) (890,121029.5598) (894,121018.7299) (898,122060.2811) (902,122360.0685) (906,121437.7338) (910,120367.7685) (914,120580.4515) (918,119605.823) (922,118799.8923) (926,119618.6116) (930,119722.9814) (934,119944.602) (938,118204.2238) (942,119532.2708) (946,117925.2194) (950,119057.191) (954,119667.2777) (958,119296.0987) (962,118024.5332) (966,116814.535) (970,117598.569) (974,119528.0446) (978,120503.8609) (982,120461.5036) (986,120352.89) (990,120457.0639) (994,120768.7188) (998,121041.5311)};
\addplot[pubNavy, opacity=0.20, line width=0.35pt] coordinates {(10,112093.1067) (14,111487.2889) (18,110686.0049) (22,111815.6373) (26,112931.2226) (30,112575.3652) (34,114659.2991) (38,113279.4031) (42,113625.4191) (46,114207.5397) (50,115929.2844) (54,115900.4535) (58,116848.7478) (62,114795.3314) (66,118006.2768) (70,118603.7659) (74,121415.5812) (78,121938.182) (82,121781.8457) (86,122895.4811) (90,125747.2907) (94,125124.256) (98,127106.0843) (102,125734.3181) (106,128678.2259) (110,128964.4647) (114,131181.8825) (118,130930.6937) (122,132160.9857) (126,132269.9393) (130,133152.3732) (134,130553.1341) (138,131625.5189) (142,130336.8294) (146,130733.9724) (150,128985.7567) (154,130416.7104) (158,128309.2766) (162,127505.4994) (166,127486.9488) (170,125215.703) (174,124279.0587) (178,122396.6931) (182,125797.7057) (186,126132.5079) (190,125815.5135) (194,128220.054) (198,128689.9092) (202,130457.2296) (206,128784.594) (210,129210.8874) (214,126618.3175) (218,127699.1759) (222,127908.7249) (226,129525.7497) (230,128738.4961) (234,129026.8327) (238,128829.7438) (242,127552.8182) (246,127391.0804) (250,125893.989) (254,126675.4974) (258,127956.7965) (262,128356.6384) (266,127254.4661) (270,127359.8217) (274,127738.9093) (278,127557.0338) (282,128465.2681) (286,129844.3439) (290,130880.6351) (294,128721.9015) (298,129966.4038) (302,131661.0224) (306,131655.0076) (310,130291.3422) (314,130115.5931) (318,129746.9584) (322,128302.2374) (326,127824.779) (330,127793.79) (334,125991.602) (338,125414.5028) (342,127047.3026) (346,124068.1818) (350,125511.7925) (354,122935.7371) (358,126050.66) (362,127251.2401) (366,127784.2113) (370,128625.2144) (374,130061.6892) (378,130738.4338) (382,129981.2796) (386,130323.2369) (390,130940.2035) (394,131815.1701) (398,132370.9862) (402,128155.4519) (406,125693.8084) (410,127115.6745) (414,128737.5409) (418,129769.9368) (422,126159.1775) (426,127482.1014) (430,126295.267) (434,128002.0897) (438,129473.486) (442,127712.2545) (446,128370.9221) (450,129674.3593) (454,129445.8112) (458,130331.1516) (462,131864.3821) (466,130495.939) (470,129511.4319) (474,131815.1358) (478,130702.3134) (482,131541.4766) (486,128990.3975) (490,131841.8262) (494,130717.6709) (498,129199.2343) (502,129704.6462) (506,130279.3412) (510,131526.3246) (514,129308.2333) (518,128214.5425) (522,130727.3078) (526,131593.0681) (530,130393.4532) (534,132020.6711) (538,134483.4258) (542,131495.9372) (546,130892.3599) (550,131607.2123) (554,133632.0433) (558,131423.9718) (562,132506.4994) (566,135854.4654) (570,135551.1423) (574,135017.3898) (578,134747.9106) (582,134284.0114) (586,133153.5749) (590,131813.3108) (594,133711.3373) (598,137408.7863) (602,138219.5032) (606,136212.7875) (610,136798.7658) (614,135967.8931) (618,136120.361) (622,135222.8469) (626,134872.6548) (630,134919.0972) (634,133998.2863) (638,134025.1979) (642,134156.8414) (646,131030.0665) (650,128381.347) (654,127839.5557) (658,128298.1075) (662,130101.0744) (666,127543.7) (670,129160.6527) (674,132680.9665) (678,131353.8536) (682,134192.2148) (686,134981.4159) (690,134080.5243) (694,136075.003) (698,136084.207) (702,135129.2513) (706,134030.1004) (710,134393.2936) (714,132160.9581) (718,135534.0195) (722,132725.933) (726,132620.6826) (730,130611.7016) (734,130146.945) (738,129186.85) (742,127709.4847) (746,124843.5866) (750,124811.7259) (754,123953.7173) (758,122522.7285) (762,120995.4128) (766,122459.8099) (770,122986.3179) (774,121723.3185) (778,120675.0666) (782,120899.4142) (786,123930.6346) (790,125739.9276) (794,125651.1586) (798,128161.9638) (802,132574.9387) (806,133870.1509) (810,135584.6696) (814,136194.343) (818,136603.4761) (822,136435.3022) (826,136914.9159) (830,138066.9725) (834,140049.2164) (838,139320.2242) (842,138810.2414) (846,137990.4745) (850,136537.4962) (854,136423.9189) (858,134713.7043) (862,135755.7735) (866,132595.5214) (870,132695.5244) (874,131568.0218) (878,130885.7141) (882,129413.3386) (886,126479.9901) (890,128608.1317) (894,129582.1112) (898,130883.6704) (902,132427.8712) (906,131203.2588) (910,131196.6916) (914,131891.2964) (918,132437.7828) (922,132443.7684) (926,133378.7265) (930,131742.9163) (934,131145.1387) (938,128255.6876) (942,130420.9134) (946,129436.1167) (950,128610.095) (954,128928.8305) (958,129037.1505) (962,126707.9153) (966,125806.2928) (970,125152.2283) (974,127850.7115) (978,130764.6523) (982,130688.9245) (986,131767.4817) (990,131813.4438) (994,127964.3039) (998,130159.5303)};
\addplot[pubNavy, opacity=0.20, line width=0.35pt] coordinates {(10,114446.2497) (14,114393.9723) (18,113029.5293) (22,112721.0992) (26,114864.7707) (30,113682.2842) (34,114426.0838) (38,112523.0382) (42,112836.9352) (46,113295.5262) (50,115544.5298) (54,115160.8583) (58,115860.4079) (62,113840.3167) (66,115880.8455) (70,116465.596) (74,117860.7381) (78,117856.4972) (82,117667.5329) (86,119637.5296) (90,120745.0388) (94,119775.075) (98,120179.3259) (102,119030.4142) (106,120958.2257) (110,120944.5656) (114,122814.8458) (118,122572.526) (122,123367.09) (126,123402.2299) (130,125238.2209) (134,122622.8062) (138,123002.0255) (142,123294.6477) (146,124258.3802) (150,122221.7908) (154,123461.3917) (158,121567.1328) (162,121495.6878) (166,121757.5153) (170,118962.5808) (174,118986.101) (178,117612.9786) (182,119411.5158) (186,120153.0631) (190,119634.3926) (194,121109.713) (198,121542.7317) (202,122345.5728) (206,120699.6708) (210,120850.8831) (214,117631.9534) (218,119097.6731) (222,120156.912) (226,121438.9122) (230,119810.9798) (234,119774.0096) (238,119213.1807) (242,119148.893) (246,118956.5721) (250,119244.2864) (254,120506.7531) (258,121704.0292) (262,121021.8023) (266,121559.8738) (270,121072.4068) (274,120751.2648) (278,121348.71) (282,122796.8859) (286,123387.0794) (290,125009.928) (294,122833.0751) (298,124038.8084) (302,125965.6628) (306,126539.4434) (310,125441.1713) (314,125511.9056) (318,125328.6406) (322,124637.5183) (326,124372.5366) (330,123729.5554) (334,123403.0835) (338,122757.0796) (342,125253.0354) (346,123123.1611) (350,122107.7373) (354,118215.3785) (358,120398.3654) (362,122144.0755) (366,122098.8884) (370,122835.1996) (374,122219.1024) (378,123262.6654) (382,121944.6644) (386,122731.6026) (390,122606.3973) (394,123207.6886) (398,124320.192) (402,122468.525) (406,120736.082) (410,122769.6219) (414,123129.7716) (418,123999.0757) (422,123111.176) (426,124343.6361) (430,123972.8068) (434,123931.1697) (438,124039.6276) (442,123402.9439) (446,124637.203) (450,125191.2595) (454,125312.0836) (458,125695.5573) (462,127049.7046) (466,125908.1399) (470,123897.8887) (474,125698.459) (478,124333.7882) (482,125703.6757) (486,125156.0944) (490,125260.5615) (494,123865.9094) (498,123838.8307) (502,124458.0634) (506,124105.5468) (510,122540.5102) (514,120439.187) (518,120546.48) (522,122529.1971) (526,121784.0935) (530,120467.5784) (534,120558.4234) (538,121377.1192) (542,119066.8359) (546,119911.7765) (550,119336.6538) (554,120996.5764) (558,120138.4076) (562,120707.6264) (566,121754.7814) (570,121231.1163) (574,119401.6642) (578,118701.9878) (582,119068.3361) (586,118644.5105) (590,118196.3587) (594,118884.3266) (598,120291.6514) (602,118591.2452) (606,117321.9594) (610,118906.5657) (614,118944.9055) (618,118908.5273) (622,119393.4323) (626,121839.1545) (630,121954.941) (634,122162.9986) (638,122697.622) (642,122508.7942) (646,121371.6617) (650,121629.7345) (654,122608.2709) (658,122570.1465) (662,122925.8441) (666,122719.4059) (670,122889.9071) (674,122310.9499) (678,120226.998) (682,121255.9681) (686,121332.7756) (690,121364.2268) (694,120571.8599) (698,119269.076) (702,119832.7871) (706,119533.5111) (710,120806.8121) (714,118186.4891) (718,119111.4898) (722,119156.7094) (726,119356.1947) (730,119697.4741) (734,119634.0104) (738,119550.3533) (742,119755.58) (746,120284.708) (750,119766.1071) (754,117687.896) (758,117438.4826) (762,117786.4315) (766,119737.2552) (770,119998.0881) (774,120151.3397) (778,119596.8377) (782,119519.4001) (786,119738.1942) (790,118654.7258) (794,118429.396) (798,118762.8209) (802,122809.9287) (806,121484.5936) (810,121715.4062) (814,121132.58) (818,121742.5212) (822,120361.3516) (826,119955.488) (830,120470.3045) (834,120496.1903) (838,120767.4711) (842,121331.0043) (846,120961.8043) (850,119476.6209) (854,118552.6953) (858,119224.2699) (862,119770.0835) (866,118050.1577) (870,118750.6255) (874,117731.4502) (878,116275.486) (882,115869.0514) (886,115196.2112) (890,116253.27) (894,117473.5653) (898,117383.8999) (902,119192.1018) (906,118449.9973) (910,119298.0267) (914,118896.7746) (918,120292.0713) (922,120607.1262) (926,122121.5954) (930,122296.2488) (934,121303.1068) (938,119964.0106) (942,120298.8086) (946,119408.2834) (950,119842.9331) (954,119881.8508) (958,119034.0769) (962,117567.2495) (966,116734.3137) (970,115850.0843) (974,118470.0872) (978,120500.701) (982,120846.506) (986,122961.43) (990,122698.3781) (994,122170.1591) (998,123402.7805)};
\addplot[pubNavy, opacity=0.20, line width=0.35pt] coordinates {(10,113994.4459) (14,112651.6921) (18,110864.3597) (22,111419.8481) (26,114497.6642) (30,113664.513) (34,114041.0556) (38,112704.9183) (42,113069.9771) (46,112584.4487) (50,115118.7718) (54,113863.6403) (58,114342.585) (62,111988.0722) (66,114605.965) (70,114296.025) (74,115624.2634) (78,114810.712) (82,115782.5231) (86,117506.5913) (90,120731.0391) (94,120066.3511) (98,122201.5515) (102,118416.241) (106,120878.1679) (110,122483.2934) (114,123961.2312) (118,124002.6073) (122,124135.1517) (126,123395.255) (130,124531.5882) (134,121200.7897) (138,120849.9572) (142,118242.7724) (146,119339.3972) (150,118543.6174) (154,119087.6931) (158,118311.2395) (162,118188.218) (166,118337.4506) (170,114891.8754) (174,115367.6156) (178,115870.4711) (182,117825.1812) (186,118531.9516) (190,118051.6212) (194,118791.8016) (198,118383.6762) (202,118118.5407) (206,117234.0861) (210,117734.8163) (214,115432.2014) (218,116909.422) (222,116669.4794) (226,116824.8455) (230,114573.5752) (234,113533.9375) (238,113395.6208) (242,113601.6056) (246,114174.0607) (250,114525.0538) (254,115978.14) (258,117510.9532) (262,116939.5365) (266,115751.2921) (270,115899.9537) (274,117182.7749) (278,118254.6364) (282,120207.0438) (286,121314.1585) (290,121389.6029) (294,119471.9116) (298,117891.9049) (302,118352.2634) (306,117158.4104) (310,115568.6372) (314,116566.9312) (318,116843.4708) (322,116146.5376) (326,115396.4054) (330,114845.1105) (334,114037.235) (338,113451.1696) (342,117070.5853) (346,116832.1296) (350,116168.9566) (354,115721.0289) (358,116742.3736) (362,116539.0046) (366,115051.5834) (370,114386.9307) (374,116820.6899) (378,116831.5237) (382,117175.5422) (386,117177.5844) (390,115136.5529) (394,115767.768) (398,117484.1191) (402,114571.2509) (406,115552.4338) (410,117275.6559) (414,118591.2583) (418,120923.8866) (422,119065.2178) (426,121656.3921) (430,120631.3623) (434,121320.5489) (438,121394.884) (442,119761.6806) (446,120555.9342) (450,122499.8564) (454,121717.859) (458,121329.2288) (462,122581.8421) (466,120600.0482) (470,120551.2918) (474,122717.4072) (478,119521.5027) (482,119173.1235) (486,118103.5107) (490,120447.9514) (494,121544.9147) (498,120702.9982) (502,120667.5661) (506,120535.9779) (510,120327.2329) (514,119938.5441) (518,118387.2446) (522,120697.8403) (526,119280.3529) (530,120805.4731) (534,122533.9852) (538,123608.3075) (542,120993.6253) (546,120230.6948) (550,120876.7909) (554,125073.2255) (558,125058.2366) (562,124591.7371) (566,125971.7464) (570,127387.8795) (574,126944.9948) (578,128004.736) (582,126744.6846) (586,127201.759) (590,125822.8464) (594,125950.8252) (598,128401.1322) (602,127911.8939) (606,125100.1807) (610,127191.7594) (614,127386.0257) (618,127407.6548) (622,126901.4771) (626,126780.1764) (630,125672.8173) (634,124297.478) (638,122980.5422) (642,124984.5619) (646,125152.7731) (650,123999.4773) (654,124660.5495) (658,125812.1134) (662,128407.8506) (666,126372.6881) (670,126841.7293) (674,128309.136) (678,127690.4552) (682,129833.0032) (686,129564.0522) (690,129756.0797) (694,128883.7967) (698,128530.2078) (702,128363.3795) (706,125787.8673) (710,127169.4808) (714,122889.7777) (718,125501.6661) (722,123996.6519) (726,123737.6159) (730,123981.4468) (734,123741.2512) (738,124715.6585) (742,124375.5924) (746,123934.899) (750,124588.6833) (754,121938.6766) (758,122401.8809) (762,122057.3718) (766,122176.4464) (770,123154.1816) (774,122194.4309) (778,120533.1862) (782,123135.9035) (786,124189.6107) (790,123847.9537) (794,123286.0435) (798,123191.9624) (802,125983.4913) (806,125899.9367) (810,125381.6567) (814,125223.0028) (818,125981.1879) (822,126454.2397) (826,127472.9148) (830,127171.1831) (834,126995.3881) (838,127548.2641) (842,126796.6701) (846,126309.5163) (850,124764.9738) (854,124829.1561) (858,124623.7953) (862,124534.7234) (866,124016.5194) (870,124090.2412) (874,124036.6113) (878,123133.3909) (882,121903.1266) (886,119272.4163) (890,119504.007) (894,120812.6974) (898,122490.9754) (902,124365.0787) (906,123456.9521) (910,124331.5366) (914,125496.0442) (918,125203.4827) (922,124185.8873) (926,125193.2481) (930,123939.3169) (934,125369.8805) (938,123333.8003) (942,125705.5172) (946,123939.3446) (950,123722.9492) (954,123580.529) (958,124216.2576) (962,122803.8844) (966,122196.4914) (970,120630.0803) (974,121658.0554) (978,125279.398) (982,126091.9459) (986,125700.0673) (990,124293.3711) (994,123753.7858) (998,124520.0284)};
\addplot[pubNavy, opacity=0.20, line width=0.35pt] coordinates {(10,109998.4873) (14,109537.4952) (18,109833.2841) (22,110382.1379) (26,111108.0042) (30,111345.5977) (34,112039.3283) (38,110524.9661) (42,110205.9363) (46,110425.3325) (50,111180.6544) (54,110708.8232) (58,112244.9579) (62,111798.2396) (66,113176.009) (70,113116.6196) (74,114251.7809) (78,114397.0709) (82,114442.5963) (86,115344.785) (90,117472.2859) (94,117061.6887) (98,118988.8185) (102,118766.1023) (106,120664.744) (110,120078.2514) (114,120624.6277) (118,122146.3791) (122,123209.9347) (126,123830.0826) (130,124784.3164) (134,124174.6607) (138,124797.3196) (142,126253.4852) (146,126427.261) (150,125683.6832) (154,125652.4072) (158,124256.9307) (162,124409.5879) (166,123791.6217) (170,122171.1194) (174,121787.3594) (178,122121.4918) (182,122127.267) (186,122012.7089) (190,120845.2855) (194,120625.1007) (198,120406.9407) (202,121846.7381) (206,121349.0562) (210,122495.8801) (214,122664.4986) (218,125011.7234) (222,125016.0496) (226,125936.8478) (230,124190.1058) (234,124897.4865) (238,124734.5674) (242,125310.0368) (246,127037.1686) (250,126751.1963) (254,127802.549) (258,126766.8906) (262,126026.282) (266,123295.3161) (270,122442.3344) (274,123342.4527) (278,123797.8415) (282,123922.7728) (286,124535.5915) (290,126280.6296) (294,124288.438) (298,124031.6392) (302,122883.354) (306,123618.1688) (310,122831.6875) (314,123981.5444) (318,124762.2133) (322,124844.0574) (326,124352.0878) (330,123974.4276) (334,123945.1541) (338,122763.5626) (342,124487.4108) (346,122846.9438) (350,123034.7823) (354,122033.64) (358,122018.3071) (362,122438.1934) (366,122576.1291) (370,123475.2616) (374,124940.6902) (378,124499.9211) (382,123994.9764) (386,124298.3042) (390,123263.9738) (394,123481.2981) (398,125083.4661) (402,124285.6311) (406,124308.4002) (410,125464.7297) (414,124819.0679) (418,126917.4305) (422,124206.0494) (426,125065.5613) (430,124114.8096) (434,124302.0931) (438,124063.9145) (442,122878.0138) (446,123259.8914) (450,123978.4151) (454,122637.8764) (458,122351.5763) (462,122294.6281) (466,120107.2978) (470,119474.673) (474,121413.6683) (478,119684.9471) (482,120800.9416) (486,120529.3498) (490,121846.3371) (494,121414.2129) (498,120638.9994) (502,121749.9004) (506,123401.4724) (510,123283.4823) (514,122987.7519) (518,122709.4427) (522,123609.9622) (526,125005.8457) (530,124653.6174) (534,125427.0223) (538,127379.138) (542,126037.3506) (546,126221.0422) (550,125745.3289) (554,126239.9796) (558,127108.1077) (562,126846.2383) (566,127708.4412) (570,128419.468) (574,127545.03) (578,126859.005) (582,126563.1178) (586,125358.2039) (590,123606.2739) (594,123157.1446) (598,125435.2881) (602,125635.399) (606,123151.9144) (610,124528.7191) (614,125329.4543) (618,125589.44) (622,124686.7324) (626,125783.9321) (630,125840.418) (634,126579.3684) (638,126741.8591) (642,128532.9872) (646,127372.1488) (650,127761.3353) (654,128396.4708) (658,128968.2482) (662,128870.0806) (666,127714.9559) (670,128721.236) (674,128569.3268) (678,127980.4158) (682,130486.2972) (686,130874.7232) (690,129740.365) (694,128761.1709) (698,126895.029) (702,127030.4061) (706,126808.4691) (710,128512.1412) (714,127324.3598) (718,127397.1308) (722,125809.2779) (726,126484.1622) (730,125159.8767) (734,125366.6802) (738,125993.9701) (742,127325.3596) (746,125672.7062) (750,127137.0096) (754,125824.7467) (758,125621.6116) (762,125651.0338) (766,126751.0481) (770,127980.2518) (774,127032.3409) (778,126751.9897) (782,127205.1359) (786,127241.8197) (790,127976.3518) (794,128509.9799) (798,129233.9311) (802,130225.5084) (806,130176.1039) (810,130537.7962) (814,131317.7602) (818,131671.2452) (822,130791.0045) (826,130109.0833) (830,130633.4169) (834,131744.6063) (838,130618.8933) (842,130664.1471) (846,129967.4304) (850,128886.3341) (854,130097.306) (858,130566.0133) (862,130257.141) (866,129291.0862) (870,129964.9916) (874,128747.5238) (878,127527.5055) (882,126745.4991) (886,125292.8659) (890,125879.4612) (894,124881.8551) (898,125777.2974) (902,127243.0769) (906,125772.0703) (910,124486.668) (914,124889.0343) (918,124603.552) (922,123495.794) (926,125669.0529) (930,125433.4641) (934,125822.2198) (938,124406.3638) (942,128381.3369) (946,127880.5416) (950,129601.4917) (954,130111.8162) (958,129087.6167) (962,126878.5199) (966,126750.986) (970,127431.8894) (974,129613.7136) (978,129684.3172) (982,129458.0987) (986,129178.7633) (990,127447.0646) (994,125600.7302) (998,125351.59)};
\addplot[pubNavy, opacity=0.20, line width=0.35pt] coordinates {(10,112795.4096) (14,112774.445) (18,112890.1183) (22,113375.4378) (26,115300.8436) (30,115551.1663) (34,115293.4445) (38,113769.6669) (42,112906.8203) (46,111839.7516) (50,114069.1994) (54,112304.9716) (58,113055.5699) (62,110906.0929) (66,112331.4044) (70,110615.2081) (74,111422.895) (78,109532.9685) (82,109643.0131) (86,112104.9273) (90,114651.9984) (94,114440.2887) (98,117223.9849) (102,113823.4217) (106,115176.3664) (110,116833.4089) (114,117746.4409) (118,118694.7231) (122,119384.9016) (126,120603.359) (130,122652.4482) (134,120180.0368) (138,120669.5312) (142,119146.8966) (146,120290.6951) (150,120346.6823) (154,120724.6312) (158,119462.8298) (162,119018.7441) (166,118012.2118) (170,116174.5621) (174,115783.8267) (178,114593.4081) (182,116298.629) (186,116440.5925) (190,115610.1824) (194,117103.6756) (198,117932.1848) (202,118801.4043) (206,118400.9018) (210,119985.4706) (214,120593.4713) (218,121534.4591) (222,122150.4704) (226,124534.2835) (230,123401.4154) (234,123211.4676) (238,123218.3227) (242,123579.4891) (246,123916.5355) (250,122808.4288) (254,125602.9012) (258,125703.8797) (262,124572.5981) (266,124303.0666) (270,122969.7371) (274,122957.698) (278,123140.112) (282,123975.1788) (286,123930.6753) (290,124557.3897) (294,122070.7794) (298,121936.1001) (302,121867.0705) (306,122120.2077) (310,120309.7443) (314,121782.1196) (318,122513.1045) (322,123260.1859) (326,122887.9266) (330,122698.5181) (334,123223.9463) (338,122897.6134) (342,125685.2799) (346,123802.0362) (350,123851.5581) (354,122477.3565) (358,123326.493) (362,125715.2186) (366,125618.194) (370,124396.1208) (374,125066.658) (378,125513.7061) (382,125159.2539) (386,125196.1636) (390,124650.6868) (394,125981.8187) (398,128083.2757) (402,125546.5845) (406,124669.4711) (410,125532.6178) (414,125119.6589) (418,128089.0378) (422,127194.4643) (426,128502.3986) (430,127080.648) (434,127712.4552) (438,128470.9077) (442,127773.5523) (446,127579.4345) (450,128158.8612) (454,128002.7135) (458,128318.0443) (462,131377.0666) (466,127834.0065) (470,126012.7975) (474,130247.0242) (478,127769.8116) (482,129769.4271) (486,127997.1862) (490,129570.6053) (494,130923.0283) (498,129774.0552) (502,129201.4623) (506,130203.4659) (510,127830.6436) (514,126871.8357) (518,125233.2976) (522,126370.4704) (526,126466.1229) (530,125872.0989) (534,125597.6979) (538,127204.9574) (542,123413.9482) (546,122217.9203) (550,122434.2068) (554,124635.4315) (558,123986.551) (562,123988.1482) (566,126283.809) (570,126282.2949) (574,124376.1202) (578,123412.8282) (582,121974.737) (586,121901.1318) (590,122027.4487) (594,121287.5418) (598,124028.2552) (602,124897.9488) (606,122829.023) (610,124365.0945) (614,123988.7136) (618,124293.5151) (622,124976.5845) (626,126470.2358) (630,126468.2919) (634,126170.206) (638,125368.9432) (642,127516.2647) (646,125162.2264) (650,123438.0155) (654,123704.9581) (658,124254.4741) (662,125588.2864) (666,124207.067) (670,125315.0707) (674,128175.1983) (678,126546.9229) (682,127021.9684) (686,127899.4459) (690,127594.7988) (694,127682.671) (698,127009.6588) (702,126665.7516) (706,126161.3327) (710,125596.6243) (714,123865.503) (718,125092.2617) (722,122954.4296) (726,122611.9633) (730,122992.5901) (734,122945.4985) (738,122736.3714) (742,122740.6741) (746,120972.3162) (750,121171.6993) (754,120489.0489) (758,119156.7277) (762,119918.3822) (766,120917.4863) (770,122664.8295) (774,121034.1747) (778,120064.3501) (782,119457.7584) (786,120969.1365) (790,120736.6232) (794,120805.9788) (798,121871.6173) (802,123861.2925) (806,124068.2036) (810,124171.0208) (814,123462.0222) (818,123887.5437) (822,124198.9909) (826,125036.0028) (830,125627.0635) (834,128545.249) (838,127851.0836) (842,128496.1121) (846,129089.9488) (850,128182.2505) (854,128085.803) (858,126983.3205) (862,127244.3454) (866,125027.8279) (870,123870.9449) (874,123588.1623) (878,122917.7491) (882,120510.7563) (886,117560.9854) (890,118995.8076) (894,119312.9292) (898,120420.6598) (902,121852.0988) (906,122292.1035) (910,122863.291) (914,123526.6769) (918,124772.7147) (922,124601.8219) (926,124072.3305) (930,123342.0566) (934,122692.877) (938,120377.3396) (942,122136.828) (946,120444.0491) (950,121427.3223) (954,121660.0102) (958,122023.8462) (962,121026.3345) (966,119853.8644) (970,119216.4964) (974,120819.893) (978,122933.1295) (982,123864.7429) (986,125049.0901) (990,125226.7359) (994,124415.1411) (998,124578.9968)};
\addplot[pubNavy, opacity=0.20, line width=0.35pt] coordinates {(10,112196.5344) (14,111103.1993) (18,112632.2989) (22,112670.9221) (26,114150.4401) (30,113941.2285) (34,114622.0303) (38,113450.6878) (42,114010.9628) (46,114071.751) (50,115439.0803) (54,115207.7941) (58,115362.4475) (62,114373.7881) (66,116145.0955) (70,116410.9902) (74,117935.924) (78,117517.6221) (82,117494.7902) (86,118643.5917) (90,120476.0028) (94,119210.9546) (98,120165.7906) (102,119253.4635) (106,121358.3291) (110,120315.4029) (114,121487.0389) (118,121595.6266) (122,122320.1731) (126,122064.058) (130,123230.401) (134,121845.0854) (138,121760.6099) (142,122189.0127) (146,122382.1488) (150,121767.999) (154,123443.2874) (158,121643.3508) (162,123042.0042) (166,122208.4114) (170,120704.1364) (174,121030.5583) (178,120241.3576) (182,121869.8747) (186,121686.0822) (190,120462.8303) (194,120977.9127) (198,121061.5596) (202,120410.5412) (206,118096.4108) (210,119404.3242) (214,116298.5285) (218,117950.2783) (222,117311.4016) (226,117515.1726) (230,116099.8073) (234,115637.4556) (238,115333.7396) (242,115308.3836) (246,116826.426) (250,115645.7134) (254,117960.6275) (258,118884.7945) (262,118270.8804) (266,117306.7391) (270,116470.2983) (274,117392.2094) (278,117064.2086) (282,117124.1425) (286,117264.9612) (290,117377.3169) (294,115311.6828) (298,114634.4016) (302,115898.4687) (306,116053.0653) (310,115742.8504) (314,117215.0822) (318,118455.2139) (322,117769.1991) (326,118438.4457) (330,117837.7669) (334,118324.771) (338,118815.5736) (342,121616.3488) (346,121307.715) (350,120790.5201) (354,118069.7025) (358,118290.0177) (362,119140.4576) (366,118087.9899) (370,118522.0303) (374,120697.8219) (378,121014.39) (382,121660.284) (386,120833.7852) (390,119539.7847) (394,119682.2768) (398,121036.9963) (402,119736.8744) (406,120185.6119) (410,121581.9976) (414,121209.09) (418,122733.4916) (422,120810.895) (426,121150.5885) (430,120423.7758) (434,121412.0306) (438,121418.1389) (442,119959.1906) (446,121115.7014) (450,122624.3531) (454,121881.2646) (458,122897.3014) (462,124126.6096) (466,122615.5815) (470,121562.9452) (474,122835.8778) (478,122090.7639) (482,123437.3611) (486,123202.8983) (490,124744.4403) (494,124420.4786) (498,123757.6931) (502,124115.2644) (506,124095.4231) (510,122864.5674) (514,121221.2013) (518,121627.6187) (522,122427.462) (526,122572.5284) (530,121796.3014) (534,121897.625) (538,121653.8649) (542,121305.2455) (546,120997.3469) (550,121477.8206) (554,121924.1127) (558,122894.5472) (562,123222.7132) (566,122826.6313) (570,123056.9697) (574,121021.008) (578,120952.0805) (582,120742.8405) (586,120642.262) (590,119528.6919) (594,119030.875) (598,120080.919) (602,119325.7436) (606,119526.0303) (610,119590.0463) (614,120388.4257) (618,120972.8261) (622,121236.4011) (626,121880.7944) (630,122341.29) (634,121301.5233) (638,121561.2128) (642,122429.4156) (646,121989.7063) (650,122419.3248) (654,121421.0915) (658,121292.2678) (662,122063.6013) (666,121354.9784) (670,120602.3904) (674,120223.3188) (678,119470.3752) (682,121372.3703) (686,121579.0742) (690,121973.5525) (694,121736.6756) (698,121902.0232) (702,121085.0944) (706,121195.6053) (710,121974.7249) (714,119141.4045) (718,120431.0914) (722,121255.1677) (726,121687.3493) (730,121306.5282) (734,121423.9578) (738,121853.5826) (742,122540.7586) (746,121479.2137) (750,120552.9753) (754,119150.4423) (758,119378.3357) (762,120635.2876) (766,121708.4713) (770,121747.2383) (774,120104.3936) (778,119789.3248) (782,119994.5746) (786,120216.1546) (790,119023.8538) (794,118764.5122) (798,119134.5795) (802,121137.5136) (806,120821.7913) (810,120012.5091) (814,119219.5259) (818,119837.9155) (822,120142.4218) (826,120958.873) (830,121846.1283) (834,123837.2758) (838,123010.3824) (842,123702.3551) (846,123757.0735) (850,124059.4416) (854,123652.0167) (858,123363.3458) (862,124196.3253) (866,123252.9503) (870,123360.7778) (874,122545.8291) (878,121117.8091) (882,119555.1138) (886,117612.5388) (890,118636.115) (894,118357.3396) (898,118576.095) (902,119173.6608) (906,119295.9867) (910,119663.6864) (914,119893.3622) (918,120183.9222) (922,119599.4945) (926,119575.7511) (930,120293.3082) (934,121268.4906) (938,119315.0374) (942,120143.5087) (946,118888.5788) (950,120205.4236) (954,120853.681) (958,121174.7129) (962,119831.3402) (966,119506.1039) (970,119491.8783) (974,120161.8753) (978,121182.1451) (982,120770.2062) (986,120381.1706) (990,120312.1201) (994,121154.1206) (998,121507.6298)};
\addplot[pubNavy, opacity=0.20, line width=0.35pt] coordinates {(10,111571.977) (14,113201.1552) (18,112569.0227) (22,113870.6167) (26,115740.5692) (30,113965.3783) (34,115107.3384) (38,113342.734) (42,113381.2259) (46,112471.3323) (50,114939.8581) (54,113517.7585) (58,114505.3136) (62,112730.839) (66,114142.3723) (70,112927.4747) (74,113880.4021) (78,113416.4028) (82,113823.2844) (86,116258.1198) (90,119026.8782) (94,119121.3347) (98,123077.0698) (102,120007.2468) (106,122277.9869) (110,123325.1609) (114,125550.8019) (118,126355.7013) (122,127627.2288) (126,128944.5369) (130,131073.9134) (134,127885.5747) (138,129246.009) (142,126748.6977) (146,127727.1443) (150,126464.1661) (154,127175.8813) (158,125401.9394) (162,125008.1563) (166,124267.6094) (170,121804.2903) (174,119295.7429) (178,118083.0383) (182,119149.463) (186,119259.4258) (190,118442.5216) (194,121318.895) (198,121616.421) (202,122344.6768) (206,120835.9401) (210,122583.7173) (214,121335.2888) (218,122806.8069) (222,123710.3123) (226,125740.4651) (230,125514.5787) (234,125383.2524) (238,126019.6372) (242,123851.2831) (246,124341.2382) (250,123610.3832) (254,125692.2174) (258,126449.4372) (262,126695.1626) (266,126448.1995) (270,125674.0885) (274,126566.4411) (278,125862.3923) (282,127058.2194) (286,127069.7648) (290,128834.8139) (294,127087.9154) (298,126714.2059) (302,128384.7403) (306,128311.6205) (310,127086.7076) (314,127960.1968) (318,127224.1273) (322,128273.2075) (326,127277.6332) (330,128985.5217) (334,129603.6787) (338,129334.2407) (342,131105.1653) (346,129446.2031) (350,129095.8978) (354,125001.7045) (358,126321.683) (362,128821.019) (366,129167.8894) (370,128453.4949) (374,130154.3249) (378,130231.799) (382,129619.9705) (386,129444.9157) (390,128833.1564) (394,129947.3861) (398,131051.7385) (402,127898.5321) (406,127738.8195) (410,129087.1949) (414,129078.281) (418,130930.8907) (422,128432.0926) (426,130863.7084) (430,127884.6882) (434,128558.096) (438,130181.475) (442,128537.7343) (446,129377.085) (450,129973.0113) (454,129621.6784) (458,129798.7821) (462,130575.7321) (466,128996.1039) (470,127658.9658) (474,130450.812) (478,127797.9979) (482,129444.3011) (486,126693.6479) (490,129476.2881) (494,129930.1984) (498,129090.9983) (502,128833.8279) (506,130315.1122) (510,131222.1015) (514,129366.6259) (518,128547.8269) (522,130918.4823) (526,130778.1534) (530,131407.135) (534,131460.5866) (538,133544.6722) (542,129525.7473) (546,128336.3509) (550,128486.2577) (554,129840.8195) (558,127210.2907) (562,127882.135) (566,130634.499) (570,129231.2769) (574,127976.1802) (578,127177.0618) (582,126195.9579) (586,125994.3093) (590,126390.7182) (594,127930.1519) (598,131511.0561) (602,133512.5005) (606,132720.062) (610,134350.0807) (614,134731.8514) (618,134199.0187) (622,134160.3719) (626,134901.5006) (630,134656.3386) (634,132985.3409) (638,132418.341) (642,132324.1634) (646,129437.7179) (650,128505.3887) (654,127820.1078) (658,127387.0923) (662,129782.6713) (666,126972.8926) (670,128717.4947) (674,130588.3467) (678,128722.6784) (682,131032.4342) (686,131682.5305) (690,132276.3438) (694,133373.8503) (698,132118.2718) (702,130747.6997) (706,129613.5626) (710,130133.1331) (714,127570.2236) (718,129824.5247) (722,128009.1585) (726,127634.1884) (730,127590.9886) (734,127431.6864) (738,127770.5516) (742,125647.3262) (746,124209.546) (750,124477.0187) (754,123889.5521) (758,123461.6583) (762,122710.4436) (766,122805.075) (770,124064.1799) (774,122336.7605) (778,120939.9555) (782,121815.2005) (786,123610.0795) (790,122821.3878) (794,123757.1924) (798,124705.254) (802,127805.3817) (806,127095.6171) (810,128619.2231) (814,129107.5555) (818,130104.5061) (822,131732.847) (826,132639.4555) (830,132227.6419) (834,133242.7467) (838,134204.2277) (842,134556.4274) (846,134817.7935) (850,133206.1553) (854,133835.3593) (858,132650.847) (862,133085.1587) (866,131499.5048) (870,130071.7131) (874,128418.448) (878,127828.5318) (882,127164.8717) (886,123605.6788) (890,123952.1055) (894,123421.2489) (898,124284.923) (902,126142.2889) (906,124359.3489) (910,124797.1614) (914,124927.6371) (918,124818.4201) (922,125095.1319) (926,126192.7282) (930,125510.9008) (934,126598.8804) (938,125422.9513) (942,129080.2573) (946,127699.0436) (950,126487.0369) (954,126976.0922) (958,126832.387) (962,126298.43) (966,125851.5547) (970,125574.2677) (974,127073.8011) (978,129501.5475) (982,129491.592) (986,129494.4641) (990,127163.3572) (994,126039.9687) (998,127177.0318)};
\addplot[pubNavy, opacity=0.20, line width=0.35pt] coordinates {(10,113937.3014) (14,111831.9833) (18,110112.3707) (22,109299.2365) (26,109511.2349) (30,109989.5149) (34,112040.3853) (38,111063.9209) (42,111091.1225) (46,111196.2523) (50,113461.3202) (54,112424.8835) (58,113195.0006) (62,111184.2504) (66,114114.6891) (70,114764.5135) (74,117107.2117) (78,116859.4418) (82,117090.7936) (86,118932.7074) (90,120998.9509) (94,120536.4002) (98,122783.9829) (102,121800.526) (106,124313.9143) (110,124681.9096) (114,126302.8182) (118,126300.8706) (122,127156.7727) (126,127300.1096) (130,127989.3527) (134,125250.2199) (138,125833.638) (142,125864.1882) (146,125529.8105) (150,123249.9317) (154,123785.8302) (158,121958.6198) (162,122461.0542) (166,120938.6629) (170,119406.2023) (174,118951.7035) (178,118304.2449) (182,119757.197) (186,120541.8114) (190,120215.0732) (194,121042.8651) (198,121482.7995) (202,122835.3678) (206,122781.7621) (210,124867.0525) (214,123995.7369) (218,126079.1054) (222,125981.7191) (226,126890.4145) (230,125822.6159) (234,125456.1428) (238,124930.5888) (242,124451.8783) (246,125247.7736) (250,124774.7796) (254,126196.5256) (258,124543.3099) (262,123696.6938) (266,122297.9499) (270,121478.4476) (274,121831.0111) (278,122128.0022) (282,122556.0305) (286,122401.1298) (290,124320.1689) (294,122769.3708) (298,122816.011) (302,122962.16) (306,123485.885) (310,122972.5756) (314,123921.806) (318,123609.5725) (322,123756.9328) (326,123216.0237) (330,123141.8099) (334,123138.5517) (338,122358.8748) (342,124548.1769) (346,122956.9224) (350,121103.3904) (354,119699.217) (358,119788.9512) (362,120615.0209) (366,119542.0616) (370,119928.0048) (374,121457.437) (378,120800.6767) (382,120164.884) (386,120671.3327) (390,119391.2009) (394,119872.2391) (398,122703.4735) (402,122045.0084) (406,121582.3943) (410,123140.8598) (414,123789.3935) (418,125877.2711) (422,123983.6956) (426,125997.4331) (430,124457.5891) (434,124842.0445) (438,126034.125) (442,125222.7958) (446,125986.178) (450,126940.2986) (454,126442.0694) (458,127178.997) (462,129519.8118) (466,127791.6282) (470,127341.2809) (474,129205.9564) (478,126860.5895) (482,128170.1283) (486,127295.4126) (490,128245.8186) (494,128092.4152) (498,127987.2237) (502,128217.5315) (506,128372.3529) (510,127852.8584) (514,125536.9048) (518,124872.6602) (522,127456.192) (526,127300.4113) (530,126613.4188) (534,127957.9813) (538,128978.0656) (542,125951.3299) (546,126415.8341) (550,126729.8882) (554,129207.426) (558,127119.6887) (562,126065.9215) (566,127263.7851) (570,127027.9422) (574,126084.2081) (578,126114.1111) (582,125327.6635) (586,124634.3127) (590,124530.7477) (594,124594.7553) (598,126394.9471) (602,126165.8198) (606,124396.5475) (610,126828.7056) (614,127814.0716) (618,128699.914) (622,127505.2539) (626,128775.5211) (630,128522.9264) (634,128094.8162) (638,127886.5575) (642,130129.4453) (646,129023.3615) (650,127942.2045) (654,127857.7893) (658,127975.7491) (662,129231.5134) (666,126771.3942) (670,126820.1563) (674,127585.4745) (678,126137.3112) (682,127184.6907) (686,126952.3493) (690,127561.2586) (694,127089.0453) (698,128103.3328) (702,128359.0402) (706,128377.5233) (710,129222.5153) (714,127065.0358) (718,127451.0656) (722,127379.9988) (726,128200.6484) (730,127883.6862) (734,127341.0657) (738,128205.6781) (742,128575.1174) (746,126006.896) (750,125158.0332) (754,123514.2811) (758,124676.1368) (762,124193.8379) (766,126269.5986) (770,126512.7313) (774,126703.8498) (778,125331.6591) (782,127668.1877) (786,127698.7106) (790,125047.2697) (794,124773.2072) (798,124779.3405) (802,126073.9903) (806,126157.2697) (810,125466.0511) (814,125947.5505) (818,125955.7852) (822,125345.0389) (826,126698.1697) (830,126417.6998) (834,126431.8059) (838,127087.8733) (842,128085.2536) (846,127382.9444) (850,127872.2927) (854,127998.3065) (858,127042.0022) (862,128011.4068) (866,126814.4169) (870,126799.7622) (874,127409.3512) (878,126227.8483) (882,124688.8567) (886,122837.4333) (890,123880.0131) (894,124005.5718) (898,123439.3944) (902,124352.2835) (906,124180.09) (910,123465.5491) (914,123245.0675) (918,124175.3773) (922,123068.8675) (926,122699.6639) (930,123038.7916) (934,124606.8627) (938,122625.0119) (942,124425.3738) (946,123338.4906) (950,125041.6934) (954,125196.1884) (958,124942.2475) (962,124396.5582) (966,122912.2187) (970,122788.782) (974,125080.7941) (978,128501.5985) (982,128058.3896) (986,128643.3315) (990,128382.4958) (994,128559.5684) (998,130086.6505)};
\addplot[pubNavy, opacity=0.20, line width=0.35pt] coordinates {(10,119104.4588) (14,117395.3597) (18,113526.7939) (22,112795.2035) (26,113552.2972) (30,112303.3304) (34,113451.6103) (38,112717.6835) (42,112896.4759) (46,112975.0506) (50,113878.9857) (54,113414.9842) (58,113486.4195) (62,111621.2089) (66,114433.6875) (70,116031.9526) (74,118597.5788) (78,118907.3407) (82,119565.5726) (86,120387.2913) (90,122045.1632) (94,122006.3755) (98,123314.8559) (102,123691.7384) (106,126435.4917) (110,125917.9212) (114,126984.3167) (118,126974.8549) (122,127388.7728) (126,126784.3803) (130,127870.6279) (134,125756.4391) (138,126979.2004) (142,127883.7206) (146,128628.0462) (150,125812.7149) (154,127616.9852) (158,125345.836) (162,126511.4657) (166,126985.5655) (170,125422.2121) (174,125407.16) (178,124767.3469) (182,125983.2165) (186,125347.0691) (190,124969.682) (194,125953.6326) (198,126244.0176) (202,127494.8627) (206,125618.347) (210,125734.1825) (214,123061.0717) (218,124288.771) (222,124128.2596) (226,125226.4413) (230,123729.9575) (234,124206.3992) (238,124018.5987) (242,123614.3291) (246,123080.771) (250,122845.491) (254,124171.3084) (258,124811.2795) (262,123849.0847) (266,122202.5097) (270,121490.5122) (274,121510.3806) (278,121000.5114) (282,122408.6579) (286,123490.4851) (290,124158.5184) (294,122471.4304) (298,123147.3118) (302,124766.3477) (306,124888.432) (310,124464.937) (314,124831.2916) (318,125633.529) (322,125746.4977) (326,125589.9842) (330,125474.7329) (334,124963.8142) (338,124781.3145) (342,127094.4357) (346,125026.0811) (350,125053.3624) (354,121417.0491) (358,121540.5122) (362,122575.9929) (366,122556.5795) (370,123848.3817) (374,125941.1564) (378,125362.9751) (382,124332.3734) (386,123634.2077) (390,122270.3129) (394,120901.7537) (398,122242.8024) (402,121451.7715) (406,120449.095) (410,122142.0293) (414,123279.2665) (418,123962.9252) (422,120867.7972) (426,121553.7756) (430,121535.1844) (434,123332.23) (438,123232.7367) (442,122194.926) (446,124299.6958) (450,126025.0189) (454,125685.6635) (458,126743.2963) (462,127255.6092) (466,125358.135) (470,124527.5862) (474,125236.2624) (478,123772.0247) (482,123376.4896) (486,122617.7319) (490,124397.1494) (494,124514.6206) (498,122651.4857) (502,123776.0214) (506,124194.8855) (510,123756.0325) (514,122039.9886) (518,121842.2071) (522,123732.4707) (526,123627.986) (530,123723.097) (534,124369.1651) (538,126421.5589) (542,124020.0862) (546,123429.7373) (550,124238.5957) (554,125301.3552) (558,125289.4049) (562,125529.9113) (566,126510.1823) (570,126978.2581) (574,127451.0724) (578,127239.9261) (582,126974.6936) (586,126566.7391) (590,124551.0361) (594,125899.2691) (598,127149.421) (602,124892.5372) (606,123489.2253) (610,124678.417) (614,125095.6676) (618,126547.5533) (622,124981.8012) (626,124702.1238) (630,125306.997) (634,124991.8721) (638,125207.7447) (642,127156.0249) (646,125620.6538) (650,126635.913) (654,129584.5074) (658,131839.728) (662,133169.9602) (666,130515.2222) (670,130528.6722) (674,131243.4316) (678,132002.7236) (682,135558.6313) (686,135291.7639) (690,134812.0693) (694,134520.4894) (698,133020.9592) (702,131382.1277) (706,129546.2739) (710,130689.5922) (714,128055.8515) (718,129636.0178) (722,129467.9535) (726,128280.0942) (730,126430.6048) (734,124721.3899) (738,125451.29) (742,124854.3297) (746,122965.4817) (750,124300.1712) (754,123066.4908) (758,122750.4736) (762,122679.9601) (766,124787.9955) (770,124820.1326) (774,125316.6484) (778,126000.7014) (782,126133.479) (786,128327.5528) (790,128165.6042) (794,128589.9292) (798,128879.204) (802,131953.7632) (806,132811.3773) (810,133396.3178) (814,133694.278) (818,133878.2315) (822,133480.7254) (826,132698.7302) (830,132872.4255) (834,133791.658) (838,134592.2541) (842,134930.2188) (846,134579.0543) (850,133218.2846) (854,131459.1454) (858,129852.0022) (862,130788.0562) (866,128700.5324) (870,129668.9583) (874,128719.158) (878,127713.396) (882,126416.1212) (886,123915.3074) (890,123924.38) (894,124252.6807) (898,126097.0007) (902,127378.8325) (906,128162.1878) (910,126266.4223) (914,125808.6173) (918,127209.9849) (922,126986.5343) (926,127920.7259) (930,128789.3271) (934,128698.8928) (938,126352.2329) (942,127555.6987) (946,126341.2789) (950,127332.4051) (954,127576.2347) (958,127723.4053) (962,126250.6544) (966,125856.891) (970,125130.114) (974,127083.247) (978,129843.2968) (982,128965.8375) (986,129065.2291) (990,128534.8447) (994,125958.9941) (998,127082.7942)};
\addplot[pubNavy, opacity=0.20, line width=0.35pt] coordinates {(10,116775.4857) (14,116142.6485) (18,112449.9319) (22,112724.7326) (26,114119.1148) (30,112663.7579) (34,114179.8861) (38,114418.9313) (42,113531.704) (46,113414.1107) (50,114915.1576) (54,113440.2311) (58,113665.853) (62,111246.5798) (66,114264.0029) (70,114957.355) (74,116382.4361) (78,116345.9136) (82,116378.5933) (86,117194.413) (90,118820.8055) (94,118495.5165) (98,119952.0536) (102,119113.6119) (106,121270.4567) (110,121511.0456) (114,122974.7854) (118,122196.1797) (122,122726.7195) (126,122707.0766) (130,124123.6013) (134,121430.6709) (138,121639.5944) (142,121002.2967) (146,121820.4623) (150,120057.3973) (154,122185.7612) (158,119983.6495) (162,119107.5028) (166,120059.4978) (170,118174.6897) (174,117775.2502) (178,116372.948) (182,118323.5222) (186,119035.3479) (190,118856.759) (194,120422.0804) (198,120320.1435) (202,119880.4759) (206,117485.8178) (210,118950.9633) (214,117035.4873) (218,118541.6073) (222,118764.9854) (226,119276.369) (230,118588.8786) (234,118128.2715) (238,117841.5229) (242,116876.2014) (246,118977.5087) (250,119174.9181) (254,121265.9031) (258,121935.4759) (262,120946.8423) (266,118986.2893) (270,117904.1542) (274,118668.0134) (278,118209.2455) (282,119413.9254) (286,118408.8617) (290,118916.7266) (294,116204.659) (298,114779.819) (302,114241.2289) (306,113711.7005) (310,113314.8874) (314,114631.8334) (318,115492.037) (322,115608.7007) (326,114875.1605) (330,115368.4902) (334,116113.808) (338,117061.9648) (342,119664.3127) (346,118884.4461) (350,117841.4523) (354,118072.1547) (358,119314.9258) (362,119824.7951) (366,119559.1676) (370,120105.627) (374,121219.4192) (378,121280.4147) (382,120604.1899) (386,120002.5857) (390,118174.8843) (394,117724.5029) (398,119892.6678) (402,118675.2054) (406,117286.7614) (410,119291.8637) (414,120308.1783) (418,121752.875) (422,119971.2194) (426,121491.3141) (430,120553.6045) (434,121383.4303) (438,121265.8266) (442,120229.7878) (446,121446.6914) (450,121603.3913) (454,120838.2385) (458,122836.7008) (462,123543.5455) (466,122687.1748) (470,121074.6937) (474,121323.9166) (478,119599.5642) (482,120747.6844) (486,119486.5722) (490,119438.4285) (494,120323.5441) (498,119587.5563) (502,120946.8097) (506,120441.2075) (510,119004.5643) (514,116515.8864) (518,116351.7257) (522,119229.087) (526,119792.3904) (530,120389.9419) (534,121364.4139) (538,124859.6231) (542,122040.6047) (546,121441.3714) (550,122469.9862) (554,125539.6399) (558,125332.5743) (562,125792.4529) (566,127741.9733) (570,126090.4485) (574,126168.6258) (578,125228.3224) (582,125815.7369) (586,124420.2596) (590,124317.958) (594,126487.4211) (598,128841.4514) (602,129018.2694) (606,126908.0704) (610,128272.1974) (614,128745.0568) (618,129092.0473) (622,129674.8349) (626,129926.8775) (630,128941.3632) (634,128751.9932) (638,128356.8236) (642,129998.6829) (646,127750.1423) (650,127555.8755) (654,128317.9251) (658,129480.7319) (662,131158.1953) (666,128157.8599) (670,129084.0416) (674,129904.2524) (678,129438.3104) (682,130582.164) (686,132215.1156) (690,131789.3409) (694,130586.6262) (698,130205.6952) (702,128250.9479) (706,126019.8481) (710,126851.3994) (714,123207.9833) (718,125630.7528) (722,125493.0852) (726,125192.3164) (730,124459.1216) (734,124411.3155) (738,124046.5315) (742,123939.1783) (746,122372.1873) (750,122123.8553) (754,121303.6774) (758,121807.9492) (762,121468.1609) (766,122155.1997) (770,122411.8287) (774,121119.9196) (778,119960.7477) (782,120913.5562) (786,121835.3055) (790,121778.6596) (794,121378.2309) (798,122687.7121) (802,127546.7986) (806,128519.8682) (810,128969.3248) (814,129558.3244) (818,129925.3835) (822,129908.5847) (826,129836.9124) (830,128955.656) (834,131291.007) (838,132001.277) (842,131507.5786) (846,131200.2186) (850,129684.8764) (854,128478.3488) (858,125958.3147) (862,128332.3641) (866,125560.5144) (870,127208.9823) (874,126732.2767) (878,126085.3175) (882,124922.0104) (886,121449.1461) (890,122937.2096) (894,123762.594) (898,124701.0282) (902,125687.9067) (906,125270.6107) (910,124343.0468) (914,124418.9493) (918,124919.5913) (922,124196.141) (926,126044.1241) (930,126370.9362) (934,126917.5512) (938,124106.3568) (942,126383.991) (946,124117.5665) (950,124540.7076) (954,124626.1128) (958,124950.3458) (962,124357.5156) (966,122875.5526) (970,121886.2223) (974,123450.9011) (978,125543.5031) (982,123666.6759) (986,125809.685) (990,125686.2003) (994,123748.7927) (998,124483.1547)};
\addplot[pubNavy, opacity=0.20, line width=0.35pt] coordinates {(10,116810.4541) (14,114488.4692) (18,112517.1053) (22,111688.1843) (26,113282.5166) (30,112953.0617) (34,114885.6727) (38,112738.0839) (42,112419.5724) (46,111412.7396) (50,112583.0006) (54,111505.2542) (58,112067.1372) (62,111221.3191) (66,112727.1116) (70,112004.5385) (74,113142.5815) (78,112868.1538) (82,112741.9707) (86,115069.7617) (90,117744.7432) (94,117753.1456) (98,120263.7687) (102,118479.9927) (106,119560.4486) (110,119258.1091) (114,120099.3869) (118,120685.4875) (122,120446.212) (126,120824.0255) (130,121809.6641) (134,118181.9255) (138,117552.4704) (142,116831.1311) (146,117516.4819) (150,116052.4414) (154,117249.8654) (158,116821.8592) (162,117121.6944) (166,116630.1242) (170,115941.9935) (174,116525.4792) (178,116210.5434) (182,118250.2681) (186,118918.3425) (190,118257.2968) (194,119583.0081) (198,119924.4625) (202,120333.2789) (206,118351.0391) (210,118988.9173) (214,119366.9085) (218,120029.6658) (222,119938.5868) (226,120029.6896) (230,119455.8304) (234,119029.9008) (238,119714.9348) (242,119039.0937) (246,121030.7199) (250,120588.7698) (254,121830.7507) (258,122685.5561) (262,121124.0957) (266,119739.4382) (270,118734.8423) (274,118713.4991) (278,118764.7329) (282,119646.9099) (286,118971.6496) (290,119566.6063) (294,118407.6042) (298,116929.1519) (302,117357.8007) (306,117586.8358) (310,117478.5356) (314,118434.0923) (318,119882.4747) (322,119669.646) (326,120276.0697) (330,120817.0178) (334,120267.1493) (338,119268.1061) (342,121597.1645) (346,121338.8161) (350,121180.5777) (354,118592.5751) (358,119674.3622) (362,119312.6415) (366,118827.1149) (370,118353.5945) (374,119332.79) (378,120357.3436) (382,120146.8402) (386,119402.6146) (390,118151.0903) (394,117671.5533) (398,118941.6261) (402,118103.779) (406,116755.6353) (410,116623.5693) (414,118665.4975) (418,119929.8863) (422,118479.5013) (426,119405.8678) (430,118473.3087) (434,118818.7248) (438,118752.03) (442,117765.4061) (446,119198.2978) (450,120218.5627) (454,119555.6913) (458,120698.5768) (462,123166.6214) (466,121421.4497) (470,120557.0223) (474,121898.8807) (478,119920.5733) (482,120629.3808) (486,120416.804) (490,121727.2161) (494,122368.7159) (498,121920.4131) (502,122389.6822) (506,122576.5347) (510,121896.4292) (514,119398.9731) (518,118925.1135) (522,119164.89) (526,118592.9751) (530,119724.1056) (534,119778.3329) (538,122214.527) (542,117943.2398) (546,116633.6713) (550,116537.4307) (554,117708.3609) (558,117947.2907) (562,120380.0727) (566,121991.0293) (570,122212.3008) (574,123875.3999) (578,121735.5159) (582,121066.9361) (586,120480.0881) (590,119940.4293) (594,120395.1475) (598,123084.3614) (602,121729.644) (606,120458.6593) (610,120833.2861) (614,119695.5201) (618,119325.2747) (622,118559.2167) (626,119102.8036) (630,119821.0546) (634,119483.8877) (638,119670.8804) (642,120291.2403) (646,118811.2818) (650,118466.6617) (654,119359.579) (658,119398.7302) (662,121294.1464) (666,120016.3188) (670,121014.9992) (674,121465.5563) (678,120638.8932) (682,122776.8591) (686,124507.866) (690,123940.2344) (694,125971.8729) (698,124629.5396) (702,123760.4934) (706,124271.5075) (710,124404.2833) (714,122083.2618) (718,123707.0306) (722,123248.1189) (726,122681.5646) (730,121327.9883) (734,119948.4118) (738,119030.2967) (742,117734.8883) (746,116539.8094) (750,117030.1594) (754,117121.4587) (758,117687.1677) (762,117766.4844) (766,118611.6197) (770,118403.5979) (774,118091.0268) (778,117941.8844) (782,117866.0733) (786,120596.8295) (790,121013.1336) (794,122200.94) (798,121926.3733) (802,122854.0722) (806,122873.8214) (810,122877.4272) (814,123553.6068) (818,123934.5993) (822,123457.2325) (826,123183.8514) (830,124647.9376) (834,125133.7096) (838,124042.2851) (842,124217.3044) (846,124578.8129) (850,124476.0064) (854,124366.25) (858,124216.5183) (862,125035.4416) (866,123420.0958) (870,124782.0919) (874,124285.2647) (878,123779.2477) (882,121485.7516) (886,119792.3584) (890,121381.4115) (894,120455.3535) (898,120648.6211) (902,121496.0619) (906,120308.7648) (910,120013.9034) (914,120051.465) (918,120189.7948) (922,120388.9766) (926,121736.195) (930,122788.7573) (934,123086.6973) (938,121061.0849) (942,122021.1491) (946,120656.3574) (950,120313.1548) (954,121629.6912) (958,121911.3133) (962,120953.6001) (966,119760.7389) (970,120039.9232) (974,121125.2787) (978,121604.9804) (982,120848.0863) (986,121073.6248) (990,121997.9351) (994,122306.5075) (998,123596.2803)};
\addplot[pubNavy, opacity=0.20, line width=0.35pt] coordinates {(10,111652.2202) (14,113380.6144) (18,112200.2797) (22,111964.0671) (26,113030.6383) (30,111613.4482) (34,113065.0795) (38,111781.7219) (42,111580.9677) (46,111997.0334) (50,113238.6026) (54,112956.9183) (58,113191.9615) (62,111526.9615) (66,113231.4109) (70,113246.2579) (74,114416.6938) (78,114326.3464) (82,114442.3771) (86,115418.8999) (90,117027.5049) (94,115780.6977) (98,117095.2969) (102,116225.0307) (106,117906.7367) (110,117597.5922) (114,118299.4734) (118,118662.1876) (122,119275.0735) (126,119826.9662) (130,120756.2498) (134,119838.8211) (138,119530.4802) (142,120506.6897) (146,120833.288) (150,120005.1804) (154,120436.0876) (158,119429.1811) (162,120401.0596) (166,119825.91) (170,119502.2045) (174,120828.7141) (178,120463.6687) (182,121036.8246) (186,121469.6744) (190,122053.0953) (194,123629.353) (198,124438.0174) (202,126461.1949) (206,124858.2798) (210,125740.1767) (214,124564.7602) (218,126278.9316) (222,125667.3244) (226,125989.1985) (230,125447.9801) (234,125452.1865) (238,124555.5216) (242,123827.3948) (246,123643.986) (250,123167.9141) (254,123525.1341) (258,124693.1929) (262,124218.7658) (266,123096.3198) (270,122909.3393) (274,123464.597) (278,124062.5377) (282,124455.1321) (286,126177.2166) (290,127221.0058) (294,124537.5584) (298,124174.0484) (302,125062.8675) (306,125564.9225) (310,124657.0083) (314,123868.7372) (318,123479.7601) (322,122565.0133) (326,122274.1041) (330,122358.3743) (334,123401.3173) (338,123414.6726) (342,125748.4931) (346,124737.3149) (350,123826.7628) (354,121772.6556) (358,123249.4038) (362,124029.8499) (366,124554.6845) (370,126342.7691) (374,126400.9298) (378,125460.7908) (382,123927.3354) (386,122797.3833) (390,121874.3749) (394,122512.7201) (398,124967.3925) (402,123127.5138) (406,123302.584) (410,124492.8869) (414,125404.18) (418,126355.5079) (422,123530.36) (426,124863.2241) (430,123599.4328) (434,123573.5469) (438,124366.0074) (442,123037.389) (446,124152.6227) (450,124175.5396) (454,122930.555) (458,122473.169) (462,122816.0181) (466,121845.6967) (470,121384.4126) (474,123732.9734) (478,122676.2069) (482,123914.5853) (486,122763.5033) (490,122812.3528) (494,122964.8503) (498,123809.4117) (502,123796.5698) (506,125784.9015) (510,126033.9594) (514,123862.9307) (518,123578.0733) (522,124898.5877) (526,125113.9174) (530,124253.7335) (534,124582.7673) (538,126546.0789) (542,124054.6862) (546,124639.3082) (550,124816.016) (554,125579.2406) (558,124977.2139) (562,124795.4908) (566,126493.4297) (570,125767.8582) (574,125726.0327) (578,127031.1789) (582,127270.1434) (586,126049.323) (590,126208.2137) (594,124906.3386) (598,126967.8172) (602,125729.0983) (606,124942.8398) (610,126294.0814) (614,127455.6559) (618,127707.4635) (622,127920.2906) (626,129171.6819) (630,128149.6254) (634,128206.7529) (638,128243.8141) (642,128632.128) (646,128454.1249) (650,127419.2536) (654,127945.0545) (658,128310.8294) (662,128669.6281) (666,127319.3636) (670,127192.1199) (674,125954.6254) (678,124861.7969) (682,125718.9319) (686,125592.723) (690,126104.2919) (694,125558.0788) (698,125086.2381) (702,125015.3748) (706,124350.3106) (710,124663.3382) (714,123063.0191) (718,123327.0248) (722,122800.7189) (726,123650.1441) (730,123462.6126) (734,124371.5892) (738,124438.1904) (742,124469.0665) (746,123987.7329) (750,124107.9181) (754,123110.0362) (758,122526.0477) (762,122713.4537) (766,123701.609) (770,124374.7249) (774,123662.932) (778,123923.7734) (782,123449.9571) (786,123128.341) (790,123195.5203) (794,122962.0002) (798,122244.6944) (802,124514.5658) (806,124787.4573) (810,126317.6854) (814,128157.2263) (818,128318.5773) (822,127981.8965) (826,129195.4689) (830,130225.0843) (834,131155.717) (838,131453.7529) (842,133395.7029) (846,133577.2057) (850,132542.4488) (854,132988.7975) (858,132069.8801) (862,132202.0337) (866,130538.071) (870,130154.5194) (874,129312.8246) (878,127816.438) (882,125899.0839) (886,124397.0024) (890,124826.3121) (894,123933.6462) (898,124219.7496) (902,125698.7882) (906,123842.1702) (910,123114.8025) (914,123063.0873) (918,124693.8187) (922,123781.7329) (926,123514.3165) (930,123467.9118) (934,122963.3166) (938,122219.4844) (942,122185.9109) (946,122537.2687) (950,123676.5409) (954,123738.0665) (958,123879.4206) (962,122745.0236) (966,120594.6994) (970,120971.7339) (974,123566.0958) (978,126724.8498) (982,126518.1242) (986,126184.1975) (990,125476.1886) (994,125299.201) (998,126136.1392)};
\addplot[pubNavy, opacity=0.20, line width=0.35pt] coordinates {(10,114092.2809) (14,113748.582) (18,113649.6944) (22,112637.1209) (26,114554.2139) (30,113721.0299) (34,114691.9988) (38,113359.8603) (42,113428.6423) (46,113998.3885) (50,115380.7326) (54,114959.8586) (58,116114.6655) (62,114217.2252) (66,116569.7982) (70,116948.139) (74,118128.0204) (78,117673.1827) (82,117221.0286) (86,118410.1154) (90,120679.5351) (94,120159.8403) (98,120135.4829) (102,118705.9905) (106,120474.0824) (110,119681.9455) (114,120906.5344) (118,120913.186) (122,120786.5933) (126,121425.924) (130,123734.3418) (134,120771.7243) (138,121145.8558) (142,120334.6772) (146,122430.9648) (150,121342.0442) (154,123365.071) (158,121976.7483) (162,121494.0786) (166,121240.261) (170,120723.3657) (174,120521.422) (178,119294.6574) (182,120964.8631) (186,121109.1375) (190,120193.3642) (194,121430.3502) (198,121542.9934) (202,122054.558) (206,120841.3399) (210,121680.6876) (214,120649.6038) (218,121666.9989) (222,121588.1694) (226,122932.3904) (230,122159.6715) (234,122175.7732) (238,122215.4141) (242,122402.3512) (246,122848.4374) (250,122991.0578) (254,124959.987) (258,125160.3448) (262,125254.0709) (266,123674.4009) (270,123489.0902) (274,124328.5487) (278,124499.6379) (282,125250.5969) (286,126104.5814) (290,128373.1307) (294,126290.941) (298,125684.7222) (302,126547.8956) (306,126300.636) (310,125379.9306) (314,126324.1916) (318,126111.4552) (322,125412.5139) (326,124790.3595) (330,124501.6504) (334,123978.1164) (338,123398.8329) (342,126007.7698) (346,123983.936) (350,124082.2895) (354,120948.5157) (358,122891.903) (362,123516.5987) (366,123912.9043) (370,125647.4044) (374,126241.3467) (378,127395.1416) (382,128250.1272) (386,127897.9282) (390,126043.7794) (394,127275.9868) (398,129246.0099) (402,127172.5054) (406,125543.4543) (410,126891.9761) (414,127897.7839) (418,129416.945) (422,127755.9141) (426,130014.4446) (430,128092.7484) (434,128156.8313) (438,128683.794) (442,127220.1641) (446,127524.7609) (450,127867.4313) (454,127842.6272) (458,129566.6541) (462,130891.4716) (466,129623.8788) (470,127260.726) (474,128306.855) (478,126915.6559) (482,128985.7553) (486,128124.9094) (490,130781.0543) (494,130594.8758) (498,130946.0989) (502,131578.2795) (506,131744.8126) (510,130315.0542) (514,128309.3751) (518,127697.5246) (522,129184.3659) (526,129993.7493) (530,129645.6034) (534,130062.0271) (538,133029.3849) (542,127242.8017) (546,126379.4617) (550,126567.5804) (554,127558.0671) (558,127730.4527) (562,126708.5884) (566,129477.0852) (570,129266.8267) (574,129532.6982) (578,127568.057) (582,126540.9131) (586,124794.5717) (590,124333.3762) (594,125915.1184) (598,129303.3107) (602,129642.4763) (606,127460.5485) (610,127707.1665) (614,128449.8902) (618,128621.5255) (622,127997.6128) (626,128423.7319) (630,129291.2016) (634,129558.6483) (638,129593.4749) (642,131075.1819) (646,128915.9422) (650,126416.2452) (654,127894.7943) (658,129307.445) (662,131736.0817) (666,128820.7235) (670,129899.312) (674,131251.2401) (678,129177.0406) (682,131028.6447) (686,131336.4153) (690,130834.2843) (694,131492.225) (698,132376.0598) (702,131712.9129) (706,130855.4777) (710,131542.2075) (714,129084.2446) (718,131718.8387) (722,130277.914) (726,130358.2342) (730,129819.3175) (734,128741.2015) (738,128780.6216) (742,128141.1985) (746,125227.3719) (750,124015.2627) (754,122685.5653) (758,121955.0237) (762,120027.53) (766,120388.6982) (770,120425.1432) (774,119729.8251) (778,118986.8303) (782,119422.881) (786,121075.9777) (790,120797.104) (794,121311.8706) (798,122321.8019) (802,125289.8897) (806,125550.3933) (810,126091.7706) (814,126858.56) (818,127546.7528) (822,126669.9592) (826,127119.075) (830,127504.5805) (834,128317.5977) (838,130311.9375) (842,130416.1731) (846,130012.2289) (850,128847.3375) (854,128537.5533) (858,129035.2084) (862,129750.8409) (866,128523.1085) (870,128537.5804) (874,128062.7256) (878,127006.807) (882,125500.8371) (886,121846.5967) (890,122948.519) (894,123424.5052) (898,124246.2827) (902,125095.1798) (906,123826.7015) (910,122871.5248) (914,123489.4565) (918,123917.7171) (922,124304.1547) (926,126375.8404) (930,126991.3028) (934,127906.6726) (938,125467.0621) (942,126894.8202) (946,125970.3273) (950,126238.4126) (954,126507.1702) (958,126413.5165) (962,124713.2115) (966,123887.6924) (970,122843.9944) (974,124809.123) (978,126641.3255) (982,125733.8411) (986,127406.6718) (990,127546.0459) (994,126275.7661) (998,127574.6751)};
\addplot[pubNavy, opacity=0.20, line width=0.35pt] coordinates {(10,110388.2701) (14,111718.5729) (18,109089.2691) (22,109820.9757) (26,110418.1669) (30,109580.9765) (34,111120.5841) (38,110561.9111) (42,109975.7671) (46,110087.9857) (50,111840.9413) (54,111680.2497) (58,112577.6376) (62,111151.4163) (66,114007.749) (70,114536.6645) (74,116734.0317) (78,117512.2199) (82,117855.6891) (86,118546.3157) (90,120626.1269) (94,121140.5391) (98,122558.6955) (102,120791.0054) (106,122821.6056) (110,122640.2893) (114,123572.8463) (118,123488.1626) (122,122500.4028) (126,122404.3161) (130,123663.8923) (134,120758.1669) (138,120834.8458) (142,120471.3006) (146,120774.7656) (150,118616.2669) (154,119039.5312) (158,117949.4533) (162,117890.8359) (166,118007.2227) (170,116608.4735) (174,117418.6385) (178,116857.1389) (182,118111.2299) (186,119224.5657) (190,117714.0727) (194,119355.3089) (198,119545.0433) (202,120900.1765) (206,118980.8532) (210,119232.3563) (214,118014.5896) (218,118216.5167) (222,118374.4629) (226,118673.9246) (230,118028.0973) (234,117851.9986) (238,118679.2202) (242,118419.489) (246,119770.7721) (250,119877.7244) (254,121385.133) (258,123345.4523) (262,122336.798) (266,121523.3863) (270,120819.3186) (274,121429.1049) (278,121888.9922) (282,122143.3439) (286,121367.6813) (290,121868.866) (294,120011.2554) (298,118044.3734) (302,118527.1567) (306,116676.5367) (310,116094.1247) (314,116920.5229) (318,116988.2065) (322,117219.2273) (326,116615.8681) (330,116820.7122) (334,118228.0632) (338,117996.2718) (342,119405.0245) (346,120162.4949) (350,120008.9258) (354,117919.9438) (358,118073.9317) (362,118671.7096) (366,118729.3493) (370,118682.4302) (374,119821.7281) (378,119630.6486) (382,118823.2112) (386,117891.3412) (390,116988.667) (394,116523.5235) (398,117943.4532) (402,116953.1191) (406,117539.7936) (410,118872.035) (414,120258.1062) (418,121844.9207) (422,120300.0127) (426,121779.6731) (430,121393.041) (434,122984.1198) (438,122544.139) (442,122233.995) (446,124491.7287) (450,126102.4269) (454,126366.6318) (458,126967.5393) (462,128101.0053) (466,125547.9489) (470,125320.8294) (474,126723.3926) (478,125718.4218) (482,126746.5286) (486,126462.1426) (490,127001.2755) (494,126825.7269) (498,125718.9096) (502,126550.1629) (506,127341.9785) (510,126767.935) (514,124609.2123) (518,124746.783) (522,126534.4626) (526,125437.3961) (530,125938.4848) (534,125693.7163) (538,127708.7837) (542,127385.4688) (546,126551.1052) (550,127104.5286) (554,127277.6466) (558,126848.1456) (562,127060.5318) (566,128129.3853) (570,127313.8806) (574,126892.5051) (578,125620.5876) (582,125505.6626) (586,125011.2994) (590,122192.4486) (594,122985.2651) (598,125288.8043) (602,124053.0639) (606,123045.1201) (610,124306.3051) (614,124984.487) (618,125680.0263) (622,125524.3459) (626,127107.026) (630,126949.8474) (634,126415.4416) (638,125512.633) (642,127237.3285) (646,127039.9527) (650,126387.21) (654,127790.7911) (658,128791.2473) (662,129217.6361) (666,127892.5614) (670,127931.3353) (674,128047.5698) (678,128973.6913) (682,131092.9433) (686,132085.5205) (690,132555.86) (694,130688.7435) (698,129491.0762) (702,128571.4499) (706,127712.6859) (710,129150.45) (714,126322.3938) (718,126825.1651) (722,127657.9042) (726,126391.567) (730,125268.5393) (734,123340.4496) (738,123087.1247) (742,124402.0428) (746,123036.4498) (750,122957.2286) (754,121320.2747) (758,121386.3552) (762,121183.8066) (766,122804.2906) (770,122591.1058) (774,121862.3059) (778,120792.6886) (782,122676.4596) (786,124363.9305) (790,124296.205) (794,123977.1893) (798,123808.1003) (802,126804.3903) (806,126564.2297) (810,128010.0228) (814,128663.4338) (818,129779.0728) (822,130944.9289) (826,130337.5479) (830,130639.7015) (834,132127.3821) (838,131318.2073) (842,131915.26) (846,132107.1041) (850,130646.094) (854,129774.2494) (858,128963.9557) (862,129716.6997) (866,126857.9088) (870,127445.1377) (874,126002.9269) (878,125052.0604) (882,123134.7422) (886,120301.1156) (890,120700.7872) (894,120688.5971) (898,121248.0515) (902,122670.3504) (906,122126.2543) (910,121173.7382) (914,121595.3956) (918,121851.324) (922,122050.9025) (926,122729.3555) (930,123060.4767) (934,123274.349) (938,121709.6021) (942,123379.8484) (946,121991.4256) (950,122259.6865) (954,122409.8414) (958,122648.0318) (962,122162.6945) (966,121614.3553) (970,120989.8979) (974,123042.3946) (978,125248.7195) (982,124500.2325) (986,126031.36) (990,125865.2356) (994,126072.0527) (998,126814.8549)};
\addplot[pubNavy, opacity=0.20, line width=0.35pt] coordinates {(10,114406.0385) (14,114917.0214) (18,113364.5849) (22,113169.8354) (26,114419.8587) (30,112954.4185) (34,113850.037) (38,112962.065) (42,112746.2648) (46,112232.5841) (50,113095.7406) (54,112639.0548) (58,113383.8552) (62,111425.7479) (66,113491.5457) (70,113288.1433) (74,115009.9167) (78,115240.8806) (82,115389.7382) (86,116536.7255) (90,118411.6318) (94,118014.5465) (98,118716.8026) (102,118451.519) (106,120443.4692) (110,119349.1213) (114,121209.039) (118,121675.6697) (122,121917.2303) (126,122009.7019) (130,123018.3473) (134,120204.5651) (138,120617.658) (142,120669.3344) (146,121243.6209) (150,119693.4834) (154,121244.2458) (158,119456.3785) (162,119606.834) (166,120105.3893) (170,118740.6557) (174,119297.172) (178,118743.499) (182,120030.9662) (186,120794.5946) (190,120661.9008) (194,122047.8556) (198,123436.8622) (202,124261.4364) (206,122896.2696) (210,124077.0263) (214,122045.6787) (218,123213.505) (222,123006.2233) (226,124349.308) (230,123388.9652) (234,123690.6971) (238,123305.0694) (242,122221.9399) (246,122423.2385) (250,121645.1647) (254,123370.876) (258,123489.2286) (262,122616.7207) (266,120243.2913) (270,119051.4957) (274,119102.5344) (278,119152.2422) (282,120189.3165) (286,120497.2786) (290,120681.0368) (294,119197.9802) (298,118746.9431) (302,119858.4566) (306,119845.6628) (310,119728.6239) (314,120057.3941) (318,120628.9879) (322,120886.0481) (326,120549.1106) (330,121319.2112) (334,121695.8638) (338,122518.2221) (342,124097.3592) (346,122100.3686) (350,122646.7495) (354,119620.2623) (358,121072.9453) (362,121784.4043) (366,122882.7908) (370,122906.4247) (374,124215.4454) (378,124203.4478) (382,123441.2371) (386,121833.3495) (390,120386.4519) (394,120104.7681) (398,120834.5458) (402,119964.7861) (406,118156.0097) (410,119215.2581) (414,119665.2413) (418,120928.5562) (422,119938.3444) (426,120135.9522) (430,119672.6266) (434,120168.5454) (438,120718.4698) (442,119978.0538) (446,121667.6804) (450,121579.285) (454,121391.6329) (458,121825.5161) (462,123187.6171) (466,120324.2472) (470,118253.818) (474,119844.8243) (478,118960.2593) (482,119995.9759) (486,119303.2545) (490,120871.1309) (494,120954.3123) (498,120483.0323) (502,120669.1997) (506,120550.6464) (510,120112.3536) (514,118470.927) (518,118690.4426) (522,120252.8836) (526,119948.5305) (530,118682.2542) (534,119796.1131) (538,121083.3307) (542,118390.34) (546,117579.2898) (550,119174.2831) (554,121710.0352) (558,122434.9108) (562,122688.5543) (566,124708.3201) (570,124956.8386) (574,125532.242) (578,125024.8883) (582,124965.6906) (586,124006.5002) (590,123978.6476) (594,125420.1816) (598,127191.2628) (602,125787.4355) (606,124341.6384) (610,124485.182) (614,124820.321) (618,124522.9462) (622,123949.7424) (626,124542.4055) (630,125338.1085) (634,125882.4603) (638,125546.1753) (642,127338.1549) (646,126389.0843) (650,125918.941) (654,125810.9096) (658,125699.5968) (662,126809.0454) (666,125630.5984) (670,126453.862) (674,126414.1291) (678,125984.8301) (682,126438.1202) (686,126830.1118) (690,126518.465) (694,125713.0891) (698,126986.1001) (702,127049.4189) (706,127219.6393) (710,129771.0131) (714,126808.8849) (718,129218.4623) (722,129460.1554) (726,129990.2959) (730,129540.507) (734,128919.1221) (738,128818.929) (742,129322.5381) (746,127451.5895) (750,125969.3198) (754,124368.8343) (758,124817.7464) (762,123386.2328) (766,124521.8795) (770,124113.7037) (774,122802.4785) (778,122483.5643) (782,122129.5776) (786,122496.4682) (790,123308.5037) (794,122849.3491) (798,123501.416) (802,127125.595) (806,127416.7746) (810,128307.6152) (814,128705.204) (818,129110.6528) (822,129060.0299) (826,129524.1239) (830,130328.4925) (834,131274.9112) (838,131474.6506) (842,131233.629) (846,130350.6063) (850,129307.292) (854,127950.3352) (858,127427.7603) (862,126864.5114) (866,124466.2354) (870,125276.4019) (874,123983.1943) (878,122468.1963) (882,122218.7476) (886,119583.7267) (890,120737.2926) (894,120825.6444) (898,121736.3094) (902,122712.109) (906,122271.551) (910,122020.345) (914,123297.0671) (918,124126.6191) (922,124187.1774) (926,124906.9554) (930,126114.0701) (934,127143.4123) (938,124065.4889) (942,125506.9034) (946,125150.3668) (950,126336.1392) (954,126868.5904) (958,127092.4722) (962,125472.1428) (966,122983.2149) (970,122398.5473) (974,123648.6621) (978,126312.998) (982,124293.5515) (986,125802.2847) (990,125833.781) (994,124308.2214) (998,124938.8572)};
\addplot[pubNavy, opacity=0.20, line width=0.35pt] coordinates {(10,112977.5165) (14,112599.4242) (18,112126.3985) (22,111477.9151) (26,112640.1436) (30,112263.8689) (34,113197.3533) (38,112606.5303) (42,112559.868) (46,112064.2502) (50,113941.575) (54,112656.2185) (58,113050.3484) (62,110744.3546) (66,113182.6071) (70,113216.8433) (74,115008.2797) (78,115242.8493) (82,115864.2749) (86,116877.0002) (90,119845.885) (94,119631.8629) (98,123399.3443) (102,122553.878) (106,125445.3391) (110,126764.3017) (114,128745.9227) (118,129512.7537) (122,131232.879) (126,131029.9143) (130,132534.757) (134,131376.4903) (138,132750.1365) (142,131927.6093) (146,132070.8885) (150,130147.2258) (154,131456.8835) (158,129700.3689) (162,129284.4913) (166,129492.2082) (170,127325.2597) (174,126330.3267) (178,125654.402) (182,125701.8952) (186,125072.4661) (190,125157.1898) (194,126920.2625) (198,127135.8381) (202,129057.9859) (206,128475.0884) (210,129583.3666) (214,128240.6831) (218,130052.433) (222,130402.8241) (226,132550.101) (230,132032.0974) (234,132529.5789) (238,132258.6884) (242,131028.8645) (246,131500.1215) (250,131069.1751) (254,132487.1549) (258,131550.4674) (262,130993.9256) (266,129471.6387) (270,129367.4329) (274,129677.8233) (278,129613.5097) (282,130187.2636) (286,131459.3648) (290,132918.8238) (294,130747.6926) (298,131443.3999) (302,133612.9074) (306,133921.3473) (310,132846.9334) (314,132938.6913) (318,133064.9988) (322,132764.7992) (326,131891.7041) (330,131997.2279) (334,131807.8283) (338,131295.5101) (342,133493.9813) (346,131422.6749) (350,131344.5829) (354,126382.8193) (358,127403.8071) (362,129430.8964) (366,129711.9591) (370,130694.6228) (374,132261.9307) (378,132286.502) (382,130495.5302) (386,130666.9418) (390,130694.1693) (394,131004.7331) (398,132865.264) (402,131508.3199) (406,130222.5755) (410,131584.1658) (414,132628.6307) (418,135118.1933) (422,131731.7429) (426,132867.7524) (430,131445.8254) (434,133138.7238) (438,133583.4364) (442,132361.0551) (446,132924.7852) (450,133129.4412) (454,132583.8711) (458,134526.954) (462,135370.3547) (466,133061.1972) (470,131368.4355) (474,133931.7491) (478,132791.2689) (482,134460.9537) (486,132330.1735) (490,134305.9926) (494,134131.927) (498,133333.1249) (502,133597.154) (506,134880.7338) (510,134793.8064) (514,132745.2852) (518,131596.961) (522,133913.8159) (526,135301.6463) (530,132618.7447) (534,133423.0269) (538,136082.0278) (542,132613.8631) (546,131622.7365) (550,131941.8519) (554,133938.9027) (558,132935.8676) (562,133192.9274) (566,135427.3502) (570,134697.5393) (574,133186.0774) (578,133344.7953) (582,134063.678) (586,132745.4739) (590,131665.9613) (594,133221.9883) (598,135844.8747) (602,135808.7472) (606,133557.0952) (610,134396.8819) (614,134589.3106) (618,135110.3974) (622,134574.6052) (626,135175.3202) (630,135451.9313) (634,135055.1465) (638,135205.7222) (642,136784.7673) (646,133281.2989) (650,132461.0452) (654,131557.3033) (658,131767.8977) (662,133813.4219) (666,131939.9799) (670,133076.0205) (674,134690.336) (678,133334.3478) (682,135532.4218) (686,136315.1249) (690,135359.0423) (694,134787.3208) (698,135426.7688) (702,134911.53) (706,134575.4921) (710,135257.9406) (714,132856.5297) (718,134119.6079) (722,131533.2981) (726,131876.3928) (730,130829.8986) (734,128479.9267) (738,128221.1789) (742,128658.3853) (746,126326.1706) (750,126560.9896) (754,125673.2796) (758,125146.588) (762,125848) (766,126074.4134) (770,127141.159) (774,126675.7612) (778,125719.6938) (782,127358.2114) (786,128734.3556) (790,129346.0132) (794,130091.6997) (798,130553.1201) (802,131837.3913) (806,132207.5922) (810,134048.6276) (814,134838.8041) (818,135297.6254) (822,135593.6737) (826,136501.4546) (830,137674.4539) (834,139504.3133) (838,139020.2801) (842,138169.8869) (846,138294.7782) (850,137809.8277) (854,139998.857) (858,138580.4935) (862,138858.5086) (866,137152.0377) (870,136152.4234) (874,134379.3913) (878,131487.3145) (882,129573.8629) (886,126998.3693) (890,128198.4544) (894,127862.9848) (898,128909.5175) (902,129293.4846) (906,127238.1911) (910,126482.8817) (914,127654.3144) (918,128764.4044) (922,128175.1005) (926,128869.9112) (930,129085.9326) (934,129086.354) (938,127231.8261) (942,130026.5989) (946,129170.9216) (950,128612.9886) (954,128962.8662) (958,129370.8062) (962,127079.5543) (966,124631.9062) (970,125044.1042) (974,127440.7945) (978,130083.9851) (982,130271.2987) (986,130171.8005) (990,129397.9272) (994,127593.1906) (998,128836.0197)};
\addplot[pubNavy, opacity=0.20, line width=0.35pt] coordinates {(10,108024.3985) (14,112498.6977) (18,111259.0489) (22,110569.0826) (26,113250.7407) (30,112677.5748) (34,114188.9057) (38,112899.1763) (42,112200.4088) (46,112055.5187) (50,113275.5159) (54,113169.1823) (58,114289.8844) (62,112429.3415) (66,113675.6037) (70,113815.1869) (74,114612.8321) (78,114407.9495) (82,114099.4847) (86,115527.1637) (90,117957.7674) (94,117958.4503) (98,120060.4191) (102,119575.3442) (106,121993.0085) (110,122399.9972) (114,124275.6015) (118,124903.3743) (122,125672.0224) (126,125743.0606) (130,128241.9312) (134,126244.5528) (138,125707.4994) (142,126243.122) (146,126713.6647) (150,123655.8735) (154,124698.9997) (158,122961.6545) (162,122973.7182) (166,123238.5731) (170,120318.4142) (174,120206.6378) (178,120432.552) (182,120888.6294) (186,120987.2775) (190,120382.6675) (194,119922.5013) (198,118977.441) (202,118921.508) (206,117660.0303) (210,117077.6972) (214,113677.0143) (218,115547.834) (222,116856.1339) (226,117565.7287) (230,115417.6483) (234,114913.9778) (238,114784.4374) (242,115987.198) (246,117882.1854) (250,119124.9327) (254,120814.9112) (258,122265.4101) (262,121673.0639) (266,119650.5559) (270,119416.0422) (274,120154.0324) (278,119599.0447) (282,120533.9488) (286,121238.7402) (290,122038.0734) (294,118896.4139) (298,117450.0484) (302,117189.3715) (306,115943.8267) (310,115380.9383) (314,115845.7018) (318,115889.9745) (322,115742.0387) (326,114300.5352) (330,113541.7016) (334,113502.6783) (338,114089.4763) (342,116948.3314) (346,115947.2406) (350,116337.7253) (354,115264.4603) (358,115946.0944) (362,117541.1331) (366,117018.5579) (370,117181.8468) (374,119188.9855) (378,119478.7899) (382,118299.9902) (386,118322.7865) (390,116430.0232) (394,116260.7903) (398,119005.2172) (402,118602.1006) (406,118279.2962) (410,119462.3089) (414,120409.5084) (418,123499.4235) (422,121132.2668) (426,122247.0698) (430,122523.7031) (434,124497.2405) (438,123873.4047) (442,123131.3714) (446,123625.1799) (450,123642.4175) (454,123453.5595) (458,125436.3302) (462,127526.3205) (466,125467.9478) (470,123941.0253) (474,126442.6585) (478,125043.3514) (482,127159.5401) (486,126481.4016) (490,127568.1913) (494,127613.8871) (498,128265.0334) (502,129471.2345) (506,128632.0038) (510,128539.1486) (514,126055.3396) (518,125582.8978) (522,127207.5857) (526,126896.7043) (530,125991.1911) (534,126865.6165) (538,128683.4922) (542,128308.9221) (546,127069.2767) (550,126404.3026) (554,128207.9776) (558,127928.5869) (562,127825.0269) (566,130052.8684) (570,130624.1067) (574,130537.3011) (578,129330.9106) (582,128818.8105) (586,126822.097) (590,124360.8564) (594,124877.7504) (598,127465.1629) (602,126974.5148) (606,124357.2863) (610,125911.6325) (614,126004.8296) (618,125643.1854) (622,124849.2175) (626,124708.9742) (630,125337.1621) (634,126692.325) (638,126678.368) (642,128095.8224) (646,127250.5534) (650,126062.521) (654,127908.877) (658,129190.0729) (662,130348.5533) (666,128567.6226) (670,128011.7703) (674,129892.8945) (678,129468.214) (682,131349.7665) (686,131074.4594) (690,130311.4614) (694,129578.3802) (698,129369.8288) (702,129171.4798) (706,128827.8762) (710,129799.4701) (714,128162.8869) (718,129197.4895) (722,128886.8579) (726,128156.3446) (730,127329.3261) (734,126675.166) (738,127019.2267) (742,126607.2614) (746,124277.6204) (750,124755.9103) (754,123125.0397) (758,121531.0053) (762,120208.7374) (766,122578.9286) (770,123209.5575) (774,121945.3084) (778,121028.155) (782,122430.2823) (786,122991.0239) (790,123724.8049) (794,124592.9638) (798,125364.0505) (802,128821.7234) (806,128815.5891) (810,129550.9827) (814,130016.5161) (818,131748.1158) (822,132326.1437) (826,132644.6947) (830,133173.5411) (834,135501.5869) (838,135305.532) (842,136204.7186) (846,134787.2689) (850,134037.754) (854,133424.8142) (858,133331.1163) (862,133341.4953) (866,130596.3703) (870,130092.7835) (874,128794.9915) (878,127845.2858) (882,126724.6092) (886,124010.7719) (890,123665.8318) (894,123800.1666) (898,125887.7161) (902,125841.2622) (906,125170.1627) (910,124180.1543) (914,125192.8088) (918,125404.481) (922,124128.5898) (926,124659.214) (930,124539.4317) (934,125081.6916) (938,123341.314) (942,125191.2509) (946,123135.557) (950,124255.3703) (954,125103.9912) (958,124884.4099) (962,123712.3629) (966,122522.0398) (970,123875.5437) (974,125714.269) (978,129413.2608) (982,129101.5273) (986,129396.3878) (990,128273.5131) (994,127137.3457) (998,128813.4757)};
\addplot[pubNavy, ultra thick] coordinates {(10,113965.8736) (14,113290.8848) (18,112483.5186) (22,112451.041) (26,113835.706) (30,112670.6663) (34,114110.4709) (38,112727.8837) (42,112791.6) (46,112148.4172) (50,113910.2804) (54,113063.0503) (58,113435.1373) (62,111574.0852) (66,114061.219) (70,114416.3447) (74,115694.6067) (78,115456.2851) (82,116121.4341) (86,117117.063) (90,119506.696) (94,119166.1447) (98,120221.5473) (102,119311.4849) (106,121486.5618) (110,121777.9586) (114,122945.1821) (118,123145.884) (122,123553.7074) (126,123616.1563) (130,124704.6526) (134,122991.0229) (138,123460.7376) (142,123246.2509) (146,124060.0781) (150,122359.9487) (154,123452.3396) (158,121800.9853) (162,121978.371) (166,121498.8882) (170,119910.3093) (174,119751.9049) (178,119019.0782) (182,120459.7978) (186,120890.9361) (190,120298.8704) (194,121076.2891) (198,121512.7656) (202,122199.6174) (206,121095.1981) (210,122088.2838) (214,120992.4463) (218,122236.9029) (222,122578.3468) (226,124441.7958) (230,123395.1903) (234,123451.0823) (238,123261.696) (242,122903.7933) (246,122964.6042) (250,122826.9599) (254,123848.2212) (258,124618.2514) (262,123772.8892) (266,122250.2298) (270,121707.2108) (274,122394.3546) (278,122634.0571) (282,123359.8294) (286,123710.5802) (290,124783.6588) (294,122801.2229) (298,123352.876) (302,123795.5887) (306,124219.9329) (310,123718.7563) (314,123951.6752) (318,124046.4604) (322,124147.5462) (326,123452.3971) (330,123435.6826) (334,123402.2004) (338,122830.588) (342,125469.1576) (346,123462.5987) (350,123430.7725) (354,121122.7013) (358,121779.4096) (362,122507.0931) (366,122729.4599) (370,123433.2494) (374,124829.9069) (378,124728.2046) (382,123961.1559) (386,123215.7955) (390,122438.3551) (394,122860.2043) (398,124366.4832) (402,122256.7667) (406,121159.2382) (410,122955.2409) (414,123534.33) (418,124938.1734) (422,123320.768) (426,124603.4301) (430,123786.1198) (434,124116.6314) (438,124051.7711) (442,123084.3802) (446,124395.7123) (450,125608.1392) (454,125498.8736) (458,126219.4268) (462,127390.9648) (466,125507.9483) (470,124234.3057) (474,126070.5588) (478,124688.5698) (482,126225.1021) (486,125809.1185) (490,126130.9185) (494,125670.1738) (498,124778.8701) (502,125504.1131) (506,126563.44) (510,126400.9472) (514,124236.0715) (518,124162.4281) (522,125634.5291) (526,125275.6568) (530,125262.8581) (534,125512.3601) (538,127292.0477) (542,125003.0081) (546,125430.1752) (550,125280.6724) (554,125909.6101) (558,126090.36) (562,125929.1872) (566,127486.1131) (570,127170.9114) (574,126918.75) (578,126945.0919) (582,126368.4355) (586,125367.3834) (590,124389.3623) (594,125163.2601) (598,127170.3419) (602,125976.6276) (606,124376.9169) (610,126102.8569) (614,126695.4277) (618,126977.6041) (622,126212.9115) (626,126943.6012) (630,126709.0696) (634,126635.8467) (638,126578.1432) (642,127806.0435) (646,127145.2531) (650,126401.7276) (654,127829.8317) (658,128136.9283) (662,129043.8583) (666,127146.1281) (670,127561.7276) (674,128340.2198) (678,127835.4355) (682,130159.6502) (686,130219.3877) (690,129748.2224) (694,128901.057) (698,128678.8172) (702,128304.9941) (706,127014.0542) (710,128831.2956) (714,126565.6393) (718,127111.1479) (722,126594.6383) (726,126437.8646) (730,125214.208) (734,124566.3527) (738,125083.4743) (742,124909.3656) (746,123961.316) (750,124204.0446) (754,122876.028) (758,122178.4736) (762,121762.7663) (766,122691.6096) (770,123070.2497) (774,122069.8696) (778,120866.322) (782,122279.93) (786,123369.2103) (790,123252.012) (794,123483.6689) (798,123654.7581) (802,126439.1903) (806,126360.7497) (810,127405.2424) (814,128410.33) (818,129399.0421) (822,129385.1768) (826,129680.5181) (830,130276.7884) (834,131282.9591) (838,131385.9801) (842,131370.6038) (846,130775.4124) (850,129496.0842) (854,129155.9014) (858,128999.5821) (862,129733.7703) (866,127456.8275) (870,127423.5426) (874,127231.8579) (878,126156.5829) (882,124805.4335) (886,121647.8714) (890,122942.8643) (894,123422.877) (898,124233.0162) (902,125174.6138) (906,124011.1301) (910,123550.0472) (914,123957.6001) (918,124733.2667) (922,124186.5323) (926,124969.8507) (930,124396.6861) (934,125225.7861) (938,123337.5571) (942,125349.0772) (946,123638.9176) (950,124398.039) (954,124865.052) (958,124913.3287) (962,124034.9393) (966,122698.7962) (970,122142.3848) (974,123607.3789) (978,126477.1618) (982,125912.8935) (986,126795.4346) (990,126514.2964) (994,125757.8615) (998,126286.4303)};
\nextgroupplot[title={Rolling turnover}, xlabel={Pre-training episode}]
\addplot[pubNavy, opacity=0.20, line width=0.35pt] coordinates {(10,1.058785449) (14,1.074601967) (18,1.07542363) (22,1.076739998) (26,1.074403501) (30,1.073985917) (34,1.075713621) (38,1.071169697) (42,1.068497554) (46,0.9950157484) (50,0.9320892851) (54,0.8671892234) (58,0.7980811277) (62,0.7296039483) (66,0.6551682563) (70,0.5845794863) (74,0.5162712527) (78,0.4522817834) (82,0.382582547) (86,0.3141556812) (90,0.2394546758) (94,0.209948694) (98,0.2113613449) (102,0.2139491865) (106,0.2153464804) (110,0.2200207301) (114,0.2261862079) (118,0.2270965532) (122,0.2297148215) (126,0.2236758936) (130,0.2238708999) (134,0.2195334367) (138,0.2175842196) (142,0.2153959986) (146,0.2115066665) (150,0.2040271171) (154,0.2005077184) (158,0.205268724) (162,0.2091976303) (166,0.2132441095) (170,0.2174264171) (174,0.2299919979) (178,0.2380514033) (182,0.2512568335) (186,0.2694034515) (190,0.2920379919) (194,0.3145249593) (198,0.3357394333) (202,0.3614994607) (206,0.3759234064) (210,0.3908748129) (214,0.4028297295) (218,0.4128048979) (222,0.4271282084) (226,0.4347286335) (230,0.444094486) (234,0.4402535373) (238,0.4409492651) (242,0.4440460359) (246,0.4555743404) (250,0.4567739027) (254,0.4566154282) (258,0.4616084705) (262,0.4594728575) (266,0.4660675288) (270,0.4637428108) (274,0.471672605) (278,0.4729281955) (282,0.4863689605) (286,0.4956444156) (290,0.501455335) (294,0.4927541103) (298,0.4941855906) (302,0.4959164476) (306,0.5029410867) (310,0.5047779889) (314,0.5185999043) (318,0.5189313333) (322,0.5197886391) (326,0.5250674045) (330,0.5332987963) (334,0.5388248131) (338,0.5321270247) (342,0.5403446096) (346,0.5471625644) (350,0.5514073469) (354,0.5559004489) (358,0.5648768501) (362,0.5612333047) (366,0.5665628837) (370,0.5721188959) (374,0.5843255168) (378,0.5864234714) (382,0.5983148361) (386,0.5985206719) (390,0.6075334173) (394,0.6145215775) (398,0.6158136788) (402,0.6202330064) (406,0.6191831158) (410,0.6233292688) (414,0.6195112481) (418,0.6228893087) (422,0.6295675961) (426,0.6258244013) (430,0.6237694982) (434,0.6282585348) (438,0.6342375807) (442,0.6420018106) (446,0.646913984) (450,0.6473517594) (454,0.655113128) (458,0.6536159665) (462,0.6580746335) (466,0.6607422103) (470,0.6695310551) (474,0.680709456) (478,0.6864830887) (482,0.6960696357) (486,0.6963412983) (490,0.6886323777) (494,0.6951544034) (498,0.7097782181) (502,0.7137145409) (506,0.7162932347) (510,0.7225261826) (514,0.738670984) (518,0.738608555) (522,0.7356396308) (526,0.731195096) (530,0.7223930431) (534,0.7103246773) (538,0.7255026004) (542,0.7259925682) (546,0.7210090749) (550,0.7130050775) (554,0.7168316809) (558,0.7108597204) (562,0.7044409294) (566,0.695365497) (570,0.7016165561) (574,0.6966038736) (578,0.7004867453) (582,0.7196695685) (586,0.7227703281) (590,0.7240691701) (594,0.7214533073) (598,0.7275766935) (602,0.7390767915) (606,0.7336139539) (610,0.737260144) (614,0.7453814888) (618,0.7431922399) (622,0.7462519323) (626,0.7469491656) (630,0.7289963761) (634,0.7089804046) (638,0.7050184753) (642,0.6901955293) (646,0.6826206871) (650,0.6711717869) (654,0.673717866) (658,0.6779384618) (662,0.6781829728) (666,0.6748280405) (670,0.6754387778) (674,0.6804419873) (678,0.6788443958) (682,0.6839760508) (686,0.674524248) (690,0.6704361425) (694,0.6603275571) (698,0.6566772179) (702,0.6327993185) (706,0.6369840505) (710,0.6311751021) (714,0.6222545716) (718,0.6161872549) (722,0.6066215231) (726,0.608171207) (730,0.607197556) (734,0.6112132605) (738,0.6185027559) (742,0.6295822366) (746,0.6423123209) (750,0.6536923446) (754,0.6588249463) (758,0.6568001773) (762,0.6661777785) (766,0.6809121316) (770,0.6860709674) (774,0.6887709692) (778,0.6904268924) (782,0.7039274161) (786,0.7099970666) (790,0.6926805727) (794,0.6890432894) (798,0.6890304214) (802,0.6789183194) (806,0.6766238315) (810,0.6667075259) (814,0.6616239933) (818,0.6665173701) (822,0.6693888029) (826,0.6736635117) (830,0.6662905723) (834,0.6687132195) (838,0.6767000344) (842,0.6898011031) (846,0.6978120916) (850,0.685906678) (854,0.6841005195) (858,0.6877291161) (862,0.6902416035) (866,0.6905531592) (870,0.6763783227) (874,0.6663306116) (878,0.6518167549) (882,0.6418007864) (886,0.634453705) (890,0.6323334021) (894,0.6309595716) (898,0.6282451223) (902,0.6464700449) (906,0.650719063) (910,0.6405738581) (914,0.637135222) (918,0.645326979) (922,0.6523734319) (926,0.6492992356) (930,0.6531393348) (934,0.6733970994) (938,0.6803219287) (942,0.6791796183) (946,0.6765629375) (950,0.6823610121) (954,0.6763964984) (958,0.6819986427) (962,0.6803776317) (966,0.680741426) (970,0.6789601192) (974,0.6773402872) (978,0.6865247291) (982,0.6906103144) (986,0.6722852926) (990,0.6647421569) (994,0.6567724836) (998,0.6557290946)};
\addplot[pubNavy, opacity=0.20, line width=0.35pt] coordinates {(10,1.069545483) (14,1.102535665) (18,1.09169062) (22,1.084770847) (26,1.086930163) (30,1.089352812) (34,1.091482277) (38,1.09043775) (42,1.082769385) (46,1.008432803) (50,0.9463943162) (54,0.8755100598) (58,0.8077126033) (62,0.7341994558) (66,0.662889345) (70,0.5970739502) (74,0.5248463005) (78,0.4513337956) (82,0.3821328546) (86,0.3136028534) (90,0.2388046284) (94,0.2067438444) (98,0.2048969485) (102,0.2049056713) (106,0.2046035749) (110,0.2076270509) (114,0.2126372015) (118,0.2189123961) (122,0.232752314) (126,0.247498314) (130,0.2697515343) (134,0.2866629946) (138,0.2991472672) (142,0.3136721847) (146,0.326985801) (150,0.3385531833) (154,0.3483644192) (158,0.3728121792) (162,0.3887767506) (166,0.4051995577) (170,0.4046587366) (174,0.4115756486) (178,0.3992004267) (182,0.4029396249) (186,0.4091711853) (190,0.4191621619) (194,0.4244116332) (198,0.4248403482) (202,0.4345041958) (206,0.4266746119) (210,0.42578413) (214,0.4160389135) (218,0.4145233536) (222,0.4230793522) (226,0.4333303467) (230,0.4378105177) (234,0.4254218048) (238,0.4223212639) (242,0.4295101152) (246,0.4363562963) (250,0.4448454626) (254,0.4508791253) (258,0.4557097949) (262,0.4610362589) (266,0.4701701027) (270,0.4643418727) (274,0.4840830016) (278,0.4843451488) (282,0.4942907186) (286,0.4984118153) (290,0.5014269584) (294,0.5007248329) (298,0.4956626477) (302,0.4950228683) (306,0.5009741133) (310,0.5044150682) (314,0.5152070425) (318,0.5241101345) (322,0.5292662079) (326,0.5242197372) (330,0.5256133103) (334,0.5173487447) (338,0.5171016106) (342,0.5256073217) (346,0.5264363633) (350,0.5375705667) (354,0.5317804081) (358,0.5287326391) (362,0.5248159808) (366,0.5246385247) (370,0.5185371184) (374,0.5181289063) (378,0.5218281161) (382,0.5361701917) (386,0.5301227741) (390,0.5405437374) (394,0.5457544499) (398,0.5505803674) (402,0.5542466889) (406,0.5576612324) (410,0.5534233814) (414,0.5515043104) (418,0.55593487) (422,0.5592702317) (426,0.5596868817) (430,0.5622404247) (434,0.5623678818) (438,0.5650445951) (442,0.5660148258) (446,0.5669143195) (450,0.5540170016) (454,0.5589213713) (458,0.5601875626) (462,0.5683992477) (466,0.5538033679) (470,0.5558862573) (474,0.5564744575) (478,0.5609177226) (482,0.5618740363) (486,0.5546547976) (490,0.5403127419) (494,0.5408620771) (498,0.541118748) (502,0.5308712314) (506,0.5201738132) (510,0.5146220801) (514,0.5194395893) (518,0.5174657736) (522,0.5138318201) (526,0.4988041509) (530,0.494672388) (534,0.4932161859) (538,0.5010581897) (542,0.5037446391) (546,0.5037124493) (550,0.5028234465) (554,0.5069737262) (558,0.4986699383) (562,0.4942126445) (566,0.4926123824) (570,0.4983034754) (574,0.5023930984) (578,0.4980422904) (582,0.5008439008) (586,0.4970223284) (590,0.4926831841) (594,0.4915968148) (598,0.492572503) (602,0.4985550907) (606,0.4996002978) (610,0.5010025649) (614,0.4986439781) (618,0.4997061608) (622,0.4995998357) (626,0.5012641163) (630,0.4951120089) (634,0.4989831636) (638,0.4977208659) (642,0.5027795699) (646,0.5067694479) (650,0.4906875404) (654,0.4904511016) (658,0.4988251298) (662,0.5084308051) (666,0.5160537651) (670,0.5099219753) (674,0.5205357505) (678,0.5162335617) (682,0.5218992811) (686,0.5221004047) (690,0.533363967) (694,0.5395893834) (698,0.5496278585) (702,0.5461203005) (706,0.5542845551) (710,0.5691191266) (714,0.5704250236) (718,0.5739579626) (722,0.5726844901) (726,0.5743029781) (730,0.5825087946) (734,0.5878565957) (738,0.5894659427) (742,0.591082437) (746,0.5840769891) (750,0.5935272007) (754,0.5874270382) (758,0.5739181087) (762,0.5616809151) (766,0.5746774038) (770,0.5815341445) (774,0.5833512336) (778,0.5789881798) (782,0.5814988821) (786,0.5820072643) (790,0.5785305323) (794,0.5652196814) (798,0.5699731445) (802,0.570566548) (806,0.5697817557) (810,0.5659372999) (814,0.5577528409) (818,0.5560517286) (822,0.5491856898) (826,0.5465633092) (830,0.5368023068) (834,0.5337198515) (838,0.5466088538) (842,0.5494126004) (846,0.5485351759) (850,0.5417928179) (854,0.5424195633) (858,0.54720026) (862,0.5553926947) (866,0.5473469994) (870,0.5518256022) (874,0.5467607147) (878,0.5465881343) (882,0.5357513283) (886,0.5282163469) (890,0.5247639965) (894,0.5209659367) (898,0.5231828808) (902,0.5147365908) (906,0.5162108375) (910,0.5153357717) (914,0.5095910682) (918,0.5151979401) (922,0.5142689472) (926,0.5084691747) (930,0.5230211661) (934,0.5464652707) (938,0.5538643374) (942,0.5624547192) (946,0.5609072521) (950,0.5664937043) (954,0.5796657022) (958,0.5789826898) (962,0.5797041977) (966,0.5821953657) (970,0.5951388405) (974,0.6018881058) (978,0.6192388828) (982,0.605861568) (986,0.6038130768) (990,0.6055089402) (994,0.6022273873) (998,0.6169840745)};
\addplot[pubNavy, opacity=0.20, line width=0.35pt] coordinates {(10,1.022851915) (14,1.032171856) (18,1.043468781) (22,1.04774592) (26,1.054621605) (30,1.059741893) (34,1.058249663) (38,1.055110096) (42,1.053812872) (46,0.9864883311) (50,0.928506509) (54,0.8644842899) (58,0.8049056387) (62,0.7424784431) (66,0.6783769949) (70,0.617541286) (74,0.5571717723) (78,0.491550451) (82,0.4257585622) (86,0.3607910245) (90,0.2926365242) (94,0.2552720855) (98,0.2467206024) (102,0.2359683439) (106,0.2245067252) (110,0.2126295822) (114,0.1978131894) (118,0.185087311) (122,0.1698002052) (126,0.1559760472) (130,0.147090672) (134,0.1430805969) (138,0.1390710432) (142,0.1409259573) (146,0.144114694) (150,0.1565983643) (154,0.1728283903) (158,0.1918931317) (162,0.2031459735) (166,0.2122825869) (170,0.2206702481) (174,0.2313814881) (178,0.2389984618) (182,0.2519169648) (186,0.2663403044) (190,0.2757580705) (194,0.2875170052) (198,0.2942611728) (202,0.2928979897) (206,0.2900259224) (210,0.2923323787) (214,0.2948212959) (218,0.2974723377) (222,0.3087125653) (226,0.3148113939) (230,0.3175180578) (234,0.3090747481) (238,0.3043307456) (242,0.3006204737) (246,0.300007089) (250,0.2954939189) (254,0.2940802317) (258,0.2916954255) (262,0.2904188411) (266,0.2902429773) (270,0.2782972198) (274,0.2761615003) (278,0.2692144379) (282,0.2667004342) (286,0.27093873) (290,0.2709875513) (294,0.2617755723) (298,0.2571460729) (302,0.2545088781) (306,0.2534035185) (310,0.2483584138) (314,0.2479469408) (318,0.2463962095) (322,0.246560099) (326,0.2441709931) (330,0.2438309952) (334,0.2373293776) (338,0.2334743072) (342,0.2309899019) (346,0.2298392117) (350,0.2286629838) (354,0.2278528398) (358,0.2291220125) (362,0.2252749306) (366,0.2247359598) (370,0.2196735481) (374,0.2207802617) (378,0.2222864281) (382,0.219357706) (386,0.2210221928) (390,0.2213358531) (394,0.2247922393) (398,0.2283301003) (402,0.2369959734) (406,0.2442798737) (410,0.2530393061) (414,0.2660895435) (418,0.2844034744) (422,0.292017414) (426,0.3055698104) (430,0.3235850655) (434,0.3429947376) (438,0.3671524799) (442,0.3855238619) (446,0.4054639827) (450,0.4248932367) (454,0.4373875644) (458,0.4554104563) (462,0.467708913) (466,0.4798119712) (470,0.4915762044) (474,0.5130810385) (478,0.5274648997) (482,0.5318698205) (486,0.5258612842) (490,0.5298286113) (494,0.5414261104) (498,0.5493853081) (502,0.5470161104) (506,0.5543416544) (510,0.5468868184) (514,0.5504257655) (518,0.5488397152) (522,0.539291411) (526,0.5316908519) (530,0.5311780503) (534,0.5379965087) (538,0.5417213197) (542,0.5344826309) (546,0.5278196023) (550,0.5146120599) (554,0.514280381) (558,0.5089859673) (562,0.5059325196) (566,0.4938398532) (570,0.4929776656) (574,0.4863873244) (578,0.4739191243) (582,0.4638421527) (586,0.4525121653) (590,0.4406203117) (594,0.4327327667) (598,0.431971171) (602,0.4341173497) (606,0.4157381783) (610,0.406661037) (614,0.401230778) (618,0.3964823871) (622,0.400054723) (626,0.3928734358) (630,0.3902169345) (634,0.3865640337) (638,0.3857204482) (642,0.3835509527) (646,0.3798331569) (650,0.3731088872) (654,0.376190034) (658,0.3805042786) (662,0.3923856983) (666,0.3922380972) (670,0.3882743329) (674,0.3878268541) (678,0.381119792) (682,0.3822428049) (686,0.3832023593) (690,0.3856598485) (694,0.3886369403) (698,0.3949550233) (702,0.3931607307) (706,0.3906819289) (710,0.38619617) (714,0.3874308482) (718,0.3871233798) (722,0.3832032828) (726,0.3905645621) (730,0.3918882723) (734,0.3968005179) (738,0.3922588106) (742,0.3918693778) (746,0.3913373218) (750,0.3903268922) (754,0.3923109335) (758,0.3917673434) (762,0.3948988004) (766,0.4047545413) (770,0.4088818568) (774,0.4103309032) (778,0.4107567398) (782,0.4132673212) (786,0.4211050303) (790,0.4162949034) (794,0.4107791851) (798,0.4102116424) (802,0.4113405638) (806,0.4030146895) (810,0.4033792458) (814,0.3962026856) (818,0.3914423464) (822,0.3907491554) (826,0.382948435) (830,0.3829820958) (834,0.379395214) (838,0.3813596399) (842,0.3847904695) (846,0.3862742053) (850,0.3854787001) (854,0.3837012178) (858,0.3853731895) (862,0.3840967694) (866,0.392803438) (870,0.3842072128) (874,0.3817873592) (878,0.3815323196) (882,0.3772997077) (886,0.374036671) (890,0.3756348416) (894,0.386359317) (898,0.3774711307) (902,0.3822565973) (906,0.3927755774) (910,0.3997592779) (914,0.4000934358) (918,0.3995790228) (922,0.3961740474) (926,0.3931740479) (930,0.3981948303) (934,0.4046838) (938,0.4086311448) (942,0.4242318485) (946,0.4241993105) (950,0.4234659802) (954,0.4243674258) (958,0.4212581901) (962,0.4148732772) (966,0.4214196866) (970,0.426040538) (974,0.4285009752) (978,0.4346405457) (982,0.441675461) (986,0.4339588774) (990,0.4284632805) (994,0.4115813747) (998,0.4127622704)};
\addplot[pubNavy, opacity=0.20, line width=0.35pt] coordinates {(10,1.082013027) (14,1.07671832) (18,1.079068969) (22,1.077463089) (26,1.083563238) (30,1.0769917) (34,1.080168047) (38,1.076015631) (42,1.070908789) (46,1.004350351) (50,0.9474396784) (54,0.8815041834) (58,0.8143797686) (62,0.7483626529) (66,0.6834038269) (70,0.6173574566) (74,0.5457214052) (78,0.4818828601) (82,0.4172121234) (86,0.3512896328) (90,0.2813380164) (94,0.2513072435) (98,0.2481498438) (102,0.2482433597) (106,0.2569360533) (110,0.2694889011) (114,0.2846876083) (118,0.302253318) (122,0.3240007951) (126,0.3391563507) (130,0.3692115475) (134,0.3900023311) (138,0.412968381) (142,0.4323246279) (146,0.4635694021) (150,0.4801383721) (154,0.497958574) (158,0.5205805444) (162,0.528505264) (166,0.544902618) (170,0.5345193518) (174,0.5383059778) (178,0.5307543788) (182,0.5413341465) (186,0.5592850091) (190,0.5725514441) (194,0.5800128002) (198,0.5899452811) (202,0.5880986826) (206,0.5816211624) (210,0.5866086751) (214,0.5763409121) (218,0.5841746616) (222,0.6099256416) (226,0.6227140219) (230,0.61031258) (234,0.6003408567) (238,0.5941476009) (242,0.5882833898) (246,0.5954722973) (250,0.603110761) (254,0.6159633345) (258,0.6266945346) (262,0.6274031035) (266,0.6343028134) (270,0.621080027) (274,0.6270540153) (278,0.6264802291) (282,0.6345577829) (286,0.6216063331) (290,0.6169844676) (294,0.6061436455) (298,0.5859168318) (302,0.585263151) (306,0.5681898728) (310,0.558074526) (314,0.5584093735) (318,0.5469485613) (322,0.5545062796) (326,0.5535297081) (330,0.5583286352) (334,0.56477495) (338,0.5551751653) (342,0.5682253893) (346,0.5672858292) (350,0.5689154551) (354,0.5672088465) (358,0.5623918871) (362,0.5673700473) (366,0.5717663283) (370,0.5551149079) (374,0.5573530569) (378,0.5641416711) (382,0.5648717079) (386,0.5602463934) (390,0.5599083356) (394,0.5723336256) (398,0.5698163397) (402,0.5649152455) (406,0.5558638636) (410,0.5620475097) (414,0.5501086084) (418,0.5505004797) (422,0.5491368686) (426,0.5394832213) (430,0.529396434) (434,0.5335507211) (438,0.5360521483) (442,0.5412585398) (446,0.5442255986) (450,0.5472087379) (454,0.5557940569) (458,0.5623525852) (462,0.5714038708) (466,0.5733269278) (470,0.5739066298) (474,0.5784720333) (478,0.5767589705) (482,0.5869944253) (486,0.5755424539) (490,0.5524914357) (494,0.5489920482) (498,0.5471084429) (502,0.5464380518) (506,0.5456450335) (510,0.5458033617) (514,0.5434289395) (518,0.535988575) (522,0.543806853) (526,0.5439637377) (530,0.5537246877) (534,0.5578121891) (538,0.5831277247) (542,0.5852671313) (546,0.6024215803) (550,0.6049415655) (554,0.6229253752) (558,0.6222647661) (562,0.6340828404) (566,0.6449516477) (570,0.65796132) (574,0.6522729435) (578,0.6570504095) (582,0.6618640925) (586,0.6563761625) (590,0.6629497828) (594,0.6745388872) (598,0.6852242696) (602,0.6917048339) (606,0.6828597083) (610,0.6818751617) (614,0.6819963579) (618,0.6851897611) (622,0.7006416457) (626,0.7054726376) (630,0.710908367) (634,0.7167979149) (638,0.7149722352) (642,0.712400718) (646,0.7110099565) (650,0.6978112029) (654,0.7093050241) (658,0.7258387766) (662,0.746671283) (666,0.737075127) (670,0.7364563393) (674,0.7470114826) (678,0.7319071995) (682,0.7318453142) (686,0.7190600828) (690,0.7293952763) (694,0.7263773024) (698,0.7341181035) (702,0.7226642743) (706,0.7137582554) (710,0.7124366208) (714,0.6982174874) (718,0.7064395476) (722,0.6984742907) (726,0.7013256849) (730,0.7108401809) (734,0.7279773593) (738,0.727830609) (742,0.7406051326) (746,0.7382017853) (750,0.7392745022) (754,0.7351958495) (758,0.7315154426) (762,0.7263140795) (766,0.7486885966) (770,0.7628254723) (774,0.7576227302) (778,0.760251974) (782,0.7621587572) (786,0.7555610473) (790,0.7492531013) (794,0.7378103194) (798,0.746157662) (802,0.7554240638) (806,0.7394800924) (810,0.7326400475) (814,0.7107113294) (818,0.7059594363) (822,0.7067952882) (826,0.7125712668) (830,0.6980068645) (834,0.7017817931) (838,0.7096352664) (842,0.7062499583) (846,0.6985519588) (850,0.6844751292) (854,0.6731326444) (858,0.676763284) (862,0.681628606) (866,0.675797485) (870,0.6662252896) (874,0.6547444113) (878,0.6575270129) (882,0.6493056079) (886,0.6344980821) (890,0.642663545) (894,0.6612442013) (898,0.6559340916) (902,0.6707382844) (906,0.6857602971) (910,0.6900041857) (914,0.6953454411) (918,0.6924444147) (922,0.6889160087) (926,0.6786954706) (930,0.6734455835) (934,0.6940639882) (938,0.6849424807) (942,0.6811763617) (946,0.6784522438) (950,0.6793376414) (954,0.6726811977) (958,0.6657663683) (962,0.6634376961) (966,0.6645623541) (970,0.6646392151) (974,0.6712018648) (978,0.679357116) (982,0.6768069144) (986,0.6702727129) (990,0.6682082542) (994,0.6521683956) (998,0.657782136)};
\addplot[pubNavy, opacity=0.20, line width=0.35pt] coordinates {(10,1.087494931) (14,1.079838293) (18,1.076374391) (22,1.075520055) (26,1.078108429) (30,1.082976615) (34,1.078725056) (38,1.076895573) (42,1.071706702) (46,0.996918827) (50,0.9317229907) (54,0.8612173372) (58,0.7906602126) (62,0.7290914478) (66,0.6688499482) (70,0.6062308396) (74,0.5386052815) (78,0.4707991667) (82,0.404984626) (86,0.3482173264) (90,0.2852824827) (94,0.2670287956) (98,0.2735856257) (102,0.2855926715) (106,0.2881371994) (110,0.2901383413) (114,0.2864186439) (118,0.2878877182) (122,0.2919813886) (126,0.2967649363) (130,0.3087761126) (134,0.3140010913) (138,0.3152392229) (142,0.3167021261) (146,0.3120680272) (150,0.3095749156) (154,0.3150722402) (158,0.3165540355) (162,0.3205267611) (166,0.3241362291) (170,0.3259762253) (174,0.3318053135) (178,0.3344945386) (182,0.3367742035) (186,0.34604005) (190,0.3550615195) (194,0.3686676143) (198,0.3838868925) (202,0.3917637514) (206,0.3953458025) (210,0.4001698684) (214,0.405382297) (218,0.4155386213) (222,0.4247536975) (226,0.4252381541) (230,0.420825014) (234,0.4286386974) (238,0.4212492792) (242,0.417262769) (246,0.4248505377) (250,0.4305462483) (254,0.43916142) (258,0.4531229979) (262,0.4642877063) (266,0.471537285) (270,0.4830542261) (274,0.4824100549) (278,0.4887106331) (282,0.4914718753) (286,0.4845079658) (290,0.4933218255) (294,0.4940442694) (298,0.4786582726) (302,0.4794902753) (306,0.4759444298) (310,0.4659857191) (314,0.464589604) (318,0.4523646552) (322,0.4513738093) (326,0.4468681977) (330,0.4486084593) (334,0.469346281) (338,0.4672135044) (342,0.472500553) (346,0.4741145069) (350,0.4777329524) (354,0.4696627441) (358,0.4716012238) (362,0.4764482828) (366,0.4726313693) (370,0.463059462) (374,0.467467049) (378,0.4780889978) (382,0.4733740394) (386,0.459994692) (390,0.4730884793) (394,0.4738511407) (398,0.4701047914) (402,0.4790811652) (406,0.4779299641) (410,0.4891579736) (414,0.4940006287) (418,0.4987406357) (422,0.5136265682) (426,0.4992099452) (430,0.4981499893) (434,0.4965968214) (438,0.504608332) (442,0.5073103123) (446,0.5079504028) (450,0.4998191937) (454,0.5148391567) (458,0.5124934466) (462,0.5114405809) (466,0.5060749064) (470,0.5188404022) (474,0.5368406187) (478,0.5425268895) (482,0.5596125489) (486,0.5665286962) (490,0.5527989289) (494,0.5552163348) (498,0.5696121103) (502,0.563352553) (506,0.5575124129) (510,0.5609579163) (514,0.5580261271) (518,0.5435542441) (522,0.5237789708) (526,0.5207306135) (530,0.5224851739) (534,0.5183429526) (538,0.5119401882) (542,0.511841349) (546,0.5088910224) (550,0.5084438991) (554,0.5177476624) (558,0.5053126036) (562,0.5035226213) (566,0.5110060932) (570,0.5211568951) (574,0.5213683972) (578,0.5105000822) (582,0.506973377) (586,0.503956855) (590,0.5067655731) (594,0.5031982482) (598,0.5067784316) (602,0.5065694658) (606,0.5068239596) (610,0.504384983) (614,0.5140375781) (618,0.5142306644) (622,0.5229834286) (626,0.5315095456) (630,0.5359574152) (634,0.5428235818) (638,0.5441133327) (642,0.5450771727) (646,0.5565464162) (650,0.5487413467) (654,0.5480408452) (658,0.554575561) (662,0.5629489882) (666,0.5609285426) (670,0.5535955949) (674,0.5517584318) (678,0.5551055751) (682,0.5536367653) (686,0.5500781164) (690,0.5526523803) (694,0.5419307451) (698,0.538456423) (702,0.5291619324) (706,0.5330713166) (710,0.5355543376) (714,0.5345434696) (718,0.5371393893) (722,0.5276884884) (726,0.5315554645) (730,0.5266243077) (734,0.5342235407) (738,0.5292695371) (742,0.5314253566) (746,0.5335510152) (750,0.545782369) (754,0.5495243971) (758,0.545913191) (762,0.5382676239) (766,0.5524617234) (770,0.5521086253) (774,0.5457900145) (778,0.5420636666) (782,0.5454627924) (786,0.5349860925) (790,0.528850171) (794,0.523899042) (798,0.5228946807) (802,0.5169073762) (806,0.5088809718) (810,0.4976595621) (814,0.4858504909) (818,0.4912217781) (822,0.5024151232) (826,0.5084620912) (830,0.5023586243) (834,0.505155865) (838,0.5159336514) (842,0.5256670387) (846,0.5366610553) (850,0.5335158432) (854,0.5289856617) (858,0.5379095596) (862,0.5439319621) (866,0.547768719) (870,0.534421884) (874,0.521765076) (878,0.5167576199) (882,0.5223644284) (886,0.5248844548) (890,0.518639996) (894,0.5203165612) (898,0.5075970585) (902,0.5133834992) (906,0.5205880518) (910,0.5223743488) (914,0.5202535269) (918,0.5328612499) (922,0.5393434257) (926,0.5363714821) (930,0.5363361964) (934,0.5419781606) (938,0.5453757135) (942,0.5518062399) (946,0.5460226517) (950,0.5449913968) (954,0.5531903049) (958,0.5529037959) (962,0.5555949223) (966,0.5548502893) (970,0.5473107905) (974,0.5503457336) (978,0.5489197122) (982,0.5557832282) (986,0.5611657852) (990,0.5670909749) (994,0.5574734787) (998,0.5708165917)};
\addplot[pubNavy, opacity=0.20, line width=0.35pt] coordinates {(10,1.059148013) (14,1.071460453) (18,1.064671035) (22,1.067154382) (26,1.064801761) (30,1.067373904) (34,1.07425276) (38,1.077615855) (42,1.073380658) (46,1.003960446) (50,0.9441109533) (54,0.8799893798) (58,0.8086173311) (62,0.7400627789) (66,0.6680044278) (70,0.6050841185) (74,0.5346074985) (78,0.4683530404) (82,0.394161355) (86,0.3144812108) (90,0.2433516388) (94,0.204763247) (98,0.1999812822) (102,0.2006853335) (106,0.2055061854) (110,0.2150375209) (114,0.231625703) (118,0.2434000979) (122,0.2716751423) (126,0.2959591607) (130,0.3188818562) (134,0.3452620713) (138,0.3715184735) (142,0.3912849314) (146,0.4122462964) (150,0.4259312379) (154,0.4410894137) (158,0.46010452) (162,0.4735948838) (166,0.4824373319) (170,0.4795804542) (174,0.4722708212) (178,0.4635799582) (182,0.4664240962) (186,0.4724523723) (190,0.4792820251) (194,0.4774389251) (198,0.4810177999) (202,0.4824507017) (206,0.4793132732) (210,0.4737488514) (214,0.4703526534) (218,0.4680425316) (222,0.4721240123) (226,0.4800827502) (230,0.4805842604) (234,0.4779581945) (238,0.4734658515) (242,0.4713277891) (246,0.4784840128) (250,0.4807282703) (254,0.4799241322) (258,0.4822368061) (262,0.4839333666) (266,0.4858019024) (270,0.476613572) (274,0.4781006967) (278,0.4755583511) (282,0.4822134677) (286,0.4764614067) (290,0.4800546554) (294,0.4631337832) (298,0.4583685185) (302,0.4508502999) (306,0.4533890635) (310,0.4479396812) (314,0.4543738511) (318,0.4556498728) (322,0.4559860754) (326,0.4485362161) (330,0.4518289074) (334,0.4520785676) (338,0.4472013334) (342,0.4423373566) (346,0.4540094087) (350,0.4526337841) (354,0.4378024119) (358,0.4458872487) (362,0.4451874854) (366,0.4493436849) (370,0.448229028) (374,0.4578993011) (378,0.4547915672) (382,0.4650307518) (386,0.4636937812) (390,0.4685214785) (394,0.4896584066) (398,0.4920562901) (402,0.4977302459) (406,0.502909207) (410,0.4910457471) (414,0.4880576808) (418,0.4913915286) (422,0.4800562083) (426,0.4731004948) (430,0.4739827075) (434,0.4692035522) (438,0.476864482) (442,0.4807310518) (446,0.4724057876) (450,0.472830377) (454,0.4782217349) (458,0.4899146063) (462,0.4910284602) (466,0.4923806813) (470,0.5038148458) (474,0.5238459306) (478,0.5339393997) (482,0.5334722499) (486,0.5297541019) (490,0.5146923736) (494,0.5119722672) (498,0.5277060157) (502,0.5180478269) (506,0.521630999) (510,0.5200747501) (514,0.5227122571) (518,0.5165665453) (522,0.5067552791) (526,0.4847201039) (530,0.4748980077) (534,0.472082345) (538,0.4860481629) (542,0.4834558155) (546,0.4772304593) (550,0.4733517446) (554,0.477639019) (558,0.4788959767) (562,0.4840971692) (566,0.4821147629) (570,0.4836919452) (574,0.4890082333) (578,0.4814394885) (582,0.4819020214) (586,0.4762841081) (590,0.4773027077) (594,0.4796128554) (598,0.4747337465) (602,0.4727026191) (606,0.4695685263) (610,0.4639780396) (614,0.4557632265) (618,0.457628059) (622,0.4596136993) (626,0.4620539525) (630,0.4623676364) (634,0.4633584547) (638,0.4546943118) (642,0.4533009109) (646,0.4673663518) (650,0.4644524018) (654,0.4659380561) (658,0.4763725995) (662,0.493647498) (666,0.5018012413) (670,0.4949066172) (674,0.5045012541) (678,0.5098237587) (682,0.5219320544) (686,0.5240426433) (690,0.5282762367) (694,0.5268434877) (698,0.5220560328) (702,0.5126326972) (706,0.5172892951) (710,0.5189151634) (714,0.5127164394) (718,0.5090055847) (722,0.5107290868) (726,0.5021969654) (730,0.5066622641) (734,0.5033463274) (738,0.5128886523) (742,0.5149412131) (746,0.525492003) (750,0.5401669003) (754,0.5404110099) (758,0.5314446992) (762,0.5226223275) (766,0.5298802452) (770,0.5394115946) (774,0.5422235408) (778,0.5393720762) (782,0.5455940889) (786,0.5459428043) (790,0.5344540668) (794,0.5370909267) (798,0.5291976215) (802,0.527120227) (806,0.5309121516) (810,0.523496223) (814,0.5289444362) (818,0.5242638804) (822,0.5346122707) (826,0.5407418684) (830,0.5406270566) (834,0.5374421008) (838,0.5529142519) (842,0.5441956716) (846,0.5508965562) (850,0.5499226581) (854,0.5510536227) (858,0.5527972644) (862,0.5535589112) (866,0.5479160412) (870,0.550086616) (874,0.5418024384) (878,0.5307931872) (882,0.5200924425) (886,0.5134236147) (890,0.5106411018) (894,0.5090843672) (898,0.5020655369) (902,0.5096026195) (906,0.5091866705) (910,0.5155605397) (914,0.5131422117) (918,0.5123511769) (922,0.5134440385) (926,0.5138390201) (930,0.5178432181) (934,0.5281212749) (938,0.5349004431) (942,0.5318003742) (946,0.5205430952) (950,0.5198490693) (954,0.5175307504) (958,0.5210633874) (962,0.511592551) (966,0.5153880354) (970,0.5196899908) (974,0.5147357453) (978,0.5286842197) (982,0.524490282) (986,0.5203609365) (990,0.5195468908) (994,0.52054702) (998,0.5272167446)};
\addplot[pubNavy, opacity=0.20, line width=0.35pt] coordinates {(10,1.098758986) (14,1.095816457) (18,1.0864823) (22,1.079447217) (26,1.077640644) (30,1.073652441) (34,1.068818053) (38,1.067479826) (42,1.067033004) (46,0.9958821173) (50,0.9367963149) (54,0.8726142308) (58,0.8055900912) (62,0.732762607) (66,0.6682056506) (70,0.5987153662) (74,0.5372978037) (78,0.4724077361) (82,0.4106534909) (86,0.348375932) (90,0.2801265988) (94,0.2487781765) (98,0.2482272148) (102,0.2472232587) (106,0.2488897831) (110,0.2487718648) (114,0.2480656101) (118,0.2479516035) (122,0.2467738313) (126,0.2441927931) (130,0.242135448) (134,0.2405527884) (138,0.2388579919) (142,0.2374110876) (146,0.2362496251) (150,0.2332854712) (154,0.2267101835) (158,0.2235409444) (162,0.2297451657) (166,0.2285819148) (170,0.2315281845) (174,0.2320520047) (178,0.2346663717) (182,0.2376776522) (186,0.2409967096) (190,0.2464935332) (194,0.252842382) (198,0.2624196306) (202,0.2702837841) (206,0.2813674355) (210,0.2863369812) (214,0.2899548487) (218,0.2965104887) (222,0.3014182576) (226,0.3074074852) (230,0.3156479816) (234,0.3165870226) (238,0.3208715324) (242,0.3254617979) (246,0.3254124695) (250,0.3304824348) (254,0.3280930533) (258,0.3348556079) (262,0.3403175783) (266,0.339132738) (270,0.3371123349) (274,0.3500227677) (278,0.3546386634) (282,0.36253816) (286,0.364784221) (290,0.3712520685) (294,0.3756688727) (298,0.3778195194) (302,0.3775069708) (306,0.3731445774) (310,0.3648476575) (314,0.3739974149) (318,0.374284357) (322,0.3730130853) (326,0.3615294977) (330,0.3595682372) (334,0.3576465675) (338,0.3540546527) (342,0.3401310886) (346,0.3366858213) (350,0.3451487614) (354,0.3435813348) (358,0.347150708) (362,0.3451371046) (366,0.3487585643) (370,0.341365133) (374,0.3483440143) (378,0.3477991286) (382,0.3387935054) (386,0.3403626891) (390,0.353343133) (394,0.3467783051) (398,0.3453571366) (402,0.3475419162) (406,0.3458686981) (410,0.3428510167) (414,0.34237339) (418,0.3402592811) (422,0.3391296415) (426,0.3321726861) (430,0.3384654447) (434,0.3329210185) (438,0.3269763847) (442,0.3350674109) (446,0.3366390574) (450,0.332673852) (454,0.3303112652) (458,0.3244648082) (462,0.3206223394) (466,0.3194254951) (470,0.3175743413) (474,0.3128438787) (478,0.3173047379) (482,0.3197911462) (486,0.320864307) (490,0.309674919) (494,0.2931307293) (498,0.3001464344) (502,0.2963660236) (506,0.2995293818) (510,0.3046445183) (514,0.3176915101) (518,0.32284687) (522,0.3207067417) (526,0.32605172) (530,0.3228368264) (534,0.3137468228) (538,0.3207433954) (542,0.3382808436) (546,0.3340767215) (550,0.3356176171) (554,0.3370536134) (558,0.3225848449) (562,0.321788811) (566,0.3096822137) (570,0.3128606156) (574,0.3061836768) (578,0.3102553581) (582,0.3085930583) (586,0.3093569218) (590,0.3089083856) (594,0.3258063652) (598,0.3133060539) (602,0.3157192785) (606,0.3256157364) (610,0.3235213805) (614,0.3205379104) (618,0.3335520257) (622,0.3391714121) (626,0.3363159168) (630,0.3301785836) (634,0.3205919557) (638,0.3162042558) (642,0.3137796412) (646,0.3160603396) (650,0.320228121) (654,0.3289555153) (658,0.3294166693) (662,0.3354028762) (666,0.3342160126) (670,0.3328588994) (674,0.3412019896) (678,0.3512083192) (682,0.3750399425) (686,0.3839284787) (690,0.397644782) (694,0.3894016274) (698,0.4043954703) (702,0.3873119386) (706,0.389523909) (710,0.3911160143) (714,0.3819550764) (718,0.3819138986) (722,0.3785953755) (726,0.3780792825) (730,0.3611367364) (734,0.3617453648) (738,0.3481135733) (742,0.3589766579) (746,0.3582266929) (750,0.3460887526) (754,0.3513683852) (758,0.3434601014) (762,0.3387236921) (766,0.3638404102) (770,0.3668588804) (774,0.3565819818) (778,0.3617010071) (782,0.3678766846) (786,0.3720737374) (790,0.3743633567) (794,0.3675220128) (798,0.3716763551) (802,0.3622694482) (806,0.356908486) (810,0.3656834214) (814,0.3635794981) (818,0.3603682587) (822,0.3798584963) (826,0.3856519766) (830,0.3810356607) (834,0.3765703844) (838,0.3895598196) (842,0.3824458144) (846,0.3813917296) (850,0.3963120026) (854,0.3901153214) (858,0.3859960293) (862,0.3796449736) (866,0.385338538) (870,0.3617100872) (874,0.3617421116) (878,0.3710721557) (882,0.3701665249) (886,0.3696099412) (890,0.3680966925) (894,0.3717386377) (898,0.3760260051) (902,0.3904942964) (906,0.4034841579) (910,0.4170941312) (914,0.4186832285) (918,0.4231155702) (922,0.4328743281) (926,0.4248616621) (930,0.4211098982) (934,0.4385638377) (938,0.4327566162) (942,0.4414286143) (946,0.4410063296) (950,0.4434581708) (954,0.4368044082) (958,0.4335859315) (962,0.4350539674) (966,0.4417957122) (970,0.4451271354) (974,0.4391333866) (978,0.4396868233) (982,0.4454785538) (986,0.4384999095) (990,0.4423762623) (994,0.4355501136) (998,0.4335823029)};
\addplot[pubNavy, opacity=0.20, line width=0.35pt] coordinates {(10,1.05035823) (14,1.055375609) (18,1.049327914) (22,1.056839656) (26,1.061959033) (30,1.060234296) (34,1.0593705) (38,1.067092508) (42,1.062736153) (46,0.994381443) (50,0.9448879107) (54,0.893977745) (58,0.8415209802) (62,0.7877736817) (66,0.7304267667) (70,0.6756671862) (74,0.6104244046) (78,0.5469908743) (82,0.4870320274) (86,0.4202502636) (90,0.3496699754) (94,0.3218571204) (98,0.320225213) (102,0.308242209) (106,0.3067428719) (110,0.2968692533) (114,0.2942898546) (118,0.2884927031) (122,0.2958275198) (126,0.3008319611) (130,0.3110890701) (134,0.3267831101) (138,0.3384974957) (142,0.3473671596) (146,0.358130679) (150,0.3743827859) (154,0.3900671987) (158,0.4050940719) (162,0.4247822722) (166,0.4423280172) (170,0.4534078953) (174,0.4616891275) (178,0.4665619952) (182,0.4690957031) (186,0.4734367582) (190,0.4864902666) (194,0.4961940985) (198,0.5035926355) (202,0.5105766401) (206,0.5127480369) (210,0.5120463794) (214,0.5053849801) (218,0.5055377611) (222,0.5086377298) (226,0.5068924428) (230,0.5100888927) (234,0.5103849631) (238,0.5064247546) (242,0.5040492299) (246,0.5073425912) (250,0.5052261478) (254,0.4890583049) (258,0.487496334) (262,0.4874951963) (266,0.4901630556) (270,0.4922397501) (274,0.4979387688) (278,0.5006986834) (282,0.5199902343) (286,0.5269415085) (290,0.5427614727) (294,0.5260935803) (298,0.529998867) (302,0.5318843648) (306,0.5413774067) (310,0.5496530987) (314,0.5570497801) (318,0.5684921351) (322,0.5758395474) (326,0.5933136072) (330,0.6195957507) (334,0.6097752136) (338,0.6038470971) (342,0.6066088411) (346,0.6226679857) (350,0.6230148093) (354,0.6260759353) (358,0.628370898) (362,0.6161982007) (366,0.6216321053) (370,0.6241671172) (374,0.6194056422) (378,0.6094796157) (382,0.6129486662) (386,0.6147408566) (390,0.6200960516) (394,0.6207452684) (398,0.6189177983) (402,0.6288805945) (406,0.6395284595) (410,0.6357288902) (414,0.6440104671) (418,0.6646986281) (422,0.6766171482) (426,0.6768220732) (430,0.6982630539) (434,0.6955619196) (438,0.7060516233) (442,0.7124594433) (446,0.719379449) (450,0.7390846303) (454,0.7547405231) (458,0.7546836314) (462,0.7545412727) (466,0.747626026) (470,0.751643217) (474,0.751814176) (478,0.7584743195) (482,0.7746475967) (486,0.7818503868) (490,0.7891863889) (494,0.7957372733) (498,0.8096854243) (502,0.8161626564) (506,0.8346458464) (510,0.8344912879) (514,0.8596897984) (518,0.873316594) (522,0.8596927005) (526,0.8494820251) (530,0.8361385361) (534,0.8420748608) (538,0.852135797) (542,0.8590286538) (546,0.8510262915) (550,0.8546096435) (554,0.8401110282) (558,0.8353543912) (562,0.8297195728) (566,0.8159175218) (570,0.8227110404) (574,0.8292621312) (578,0.8075241278) (582,0.7889413895) (586,0.7644116811) (590,0.747959138) (594,0.7335425638) (598,0.7268057738) (602,0.7096645344) (606,0.6981242286) (610,0.682229998) (614,0.6802343229) (618,0.6740112221) (622,0.6709304846) (626,0.6727708918) (630,0.6653840993) (634,0.6758183523) (638,0.6667400825) (642,0.6726203451) (646,0.6673207175) (650,0.6577524391) (654,0.6612962933) (658,0.6675636582) (662,0.6761130096) (666,0.6691777141) (670,0.6528633945) (674,0.6553395141) (678,0.6494526456) (682,0.6498331743) (686,0.6390738881) (690,0.6552249717) (694,0.6376059945) (698,0.6375504767) (702,0.6230698568) (706,0.6189053836) (710,0.6247975501) (714,0.6101479047) (718,0.6136036269) (722,0.6166882295) (726,0.6160346891) (730,0.6293161395) (734,0.6404072768) (738,0.6470332358) (742,0.6591984484) (746,0.6615699833) (750,0.6683765191) (754,0.6768788004) (758,0.6718705957) (762,0.6769149842) (766,0.6932492852) (770,0.7054060391) (774,0.7063151358) (778,0.7096244557) (782,0.7113959595) (786,0.7135274902) (790,0.7046073367) (794,0.7004529507) (798,0.705759255) (802,0.7009566728) (806,0.702457934) (810,0.6984757392) (814,0.687207228) (818,0.6859404222) (822,0.6898501116) (826,0.6942733145) (830,0.67750923) (834,0.6731869676) (838,0.6738734461) (842,0.677012983) (846,0.6666387929) (850,0.6711557608) (854,0.6667146645) (858,0.6669149589) (862,0.6606176316) (866,0.6595880872) (870,0.6498034943) (874,0.649260146) (878,0.6455374521) (882,0.6485810308) (886,0.6403812776) (890,0.6401507844) (894,0.638048214) (898,0.6379723045) (902,0.636316856) (906,0.6343237967) (910,0.6307978226) (914,0.6339969841) (918,0.6325754503) (922,0.6339840243) (926,0.6274799789) (930,0.6347119877) (934,0.6366860753) (938,0.6333679179) (942,0.6280455284) (946,0.619785982) (950,0.615626761) (954,0.6120748526) (958,0.6192518901) (962,0.618522005) (966,0.6158539233) (970,0.6221056243) (974,0.617668941) (978,0.6111256957) (982,0.6111356071) (986,0.6051760617) (990,0.6162631999) (994,0.6196104157) (998,0.621850053)};
\addplot[pubNavy, opacity=0.20, line width=0.35pt] coordinates {(10,1.057971364) (14,1.0597948) (18,1.04751989) (22,1.043298132) (26,1.044138859) (30,1.044135086) (34,1.04494577) (38,1.051893174) (42,1.050688268) (46,0.9811082734) (50,0.9224035013) (54,0.8538577464) (58,0.7896841307) (62,0.7201227067) (66,0.657989974) (70,0.5950590168) (74,0.5363227755) (78,0.4727556956) (82,0.4127843203) (86,0.354560539) (90,0.2857012739) (94,0.2597611887) (98,0.2573822721) (102,0.2593819332) (106,0.2604646527) (110,0.2638948078) (114,0.2663555963) (118,0.2712097076) (122,0.2746581888) (126,0.2722313726) (130,0.2774079699) (134,0.2855455976) (138,0.2897515688) (142,0.2950593644) (146,0.3054782393) (150,0.3171537351) (154,0.328987936) (158,0.342819131) (162,0.3672230055) (166,0.3779398998) (170,0.3906664534) (174,0.4012114334) (178,0.4084162823) (182,0.4166732858) (186,0.4225461126) (190,0.4302384897) (194,0.4362527316) (198,0.4421836549) (202,0.4486381901) (206,0.4424769556) (210,0.4375932233) (214,0.4302119538) (218,0.4223893036) (222,0.4156338338) (226,0.4171673665) (230,0.4097353526) (234,0.3996733641) (238,0.3964838659) (242,0.3927846192) (246,0.3872109758) (250,0.3771006993) (254,0.3717007242) (258,0.3710563159) (262,0.3727650653) (266,0.3680198618) (270,0.3663408569) (274,0.358522057) (278,0.3535897924) (282,0.3503883246) (286,0.3491311917) (290,0.3454024757) (294,0.3361905736) (298,0.3277630912) (302,0.3252786844) (306,0.3186759049) (310,0.3118408127) (314,0.3075236934) (318,0.3043136474) (322,0.2998572563) (326,0.3000235922) (330,0.3022781199) (334,0.2934434624) (338,0.2926796954) (342,0.288658392) (346,0.300174817) (350,0.3008159651) (354,0.3024474393) (358,0.3077140372) (362,0.3087767964) (366,0.3123488278) (370,0.3130835446) (374,0.3207372172) (378,0.3273723588) (382,0.3485564835) (386,0.3468406292) (390,0.3496624239) (394,0.3548170404) (398,0.3546301938) (402,0.3628785168) (406,0.3730950467) (410,0.3647664006) (414,0.3668940295) (418,0.3731066181) (422,0.371295384) (426,0.3618865139) (430,0.3466286122) (434,0.3377621879) (438,0.327559826) (442,0.3278627117) (446,0.3330912013) (450,0.3265315782) (454,0.3212276868) (458,0.3143099265) (462,0.3056547576) (466,0.3048576516) (470,0.3140917075) (474,0.3253466965) (478,0.3452071214) (482,0.3501822237) (486,0.3700647838) (490,0.3720351944) (494,0.3757024792) (498,0.3841687545) (502,0.3889924669) (506,0.3929326958) (510,0.3994950155) (514,0.4191079174) (518,0.4384184082) (522,0.417731194) (526,0.4055372959) (530,0.3904641341) (534,0.3937032845) (538,0.3937468573) (542,0.4025815734) (546,0.3935389645) (550,0.3874161618) (554,0.372418944) (558,0.3844082953) (562,0.376761468) (566,0.3699267582) (570,0.3719077922) (574,0.3870634185) (578,0.3854630316) (582,0.3886640714) (586,0.3877754543) (590,0.3874885715) (594,0.3783608693) (598,0.3666477495) (602,0.3792463891) (606,0.3941090903) (610,0.3910475242) (614,0.3775177159) (618,0.3851252992) (622,0.3823352727) (626,0.3935493108) (630,0.3926106362) (634,0.4033111017) (638,0.3958197408) (642,0.4012175355) (646,0.4227241467) (650,0.4214122829) (654,0.4149161459) (658,0.40432041) (662,0.4150275271) (666,0.433225808) (670,0.4284334178) (674,0.4195685078) (678,0.4108181925) (682,0.4071954131) (686,0.4144456144) (690,0.4129475977) (694,0.4108189676) (698,0.4054611751) (702,0.4105568669) (706,0.4176393537) (710,0.4145679848) (714,0.4152040076) (718,0.4018380663) (722,0.4038876293) (726,0.4049706389) (730,0.4217192771) (734,0.4145661475) (738,0.4013244489) (742,0.3996725117) (746,0.4011677652) (750,0.3859971042) (754,0.3877729812) (758,0.3820766187) (762,0.3920742459) (766,0.3962896056) (770,0.3992541023) (774,0.4011449674) (778,0.3899182252) (782,0.3873281095) (786,0.388134753) (790,0.3841223541) (794,0.3787342269) (798,0.3787940779) (802,0.3743370432) (806,0.3715114617) (810,0.3652585079) (814,0.3599456666) (818,0.3566962158) (822,0.3571353533) (826,0.3457974042) (830,0.344768754) (834,0.3539917573) (838,0.3477664186) (842,0.3521894857) (846,0.35287115) (850,0.3670568571) (854,0.3697033401) (858,0.3715552462) (862,0.3771710004) (866,0.3744491636) (870,0.3695186892) (874,0.377541496) (878,0.3803772144) (882,0.374049139) (886,0.383782668) (890,0.3874994545) (894,0.3743658576) (898,0.3777232556) (902,0.3678069755) (906,0.3635801531) (910,0.3710190689) (914,0.3601353391) (918,0.3686817515) (922,0.3736828594) (926,0.3816102896) (930,0.4020660708) (934,0.4085516697) (938,0.4106981988) (942,0.4012329575) (946,0.4038917109) (950,0.3922013558) (954,0.406734871) (958,0.4013480356) (962,0.392750391) (966,0.3932746892) (970,0.3881437198) (974,0.3783859915) (978,0.3675206324) (982,0.3645903058) (986,0.3644746144) (990,0.3665050307) (994,0.3775314299) (998,0.3730628432)};
\addplot[pubNavy, opacity=0.20, line width=0.35pt] coordinates {(10,1.051387425) (14,1.042500424) (18,1.044861427) (22,1.048950692) (26,1.05045949) (30,1.047179587) (34,1.054427441) (38,1.055718501) (42,1.051123358) (46,0.986241518) (50,0.9277664157) (54,0.8661186207) (58,0.8001228656) (62,0.7349882005) (66,0.6756946894) (70,0.6146645289) (74,0.5601121498) (78,0.5019107218) (82,0.4449950383) (86,0.3844143273) (90,0.3212081652) (94,0.3045360631) (98,0.3210455267) (102,0.3427569578) (106,0.3703770668) (110,0.3957894576) (114,0.4202853054) (118,0.4427324804) (122,0.4588063649) (126,0.4721800096) (130,0.4940176693) (134,0.5096723954) (138,0.5195517048) (142,0.5366271215) (146,0.5502146279) (150,0.5546370762) (154,0.5517165085) (158,0.5704087313) (162,0.5785467507) (166,0.5839172163) (170,0.5958833539) (174,0.5981649055) (178,0.5914961301) (182,0.6133143321) (186,0.6273014961) (190,0.64206811) (194,0.6551503307) (198,0.6635667338) (202,0.6809901216) (206,0.6717480791) (210,0.6672389063) (214,0.6646262051) (218,0.6560915302) (222,0.662424272) (226,0.6699816095) (230,0.662243647) (234,0.6411376939) (238,0.6194392286) (242,0.6111174224) (246,0.6022367671) (250,0.5972249087) (254,0.5916197434) (258,0.5911292134) (262,0.5895605901) (266,0.5876290821) (270,0.5817020694) (274,0.5713721114) (278,0.569758093) (282,0.5734049762) (286,0.5792765226) (290,0.5811239876) (294,0.5752695642) (298,0.5648693145) (302,0.5491083025) (306,0.5417160235) (310,0.5273226474) (314,0.5114659823) (318,0.5097878952) (322,0.4996896322) (326,0.4953081792) (330,0.4986621237) (334,0.5035628601) (338,0.4965897715) (342,0.4985579509) (346,0.5037414321) (350,0.5141545965) (354,0.52256616) (358,0.532405115) (362,0.5427607389) (366,0.5478410741) (370,0.5535886568) (374,0.5655379694) (378,0.574384566) (382,0.5773658524) (386,0.5816448426) (390,0.5948216864) (394,0.5997152316) (398,0.6025018143) (402,0.6188583786) (406,0.6203587975) (410,0.6240542023) (414,0.6313183994) (418,0.6474107308) (422,0.6582006503) (426,0.6520608358) (430,0.6604321546) (434,0.6631500111) (438,0.6626576677) (442,0.6624232295) (446,0.6675303542) (450,0.6620427708) (454,0.6614104087) (458,0.6658074332) (462,0.6711417325) (466,0.6699921151) (470,0.6780237931) (474,0.6769411208) (478,0.6901008806) (482,0.6996405535) (486,0.6970806394) (490,0.6851596439) (494,0.6950235872) (498,0.7031652257) (502,0.70948562) (506,0.719446637) (510,0.7214726893) (514,0.7300594507) (518,0.7206506734) (522,0.7136450234) (526,0.709636313) (530,0.7078549283) (534,0.6955256003) (538,0.7152517035) (542,0.7165996079) (546,0.7079043821) (550,0.6990074548) (554,0.69580945) (558,0.6972445342) (562,0.6957584776) (566,0.6966883817) (570,0.7053403477) (574,0.7076035977) (578,0.6951150596) (582,0.7021410386) (586,0.7047766318) (590,0.6992790146) (594,0.7012827007) (598,0.7181695278) (602,0.7230739339) (606,0.7142980854) (610,0.7106907287) (614,0.7165871586) (618,0.7178430564) (622,0.7270654758) (626,0.7392891235) (630,0.7407991847) (634,0.7480115264) (638,0.7380974075) (642,0.7497869951) (646,0.7502153754) (650,0.7396504916) (654,0.7443565884) (658,0.7728893457) (662,0.7938003891) (666,0.7892058965) (670,0.77326265) (674,0.7778527347) (678,0.7682934947) (682,0.7647094965) (686,0.7594242534) (690,0.7596767729) (694,0.7568039407) (698,0.7508353096) (702,0.747819462) (706,0.7362802177) (710,0.7250375852) (714,0.7164912534) (718,0.7172935156) (722,0.7052606786) (726,0.7083377283) (730,0.7191769718) (734,0.7284095366) (738,0.7335057811) (742,0.7425212619) (746,0.7480737281) (750,0.7653430178) (754,0.7624236044) (758,0.7633168538) (762,0.7650738704) (766,0.7829548234) (770,0.8033095646) (774,0.8081637839) (778,0.8015856377) (782,0.7887068498) (786,0.7846860087) (790,0.7842696342) (794,0.7757507539) (798,0.7628314352) (802,0.7562572439) (806,0.7436586422) (810,0.7244776369) (814,0.7267377153) (818,0.711747193) (822,0.7126669928) (826,0.7121337257) (830,0.7078255365) (834,0.7050103529) (838,0.7155843373) (842,0.7085774934) (846,0.7114291492) (850,0.7040888708) (854,0.7033419594) (858,0.7030987723) (862,0.7120568149) (866,0.7044722129) (870,0.7101660968) (874,0.7103251452) (878,0.7042521145) (882,0.6935759098) (886,0.6765150674) (890,0.6734144051) (894,0.6709670707) (898,0.6769945911) (902,0.6647712387) (906,0.6698560263) (910,0.6585612448) (914,0.6548487496) (918,0.6545964439) (922,0.6488658161) (926,0.6432544953) (930,0.6567088942) (934,0.6677840499) (938,0.6622573919) (942,0.6485454867) (946,0.6421621252) (950,0.6405252759) (954,0.6498439488) (958,0.6535159993) (962,0.6468845704) (966,0.6500191744) (970,0.6633349154) (974,0.6675458383) (978,0.6786523482) (982,0.6689826202) (986,0.6674393208) (990,0.6649763385) (994,0.6661895158) (998,0.6821617076)};
\addplot[pubNavy, opacity=0.20, line width=0.35pt] coordinates {(10,1.086211579) (14,1.078255284) (18,1.074501401) (22,1.079539561) (26,1.072485434) (30,1.070464325) (34,1.073018411) (38,1.066748609) (42,1.05719507) (46,0.9919199582) (50,0.9319060262) (54,0.868872406) (58,0.8038382575) (62,0.7335362618) (66,0.6702721406) (70,0.6011436027) (74,0.5413080841) (78,0.4746337186) (82,0.4055258742) (86,0.3432659136) (90,0.2829766729) (94,0.2592702447) (98,0.2622675472) (102,0.2698977994) (106,0.2752833385) (110,0.281926412) (114,0.2953141079) (118,0.3110009356) (122,0.3334541323) (126,0.3526321424) (130,0.3838237956) (134,0.4212276281) (138,0.4486164794) (142,0.4732003325) (146,0.4997549623) (150,0.5309706075) (154,0.5539579114) (158,0.6009208185) (162,0.6328645783) (166,0.6606452784) (170,0.6847861063) (174,0.6962500017) (178,0.7064377177) (182,0.7161981516) (186,0.7233576298) (190,0.7276491508) (194,0.7358283917) (198,0.7313737796) (202,0.7279529781) (206,0.7179376794) (210,0.6998148267) (214,0.6675489526) (218,0.6510503033) (222,0.6320564356) (226,0.6093454432) (230,0.597489822) (234,0.5706064683) (238,0.5533172277) (242,0.5444047664) (246,0.5336660104) (250,0.5187311713) (254,0.499045164) (258,0.4887595644) (262,0.4923297696) (266,0.4986280489) (270,0.4942178967) (274,0.5038597263) (278,0.5126172252) (282,0.526878142) (286,0.5310341976) (290,0.5374374621) (294,0.5199904346) (298,0.5225921347) (302,0.5272709725) (306,0.5319954346) (310,0.5261863869) (314,0.523134992) (318,0.5197819782) (322,0.5138136086) (326,0.5018428522) (330,0.4964000556) (334,0.4891732438) (338,0.4709447317) (342,0.4674378076) (346,0.4772246615) (350,0.4798566591) (354,0.4757266072) (358,0.4775485892) (362,0.4759605162) (366,0.4789019506) (370,0.4766325931) (374,0.4829613479) (378,0.4800489868) (382,0.4793357088) (386,0.4815774459) (390,0.4931920966) (394,0.4901028762) (398,0.489200007) (402,0.4916344415) (406,0.4964179611) (410,0.4858098636) (414,0.4817666774) (418,0.4834817444) (422,0.4736871135) (426,0.4677372064) (430,0.4627828462) (434,0.4617187411) (438,0.4540787701) (442,0.4566984592) (446,0.4641853474) (450,0.4579890739) (454,0.4568186882) (458,0.457624704) (462,0.4581234993) (466,0.4594887132) (470,0.4671478981) (474,0.4651995008) (478,0.4823380547) (482,0.4760126495) (486,0.4875928554) (490,0.4895208598) (494,0.4754083193) (498,0.4825169893) (502,0.4885926835) (506,0.4952766841) (510,0.5042942377) (514,0.5147904212) (518,0.527650478) (522,0.5269879471) (526,0.5268605961) (530,0.5155660271) (534,0.5010374677) (538,0.50656545) (542,0.5093788444) (546,0.508244097) (550,0.4990466603) (554,0.4918844845) (558,0.4863922881) (562,0.4747382274) (566,0.4612967336) (570,0.4705575099) (574,0.4652634567) (578,0.4776065701) (582,0.4821660891) (586,0.4769806362) (590,0.4784402608) (594,0.4953945868) (598,0.489386486) (602,0.5004200391) (606,0.4977429405) (610,0.4986525311) (614,0.498516338) (618,0.500212579) (622,0.4955830725) (626,0.5000521521) (630,0.4992271936) (634,0.5115009457) (638,0.5217432874) (642,0.5204399566) (646,0.5219422542) (650,0.5057380258) (654,0.5137390184) (658,0.5126459125) (662,0.5178600345) (666,0.5201926867) (670,0.5169097871) (674,0.5309598362) (678,0.5269876506) (682,0.5253430596) (686,0.5227691269) (690,0.5113904121) (694,0.5104721274) (698,0.5225666427) (702,0.5153941289) (706,0.518229326) (710,0.5153338844) (714,0.5070547452) (718,0.4983944503) (722,0.4987063648) (726,0.4958126125) (730,0.4949953898) (734,0.499053561) (738,0.5027566475) (742,0.506960107) (746,0.4903443072) (750,0.4906336026) (754,0.5000187505) (758,0.4942830302) (762,0.504532149) (766,0.5224195448) (770,0.5233910427) (774,0.5315423602) (778,0.5180405716) (782,0.5042646319) (786,0.5081773445) (790,0.5074964083) (794,0.5041704183) (798,0.5040409185) (802,0.4908961056) (806,0.4953265749) (810,0.4899385654) (814,0.4943434748) (818,0.488274135) (822,0.4977998959) (826,0.4949636632) (830,0.5006815867) (834,0.4920394729) (838,0.4931754404) (842,0.4864223604) (846,0.4890096256) (850,0.4937617494) (854,0.491931022) (858,0.4870926258) (862,0.4777566055) (866,0.4771673401) (870,0.4704864149) (874,0.4643420225) (878,0.4727737376) (882,0.4605321178) (886,0.4698699369) (890,0.4722227723) (894,0.4643323687) (898,0.4631709285) (902,0.4621995657) (906,0.4552560511) (910,0.4520036741) (914,0.4433951289) (918,0.4405591594) (922,0.4507773434) (926,0.4547437925) (930,0.4458572904) (934,0.4589302555) (938,0.4537261421) (942,0.4477694884) (946,0.4575006428) (950,0.4550045093) (954,0.4553948064) (958,0.4592966709) (962,0.4652922331) (966,0.488199401) (970,0.4901855193) (974,0.479961023) (978,0.481685051) (982,0.4902981391) (986,0.4738605358) (990,0.4836147462) (994,0.4797989795) (998,0.4913954829)};
\addplot[pubNavy, opacity=0.20, line width=0.35pt] coordinates {(10,1.076581584) (14,1.082034682) (18,1.080205035) (22,1.071921075) (26,1.064138603) (30,1.062709096) (34,1.063603132) (38,1.062225905) (42,1.052941816) (46,0.9842223977) (50,0.9283148557) (54,0.8660355483) (58,0.8002745526) (62,0.7315866281) (66,0.6675381303) (70,0.6015842811) (74,0.5421861101) (78,0.4792984559) (82,0.4142522623) (86,0.3525220318) (90,0.2932618278) (94,0.2667708872) (98,0.2644452389) (102,0.2553353116) (106,0.2446927156) (110,0.2356378821) (114,0.2252835769) (118,0.2157724033) (122,0.2038578046) (126,0.1924977724) (130,0.1865872376) (134,0.1796373688) (138,0.1737108368) (142,0.1734631184) (146,0.1731782222) (150,0.174975301) (154,0.1783109841) (158,0.186436438) (162,0.2062671284) (166,0.2215109645) (170,0.2359593541) (174,0.2534820822) (178,0.2664585605) (182,0.2755655454) (186,0.2915562729) (190,0.300501268) (194,0.3094524488) (198,0.324131468) (202,0.3373181958) (206,0.3499296036) (210,0.3611900403) (214,0.3608826234) (218,0.3738723973) (222,0.3845961176) (226,0.3961770689) (230,0.4079620521) (234,0.4100928329) (238,0.4192498252) (242,0.430533193) (246,0.4502732008) (250,0.4562400595) (254,0.4540801237) (258,0.4559838251) (262,0.4685524669) (266,0.4753706287) (270,0.4703855006) (274,0.4767739165) (278,0.4850834845) (282,0.5016560613) (286,0.5024602146) (290,0.5079841238) (294,0.4832377997) (298,0.4809170311) (302,0.4872944069) (306,0.4968984866) (310,0.4878733082) (314,0.4930317606) (318,0.4957214539) (322,0.5037272209) (326,0.4986060073) (330,0.5027448216) (334,0.5001392216) (338,0.5073671338) (342,0.5134623538) (346,0.5369957231) (350,0.5375576042) (354,0.5328426474) (358,0.5488471889) (362,0.5544750761) (366,0.5724252996) (370,0.5673103631) (374,0.5813769501) (378,0.5790148828) (382,0.5827358655) (386,0.5896501424) (390,0.601552921) (394,0.6086306112) (398,0.6172139728) (402,0.6293258994) (406,0.6344247561) (410,0.6389490706) (414,0.6417922085) (418,0.6513390182) (422,0.6532823193) (426,0.6485724129) (430,0.6557734257) (434,0.6484019215) (438,0.6542539825) (442,0.6637507151) (446,0.6589139892) (450,0.6570969373) (454,0.6654495848) (458,0.6691645839) (462,0.6750408593) (466,0.6749901728) (470,0.6863946531) (474,0.6975693957) (478,0.7224025154) (482,0.7404330609) (486,0.7313204373) (490,0.7141690688) (494,0.7241517578) (498,0.7424194848) (502,0.7345853775) (506,0.7289581591) (510,0.7383709627) (514,0.7425681262) (518,0.7439697243) (522,0.7378641813) (526,0.7247236278) (530,0.7166751437) (534,0.7040850694) (538,0.7183743029) (542,0.7125204162) (546,0.7060616281) (550,0.6901677074) (554,0.6911736193) (558,0.6806410536) (562,0.6726077003) (566,0.6508865131) (570,0.6567178026) (574,0.6464552183) (578,0.6411436621) (582,0.6316549926) (586,0.628812968) (590,0.628925988) (594,0.6310765649) (598,0.6292260651) (602,0.6332114054) (606,0.6298986014) (610,0.6224580673) (614,0.6157788287) (618,0.6233201364) (622,0.6223941632) (626,0.6141343868) (630,0.6103696709) (634,0.6077769032) (638,0.5910923516) (642,0.5994672163) (646,0.592201113) (650,0.5792131459) (654,0.5742606861) (658,0.5737829239) (662,0.5880441639) (666,0.5872563047) (670,0.5803417176) (674,0.5965307716) (678,0.5948887495) (682,0.6016191928) (686,0.6047398999) (690,0.6174121188) (694,0.6141202974) (698,0.6368945044) (702,0.6349487581) (706,0.6454601501) (710,0.6586028765) (714,0.6543761214) (718,0.6736852258) (722,0.6634150169) (726,0.6613321053) (730,0.6634829099) (734,0.6746576015) (738,0.6782657754) (742,0.7036173136) (746,0.6927777066) (750,0.679851819) (754,0.6755040648) (758,0.6654262766) (762,0.6573307942) (766,0.6702748274) (770,0.6781363739) (774,0.6761210709) (778,0.6647645027) (782,0.6707820083) (786,0.6731805026) (790,0.6671685184) (794,0.6521968941) (798,0.644927545) (802,0.6536817348) (806,0.6519976199) (810,0.6465681651) (814,0.6304291724) (818,0.6165721582) (822,0.614852585) (826,0.6219119097) (830,0.6127670708) (834,0.6024502596) (838,0.6005187116) (842,0.5903426382) (846,0.5864514884) (850,0.5850825409) (854,0.5685616388) (858,0.5579263641) (862,0.5598212058) (866,0.5549629815) (870,0.5493522045) (874,0.5481868662) (878,0.5446777836) (882,0.5333727825) (886,0.5337935348) (890,0.5254222133) (894,0.5252656474) (898,0.5235762446) (902,0.5281175308) (906,0.5295429904) (910,0.5340456452) (914,0.5394003414) (918,0.5417230441) (922,0.5508190132) (926,0.5435909433) (930,0.5515269549) (934,0.5721216038) (938,0.5582206505) (942,0.579815254) (946,0.5858341362) (950,0.5770976109) (954,0.5791490012) (958,0.5873492624) (962,0.5821838471) (966,0.5893563295) (970,0.5913095602) (974,0.5988146021) (978,0.6191287964) (982,0.6263843652) (986,0.6211424943) (990,0.6241143189) (994,0.6148220198) (998,0.619431853)};
\addplot[pubNavy, opacity=0.20, line width=0.35pt] coordinates {(10,1.074469417) (14,1.07112408) (18,1.069546642) (22,1.082464925) (26,1.078715523) (30,1.080493402) (34,1.080294608) (38,1.074514955) (42,1.06436949) (46,0.9993076861) (50,0.9446274739) (54,0.8873054707) (58,0.8243885432) (62,0.7658858887) (66,0.7039646229) (70,0.6378812922) (74,0.5767847074) (78,0.5237231853) (82,0.4593008239) (86,0.4009838468) (90,0.3494798055) (94,0.3280362161) (98,0.3331698364) (102,0.334251403) (106,0.3397431414) (110,0.3448794515) (114,0.3458128801) (118,0.3456292221) (122,0.3536689459) (126,0.3576034729) (130,0.3597635386) (134,0.3756928004) (138,0.3846834073) (142,0.3900950544) (146,0.3987557205) (150,0.4068832002) (154,0.4081129131) (158,0.4280627524) (162,0.445036712) (166,0.4735651383) (170,0.4911126872) (174,0.5062155803) (178,0.507615898) (182,0.5157120187) (186,0.532009366) (190,0.5416182782) (194,0.5488724119) (198,0.5489411594) (202,0.5551639905) (206,0.5572780772) (210,0.5514209725) (214,0.5408056518) (218,0.5356556966) (222,0.531075498) (226,0.5372069196) (230,0.5312222603) (234,0.5200057926) (238,0.5031889041) (242,0.4983106276) (246,0.4977954911) (250,0.5053178542) (254,0.4937931679) (258,0.4833450936) (262,0.4787405313) (266,0.4770986775) (270,0.476156006) (274,0.4870888994) (278,0.4749275846) (282,0.476706319) (286,0.4771972784) (290,0.478563383) (294,0.4657088233) (298,0.454586711) (302,0.4492974772) (306,0.4529788421) (310,0.4512851511) (314,0.450667783) (318,0.4445863002) (322,0.4438678005) (326,0.4367415896) (330,0.4420749289) (334,0.4468859601) (338,0.4422561712) (342,0.4402519686) (346,0.4506877391) (350,0.4517852338) (354,0.4501648415) (358,0.4544175402) (362,0.4498546955) (366,0.4511694535) (370,0.4516542876) (374,0.4511798895) (378,0.45275491) (382,0.446637607) (386,0.4423538587) (390,0.4422242927) (394,0.4381693888) (398,0.4358998515) (402,0.4508158671) (406,0.4468876927) (410,0.444756579) (414,0.4471928132) (418,0.4508535083) (422,0.4378956437) (426,0.4273723105) (430,0.4306775835) (434,0.4433966409) (438,0.4575448755) (442,0.4766733579) (446,0.4957439361) (450,0.4878726657) (454,0.4962577787) (458,0.506474374) (462,0.5089794125) (466,0.5211870258) (470,0.5297343735) (474,0.5452709148) (478,0.5591213377) (482,0.5722534625) (486,0.5704211606) (490,0.5604090702) (494,0.5578962548) (498,0.5777961313) (502,0.5922902933) (506,0.595520084) (510,0.6023976244) (514,0.597971194) (518,0.6122133158) (522,0.6154579886) (526,0.621574112) (530,0.6068654607) (534,0.6050691943) (538,0.6142379877) (542,0.6206792807) (546,0.6218403682) (550,0.6195533063) (554,0.6264982791) (558,0.6339826409) (562,0.6398913052) (566,0.6361024816) (570,0.6366210734) (574,0.6399678568) (578,0.6503635085) (582,0.6607289143) (586,0.6590639558) (590,0.6592651594) (594,0.6749408123) (598,0.6772072638) (602,0.6805047784) (606,0.6641305178) (610,0.6514990801) (614,0.6493212861) (618,0.6519526201) (622,0.6605150633) (626,0.6640075788) (630,0.6510158249) (634,0.6527826083) (638,0.6618323195) (642,0.6434616785) (646,0.6393654649) (650,0.6309320921) (654,0.6404351804) (658,0.6580008608) (662,0.6714759545) (666,0.670866366) (670,0.6746694528) (674,0.6752439998) (678,0.6610136044) (682,0.6749317884) (686,0.6705395585) (690,0.6719943636) (694,0.6860818208) (698,0.6919051667) (702,0.6829129284) (706,0.6919990746) (710,0.6944915356) (714,0.6744191765) (718,0.6673554075) (722,0.6692464491) (726,0.6777632376) (730,0.6669507563) (734,0.6692594037) (738,0.6635479031) (742,0.6652151892) (746,0.6663370195) (750,0.6614233345) (754,0.659805515) (758,0.6437023074) (762,0.6379660211) (766,0.6643117653) (770,0.6686264189) (774,0.6570101807) (778,0.6564514431) (782,0.6458967208) (786,0.6391993215) (790,0.6313148856) (794,0.6257429496) (798,0.6151943978) (802,0.5987804146) (806,0.5822592912) (810,0.5808490704) (814,0.5765386574) (818,0.5618204543) (822,0.5741045166) (826,0.5651860131) (830,0.5569416757) (834,0.5447798226) (838,0.560250739) (842,0.5587486512) (846,0.562098392) (850,0.5673825892) (854,0.5683591053) (858,0.5708643019) (862,0.5677391413) (866,0.5615671954) (870,0.5527624407) (874,0.5524943764) (878,0.5547972772) (882,0.554820922) (886,0.5460160395) (890,0.5509756089) (894,0.5588853829) (898,0.5526136944) (902,0.5640785045) (906,0.5694127542) (910,0.5786426878) (914,0.5753198736) (918,0.582608491) (922,0.5847921746) (926,0.5876169458) (930,0.5828224304) (934,0.6023837196) (938,0.6017762123) (942,0.6043074175) (946,0.5915765361) (950,0.5886761662) (954,0.5878595238) (958,0.5947491212) (962,0.5925668058) (966,0.6016967983) (970,0.6104448794) (974,0.6150491509) (978,0.6288217086) (982,0.6315735982) (986,0.6258500809) (990,0.6176442337) (994,0.6091640427) (998,0.6209164754)};
\addplot[pubNavy, opacity=0.20, line width=0.35pt] coordinates {(10,1.093200367) (14,1.085252192) (18,1.077834759) (22,1.082199436) (26,1.082117706) (30,1.075642248) (34,1.079297055) (38,1.080942742) (42,1.078426388) (46,1.021075558) (50,0.9710879614) (54,0.918792401) (58,0.8559493607) (62,0.7889501143) (66,0.7349365872) (70,0.6799390579) (74,0.620157051) (78,0.5670091942) (82,0.512999486) (86,0.4566955369) (90,0.399721742) (94,0.3751402472) (98,0.3782631906) (102,0.3774807738) (106,0.3805638506) (110,0.3871886736) (114,0.3910684245) (118,0.3925753423) (122,0.3892781567) (126,0.3803438385) (130,0.3718388785) (134,0.3654077197) (138,0.355105449) (142,0.3445435274) (146,0.3305525861) (150,0.3159304035) (154,0.3043193739) (158,0.2884216929) (162,0.2747731652) (166,0.258811742) (170,0.2558910878) (174,0.2546141058) (178,0.2494505522) (182,0.2492717181) (186,0.2523300483) (190,0.2545026162) (194,0.2584119324) (198,0.2659340808) (202,0.2741788197) (206,0.286649837) (210,0.304956275) (214,0.31974238) (218,0.3306144683) (222,0.3429405592) (226,0.3618671517) (230,0.3712055833) (234,0.3755617328) (238,0.3786570815) (242,0.3837861963) (246,0.389060844) (250,0.3976256226) (254,0.399312) (258,0.4004033379) (262,0.3984626271) (266,0.3975074708) (270,0.3910879134) (274,0.3923545157) (278,0.3907948377) (282,0.3907430049) (286,0.3948832474) (290,0.405220091) (294,0.396553822) (298,0.3935069558) (302,0.3936060785) (306,0.3933536759) (310,0.3898844631) (314,0.3938340306) (318,0.3977183093) (322,0.4027901971) (326,0.4002341187) (330,0.4089863595) (334,0.4108350164) (338,0.411062238) (342,0.4175355527) (346,0.4247963532) (350,0.4306656169) (354,0.4378963783) (358,0.4498728774) (362,0.4599718368) (366,0.4683653929) (370,0.4748289049) (374,0.4930447484) (378,0.5046465681) (382,0.5195876542) (386,0.5270096662) (390,0.5364950348) (394,0.552573575) (398,0.5656591259) (402,0.5713884049) (406,0.5803208592) (410,0.5822133291) (414,0.5848727929) (418,0.5963444922) (422,0.5949147802) (426,0.592024509) (430,0.5864401606) (434,0.5879017337) (438,0.5878480904) (442,0.5952032321) (446,0.5946365091) (450,0.5963678221) (454,0.6042294131) (458,0.6053770688) (462,0.6083077566) (466,0.6170911283) (470,0.6238751358) (474,0.640022039) (478,0.6552310703) (482,0.663912895) (486,0.6641780138) (490,0.6613773065) (494,0.6622121719) (498,0.670847143) (502,0.6669487816) (506,0.6662889086) (510,0.6694506901) (514,0.6651160816) (518,0.6561034847) (522,0.6464436021) (526,0.630892207) (530,0.6162514136) (534,0.6151359757) (538,0.6235103675) (542,0.6230589649) (546,0.6228707235) (550,0.6175352227) (554,0.6288263817) (558,0.6254764518) (562,0.6306602305) (566,0.6292838709) (570,0.6329742939) (574,0.64162794) (578,0.6450573874) (582,0.6443854548) (586,0.6454399426) (590,0.6407653502) (594,0.6416382184) (598,0.632570684) (602,0.6313665793) (606,0.6224183488) (610,0.619469389) (614,0.611502999) (618,0.6149823639) (622,0.6107725207) (626,0.6090961315) (630,0.6117970313) (634,0.6055092689) (638,0.6008676547) (642,0.6022069186) (646,0.6134704842) (650,0.6058740997) (654,0.6054654712) (658,0.6110733617) (662,0.611805228) (666,0.6168901887) (670,0.6145442123) (674,0.6235143792) (678,0.6129424224) (682,0.6160502452) (686,0.6198881495) (690,0.6251138387) (694,0.6211227149) (698,0.6249319599) (702,0.6240951639) (706,0.6310415084) (710,0.6477676397) (714,0.6377222072) (718,0.6389848273) (722,0.6364522451) (726,0.6338877185) (730,0.6329995845) (734,0.6307978359) (738,0.6176524178) (742,0.6230795521) (746,0.6162341456) (750,0.6088639282) (754,0.6048223911) (758,0.5912914551) (762,0.5799452084) (766,0.5748516416) (770,0.5760114921) (774,0.5721129975) (778,0.5699095338) (782,0.5719644173) (786,0.5729678923) (790,0.5831247386) (794,0.5769884518) (798,0.575642254) (802,0.5724224399) (806,0.5705681814) (810,0.5606147111) (814,0.5617648957) (818,0.5669562109) (822,0.577492357) (826,0.5811725615) (830,0.573199982) (834,0.5793455269) (838,0.5763722611) (842,0.5821647668) (846,0.5721071722) (850,0.574122102) (854,0.5680558241) (858,0.5704510927) (862,0.5743278056) (866,0.5745612601) (870,0.5569811019) (874,0.5472111983) (878,0.5397089888) (882,0.5282681348) (886,0.5222255343) (890,0.5138641875) (894,0.5174710711) (898,0.5184478991) (902,0.5231941407) (906,0.5240738901) (910,0.5166183695) (914,0.5130868138) (918,0.5190129479) (922,0.5189530631) (926,0.5246944337) (930,0.533367349) (934,0.5427621399) (938,0.5386494064) (942,0.5406700833) (946,0.5428907515) (950,0.5410077954) (954,0.543108006) (958,0.5515267042) (962,0.5545075297) (966,0.5496801726) (970,0.5391971497) (974,0.5357793838) (978,0.5299148107) (982,0.529550087) (986,0.5252283585) (990,0.5288201557) (994,0.5245204328) (998,0.5239427886)};
\addplot[pubNavy, opacity=0.20, line width=0.35pt] coordinates {(10,1.079780923) (14,1.073193284) (18,1.081616345) (22,1.081035341) (26,1.081901554) (30,1.086852312) (34,1.088543418) (38,1.087536779) (42,1.083391575) (46,1.023363685) (50,0.9738132809) (54,0.9166488773) (58,0.8608165246) (62,0.8019504941) (66,0.744131312) (70,0.6826907628) (74,0.6222986032) (78,0.5554680433) (82,0.4884561513) (86,0.4271292843) (90,0.3668019994) (94,0.3360789313) (98,0.328368163) (102,0.3270700824) (106,0.3280218951) (110,0.3247961398) (114,0.3268041464) (118,0.3255065873) (122,0.3302469865) (126,0.3325727935) (130,0.3426746026) (134,0.34555084) (138,0.3478320046) (142,0.3493597959) (146,0.3572833618) (150,0.3619409335) (154,0.3621371897) (158,0.3719510504) (162,0.3783683847) (166,0.3907729267) (170,0.4020812767) (174,0.415917754) (178,0.4219220494) (182,0.441673399) (186,0.4563446836) (190,0.4739864348) (194,0.4982680373) (198,0.5123228473) (202,0.5287948584) (206,0.5383346006) (210,0.5447376551) (214,0.545296737) (218,0.5540876026) (222,0.5558895128) (226,0.5677509501) (230,0.5591164839) (234,0.5517785858) (238,0.5489928108) (242,0.5565728907) (246,0.5635234272) (250,0.56322171) (254,0.5513593722) (258,0.55559461) (262,0.5705573551) (266,0.5768378975) (270,0.5789401544) (274,0.5906858987) (278,0.5850359866) (282,0.5918111408) (286,0.6091977059) (290,0.6190323631) (294,0.5929909946) (298,0.5857416096) (302,0.589326197) (306,0.5783919317) (310,0.5681998058) (314,0.5500231745) (318,0.5420240436) (322,0.5356406768) (326,0.5255668051) (330,0.5305818969) (334,0.5305643948) (338,0.5198842371) (342,0.5107257953) (346,0.5112731727) (350,0.5222042499) (354,0.5185553473) (358,0.537181674) (362,0.5534434777) (366,0.5561959733) (370,0.5542865378) (374,0.5587699441) (378,0.5564277902) (382,0.5479109663) (386,0.5517165565) (390,0.5579391117) (394,0.5702596348) (398,0.5625621393) (402,0.5690079994) (406,0.5560828542) (410,0.5548438144) (414,0.5596532587) (418,0.5624702675) (422,0.5572710937) (426,0.5573812229) (430,0.5562396005) (434,0.554276101) (438,0.558637084) (442,0.5600109912) (446,0.5729763153) (450,0.55859584) (454,0.5664411193) (458,0.5826734601) (462,0.5808912684) (466,0.5915722963) (470,0.6008107028) (474,0.6142436202) (478,0.6358589286) (482,0.6496059268) (486,0.639297128) (490,0.6327552319) (494,0.6207532258) (498,0.6408500803) (502,0.6415974659) (506,0.6344671225) (510,0.6377731583) (514,0.6415671252) (518,0.6443256136) (522,0.6444914632) (526,0.634215223) (530,0.6264887613) (534,0.6153868187) (538,0.6229033419) (542,0.6326713614) (546,0.6317198458) (550,0.6232642053) (554,0.6357548588) (558,0.6276894924) (562,0.6216163136) (566,0.6096637241) (570,0.6101947429) (574,0.6121007517) (578,0.6061725195) (582,0.6145075079) (586,0.6159593621) (590,0.617782362) (594,0.6139718167) (598,0.6035223522) (602,0.6102213701) (606,0.6146942866) (610,0.6074423758) (614,0.6053355414) (618,0.6077556681) (622,0.615560282) (626,0.6052959614) (630,0.6019151827) (634,0.6006674016) (638,0.5993157425) (642,0.6021621202) (646,0.6096738601) (650,0.6032836427) (654,0.6030429045) (658,0.6119251934) (662,0.6162535288) (666,0.6284453644) (670,0.6195150888) (674,0.6289886845) (678,0.6229182737) (682,0.6339831013) (686,0.6248431804) (690,0.6229369744) (694,0.6056201768) (698,0.6161687888) (702,0.5977970362) (706,0.5839209232) (710,0.5749453643) (714,0.5542789625) (718,0.5421629367) (722,0.532049522) (726,0.5251380047) (730,0.5177651876) (734,0.5065334762) (738,0.5092492644) (742,0.5244005845) (746,0.5171243218) (750,0.513777319) (754,0.5156843819) (758,0.520359903) (762,0.5203342947) (766,0.5491027022) (770,0.5582252023) (774,0.5617296886) (778,0.5627046646) (782,0.5702098705) (786,0.5833857777) (790,0.5755110643) (794,0.57292088) (798,0.575124988) (802,0.5757551979) (806,0.5728370885) (810,0.5650014942) (814,0.5625192278) (818,0.5564924921) (822,0.5593670371) (826,0.5621443795) (830,0.5471684164) (834,0.5302100172) (838,0.5534003268) (842,0.5403353061) (846,0.5397202175) (850,0.5334903327) (854,0.5152684994) (858,0.5015756476) (862,0.5082525698) (866,0.51020123) (870,0.4996381139) (874,0.501066103) (878,0.4911750857) (882,0.4874417648) (886,0.4817077825) (890,0.4843311611) (894,0.4859349925) (898,0.4929990827) (902,0.5062647936) (906,0.5185662509) (910,0.5344595147) (914,0.5262811372) (918,0.5295403023) (922,0.5234715698) (926,0.5145768608) (930,0.5244840165) (934,0.5507227482) (938,0.5417014455) (942,0.5286985187) (946,0.5263939345) (950,0.5265627791) (954,0.5217730791) (958,0.5114657784) (962,0.5048154028) (966,0.5014411489) (970,0.5105743823) (974,0.5159345175) (978,0.529505094) (982,0.5128388743) (986,0.5045023155) (990,0.5115982585) (994,0.5099305719) (998,0.5100033042)};
\addplot[pubNavy, opacity=0.20, line width=0.35pt] coordinates {(10,1.093432383) (14,1.074446942) (18,1.072436567) (22,1.072930824) (26,1.07399564) (30,1.07826009) (34,1.079317186) (38,1.071989586) (42,1.07119722) (46,1.005715685) (50,0.9470405994) (54,0.8769651676) (58,0.8095212423) (62,0.7516875273) (66,0.693719458) (70,0.6316549282) (74,0.5750382215) (78,0.5185227588) (82,0.4626458909) (86,0.4097493733) (90,0.3561782843) (94,0.3257112699) (98,0.3262381206) (102,0.3263669664) (106,0.3312403609) (110,0.334535386) (114,0.3389573527) (118,0.3361932497) (122,0.3371823) (126,0.3343108646) (130,0.3339126713) (134,0.3317962753) (138,0.3292618186) (142,0.32884066) (146,0.3394064733) (150,0.3564998633) (154,0.3718121417) (158,0.3937952443) (162,0.4079251015) (166,0.4150966243) (170,0.4277848552) (174,0.4361891649) (178,0.4429369567) (182,0.4513514257) (186,0.4627876856) (190,0.4766048004) (194,0.4841829832) (198,0.4984033429) (202,0.4957333057) (206,0.4890901206) (210,0.4883567523) (214,0.4951796777) (218,0.4922751726) (222,0.5027259396) (226,0.5180650749) (230,0.5194681161) (234,0.5173551071) (238,0.5141331905) (242,0.5264619504) (246,0.5349901495) (250,0.5359676222) (254,0.5432944756) (258,0.5501173961) (262,0.5521227571) (266,0.560825242) (270,0.5528278616) (274,0.5616207498) (278,0.5632623417) (282,0.5730520665) (286,0.580369889) (290,0.5819543224) (294,0.5593548667) (298,0.5507316946) (302,0.5521102881) (306,0.5489705705) (310,0.5505541498) (314,0.559341038) (318,0.5679602737) (322,0.5760513303) (326,0.5671447515) (330,0.5758996141) (334,0.5649827386) (338,0.5589125234) (342,0.5664686018) (346,0.5736619021) (350,0.5749898711) (354,0.5726950304) (358,0.5763744484) (362,0.5717882066) (366,0.5656987094) (370,0.5600567623) (374,0.5675437222) (378,0.5795920251) (382,0.5882759168) (386,0.5954608518) (390,0.6097471948) (394,0.6272811611) (398,0.6352355654) (402,0.6521275391) (406,0.6482395973) (410,0.6540084866) (414,0.6550551764) (418,0.6825878071) (422,0.6933464246) (426,0.6885023529) (430,0.6910758021) (434,0.6881182029) (438,0.7030588796) (442,0.7178892322) (446,0.7229652051) (450,0.7104922436) (454,0.7316464412) (458,0.744000977) (462,0.7585066089) (466,0.7606250864) (470,0.7632253608) (474,0.773874959) (478,0.786224865) (482,0.7959029727) (486,0.7953556875) (490,0.7747430983) (494,0.7767683957) (498,0.7770122436) (502,0.7803555053) (506,0.7714978739) (510,0.7649698702) (514,0.7556534241) (518,0.7491127928) (522,0.7448038927) (526,0.7310864974) (530,0.7231659837) (534,0.7053433009) (538,0.7121964602) (542,0.7100917355) (546,0.7164953668) (550,0.7110626068) (554,0.7296303277) (558,0.7162546102) (562,0.7057196363) (566,0.7203509536) (570,0.7229648432) (574,0.7151510541) (578,0.7111692656) (582,0.7261159407) (586,0.7150377509) (590,0.7070640279) (594,0.7067366362) (598,0.7121346142) (602,0.7120514853) (606,0.701058998) (610,0.6906774333) (614,0.6924562491) (618,0.68088806) (622,0.688366201) (626,0.6975619291) (630,0.6822459032) (634,0.6690400974) (638,0.6706740021) (642,0.6708272143) (646,0.6681431216) (650,0.6440366846) (654,0.6488699506) (658,0.6548109613) (662,0.6684749692) (666,0.6606664779) (670,0.6622006759) (674,0.6660187454) (678,0.6563498302) (682,0.6651146058) (686,0.6552277001) (690,0.6677226477) (694,0.663890353) (698,0.6723782311) (702,0.6605386143) (706,0.6638668992) (710,0.666377137) (714,0.6521218916) (718,0.6720707281) (722,0.6501255264) (726,0.657890629) (730,0.6614466262) (734,0.6790113891) (738,0.6792417625) (742,0.6849370883) (746,0.6852050158) (750,0.7007161953) (754,0.692524645) (758,0.6782091411) (762,0.673860965) (766,0.6839699022) (770,0.6943649972) (774,0.6969945867) (778,0.6803779398) (782,0.6829360536) (786,0.6886937398) (790,0.6856138364) (794,0.6862974633) (798,0.6747965778) (802,0.6796280759) (806,0.6715301157) (810,0.6641609832) (814,0.6700624385) (818,0.6604667141) (822,0.6583258395) (826,0.6620295819) (830,0.6542993334) (834,0.6550437535) (838,0.652145862) (842,0.6521260744) (846,0.6491627192) (850,0.6374675959) (854,0.6288811922) (858,0.6387674632) (862,0.6305385798) (866,0.6372747505) (870,0.6279100151) (874,0.6198042626) (878,0.6106702974) (882,0.6003297954) (886,0.5903861005) (890,0.6042888458) (894,0.6110688086) (898,0.601860871) (902,0.6171680205) (906,0.6315670824) (910,0.6314625123) (914,0.6375486652) (918,0.6381021993) (922,0.6535585308) (926,0.6572860444) (930,0.6454435986) (934,0.6658176435) (938,0.654850345) (942,0.6449779872) (946,0.6449880055) (950,0.6375989887) (954,0.6248332098) (958,0.6243058991) (962,0.6111021959) (966,0.6095741711) (970,0.600326744) (974,0.5924948674) (978,0.6087503083) (982,0.6267248717) (986,0.6198191204) (990,0.6228210376) (994,0.6171277487) (998,0.6281152347)};
\addplot[pubNavy, opacity=0.20, line width=0.35pt] coordinates {(10,1.03863969) (14,1.046308152) (18,1.063753748) (22,1.069303921) (26,1.061744361) (30,1.055953471) (34,1.059450382) (38,1.067140082) (42,1.05823885) (46,0.9904845652) (50,0.931071408) (54,0.870043442) (58,0.8046621883) (62,0.7411474873) (66,0.6752512635) (70,0.6084581256) (74,0.5427877034) (78,0.4909561485) (82,0.4330871737) (86,0.3640660561) (90,0.2994351462) (94,0.2653120712) (98,0.2604868608) (102,0.2606206571) (106,0.2621836688) (110,0.2636547297) (114,0.2630932704) (118,0.2729790908) (122,0.2861222678) (126,0.2934985437) (130,0.3029935425) (134,0.3139915916) (138,0.3225760087) (142,0.3414839556) (146,0.3510565033) (150,0.358203521) (154,0.368148055) (158,0.3881734233) (162,0.396432579) (166,0.4051179328) (170,0.3953432) (174,0.397013362) (178,0.3834319363) (182,0.3871795796) (186,0.3955727019) (190,0.397838684) (194,0.4065938583) (198,0.4215363225) (202,0.4269953239) (206,0.4242591026) (210,0.4314554759) (214,0.4291706376) (218,0.4376477188) (222,0.4551800059) (226,0.4733099258) (230,0.4838169818) (234,0.4814338114) (238,0.4830914167) (242,0.4899091048) (246,0.494570865) (250,0.5043102335) (254,0.5135974902) (258,0.5259012236) (262,0.5282151992) (266,0.5406052798) (270,0.5357849452) (274,0.5508572014) (278,0.5534399742) (282,0.5566220617) (286,0.5541813363) (290,0.5606974976) (294,0.555873115) (298,0.5418990386) (302,0.5426468573) (306,0.5336087748) (310,0.5312906546) (314,0.5325611054) (318,0.5337190452) (322,0.5295972758) (326,0.5160061311) (330,0.5316225808) (334,0.5318615387) (338,0.5221271934) (342,0.5293896883) (346,0.5352809408) (350,0.5438985482) (354,0.5417832237) (358,0.536700869) (362,0.529564483) (366,0.5341591066) (370,0.5230130202) (374,0.5340693619) (378,0.5289703133) (382,0.5247270955) (386,0.5158094406) (390,0.5102792831) (394,0.5132734757) (398,0.5071468284) (402,0.5095207093) (406,0.5128040169) (410,0.5094804167) (414,0.504615898) (418,0.5062936648) (422,0.5042778667) (426,0.4923741974) (430,0.4903129395) (434,0.4956105913) (438,0.4935901201) (442,0.4963802767) (446,0.5009190211) (450,0.4943605955) (454,0.4985798151) (458,0.4995392477) (462,0.5014791676) (466,0.4997710035) (470,0.4989013545) (474,0.4948498879) (478,0.4914026273) (482,0.488135598) (486,0.482352116) (490,0.4657714202) (494,0.4543451198) (498,0.4523311746) (502,0.4397316099) (506,0.4366201113) (510,0.4392882719) (514,0.4277815129) (518,0.4239386885) (522,0.428484676) (526,0.4225142868) (530,0.4282720235) (534,0.4177430289) (538,0.4215010931) (542,0.4304954175) (546,0.4312678399) (550,0.4322869095) (554,0.4401783119) (558,0.4370732944) (562,0.4337146817) (566,0.4369332709) (570,0.4365419398) (574,0.4343413019) (578,0.4309527942) (582,0.4354238457) (586,0.4435762939) (590,0.4372608553) (594,0.4444190244) (598,0.4431010629) (602,0.444223699) (606,0.4351986007) (610,0.4311878205) (614,0.4343815141) (618,0.4343974253) (622,0.4386840973) (626,0.4418580497) (630,0.4389048592) (634,0.4369442308) (638,0.4344793836) (642,0.4300296987) (646,0.4315117595) (650,0.4280335222) (654,0.4310293763) (658,0.4359043872) (662,0.4443282523) (666,0.4528534321) (670,0.4461227086) (674,0.4544232636) (678,0.4578944803) (682,0.4701733861) (686,0.4648265871) (690,0.4740393462) (694,0.477592514) (698,0.4842452855) (702,0.4897902667) (706,0.4939370959) (710,0.5010004472) (714,0.4959701612) (718,0.5057185335) (722,0.5066355285) (726,0.5031410168) (730,0.5012227611) (734,0.5019211704) (738,0.5105813691) (742,0.5104968741) (746,0.5101901703) (750,0.5080191901) (754,0.5094945336) (758,0.5047204254) (762,0.4960629414) (766,0.5046396856) (770,0.498772616) (774,0.499733091) (778,0.5031488786) (782,0.5077071456) (786,0.5076178297) (790,0.4997068222) (794,0.4969672851) (798,0.4941836843) (802,0.4902105446) (806,0.4876213969) (810,0.4893855479) (814,0.4869486301) (818,0.489169985) (822,0.4938452989) (826,0.5002654054) (830,0.4895442897) (834,0.4856323038) (838,0.4982645721) (842,0.4984649568) (846,0.5065165718) (850,0.5053473901) (854,0.4987999438) (858,0.5050883151) (862,0.505371869) (866,0.5133355862) (870,0.5130701077) (874,0.5090523356) (878,0.5073293542) (882,0.5033975735) (886,0.4992905852) (890,0.5014660365) (894,0.5068026817) (898,0.5047070441) (902,0.5184834733) (906,0.520581764) (910,0.5180693374) (914,0.5172748976) (918,0.5185470544) (922,0.514268166) (926,0.5125612317) (930,0.5160592612) (934,0.5385890182) (938,0.5378789196) (942,0.538499827) (946,0.5351103075) (950,0.5380271689) (954,0.5345040692) (958,0.5417037285) (962,0.5429840653) (966,0.5431303758) (970,0.5447326473) (974,0.5402517952) (978,0.537217112) (982,0.5300820707) (986,0.5239369331) (990,0.5247877376) (994,0.5161497806) (998,0.5113921611)};
\addplot[pubNavy, opacity=0.20, line width=0.35pt] coordinates {(10,1.099182344) (14,1.085752468) (18,1.080643217) (22,1.079385027) (26,1.082819801) (30,1.076912507) (34,1.081212057) (38,1.079392328) (42,1.072297367) (46,1.010005782) (50,0.956645265) (54,0.8976464518) (58,0.8358686714) (62,0.7781818072) (66,0.716479873) (70,0.6573603092) (74,0.5964093293) (78,0.5339846215) (82,0.4701303347) (86,0.4088174503) (90,0.3421223191) (94,0.3133563926) (98,0.3088532914) (102,0.3054565685) (106,0.3009276479) (110,0.2914108515) (114,0.2860953731) (118,0.2816008333) (122,0.2715733626) (126,0.2642772594) (130,0.2536464588) (134,0.2482736409) (138,0.2480120082) (142,0.2513390001) (146,0.2496158661) (150,0.2507633103) (154,0.2572625312) (158,0.2705526749) (162,0.2877142627) (166,0.3011096276) (170,0.3142296286) (174,0.3358128281) (178,0.346419175) (182,0.3695089519) (186,0.3817083871) (190,0.4006281035) (194,0.4146839758) (198,0.4281133574) (202,0.4427934989) (206,0.4440050819) (210,0.4529464269) (214,0.4472549415) (218,0.4402513854) (222,0.4483543091) (226,0.4562570573) (230,0.4609873247) (234,0.4559616306) (238,0.4612726995) (242,0.4624127478) (246,0.4652378233) (250,0.4756909158) (254,0.4716418506) (258,0.4683534141) (262,0.4672458685) (266,0.4780966965) (270,0.4698201028) (274,0.4745060351) (278,0.4753031246) (282,0.4861139602) (286,0.4890694833) (290,0.5069819601) (294,0.4840832035) (298,0.4841376414) (302,0.4887053421) (306,0.4926565511) (310,0.4898646434) (314,0.4933537866) (318,0.4950784366) (322,0.5101110552) (326,0.5063305424) (330,0.5185557694) (334,0.527392405) (338,0.5188740106) (342,0.5279429168) (346,0.5381646579) (350,0.5456577817) (354,0.5340520082) (358,0.5406886565) (362,0.5503238042) (366,0.563849498) (370,0.5650734221) (374,0.5860844284) (378,0.5879799696) (382,0.5916021256) (386,0.5848877835) (390,0.5817500703) (394,0.580760542) (398,0.5701682212) (402,0.5804281521) (406,0.5693000551) (410,0.5686704355) (414,0.5687322229) (418,0.5621372045) (422,0.550260282) (426,0.5382079806) (430,0.5267672574) (434,0.5302568936) (438,0.5365631493) (442,0.5463016879) (446,0.5582182513) (450,0.5542682254) (454,0.5528073527) (458,0.5626050097) (462,0.5671355939) (466,0.5628021614) (470,0.5657562152) (474,0.5728094785) (478,0.5767999353) (482,0.5751715458) (486,0.5674043713) (490,0.5590607853) (494,0.549702054) (498,0.5596262925) (502,0.5636931151) (506,0.5708466992) (510,0.5616316574) (514,0.5697787333) (518,0.5789118865) (522,0.5807822841) (526,0.5807513118) (530,0.5746771772) (534,0.5807508789) (538,0.5835471077) (542,0.581897753) (546,0.5770652873) (550,0.567946654) (554,0.5662611828) (558,0.5778578302) (562,0.5856519361) (566,0.5826949701) (570,0.5761647424) (574,0.5717884221) (578,0.5781456144) (582,0.5768489806) (586,0.5822673699) (590,0.5872357748) (594,0.6068483778) (598,0.6117904918) (602,0.6120858867) (606,0.6038446348) (610,0.5888718707) (614,0.5833855825) (618,0.5863182512) (622,0.6064240903) (626,0.6164336255) (630,0.610457582) (634,0.6054926366) (638,0.6049874024) (642,0.6001787839) (646,0.59902507) (650,0.5972872828) (654,0.6124346823) (658,0.6380091178) (662,0.654857015) (666,0.650902238) (670,0.6407505001) (674,0.6343260215) (678,0.6160773521) (682,0.6319938787) (686,0.6295876055) (690,0.6277018319) (694,0.6420970002) (698,0.6421102616) (702,0.6441788491) (706,0.6366074801) (710,0.6261407247) (714,0.60311765) (718,0.611109981) (722,0.6080225593) (726,0.6209314522) (730,0.6298351066) (734,0.6521658812) (738,0.6543176739) (742,0.6516534677) (746,0.6594422155) (750,0.6703154592) (754,0.6663936267) (758,0.6610662252) (762,0.6747155418) (766,0.7051116394) (770,0.7162740601) (774,0.7240652968) (778,0.7111388363) (782,0.6887346594) (786,0.6818418319) (790,0.665495262) (794,0.6635119843) (798,0.6582615581) (802,0.6577733677) (806,0.6420465889) (810,0.6240341897) (814,0.61815256) (818,0.6026978692) (822,0.613733621) (826,0.6081431898) (830,0.6220702038) (834,0.6075182502) (838,0.6320907017) (842,0.6247191875) (846,0.617697075) (850,0.6003284856) (854,0.5872500278) (858,0.5933298987) (862,0.5987045583) (866,0.5928395307) (870,0.5900403022) (874,0.5875075946) (878,0.5860119581) (882,0.5780837689) (886,0.564246313) (890,0.560136219) (894,0.576988682) (898,0.5737841251) (902,0.5881378553) (906,0.6106620088) (910,0.6215821065) (914,0.6284011464) (918,0.6272301645) (922,0.6142551486) (926,0.605165564) (930,0.6014114622) (934,0.6306669993) (938,0.634136035) (942,0.6256073253) (946,0.635430669) (950,0.6316166535) (954,0.6094726964) (958,0.5919246907) (962,0.572950027) (966,0.582117961) (970,0.5844019445) (974,0.5854921262) (978,0.6096465953) (982,0.6094337905) (986,0.5905193458) (990,0.5839200299) (994,0.5697095158) (998,0.580795452)};
\addplot[pubNavy, opacity=0.20, line width=0.35pt] coordinates {(10,1.07771205) (14,1.071128207) (18,1.0745763) (22,1.068412753) (26,1.067232234) (30,1.067055156) (34,1.073791375) (38,1.065878843) (42,1.056452834) (46,0.9779059856) (50,0.9133487759) (54,0.8396750749) (58,0.7657521898) (62,0.6890808225) (66,0.6165928496) (70,0.5494183274) (74,0.4830458058) (78,0.416082173) (82,0.3460666505) (86,0.2792154894) (90,0.2196115465) (94,0.2008313815) (98,0.2113946268) (102,0.223674723) (106,0.23953479) (110,0.2599638793) (114,0.2711491977) (118,0.2802208412) (122,0.2906058844) (126,0.297745684) (130,0.3067337754) (134,0.3136483543) (138,0.3108893372) (142,0.3121950106) (146,0.3166584852) (150,0.3227250409) (154,0.325174775) (158,0.3320707836) (162,0.3428504988) (166,0.3551928957) (170,0.358252626) (174,0.3728056373) (178,0.3797365026) (182,0.3947624597) (186,0.4189560872) (190,0.4427413775) (194,0.4624233328) (198,0.4809607651) (202,0.4893248955) (206,0.5041059898) (210,0.5141132109) (214,0.511436119) (218,0.5136606906) (222,0.5250972049) (226,0.5342011094) (230,0.5364936295) (234,0.5280998474) (238,0.5147651877) (242,0.5109721217) (246,0.5108555853) (250,0.5169786915) (254,0.5022833806) (258,0.5034377644) (262,0.5013762958) (266,0.50178917) (270,0.4917067873) (274,0.4907005919) (278,0.4814313632) (282,0.4853086301) (286,0.4913253241) (290,0.5091180295) (294,0.4846226365) (298,0.4772714397) (302,0.4719167156) (306,0.4720443965) (310,0.4641296406) (314,0.4676815651) (318,0.4725488466) (322,0.4794355738) (326,0.4796647242) (330,0.4874244893) (334,0.4898216057) (338,0.4686542437) (342,0.4754800285) (346,0.4795709144) (350,0.4820969567) (354,0.4740696754) (358,0.478014137) (362,0.4776082067) (366,0.4728639399) (370,0.4632843663) (374,0.4685084106) (378,0.469269025) (382,0.4596278351) (386,0.455310401) (390,0.4519761321) (394,0.4561784676) (398,0.4494859123) (402,0.4616671218) (406,0.452378766) (410,0.4555248117) (414,0.4569773632) (418,0.4643334446) (422,0.4506164202) (426,0.4410984545) (430,0.4441553283) (434,0.4520403635) (438,0.4639536541) (442,0.4700932358) (446,0.4790184189) (450,0.4776772355) (454,0.4835198438) (458,0.495341793) (462,0.4961692434) (466,0.4980298668) (470,0.5030203997) (474,0.5175879097) (478,0.52735554) (482,0.5349732522) (486,0.5373654995) (490,0.5275434392) (494,0.5259594241) (498,0.5286686652) (502,0.5285990184) (506,0.5261906725) (510,0.5253422661) (514,0.5256718277) (518,0.5278847601) (522,0.5358165243) (526,0.5312254383) (530,0.5353463653) (534,0.5285689021) (538,0.5336286157) (542,0.5411403771) (546,0.5359879183) (550,0.5405934757) (554,0.5452882191) (558,0.5521742919) (562,0.5588840667) (566,0.5673570338) (570,0.5650800162) (574,0.5694669153) (578,0.5621396941) (582,0.5641902614) (586,0.5630375521) (590,0.5613643348) (594,0.5679882797) (598,0.5861924024) (602,0.5845395528) (606,0.5787553937) (610,0.5742899879) (614,0.574091351) (618,0.5819727781) (622,0.5931226541) (626,0.5983049935) (630,0.5998301325) (634,0.602567743) (638,0.6023509366) (642,0.6068333437) (646,0.6057680875) (650,0.5970612744) (654,0.5999299814) (658,0.6150757455) (662,0.6257991561) (666,0.6227177109) (670,0.6132244722) (674,0.6063869246) (678,0.5929140488) (682,0.5963654225) (686,0.5875155331) (690,0.5860651421) (694,0.5843491137) (698,0.5876829631) (702,0.5800284541) (706,0.5808135893) (710,0.5843766892) (714,0.5719237058) (718,0.5822716352) (722,0.5843239302) (726,0.5996551042) (730,0.6076432882) (734,0.618731462) (738,0.6224706683) (742,0.6321747394) (746,0.6356840642) (750,0.6486088799) (754,0.6461536115) (758,0.6356971303) (762,0.6459705789) (766,0.6738652891) (770,0.6868865473) (774,0.6901417698) (778,0.6854539831) (782,0.6830804833) (786,0.6865364251) (790,0.6847500113) (794,0.6805169688) (798,0.6811657761) (802,0.6813861932) (806,0.6636846374) (810,0.6645232554) (814,0.6554922455) (818,0.6366922021) (822,0.6491302168) (826,0.6438999439) (830,0.648003056) (834,0.6359731525) (838,0.6528563108) (842,0.6455763716) (846,0.6447238643) (850,0.6303990209) (854,0.6189689532) (858,0.6127406203) (862,0.6056784837) (866,0.5985455867) (870,0.5885583016) (874,0.5671886901) (878,0.5606724694) (882,0.546051675) (886,0.5387013798) (890,0.519396332) (894,0.523542157) (898,0.5229903055) (902,0.5301500753) (906,0.5430428583) (910,0.545437347) (914,0.5439717536) (918,0.5443130536) (922,0.5432793104) (926,0.5423256213) (930,0.5481342771) (934,0.5709024689) (938,0.571972886) (942,0.5732816559) (946,0.5757130811) (950,0.5683642007) (954,0.5582883597) (958,0.5525698834) (962,0.5467386198) (966,0.5630366756) (970,0.5786645974) (974,0.5738119556) (978,0.5918175952) (982,0.5962815155) (986,0.5943808438) (990,0.5959597018) (994,0.5920221573) (998,0.6018902124)};
\addplot[pubNavy, opacity=0.20, line width=0.35pt] coordinates {(10,1.086755537) (14,1.086181764) (18,1.08947761) (22,1.085286538) (26,1.080338181) (30,1.086760277) (34,1.083596111) (38,1.083439803) (42,1.073971751) (46,1.005408625) (50,0.9461236045) (54,0.8887242774) (58,0.8250437119) (62,0.7642689987) (66,0.7023418308) (70,0.6430642326) (74,0.588379752) (78,0.5266369207) (82,0.4695361914) (86,0.4217209542) (90,0.3669241094) (94,0.3513906195) (98,0.3573093173) (102,0.365679688) (106,0.3663021118) (110,0.3660816888) (114,0.3659399693) (118,0.3685634859) (122,0.369321733) (126,0.3653470466) (130,0.3653135138) (134,0.3632966164) (138,0.3557156661) (142,0.3482773487) (146,0.3499618641) (150,0.3452931136) (154,0.3495533382) (158,0.3594853353) (162,0.3694406418) (166,0.3855648092) (170,0.3952321534) (174,0.4094289154) (178,0.4189159975) (182,0.4260892542) (186,0.4498439767) (190,0.4677610511) (194,0.4873712856) (198,0.5104044777) (202,0.5275882985) (206,0.5445872467) (210,0.5607780789) (214,0.5687792246) (218,0.5762253636) (222,0.5851758328) (226,0.591261181) (230,0.5925387012) (234,0.5927367295) (238,0.5886517397) (242,0.5895583098) (246,0.5851729747) (250,0.5793088886) (254,0.5699979421) (258,0.5652737886) (262,0.5604688305) (266,0.5563069317) (270,0.5396453034) (274,0.5423781882) (278,0.5476102171) (282,0.5543544335) (286,0.5513201772) (290,0.5545868044) (294,0.5422429568) (298,0.5409336886) (302,0.5439571583) (306,0.5494301478) (310,0.5444197947) (314,0.5569972357) (318,0.5620745453) (322,0.5674681817) (326,0.557522512) (330,0.5534889453) (334,0.5536882937) (338,0.5511148183) (342,0.5447036977) (346,0.5512161189) (350,0.5438347224) (354,0.5294870839) (358,0.5292580103) (362,0.5206703825) (366,0.524108267) (370,0.5110865484) (374,0.5202720984) (378,0.5146483492) (382,0.5202301525) (386,0.516128704) (390,0.52537776) (394,0.5375542283) (398,0.5422336422) (402,0.5515745224) (406,0.5523168905) (410,0.5485694692) (414,0.5456850763) (418,0.5414556239) (422,0.5381144832) (426,0.5370479556) (430,0.5431846159) (434,0.5420703644) (438,0.5302455821) (442,0.5328530194) (446,0.5259717371) (450,0.5259339854) (454,0.5267757277) (458,0.5285941644) (462,0.5269780182) (466,0.5270010645) (470,0.5350973473) (474,0.5365176967) (478,0.550063768) (482,0.55467523) (486,0.5576023804) (490,0.5505712019) (494,0.55696266) (498,0.5692113581) (502,0.5655006185) (506,0.5726815444) (510,0.5882095968) (514,0.5869979667) (518,0.5983250077) (522,0.5918295049) (526,0.5834251693) (530,0.5753017262) (534,0.5684736163) (538,0.5749641083) (542,0.5807695487) (546,0.5770815588) (550,0.5834184928) (554,0.5870223017) (558,0.5789765322) (562,0.575400272) (566,0.5718891198) (570,0.5731889066) (574,0.5743208286) (578,0.5854350722) (582,0.5953893284) (586,0.5974586007) (590,0.5937628578) (594,0.6014484734) (598,0.6007454451) (602,0.5995609056) (606,0.6031855237) (610,0.5941591216) (614,0.5974220308) (618,0.6067708537) (622,0.6179451031) (626,0.6210063843) (630,0.6064880514) (634,0.5984084608) (638,0.596217187) (642,0.6046360162) (646,0.6022090614) (650,0.588728178) (654,0.5858192487) (658,0.5979735647) (662,0.6051664857) (666,0.6032177449) (670,0.5870588213) (674,0.5945324603) (678,0.5937360708) (682,0.6002078412) (686,0.585884633) (690,0.5873895568) (694,0.5828299879) (698,0.5917229641) (702,0.5889955026) (706,0.5918976175) (710,0.6014115584) (714,0.5940489537) (718,0.6076410505) (722,0.6092817738) (726,0.6106196609) (730,0.6226138665) (734,0.6383116991) (738,0.6554008063) (742,0.6699362607) (746,0.6723381539) (750,0.6786053855) (754,0.6725812543) (758,0.6676703054) (762,0.6622677568) (766,0.6726637757) (770,0.6675817166) (774,0.6639609319) (778,0.656632483) (782,0.6563688988) (786,0.6515259939) (790,0.646000818) (794,0.6402062105) (798,0.6320400106) (802,0.6345866327) (806,0.6315632019) (810,0.6287216995) (814,0.6167744391) (818,0.6104125023) (822,0.6144186429) (826,0.6231319339) (830,0.6045535953) (834,0.5945960515) (838,0.6110907539) (842,0.599636126) (846,0.599680973) (850,0.5959905427) (854,0.5809027481) (858,0.5839580619) (862,0.5882794333) (866,0.5960025755) (870,0.590325001) (874,0.5819777687) (878,0.5844627231) (882,0.5876902533) (886,0.5765043455) (890,0.5683399612) (894,0.577290413) (898,0.5716231153) (902,0.5828605598) (906,0.59609727) (910,0.6028416513) (914,0.5909558754) (918,0.5929265914) (922,0.6018218291) (926,0.5910933028) (930,0.589747336) (934,0.60915361) (938,0.6055910555) (942,0.5975935595) (946,0.5989189222) (950,0.6060425015) (954,0.5986398841) (958,0.5869123112) (962,0.5956761862) (966,0.5924067959) (970,0.5903528865) (974,0.5853088252) (978,0.5966105351) (982,0.5998310197) (986,0.5907281287) (990,0.6031824332) (994,0.5897550683) (998,0.5848637901)};
\addplot[pubNavy, ultra thick] coordinates {(10,1.077146817) (14,1.074524454) (18,1.074999965) (22,1.076130027) (26,1.074199571) (30,1.073819179) (34,1.074983191) (38,1.071579641) (42,1.067765279) (46,0.9964004721) (50,0.9404536341) (54,0.8740621453) (58,0.8066513473) (62,0.7406051331) (66,0.6754729765) (70,0.6115613273) (74,0.5442545543) (78,0.4864195043) (82,0.4214853428) (86,0.3576757817) (90,0.292949176) (94,0.2660414792) (98,0.2633563931) (102,0.2652592282) (106,0.2687335037) (110,0.2757076566) (114,0.2853914907) (118,0.2847442757) (122,0.2912936365) (126,0.2972553102) (130,0.3099325913) (134,0.3203921007) (138,0.3259189137) (142,0.3351623078) (146,0.3349795297) (150,0.3419231485) (154,0.3489588787) (158,0.3657181929) (162,0.3739045132) (166,0.388168868) (170,0.3952876767) (174,0.4053201744) (178,0.4038083545) (182,0.4098064554) (186,0.4207510999) (190,0.4364899336) (194,0.4493380322) (198,0.46157221) (202,0.4655444459) (206,0.4616591775) (210,0.4633476391) (214,0.4588037974) (218,0.4541469585) (222,0.4636520091) (226,0.476696338) (230,0.4822006211) (234,0.479696003) (238,0.4782786341) (242,0.480618447) (246,0.4865274389) (250,0.4925192519) (254,0.4844912185) (258,0.4827909498) (262,0.481336949) (266,0.4819492995) (270,0.4798338991) (274,0.4855859505) (278,0.4847143167) (282,0.4928812969) (286,0.4970281154) (290,0.5074830419) (294,0.4933991899) (298,0.489161616) (302,0.4918641052) (306,0.4989362999) (310,0.4971398558) (314,0.5024098845) (318,0.5027546745) (322,0.506919138) (326,0.5002244297) (330,0.5007034726) (334,0.5018510408) (338,0.5019784527) (342,0.5046418731) (346,0.5075073024) (350,0.5181794232) (354,0.5205607537) (358,0.5289953247) (362,0.5227431817) (366,0.5243733959) (370,0.5148118334) (374,0.5192005023) (378,0.5182382327) (382,0.522478624) (386,0.5215691851) (390,0.5309363974) (394,0.5416543391) (398,0.5464070048) (402,0.5529106057) (406,0.554090377) (410,0.5509964253) (414,0.5478968423) (418,0.5459780518) (422,0.5436256759) (426,0.5376279681) (430,0.5280818457) (434,0.5319038074) (438,0.5331488652) (442,0.5370557796) (446,0.5350986679) (450,0.5365713617) (454,0.5397915402) (458,0.5443908635) (462,0.5470568061) (466,0.5404022162) (470,0.5454918023) (474,0.5508726861) (478,0.5600195302) (482,0.5670637494) (486,0.5669665337) (490,0.5526451823) (494,0.5524591944) (498,0.5644188253) (502,0.563522834) (506,0.564179556) (510,0.5612947869) (514,0.5639024302) (518,0.5638758008) (522,0.5622945685) (526,0.5623575247) (530,0.5642009324) (534,0.5631429027) (538,0.5790459165) (542,0.5813336508) (546,0.5770734231) (550,0.5756825734) (554,0.5766417423) (558,0.5784171812) (562,0.5805261041) (566,0.5772920449) (570,0.5746768245) (574,0.5730546253) (578,0.5817903433) (582,0.5861191545) (586,0.5898629853) (590,0.5904993163) (594,0.6041484256) (598,0.6021338987) (602,0.6048911379) (606,0.6035150792) (610,0.5915154961) (614,0.5904038067) (618,0.5965445525) (622,0.6085983055) (626,0.6071960464) (630,0.604201617) (634,0.6016175723) (638,0.5977664648) (642,0.6011704521) (646,0.6006170657) (650,0.5928947262) (654,0.5928746151) (658,0.6045234632) (662,0.6084858568) (666,0.6100539668) (670,0.6001416467) (674,0.6014588481) (678,0.5943124102) (682,0.600913517) (686,0.5961277165) (690,0.6024008378) (694,0.5949846453) (698,0.6039458764) (702,0.5933962694) (706,0.5879092703) (710,0.5928941238) (714,0.5829863297) (718,0.5949563428) (722,0.5954727267) (726,0.6039131556) (730,0.6074204221) (734,0.6149723613) (738,0.6180775868) (742,0.6263308943) (746,0.6259591049) (750,0.6287364041) (754,0.6254880013) (758,0.6134942927) (762,0.6089556148) (766,0.6195817034) (770,0.6245579305) (774,0.6201807072) (778,0.6177198115) (782,0.6136978015) (786,0.6112925496) (790,0.6072198121) (794,0.6013657007) (798,0.5954183259) (802,0.5872678063) (806,0.5775481899) (810,0.5733931851) (814,0.5695289426) (818,0.5643883326) (822,0.5757984368) (826,0.5731792873) (830,0.5650708289) (834,0.5620626747) (838,0.5683115001) (842,0.570456709) (846,0.5671027821) (850,0.5707523456) (854,0.5682074647) (858,0.5641887284) (862,0.5637801736) (866,0.5582650885) (870,0.5522940214) (874,0.5476990322) (878,0.5456329589) (882,0.5345620554) (886,0.5310049409) (890,0.5220801643) (894,0.5222540469) (898,0.5230865932) (902,0.5256558358) (906,0.5268084402) (910,0.5342525799) (914,0.5328407393) (918,0.537292147) (922,0.541311368) (926,0.5393485517) (930,0.5422352367) (934,0.5608126085) (938,0.556042494) (942,0.5678681875) (946,0.5683101666) (950,0.5674289525) (954,0.5687186804) (958,0.5659432428) (962,0.5642724746) (966,0.5725773183) (970,0.581533271) (974,0.5795603904) (978,0.5942140652) (982,0.5980562676) (986,0.5906237372) (990,0.5899398658) (994,0.5797322921) (998,0.582829621)};
\nextgroupplot[title={Rolling gross exposure}, xlabel={Pre-training episode}]
\addplot[pubNavy, opacity=0.20, line width=0.35pt] coordinates {(10,1.246422338) (14,1.245952206) (18,1.248023102) (22,1.246011168) (26,1.248335818) (30,1.24787593) (34,1.248961146) (38,1.247924307) (42,1.246642127) (46,1.234712772) (50,1.225370407) (54,1.21645508) (58,1.208716009) (62,1.203157148) (66,1.197402147) (70,1.193995969) (74,1.191541645) (78,1.19080243) (82,1.192123135) (86,1.193501999) (90,1.196066598) (94,1.20873917) (98,1.225890701) (102,1.242843093) (106,1.259947992) (110,1.277286011) (114,1.293935179) (118,1.310125487) (122,1.323715182) (126,1.334601757) (130,1.34227886) (134,1.351057128) (138,1.358615175) (142,1.364478243) (146,1.36650434) (150,1.367219306) (154,1.366254052) (158,1.365229733) (162,1.364120675) (166,1.361193198) (170,1.360347342) (174,1.360170168) (178,1.360376647) (182,1.360594547) (186,1.360988245) (190,1.360822588) (194,1.363075244) (198,1.364991044) (202,1.367976505) (206,1.371431323) (210,1.37444993) (214,1.376352721) (218,1.378833326) (222,1.379786457) (226,1.38225075) (230,1.382546821) (234,1.383624229) (238,1.38532721) (242,1.38620332) (246,1.387127079) (250,1.386937092) (254,1.387922146) (258,1.388313132) (262,1.388339274) (266,1.387939229) (270,1.386793174) (274,1.388325391) (278,1.388082209) (282,1.389056917) (286,1.389757587) (290,1.390232378) (294,1.390822336) (298,1.390885854) (302,1.391418264) (306,1.391496151) (310,1.392348886) (314,1.392698581) (318,1.393854448) (322,1.394063789) (326,1.393770102) (330,1.393622536) (334,1.393850206) (338,1.393845878) (342,1.393454992) (346,1.394638854) (350,1.394653146) (354,1.394765207) (358,1.395152612) (362,1.395022408) (366,1.394593908) (370,1.393921845) (374,1.393964023) (378,1.393857214) (382,1.394329825) (386,1.394096291) (390,1.394390602) (394,1.394326726) (398,1.394514596) (402,1.394257389) (406,1.394402739) (410,1.394390239) (414,1.395310123) (418,1.395505207) (422,1.396321913) (426,1.395728196) (430,1.395942555) (434,1.39626328) (438,1.396267079) (442,1.395942399) (446,1.396162424) (450,1.39598784) (454,1.396144106) (458,1.395611706) (462,1.395381151) (466,1.395625738) (470,1.395413728) (474,1.395607367) (478,1.396534599) (482,1.397020062) (486,1.396690854) (490,1.396341482) (494,1.396517089) (498,1.396746188) (502,1.396728943) (506,1.396893095) (510,1.396705679) (514,1.396383878) (518,1.396081539) (522,1.39585384) (526,1.395447209) (530,1.395458742) (534,1.395428815) (538,1.395789766) (542,1.395263662) (546,1.395304414) (550,1.394995881) (554,1.39511188) (558,1.395602445) (562,1.396001601) (566,1.396154619) (570,1.396094751) (574,1.396309479) (578,1.396462178) (582,1.396088534) (586,1.395805611) (590,1.395877656) (594,1.396166099) (598,1.39635175) (602,1.396630591) (606,1.396796136) (610,1.396163955) (614,1.395965817) (618,1.396413968) (622,1.396472731) (626,1.396391284) (630,1.396653963) (634,1.39733103) (638,1.396804422) (642,1.397598387) (646,1.397864249) (650,1.397170769) (654,1.396741027) (658,1.396435997) (662,1.396795515) (666,1.397379732) (670,1.396515927) (674,1.396750247) (678,1.397058101) (682,1.396738287) (686,1.396266567) (690,1.396776229) (694,1.396497025) (698,1.396372403) (702,1.396718091) (706,1.397138777) (710,1.398039872) (714,1.397803572) (718,1.397889447) (722,1.39799559) (726,1.398006568) (730,1.3979354) (734,1.39808013) (738,1.398108411) (742,1.398525794) (746,1.398266506) (750,1.3979204) (754,1.398214921) (758,1.397179781) (762,1.396758488) (766,1.397311089) (770,1.397113393) (774,1.396996933) (778,1.396927987) (782,1.397169622) (786,1.397012457) (790,1.396396357) (794,1.39621774) (798,1.396281491) (802,1.3966949) (806,1.397028739) (810,1.397032925) (814,1.396578583) (818,1.396855718) (822,1.397170369) (826,1.396886175) (830,1.396562703) (834,1.396489991) (838,1.397198461) (842,1.397233822) (846,1.397277209) (850,1.397191867) (854,1.397197471) (858,1.397064424) (862,1.397792358) (866,1.398017085) (870,1.398291404) (874,1.398496835) (878,1.398801615) (882,1.398749821) (886,1.398835534) (890,1.39918075) (894,1.398687148) (898,1.398762707) (902,1.399214689) (906,1.398738076) (910,1.398616201) (914,1.398375468) (918,1.398344211) (922,1.398076603) (926,1.397581285) (930,1.397416809) (934,1.397568179) (938,1.397303114) (942,1.397600566) (946,1.397612882) (950,1.396926875) (954,1.397509329) (958,1.397692675) (962,1.39701359) (966,1.397516466) (970,1.397692924) (974,1.397801343) (978,1.397939795) (982,1.397978645) (986,1.397946601) (990,1.397831225) (994,1.397883173) (998,1.397818904)};
\addplot[pubNavy, opacity=0.20, line width=0.35pt] coordinates {(10,1.241960485) (14,1.250606296) (18,1.252898575) (22,1.253818106) (26,1.254373379) (30,1.255018738) (34,1.253096446) (38,1.252996531) (42,1.250212508) (46,1.236414288) (50,1.225710536) (54,1.21548759) (58,1.205658429) (62,1.195427306) (66,1.185556647) (70,1.175788818) (74,1.169437001) (78,1.161893001) (82,1.158201076) (86,1.156360634) (90,1.155001008) (94,1.165134265) (98,1.17872049) (102,1.193112701) (106,1.20878233) (110,1.225948234) (114,1.241108803) (118,1.25688365) (122,1.272727077) (126,1.287813243) (130,1.301214691) (134,1.313848278) (138,1.326315815) (142,1.335684531) (146,1.34235566) (150,1.349663443) (154,1.353723168) (158,1.358732743) (162,1.363862959) (166,1.365941121) (170,1.369989654) (174,1.37167296) (178,1.372602282) (182,1.374567787) (186,1.376225773) (190,1.377991811) (194,1.380032754) (198,1.381605948) (202,1.384340526) (206,1.385793114) (210,1.386081346) (214,1.388800681) (218,1.389754376) (222,1.390144136) (226,1.39165476) (230,1.391891047) (234,1.39326945) (238,1.392089754) (242,1.391211029) (246,1.39009015) (250,1.388938965) (254,1.388701461) (258,1.388302289) (262,1.38709823) (266,1.386033667) (270,1.385043353) (274,1.386051231) (278,1.385605897) (282,1.385366067) (286,1.384897064) (290,1.385236331) (294,1.386125871) (298,1.386535532) (302,1.38729172) (306,1.388173857) (310,1.389029708) (314,1.390533041) (318,1.391718413) (322,1.391765289) (326,1.391892229) (330,1.392125685) (334,1.392341499) (338,1.392320965) (342,1.391616318) (346,1.393158851) (350,1.39295097) (354,1.392645725) (358,1.392674182) (362,1.391838326) (366,1.390996462) (370,1.390686262) (374,1.39120566) (378,1.390983016) (382,1.391772686) (386,1.391183885) (390,1.391590389) (394,1.39159436) (398,1.392443256) (402,1.392124274) (406,1.393124333) (410,1.392858979) (414,1.393457808) (418,1.393688561) (422,1.394930131) (426,1.39412314) (430,1.394647483) (434,1.394926209) (438,1.395254136) (442,1.395530361) (446,1.395796645) (450,1.396167686) (454,1.396294127) (458,1.39574645) (462,1.395712612) (466,1.395397802) (470,1.395351629) (474,1.395447862) (478,1.39592978) (482,1.396079686) (486,1.395807825) (490,1.395751561) (494,1.395961591) (498,1.395958113) (502,1.396008309) (506,1.396288518) (510,1.396609891) (514,1.397342039) (518,1.397545925) (522,1.397336593) (526,1.39749302) (530,1.397801627) (534,1.397558078) (538,1.397855503) (542,1.397760554) (546,1.398069301) (550,1.397861426) (554,1.39808418) (558,1.398082384) (562,1.397569117) (566,1.397521773) (570,1.397966653) (574,1.397750503) (578,1.397710234) (582,1.39762116) (586,1.396793843) (590,1.396817974) (594,1.396258186) (598,1.395958996) (602,1.396104926) (606,1.395981903) (610,1.39629811) (614,1.396663725) (618,1.396428377) (622,1.396163693) (626,1.396619685) (630,1.396108783) (634,1.396357543) (638,1.396129035) (642,1.396561788) (646,1.397124267) (650,1.396487147) (654,1.396268193) (658,1.396077163) (662,1.396213982) (666,1.396239518) (670,1.396054024) (674,1.396152717) (678,1.396458904) (682,1.397024893) (686,1.396599249) (690,1.396999592) (694,1.396294814) (698,1.396357578) (702,1.396435555) (706,1.396692752) (710,1.397158516) (714,1.396497561) (718,1.396319018) (722,1.396119713) (726,1.39656647) (730,1.396589207) (734,1.396667083) (738,1.397044144) (742,1.397179758) (746,1.397750168) (750,1.398178159) (754,1.397631495) (758,1.397635884) (762,1.397870285) (766,1.398004499) (770,1.397941016) (774,1.397767346) (778,1.397654738) (782,1.397008012) (786,1.396325068) (790,1.395064991) (794,1.394632652) (798,1.394435389) (802,1.394836325) (806,1.39528781) (810,1.394671154) (814,1.394267635) (818,1.394359765) (822,1.394534583) (826,1.394300694) (830,1.394539633) (834,1.39472188) (838,1.395349162) (842,1.395933695) (846,1.396452247) (850,1.396148188) (854,1.396001675) (858,1.396133964) (862,1.397179864) (866,1.397534673) (870,1.397677218) (874,1.398147967) (878,1.398484211) (882,1.398666269) (886,1.39853004) (890,1.398768225) (894,1.39828097) (898,1.398737109) (902,1.398544134) (906,1.398024605) (910,1.397520121) (914,1.39773784) (918,1.397234537) (922,1.397318317) (926,1.396713371) (930,1.396833257) (934,1.397294941) (938,1.397245214) (942,1.397363427) (946,1.396728148) (950,1.396743889) (954,1.39748307) (958,1.397708973) (962,1.397392755) (966,1.397741422) (970,1.397976458) (974,1.39849801) (978,1.398393076) (982,1.398488353) (986,1.39774017) (990,1.398065007) (994,1.398744945) (998,1.398464425)};
\addplot[pubNavy, opacity=0.20, line width=0.35pt] coordinates {(10,1.239075817) (14,1.24407281) (18,1.245558984) (22,1.243920992) (26,1.245699005) (30,1.245396664) (34,1.245229773) (38,1.244262508) (42,1.243381067) (46,1.233159923) (50,1.226553646) (54,1.222431293) (58,1.221075699) (62,1.221683463) (66,1.226342552) (70,1.234713921) (74,1.245563443) (78,1.258334446) (82,1.270577707) (86,1.28517837) (90,1.296961921) (94,1.31589929) (98,1.337086396) (102,1.356438373) (106,1.37239707) (110,1.385266171) (114,1.394121386) (118,1.399641672) (122,1.401788718) (126,1.401942563) (130,1.398586645) (134,1.396772772) (138,1.394508402) (142,1.390776047) (146,1.38675792) (150,1.382207223) (154,1.378508541) (158,1.37575842) (162,1.373365351) (166,1.371496433) (170,1.370783634) (174,1.369625445) (178,1.369018384) (182,1.369806853) (186,1.370294225) (190,1.372387935) (194,1.375406053) (198,1.37747521) (202,1.380816947) (206,1.382809107) (210,1.384975962) (214,1.386769659) (218,1.387965266) (222,1.388915976) (226,1.390134455) (230,1.390084297) (234,1.391838614) (238,1.391996965) (242,1.39231782) (246,1.392799558) (250,1.392555333) (254,1.39270448) (258,1.393102397) (262,1.393098149) (266,1.393152868) (270,1.392837224) (274,1.393509908) (278,1.394223915) (282,1.394379437) (286,1.394500752) (290,1.395038086) (294,1.395267341) (298,1.39538746) (302,1.39595031) (306,1.396223716) (310,1.396606399) (314,1.396517542) (318,1.396802161) (322,1.396404362) (326,1.396074351) (330,1.395571429) (334,1.39565457) (338,1.395273696) (342,1.394680793) (346,1.395159928) (350,1.39493594) (354,1.394227147) (358,1.394169098) (362,1.394269309) (366,1.394134332) (370,1.394287822) (374,1.394529139) (378,1.394616482) (382,1.39478896) (386,1.394918977) (390,1.394900925) (394,1.394554149) (398,1.394311831) (402,1.394373856) (406,1.393978143) (410,1.39370674) (414,1.39371966) (418,1.393334138) (422,1.393625493) (426,1.392851156) (430,1.392973017) (434,1.392840029) (438,1.393216694) (442,1.393255288) (446,1.393713369) (450,1.394385809) (454,1.395563697) (458,1.395815539) (462,1.39663584) (466,1.396711714) (470,1.397331597) (474,1.397914875) (478,1.398958383) (482,1.399658938) (486,1.399728901) (490,1.399861637) (494,1.400121283) (498,1.400050878) (502,1.399417937) (506,1.399848244) (510,1.399917781) (514,1.399942067) (518,1.400323249) (522,1.399583435) (526,1.399220107) (530,1.398426529) (534,1.398561751) (538,1.398803127) (542,1.398999434) (546,1.398986039) (550,1.399067972) (554,1.399026241) (558,1.399347805) (562,1.399585824) (566,1.399363312) (570,1.399364108) (574,1.399432335) (578,1.399554025) (582,1.39973959) (586,1.399979027) (590,1.399788022) (594,1.399925041) (598,1.39994645) (602,1.400030993) (606,1.399593302) (610,1.399014745) (614,1.399009372) (618,1.398903438) (622,1.398529379) (626,1.398824677) (630,1.399371355) (634,1.399588838) (638,1.399244344) (642,1.399182666) (646,1.399491113) (650,1.398581999) (654,1.398517879) (658,1.398970922) (662,1.399068291) (666,1.39957044) (670,1.399431882) (674,1.399871091) (678,1.399168126) (682,1.399131898) (686,1.39871952) (690,1.398264403) (694,1.398092232) (698,1.398260524) (702,1.397621695) (706,1.397575678) (710,1.397408326) (714,1.397268341) (718,1.397181193) (722,1.397324081) (726,1.397505596) (730,1.397485801) (734,1.397561683) (738,1.397874085) (742,1.398527783) (746,1.398837328) (750,1.398956093) (754,1.39944853) (758,1.398013191) (762,1.398481931) (766,1.398705568) (770,1.399224149) (774,1.399021931) (778,1.398047215) (782,1.39788736) (786,1.397897121) (790,1.397247441) (794,1.397138252) (798,1.396469735) (802,1.396782632) (806,1.398070154) (810,1.39798959) (814,1.397880616) (818,1.397692121) (822,1.39784771) (826,1.398241227) (830,1.398233635) (834,1.398019758) (838,1.39826032) (842,1.398285688) (846,1.398365505) (850,1.398882374) (854,1.39888332) (858,1.398703881) (862,1.398526036) (866,1.398378428) (870,1.397889895) (874,1.39822956) (878,1.398187641) (882,1.398261887) (886,1.398213366) (890,1.398377806) (894,1.397804499) (898,1.397635347) (902,1.39748272) (906,1.397244569) (910,1.397588913) (914,1.397341664) (918,1.397292765) (922,1.398133241) (926,1.398207534) (930,1.398236725) (934,1.398334828) (938,1.397945012) (942,1.397543386) (946,1.398210716) (950,1.397814024) (954,1.397355894) (958,1.397511283) (962,1.397252215) (966,1.39765837) (970,1.397637116) (974,1.397350994) (978,1.397330864) (982,1.397588458) (986,1.397495014) (990,1.398480528) (994,1.39867717) (998,1.398522119)};
\addplot[pubNavy, opacity=0.20, line width=0.35pt] coordinates {(10,1.264833921) (14,1.262620253) (18,1.260993307) (22,1.25895435) (26,1.256555443) (30,1.253634854) (34,1.253919486) (38,1.25439397) (42,1.252801215) (46,1.238992467) (50,1.22791021) (54,1.21622639) (58,1.204724826) (62,1.193147433) (66,1.184469189) (70,1.177791254) (74,1.171331273) (78,1.169125885) (82,1.166728708) (86,1.16485352) (90,1.164890121) (94,1.174526888) (98,1.18944624) (102,1.205538135) (106,1.219487715) (110,1.234119467) (114,1.245399147) (118,1.258585432) (122,1.272774373) (126,1.285368363) (130,1.29514725) (134,1.306421358) (138,1.316910058) (142,1.326128332) (146,1.331408391) (150,1.337671648) (154,1.343256783) (158,1.347919038) (162,1.354195217) (166,1.360594406) (170,1.365780416) (174,1.367057225) (178,1.36886344) (182,1.370819359) (186,1.371748065) (190,1.373673765) (194,1.376184597) (198,1.376977837) (202,1.379288597) (206,1.380976099) (210,1.381467464) (214,1.383059329) (218,1.384804715) (222,1.385162129) (226,1.387223168) (230,1.387781212) (234,1.388446626) (238,1.388141145) (242,1.388176491) (246,1.389061802) (250,1.388419979) (254,1.388845141) (258,1.389180059) (262,1.390394551) (266,1.390444722) (270,1.390044132) (274,1.391268932) (278,1.391602712) (282,1.392271151) (286,1.392509962) (290,1.392767807) (294,1.392122797) (298,1.392387928) (302,1.391551512) (306,1.391706071) (310,1.391587154) (314,1.392444417) (318,1.392947498) (322,1.393087855) (326,1.392508246) (330,1.393063375) (334,1.392494876) (338,1.392009219) (342,1.391993591) (346,1.393127289) (350,1.393964011) (354,1.394214956) (358,1.394461624) (362,1.393256852) (366,1.392377428) (370,1.392072449) (374,1.392379026) (378,1.391882805) (382,1.392988091) (386,1.393028423) (390,1.393335833) (394,1.393593528) (398,1.393915391) (402,1.393559689) (406,1.393590079) (410,1.393994006) (414,1.39542954) (418,1.395625869) (422,1.396935629) (426,1.396402059) (430,1.396307873) (434,1.396305927) (438,1.396007409) (442,1.395809204) (446,1.395503787) (450,1.395197177) (454,1.394748895) (458,1.39427864) (462,1.393933369) (466,1.393960254) (470,1.393450925) (474,1.394265731) (478,1.394721454) (482,1.395619878) (486,1.395527445) (490,1.395703486) (494,1.396121941) (498,1.396512941) (502,1.397373606) (506,1.397465666) (510,1.397910633) (514,1.398076492) (518,1.397870048) (522,1.396672784) (526,1.396493124) (530,1.395966958) (534,1.395222656) (538,1.395981012) (542,1.395315388) (546,1.395377247) (550,1.395152611) (554,1.394823466) (558,1.395368576) (562,1.39511379) (566,1.395311691) (570,1.39585521) (574,1.396205442) (578,1.395587688) (582,1.395772566) (586,1.395628523) (590,1.395765925) (594,1.395464325) (598,1.395610049) (602,1.395645714) (606,1.395475416) (610,1.395110784) (614,1.395276708) (618,1.395366496) (622,1.395473203) (626,1.396051275) (630,1.396265749) (634,1.396631285) (638,1.396534269) (642,1.397367406) (646,1.398091592) (650,1.397833523) (654,1.397657093) (658,1.398243331) (662,1.398047392) (666,1.398520574) (670,1.398017383) (674,1.398381197) (678,1.398295401) (682,1.398535987) (686,1.397792119) (690,1.39779374) (694,1.396680452) (698,1.396294127) (702,1.396171075) (706,1.3962967) (710,1.396707035) (714,1.396379221) (718,1.396245604) (722,1.396115989) (726,1.396036557) (730,1.396034724) (734,1.396537587) (738,1.39689046) (742,1.397486301) (746,1.397738811) (750,1.397914819) (754,1.398072726) (758,1.397620082) (762,1.397759209) (766,1.398185783) (770,1.398501294) (774,1.398687163) (778,1.398776299) (782,1.398525769) (786,1.398468427) (790,1.39778314) (794,1.397713534) (798,1.397442293) (802,1.397988448) (806,1.398002378) (810,1.397784117) (814,1.397970496) (818,1.397710246) (822,1.397611555) (826,1.398070952) (830,1.397677184) (834,1.397206925) (838,1.397620925) (842,1.397680746) (846,1.398039608) (850,1.39813862) (854,1.398238178) (858,1.397920169) (862,1.397794936) (866,1.397961397) (870,1.398359345) (874,1.398312064) (878,1.397847324) (882,1.397537848) (886,1.39793642) (890,1.398336269) (894,1.397786086) (898,1.397284153) (902,1.397497386) (906,1.397785198) (910,1.397995883) (914,1.398108873) (918,1.39781178) (922,1.398184094) (926,1.398609209) (930,1.398394036) (934,1.39903799) (938,1.398923205) (942,1.39862055) (946,1.398526441) (950,1.398025194) (954,1.397884545) (958,1.397728194) (962,1.397319888) (966,1.397335935) (970,1.397148213) (974,1.397110561) (978,1.397054182) (982,1.397185384) (986,1.396752395) (990,1.397432527) (994,1.397696739) (998,1.397905871)};
\addplot[pubNavy, opacity=0.20, line width=0.35pt] coordinates {(10,1.241358226) (14,1.244682404) (18,1.247408705) (22,1.243881974) (26,1.245079208) (30,1.247260282) (34,1.246249593) (38,1.247605971) (42,1.245957573) (46,1.231620882) (50,1.220388803) (54,1.209855451) (58,1.198287002) (62,1.190247519) (66,1.181380603) (70,1.176911301) (74,1.173900139) (78,1.170990441) (82,1.170360186) (86,1.172602351) (90,1.175063211) (94,1.188716942) (98,1.208671493) (102,1.228974604) (106,1.249827458) (110,1.270440302) (114,1.290324667) (118,1.308982981) (122,1.325574399) (126,1.340614211) (130,1.352369444) (134,1.362547598) (138,1.372200677) (142,1.377699322) (146,1.379838255) (150,1.383254956) (154,1.383627356) (158,1.384519106) (162,1.384770349) (166,1.383190946) (170,1.381209896) (174,1.377180064) (178,1.375771847) (182,1.373972217) (186,1.372819417) (190,1.372644811) (194,1.376568629) (198,1.376599897) (202,1.378300568) (206,1.379391878) (210,1.379154313) (214,1.380159427) (218,1.381974581) (222,1.383151069) (226,1.385904733) (230,1.387834664) (234,1.389053295) (238,1.39051151) (242,1.391109258) (246,1.389716051) (250,1.389406429) (254,1.387527479) (258,1.388645815) (262,1.388564325) (266,1.387697882) (270,1.386825378) (274,1.388513624) (278,1.388024688) (282,1.388344325) (286,1.38728809) (290,1.386661285) (294,1.385943888) (298,1.386171466) (302,1.387103043) (306,1.387690024) (310,1.387832905) (314,1.388990883) (318,1.389631707) (322,1.389867969) (326,1.389339406) (330,1.390303799) (334,1.391374385) (338,1.391297517) (342,1.392293359) (346,1.394240097) (350,1.394283991) (354,1.394896357) (358,1.395913147) (362,1.395656773) (366,1.395723891) (370,1.395303827) (374,1.396348439) (378,1.395593168) (382,1.396032432) (386,1.395806885) (390,1.396053484) (394,1.396243649) (398,1.396356947) (402,1.396177698) (406,1.395956558) (410,1.395443113) (414,1.3954729) (418,1.396200555) (422,1.396406579) (426,1.395657386) (430,1.395695416) (434,1.395433509) (438,1.395333352) (442,1.395854838) (446,1.39524232) (450,1.395122022) (454,1.394772572) (458,1.394551065) (462,1.394265062) (466,1.394747736) (470,1.394575353) (474,1.394947674) (478,1.396131372) (482,1.396799896) (486,1.396163864) (490,1.396308542) (494,1.39647158) (498,1.397640822) (502,1.398297836) (506,1.398279517) (510,1.398611158) (514,1.39896497) (518,1.398617041) (522,1.39804959) (526,1.397828113) (530,1.397320945) (534,1.397259921) (538,1.398456917) (542,1.397744775) (546,1.397318958) (550,1.397061362) (554,1.397466297) (558,1.397695123) (562,1.397675199) (566,1.397583461) (570,1.398055516) (574,1.398360674) (578,1.397934538) (582,1.398182169) (586,1.397751308) (590,1.39769682) (594,1.397905222) (598,1.397923035) (602,1.397975724) (606,1.397523233) (610,1.397031376) (614,1.397069371) (618,1.397226614) (622,1.397389644) (626,1.397620685) (630,1.397787884) (634,1.397961958) (638,1.397825891) (642,1.398442472) (646,1.398379243) (650,1.398065828) (654,1.397959645) (658,1.398160286) (662,1.398245699) (666,1.398029138) (670,1.397754733) (674,1.398039625) (678,1.397696004) (682,1.397437613) (686,1.397143409) (690,1.396995604) (694,1.39679971) (698,1.396834961) (702,1.396465083) (706,1.396320472) (710,1.396985373) (714,1.397483875) (718,1.397747272) (722,1.397098199) (726,1.397646727) (730,1.397814822) (734,1.398346675) (738,1.398924613) (742,1.399376259) (746,1.399032722) (750,1.3994973) (754,1.399562402) (758,1.399285864) (762,1.39948232) (766,1.399687931) (770,1.399563213) (774,1.399883602) (778,1.400024215) (782,1.399827739) (786,1.39969801) (790,1.399230342) (794,1.399203499) (798,1.398963577) (802,1.399147756) (806,1.398677144) (810,1.398558862) (814,1.398264384) (818,1.398074688) (822,1.397855857) (826,1.397920722) (830,1.397301094) (834,1.39714775) (838,1.396997308) (842,1.397416911) (846,1.397506577) (850,1.39713819) (854,1.397449014) (858,1.397397744) (862,1.397625634) (866,1.397759283) (870,1.398326931) (874,1.398303406) (878,1.397678492) (882,1.398067748) (886,1.397979263) (890,1.398020695) (894,1.397930857) (898,1.397715422) (902,1.397543631) (906,1.39767049) (910,1.397318019) (914,1.397241415) (918,1.396887282) (922,1.396735218) (926,1.396906391) (930,1.396847736) (934,1.396747165) (938,1.396698912) (942,1.39681615) (946,1.396630606) (950,1.397153212) (954,1.397555983) (958,1.39780418) (962,1.397613404) (966,1.397729444) (970,1.397734819) (974,1.398132837) (978,1.397690563) (982,1.397500269) (986,1.397503745) (990,1.397701908) (994,1.397521253) (998,1.397782647)};
\addplot[pubNavy, opacity=0.20, line width=0.35pt] coordinates {(10,1.257098936) (14,1.253744352) (18,1.249873471) (22,1.247251982) (26,1.247349782) (30,1.249746576) (34,1.25016524) (38,1.251300604) (42,1.249576497) (46,1.235329397) (50,1.223797261) (54,1.210919866) (58,1.198120624) (62,1.187543566) (66,1.179026606) (70,1.171015842) (74,1.163084216) (78,1.156744754) (82,1.148243485) (86,1.141431916) (90,1.136586352) (94,1.138617488) (98,1.145991847) (102,1.154670196) (106,1.16438149) (110,1.174903143) (114,1.184669013) (118,1.195024341) (122,1.206126908) (126,1.219161057) (130,1.232532062) (134,1.247342708) (138,1.263670302) (142,1.278290895) (146,1.291864812) (150,1.306768409) (154,1.319406088) (158,1.331636647) (162,1.342547573) (166,1.352650502) (170,1.361971662) (174,1.367254529) (178,1.370452444) (182,1.373891838) (186,1.372884433) (190,1.37307053) (194,1.374591566) (198,1.374769023) (202,1.37617387) (206,1.375900954) (210,1.376355477) (214,1.377299282) (218,1.378495307) (222,1.378959559) (226,1.380693813) (230,1.381049577) (234,1.383503363) (238,1.384684774) (242,1.385624776) (246,1.387684428) (250,1.386337058) (254,1.387242469) (258,1.388432907) (262,1.388686506) (266,1.388238461) (270,1.387755856) (274,1.389995452) (278,1.389528697) (282,1.388472392) (286,1.388394528) (290,1.387753283) (294,1.388711918) (298,1.389062274) (302,1.389796148) (306,1.389927268) (310,1.390001035) (314,1.390965977) (318,1.391068687) (322,1.390812589) (326,1.390751054) (330,1.391355819) (334,1.39135725) (338,1.391359631) (342,1.390709218) (346,1.391395091) (350,1.390790367) (354,1.391056581) (358,1.390901832) (362,1.390927612) (366,1.390592822) (370,1.390200439) (374,1.390123977) (378,1.390071952) (382,1.390516579) (386,1.389841809) (390,1.390655587) (394,1.390420258) (398,1.391003727) (402,1.390722087) (406,1.391403416) (410,1.390851535) (414,1.391021415) (418,1.391727471) (422,1.392414225) (426,1.391412175) (430,1.391849224) (434,1.392436599) (438,1.392088036) (442,1.39207666) (446,1.391796033) (450,1.391625487) (454,1.39146285) (458,1.391139463) (462,1.390878137) (466,1.390867577) (470,1.390560551) (474,1.390540748) (478,1.391373413) (482,1.391339122) (486,1.391427559) (490,1.391418198) (494,1.391380551) (498,1.391410705) (502,1.391865938) (506,1.392051823) (510,1.392068853) (514,1.392228863) (518,1.392286415) (522,1.391703189) (526,1.39202472) (530,1.392144265) (534,1.391710874) (538,1.391934564) (542,1.392433668) (546,1.392171827) (550,1.392076923) (554,1.392093525) (558,1.392599023) (562,1.393248785) (566,1.393330736) (570,1.394066563) (574,1.393702223) (578,1.393179977) (582,1.393306265) (586,1.39290937) (590,1.392755586) (594,1.39313036) (598,1.393099777) (602,1.393053137) (606,1.392685828) (610,1.392591801) (614,1.391926544) (618,1.391646792) (622,1.391136521) (626,1.391789842) (630,1.392563037) (634,1.393179776) (638,1.392868803) (642,1.3933488) (646,1.393691026) (650,1.39340219) (654,1.393597887) (658,1.394040404) (662,1.39437338) (666,1.394616795) (670,1.394533548) (674,1.394820033) (678,1.394761594) (682,1.394103826) (686,1.394015942) (690,1.394440821) (694,1.39411003) (698,1.394094397) (702,1.394596496) (706,1.395039412) (710,1.395173101) (714,1.395369873) (718,1.395199055) (722,1.394919263) (726,1.394257591) (730,1.394507881) (734,1.394788459) (738,1.395163534) (742,1.395222878) (746,1.394473566) (750,1.394356983) (754,1.393981646) (758,1.393795947) (762,1.393352557) (766,1.393583947) (770,1.394002319) (774,1.394229131) (778,1.394383716) (782,1.394112478) (786,1.394133545) (790,1.39414391) (794,1.394226641) (798,1.394299751) (802,1.394333797) (806,1.39437427) (810,1.394397906) (814,1.394764599) (818,1.395098796) (822,1.394495592) (826,1.394677901) (830,1.394680479) (834,1.39455249) (838,1.393695045) (842,1.393992416) (846,1.393601686) (850,1.394302094) (854,1.394334873) (858,1.393602657) (862,1.393677025) (866,1.393369451) (870,1.39369603) (874,1.39369134) (878,1.39426316) (882,1.394223608) (886,1.394747723) (890,1.395536754) (894,1.394620312) (898,1.394991365) (902,1.394615876) (906,1.394368108) (910,1.394009254) (914,1.393924163) (918,1.393624453) (922,1.394183817) (926,1.3935565) (930,1.393508327) (934,1.393922323) (938,1.393539397) (942,1.393721314) (946,1.393578332) (950,1.393390055) (954,1.39366757) (958,1.394300155) (962,1.394274465) (966,1.394532462) (970,1.394142589) (974,1.394056067) (978,1.393642326) (982,1.393742203) (986,1.393287988) (990,1.392781451) (994,1.393304318) (998,1.39345732)};
\addplot[pubNavy, opacity=0.20, line width=0.35pt] coordinates {(10,1.24905523) (14,1.250731961) (18,1.251755713) (22,1.250293297) (26,1.249830521) (30,1.248060304) (34,1.248872606) (38,1.24710756) (42,1.24490624) (46,1.2321198) (50,1.2223017) (54,1.212353619) (58,1.204498677) (62,1.194935495) (66,1.187000551) (70,1.182806453) (74,1.181392348) (78,1.180808876) (82,1.182982808) (86,1.188048022) (90,1.196298208) (94,1.212698187) (98,1.234507718) (102,1.25713388) (106,1.27900342) (110,1.299936963) (114,1.319635121) (118,1.337701931) (122,1.353503568) (126,1.36749661) (130,1.379084036) (134,1.387621577) (138,1.394215836) (142,1.397921287) (146,1.399976869) (150,1.400392246) (154,1.399064303) (158,1.398664893) (162,1.396520419) (166,1.395251848) (170,1.393277566) (174,1.388647091) (178,1.386575639) (182,1.386416873) (186,1.384485882) (190,1.38422592) (194,1.385102387) (198,1.38413521) (202,1.383809359) (206,1.381078069) (210,1.379503778) (214,1.379455448) (218,1.380271925) (222,1.382067251) (226,1.383893111) (230,1.383399954) (234,1.385125996) (238,1.384049978) (242,1.383807339) (246,1.384818882) (250,1.384691565) (254,1.386361762) (258,1.389284486) (262,1.390712933) (266,1.391390342) (270,1.392484743) (274,1.39379732) (278,1.393544138) (282,1.393678599) (286,1.394044711) (290,1.394770237) (294,1.393768758) (298,1.393270504) (302,1.39249955) (306,1.393371799) (310,1.39278038) (314,1.392271011) (318,1.39103643) (322,1.390713409) (326,1.391065975) (330,1.391203001) (334,1.39041617) (338,1.389272901) (342,1.389984959) (346,1.390394836) (350,1.38977051) (354,1.389923765) (358,1.389385567) (362,1.389158229) (366,1.389204395) (370,1.390274375) (374,1.390142823) (378,1.389466794) (382,1.390667119) (386,1.390696339) (390,1.390568612) (394,1.39144114) (398,1.392078912) (402,1.393442558) (406,1.393886944) (410,1.394286569) (414,1.394768436) (418,1.395318901) (422,1.395437128) (426,1.396435562) (430,1.396550106) (434,1.397286537) (438,1.397680726) (442,1.397869648) (446,1.39774992) (450,1.39689303) (454,1.397170359) (458,1.397012406) (462,1.396843471) (466,1.397148097) (470,1.397302676) (474,1.397543175) (478,1.397204246) (482,1.397131741) (486,1.396949262) (490,1.39685342) (494,1.397075481) (498,1.39772783) (502,1.398087474) (506,1.397922413) (510,1.398195718) (514,1.397664999) (518,1.397737846) (522,1.397325847) (526,1.397306226) (530,1.396920387) (534,1.3971874) (538,1.397663554) (542,1.396774765) (546,1.396367751) (550,1.396833651) (554,1.396527559) (558,1.396779636) (562,1.39663801) (566,1.397143625) (570,1.396998137) (574,1.397543815) (578,1.397772985) (582,1.398017349) (586,1.397536535) (590,1.39793069) (594,1.397531036) (598,1.397559079) (602,1.397421726) (606,1.397292156) (610,1.397488344) (614,1.397713496) (618,1.397046266) (622,1.39722403) (626,1.397674751) (630,1.397538819) (634,1.397704917) (638,1.397959119) (642,1.397933622) (646,1.398599441) (650,1.398503273) (654,1.398788703) (658,1.398771022) (662,1.398953647) (666,1.399212332) (670,1.399210172) (674,1.399267806) (678,1.399112323) (682,1.39873841) (686,1.398491208) (690,1.398450243) (694,1.398533973) (698,1.398111439) (702,1.398257799) (706,1.398575659) (710,1.398812482) (714,1.398895985) (718,1.398913445) (722,1.3988198) (726,1.39879799) (730,1.399314785) (734,1.399971664) (738,1.400103329) (742,1.399851824) (746,1.400099374) (750,1.400176172) (754,1.400061401) (758,1.400036612) (762,1.399890169) (766,1.399867415) (770,1.399833037) (774,1.40011139) (778,1.400205892) (782,1.400141976) (786,1.399952455) (790,1.399715182) (794,1.399560221) (798,1.399524277) (802,1.40022853) (806,1.399929947) (810,1.399535207) (814,1.39925247) (818,1.399609409) (822,1.399117907) (826,1.398917309) (830,1.398698567) (834,1.398729063) (838,1.398519198) (842,1.398772116) (846,1.398738848) (850,1.39873641) (854,1.398995554) (858,1.39939051) (862,1.399741333) (866,1.399757976) (870,1.400367405) (874,1.400412873) (878,1.400243887) (882,1.400500896) (886,1.400661884) (890,1.400969275) (894,1.400761415) (898,1.400869617) (902,1.400598033) (906,1.400208225) (910,1.399953082) (914,1.399844663) (918,1.399667348) (922,1.399511846) (926,1.39954344) (930,1.399668238) (934,1.399587179) (938,1.39972544) (942,1.39936636) (946,1.399176538) (950,1.399273185) (954,1.399667894) (958,1.399717733) (962,1.399299986) (966,1.399285265) (970,1.399215759) (974,1.399340538) (978,1.39898268) (982,1.399069574) (986,1.398841147) (990,1.399048486) (994,1.399005479) (998,1.399133472)};
\addplot[pubNavy, opacity=0.20, line width=0.35pt] coordinates {(10,1.237601024) (14,1.241317737) (18,1.245319319) (22,1.245977161) (26,1.24725479) (30,1.24682377) (34,1.246277933) (38,1.248474141) (42,1.249867709) (46,1.236253722) (50,1.22584517) (54,1.215778142) (58,1.207461455) (62,1.199994684) (66,1.191303356) (70,1.184618181) (74,1.178938366) (78,1.173823073) (82,1.1708989) (86,1.165666496) (90,1.15564378) (94,1.15531801) (98,1.159674134) (102,1.165503729) (106,1.173214433) (110,1.184435427) (114,1.197300606) (118,1.212382905) (122,1.228046522) (126,1.243556028) (130,1.257426785) (134,1.271472732) (138,1.288972106) (142,1.306216796) (146,1.322903934) (150,1.338334356) (154,1.350453045) (158,1.360812236) (162,1.366363137) (166,1.368029046) (170,1.366762551) (174,1.364428919) (178,1.363216481) (182,1.363370118) (186,1.362115785) (190,1.361266556) (194,1.362631023) (198,1.363216581) (202,1.366444336) (206,1.367629836) (210,1.368907951) (214,1.371378525) (218,1.3732365) (222,1.376336684) (226,1.379464205) (230,1.379580352) (234,1.38024452) (238,1.381446549) (242,1.382637399) (246,1.383356268) (250,1.382589493) (254,1.383406803) (258,1.383459259) (262,1.383240547) (266,1.383730021) (270,1.382523003) (274,1.383494964) (278,1.383756294) (282,1.383914529) (286,1.383850389) (290,1.382719721) (294,1.382900261) (298,1.383572024) (302,1.382790152) (306,1.382367625) (310,1.383246689) (314,1.384174835) (318,1.38450968) (322,1.384810141) (326,1.385219944) (330,1.386566825) (334,1.387335667) (338,1.387521036) (342,1.388319795) (346,1.389083919) (350,1.389231854) (354,1.388869679) (358,1.388830227) (362,1.388601819) (366,1.388486036) (370,1.388546931) (374,1.388639664) (378,1.388183917) (382,1.388564414) (386,1.388120563) (390,1.388292722) (394,1.388178075) (398,1.388220905) (402,1.388478771) (406,1.389277479) (410,1.389809798) (414,1.390651095) (418,1.391279948) (422,1.392460347) (426,1.391824808) (430,1.392377567) (434,1.391960921) (438,1.39227828) (442,1.392937282) (446,1.393385447) (450,1.393423234) (454,1.394040585) (458,1.39409502) (462,1.393835388) (466,1.394085057) (470,1.393903779) (474,1.39441878) (478,1.395228326) (482,1.396858397) (486,1.396756329) (490,1.396924824) (494,1.396969482) (498,1.397323043) (502,1.398005053) (506,1.398107315) (510,1.397886649) (514,1.397768278) (518,1.397757539) (522,1.397525097) (526,1.397061735) (530,1.396395803) (534,1.396024801) (538,1.396692882) (542,1.396529287) (546,1.396728742) (550,1.396916645) (554,1.39689866) (558,1.397438463) (562,1.397414051) (566,1.397490598) (570,1.397178948) (574,1.396965895) (578,1.396889905) (582,1.397561626) (586,1.397112639) (590,1.397116593) (594,1.396898074) (598,1.396739765) (602,1.396593455) (606,1.396001669) (610,1.396193924) (614,1.396021895) (618,1.395948896) (622,1.396261701) (626,1.396456369) (630,1.396516875) (634,1.396222668) (638,1.396517515) (642,1.396821825) (646,1.39695785) (650,1.396900252) (654,1.396731889) (658,1.396998) (662,1.396584709) (666,1.39698687) (670,1.39714301) (674,1.39753634) (678,1.397072514) (682,1.396952222) (686,1.397267243) (690,1.396457787) (694,1.39638709) (698,1.396134698) (702,1.395713876) (706,1.395935677) (710,1.396247657) (714,1.39584919) (718,1.395421534) (722,1.394565022) (726,1.395075886) (730,1.394609759) (734,1.395245095) (738,1.395838314) (742,1.396008469) (746,1.397235636) (750,1.397242153) (754,1.397449368) (758,1.397481588) (762,1.396796256) (766,1.39712823) (770,1.397353299) (774,1.3978404) (778,1.397678863) (782,1.397365038) (786,1.397284253) (790,1.396720903) (794,1.396172927) (798,1.395891333) (802,1.395931488) (806,1.395799164) (810,1.395495272) (814,1.395806957) (818,1.39601938) (822,1.396191748) (826,1.395978492) (830,1.395652934) (834,1.395127307) (838,1.39504015) (842,1.395378058) (846,1.395370776) (850,1.395064983) (854,1.395377653) (858,1.395604815) (862,1.395904591) (866,1.396365283) (870,1.395952965) (874,1.395769346) (878,1.395215191) (882,1.395379854) (886,1.395012938) (890,1.39520498) (894,1.394894882) (898,1.395107532) (902,1.394853591) (906,1.394687626) (910,1.394584763) (914,1.394344772) (918,1.393874166) (922,1.394640812) (926,1.395194514) (930,1.395402377) (934,1.396422874) (938,1.39655738) (942,1.396302702) (946,1.396362915) (950,1.39641154) (954,1.396311173) (958,1.396902396) (962,1.396702392) (966,1.396899468) (970,1.396670444) (974,1.396275411) (978,1.396181179) (982,1.396257029) (986,1.395773137) (990,1.395951654) (994,1.395935055) (998,1.395667223)};
\addplot[pubNavy, opacity=0.20, line width=0.35pt] coordinates {(10,1.245480499) (14,1.245694921) (18,1.245946593) (22,1.247267774) (26,1.247581191) (30,1.245371507) (34,1.246169879) (38,1.24690657) (42,1.245106266) (46,1.230698381) (50,1.218929597) (54,1.206953824) (58,1.194453494) (62,1.183161475) (66,1.17198698) (70,1.161366856) (74,1.152643601) (78,1.145777416) (82,1.142774797) (86,1.141328795) (90,1.143218833) (94,1.156985263) (98,1.178251153) (102,1.201015081) (106,1.22460172) (110,1.247525954) (114,1.268656719) (118,1.289106211) (122,1.307209639) (126,1.324203313) (130,1.338424511) (134,1.348115399) (138,1.355278306) (142,1.359143545) (146,1.361471602) (150,1.362319647) (154,1.360677574) (158,1.357834814) (162,1.357987368) (166,1.358851056) (170,1.35917689) (174,1.360791298) (178,1.361222667) (182,1.364345583) (186,1.365985662) (190,1.369576708) (194,1.37338111) (198,1.372872918) (202,1.374931479) (206,1.375105115) (210,1.3778334) (214,1.377280965) (218,1.380028558) (222,1.381225415) (226,1.382556291) (230,1.383717288) (234,1.385819819) (238,1.387308143) (242,1.386628358) (246,1.388270191) (250,1.387168908) (254,1.388176004) (258,1.387285009) (262,1.387878687) (266,1.387896763) (270,1.386100429) (274,1.387585053) (278,1.387343416) (282,1.387557944) (286,1.38794126) (290,1.387702115) (294,1.388613987) (298,1.389286856) (302,1.390363036) (306,1.391173058) (310,1.391527548) (314,1.391783425) (318,1.392405067) (322,1.393148791) (326,1.392052346) (330,1.390991368) (334,1.391480503) (338,1.391616484) (342,1.391143838) (346,1.390243591) (350,1.390578798) (354,1.389995376) (358,1.389720913) (362,1.389789227) (366,1.389189598) (370,1.389106139) (374,1.388889701) (378,1.388565851) (382,1.387505219) (386,1.387847368) (390,1.388049169) (394,1.389363861) (398,1.389694862) (402,1.388647579) (406,1.389220316) (410,1.389910607) (414,1.390519994) (418,1.389947251) (422,1.390179837) (426,1.390866165) (430,1.392482309) (434,1.393717928) (438,1.39396927) (442,1.393533726) (446,1.39376244) (450,1.393755857) (454,1.394288879) (458,1.395030765) (462,1.394851461) (466,1.396012719) (470,1.396597908) (474,1.397024132) (478,1.39692667) (482,1.397048789) (486,1.396475249) (490,1.395752319) (494,1.396604768) (498,1.397075811) (502,1.396921472) (506,1.396789164) (510,1.396687927) (514,1.396722678) (518,1.396264884) (522,1.396255381) (526,1.396439702) (530,1.396372413) (534,1.396560179) (538,1.397904651) (542,1.397878846) (546,1.397482178) (550,1.397667745) (554,1.398295354) (558,1.398648151) (562,1.398585159) (566,1.398212858) (570,1.398490673) (574,1.398276524) (578,1.398312515) (582,1.398155281) (586,1.397540645) (590,1.397053365) (594,1.397194051) (598,1.397636587) (602,1.397305505) (606,1.397277985) (610,1.396240961) (614,1.397174308) (618,1.397017141) (622,1.396728279) (626,1.397103504) (630,1.396937052) (634,1.397571616) (638,1.397147401) (642,1.397484142) (646,1.397496806) (650,1.397104204) (654,1.39654361) (658,1.397061583) (662,1.397502658) (666,1.397362981) (670,1.397249601) (674,1.397396916) (678,1.397556828) (682,1.397378784) (686,1.396971721) (690,1.397963099) (694,1.397623224) (698,1.396870716) (702,1.396860854) (706,1.397392485) (710,1.397477145) (714,1.397488552) (718,1.397692779) (722,1.397623353) (726,1.397520629) (730,1.39745447) (734,1.397229216) (738,1.398111111) (742,1.398224377) (746,1.397892139) (750,1.398810875) (754,1.398924881) (758,1.398784255) (762,1.398458717) (766,1.398807261) (770,1.399170198) (774,1.399124467) (778,1.399529717) (782,1.399321814) (786,1.399335209) (790,1.39890859) (794,1.398844354) (798,1.398762268) (802,1.398759086) (806,1.398947035) (810,1.398700126) (814,1.399068232) (818,1.398720808) (822,1.398601605) (826,1.398553626) (830,1.397978199) (834,1.397766279) (838,1.397909157) (842,1.398136384) (846,1.398522175) (850,1.398358444) (854,1.398212989) (858,1.398412304) (862,1.3982956) (866,1.398709752) (870,1.399026791) (874,1.398770407) (878,1.398925242) (882,1.399289441) (886,1.398925042) (890,1.399032473) (894,1.398898592) (898,1.398607166) (902,1.399060662) (906,1.39892722) (910,1.398617173) (914,1.398346674) (918,1.397642831) (922,1.397621495) (926,1.397541171) (930,1.397239832) (934,1.397271187) (938,1.397184072) (942,1.396979165) (946,1.396807136) (950,1.397250552) (954,1.397137651) (958,1.397302569) (962,1.397264831) (966,1.397675918) (970,1.398038267) (974,1.39836368) (978,1.398349072) (982,1.398609362) (986,1.398388418) (990,1.398916539) (994,1.39899464) (998,1.398840497)};
\addplot[pubNavy, opacity=0.20, line width=0.35pt] coordinates {(10,1.249080429) (14,1.245189665) (18,1.240189499) (22,1.238824412) (26,1.243653379) (30,1.242318151) (34,1.244257198) (38,1.245072564) (42,1.244174374) (46,1.23022642) (50,1.219299432) (54,1.208107335) (58,1.19754157) (62,1.187652392) (66,1.180999201) (70,1.17644885) (74,1.171028865) (78,1.167660385) (82,1.166346315) (86,1.165961476) (90,1.166831473) (94,1.177962055) (98,1.194056412) (102,1.211381562) (106,1.231042778) (110,1.251838255) (114,1.270734375) (118,1.289323983) (122,1.303855746) (126,1.319147) (130,1.33158318) (134,1.339991926) (138,1.346809094) (142,1.351985634) (146,1.354990352) (150,1.360473182) (154,1.361753562) (158,1.360651712) (162,1.359649046) (166,1.35823964) (170,1.358034775) (174,1.356080265) (178,1.354067455) (182,1.356330474) (186,1.356846248) (190,1.361191286) (194,1.364947886) (198,1.366174237) (202,1.368168732) (206,1.369681097) (210,1.370225845) (214,1.372880259) (218,1.374524051) (222,1.376830876) (226,1.3800163) (230,1.380368778) (234,1.38126027) (238,1.382702656) (242,1.382797136) (246,1.384257492) (250,1.38336645) (254,1.384495278) (258,1.385221777) (262,1.385989665) (266,1.386024343) (270,1.385201883) (274,1.386633543) (278,1.386860477) (282,1.387500813) (286,1.387372389) (290,1.386916118) (294,1.387928832) (298,1.38811851) (302,1.38846869) (306,1.388435363) (310,1.389494637) (314,1.390562292) (318,1.391108683) (322,1.390801719) (326,1.390586342) (330,1.390214595) (334,1.389832558) (338,1.390151074) (342,1.390233067) (346,1.391238181) (350,1.391427753) (354,1.391441164) (358,1.391004557) (362,1.391027781) (366,1.39050125) (370,1.390001038) (374,1.39021225) (378,1.390442858) (382,1.392006699) (386,1.391562666) (390,1.391486856) (394,1.392290888) (398,1.392738811) (402,1.392485033) (406,1.393028857) (410,1.393393662) (414,1.393154931) (418,1.39386021) (422,1.394378264) (426,1.393839527) (430,1.394044457) (434,1.394627754) (438,1.394684382) (442,1.39453433) (446,1.394340579) (450,1.393858222) (454,1.393947748) (458,1.393678106) (462,1.393482646) (466,1.393579576) (470,1.393783918) (474,1.394280814) (478,1.39461352) (482,1.394897322) (486,1.393774475) (490,1.393497963) (494,1.394107408) (498,1.394889137) (502,1.395266354) (506,1.394369704) (510,1.394640676) (514,1.394517459) (518,1.39440408) (522,1.393808506) (526,1.393630131) (530,1.393304496) (534,1.393464041) (538,1.394012016) (542,1.393500032) (546,1.39308779) (550,1.392749585) (554,1.393562375) (558,1.393644145) (562,1.393107471) (566,1.393234135) (570,1.393624316) (574,1.393542587) (578,1.394128676) (582,1.39506823) (586,1.394989894) (590,1.395711871) (594,1.395644394) (598,1.39528185) (602,1.395656149) (606,1.39553912) (610,1.395204058) (614,1.395511551) (618,1.39544107) (622,1.395751896) (626,1.395514526) (630,1.394901821) (634,1.394494536) (638,1.394031665) (642,1.394819345) (646,1.395300181) (650,1.395119064) (654,1.395025216) (658,1.395745197) (662,1.39599981) (666,1.396680808) (670,1.396344542) (674,1.396547871) (678,1.396580561) (682,1.396826636) (686,1.396234615) (690,1.396732501) (694,1.395742375) (698,1.395647915) (702,1.395599307) (706,1.395576557) (710,1.396027093) (714,1.395629547) (718,1.395317089) (722,1.395523584) (726,1.394989958) (730,1.395095544) (734,1.395848504) (738,1.396037576) (742,1.395909346) (746,1.396037403) (750,1.396241455) (754,1.396108727) (758,1.395570866) (762,1.395418422) (766,1.395104765) (770,1.394991357) (774,1.395725239) (778,1.396074109) (782,1.395522524) (786,1.39565809) (790,1.395756014) (794,1.396335408) (798,1.396272361) (802,1.396628013) (806,1.396555312) (810,1.395837247) (814,1.396369273) (818,1.396301343) (822,1.395801764) (826,1.395318296) (830,1.394764253) (834,1.394701679) (838,1.394190665) (842,1.39440961) (846,1.394440087) (850,1.394424121) (854,1.3944073) (858,1.395200663) (862,1.39539185) (866,1.395644073) (870,1.396051023) (874,1.395812445) (878,1.39667801) (882,1.396714784) (886,1.396530921) (890,1.397275671) (894,1.397256436) (898,1.396918861) (902,1.39749271) (906,1.39756409) (910,1.397647652) (914,1.397860564) (918,1.39802348) (922,1.397668852) (926,1.397464743) (930,1.396952811) (934,1.397232869) (938,1.397068179) (942,1.397775077) (946,1.39697407) (950,1.397244559) (954,1.397140802) (958,1.396632783) (962,1.396206899) (966,1.395677063) (970,1.39584814) (974,1.396250066) (978,1.395974634) (982,1.395847797) (986,1.396008322) (990,1.396023082) (994,1.396261305) (998,1.396312968)};
\addplot[pubNavy, opacity=0.20, line width=0.35pt] coordinates {(10,1.244413898) (14,1.244958944) (18,1.245916258) (22,1.245805517) (26,1.246823439) (30,1.249516516) (34,1.249589357) (38,1.250058746) (42,1.248329769) (46,1.234769029) (50,1.224282914) (54,1.213307234) (58,1.204493541) (62,1.195459572) (66,1.188288966) (70,1.181147141) (74,1.17568654) (78,1.168076868) (82,1.162785977) (86,1.158875704) (90,1.156157914) (94,1.163061509) (98,1.175702153) (102,1.190172294) (106,1.20562355) (110,1.222135333) (114,1.239181006) (118,1.257154699) (122,1.274142007) (126,1.291487851) (130,1.307881614) (134,1.322957094) (138,1.337729984) (142,1.35069188) (146,1.361000518) (150,1.370038572) (154,1.376940202) (158,1.380683585) (162,1.384205303) (166,1.385492935) (170,1.386155942) (174,1.385542929) (178,1.384650306) (182,1.38551909) (186,1.38581009) (190,1.386203907) (194,1.387085756) (198,1.387033291) (202,1.387632557) (206,1.388596635) (210,1.388633307) (214,1.389620427) (218,1.389936167) (222,1.389535805) (226,1.390885504) (230,1.39091327) (234,1.390312197) (238,1.390425098) (242,1.390154361) (246,1.389873055) (250,1.38997301) (254,1.389534706) (258,1.389901538) (262,1.389658758) (266,1.389415093) (270,1.388633229) (274,1.390474538) (278,1.390271667) (282,1.390792189) (286,1.390696796) (290,1.390098997) (294,1.389965307) (298,1.389624445) (302,1.389083228) (306,1.389796723) (310,1.389946374) (314,1.390469252) (318,1.39075991) (322,1.390522746) (326,1.389329573) (330,1.389011649) (334,1.388131391) (338,1.388067807) (342,1.387789574) (346,1.388747327) (350,1.388848343) (354,1.38756572) (358,1.387537779) (362,1.387263913) (366,1.386987128) (370,1.386118527) (374,1.386901749) (378,1.386738412) (382,1.387132726) (386,1.386447755) (390,1.38608101) (394,1.385309737) (398,1.385533692) (402,1.38590055) (406,1.386684435) (410,1.386186487) (414,1.386261765) (418,1.386111334) (422,1.386926414) (426,1.386666799) (430,1.387210631) (434,1.387738697) (438,1.387559885) (442,1.388324763) (446,1.388551637) (450,1.387765193) (454,1.387776886) (458,1.38825225) (462,1.387969284) (466,1.388342917) (470,1.388415471) (474,1.3882723) (478,1.388525015) (482,1.38909668) (486,1.388694706) (490,1.389785465) (494,1.389698999) (498,1.390337364) (502,1.391019041) (506,1.391062661) (510,1.390889262) (514,1.390664543) (518,1.390731645) (522,1.389921451) (526,1.389608664) (530,1.389144171) (534,1.389549725) (538,1.38939111) (542,1.389231089) (546,1.389111051) (550,1.389523793) (554,1.38927948) (558,1.389562291) (562,1.38959606) (566,1.389281511) (570,1.389251802) (574,1.38892717) (578,1.388838604) (582,1.388707822) (586,1.388358521) (590,1.38813547) (594,1.387731942) (598,1.387530483) (602,1.386800227) (606,1.386460625) (610,1.385726773) (614,1.386019937) (618,1.385930547) (622,1.385461115) (626,1.385603782) (630,1.386462597) (634,1.386872335) (638,1.386825379) (642,1.387415492) (646,1.388541929) (650,1.388584641) (654,1.389003339) (658,1.390228785) (662,1.390490796) (666,1.390843199) (670,1.391271166) (674,1.392256029) (678,1.392858554) (682,1.393278929) (686,1.393383681) (690,1.393681361) (694,1.393068441) (698,1.393434078) (702,1.393665205) (706,1.393813069) (710,1.393859998) (714,1.393518364) (718,1.393809807) (722,1.393413189) (726,1.393120105) (730,1.39308075) (734,1.393278477) (738,1.393687167) (742,1.394394327) (746,1.394499572) (750,1.39434464) (754,1.394649859) (758,1.39463912) (762,1.394706732) (766,1.395108067) (770,1.395095023) (774,1.395476846) (778,1.395798403) (782,1.395598192) (786,1.395573337) (790,1.395466161) (794,1.395247524) (798,1.395180006) (802,1.395135428) (806,1.395263299) (810,1.395376248) (814,1.395783226) (818,1.395852303) (822,1.395859857) (826,1.396310398) (830,1.396000297) (834,1.395892195) (838,1.395640297) (842,1.395477629) (846,1.396214815) (850,1.395877959) (854,1.396322362) (858,1.395972274) (862,1.396082155) (866,1.396628896) (870,1.396450276) (874,1.396641591) (878,1.396277478) (882,1.39631176) (886,1.39646641) (890,1.396710594) (894,1.396162109) (898,1.395784325) (902,1.396160027) (906,1.396725683) (910,1.396316129) (914,1.39641349) (918,1.396518952) (922,1.396389608) (926,1.396311011) (930,1.396554397) (934,1.396771767) (938,1.396711615) (942,1.396909342) (946,1.396655618) (950,1.39707486) (954,1.3966512) (958,1.396897308) (962,1.396384847) (966,1.395631395) (970,1.395345506) (974,1.39506733) (978,1.39456849) (982,1.394830584) (986,1.394113858) (990,1.394282587) (994,1.393906227) (998,1.393563792)};
\addplot[pubNavy, opacity=0.20, line width=0.35pt] coordinates {(10,1.243934376) (14,1.251473139) (18,1.25020825) (22,1.247313892) (26,1.248899473) (30,1.249251193) (34,1.24881756) (38,1.247699219) (42,1.248371464) (46,1.233939562) (50,1.222685198) (54,1.21234147) (58,1.2016289) (62,1.192086166) (66,1.183922832) (70,1.179736984) (74,1.177514413) (78,1.177629583) (82,1.181116843) (86,1.187323219) (90,1.194962838) (94,1.211999495) (98,1.235657664) (102,1.259608352) (106,1.282657851) (110,1.304061889) (114,1.324122096) (118,1.342213903) (122,1.357495086) (126,1.369862731) (130,1.379935339) (134,1.386387963) (138,1.390533874) (142,1.393374068) (146,1.393427479) (150,1.393063255) (154,1.391508291) (158,1.390401164) (162,1.388686978) (166,1.38745903) (170,1.386539912) (174,1.384273579) (178,1.3829202) (182,1.381276986) (186,1.380791017) (190,1.379671885) (194,1.379230954) (198,1.379504594) (202,1.379915954) (206,1.379619144) (210,1.379364395) (214,1.381276596) (218,1.382472891) (222,1.383416101) (226,1.385869118) (230,1.385759265) (234,1.387210414) (238,1.388003221) (242,1.389193029) (246,1.390252221) (250,1.389504327) (254,1.390198163) (258,1.390933901) (262,1.391294019) (266,1.391912519) (270,1.391903733) (274,1.393011929) (278,1.392955356) (282,1.392949021) (286,1.393384132) (290,1.39321191) (294,1.393348979) (298,1.393222364) (302,1.393177876) (306,1.393094648) (310,1.394131046) (314,1.395065658) (318,1.395266864) (322,1.394641389) (326,1.394770294) (330,1.396443652) (334,1.395000669) (338,1.39459307) (342,1.395317003) (346,1.396589215) (350,1.396969352) (354,1.397528706) (358,1.396782027) (362,1.396079244) (366,1.396054131) (370,1.39613616) (374,1.396189546) (378,1.39572228) (382,1.397370839) (386,1.396237608) (390,1.395755515) (394,1.395480198) (398,1.395675612) (402,1.39536457) (406,1.395969039) (410,1.395586362) (414,1.395915375) (418,1.396147863) (422,1.396670853) (426,1.395805153) (430,1.395699167) (434,1.396408404) (438,1.396566026) (442,1.396890152) (446,1.396757771) (450,1.395901537) (454,1.396778125) (458,1.395992876) (462,1.395044565) (466,1.395144475) (470,1.395019465) (474,1.395587703) (478,1.395397467) (482,1.395949636) (486,1.395611837) (490,1.395648853) (494,1.395535814) (498,1.396595318) (502,1.396825837) (506,1.397299729) (510,1.397899115) (514,1.398807529) (518,1.399699817) (522,1.399287833) (526,1.399172952) (530,1.397966947) (534,1.397972886) (538,1.398552457) (542,1.398036363) (546,1.398110637) (550,1.398393517) (554,1.398248238) (558,1.39886789) (562,1.398976886) (566,1.398432691) (570,1.398321436) (574,1.397954482) (578,1.398431354) (582,1.398866922) (586,1.3983647) (590,1.399224297) (594,1.399283269) (598,1.398814498) (602,1.399038304) (606,1.399014911) (610,1.397822128) (614,1.398420486) (618,1.398528259) (622,1.398810464) (626,1.399190399) (630,1.399742605) (634,1.400170149) (638,1.399707022) (642,1.399429551) (646,1.400136092) (650,1.399681278) (654,1.399383674) (658,1.400054781) (662,1.399739262) (666,1.400174633) (670,1.399764084) (674,1.399425346) (678,1.398963562) (682,1.398469898) (686,1.398698635) (690,1.398327394) (694,1.398028296) (698,1.397500805) (702,1.397635256) (706,1.397529003) (710,1.397666147) (714,1.39758767) (718,1.39725877) (722,1.397490741) (726,1.398137357) (730,1.398469207) (734,1.398352877) (738,1.398345665) (742,1.398723801) (746,1.399375198) (750,1.399615742) (754,1.400069553) (758,1.400001129) (762,1.400576402) (766,1.400430367) (770,1.400717975) (774,1.400911159) (778,1.401036508) (782,1.400934447) (786,1.400963607) (790,1.400774515) (794,1.40065198) (798,1.400766984) (802,1.400767903) (806,1.400875652) (810,1.400187) (814,1.400289825) (818,1.40048856) (822,1.400052197) (826,1.399844734) (830,1.399552281) (834,1.399656999) (838,1.399219988) (842,1.399424804) (846,1.399161907) (850,1.398970178) (854,1.399542597) (858,1.399677687) (862,1.400092881) (866,1.400005892) (870,1.400153715) (874,1.399918682) (878,1.40000814) (882,1.400117283) (886,1.400177189) (890,1.40058135) (894,1.400478758) (898,1.400512143) (902,1.400129804) (906,1.39984407) (910,1.399551852) (914,1.399511318) (918,1.39888229) (922,1.398857989) (926,1.399079684) (930,1.399320923) (934,1.399096346) (938,1.39921659) (942,1.398889429) (946,1.398374012) (950,1.39837393) (954,1.398635337) (958,1.398977152) (962,1.399026919) (966,1.398785724) (970,1.399549948) (974,1.399577128) (978,1.399151091) (982,1.398953252) (986,1.399135741) (990,1.399225651) (994,1.399508441) (998,1.399663794)};
\addplot[pubNavy, opacity=0.20, line width=0.35pt] coordinates {(10,1.256841015) (14,1.25824627) (18,1.257502924) (22,1.254507085) (26,1.253527083) (30,1.251116191) (34,1.25023182) (38,1.248362051) (42,1.24573654) (46,1.232477759) (50,1.222431162) (54,1.210245442) (58,1.201911634) (62,1.194036836) (66,1.187389021) (70,1.180906617) (74,1.176348788) (78,1.174939521) (82,1.173131792) (86,1.17392828) (90,1.178572785) (94,1.189913203) (98,1.207234405) (102,1.225132761) (106,1.243606469) (110,1.261549992) (114,1.278301431) (118,1.293242316) (122,1.307307523) (126,1.319961705) (130,1.330467187) (134,1.341228197) (138,1.350107022) (142,1.356205656) (146,1.358997229) (150,1.361287193) (154,1.36242416) (158,1.363857893) (162,1.363509065) (166,1.363287233) (170,1.364844563) (174,1.364642742) (178,1.366273564) (182,1.367795093) (186,1.368984074) (190,1.370893455) (194,1.372660288) (198,1.374530876) (202,1.377084953) (206,1.379067163) (210,1.380107499) (214,1.383671447) (218,1.385427139) (222,1.387967805) (226,1.39019902) (230,1.390464996) (234,1.390901278) (238,1.390937408) (242,1.391736347) (246,1.393334107) (250,1.393319501) (254,1.394215172) (258,1.395318675) (262,1.394391264) (266,1.39371208) (270,1.393031679) (274,1.392566581) (278,1.390860591) (282,1.390552482) (286,1.39014772) (290,1.389124658) (294,1.39011379) (298,1.390572432) (302,1.390167421) (306,1.389467572) (310,1.38984491) (314,1.390508175) (318,1.390712276) (322,1.390668999) (326,1.392211943) (330,1.393422956) (334,1.394053562) (338,1.394059974) (342,1.393840095) (346,1.394388888) (350,1.394485571) (354,1.394753569) (358,1.395008998) (362,1.394867512) (366,1.394818978) (370,1.394676542) (374,1.395050597) (378,1.395209772) (382,1.395386631) (386,1.395778332) (390,1.396207078) (394,1.396511301) (398,1.396775601) (402,1.396450815) (406,1.396929656) (410,1.396772778) (414,1.397028856) (418,1.397296858) (422,1.397749381) (426,1.397868101) (430,1.398532248) (434,1.398409709) (438,1.39886968) (442,1.398626271) (446,1.398487589) (450,1.398745983) (454,1.39832015) (458,1.398071941) (462,1.397595288) (466,1.39788083) (470,1.397556965) (474,1.397692739) (478,1.397517212) (482,1.397694049) (486,1.396872633) (490,1.396405922) (494,1.396474949) (498,1.396505155) (502,1.396459) (506,1.396362492) (510,1.396315745) (514,1.395772141) (518,1.395392473) (522,1.394782461) (526,1.394788294) (530,1.394488326) (534,1.394278828) (538,1.395097589) (542,1.394885355) (546,1.394609481) (550,1.394281318) (554,1.394375023) (558,1.394503455) (562,1.395053011) (566,1.395609538) (570,1.395892701) (574,1.396109267) (578,1.396367974) (582,1.396669015) (586,1.396782219) (590,1.397115922) (594,1.3976146) (598,1.397899873) (602,1.39759137) (606,1.397498094) (610,1.3972381) (614,1.397088107) (618,1.397269408) (622,1.3973248) (626,1.397301135) (630,1.397099613) (634,1.396945714) (638,1.396520119) (642,1.396374247) (646,1.396168903) (650,1.396086265) (654,1.396044729) (658,1.396360912) (662,1.396225041) (666,1.396349891) (670,1.396378428) (674,1.396504024) (678,1.396385415) (682,1.396508367) (686,1.396026552) (690,1.396398683) (694,1.396083441) (698,1.39596937) (702,1.395337042) (706,1.395185979) (710,1.395561556) (714,1.395215244) (718,1.394954208) (722,1.395014796) (726,1.395427107) (730,1.39537852) (734,1.395825457) (738,1.396370556) (742,1.396921592) (746,1.397103738) (750,1.397468708) (754,1.398016811) (758,1.39804518) (762,1.398466205) (766,1.398925617) (770,1.399136256) (774,1.39923602) (778,1.399495934) (782,1.399549306) (786,1.399680244) (790,1.398539294) (794,1.398408977) (798,1.398207337) (802,1.398690377) (806,1.398823338) (810,1.398235933) (814,1.398130948) (818,1.397662735) (822,1.39727859) (826,1.397725173) (830,1.3969618) (834,1.396550611) (838,1.396705023) (842,1.396366961) (846,1.396637971) (850,1.396840276) (854,1.396947452) (858,1.397075489) (862,1.397121117) (866,1.39696625) (870,1.397259532) (874,1.396667299) (878,1.396656596) (882,1.396653999) (886,1.396404547) (890,1.396741297) (894,1.397048602) (898,1.39708001) (902,1.396522001) (906,1.396234763) (910,1.39606124) (914,1.39605598) (918,1.396079505) (922,1.396665987) (926,1.396633464) (930,1.395962305) (934,1.396895554) (938,1.397047949) (942,1.396910383) (946,1.396790741) (950,1.397210247) (954,1.397136242) (958,1.397535388) (962,1.397662376) (966,1.398107614) (970,1.398032922) (974,1.398043344) (978,1.398050686) (982,1.398543136) (986,1.39817092) (990,1.398006051) (994,1.39757961) (998,1.397357482)};
\addplot[pubNavy, opacity=0.20, line width=0.35pt] coordinates {(10,1.263329536) (14,1.257953802) (18,1.253234986) (22,1.255354726) (26,1.253113649) (30,1.250002444) (34,1.24888604) (38,1.249480632) (42,1.248223771) (46,1.234177303) (50,1.223077681) (54,1.210749814) (58,1.199082119) (62,1.189167462) (66,1.181781869) (70,1.173935324) (74,1.167527244) (78,1.163849917) (82,1.163269546) (86,1.162495311) (90,1.162447652) (94,1.173768667) (98,1.191704361) (102,1.210779329) (106,1.230573285) (110,1.251479092) (114,1.273004277) (118,1.293234576) (122,1.312081317) (126,1.329475611) (130,1.344673269) (134,1.358337974) (138,1.370582928) (142,1.38032378) (146,1.387740225) (150,1.393306985) (154,1.397281087) (158,1.39937901) (162,1.399220469) (166,1.398416121) (170,1.39759375) (174,1.395489054) (178,1.393560652) (182,1.392419174) (186,1.390633057) (190,1.389354491) (194,1.387632269) (198,1.38518143) (202,1.382915217) (206,1.381151896) (210,1.378667484) (214,1.377957582) (218,1.376890877) (222,1.376278617) (226,1.375414792) (230,1.373498711) (234,1.372250488) (238,1.372718989) (242,1.373043255) (246,1.373948705) (250,1.37495314) (254,1.376302845) (258,1.377059214) (262,1.378349974) (266,1.379506105) (270,1.379325167) (274,1.38235212) (278,1.3853544) (282,1.387199667) (286,1.388922598) (290,1.389898314) (294,1.390805609) (298,1.391457592) (302,1.391775252) (306,1.392734773) (310,1.393045585) (314,1.39422203) (318,1.394901295) (322,1.394915388) (326,1.394344558) (330,1.394930765) (334,1.393898518) (338,1.393758704) (342,1.392875425) (346,1.393288913) (350,1.392519237) (354,1.392002946) (358,1.391918811) (362,1.391368152) (366,1.390780641) (370,1.390633025) (374,1.390384538) (378,1.390433524) (382,1.390549261) (386,1.390426763) (390,1.390173351) (394,1.389894744) (398,1.390299968) (402,1.390441218) (406,1.389557047) (410,1.389098364) (414,1.388671571) (418,1.388984786) (422,1.389592317) (426,1.387935593) (430,1.38783522) (434,1.388036508) (438,1.388393334) (442,1.388050474) (446,1.387264621) (450,1.386855476) (454,1.38685966) (458,1.386004632) (462,1.385818158) (466,1.386429558) (470,1.38599387) (474,1.386001971) (478,1.386652883) (482,1.388177381) (486,1.387140429) (490,1.387005215) (494,1.387817476) (498,1.389501813) (502,1.389663202) (506,1.390056678) (510,1.391599948) (514,1.391597812) (518,1.392003574) (522,1.391242062) (526,1.391569689) (530,1.390891071) (534,1.390839268) (538,1.391696971) (542,1.390316543) (546,1.389999764) (550,1.390016704) (554,1.39084531) (558,1.390855866) (562,1.390878319) (566,1.390753488) (570,1.390405397) (574,1.390705021) (578,1.390130378) (582,1.390334849) (586,1.3904073) (590,1.391068656) (594,1.391044542) (598,1.390984569) (602,1.390800992) (606,1.390869993) (610,1.389538632) (614,1.390115654) (618,1.390029975) (622,1.390345971) (626,1.39096356) (630,1.39190595) (634,1.39263411) (638,1.392539548) (642,1.393163949) (646,1.393903555) (650,1.394060699) (654,1.393905274) (658,1.395439136) (662,1.395820756) (666,1.396299717) (670,1.396509212) (674,1.397093062) (678,1.396976976) (682,1.396771567) (686,1.395785689) (690,1.395841449) (694,1.396178971) (698,1.395500416) (702,1.395273161) (706,1.395730005) (710,1.394979741) (714,1.395159475) (718,1.394507642) (722,1.394673374) (726,1.39493939) (730,1.39511316) (734,1.395307383) (738,1.395882378) (742,1.396221144) (746,1.39643862) (750,1.396753186) (754,1.397516527) (758,1.39754047) (762,1.397908257) (766,1.398947213) (770,1.399242969) (774,1.399673928) (778,1.399751902) (782,1.399930598) (786,1.399930094) (790,1.400171277) (794,1.400229931) (798,1.400045791) (802,1.400062801) (806,1.399778029) (810,1.399684734) (814,1.399836449) (818,1.399396731) (822,1.399193073) (826,1.39899187) (830,1.398929259) (834,1.399126558) (838,1.398684111) (842,1.398511071) (846,1.398622444) (850,1.398642947) (854,1.398922744) (858,1.399048217) (862,1.399118501) (866,1.399179786) (870,1.399400794) (874,1.3993838) (878,1.399456252) (882,1.399071852) (886,1.399475686) (890,1.399850131) (894,1.400044896) (898,1.399988861) (902,1.399831308) (906,1.399761158) (910,1.399682127) (914,1.399526017) (918,1.399355963) (922,1.399717657) (926,1.399484504) (930,1.399158495) (934,1.399462552) (938,1.399321248) (942,1.39882042) (946,1.398282323) (950,1.398089771) (954,1.397895151) (958,1.39840473) (962,1.398503632) (966,1.398572083) (970,1.398745624) (974,1.398724475) (978,1.398651949) (982,1.39873977) (986,1.39854092) (990,1.398458236) (994,1.397973384) (998,1.398782385)};
\addplot[pubNavy, opacity=0.20, line width=0.35pt] coordinates {(10,1.25177483) (14,1.249268707) (18,1.257578652) (22,1.257016974) (26,1.257243334) (30,1.256180885) (34,1.255722305) (38,1.255954194) (42,1.25552692) (46,1.24173082) (50,1.231349021) (54,1.221512668) (58,1.214105635) (62,1.207068393) (66,1.200647534) (70,1.197479573) (74,1.196665162) (78,1.196036347) (82,1.200422194) (86,1.204095689) (90,1.209728006) (94,1.225179483) (98,1.24663432) (102,1.268654078) (106,1.289667976) (110,1.308789717) (114,1.326274356) (118,1.341339848) (122,1.354462028) (126,1.366114814) (130,1.375449409) (134,1.382171234) (138,1.388485604) (142,1.390895431) (146,1.392366944) (150,1.393395967) (154,1.392300037) (158,1.390913966) (162,1.389289177) (166,1.389260296) (170,1.38818371) (174,1.383240626) (178,1.380363746) (182,1.38035446) (186,1.378557898) (190,1.377715466) (194,1.378036316) (198,1.376270126) (202,1.37595671) (206,1.375082316) (210,1.374340715) (214,1.374791026) (218,1.375961322) (222,1.378440395) (226,1.382287513) (230,1.382643888) (234,1.384784038) (238,1.385396451) (242,1.386401422) (246,1.388661139) (250,1.389719315) (254,1.391623212) (258,1.392750899) (262,1.392892986) (266,1.392322928) (270,1.392060999) (274,1.392788248) (278,1.391984617) (282,1.392166278) (286,1.391812605) (290,1.39112264) (294,1.391129052) (298,1.390106391) (302,1.390180091) (306,1.390906558) (310,1.391438126) (314,1.391584502) (318,1.391536725) (322,1.391402538) (326,1.391544282) (330,1.392420742) (334,1.39299654) (338,1.393108983) (342,1.39331784) (346,1.393823476) (350,1.393629243) (354,1.393168197) (358,1.393436104) (362,1.392973306) (366,1.392718543) (370,1.392447703) (374,1.392823937) (378,1.392711582) (382,1.392588766) (386,1.391873892) (390,1.390096203) (394,1.390293624) (398,1.390658911) (402,1.390197459) (406,1.391225266) (410,1.391427501) (414,1.391832731) (418,1.391347616) (422,1.392248298) (426,1.391575798) (430,1.391552315) (434,1.39161625) (438,1.3922923) (442,1.392519032) (446,1.392652627) (450,1.393075631) (454,1.393439473) (458,1.392733036) (462,1.393064956) (466,1.393562762) (470,1.394046607) (474,1.393908328) (478,1.394368289) (482,1.394949053) (486,1.395193521) (490,1.396022178) (494,1.396346159) (498,1.396413138) (502,1.396380537) (506,1.396257479) (510,1.396058575) (514,1.396081598) (518,1.395532756) (522,1.395773933) (526,1.395764082) (530,1.395677888) (534,1.395299635) (538,1.395514185) (542,1.394817067) (546,1.395289489) (550,1.395710425) (554,1.395971141) (558,1.396699237) (562,1.397035154) (566,1.397067998) (570,1.39709751) (574,1.396758127) (578,1.39697276) (582,1.397393103) (586,1.397284926) (590,1.397415859) (594,1.396971808) (598,1.396764886) (602,1.396171653) (606,1.395826371) (610,1.39579313) (614,1.395839959) (618,1.395515171) (622,1.395187944) (626,1.394813007) (630,1.394797351) (634,1.394881854) (638,1.395364058) (642,1.395708232) (646,1.395820496) (650,1.395440228) (654,1.395105851) (658,1.395590016) (662,1.395637275) (666,1.395579919) (670,1.395457033) (674,1.395898754) (678,1.395845062) (682,1.395532527) (686,1.394974302) (690,1.395415607) (694,1.39545397) (698,1.39501534) (702,1.395798865) (706,1.396039575) (710,1.395549915) (714,1.395263996) (718,1.39522812) (722,1.395370687) (726,1.395499087) (730,1.395890028) (734,1.395796397) (738,1.395721682) (742,1.395681675) (746,1.395612079) (750,1.395743709) (754,1.395726646) (758,1.395700003) (762,1.395687222) (766,1.396048318) (770,1.396442061) (774,1.395854876) (778,1.396021395) (782,1.396156252) (786,1.396083771) (790,1.396135611) (794,1.395645379) (798,1.395567659) (802,1.395450101) (806,1.395347237) (810,1.395364874) (814,1.395476684) (818,1.395696787) (822,1.395359154) (826,1.395412474) (830,1.395079038) (834,1.394872108) (838,1.394245771) (842,1.3940011) (846,1.394325106) (850,1.394879114) (854,1.395735861) (858,1.396006079) (862,1.395202737) (866,1.394954259) (870,1.395057456) (874,1.39491291) (878,1.395434105) (882,1.395797245) (886,1.396181117) (890,1.396787438) (894,1.396850135) (898,1.396465934) (902,1.396220001) (906,1.39591819) (910,1.395988679) (914,1.396397713) (918,1.396101377) (922,1.396521926) (926,1.3964079) (930,1.395204469) (934,1.395392144) (938,1.395244317) (942,1.396146086) (946,1.396511367) (950,1.396953415) (954,1.397029815) (958,1.397156955) (962,1.397456309) (966,1.397546408) (970,1.397356891) (974,1.396848161) (978,1.396762824) (982,1.397709227) (986,1.397808398) (990,1.398130231) (994,1.397604973) (998,1.397257844)};
\addplot[pubNavy, opacity=0.20, line width=0.35pt] coordinates {(10,1.26321657) (14,1.259604616) (18,1.262254975) (22,1.261454548) (26,1.262586871) (30,1.261374982) (34,1.259235821) (38,1.255685629) (42,1.255009909) (46,1.241639398) (50,1.231333178) (54,1.220039501) (58,1.209361315) (62,1.201522578) (66,1.192508006) (70,1.184865167) (74,1.17831848) (78,1.173895866) (82,1.172199592) (86,1.174584818) (90,1.178621213) (94,1.191340093) (98,1.210483249) (102,1.230508535) (106,1.250522735) (110,1.269376238) (114,1.287444598) (118,1.305559164) (122,1.321786543) (126,1.337260028) (130,1.350458533) (134,1.360830969) (138,1.369171289) (142,1.375199846) (146,1.378146125) (150,1.379391646) (154,1.379468638) (158,1.380225487) (162,1.381167363) (166,1.380588189) (170,1.381100918) (174,1.379747759) (178,1.380487418) (182,1.380934645) (186,1.381698559) (190,1.381914483) (194,1.383281604) (198,1.384149984) (202,1.385765252) (206,1.386276031) (210,1.386706408) (214,1.388342677) (218,1.389269625) (222,1.389031227) (226,1.391145178) (230,1.390444827) (234,1.390669955) (238,1.391291451) (242,1.391560977) (246,1.392424716) (250,1.392029688) (254,1.392258483) (258,1.39270771) (262,1.393183708) (266,1.392656829) (270,1.392067254) (274,1.394390807) (278,1.39375021) (282,1.394396537) (286,1.394217242) (290,1.393780118) (294,1.393767523) (298,1.393617082) (302,1.393619842) (306,1.3942454) (310,1.394367783) (314,1.394656357) (318,1.394863724) (322,1.394551492) (326,1.394155973) (330,1.39466274) (334,1.395385741) (338,1.395398727) (342,1.395093781) (346,1.395976794) (350,1.395369176) (354,1.39489806) (358,1.394967223) (362,1.395154401) (366,1.39539868) (370,1.395389367) (374,1.395526209) (378,1.395370392) (382,1.394947243) (386,1.395121779) (390,1.394799689) (394,1.39532752) (398,1.395737634) (402,1.396110503) (406,1.396378075) (410,1.396532454) (414,1.396371741) (418,1.396359501) (422,1.396826488) (426,1.396390618) (430,1.39632513) (434,1.396763672) (438,1.396750234) (442,1.395939215) (446,1.396030672) (450,1.396195258) (454,1.396717969) (458,1.39658166) (462,1.396370094) (466,1.396621911) (470,1.396708932) (474,1.396829712) (478,1.39732121) (482,1.398041571) (486,1.39813664) (490,1.398419387) (494,1.398927205) (498,1.398720205) (502,1.398918592) (506,1.39883035) (510,1.398776539) (514,1.398498988) (518,1.398703511) (522,1.39857312) (526,1.398245944) (530,1.397490808) (534,1.397407215) (538,1.397491819) (542,1.397139959) (546,1.397162432) (550,1.397246561) (554,1.397510118) (558,1.397735061) (562,1.398144907) (566,1.398182046) (570,1.398031933) (574,1.398211862) (578,1.398126554) (582,1.398345092) (586,1.398125164) (590,1.397891503) (594,1.397468226) (598,1.397112316) (602,1.397138513) (606,1.396785046) (610,1.396219829) (614,1.396497736) (618,1.396313127) (622,1.39629467) (626,1.395877529) (630,1.396097518) (634,1.396409443) (638,1.396333431) (642,1.396435638) (646,1.397213765) (650,1.397116604) (654,1.396717525) (658,1.39711754) (662,1.39701699) (666,1.396658333) (670,1.396064358) (674,1.396571691) (678,1.396480026) (682,1.396070493) (686,1.395780974) (690,1.395735309) (694,1.395467022) (698,1.395214101) (702,1.395318368) (706,1.395993488) (710,1.396070844) (714,1.396349349) (718,1.396234815) (722,1.396148187) (726,1.396416985) (730,1.396556058) (734,1.396503464) (738,1.397133463) (742,1.397137135) (746,1.397052554) (750,1.397108652) (754,1.396896826) (758,1.396975355) (762,1.396836023) (766,1.396791227) (770,1.397272458) (774,1.397504254) (778,1.397654508) (782,1.397623665) (786,1.397101185) (790,1.396929516) (794,1.396992167) (798,1.396770459) (802,1.396833985) (806,1.396846335) (810,1.396189359) (814,1.396275017) (818,1.396178236) (822,1.395749141) (826,1.395168221) (830,1.39474029) (834,1.394813155) (838,1.395208649) (842,1.39491076) (846,1.394861903) (850,1.394761688) (854,1.394441479) (858,1.394161956) (862,1.394854386) (866,1.395042891) (870,1.395932459) (874,1.396263345) (878,1.396138168) (882,1.396142721) (886,1.395987229) (890,1.396095836) (894,1.395534702) (898,1.395636039) (902,1.396415181) (906,1.396395908) (910,1.3958896) (914,1.39584319) (918,1.395016894) (922,1.394528299) (926,1.394127181) (930,1.394559882) (934,1.395489108) (938,1.396397403) (942,1.396470286) (946,1.396924487) (950,1.396169026) (954,1.396138887) (958,1.396959982) (962,1.396781352) (966,1.396977135) (970,1.397928719) (974,1.39773759) (978,1.39730602) (982,1.397110235) (986,1.396623979) (990,1.396825533) (994,1.396723332) (998,1.396898528)};
\addplot[pubNavy, opacity=0.20, line width=0.35pt] coordinates {(10,1.258736852) (14,1.260837839) (18,1.256824943) (22,1.258425047) (26,1.255178047) (30,1.253914474) (34,1.255136662) (38,1.254940451) (42,1.25207697) (46,1.237845069) (50,1.225918406) (54,1.212449668) (58,1.198760046) (62,1.184537088) (66,1.174055409) (70,1.165041545) (74,1.156988538) (78,1.15356488) (82,1.150870203) (86,1.152103312) (90,1.155714861) (94,1.169836372) (98,1.190710369) (102,1.212857466) (106,1.234648811) (110,1.256741957) (114,1.276515457) (118,1.295367936) (122,1.312785827) (126,1.328326971) (130,1.340240906) (134,1.350447893) (138,1.358225835) (142,1.364496841) (146,1.366083785) (150,1.366800486) (154,1.366845251) (158,1.367462877) (162,1.369756814) (166,1.370943152) (170,1.374196357) (174,1.374344132) (178,1.374563282) (182,1.375620309) (186,1.376924144) (190,1.37748097) (194,1.378235407) (198,1.379360602) (202,1.381250722) (206,1.38225856) (210,1.384004476) (214,1.384269686) (218,1.384780747) (222,1.385466738) (226,1.38669963) (230,1.387042767) (234,1.387841805) (238,1.3880514) (242,1.388317668) (246,1.388600212) (250,1.388040403) (254,1.38842795) (258,1.388362017) (262,1.388555771) (266,1.387991661) (270,1.387294033) (274,1.388526286) (278,1.388756541) (282,1.389047454) (286,1.388861507) (290,1.388852437) (294,1.388549506) (298,1.390029588) (302,1.390047901) (306,1.3907111) (310,1.39082015) (314,1.391200446) (318,1.391685013) (322,1.391149811) (326,1.390580493) (330,1.39042299) (334,1.389783382) (338,1.388865716) (342,1.388137626) (346,1.387450702) (350,1.387113128) (354,1.386508071) (358,1.386967255) (362,1.387005106) (366,1.386897918) (370,1.387377847) (374,1.38793357) (378,1.387161218) (382,1.387525377) (386,1.388197302) (390,1.388810199) (394,1.389196214) (398,1.389162216) (402,1.388893453) (406,1.388676108) (410,1.388239644) (414,1.38888187) (418,1.388436739) (422,1.388794324) (426,1.388072645) (430,1.388992759) (434,1.388804836) (438,1.388266152) (442,1.388355279) (446,1.388552379) (450,1.388942549) (454,1.389344514) (458,1.38951677) (462,1.388396974) (466,1.388674232) (470,1.388879836) (474,1.389335581) (478,1.39063156) (482,1.392113313) (486,1.391813192) (490,1.392273559) (494,1.392552784) (498,1.392809345) (502,1.393513363) (506,1.393237607) (510,1.393725439) (514,1.3945795) (518,1.395628765) (522,1.394498493) (526,1.394168002) (530,1.393226107) (534,1.392509625) (538,1.393213791) (542,1.393082536) (546,1.392460525) (550,1.392754885) (554,1.392432472) (558,1.392925648) (562,1.392748635) (566,1.392296087) (570,1.392350558) (574,1.39231142) (578,1.391895793) (582,1.392795301) (586,1.39204454) (590,1.391916972) (594,1.391602475) (598,1.391404486) (602,1.391397634) (606,1.391489499) (610,1.390415345) (614,1.390186689) (618,1.390292379) (622,1.389806824) (626,1.390223628) (630,1.390558163) (634,1.391429441) (638,1.391665505) (642,1.391719934) (646,1.392648027) (650,1.39235925) (654,1.391896482) (658,1.392337539) (662,1.39237427) (666,1.392803062) (670,1.392758606) (674,1.393413729) (678,1.393485637) (682,1.39351419) (686,1.393184878) (690,1.393615273) (694,1.393617075) (698,1.39304907) (702,1.393612773) (706,1.394934019) (710,1.395109112) (714,1.395112422) (718,1.39556298) (722,1.395123309) (726,1.395332309) (730,1.395681772) (734,1.39561192) (738,1.395890874) (742,1.39609079) (746,1.396351743) (750,1.396405373) (754,1.396644748) (758,1.396240366) (762,1.396106644) (766,1.396826353) (770,1.397175764) (774,1.397046279) (778,1.396746026) (782,1.396507396) (786,1.396305123) (790,1.396031323) (794,1.39607614) (798,1.395799132) (802,1.394959115) (806,1.395185945) (810,1.395399483) (814,1.395540768) (818,1.395654705) (822,1.39611682) (826,1.396145895) (830,1.396466786) (834,1.396430876) (838,1.396505796) (842,1.396688589) (846,1.396820281) (850,1.397149203) (854,1.397563058) (858,1.397240466) (862,1.396559) (866,1.396256493) (870,1.395350481) (874,1.395596596) (878,1.395440927) (882,1.395453434) (886,1.395216636) (890,1.395438619) (894,1.395038092) (898,1.394869934) (902,1.395553731) (906,1.395286848) (910,1.395645476) (914,1.396244002) (918,1.396517728) (922,1.396466227) (926,1.396145075) (930,1.396214483) (934,1.395840113) (938,1.396029473) (942,1.396283603) (946,1.396040106) (950,1.395891249) (954,1.39614622) (958,1.396138416) (962,1.395990087) (966,1.396136746) (970,1.396516053) (974,1.396832894) (978,1.396477209) (982,1.396898667) (986,1.397275905) (990,1.396781413) (994,1.397210864) (998,1.397402422)};
\addplot[pubNavy, opacity=0.20, line width=0.35pt] coordinates {(10,1.248792711) (14,1.247115125) (18,1.248525968) (22,1.248501893) (26,1.249275335) (30,1.250810317) (34,1.253201407) (38,1.251390023) (42,1.251296498) (46,1.237346302) (50,1.226288087) (54,1.215266699) (58,1.203877647) (62,1.19468455) (66,1.185781523) (70,1.177570947) (74,1.17135509) (78,1.167085269) (82,1.163615023) (86,1.162845239) (90,1.1644847) (94,1.175765378) (98,1.194905449) (102,1.214933351) (106,1.236265921) (110,1.258093674) (114,1.279693354) (118,1.299296831) (122,1.317883573) (126,1.333716191) (130,1.344866835) (134,1.356532323) (138,1.367444253) (142,1.373743289) (146,1.375527158) (150,1.375854585) (154,1.376410395) (158,1.375491318) (162,1.372867201) (166,1.370624767) (170,1.369180087) (174,1.366753141) (178,1.367474085) (182,1.36757995) (186,1.367775738) (190,1.367420545) (194,1.370321061) (198,1.371766178) (202,1.373641401) (206,1.375827383) (210,1.375085497) (214,1.378120097) (218,1.381451284) (222,1.382245993) (226,1.383439123) (230,1.384697548) (234,1.386650721) (238,1.387455869) (242,1.387506171) (246,1.388210637) (250,1.387809258) (254,1.388636419) (258,1.389839907) (262,1.390766166) (266,1.389737441) (270,1.390335826) (274,1.391599819) (278,1.391873373) (282,1.391803211) (286,1.391866276) (290,1.391028728) (294,1.391996815) (298,1.393095873) (302,1.393516524) (306,1.393412521) (310,1.393457926) (314,1.394002339) (318,1.394813514) (322,1.394394183) (326,1.393520523) (330,1.393789922) (334,1.395012224) (338,1.395217525) (342,1.394601351) (346,1.395135091) (350,1.395270736) (354,1.395751962) (358,1.396194116) (362,1.395856881) (366,1.395576746) (370,1.394969207) (374,1.39489746) (378,1.395281802) (382,1.39552971) (386,1.39553962) (390,1.396177794) (394,1.396665996) (398,1.39688588) (402,1.39694263) (406,1.397085338) (410,1.397310129) (414,1.397435415) (418,1.397491597) (422,1.399155024) (426,1.398581199) (430,1.398889604) (434,1.398639868) (438,1.398576437) (442,1.398601545) (446,1.398381919) (450,1.398080429) (454,1.398317887) (458,1.397795354) (462,1.397463972) (466,1.39796847) (470,1.398499229) (474,1.398480529) (478,1.39896605) (482,1.398978366) (486,1.398776065) (490,1.39825212) (494,1.398717877) (498,1.399036487) (502,1.398801541) (506,1.398167497) (510,1.398009415) (514,1.397799099) (518,1.397634194) (522,1.396976695) (526,1.396553242) (530,1.395804844) (534,1.395800842) (538,1.396181555) (542,1.395668271) (546,1.395579632) (550,1.395699534) (554,1.39618755) (558,1.39646542) (562,1.39622625) (566,1.396041893) (570,1.395798086) (574,1.39614611) (578,1.396324352) (582,1.396965208) (586,1.396607734) (590,1.397105407) (594,1.397285049) (598,1.397310857) (602,1.397048164) (606,1.397310771) (610,1.396979493) (614,1.396872957) (618,1.39720083) (622,1.397262633) (626,1.396530612) (630,1.396932092) (634,1.396983343) (638,1.396129198) (642,1.39587138) (646,1.395479739) (650,1.394890859) (654,1.394369774) (658,1.394459881) (662,1.394564433) (666,1.395291789) (670,1.395073537) (674,1.395853289) (678,1.395723569) (682,1.395465079) (686,1.395027045) (690,1.395620737) (694,1.395636383) (698,1.395343941) (702,1.395443015) (706,1.395696842) (710,1.396050699) (714,1.39604771) (718,1.395523284) (722,1.395739282) (726,1.395479298) (730,1.39563952) (734,1.395806638) (738,1.396317552) (742,1.396599842) (746,1.396802525) (750,1.397155122) (754,1.397320483) (758,1.397353295) (762,1.397168664) (766,1.397457754) (770,1.397472208) (774,1.39771435) (778,1.397817832) (782,1.397757834) (786,1.397450994) (790,1.396875895) (794,1.397052146) (798,1.397067184) (802,1.397366164) (806,1.397910785) (810,1.397503361) (814,1.397642295) (818,1.397734573) (822,1.397363252) (826,1.397068717) (830,1.39651908) (834,1.396514042) (838,1.396721068) (842,1.396393108) (846,1.396165926) (850,1.395936845) (854,1.395516276) (858,1.395168819) (862,1.394841943) (866,1.394529977) (870,1.393864176) (874,1.393299951) (878,1.392992993) (882,1.392789779) (886,1.392757649) (890,1.392805612) (894,1.39266015) (898,1.392434666) (902,1.391405429) (906,1.391140452) (910,1.391144745) (914,1.391072181) (918,1.391279246) (922,1.391937677) (926,1.392651884) (930,1.392631215) (934,1.392603333) (938,1.392216297) (942,1.391763307) (946,1.391522533) (950,1.392354855) (954,1.392993564) (958,1.393683453) (962,1.393686842) (966,1.393686324) (970,1.393582306) (974,1.39318671) (978,1.3930303) (982,1.393394824) (986,1.394027673) (990,1.394775655) (994,1.394690482) (998,1.394260868)};
\addplot[pubNavy, opacity=0.20, line width=0.35pt] coordinates {(10,1.242865425) (14,1.247095195) (18,1.24931967) (22,1.250216573) (26,1.251688459) (30,1.250071378) (34,1.252303908) (38,1.251855095) (42,1.248247476) (46,1.236716918) (50,1.228420619) (54,1.219382492) (58,1.213499777) (62,1.207664177) (66,1.199982603) (70,1.191402293) (74,1.181948027) (78,1.174532846) (82,1.167454491) (86,1.161279649) (90,1.158120658) (94,1.163567777) (98,1.174142497) (102,1.185302688) (106,1.198383017) (110,1.212331611) (114,1.227767233) (118,1.245803303) (122,1.265661854) (126,1.286081647) (130,1.30580793) (134,1.323804002) (138,1.340585823) (142,1.354389622) (146,1.366736766) (150,1.375985164) (154,1.382658692) (158,1.387767164) (162,1.390129726) (166,1.390659655) (170,1.389607347) (174,1.386023231) (178,1.384085977) (182,1.383699011) (186,1.38335312) (190,1.383175313) (194,1.383835884) (198,1.383073607) (202,1.383738599) (206,1.384977275) (210,1.384796272) (214,1.386333664) (218,1.387793937) (222,1.389208145) (226,1.391992659) (230,1.391500928) (234,1.39219274) (238,1.392469348) (242,1.393165175) (246,1.394246092) (250,1.394969064) (254,1.395170015) (258,1.395683907) (262,1.395800123) (266,1.395758702) (270,1.39584202) (274,1.397375011) (278,1.397455594) (282,1.397461702) (286,1.396655995) (290,1.396043536) (294,1.395448447) (298,1.394779373) (302,1.394160232) (306,1.393896012) (310,1.392840131) (314,1.393146076) (318,1.393492861) (322,1.39372174) (326,1.393370615) (330,1.394005616) (334,1.394161522) (338,1.394505236) (342,1.3947276) (346,1.395928738) (350,1.396245969) (354,1.396093624) (358,1.397109599) (362,1.397464316) (366,1.397503757) (370,1.396636845) (374,1.396161144) (378,1.395636299) (382,1.39609472) (386,1.396185707) (390,1.395779765) (394,1.395833286) (398,1.395521209) (402,1.394839547) (406,1.394530661) (410,1.394173633) (414,1.393871585) (418,1.393650727) (422,1.394348983) (426,1.392560917) (430,1.392370372) (434,1.392222951) (438,1.392732945) (442,1.393170401) (446,1.393014951) (450,1.393557372) (454,1.394382147) (458,1.394377875) (462,1.394822207) (466,1.395830073) (470,1.395901034) (474,1.396505311) (478,1.398751708) (482,1.399680418) (486,1.398443925) (490,1.398454716) (494,1.398883448) (498,1.39893813) (502,1.398711823) (506,1.39783304) (510,1.397696122) (514,1.397537043) (518,1.397553773) (522,1.397337402) (526,1.397458541) (530,1.397014553) (534,1.396711653) (538,1.397755814) (542,1.397065614) (546,1.397099285) (550,1.397473802) (554,1.397979108) (558,1.398677423) (562,1.398879708) (566,1.398796213) (570,1.39860116) (574,1.398771177) (578,1.398476843) (582,1.39897022) (586,1.398840072) (590,1.398873483) (594,1.398546074) (598,1.398468304) (602,1.398465803) (606,1.398104599) (610,1.39778234) (614,1.397657656) (618,1.39767265) (622,1.397062053) (626,1.396914339) (630,1.396937306) (634,1.397121343) (638,1.39733451) (642,1.397971395) (646,1.399029246) (650,1.398896) (654,1.399201217) (658,1.400204987) (662,1.400310972) (666,1.401045962) (670,1.401529584) (674,1.402508238) (678,1.403238989) (682,1.40346612) (686,1.403497275) (690,1.403611693) (694,1.402892254) (698,1.402795299) (702,1.402601116) (706,1.402193295) (710,1.402116565) (714,1.401636734) (718,1.401587082) (722,1.401282365) (726,1.400898628) (730,1.40067776) (734,1.40071414) (738,1.400892597) (742,1.401558145) (746,1.401650033) (750,1.401642681) (754,1.401964674) (758,1.40178042) (762,1.40105284) (766,1.40108982) (770,1.401224073) (774,1.401343485) (778,1.401389467) (782,1.401335872) (786,1.401337128) (790,1.400225436) (794,1.399683573) (798,1.398929248) (802,1.39923454) (806,1.399427659) (810,1.399711567) (814,1.399481516) (818,1.399476214) (822,1.399170595) (826,1.399641577) (830,1.399227786) (834,1.399169451) (838,1.39940355) (842,1.39961893) (846,1.400241681) (850,1.400293873) (854,1.400408511) (858,1.400408215) (862,1.401119924) (866,1.401502598) (870,1.401652135) (874,1.401686924) (878,1.401293565) (882,1.401557406) (886,1.401622733) (890,1.402585565) (894,1.402161153) (898,1.401777187) (902,1.402150972) (906,1.401901195) (910,1.401941051) (914,1.401521883) (918,1.401256886) (922,1.401524848) (926,1.401585267) (930,1.401748466) (934,1.401901959) (938,1.401679712) (942,1.401814836) (946,1.401853337) (950,1.402296759) (954,1.402062922) (958,1.402015686) (962,1.402051668) (966,1.402185614) (970,1.402148061) (974,1.402018797) (978,1.401356773) (982,1.401697195) (986,1.401730357) (990,1.401634727) (994,1.401247672) (998,1.401029683)};
\addplot[pubNavy, opacity=0.20, line width=0.35pt] coordinates {(10,1.243742403) (14,1.239525289) (18,1.245676054) (22,1.244697228) (26,1.244102474) (30,1.245682179) (34,1.245553871) (38,1.245238516) (42,1.24367749) (46,1.228478866) (50,1.216207436) (54,1.204539489) (58,1.191703806) (62,1.179756893) (66,1.167872143) (70,1.157862159) (74,1.149301471) (78,1.141362469) (82,1.136419387) (86,1.132356893) (90,1.130710798) (94,1.14107968) (98,1.159294236) (102,1.17925899) (106,1.201420514) (110,1.224799275) (114,1.248057475) (118,1.270041438) (122,1.290898171) (126,1.310358727) (130,1.328912096) (134,1.344417021) (138,1.357432404) (142,1.368234499) (146,1.376094338) (150,1.380842299) (154,1.383173409) (158,1.38459374) (162,1.383339709) (166,1.38282511) (170,1.382225127) (174,1.380574137) (178,1.379244594) (182,1.378400556) (186,1.378815778) (190,1.379647042) (194,1.381090581) (198,1.381498427) (202,1.382937101) (206,1.382988226) (210,1.382312096) (214,1.383821766) (218,1.384717953) (222,1.384824284) (226,1.386709974) (230,1.386522568) (234,1.387727178) (238,1.387561106) (242,1.387952856) (246,1.388530707) (250,1.38785755) (254,1.387636354) (258,1.388475538) (262,1.389235452) (266,1.388942874) (270,1.389309454) (274,1.391102996) (278,1.389443745) (282,1.389428608) (286,1.388983606) (290,1.388085537) (294,1.387738097) (298,1.387250534) (302,1.386736014) (306,1.387262698) (310,1.386420299) (314,1.386930266) (318,1.386701346) (322,1.38570679) (326,1.386402223) (330,1.387858245) (334,1.387226707) (338,1.387693432) (342,1.387505886) (346,1.389353947) (350,1.39023871) (354,1.391473156) (358,1.39210046) (362,1.392098983) (366,1.391384853) (370,1.391538805) (374,1.390731684) (378,1.391129865) (382,1.392590985) (386,1.391893516) (390,1.391424488) (394,1.391673539) (398,1.391342332) (402,1.391454564) (406,1.391785514) (410,1.39227908) (414,1.39287707) (418,1.394164211) (422,1.395486795) (426,1.395135044) (430,1.39551125) (434,1.395936237) (438,1.396457648) (442,1.396728271) (446,1.396570265) (450,1.396161778) (454,1.396002528) (458,1.395722144) (462,1.395380053) (466,1.395715368) (470,1.395310373) (474,1.396019412) (478,1.395971429) (482,1.396564542) (486,1.39624593) (490,1.396223955) (494,1.396017842) (498,1.396536626) (502,1.39723461) (506,1.397063333) (510,1.39691173) (514,1.397368844) (518,1.397155908) (522,1.397071378) (526,1.397213563) (530,1.397426137) (534,1.397389626) (538,1.39812625) (542,1.397922112) (546,1.397417507) (550,1.39741584) (554,1.397183064) (558,1.397672149) (562,1.397496089) (566,1.397391372) (570,1.397340387) (574,1.39746991) (578,1.396460843) (582,1.396472322) (586,1.395705917) (590,1.395553267) (594,1.394870605) (598,1.39506783) (602,1.395007253) (606,1.394431068) (610,1.394014529) (614,1.394579067) (618,1.394146385) (622,1.393750155) (626,1.394472489) (630,1.394646092) (634,1.395265407) (638,1.395112974) (642,1.395684674) (646,1.396314914) (650,1.396560567) (654,1.396057883) (658,1.396603799) (662,1.396627578) (666,1.396972224) (670,1.397123507) (674,1.396860451) (678,1.396179581) (682,1.396111512) (686,1.395787342) (690,1.396205607) (694,1.396245302) (698,1.395630952) (702,1.395842756) (706,1.396515762) (710,1.39615447) (714,1.396042591) (718,1.39590645) (722,1.396107387) (726,1.396548468) (730,1.396358981) (734,1.396618494) (738,1.396602727) (742,1.396477185) (746,1.396714057) (750,1.396695231) (754,1.396939034) (758,1.396582307) (762,1.397186451) (766,1.397629839) (770,1.397635916) (774,1.397887017) (778,1.398058602) (782,1.397885914) (786,1.39765316) (790,1.397917255) (794,1.39776049) (798,1.397446493) (802,1.397417934) (806,1.397420687) (810,1.397316152) (814,1.396738993) (818,1.396632915) (822,1.396689734) (826,1.396367657) (830,1.396449601) (834,1.396593543) (838,1.396071689) (842,1.396422016) (846,1.396346077) (850,1.396450976) (854,1.396830291) (858,1.397490351) (862,1.397513149) (866,1.397980183) (870,1.398078793) (874,1.398003614) (878,1.398344469) (882,1.398012887) (886,1.398370976) (890,1.398619974) (894,1.398699676) (898,1.398798919) (902,1.398837012) (906,1.398360859) (910,1.397885869) (914,1.39776917) (918,1.397525642) (922,1.397457145) (926,1.397644379) (930,1.397823456) (934,1.397697514) (938,1.397901885) (942,1.397722073) (946,1.39749919) (950,1.397291752) (954,1.397348112) (958,1.397223002) (962,1.397225402) (966,1.397191504) (970,1.397461899) (974,1.397342398) (978,1.397060709) (982,1.396974955) (986,1.397161545) (990,1.397139808) (994,1.39724485) (998,1.397411403)};
\addplot[pubNavy, ultra thick] coordinates {(10,1.247607525) (14,1.248191916) (18,1.24959657) (22,1.247907893) (26,1.249087404) (30,1.249631546) (34,1.249275251) (38,1.248977387) (42,1.248288623) (46,1.234740901) (50,1.224826661) (54,1.212878451) (58,1.204185594) (62,1.194360693) (66,1.185669085) (70,1.178764119) (74,1.17479334) (78,1.170058163) (82,1.1670916) (86,1.165260008) (90,1.16468741) (94,1.175146133) (98,1.192880386) (102,1.212119514) (106,1.232845794) (110,1.254290106) (114,1.274759867) (118,1.293238446) (122,1.30969442) (126,1.326265142) (130,1.339332709) (134,1.349281646) (138,1.357829119) (142,1.364487542) (146,1.366620553) (150,1.372946578) (154,1.376675299) (158,1.375624869) (162,1.373116276) (166,1.371219792) (170,1.372489996) (174,1.373008546) (178,1.373582782) (182,1.374270002) (186,1.374555103) (190,1.375577368) (194,1.377302473) (198,1.377226523) (202,1.379602276) (206,1.380297621) (210,1.379434086) (214,1.380718012) (218,1.382223736) (222,1.383283585) (226,1.385886925) (230,1.386140916) (234,1.387468796) (238,1.387782163) (242,1.388064674) (246,1.388630675) (250,1.388230191) (254,1.388532184) (258,1.388912937) (262,1.389447105) (266,1.389178984) (270,1.388971341) (274,1.390788767) (278,1.389900182) (282,1.389990545) (286,1.389952654) (290,1.389998656) (294,1.3904597) (298,1.390339411) (302,1.390271564) (306,1.391039808) (310,1.391482837) (314,1.391683964) (318,1.391701713) (322,1.391583914) (326,1.391972288) (330,1.392273213) (334,1.392418188) (338,1.392165092) (342,1.392143475) (346,1.393223882) (350,1.393290107) (354,1.392906961) (358,1.393055143) (362,1.392536144) (366,1.391881141) (370,1.391805627) (374,1.391792343) (378,1.391506335) (382,1.392589875) (386,1.391883704) (390,1.391538623) (394,1.391982213) (398,1.392591033) (402,1.392963796) (406,1.393357206) (410,1.393550201) (414,1.393588734) (418,1.393774385) (422,1.394654198) (426,1.393981333) (430,1.39434597) (434,1.394776982) (438,1.394969259) (442,1.395032345) (446,1.394791449) (450,1.394753916) (454,1.394760734) (458,1.394790915) (462,1.394836834) (466,1.395271139) (470,1.395164919) (474,1.395517782) (478,1.395950605) (482,1.396682219) (486,1.396204897) (490,1.396123066) (494,1.396408869) (498,1.396565972) (502,1.396873654) (506,1.396978214) (510,1.396808705) (514,1.397355442) (518,1.397350917) (522,1.39682474) (526,1.396523183) (530,1.396169686) (534,1.395912822) (538,1.396437218) (542,1.396098779) (546,1.395973692) (550,1.396272038) (554,1.396357555) (558,1.396739436) (562,1.396836582) (566,1.397105812) (570,1.397047824) (574,1.396862011) (578,1.396676041) (582,1.397179156) (586,1.396788031) (590,1.397079386) (594,1.396934941) (598,1.396752325) (602,1.396612023) (606,1.396393358) (610,1.396206877) (614,1.396259815) (618,1.396363547) (622,1.396278186) (626,1.396423827) (630,1.396391312) (634,1.396520364) (638,1.396425473) (642,1.396498713) (646,1.397041059) (650,1.396730409) (654,1.396405902) (658,1.396519898) (662,1.396606144) (666,1.396826516) (670,1.396512569) (674,1.396805349) (678,1.396778769) (682,1.396799101) (686,1.396250591) (690,1.396595144) (694,1.396270058) (698,1.396052034) (702,1.395820811) (706,1.396168138) (710,1.396201063) (714,1.39619853) (718,1.396070632) (722,1.396111688) (726,1.396226771) (730,1.396196852) (734,1.396520525) (738,1.396746594) (742,1.397029363) (746,1.397169687) (750,1.39735543) (754,1.397574011) (758,1.397511029) (762,1.39747283) (766,1.397817169) (770,1.397788466) (774,1.397863709) (778,1.397932524) (782,1.397821874) (786,1.397552077) (790,1.397088479) (794,1.397095199) (798,1.396918822) (802,1.397100074) (806,1.397665736) (810,1.397409757) (814,1.397190644) (818,1.397259226) (822,1.39722448) (826,1.396977446) (830,1.396540891) (834,1.396532327) (838,1.396713045) (842,1.396555302) (846,1.396729126) (850,1.396989233) (854,1.397072462) (858,1.397157978) (862,1.397346507) (866,1.397646978) (870,1.397783556) (874,1.398075791) (878,1.397762908) (882,1.397775368) (886,1.397957841) (890,1.398178482) (894,1.397795292) (898,1.39745975) (902,1.397495048) (906,1.39761729) (910,1.397554517) (914,1.397539752) (918,1.397263651) (922,1.397387731) (926,1.397185567) (930,1.396900274) (934,1.397252028) (938,1.397126126) (942,1.397171296) (946,1.396865812) (950,1.39718173) (954,1.397244457) (958,1.397406926) (962,1.397258523) (966,1.397531437) (970,1.39766502) (974,1.397544292) (978,1.397318442) (982,1.397544364) (986,1.397499379) (990,1.397766566) (994,1.397592292) (998,1.397597025)};
\end{groupplot}
\end{tikzpicture}

\caption{Pre-training behaviour after a 50-episode trailing mean. The panels
show reward, terminal wealth, turnover, and gross exposure. Thin transparent
lines identify individual seeds and the dark line is the cross-seed median.
Evidence class: training-prefix diagnostic.}
\label{fig:appendix_pretraining_stability}
\end{figure}

\begin{figure}[!tbp]
\centering
% AUTO-GENERATED. DO NOT EDIT BY HAND.
% Figure: Optimizer diagnostics across 20 seeds
% Evidence class: training_prefix_diagnostic
% Regenerate with publication_pipeline_draft/generate_publication_tikz.py.
\begin{tikzpicture}
\begin{groupplot}[
 group style={group size=2 by 2, horizontal sep=1.35cm, vertical sep=2.25cm},
 publication axis, width=0.49\linewidth, height=0.3\linewidth,
 title style={font=\footnotesize\bfseries, text=pubInk, yshift=1.2mm},
]
\nextgroupplot[title={Critic loss}, xlabel={Gradient update}, scaled y ticks=true]
\addplot[pubRose, opacity=0.18, line width=0.35pt] coordinates {(100,0.0105656758) (200,0.01011095103) (300,0.009823602314) (400,0.009070752189) (500,0.008596833609) (600,0.008159515293) (700,0.008011725066) (800,0.007742240094) (900,0.007380442383) (1000,0.007129694941) (1100,0.00660170028) (1200,0.00602320265) (1300,0.005501447152) (1400,0.005234840978) (1500,0.004977873014) (1600,0.004710708978) (1700,0.004357156693) (1800,0.004179835622) (1900,0.00416249081) (2000,0.004014754738) (2100,0.003817434097) (2200,0.003725661919) (2300,0.00364608583) (2400,0.003491571289) (2500,0.003387574526) (2600,0.003448004252) (2700,0.003391408594) (2800,0.003293255088) (2900,0.003185501392) (3000,0.003154943814) (3100,0.003091872041) (3200,0.003010027576) (3300,0.00289715156) (3400,0.002895994415) (3500,0.002841398446) (3600,0.002725810488) (3700,0.002702926635) (3800,0.002585107286) (3900,0.002535600809) (4000,0.002432853344) (4100,0.002345044923) (4200,0.002384536865) (4300,0.002467467252) (4400,0.002474749729) (4500,0.002488206036) (4600,0.002479947393) (4700,0.002495550353) (4800,0.002484249999) (4900,0.002483900497) (5000,0.00252584191) (5100,0.002526988648) (5200,0.002455788245) (5300,0.00234952278) (5400,0.002286668005) (5500,0.002199177025) (5600,0.002126570547) (5700,0.001992161025) (5800,0.001993048738) (5900,0.001931172714) (6000,0.001870237431) (6100,0.001874888502) (6200,0.00188767513) (6300,0.001903705916) (6400,0.001882535091) (6500,0.001914443949) (6600,0.001884095673) (6700,0.001896103332) (6800,0.001888666442) (6900,0.00187205387) (7000,0.00186061631) (7100,0.001850370131) (7200,0.001874356496) (7300,0.001885017741) (7400,0.001878026256) (7500,0.001846000378) (7600,0.001887765259) (7700,0.001883397496) (7800,0.001860160532) (7900,0.001821498619) (8000,0.001820328203) (8100,0.001824050676) (8200,0.00180732056) (8300,0.001793323783) (8400,0.001815846516) (8500,0.001826651313) (8600,0.001773983461) (8700,0.001741497812) (8800,0.001774442848) (8900,0.00183671026) (9000,0.00181471539) (9100,0.001879569481) (9200,0.001857423503) (9300,0.001881890139) (9400,0.001872981759) (9500,0.001859434193) (9600,0.001855922642) (9700,0.0018995016) (9800,0.001876132237) (9900,0.001822066971) (10000,0.001885005319) (10100,0.001853154483) (10200,0.001835363253) (10300,0.001798863558) (10400,0.001817693922) (10500,0.001801612589) (10600,0.001873438049) (10700,0.001925356861) (10800,0.001932137215) (10900,0.001944493165) (11000,0.001943787315) (11100,0.001898829232) (11200,0.001919233764) (11300,0.001961623912) (11400,0.001940452517) (11500,0.001985895378) (11600,0.001935019426) (11700,0.001942470332) (11800,0.001969908597) (11900,0.002009073331) (12000,0.001952456275) (12100,0.001954646653) (12200,0.001952574856) (12300,0.001983681426) (12400,0.001945855073) (12500,0.001900064549) (12600,0.001931644755) (12700,0.00182926046) (12800,0.001812872593) (12900,0.001814461476) (13000,0.001863067714) (13100,0.001817996847) (13200,0.001830351469) (13300,0.001777350577) (13400,0.00189867327) (13500,0.001967629674) (13600,0.002014086465) (13700,0.002094012103) (13800,0.002166892751) (13900,0.002150882035) (14000,0.002132195223) (14100,0.002210093837) (14200,0.002254525863) (14300,0.002232123783) (14400,0.002160443517) (14500,0.002142931928) (14600,0.002068378567) (14700,0.002035266627) (14800,0.00212943512) (14900,0.002155257598) (15000,0.002191626851) (15100,0.002243328944) (15200,0.002173666831) (15300,0.002195832517) (15400,0.00227762271) (15500,0.002301039698) (15600,0.002308267867) (15700,0.002336493228) (15800,0.00222148886) (15900,0.002201736835) (16000,0.00216796142) (16100,0.002114997152) (16200,0.002193605015) (16300,0.002224218333) (16400,0.00228128673) (16500,0.002249981789) (16600,0.002251756238) (16700,0.002244299324) (16800,0.002199065383) (16900,0.002217339771) (17000,0.002328982297) (17100,0.002312041819) (17200,0.002362216334) (17300,0.002350378153) (17400,0.002271388844) (17500,0.002281040396) (17600,0.002344541671) (17700,0.002421430242) (17800,0.002431534068) (17900,0.002390884643) (18000,0.002350313694) (18100,0.002370500343) (18200,0.00237883447) (18300,0.00239277595) (18400,0.002362505265) (18500,0.002316818095) (18600,0.002323978941) (18700,0.002356096439) (18800,0.002376251167) (18900,0.002391014295) (19000,0.002359198988) (19100,0.002348341804) (19200,0.00221025435) (19300,0.002265160333) (19400,0.002230136946) (19500,0.002203119302) (19600,0.002197897236) (19700,0.002086553897) (19800,0.002121880779) (19900,0.002118283405) (20000,0.002100515959) (20100,0.002180728596) (20200,0.002215735009) (20300,0.002167215478) (20400,0.002166713157) (20500,0.002208666655) (20600,0.002224294154) (20700,0.002213636774) (20800,0.002214117523) (20900,0.002246627177) (21000,0.002278461552) (21100,0.002244385204) (21200,0.002236234082) (21300,0.002186863788) (21400,0.002213330846) (21500,0.002243320853) (21600,0.002255587396) (21700,0.002319558128) (21800,0.002243084554) (21900,0.002204348589) (22000,0.002154355554) (22100,0.002112524223) (22200,0.002154438302) (22300,0.00221241453) (22400,0.002212493529) (22500,0.002202215639) (22600,0.002129352826) (22700,0.002101841441) (22800,0.002205728518) (22900,0.00224942147) (23000,0.002277727518) (23100,0.002350436803) (23200,0.002355998056) (23300,0.00229707791) (23400,0.002327963826) (23500,0.002304724185) (23600,0.002299814671) (23700,0.002337966906) (23800,0.002307462296)};
\addplot[pubRose, opacity=0.18, line width=0.35pt] coordinates {(100,0.006234288681) (200,0.005543581443) (300,0.005053848804) (400,0.004994599265) (500,0.004971128237) (600,0.004879461524) (700,0.004989525503) (800,0.005038249947) (900,0.00515579634) (1000,0.005015533557) (1100,0.004823775357) (1200,0.004786560545) (1300,0.005170766125) (1400,0.005192454765) (1500,0.005147876358) (1600,0.005244095856) (1700,0.005062735407) (1800,0.004913907428) (1900,0.004673602548) (2000,0.004730432271) (2100,0.004753975081) (2200,0.004639060632) (2300,0.00435425241) (2400,0.004125530436) (2500,0.00424863447) (2600,0.004080223897) (2700,0.004110512789) (2800,0.004189013247) (2900,0.00408615882) (3000,0.003939535283) (3100,0.003756939084) (3200,0.003741064365) (3300,0.003466490842) (3400,0.003478229907) (3500,0.003194594732) (3600,0.00304911756) (3700,0.002890179353) (3800,0.002713085175) (3900,0.002734940522) (4000,0.002761013783) (4100,0.002915869094) (4200,0.002907591313) (4300,0.002964332956) (4400,0.002877605567) (4500,0.002841267711) (4600,0.002861594851) (4700,0.002870846703) (4800,0.002824003156) (4900,0.002796648862) (5000,0.002733063255) (5100,0.002551623643) (5200,0.002578764572) (5300,0.002494616155) (5400,0.002588938386) (5500,0.002545769629) (5600,0.002589332988) (5700,0.002619358827) (5800,0.002666355227) (5900,0.002635827963) (6000,0.002627143823) (6100,0.002617104491) (6200,0.002516879491) (6300,0.002546398109) (6400,0.002522281301) (6500,0.002533742879) (6600,0.002407438878) (6700,0.002366547112) (6800,0.00230696128) (6900,0.002358034404) (7000,0.0023459833) (7100,0.002477774245) (7200,0.00246840074) (7300,0.0024425696) (7400,0.002388708468) (7500,0.002381974726) (7600,0.00244416974) (7700,0.002401135885) (7800,0.00242072579) (7900,0.002444067155) (8000,0.002488676342) (8100,0.002384432009) (8200,0.002404059051) (8300,0.002438591071) (8400,0.002531897952) (8500,0.002656867006) (8600,0.002663878375) (8700,0.002718441095) (8800,0.002730299323) (8900,0.002633039118) (9000,0.002622341737) (9100,0.002688247478) (9200,0.00270829855) (9300,0.002726421319) (9400,0.002704946185) (9500,0.002644469892) (9600,0.002712394204) (9700,0.002651918679) (9800,0.002708717273) (9900,0.002736068517) (10000,0.002701335982) (10100,0.002656506584) (10200,0.002606627531) (10300,0.00257760745) (10400,0.002647566027) (10500,0.002602939727) (10600,0.002583376248) (10700,0.002718421305) (10800,0.0026352335) (10900,0.002663510991) (11000,0.002745762677) (11100,0.002741708118) (11200,0.002717919275) (11300,0.00274416022) (11400,0.0025763529) (11500,0.002603897371) (11600,0.002701952739) (11700,0.002627044835) (11800,0.002619845828) (11900,0.002601532557) (12000,0.002534354653) (12100,0.002564989857) (12200,0.002646221209) (12300,0.002628711017) (12400,0.002661422407) (12500,0.00269673469) (12600,0.002594551397) (12700,0.002628952893) (12800,0.00275694977) (12900,0.002807443286) (13000,0.00287708817) (13100,0.002879749425) (13200,0.002848900435) (13300,0.002984407335) (13400,0.003128611972) (13500,0.003174658585) (13600,0.003227614053) (13700,0.003248360334) (13800,0.003140595276) (13900,0.003138743271) (14000,0.003162701614) (14100,0.003095975984) (14200,0.003156549763) (14300,0.003030639165) (14400,0.002939091832) (14500,0.002888232283) (14600,0.002862500725) (14700,0.00277216339) (14800,0.002755186381) (14900,0.002740107384) (15000,0.002709957119) (15100,0.002750725904) (15200,0.00271381652) (15300,0.002792069572) (15400,0.00283597575) (15500,0.002762002172) (15600,0.002752783732) (15700,0.002804947016) (15800,0.002808726253) (15900,0.002831771597) (16000,0.002784561319) (16100,0.002707317588) (16200,0.002716205409) (16300,0.002656309144) (16400,0.002596451389) (16500,0.002678377135) (16600,0.002656220738) (16700,0.002652572934) (16800,0.002651970671) (16900,0.002662844188) (17000,0.002685231133) (17100,0.002748895925) (17200,0.002746126521) (17300,0.002791977045) (17400,0.00283728647) (17500,0.002832761989) (17600,0.002835627715) (17700,0.002805687138) (17800,0.002821085602) (17900,0.002730511106) (18000,0.002725988883) (18100,0.002714570286) (18200,0.002661647089) (18300,0.00261852229) (18400,0.002586321067) (18500,0.002555597946) (18600,0.00262113167) (18700,0.002765102312) (18800,0.002791915019) (18900,0.002811459498) (19000,0.002865513624) (19100,0.002868483844) (19200,0.002985026757) (19300,0.002946543577) (19400,0.002901887964) (19500,0.002917147474) (19600,0.002861891082) (19700,0.002794161532) (19800,0.002811316657) (19900,0.00285801366) (20000,0.002739606961) (20100,0.002718413994) (20200,0.002665342344) (20300,0.00267344499) (20400,0.002640676661) (20500,0.002602713718) (20600,0.002552785375) (20700,0.002567911963) (20800,0.002511479612) (20900,0.002412343258) (21000,0.002456047735) (21100,0.00254439963) (21200,0.002523664245) (21300,0.002521419362) (21400,0.002606374072) (21500,0.002633326594) (21600,0.002704154141) (21700,0.002646720945) (21800,0.002673866786) (21900,0.002776925429) (22000,0.002795977611) (22100,0.002669213107) (22200,0.002653856319) (22300,0.00265658407) (22400,0.002716383711) (22500,0.002678006981) (22600,0.002629635902) (22700,0.002662251703) (22800,0.002632066258) (22900,0.002620414481) (23000,0.002637210093) (23100,0.002702689194) (23200,0.002670549089) (23300,0.002682932257) (23400,0.002639319748) (23500,0.002695569769) (23600,0.002713504829) (23700,0.002732696664) (23800,0.002706271247)};
\addplot[pubRose, opacity=0.18, line width=0.35pt] coordinates {(100,0.01686268859) (200,0.01783993561) (300,0.01840637314) (400,0.01609456772) (500,0.01481904387) (600,0.01427267864) (700,0.01359806689) (800,0.0132997944) (900,0.01264460488) (1000,0.01235010661) (1100,0.0113103305) (1200,0.01006296962) (1300,0.008672616398) (1400,0.008243080182) (1500,0.007834607502) (1600,0.007247196091) (1700,0.006741066091) (1800,0.005972168222) (1900,0.005518299504) (2000,0.00483026735) (2100,0.004446529015) (2200,0.004183123843) (2300,0.003892215411) (2400,0.00366131051) (2500,0.003407603432) (2600,0.003087798902) (2700,0.002834644215) (2800,0.002736943099) (2900,0.002614635578) (3000,0.002564947284) (3100,0.002504465659) (3200,0.002333316184) (3300,0.002269890392) (3400,0.00219973349) (3500,0.002032037382) (3600,0.001941492897) (3700,0.0019743043) (3800,0.001865229895) (3900,0.00186013463) (4000,0.001778627769) (4100,0.001723596465) (4200,0.001689010661) (4300,0.001660125609) (4400,0.001654533064) (4500,0.001694432169) (4600,0.001711892581) (4700,0.001656597189) (4800,0.001695196808) (4900,0.00173324825) (5000,0.001791302708) (5100,0.001806533104) (5200,0.001806451287) (5300,0.001797566249) (5400,0.001775793219) (5500,0.001782802737) (5600,0.001758411189) (5700,0.001722827437) (5800,0.001705940231) (5900,0.001666346448) (6000,0.001575844991) (6100,0.001568017888) (6200,0.001598792651) (6300,0.001610724593) (6400,0.001621065661) (6500,0.001598198316) (6600,0.001580478926) (6700,0.001586763479) (6800,0.00160892281) (6900,0.0015807883) (7000,0.001655706763) (7100,0.001676031156) (7200,0.00168001561) (7300,0.001691682218) (7400,0.001693115453) (7500,0.001667771768) (7600,0.00172816643) (7700,0.001797683025) (7800,0.001804561564) (7900,0.001882535801) (8000,0.001837508357) (8100,0.001900777419) (8200,0.001948418573) (8300,0.002030297334) (8400,0.002080114966) (8500,0.002175769431) (8600,0.002143826883) (8700,0.002139723406) (8800,0.002223528537) (8900,0.002215603681) (9000,0.002259294153) (9100,0.002252787724) (9200,0.002181915834) (9300,0.002155269857) (9400,0.002135998535) (9500,0.002058098151) (9600,0.002149927325) (9700,0.002146913169) (9800,0.002069943317) (9900,0.002041215275) (10000,0.002064025297) (10100,0.00215699129) (10200,0.002179159899) (10300,0.002135617635) (10400,0.002074687008) (10500,0.002098988765) (10600,0.002075271844) (10700,0.002055251272) (10800,0.002046367875) (10900,0.002026781044) (11000,0.002015996398) (11100,0.002002274524) (11200,0.00206078419) (11300,0.002050195192) (11400,0.002096582961) (11500,0.002156289341) (11600,0.00215471671) (11700,0.002244527033) (11800,0.002218547591) (11900,0.002247645718) (12000,0.002210819989) (12100,0.002075025998) (12200,0.002005659125) (12300,0.002035569621) (12400,0.002077437344) (12500,0.002051793004) (12600,0.002026692743) (12700,0.002005156281) (12800,0.002013652725) (12900,0.002005917998) (13000,0.001981382724) (13100,0.002038404031) (13200,0.002027639467) (13300,0.001947328751) (13400,0.001992266881) (13500,0.001998324459) (13600,0.001968762814) (13700,0.001919102459) (13800,0.001979922666) (13900,0.001983559458) (14000,0.002063110727) (14100,0.002044083574) (14200,0.002057752851) (14300,0.002116610925) (14400,0.002000784944) (14500,0.002021792321) (14600,0.002035766863) (14700,0.002048496157) (14800,0.001999115641) (14900,0.002047195612) (15000,0.002036234993) (15100,0.0020040503) (15200,0.001987913728) (15300,0.002020454628) (15400,0.002037851268) (15500,0.002038158698) (15600,0.002007684612) (15700,0.002030125179) (15800,0.002046834596) (15900,0.001993436867) (16000,0.001972699247) (16100,0.002017606434) (16200,0.002028823365) (16300,0.001966037601) (16400,0.001948627632) (16500,0.001931448968) (16600,0.001948252146) (16700,0.00190166306) (16800,0.001865489769) (16900,0.001893033751) (17000,0.00190576961) (17100,0.001943925756) (17200,0.001877186343) (17300,0.001969037938) (17400,0.002045856044) (17500,0.001985868462) (17600,0.002008389391) (17700,0.002054845064) (17800,0.002071742422) (17900,0.002033668652) (18000,0.002062763914) (18100,0.002007427707) (18200,0.002081025345) (18300,0.001962670218) (18400,0.001919101272) (18500,0.001908323076) (18600,0.001845984545) (18700,0.001796878583) (18800,0.001885410899) (18900,0.001938349067) (19000,0.001883740013) (19100,0.001845703658) (19200,0.001881359844) (19300,0.001899839449) (19400,0.001875557308) (19500,0.001926676836) (19600,0.001974026451) (19700,0.001976670325) (19800,0.001902930555) (19900,0.001848193654) (20000,0.001834599255) (20100,0.001808308158) (20200,0.001882018615) (20300,0.001982942969) (20400,0.001966051175) (20500,0.001904264616) (20600,0.001860848162) (20700,0.001784011349) (20800,0.001739079843) (20900,0.001777562511) (21000,0.001782607834) (21100,0.001877833984) (21200,0.00176162658) (21300,0.001671952358) (21400,0.001728724129) (21500,0.001736185921) (21600,0.001788895112) (21700,0.001844498364) (21800,0.001822818024) (21900,0.001786652417) (22000,0.001795829181) (22100,0.001681129588) (22200,0.001722756261) (22300,0.001727527997) (22400,0.001666237705) (22500,0.001694145321) (22600,0.001646403433) (22700,0.001623554388) (22800,0.001662503567) (22900,0.001710487611) (23000,0.001689213363) (23100,0.001751251006) (23200,0.001689933636) (23300,0.001648886641) (23400,0.001678793796) (23500,0.001652811654) (23600,0.001645450294) (23700,0.001672458253) (23800,0.001662192773)};
\addplot[pubRose, opacity=0.18, line width=0.35pt] coordinates {(100,0.004728771746) (200,0.004541365197) (300,0.004439419912) (400,0.004133016686) (500,0.004019222129) (600,0.003900147664) (700,0.00364765485) (800,0.003640314069) (900,0.003531414912) (1000,0.003415543446) (1100,0.003236538544) (1200,0.003061308758) (1300,0.002904947661) (1400,0.002894854429) (1500,0.002841686574) (1600,0.002756061032) (1700,0.00272445546) (1800,0.002666417276) (1900,0.00265129332) (2000,0.002864849358) (2100,0.002847580402) (2200,0.002788006258) (2300,0.002774402336) (2400,0.002679955796) (2500,0.002709259326) (2600,0.002686753962) (2700,0.002761722612) (2800,0.002681655879) (2900,0.002756762574) (3000,0.002651280258) (3100,0.002672000974) (3200,0.00277850884) (3300,0.00289000757) (3400,0.002932032826) (3500,0.002807837911) (3600,0.00286205723) (3700,0.002904875483) (3800,0.002871793229) (3900,0.002755548549) (4000,0.002599540516) (4100,0.002545823948) (4200,0.002532045962) (4300,0.002373706922) (4400,0.002401065873) (4500,0.002396418247) (4600,0.002318476117) (4700,0.002295853989) (4800,0.002383615915) (4900,0.002422913187) (5000,0.002541722241) (5100,0.00256519462) (5200,0.002462406131) (5300,0.002475950145) (5400,0.002398967207) (5500,0.002485585609) (5600,0.002585811657) (5700,0.002523707272) (5800,0.002513950435) (5900,0.002680017916) (6000,0.002691322938) (6100,0.002662864584) (6200,0.002662294602) (6300,0.002631077031) (6400,0.002673634281) (6500,0.002609704109) (6600,0.002622235892) (6700,0.002666381048) (6800,0.002671793359) (6900,0.002480830904) (7000,0.002482069656) (7100,0.002495870576) (7200,0.00255000497) (7300,0.002653261926) (7400,0.00268491751) (7500,0.002711681579) (7600,0.002637886419) (7700,0.002657193691) (7800,0.002710001497) (7900,0.002736861911) (8000,0.002671447024) (8100,0.002667441638) (8200,0.002582022385) (8300,0.002480000514) (8400,0.002500740625) (8500,0.002544187615) (8600,0.002549928473) (8700,0.002544016135) (8800,0.002463651216) (8900,0.002402834711) (9000,0.002397557045) (9100,0.002438319125) (9200,0.002515036799) (9300,0.002532065334) (9400,0.002481320337) (9500,0.002440966922) (9600,0.002512406651) (9700,0.002449664148) (9800,0.002454184322) (9900,0.002512393077) (10000,0.002547969623) (10100,0.00247211731) (10200,0.002574654412) (10300,0.00268777241) (10400,0.002664993354) (10500,0.002701511304) (10600,0.002625097428) (10700,0.002706217812) (10800,0.002763762954) (10900,0.002774613397) (11000,0.00274947728) (11100,0.002770670084) (11200,0.002776002581) (11300,0.002768439264) (11400,0.002794987126) (11500,0.002737926273) (11600,0.002766381158) (11700,0.002747328905) (11800,0.002719936636) (11900,0.002745675598) (12000,0.002716458356) (12100,0.002785607707) (12200,0.00273407409) (12300,0.002773836674) (12400,0.00278395405) (12500,0.002811974264) (12600,0.002901568636) (12700,0.002873826725) (12800,0.002872621594) (12900,0.002836027718) (13000,0.002916802559) (13100,0.002912895801) (13200,0.002840851876) (13300,0.00267089156) (13400,0.002669294435) (13500,0.002811780176) (13600,0.002673024172) (13700,0.002702414291) (13800,0.002740402776) (13900,0.002719083522) (14000,0.002662288724) (14100,0.002551224851) (14200,0.002605076483) (14300,0.002717872476) (14400,0.002704818547) (14500,0.002640171279) (14600,0.002730374574) (14700,0.002635475318) (14800,0.002510789328) (14900,0.002543618798) (15000,0.002528750955) (15100,0.002598169108) (15200,0.002613538073) (15300,0.002580175863) (15400,0.002634836279) (15500,0.002562034794) (15600,0.002655802818) (15700,0.002686634811) (15800,0.002791335178) (15900,0.002837589243) (16000,0.002970360499) (16100,0.002974051377) (16200,0.002983884211) (16300,0.002966259792) (16400,0.002958262037) (16500,0.002908339025) (16600,0.002749531693) (16700,0.002891620784) (16800,0.003009470808) (16900,0.002977724816) (17000,0.002988692769) (17100,0.002978238952) (17200,0.002930682921) (17300,0.002910154569) (17400,0.002913840092) (17500,0.003034651489) (17600,0.003122497536) (17700,0.003057698812) (17800,0.002914857259) (17900,0.002961875545) (18000,0.002801938588) (18100,0.002891103877) (18200,0.00291998121) (18300,0.003110911138) (18400,0.003010090184) (18500,0.002943881392) (18600,0.002887125663) (18700,0.002939862269) (18800,0.002942872862) (18900,0.002972138324) (19000,0.003034933191) (19100,0.002918150648) (19200,0.002904674062) (19300,0.002830236522) (19400,0.002888858342) (19500,0.002869278123) (19600,0.002839956316) (19700,0.002868431248) (19800,0.002845909027) (19900,0.002790724277) (20000,0.002797682444) (20100,0.002855945379) (20200,0.002880536322) (20300,0.002821809566) (20400,0.002889038203) (20500,0.003036946477) (20600,0.003069555433) (20700,0.002921252861) (20800,0.002967848582) (20900,0.002944223047) (21000,0.00296665749) (21100,0.002930026059) (21200,0.002880690363) (21300,0.00287156452) (21400,0.002968745306) (21500,0.002893599309) (21600,0.002819889155) (21700,0.002873762394) (21800,0.002826467669) (21900,0.00284775116) (22000,0.002821080643) (22100,0.002859776886) (22200,0.002852660278) (22300,0.00281892377) (22400,0.002615180938) (22500,0.00251708054) (22600,0.002643914591) (22700,0.002632662375) (22800,0.002648434462) (22900,0.002598988242) (23000,0.002636392741) (23100,0.00262344724) (23200,0.002725649416) (23300,0.002746932488) (23400,0.002813155018) (23500,0.002885098639) (23600,0.002759162639) (23700,0.002808939223) (23800,0.002861871175)};
\addplot[pubRose, opacity=0.18, line width=0.35pt] coordinates {(100,0.008483438753) (200,0.007513226476) (300,0.007224652295) (400,0.007031050045) (500,0.006478723511) (600,0.006270385425) (700,0.006061463151) (800,0.005867378262) (900,0.005800447996) (1000,0.005775175849) (1100,0.005381866219) (1200,0.00511187322) (1300,0.004886568966) (1400,0.004709725897) (1500,0.004817816964) (1600,0.004782848316) (1700,0.004630078259) (1800,0.00454537617) (1900,0.004322911566) (2000,0.004201149894) (2100,0.004040726856) (2200,0.003982217424) (2300,0.003845579596) (2400,0.003642548877) (2500,0.003426593682) (2600,0.003239149973) (2700,0.003221365134) (2800,0.003096018592) (2900,0.003131428035) (3000,0.00291969832) (3100,0.002911363891) (3200,0.002840904449) (3300,0.002782647684) (3400,0.002782307914) (3500,0.002680896502) (3600,0.002579172049) (3700,0.002436262451) (3800,0.002494811511) (3900,0.002354823484) (4000,0.002340453409) (4100,0.002251599275) (4200,0.002211927774) (4300,0.002233701234) (4400,0.002243152435) (4500,0.002289415605) (4600,0.002374098229) (4700,0.002498342702) (4800,0.002471230575) (4900,0.002510947594) (5000,0.002549585514) (5100,0.00257607922) (5200,0.002617272781) (5300,0.002630642965) (5400,0.002585380175) (5500,0.002595461067) (5600,0.002579848655) (5700,0.002506116778) (5800,0.002463180805) (5900,0.002471162356) (6000,0.00255515154) (6100,0.00263903765) (6200,0.002688045427) (6300,0.002729035984) (6400,0.002721811342) (6500,0.002698653913) (6600,0.002714300598) (6700,0.002758797514) (6800,0.002798712277) (6900,0.002795445849) (7000,0.002702511568) (7100,0.002649800549) (7200,0.00260372567) (7300,0.002499497239) (7400,0.002517544362) (7500,0.002527919179) (7600,0.002545776148) (7700,0.002536647231) (7800,0.002562529221) (7900,0.00256753806) (8000,0.002633497561) (8100,0.00270049423) (8200,0.002650376339) (8300,0.0026668265) (8400,0.00266591236) (8500,0.002659385535) (8600,0.002592491452) (8700,0.002598980349) (8800,0.002586085023) (8900,0.002578857541) (9000,0.002515716199) (9100,0.00247576863) (9200,0.002515331772) (9300,0.00246426916) (9400,0.002508225956) (9500,0.002545061044) (9600,0.002564053691) (9700,0.002580134629) (9800,0.002549181913) (9900,0.002556487534) (10000,0.002537920966) (10100,0.002510592074) (10200,0.002534720709) (10300,0.002700148616) (10400,0.002692338638) (10500,0.002743575699) (10600,0.002763931733) (10700,0.002698808094) (10800,0.002709395369) (10900,0.002644249599) (11000,0.002595687949) (11100,0.002619468456) (11200,0.002617951168) (11300,0.002549883991) (11400,0.0024985019) (11500,0.002353985386) (11600,0.002361117967) (11700,0.002466344798) (11800,0.002427706507) (11900,0.002522910922) (12000,0.00258657271) (12100,0.002511682641) (12200,0.002502753376) (12300,0.002461785753) (12400,0.002524912078) (12500,0.002566521405) (12600,0.002509924397) (12700,0.002433419926) (12800,0.002514306339) (12900,0.002437092084) (13000,0.002523636282) (13100,0.002526683942) (13200,0.002503327117) (13300,0.002474960359) (13400,0.002537793596) (13500,0.002509147022) (13600,0.002485809033) (13700,0.002404072077) (13800,0.002409820596) (13900,0.002488309902) (14000,0.002388907864) (14100,0.00237514287) (14200,0.00236843481) (14300,0.002335851558) (14400,0.002307083912) (14500,0.002326543385) (14600,0.002347867622) (14700,0.002438343689) (14800,0.002368802764) (14900,0.002290878282) (15000,0.002406239463) (15100,0.002646976756) (15200,0.00272995187) (15300,0.00274574979) (15400,0.002677229431) (15500,0.002733527555) (15600,0.002699216397) (15700,0.002681203175) (15800,0.002768028656) (15900,0.002771464305) (16000,0.002706533636) (16100,0.002502398181) (16200,0.002455320407) (16300,0.002637037099) (16400,0.002590520005) (16500,0.002491854038) (16600,0.002505177772) (16700,0.002526113857) (16800,0.002412190707) (16900,0.002422622312) (17000,0.00233501771) (17100,0.002286847448) (17200,0.002222601231) (17300,0.00209724477) (17400,0.002126468881) (17500,0.002206478151) (17600,0.002240089024) (17700,0.002246361854) (17800,0.00221032413) (17900,0.002195375005) (18000,0.002208655153) (18100,0.002227016457) (18200,0.002319813997) (18300,0.002256844391) (18400,0.002228774235) (18500,0.002198391443) (18600,0.002182373812) (18700,0.002143157984) (18800,0.002204303513) (18900,0.002237483114) (19000,0.002344433102) (19100,0.002397122397) (19200,0.002230412816) (19300,0.002293027355) (19400,0.002328873496) (19500,0.002336149244) (19600,0.002312306396) (19700,0.002575861162) (19800,0.002572934062) (19900,0.00260073446) (20000,0.002510441665) (20100,0.002479375096) (20200,0.002597446891) (20300,0.002587455732) (20400,0.002567970718) (20500,0.002573252877) (20600,0.002610662952) (20700,0.002340011159) (20800,0.002301050164) (20900,0.002350124042) (21000,0.002371670259) (21100,0.002362064435) (21200,0.00234866098) (21300,0.002309259446) (21400,0.002327283961) (21500,0.002284355392) (21600,0.002312878962) (21700,0.002348836092) (21800,0.00247046412) (21900,0.002359656664) (22000,0.002362754918) (22100,0.002367896307) (22200,0.002302516904) (22300,0.002293054014) (22400,0.002326599066) (22500,0.002420439688) (22600,0.002353235893) (22700,0.002332344721) (22800,0.00226766339) (22900,0.002337159589) (23000,0.002293955022) (23100,0.002331419941) (23200,0.002432261081) (23300,0.002439478156) (23400,0.002417851798) (23500,0.002407235187) (23600,0.002437784174) (23700,0.002439821954) (23800,0.00245675121)};
\addplot[pubRose, opacity=0.18, line width=0.35pt] coordinates {(100,0.01063700207) (200,0.01084341248) (300,0.01030653063) (400,0.00953766075) (500,0.008836438321) (600,0.008257402573) (700,0.007836483885) (800,0.00743147172) (900,0.007040528788) (1000,0.006888053846) (1100,0.006280037668) (1200,0.00570240384) (1300,0.005252071191) (1400,0.005179166421) (1500,0.005080971448) (1600,0.005052686203) (1700,0.005006585876) (1800,0.00500909565) (1900,0.005077980179) (2000,0.004869055469) (2100,0.004939806648) (2200,0.004862876097) (2300,0.004817010695) (2400,0.00467746756) (2500,0.004453120241) (2600,0.004493754264) (2700,0.004469688237) (2800,0.004438258987) (2900,0.004353412474) (3000,0.004395967163) (3100,0.004313839693) (3200,0.004268768104) (3300,0.004258235032) (3400,0.004102651449) (3500,0.004228007514) (3600,0.00403248095) (3700,0.003951700102) (3800,0.003998591169) (3900,0.004012683709) (4000,0.004012184823) (4100,0.003829996637) (4200,0.003728798893) (4300,0.003710883413) (4400,0.003665494733) (4500,0.003539412166) (4600,0.00345166896) (4700,0.003510302771) (4800,0.003352261148) (4900,0.003357798606) (5000,0.0032850018) (5100,0.003309236886) (5200,0.003337462572) (5300,0.003124596551) (5400,0.003166216938) (5500,0.003192271781) (5600,0.003200863022) (5700,0.003087213263) (5800,0.003117676917) (5900,0.003085753461) (6000,0.002949124482) (6100,0.002898576343) (6200,0.003016672097) (6300,0.003127170145) (6400,0.003185522836) (6500,0.003233277891) (6600,0.003294489719) (6700,0.00332903876) (6800,0.00329161007) (6900,0.003363997536) (7000,0.00354943322) (7100,0.003562010499) (7200,0.003507568594) (7300,0.003483268665) (7400,0.003391714394) (7500,0.003311892133) (7600,0.003267044923) (7700,0.003198005655) (7800,0.003148616035) (7900,0.002962382045) (8000,0.002854879806) (8100,0.003002983518) (8200,0.002864326257) (8300,0.00291402142) (8400,0.002932550199) (8500,0.002931022807) (8600,0.00297812121) (8700,0.002923226217) (8800,0.002935561119) (8900,0.002932571503) (9000,0.002978124749) (9100,0.002849817462) (9200,0.002842203504) (9300,0.002841444104) (9400,0.002761066705) (9500,0.002757039992) (9600,0.002677357779) (9700,0.002811504598) (9800,0.002786969091) (9900,0.002839896013) (10000,0.002835494373) (10100,0.002833358594) (10200,0.002863052022) (10300,0.002861677622) (10400,0.002877080278) (10500,0.002985355491) (10600,0.003089035349) (10700,0.003081614943) (10800,0.003182383766) (10900,0.00315089277) (11000,0.003133540764) (11100,0.00312901549) (11200,0.003249636944) (11300,0.003135331068) (11400,0.003157691099) (11500,0.003050716151) (11600,0.002950678463) (11700,0.002921128296) (11800,0.002818836877) (11900,0.002884835843) (12000,0.002838278632) (12100,0.002964694612) (12200,0.002943511913) (12300,0.003080914379) (12400,0.003017759952) (12500,0.00300858377) (12600,0.003014107165) (12700,0.00297400346) (12800,0.003015579586) (12900,0.002972725732) (13000,0.003066997253) (13100,0.00308381319) (13200,0.003022593446) (13300,0.002996516018) (13400,0.00307969423) (13500,0.003069427167) (13600,0.003127623652) (13700,0.003085207636) (13800,0.0030398587) (13900,0.003127245465) (14000,0.003074611863) (14100,0.002917330223) (14200,0.00283808962) (14300,0.002817693097) (14400,0.00274495054) (14500,0.002864218317) (14600,0.002821439877) (14700,0.002869477472) (14800,0.00287026125) (14900,0.002835808438) (15000,0.002803926473) (15100,0.002847056324) (15200,0.002931889542) (15300,0.002882487467) (15400,0.002959280624) (15500,0.002850407176) (15600,0.002903749421) (15700,0.002878637682) (15800,0.002966172667) (15900,0.002878467878) (16000,0.002844602778) (16100,0.002878838149) (16200,0.00288632852) (16300,0.002852737415) (16400,0.002788801561) (16500,0.002748707635) (16600,0.002731691673) (16700,0.002731329086) (16800,0.002705058805) (16900,0.002696732711) (17000,0.002805030998) (17100,0.002719152882) (17200,0.002635919373) (17300,0.002702407376) (17400,0.002735079452) (17500,0.002741557802) (17600,0.002674222831) (17700,0.002705808356) (17800,0.002685768437) (17900,0.002675091987) (18000,0.002662356989) (18100,0.002665768028) (18200,0.002712471271) (18300,0.002693537856) (18400,0.00265643301) (18500,0.002802053932) (18600,0.002763438714) (18700,0.002768754098) (18800,0.002781311306) (18900,0.002866298449) (19000,0.002839070861) (19100,0.002884031297) (19200,0.002866875986) (19300,0.002852031193) (19400,0.002856673673) (19500,0.002831215318) (19600,0.002959195292) (19700,0.003010251001) (19800,0.003100453969) (19900,0.00315901225) (20000,0.003188253543) (20100,0.003209619294) (20200,0.003243192262) (20300,0.003345452528) (20400,0.003314778139) (20500,0.003174768505) (20600,0.003067954653) (20700,0.002977272752) (20800,0.002862144657) (20900,0.002693147422) (21000,0.002632816508) (21100,0.002516400302) (21200,0.002492828923) (21300,0.002459132415) (21400,0.002513926127) (21500,0.002572770114) (21600,0.002675104002) (21700,0.002682773094) (21800,0.002605992695) (21900,0.002585959213) (22000,0.002550394775) (22100,0.002699408901) (22200,0.002669183642) (22300,0.002688478737) (22400,0.0026436763) (22500,0.002639991895) (22600,0.002551385469) (22700,0.00248841896) (22800,0.002554536366) (22900,0.002697410854) (23000,0.00268495928) (23100,0.002571458928) (23200,0.002612964367) (23300,0.002539046598) (23400,0.002551617473) (23500,0.002580757532) (23600,0.002583443793) (23700,0.002649796568) (23800,0.002708964096)};
\addplot[pubRose, opacity=0.18, line width=0.35pt] coordinates {(100,0.009643000551) (200,0.008389517199) (300,0.007997559539) (400,0.007391461055) (500,0.007176585402) (600,0.007173391137) (700,0.006830230222) (800,0.006789376144) (900,0.006872353661) (1000,0.006585766282) (1100,0.006099870801) (1200,0.00585460295) (1300,0.00558193028) (1400,0.005467522796) (1500,0.005322897807) (1600,0.004989463044) (1700,0.004976792075) (1800,0.004835874075) (1900,0.004447335051) (2000,0.004553927761) (2100,0.0044508697) (2200,0.004373968067) (2300,0.004314327915) (2400,0.004202969861) (2500,0.004039469711) (2600,0.004032083671) (2700,0.003939718846) (2800,0.003832629742) (2900,0.003724950808) (3000,0.003507635207) (3100,0.003464409569) (3200,0.003324500285) (3300,0.003267419548) (3400,0.003155338834) (3500,0.003094680584) (3600,0.002969332528) (3700,0.002834752714) (3800,0.002620710642) (3900,0.002614478325) (4000,0.002597519499) (4100,0.002472042758) (4200,0.002433170052) (4300,0.002326545492) (4400,0.002250864054) (4500,0.002168504265) (4600,0.002110257198) (4700,0.00201622824) (4800,0.001996667741) (4900,0.001901840628) (5000,0.001804461738) (5100,0.001804280409) (5200,0.001801920717) (5300,0.001765352173) (5400,0.001815055322) (5500,0.001766008569) (5600,0.001730315399) (5700,0.001723805233) (5800,0.001698578289) (5900,0.001677914499) (6000,0.001645579992) (6100,0.001612053148) (6200,0.001536321372) (6300,0.001456665108) (6400,0.001408503496) (6500,0.001425847539) (6600,0.001431436127) (6700,0.001452209754) (6800,0.001422673301) (6900,0.001411176589) (7000,0.001399318082) (7100,0.001373953372) (7200,0.001410143916) (7300,0.001475194842) (7400,0.001474930544) (7500,0.001446116448) (7600,0.001468829624) (7700,0.001458406122) (7800,0.001510417974) (7900,0.001549359236) (8000,0.001555199607) (8100,0.001562326518) (8200,0.001554048469) (8300,0.001485499018) (8400,0.001480151433) (8500,0.001501163805) (8600,0.001495036657) (8700,0.0014703764) (8800,0.001421945402) (8900,0.001405364193) (9000,0.001424405747) (9100,0.001401106373) (9200,0.001390378596) (9300,0.001439043204) (9400,0.001461687963) (9500,0.001449435705) (9600,0.001447331929) (9700,0.001469186344) (9800,0.001466992393) (9900,0.001467169821) (10000,0.001427903352) (10100,0.001424106699) (10200,0.001395098981) (10300,0.001377437916) (10400,0.001351203513) (10500,0.001395111857) (10600,0.001338058023) (10700,0.001329891291) (10800,0.001372687425) (10900,0.001353445754) (11000,0.001384338143) (11100,0.001374249056) (11200,0.001350877725) (11300,0.00135142178) (11400,0.001347592368) (11500,0.00130540129) (11600,0.001367368316) (11700,0.001391645253) (11800,0.001371572446) (11900,0.001412971283) (12000,0.001378407516) (12100,0.00140892237) (12200,0.001428159652) (12300,0.001433639694) (12400,0.001416175277) (12500,0.001437132864) (12600,0.00138567033) (12700,0.001357058319) (12800,0.001380640513) (12900,0.001355988916) (13000,0.001397985441) (13100,0.001407262159) (13200,0.001448029664) (13300,0.00147419424) (13400,0.001504181756) (13500,0.00161370628) (13600,0.001596616476) (13700,0.001630119735) (13800,0.001687989885) (13900,0.001656027336) (14000,0.001630452939) (14100,0.00160284146) (14200,0.001553935965) (14300,0.001543836144) (14400,0.001524658781) (14500,0.0013842165) (14600,0.001418867405) (14700,0.001468930661) (14800,0.001400504669) (14900,0.001413127908) (15000,0.001400504645) (15100,0.001422256336) (15200,0.001445817121) (15300,0.001421537914) (15400,0.00143367846) (15500,0.001465470798) (15600,0.001464167295) (15700,0.001477389841) (15800,0.001459298481) (15900,0.001515700598) (16000,0.001524336042) (16100,0.001478871854) (16200,0.001542081905) (16300,0.001566509565) (16400,0.001576629758) (16500,0.001526656747) (16600,0.001555811241) (16700,0.00147304663) (16800,0.001497036347) (16900,0.00149888536) (17000,0.001555703918) (17100,0.001574672584) (17200,0.001495170395) (17300,0.001486240572) (17400,0.001464847371) (17500,0.001507603261) (17600,0.001477567712) (17700,0.001541371713) (17800,0.001552594244) (17900,0.001501068217) (18000,0.001477490575) (18100,0.001476066117) (18200,0.001482210681) (18300,0.00144993891) (18400,0.001482515549) (18500,0.001480816992) (18600,0.001461850165) (18700,0.001362737804) (18800,0.001333405578) (18900,0.00132541951) (19000,0.001306880894) (19100,0.001318295742) (19200,0.001359747734) (19300,0.001399807085) (19400,0.001393893885) (19500,0.001370925189) (19600,0.001373867074) (19700,0.001429123175) (19800,0.001429824322) (19900,0.001440678793) (20000,0.001479581615) (20100,0.001471366931) (20200,0.001446581841) (20300,0.001403922925) (20400,0.001360985159) (20500,0.001360301033) (20600,0.001400906919) (20700,0.001344905735) (20800,0.001333268348) (20900,0.001327635802) (21000,0.001260753581) (21100,0.001297154487) (21200,0.001299818465) (21300,0.001325124293) (21400,0.001326277596) (21500,0.001325722202) (21600,0.001281439641) (21700,0.001257381751) (21800,0.00128210549) (21900,0.001278796955) (22000,0.00127793412) (22100,0.001216403395) (22200,0.001187025802) (22300,0.001203178603) (22400,0.001212179498) (22500,0.001220119069) (22600,0.001346651954) (22700,0.001393311261) (22800,0.00138433004) (22900,0.001379240828) (23000,0.001392656576) (23100,0.001416531054) (23200,0.001425362716) (23300,0.001403310255) (23400,0.001471946971) (23500,0.001470884262) (23600,0.001354245341) (23700,0.001341916842) (23800,0.001328830339)};
\addplot[pubRose, opacity=0.18, line width=0.35pt] coordinates {(100,0.01998816617) (200,0.01904482301) (300,0.01709521251) (400,0.01653078618) (500,0.01823094524) (600,0.01791876958) (700,0.01763880838) (800,0.01713126223) (900,0.01618416193) (1000,0.01596491309) (1100,0.01502517182) (1200,0.01392869968) (1300,0.01388661005) (1400,0.01379621997) (1500,0.01208696347) (1600,0.01134796608) (1700,0.01042709113) (1800,0.009696057765) (1900,0.009469899954) (2000,0.008558832342) (2100,0.008135559689) (2200,0.007831908483) (2300,0.00722575495) (2400,0.006273292936) (2500,0.00596020231) (2600,0.00565448897) (2700,0.005431752466) (2800,0.005252423789) (2900,0.00496459303) (3000,0.004979495029) (3100,0.004733799561) (3200,0.004748573597) (3300,0.004399390332) (3400,0.004291407042) (3500,0.00411259532) (3600,0.003774743341) (3700,0.003662135173) (3800,0.003600642411) (3900,0.003614491015) (4000,0.003434087546) (4100,0.003385391179) (4200,0.00322501827) (4300,0.003227285901) (4400,0.003170095803) (4500,0.003212366905) (4600,0.003291182779) (4700,0.003275800636) (4800,0.003315535141) (4900,0.003338982048) (5000,0.003331209184) (5100,0.003296383307) (5200,0.003334130556) (5300,0.003411428211) (5400,0.003440434043) (5500,0.003461539233) (5600,0.003382175276) (5700,0.003348689014) (5800,0.003300671279) (5900,0.003181721945) (6000,0.003204119578) (6100,0.003282354097) (6200,0.003320063953) (6300,0.003284418141) (6400,0.003348765778) (6500,0.003266596794) (6600,0.003373714513) (6700,0.00330822512) (6800,0.003238237463) (6900,0.003429767909) (7000,0.003423639736) (7100,0.003348701075) (7200,0.003299470386) (7300,0.003260761453) (7400,0.003255267721) (7500,0.00329914724) (7600,0.003230866953) (7700,0.003256813344) (7800,0.003263760521) (7900,0.003053777083) (8000,0.003111690027) (8100,0.003044188907) (8200,0.003050706023) (8300,0.003028756427) (8400,0.002937951451) (8500,0.00291969846) (8600,0.002851397544) (8700,0.002869895496) (8800,0.002832299075) (8900,0.002801079) (9000,0.002642068197) (9100,0.002758800075) (9200,0.002781076287) (9300,0.002747401036) (9400,0.002689502062) (9500,0.002795864455) (9600,0.002850663965) (9700,0.002881136094) (9800,0.00293799499) (9900,0.003010545857) (10000,0.003095766483) (10100,0.00298608914) (10200,0.002922141552) (10300,0.002968735411) (10400,0.0029899535) (10500,0.002893916261) (10600,0.002916299179) (10700,0.002922374662) (10800,0.00289385824) (10900,0.002951786923) (11000,0.002976977592) (11100,0.003027784149) (11200,0.003071751283) (11300,0.003027289314) (11400,0.003084986657) (11500,0.003122576931) (11600,0.003160550655) (11700,0.003140545776) (11800,0.003207466239) (11900,0.003180228313) (12000,0.003153221286) (12100,0.003117769212) (12200,0.003266023332) (12300,0.003256872087) (12400,0.003241540422) (12500,0.003111431701) (12600,0.003047409467) (12700,0.003091007425) (12800,0.003042269428) (12900,0.003000718309) (13000,0.003017346701) (13100,0.003060742328) (13200,0.002953526517) (13300,0.002974062483) (13400,0.002984394459) (13500,0.003208439867) (13600,0.00328236362) (13700,0.003258289187) (13800,0.00317059774) (13900,0.003208710416) (14000,0.003145338618) (14100,0.003140014806) (14200,0.003146902681) (14300,0.003148960904) (14400,0.003086052136) (14500,0.002923160465) (14600,0.002886029566) (14700,0.002832951024) (14800,0.002871421911) (14900,0.002923196834) (15000,0.002983538806) (15100,0.002965077991) (15200,0.002943037078) (15300,0.002900162386) (15400,0.002982381172) (15500,0.002974730264) (15600,0.003022471466) (15700,0.003095164592) (15800,0.003028576076) (15900,0.002966669295) (16000,0.002937448793) (16100,0.0029056445) (16200,0.002844161121) (16300,0.002889354178) (16400,0.002839536243) (16500,0.002862819564) (16600,0.002850814257) (16700,0.002806969266) (16800,0.002984044305) (16900,0.002920798445) (17000,0.002905176394) (17100,0.002926271688) (17200,0.002927360707) (17300,0.002837961423) (17400,0.002867672918) (17500,0.002911573858) (17600,0.002797341044) (17700,0.002825274528) (17800,0.002790853381) (17900,0.002785448846) (18000,0.002805222617) (18100,0.002811575937) (18200,0.00283718782) (18300,0.002972577699) (18400,0.003077809955) (18500,0.002991317236) (18600,0.003104923642) (18700,0.003075550986) (18800,0.002988074906) (18900,0.003007189278) (19000,0.002994722687) (19100,0.003007852379) (19200,0.002979368949) (19300,0.002911907714) (19400,0.002793842345) (19500,0.002800585446) (19600,0.00276287063) (19700,0.002747606207) (19800,0.00276620402) (19900,0.002767429082) (20000,0.002843554039) (20100,0.002835679986) (20200,0.002835698146) (20300,0.002865415276) (20400,0.002893199236) (20500,0.002865427476) (20600,0.002828524075) (20700,0.00287006232) (20800,0.002891393774) (20900,0.002912164829) (21000,0.002782812039) (21100,0.002801421098) (21200,0.002877703914) (21300,0.002826442756) (21400,0.002802848886) (21500,0.002812177688) (21600,0.00285598035) (21700,0.002786324895) (21800,0.002861819463) (21900,0.002831570478) (22000,0.002880449663) (22100,0.002890049038) (22200,0.002861254266) (22300,0.002828640328) (22400,0.002788576088) (22500,0.002796354657) (22600,0.002748869359) (22700,0.002772162645) (22800,0.002654852602) (22900,0.002717841743) (23000,0.00283385308) (23100,0.002779365028) (23200,0.00276473898) (23300,0.002748738532) (23400,0.002769569983) (23500,0.002791135013) (23600,0.002858895832) (23700,0.002903285646) (23800,0.002893184801)};
\addplot[pubRose, opacity=0.18, line width=0.35pt] coordinates {(100,0.005515711382) (200,0.004240558832) (300,0.003759652143) (400,0.003543120518) (500,0.00357218422) (600,0.003545976399) (700,0.003657121344) (800,0.003545637155) (900,0.003757631639) (1000,0.00374388271) (1100,0.003657082887) (1200,0.003719234583) (1300,0.003854208137) (1400,0.003973349952) (1500,0.00391999674) (1600,0.003875651234) (1700,0.003840279696) (1800,0.003855088819) (1900,0.003550300887) (2000,0.003471095162) (2100,0.003301598993) (2200,0.003187056119) (2300,0.00295953434) (2400,0.00284319839) (2500,0.002768763038) (2600,0.002640053839) (2700,0.00243136494) (2800,0.002344882698) (2900,0.002339008218) (3000,0.002279719873) (3100,0.002236588812) (3200,0.002184980782) (3300,0.002209943649) (3400,0.002076921239) (3500,0.00198539407) (3600,0.00199090112) (3700,0.00195170677) (3800,0.001896148792) (3900,0.001804032945) (4000,0.00178645493) (4100,0.001746316045) (4200,0.001771486225) (4300,0.001744855358) (4400,0.00177982219) (4500,0.001796059858) (4600,0.001775185927) (4700,0.001777630311) (4800,0.00181033652) (4900,0.001923232758) (5000,0.001876179373) (5100,0.00180717347) (5200,0.001760701276) (5300,0.001755422063) (5400,0.001759802538) (5500,0.001792065566) (5600,0.001856312028) (5700,0.001921195758) (5800,0.00195669675) (5900,0.001936722815) (6000,0.001956135686) (6100,0.002022708836) (6200,0.002066278446) (6300,0.002081607119) (6400,0.00208575374) (6500,0.00204881489) (6600,0.002049648424) (6700,0.002048629278) (6800,0.001994096185) (6900,0.001969446009) (7000,0.001957042958) (7100,0.001949938131) (7200,0.001892696426) (7300,0.001882662706) (7400,0.001867391623) (7500,0.001883147447) (7600,0.001849123021) (7700,0.001815478527) (7800,0.001838558959) (7900,0.001873278501) (8000,0.001891703554) (8100,0.001842707186) (8200,0.001862400956) (8300,0.001898538461) (8400,0.001925300062) (8500,0.001935649221) (8600,0.001902708737) (8700,0.001910754223) (8800,0.001893157733) (8900,0.001876296417) (9000,0.00187183019) (9100,0.001901929092) (9200,0.001904636063) (9300,0.001864583849) (9400,0.001781920751) (9500,0.001759982621) (9600,0.001836335892) (9700,0.001786482136) (9800,0.001797229692) (9900,0.001761637221) (10000,0.001821197884) (10100,0.001794045849) (10200,0.001809088525) (10300,0.001836373843) (10400,0.001868502796) (10500,0.001889439963) (10600,0.00185049877) (10700,0.001924554014) (10800,0.001961406053) (10900,0.001999241055) (11000,0.001983687363) (11100,0.001996769349) (11200,0.00196293483) (11300,0.001956935902) (11400,0.001985844213) (11500,0.001993454562) (11600,0.00196621879) (11700,0.001947185374) (11800,0.001942784246) (11900,0.001959796622) (12000,0.001906740374) (12100,0.001933720394) (12200,0.001951114065) (12300,0.00191219612) (12400,0.001940749923) (12500,0.00199156733) (12600,0.002050957608) (12700,0.002031155513) (12800,0.00202976244) (12900,0.001966095867) (13000,0.002037578681) (13100,0.002028096258) (13200,0.002010731935) (13300,0.002056301665) (13400,0.00196092769) (13500,0.001955776149) (13600,0.001922072808) (13700,0.001900824788) (13800,0.001837929105) (13900,0.001869139553) (14000,0.00187848924) (14100,0.001878459333) (14200,0.001865368662) (14300,0.001804385986) (14400,0.001870809868) (14500,0.001801697956) (14600,0.00185253456) (14700,0.001871650445) (14800,0.00194856195) (14900,0.001929611911) (15000,0.001853975502) (15100,0.001823285501) (15200,0.001832803921) (15300,0.001854738954) (15400,0.001854459301) (15500,0.001855542662) (15600,0.001759986766) (15700,0.001815729099) (15800,0.001729249465) (15900,0.00170588044) (16000,0.001694537455) (16100,0.001720398245) (16200,0.001775440155) (16300,0.001775249245) (16400,0.001760657935) (16500,0.001718353364) (16600,0.001736740279) (16700,0.001725166664) (16800,0.001739987824) (16900,0.001729697292) (17000,0.001711506711) (17100,0.001720420574) (17200,0.001674337126) (17300,0.001729636663) (17400,0.001717514684) (17500,0.001755351364) (17600,0.001777579694) (17700,0.001718252595) (17800,0.001705562463) (17900,0.001695303782) (18000,0.001717868098) (18100,0.001669920562) (18200,0.001686573145) (18300,0.001628187927) (18400,0.00165417348) (18500,0.001655571279) (18600,0.001652382466) (18700,0.001674427604) (18800,0.00169157245) (18900,0.001717810694) (19000,0.001731369819) (19100,0.001744775998) (19200,0.001775729319) (19300,0.001763407956) (19400,0.001700492215) (19500,0.001670493488) (19600,0.001622525405) (19700,0.001590096962) (19800,0.001652453886) (19900,0.001612167479) (20000,0.001570705883) (20100,0.001563267724) (20200,0.001489720878) (20300,0.001499040041) (20400,0.001537261973) (20500,0.001590537792) (20600,0.0015933302) (20700,0.001642185077) (20800,0.001569031959) (20900,0.00156663378) (21000,0.001619417674) (21100,0.001628375461) (21200,0.001657695312) (21300,0.001663548837) (21400,0.0016661425) (21500,0.001663498208) (21600,0.001750792679) (21700,0.001734266384) (21800,0.001753126259) (21900,0.001791298727) (22000,0.001758325403) (22100,0.001853127813) (22200,0.001835403312) (22300,0.001815690729) (22400,0.001850150805) (22500,0.001804647956) (22600,0.001702535816) (22700,0.001707934064) (22800,0.001712815848) (22900,0.00171144862) (23000,0.001707380475) (23100,0.001623340347) (23200,0.001653938682) (23300,0.001659821544) (23400,0.001599800691) (23500,0.001643190102) (23600,0.001689715998) (23700,0.001662629296) (23800,0.001648308535)};
\addplot[pubRose, opacity=0.18, line width=0.35pt] coordinates {(100,0.006155986339) (200,0.006481280085) (300,0.006106218323) (400,0.005781570682) (500,0.005638517253) (600,0.005485369358) (700,0.005931606383) (800,0.005911072134) (900,0.006091772837) (1000,0.00626931428) (1100,0.006173766917) (1200,0.006288023433) (1300,0.006233793637) (1400,0.006460937904) (1500,0.006593037443) (1600,0.006690932764) (1700,0.006688850047) (1800,0.006715793302) (1900,0.006421265099) (2000,0.006241432764) (2100,0.006151373312) (2200,0.005838055722) (2300,0.005829371465) (2400,0.00548146423) (2500,0.005231952132) (2600,0.005012156861) (2700,0.004585290281) (2800,0.004328864114) (2900,0.004163669143) (3000,0.003800297854) (3100,0.003588086902) (3200,0.003413303872) (3300,0.003264871216) (3400,0.003147243243) (3500,0.002992368792) (3600,0.002909775404) (3700,0.002759061707) (3800,0.002756441338) (3900,0.002685587713) (4000,0.002725194162) (4100,0.002742103161) (4200,0.002695132699) (4300,0.00254872113) (4400,0.00249409345) (4500,0.002474333718) (4600,0.002442677179) (4700,0.002397759724) (4800,0.002341924934) (4900,0.002458070731) (5000,0.002400541911) (5100,0.002412046236) (5200,0.002457497502) (5300,0.002518455894) (5400,0.002574801538) (5500,0.002632961981) (5600,0.002663878072) (5700,0.002676336281) (5800,0.002669318928) (5900,0.002553583728) (6000,0.002666265517) (6100,0.002628505835) (6200,0.0025395673) (6300,0.00256902813) (6400,0.002471114276) (6500,0.002378413547) (6600,0.002279780409) (6700,0.002216167282) (6800,0.002142254892) (6900,0.00217829654) (7000,0.002116912929) (7100,0.002140440978) (7200,0.002192507521) (7300,0.002145552146) (7400,0.002198851504) (7500,0.002211521706) (7600,0.002282054909) (7700,0.002313057869) (7800,0.00242533125) (7900,0.002405286673) (8000,0.002363688359) (8100,0.002355459495) (8200,0.002372804168) (8300,0.002391851693) (8400,0.002416657051) (8500,0.002469102223) (8600,0.002505620988) (8700,0.002568377694) (8800,0.002506698971) (8900,0.002592611196) (9000,0.002567784698) (9100,0.002636325499) (9200,0.002631360805) (9300,0.002596102003) (9400,0.002557811048) (9500,0.00252254419) (9600,0.002524289419) (9700,0.002467667754) (9800,0.002504663309) (9900,0.002411100268) (10000,0.002440646198) (10100,0.002391164377) (10200,0.002365662786) (10300,0.002517333184) (10400,0.002534164162) (10500,0.002598403208) (10600,0.002616306627) (10700,0.002594528813) (10800,0.002535503986) (10900,0.002557111788) (11000,0.002571662655) (11100,0.002557869861) (11200,0.002590031549) (11300,0.002478134306) (11400,0.002509894874) (11500,0.00254855575) (11600,0.002445986017) (11700,0.002452431468) (11800,0.002502340369) (11900,0.002570129687) (12000,0.002562813892) (12100,0.002640064724) (12200,0.002662799193) (12300,0.002730162523) (12400,0.002775621542) (12500,0.002696408669) (12600,0.002731391182) (12700,0.002772513987) (12800,0.002814298728) (12900,0.002706320724) (13000,0.002714656037) (13100,0.002649865812) (13200,0.002549673221) (13300,0.002513112733) (13400,0.002444897359) (13500,0.002453488437) (13600,0.002526134439) (13700,0.002535786573) (13800,0.002587583289) (13900,0.002672250988) (14000,0.002643468534) (14100,0.002699421882) (14200,0.002796350978) (14300,0.002781363158) (14400,0.002842070209) (14500,0.00285199366) (14600,0.002731077489) (14700,0.00275978602) (14800,0.002763361554) (14900,0.002684144559) (15000,0.002833339875) (15100,0.002799360454) (15200,0.002771631442) (15300,0.002805771912) (15400,0.002843639674) (15500,0.002804669412) (15600,0.002874684054) (15700,0.002858829335) (15800,0.00278251532) (15900,0.002855801722) (16000,0.002705907938) (16100,0.002677379339) (16200,0.002709379629) (16300,0.002716989513) (16400,0.002659006626) (16500,0.002723802975) (16600,0.002879594895) (16700,0.002905330737) (16800,0.002943455218) (16900,0.00294129902) (17000,0.002996612294) (17100,0.003039191966) (17200,0.00295127274) (17300,0.002903430257) (17400,0.002961899387) (17500,0.002972272737) (17600,0.002864797809) (17700,0.002770890924) (17800,0.002735298849) (17900,0.002710490464) (18000,0.002627667901) (18100,0.002608823404) (18200,0.002639691485) (18300,0.002615606552) (18400,0.002581297839) (18500,0.002552218153) (18600,0.002521537966) (18700,0.002590125264) (18800,0.002499993914) (18900,0.002614317904) (19000,0.002614029846) (19100,0.002674578945) (19200,0.002667895169) (19300,0.002628296288) (19400,0.002579077194) (19500,0.002575307363) (19600,0.002600034466) (19700,0.002596113738) (19800,0.0026580523) (19900,0.002505770675) (20000,0.002530614589) (20100,0.002450106665) (20200,0.002478335984) (20300,0.002579386183) (20400,0.002606553654) (20500,0.002608907851) (20600,0.002600647463) (20700,0.002617345098) (20800,0.002574074874) (20900,0.002573748073) (21000,0.00258228723) (21100,0.002598631917) (21200,0.002527198149) (21300,0.002492139605) (21400,0.002541086217) (21500,0.002556329779) (21600,0.002574864239) (21700,0.002563866763) (21800,0.002564563137) (21900,0.002720537735) (22000,0.002730589034) (22100,0.002763395151) (22200,0.002813977562) (22300,0.002855080203) (22400,0.002937697084) (22500,0.002867879416) (22600,0.002828182862) (22700,0.002828090289) (22800,0.002878433862) (22900,0.002786413953) (23000,0.002820977615) (23100,0.00284855892) (23200,0.002888865327) (23300,0.002829023637) (23400,0.002753079566) (23500,0.002781443717) (23600,0.002794974763) (23700,0.002806560788) (23800,0.002811594773)};
\addplot[pubRose, opacity=0.18, line width=0.35pt] coordinates {(100,0.005548516754) (200,0.005943406606) (300,0.00570687574) (400,0.005618137657) (500,0.00555219762) (600,0.005810107881) (700,0.005894771844) (800,0.006064803514) (900,0.005949554654) (1000,0.005770482495) (1100,0.00564202047) (1200,0.00550773521) (1300,0.005553796655) (1400,0.005473669199) (1500,0.005227943137) (1600,0.004996317625) (1700,0.004963238304) (1800,0.004677473288) (1900,0.004547151574) (2000,0.004482849059) (2100,0.004442430544) (2200,0.004285188229) (2300,0.00405368621) (2400,0.003882983001) (2500,0.00382875856) (2600,0.003721294156) (2700,0.003446137649) (2800,0.003264829051) (2900,0.003116615373) (3000,0.003060620371) (3100,0.002906735195) (3200,0.002878627041) (3300,0.002803365001) (3400,0.002854866278) (3500,0.00279983267) (3600,0.002640755114) (3700,0.002612723538) (3800,0.002548022976) (3900,0.002506732906) (4000,0.002433743409) (4100,0.00238850168) (4200,0.002323626005) (4300,0.002287692414) (4400,0.002160493354) (4500,0.002225087176) (4600,0.002239095455) (4700,0.002129435854) (4800,0.00207201076) (4900,0.002142839774) (5000,0.002188784827) (5100,0.002180327778) (5200,0.002166052139) (5300,0.002139442344) (5400,0.002140621585) (5500,0.002099303785) (5600,0.002116031479) (5700,0.002079436171) (5800,0.002062464948) (5900,0.002008978394) (6000,0.001917741064) (6100,0.001888421888) (6200,0.001811287261) (6300,0.001804861671) (6400,0.001790413551) (6500,0.001784540468) (6600,0.001720476069) (6700,0.001735520456) (6800,0.001790051488) (6900,0.001769773918) (7000,0.001791799755) (7100,0.00181176319) (7200,0.00179152661) (7300,0.001813243865) (7400,0.001781936432) (7500,0.001849929802) (7600,0.001836538373) (7700,0.001878364442) (7800,0.001888104726) (7900,0.001840047247) (8000,0.001792871405) (8100,0.001822481875) (8200,0.001872563397) (8300,0.001781608036) (8400,0.001789875084) (8500,0.001667295804) (8600,0.001684936485) (8700,0.001651937689) (8800,0.001658773026) (8900,0.001775038056) (9000,0.001784348546) (9100,0.001733446529) (9200,0.001761042816) (9300,0.001810611517) (9400,0.001769943826) (9500,0.001836524729) (9600,0.001820085978) (9700,0.001813697303) (9800,0.001777489507) (9900,0.001690039749) (10000,0.001736458507) (10100,0.001759308728) (10200,0.001691934455) (10300,0.001686604961) (10400,0.001789089933) (10500,0.001735760551) (10600,0.001754999976) (10700,0.001756127912) (10800,0.001769969892) (10900,0.001802588755) (11000,0.001768645889) (11100,0.001788358984) (11200,0.001870151109) (11300,0.001934163133) (11400,0.001855758461) (11500,0.001844667422) (11600,0.001903611876) (11700,0.00186768563) (11800,0.001893239317) (11900,0.00186184427) (12000,0.001830438082) (12100,0.00178702079) (12200,0.001741747966) (12300,0.001667912223) (12400,0.001677053247) (12500,0.001717842568) (12600,0.001656375511) (12700,0.00168108152) (12800,0.001693241659) (12900,0.00172336949) (13000,0.001782485482) (13100,0.001783692359) (13200,0.001790590037) (13300,0.001801642834) (13400,0.001815758634) (13500,0.00181764398) (13600,0.001817319333) (13700,0.001813543751) (13800,0.001846542535) (13900,0.001893893967) (14000,0.001830585499) (14100,0.001881470799) (14200,0.001910447876) (14300,0.001884279621) (14400,0.001838323439) (14500,0.001865465066) (14600,0.001854918082) (14700,0.001843966835) (14800,0.001791114197) (14900,0.001741173503) (15000,0.001801406138) (15100,0.001757990383) (15200,0.001769284136) (15300,0.001790077204) (15400,0.001886347623) (15500,0.001857526251) (15600,0.001890988043) (15700,0.001980445406) (15800,0.00196427945) (15900,0.001969419525) (16000,0.001946950157) (16100,0.00201501851) (16200,0.00196411896) (16300,0.002025242348) (16400,0.001915902051) (16500,0.001977152773) (16600,0.001982673327) (16700,0.001967543294) (16800,0.00196677983) (16900,0.001937673846) (17000,0.001990240905) (17100,0.001980087091) (17200,0.001981224399) (17300,0.001998071163) (17400,0.002034160588) (17500,0.002034777217) (17600,0.00203812886) (17700,0.002036436251) (17800,0.002004123677) (17900,0.002024534543) (18000,0.00201428748) (18100,0.001999627135) (18200,0.002022901422) (18300,0.002060790022) (18400,0.002064301365) (18500,0.002008673677) (18600,0.001958850829) (18700,0.00191257227) (18800,0.001956838847) (18900,0.001966413367) (19000,0.001922326337) (19100,0.001861059177) (19200,0.001839788863) (19300,0.001823240751) (19400,0.001853024785) (19500,0.001844837132) (19600,0.001847056404) (19700,0.001878458599) (19800,0.001876082656) (19900,0.001896998205) (20000,0.001873261645) (20100,0.001928523707) (20200,0.001912582712) (20300,0.001867289329) (20400,0.001842250198) (20500,0.001850962394) (20600,0.00183248996) (20700,0.001798308978) (20800,0.001801584871) (20900,0.001789347106) (21000,0.001794099715) (21100,0.001730643841) (21200,0.001727740176) (21300,0.001726290828) (21400,0.001742324152) (21500,0.001763869904) (21600,0.001779385994) (21700,0.001760890346) (21800,0.001753707905) (21900,0.001768748602) (22000,0.00191527803) (22100,0.001975825289) (22200,0.002087977389) (22300,0.002081566909) (22400,0.002113829763) (22500,0.002086878126) (22600,0.002166605345) (22700,0.002193304617) (22800,0.002166130138) (22900,0.002121368819) (23000,0.001996704156) (23100,0.002024384926) (23200,0.001900446031) (23300,0.00187508727) (23400,0.001884066057) (23500,0.001930653327) (23600,0.001831937465) (23700,0.001795732952) (23800,0.001872853446)};
\addplot[pubRose, opacity=0.18, line width=0.35pt] coordinates {(100,0.002143648453) (200,0.002131344401) (300,0.002262726038) (400,0.002624231332) (500,0.002438338008) (600,0.002462128876) (700,0.002494561985) (800,0.00242217473) (900,0.002428359792) (1000,0.002432769025) (1100,0.002440383798) (1200,0.002430228167) (1300,0.002372866811) (1400,0.002261663566) (1500,0.002294435666) (1600,0.00227545806) (1700,0.002232913964) (1800,0.002278738876) (1900,0.002291841537) (2000,0.002341358957) (2100,0.002522577194) (2200,0.002574609441) (2300,0.002600668697) (2400,0.002576550678) (2500,0.002581081563) (2600,0.002566710697) (2700,0.002563205501) (2800,0.002562223375) (2900,0.002517095068) (3000,0.002441595285) (3100,0.002180504764) (3200,0.002157616301) (3300,0.00219316805) (3400,0.002148184704) (3500,0.002166449104) (3600,0.002149691165) (3700,0.002158696449) (3800,0.002125720878) (3900,0.002075307874) (4000,0.002098046883) (4100,0.002143238147) (4200,0.002040153602) (4300,0.002011847938) (4400,0.002010910236) (4500,0.001983932336) (4600,0.002017042693) (4700,0.001949725079) (4800,0.001927442488) (4900,0.001947256096) (5000,0.001907765551) (5100,0.001923512807) (5200,0.002001159487) (5300,0.001920684928) (5400,0.001956345688) (5500,0.001983313693) (5600,0.001909702958) (5700,0.001965067419) (5800,0.001963854185) (5900,0.002019033511) (6000,0.001955329999) (6100,0.001956201973) (6200,0.001952194585) (6300,0.002107995446) (6400,0.002092921711) (6500,0.002068224829) (6600,0.002087181096) (6700,0.002070863859) (6800,0.002206375229) (6900,0.00210629924) (7000,0.002149511036) (7100,0.002154403343) (7200,0.002148752729) (7300,0.002000347013) (7400,0.001985813398) (7500,0.001970387774) (7600,0.002004726441) (7700,0.001987400826) (7800,0.001928079745) (7900,0.001934718736) (8000,0.001927292324) (8100,0.001909128879) (8200,0.001893492462) (8300,0.001989645115) (8400,0.001928380178) (8500,0.001971897867) (8600,0.001907834795) (8700,0.001928463182) (8800,0.001889504981) (8900,0.001999260148) (9000,0.001964169706) (9100,0.001963190094) (9200,0.001942393801) (9300,0.00188138762) (9400,0.001958095562) (9500,0.001885256637) (9600,0.001928507839) (9700,0.001897153968) (9800,0.001893727772) (9900,0.001880452351) (10000,0.001941124513) (10100,0.001903000986) (10200,0.001880934346) (10300,0.001846255513) (10400,0.001802095782) (10500,0.001802490128) (10600,0.001792635827) (10700,0.001851602853) (10800,0.001887172461) (10900,0.001856710482) (11000,0.001848318987) (11100,0.001903108088) (11200,0.001928098465) (11300,0.001935391349) (11400,0.00194324546) (11500,0.001973797148) (11600,0.002065337915) (11700,0.002004782809) (11800,0.001926608512) (11900,0.001903121802) (12000,0.001896707923) (12100,0.001867737889) (12200,0.001878776716) (12300,0.001924504281) (12400,0.001907908078) (12500,0.001861119282) (12600,0.001740771125) (12700,0.001728317572) (12800,0.001722781407) (12900,0.00168637376) (13000,0.001723876351) (13100,0.001730584458) (13200,0.00169800811) (13300,0.001675210416) (13400,0.001724734146) (13500,0.001729197975) (13600,0.001795528585) (13700,0.001877764368) (13800,0.001932749932) (13900,0.001973588276) (14000,0.001890040841) (14100,0.00189120787) (14200,0.002017413464) (14300,0.001996675925) (14400,0.001942115114) (14500,0.002059950342) (14600,0.00201449726) (14700,0.001935773401) (14800,0.001971373532) (14900,0.002004402201) (15000,0.002022837242) (15100,0.001997428818) (15200,0.001957137592) (15300,0.001960184844) (15400,0.002018978633) (15500,0.001926351641) (15600,0.001905350958) (15700,0.001922547177) (15800,0.001813881507) (15900,0.001802644285) (16000,0.001934624219) (16100,0.001959230273) (16200,0.001910964958) (16300,0.001930509706) (16400,0.001948074892) (16500,0.002058135264) (16600,0.002058708679) (16700,0.002029738959) (16800,0.002078970626) (16900,0.002047707106) (17000,0.001909109869) (17100,0.001902017696) (17200,0.001957902184) (17300,0.001919289108) (17400,0.001842819899) (17500,0.001705346233) (17600,0.001714686013) (17700,0.001715500106) (17800,0.001677882567) (17900,0.001651449583) (18000,0.001641822059) (18100,0.001673167234) (18200,0.001567283692) (18300,0.001603871991) (18400,0.001685838215) (18500,0.001789435162) (18600,0.001813632087) (18700,0.001795422565) (18800,0.001804089104) (18900,0.001817969955) (19000,0.001951502147) (19100,0.001968609518) (19200,0.0020497786) (19300,0.002057318529) (19400,0.00199507362) (19500,0.001895179891) (19600,0.001851337927) (19700,0.001969272958) (19800,0.001949084993) (19900,0.001988554327) (20000,0.001877215598) (20100,0.001875645621) (20200,0.001863500453) (20300,0.001842852263) (20400,0.001814104302) (20500,0.001839271432) (20600,0.001902423322) (20700,0.00183378082) (20800,0.001897395402) (20900,0.001991762174) (21000,0.001998136018) (21100,0.001929873449) (21200,0.001908287546) (21300,0.001882622309) (21400,0.001932450698) (21500,0.002067826281) (21600,0.002045263455) (21700,0.002009041782) (21800,0.001982438576) (21900,0.001874861831) (22000,0.00187848435) (22100,0.001936634246) (22200,0.0019621199) (22300,0.001984618767) (22400,0.002038417221) (22500,0.001965799602) (22600,0.0019585034) (22700,0.001960605814) (22800,0.001948172075) (22900,0.001952213247) (23000,0.001970300183) (23100,0.001927178609) (23200,0.001913288538) (23300,0.001967772527) (23400,0.001935457753) (23500,0.001906321256) (23600,0.001929046214) (23700,0.001949278649) (23800,0.001963574672)};
\addplot[pubRose, opacity=0.18, line width=0.35pt] coordinates {(100,0.00666039763) (200,0.005752051482) (300,0.005384561451) (400,0.005500511965) (500,0.005086076912) (600,0.004797515382) (700,0.004650275888) (800,0.004684130254) (900,0.004690886423) (1000,0.004616710125) (1100,0.004383555753) (1200,0.004207742075) (1300,0.004081370286) (1400,0.003797126701) (1500,0.003811935661) (1600,0.003816037788) (1700,0.003663112363) (1800,0.003594428115) (1900,0.003457494965) (2000,0.003492983663) (2100,0.003355931723) (2200,0.003442181507) (2300,0.003404810815) (2400,0.003363645426) (2500,0.003350394359) (2600,0.003346549324) (2700,0.003397374926) (2800,0.003177004587) (2900,0.00311116837) (3000,0.002975889645) (3100,0.002972606244) (3200,0.002813818515) (3300,0.002753159963) (3400,0.002723469236) (3500,0.002642791998) (3600,0.002566598542) (3700,0.002675011405) (3800,0.002692026412) (3900,0.002672277484) (4000,0.002582890447) (4100,0.002563713561) (4200,0.00272086272) (4300,0.002730751364) (4400,0.002767646569) (4500,0.002726981649) (4600,0.002695297543) (4700,0.002602483216) (4800,0.002658940665) (4900,0.002653890196) (5000,0.002683541062) (5100,0.002663027006) (5200,0.002457804885) (5300,0.002378425561) (5400,0.002289186697) (5500,0.002315444266) (5600,0.002245378203) (5700,0.002171445277) (5800,0.00213377272) (5900,0.002140290418) (6000,0.002123051172) (6100,0.002140294027) (6200,0.002147789381) (6300,0.002187495201) (6400,0.002227875323) (6500,0.002222200332) (6600,0.002270847373) (6700,0.002342936932) (6800,0.002329554944) (6900,0.002314949688) (7000,0.002353398362) (7100,0.002315923455) (7200,0.002498997352) (7300,0.002521002386) (7400,0.002558926865) (7500,0.002494065859) (7600,0.0025089717) (7700,0.002435978432) (7800,0.002491455362) (7900,0.002491493616) (8000,0.002556722471) (8100,0.002600375377) (8200,0.002422145591) (8300,0.002422068338) (8400,0.002425285499) (8500,0.002514720452) (8600,0.002567915409) (8700,0.002634809958) (8800,0.002683283063) (8900,0.002666993928) (9000,0.002538540214) (9100,0.00250745751) (9200,0.002592837089) (9300,0.002603537822) (9400,0.002507087588) (9500,0.002520722733) (9600,0.002584188012) (9700,0.002500739391) (9800,0.002432973008) (9900,0.002373447083) (10000,0.002460265625) (10100,0.002481938479) (10200,0.002426197007) (10300,0.00239041131) (10400,0.002447575098) (10500,0.002407736052) (10600,0.00232493747) (10700,0.0023517936) (10800,0.002397845662) (10900,0.002481808024) (11000,0.002431841241) (11100,0.002410928858) (11200,0.002358465584) (11300,0.002425665164) (11400,0.002434085624) (11500,0.002399935771) (11600,0.002317363862) (11700,0.002301531425) (11800,0.002214164496) (11900,0.002134524775) (12000,0.002106677252) (12100,0.002089623851) (12200,0.002120359114) (12300,0.002064809541) (12400,0.002114125818) (12500,0.002141082229) (12600,0.002205429273) (12700,0.002347407583) (12800,0.002368100686) (12900,0.002433577715) (13000,0.002439420344) (13100,0.00245680646) (13200,0.002450607484) (13300,0.002454560879) (13400,0.002388410759) (13500,0.00241666066) (13600,0.002382580726) (13700,0.002235981985) (13800,0.002213920956) (13900,0.002280404675) (14000,0.002294931561) (14100,0.002311983751) (14200,0.002395960176) (14300,0.002421975695) (14400,0.002483449038) (14500,0.002434088336) (14600,0.002429704461) (14700,0.002424793039) (14800,0.002441984508) (14900,0.002423412772) (15000,0.00242995075) (15100,0.002478961041) (15200,0.002502544667) (15300,0.002461309568) (15400,0.002374225482) (15500,0.002352005115) (15600,0.002346793388) (15700,0.00241598686) (15800,0.002472159138) (15900,0.002544372261) (16000,0.002528897172) (16100,0.002426462981) (16200,0.002299612446) (16300,0.002332192834) (16400,0.002335237537) (16500,0.002561555617) (16600,0.002544473344) (16700,0.002471166104) (16800,0.002386949514) (16900,0.002314376994) (17000,0.002311853692) (17100,0.002428480238) (17200,0.002543431846) (17300,0.002538451669) (17400,0.002714124485) (17500,0.002514252556) (17600,0.002564500063) (17700,0.002629779768) (17800,0.002648799657) (17900,0.002678137785) (18000,0.002764627687) (18100,0.002779004094) (18200,0.002676057757) (18300,0.00269082404) (18400,0.002550209302) (18500,0.002564761799) (18600,0.002670716343) (18700,0.002650352183) (18800,0.002640590083) (18900,0.002662264404) (19000,0.002715406509) (19100,0.002612485352) (19200,0.002695768443) (19300,0.002675811341) (19400,0.002658580104) (19500,0.002672775788) (19600,0.002604269818) (19700,0.002661690745) (19800,0.002663633134) (19900,0.002698830282) (20000,0.002657367312) (20100,0.002659171447) (20200,0.002680116636) (20300,0.002718719048) (20400,0.002748082695) (20500,0.002814889094) (20600,0.002759728441) (20700,0.002682407177) (20800,0.002710284735) (20900,0.002632369823) (21000,0.002532334323) (21100,0.002557345689) (21200,0.002475531888) (21300,0.002424280741) (21400,0.002429342526) (21500,0.002458053571) (21600,0.002471993025) (21700,0.002571719093) (21800,0.002566688322) (21900,0.002530779131) (22000,0.00252462544) (22100,0.002555957274) (22200,0.00262235906) (22300,0.002683010534) (22400,0.00264617845) (22500,0.002545891213) (22600,0.002617263841) (22700,0.002643790119) (22800,0.002654827735) (22900,0.002702773246) (23000,0.002728786948) (23100,0.002763694851) (23200,0.002727091848) (23300,0.002696559625) (23400,0.002790568769) (23500,0.002919582743) (23600,0.002886285633) (23700,0.002841021866) (23800,0.002784081199)};
\addplot[pubRose, opacity=0.18, line width=0.35pt] coordinates {(100,0.01545402966) (200,0.01352667529) (300,0.01461769765) (400,0.01363109937) (500,0.01288528759) (600,0.01250471656) (700,0.01172779261) (800,0.01118671702) (900,0.01103030926) (1000,0.01072791214) (1100,0.00979716084) (1200,0.009530185955) (1300,0.008731272491) (1400,0.008293096907) (1500,0.0080155259) (1600,0.007538795192) (1700,0.007417109981) (1800,0.007264954271) (1900,0.006845494499) (2000,0.006581692118) (2100,0.006594323879) (2200,0.006020236621) (2300,0.005665690545) (2400,0.005400080886) (2500,0.004995586816) (2600,0.004714833316) (2700,0.004553101421) (2800,0.004199916427) (2900,0.003992625303) (3000,0.003741519293) (3100,0.003423063527) (3200,0.003450543364) (3300,0.003213944635) (3400,0.003168095578) (3500,0.00314759044) (3600,0.003142317315) (3700,0.002975279419) (3800,0.002980024577) (3900,0.002814958035) (4000,0.002809968358) (4100,0.002770752576) (4200,0.002608699817) (4300,0.002519044862) (4400,0.002423545835) (4500,0.002290124702) (4600,0.002194621018) (4700,0.002175418753) (4800,0.002205853933) (4900,0.002272246056) (5000,0.00223878643) (5100,0.00215979194) (5200,0.002235227043) (5300,0.002247244574) (5400,0.002260519459) (5500,0.002347975865) (5600,0.002372891328) (5700,0.002417253924) (5800,0.002408845874) (5900,0.00238143591) (6000,0.002304465533) (6100,0.002349483885) (6200,0.002275312168) (6300,0.002269447118) (6400,0.002224313363) (6500,0.002193227794) (6600,0.002158620546) (6700,0.00209681337) (6800,0.002042283828) (6900,0.002113474195) (7000,0.002190757729) (7100,0.002205060842) (7200,0.002205677226) (7300,0.002330684674) (7400,0.002289643337) (7500,0.002250837639) (7600,0.002227920445) (7700,0.002228620171) (7800,0.002284715266) (7900,0.002168120421) (8000,0.002122506173) (8100,0.001993595017) (8200,0.001988486352) (8300,0.001879298768) (8400,0.001919571788) (8500,0.001958523074) (8600,0.001958483376) (8700,0.001873835048) (8800,0.00178179607) (8900,0.001827506744) (9000,0.001831443748) (9100,0.001878237142) (9200,0.001859633694) (9300,0.001824411203) (9400,0.001817111555) (9500,0.001819878351) (9600,0.001854344062) (9700,0.001901672315) (9800,0.001932794228) (9900,0.001893202425) (10000,0.001820469159) (10100,0.001865867293) (10200,0.001938653807) (10300,0.001942989731) (10400,0.001909076399) (10500,0.001883865544) (10600,0.001844566746) (10700,0.001845107123) (10800,0.001808398974) (10900,0.001856283017) (11000,0.001883327204) (11100,0.001844668586) (11200,0.001779759664) (11300,0.001809915935) (11400,0.00181942858) (11500,0.001856065437) (11600,0.001829067071) (11700,0.001839915395) (11800,0.001852832932) (11900,0.001794869883) (12000,0.001799733914) (12100,0.001779492002) (12200,0.001766398025) (12300,0.001736464759) (12400,0.001733918081) (12500,0.001723175438) (12600,0.001785280602) (12700,0.001770125004) (12800,0.001782854693) (12900,0.001771758834) (13000,0.001792429853) (13100,0.001847872895) (13200,0.001882627385) (13300,0.001829773339) (13400,0.001925990393) (13500,0.001928975864) (13600,0.001954760554) (13700,0.001938784041) (13800,0.001912473026) (13900,0.001954459341) (14000,0.001930472895) (14100,0.001845236483) (14200,0.001801336894) (14300,0.001830572111) (14400,0.001750634925) (14500,0.00175505752) (14600,0.001656462369) (14700,0.001660035318) (14800,0.001695188868) (14900,0.001657122944) (15000,0.001718852844) (15100,0.001775244193) (15200,0.001880160265) (15300,0.001886496122) (15400,0.002001804544) (15500,0.001973523584) (15600,0.002020870475) (15700,0.002005651721) (15800,0.001985126885) (15900,0.002008528973) (16000,0.001938506088) (16100,0.002045801526) (16200,0.002055593615) (16300,0.00212387616) (16400,0.002057092579) (16500,0.002043552697) (16600,0.002062479779) (16700,0.002135355014) (16800,0.002114597231) (16900,0.002131284459) (17000,0.002117373131) (17100,0.001961846568) (17200,0.001850703161) (17300,0.001852603711) (17400,0.001786165987) (17500,0.001734847529) (17600,0.001711481728) (17700,0.001678665343) (17800,0.001736851514) (17900,0.001718444447) (18000,0.00173503703) (18100,0.001770286588) (18200,0.001898072753) (18300,0.00184028704) (18400,0.001835628354) (18500,0.001900058531) (18600,0.001932509243) (18700,0.00194382011) (18800,0.001897176448) (18900,0.001902627514) (19000,0.002011665201) (19100,0.002041869366) (19200,0.001921927498) (19300,0.001888036006) (19400,0.001883716171) (19500,0.00185633644) (19600,0.00184023059) (19700,0.001914111921) (19800,0.001959676831) (19900,0.001979890768) (20000,0.001939059258) (20100,0.001905286987) (20200,0.002003934514) (20300,0.002036013489) (20400,0.002012489166) (20500,0.002020085009) (20600,0.002004738548) (20700,0.001946973475) (20800,0.001868016575) (20900,0.001788190042) (21000,0.001751695026) (21100,0.001777992968) (21200,0.001715856756) (21300,0.001742150914) (21400,0.001816466521) (21500,0.001827545709) (21600,0.001808996731) (21700,0.001794787613) (21800,0.001807361876) (21900,0.00188780243) (22000,0.00189231293) (22100,0.001932509115) (22200,0.001859925746) (22300,0.001805388299) (22400,0.001833701786) (22500,0.001862390467) (22600,0.001874686952) (22700,0.001868327556) (22800,0.001883425331) (22900,0.001839284936) (23000,0.001812085602) (23100,0.001796885883) (23200,0.001862303843) (23300,0.001872269565) (23400,0.001838441379) (23500,0.001771665574) (23600,0.00175900613) (23700,0.00174610574) (23800,0.001716657984)};
\addplot[pubRose, opacity=0.18, line width=0.35pt] coordinates {(100,0.005169120152) (200,0.005079503404) (300,0.005776608208) (400,0.005852050381) (500,0.005498896353) (600,0.005555149478) (700,0.005423495213) (800,0.00539612287) (900,0.005214801032) (1000,0.005385520752) (1100,0.005343822693) (1200,0.005233955686) (1300,0.004834386078) (1400,0.00456017626) (1500,0.004587554676) (1600,0.004417281528) (1700,0.004306426109) (1800,0.004043956776) (1900,0.003997676307) (2000,0.003734780615) (2100,0.003543641465) (2200,0.003398553631) (2300,0.003362178826) (2400,0.003306499613) (2500,0.003097389452) (2600,0.00293326776) (2700,0.002820897964) (2800,0.002798580867) (2900,0.002713406994) (3000,0.002520612907) (3100,0.00255185701) (3200,0.002525516064) (3300,0.002502669138) (3400,0.002502869884) (3500,0.002482616366) (3600,0.002458108123) (3700,0.002421608102) (3800,0.002370455011) (3900,0.002412207925) (4000,0.002366709721) (4100,0.002246990998) (4200,0.002302551328) (4300,0.002220529842) (4400,0.002145881078) (4500,0.002134349942) (4600,0.002067614801) (4700,0.002058467455) (4800,0.002085270814) (4900,0.001979053172) (5000,0.002056840772) (5100,0.002069100121) (5200,0.001981954114) (5300,0.001965617016) (5400,0.001954565395) (5500,0.001943231339) (5600,0.001926345029) (5700,0.001890176011) (5800,0.001940648921) (5900,0.001951473928) (6000,0.00185916815) (6100,0.001847346895) (6200,0.001814684365) (6300,0.001846901164) (6400,0.001817541698) (6500,0.001804316579) (6600,0.001955613052) (6700,0.001977988996) (6800,0.001932403364) (6900,0.00192076338) (7000,0.001880214701) (7100,0.001848837535) (7200,0.001874469547) (7300,0.001917776442) (7400,0.001911151945) (7500,0.0018962909) (7600,0.00181462029) (7700,0.001841486536) (7800,0.001777272869) (7900,0.001799156668) (8000,0.00191298374) (8100,0.001860471605) (8200,0.001862746244) (8300,0.001805911539) (8400,0.001885282528) (8500,0.001918044523) (8600,0.001849842223) (8700,0.001852393372) (8800,0.001923097472) (8900,0.001862876245) (9000,0.001813918806) (9100,0.001907004218) (9200,0.00190869316) (9300,0.001870916621) (9400,0.001808282931) (9500,0.001784573589) (9600,0.001791052916) (9700,0.001778481365) (9800,0.001821267023) (9900,0.001826038561) (10000,0.001875550323) (10100,0.001867828495) (10200,0.001846126199) (10300,0.001913328643) (10400,0.001944085979) (10500,0.001964426518) (10600,0.001945553347) (10700,0.001924152987) (10800,0.001827460818) (10900,0.001834643458) (11000,0.001762375992) (11100,0.001751226443) (11200,0.001733624598) (11300,0.001703884394) (11400,0.00179673786) (11500,0.001745259494) (11600,0.001814922655) (11700,0.00176366372) (11800,0.001851133467) (11900,0.001875965693) (12000,0.001884094672) (12100,0.00184134593) (12200,0.001871477789) (12300,0.001860635879) (12400,0.001732706581) (12500,0.001754946308) (12600,0.001743387617) (12700,0.001782597019) (12800,0.001730628079) (12900,0.001752178348) (13000,0.00174328374) (13100,0.00179350999) (13200,0.001845239953) (13300,0.001851085224) (13400,0.00182621876) (13500,0.001790504775) (13600,0.001811629639) (13700,0.001805982494) (13800,0.001805872901) (13900,0.001794686483) (14000,0.00179052871) (14100,0.00176325565) (14200,0.001700526162) (14300,0.001787328732) (14400,0.001832042297) (14500,0.001812986762) (14600,0.001777669264) (14700,0.001802150474) (14800,0.001764110278) (14900,0.001752550073) (15000,0.00174219606) (15100,0.001767692668) (15200,0.00177803674) (15300,0.001686352433) (15400,0.001647072437) (15500,0.001779032301) (15600,0.001786804176) (15700,0.001821819227) (15800,0.001852635993) (15900,0.00182898998) (16000,0.001868200197) (16100,0.001825005445) (16200,0.001797839219) (16300,0.001817074278) (16400,0.001890874724) (16500,0.001780476456) (16600,0.001701348461) (16700,0.001669717534) (16800,0.001630350726) (16900,0.001633542369) (17000,0.001604457456) (17100,0.001620068192) (17200,0.001640612003) (17300,0.001613789459) (17400,0.001590285485) (17500,0.00171674313) (17600,0.001799104828) (17700,0.001775109139) (17800,0.00177695551) (17900,0.001890349283) (18000,0.001942963735) (18100,0.001970057015) (18200,0.001941172557) (18300,0.001996379986) (18400,0.001973368134) (18500,0.001901536679) (18600,0.001896182506) (18700,0.001895241986) (18800,0.001895503676) (18900,0.001810645743) (19000,0.001753373665) (19100,0.001727287017) (19200,0.001755898644) (19300,0.001715996338) (19400,0.001746932627) (19500,0.00172203473) (19600,0.00171674228) (19700,0.001702772954) (19800,0.001724065968) (19900,0.001719292812) (20000,0.001727971714) (20100,0.001719564409) (20200,0.001707335224) (20300,0.00170816289) (20400,0.001650854026) (20500,0.001699859661) (20600,0.00166566819) (20700,0.001701717102) (20800,0.001682756026) (20900,0.001664294) (21000,0.001638386457) (21100,0.001615516387) (21200,0.001626984589) (21300,0.00162487505) (21400,0.001638262765) (21500,0.001649880828) (21600,0.001689517475) (21700,0.00168236855) (21800,0.00164799504) (21900,0.001656239491) (22000,0.001670265733) (22100,0.001696193009) (22200,0.001748658065) (22300,0.001735323644) (22400,0.001757913863) (22500,0.001681445551) (22600,0.001678339206) (22700,0.00166279556) (22800,0.001706667233) (22900,0.001697913872) (23000,0.001692896313) (23100,0.001719060773) (23200,0.001690628775) (23300,0.001670228364) (23400,0.001667855389) (23500,0.001755266835) (23600,0.001730944996) (23700,0.001752463123) (23800,0.001770378719)};
\addplot[pubRose, opacity=0.18, line width=0.35pt] coordinates {(100,0.004772347398) (200,0.003985709045) (300,0.003851987732) (400,0.003699350636) (500,0.003632029612) (600,0.003522947159) (700,0.00334083412) (800,0.003239663143) (900,0.003210720218) (1000,0.003228509682) (1100,0.003007044573) (1200,0.003008534131) (1300,0.002875755611) (1400,0.002811251115) (1500,0.00277666552) (1600,0.002817765181) (1700,0.002904353081) (1800,0.00295232737) (1900,0.002978876373) (2000,0.002881016792) (2100,0.00288813354) (2200,0.002846176364) (2300,0.00284708431) (2400,0.002917704079) (2500,0.00290832005) (2600,0.002885118569) (2700,0.002783072251) (2800,0.002740762057) (2900,0.002686155494) (3000,0.002647376666) (3100,0.002709832159) (3200,0.002691711579) (3300,0.002705453569) (3400,0.002715450129) (3500,0.002693656506) (3600,0.002672459395) (3700,0.002698570653) (3800,0.002725211089) (3900,0.002664356888) (4000,0.002709383518) (4100,0.002621068247) (4200,0.002586003533) (4300,0.002631322667) (4400,0.002491918881) (4500,0.002446656418) (4600,0.002337014873) (4700,0.002345935453) (4800,0.00225532155) (4900,0.00223855318) (5000,0.002161525085) (5100,0.002118481998) (5200,0.0020970301) (5300,0.002008609567) (5400,0.002002679394) (5500,0.002067255159) (5600,0.002081850811) (5700,0.00211635161) (5800,0.002108693519) (5900,0.002094354481) (6000,0.002087333694) (6100,0.002121499181) (6200,0.002152150311) (6300,0.002138283337) (6400,0.002186944219) (6500,0.002147716307) (6600,0.002170454455) (6700,0.002061077976) (6800,0.002109170682) (6900,0.00214041546) (7000,0.002271748148) (7100,0.002261465439) (7200,0.002249709959) (7300,0.002294409182) (7400,0.002275869576) (7500,0.002248205664) (7600,0.002215909422) (7700,0.002338773408) (7800,0.002302051079) (7900,0.002264709887) (8000,0.002166397963) (8100,0.00213342926) (8200,0.002096513531) (8300,0.002120634366) (8400,0.002093587408) (8500,0.00209016084) (8600,0.002089902957) (8700,0.002034464397) (8800,0.002017691068) (8900,0.002084442775) (9000,0.002090867551) (9100,0.002099777339) (9200,0.002138746169) (9300,0.002084011922) (9400,0.002089900186) (9500,0.002101654501) (9600,0.002232364553) (9700,0.002205850265) (9800,0.002284738387) (9900,0.002343217877) (10000,0.002324010886) (10100,0.002318373881) (10200,0.002254427678) (10300,0.00224173743) (10400,0.00225116628) (10500,0.00223263799) (10600,0.002143640921) (10700,0.002191845921) (10800,0.002131673496) (10900,0.002034287632) (11000,0.002085358638) (11100,0.002142730064) (11200,0.002168445603) (11300,0.002164604014) (11400,0.0021677988) (11500,0.002165084693) (11600,0.002226319711) (11700,0.002176868811) (11800,0.002138069924) (11900,0.002230810048) (12000,0.002237844188) (12100,0.00219980632) (12200,0.002265100495) (12300,0.002260105719) (12400,0.002290042152) (12500,0.00232696248) (12600,0.002297833061) (12700,0.002294456481) (12800,0.002411376056) (12900,0.002350478433) (13000,0.002362098917) (13100,0.002348052419) (13200,0.0023661452) (13300,0.002449559781) (13400,0.002502423979) (13500,0.002539250895) (13600,0.00249420664) (13700,0.002529326209) (13800,0.002566818788) (13900,0.002554699255) (14000,0.002513346809) (14100,0.00252518882) (14200,0.002458777744) (14300,0.002419662545) (14400,0.002305477462) (14500,0.002237496409) (14600,0.002231991175) (14700,0.002251526504) (14800,0.002135297912) (14900,0.002115240926) (15000,0.002075607609) (15100,0.002100061835) (15200,0.002105441107) (15300,0.00208044704) (15400,0.002204389987) (15500,0.002188947541) (15600,0.002185251028) (15700,0.00214329029) (15800,0.002157429862) (15900,0.002286718809) (16000,0.002317577158) (16100,0.002320647845) (16200,0.002363974066) (16300,0.002401714656) (16400,0.002298604813) (16500,0.002332445187) (16600,0.002369630849) (16700,0.002372572082) (16800,0.002354471129) (16900,0.002249933826) (17000,0.00231638574) (17100,0.002312660264) (17200,0.002301918063) (17300,0.002369855717) (17400,0.002552879509) (17500,0.002562098391) (17600,0.002500357095) (17700,0.002465632057) (17800,0.002533212525) (17900,0.00251369077) (18000,0.002528346947) (18100,0.002564752672) (18200,0.002571872051) (18300,0.002542206075) (18400,0.002407715481) (18500,0.002390941198) (18600,0.002401821921) (18700,0.002470864286) (18800,0.002435732633) (18900,0.002456298354) (19000,0.002446089941) (19100,0.002388172294) (19200,0.00238402735) (19300,0.002427765331) (19400,0.002447078982) (19500,0.002509562229) (19600,0.002565018833) (19700,0.002511983714) (19800,0.002518815151) (19900,0.002529833) (20000,0.002502733096) (20100,0.002591678244) (20200,0.002631093259) (20300,0.002573159873) (20400,0.002497975761) (20500,0.002536722505) (20600,0.002519659325) (20700,0.002633577725) (20800,0.002583082742) (20900,0.002640645485) (21000,0.002661215561) (21100,0.002534450963) (21200,0.002477557305) (21300,0.002478256682) (21400,0.002482707705) (21500,0.002375175525) (21600,0.002385689877) (21700,0.002342624217) (21800,0.002393907052) (21900,0.002372917812) (22000,0.002319834218) (22100,0.002422433556) (22200,0.002417013352) (22300,0.00238614378) (22400,0.002446928248) (22500,0.002379337407) (22600,0.002354450512) (22700,0.002297836763) (22800,0.002234691964) (22900,0.002244612086) (23000,0.002227467764) (23100,0.002190505154) (23200,0.002169682086) (23300,0.00216204233) (23400,0.002174545149) (23500,0.00223013506) (23600,0.002214301133) (23700,0.002216810442) (23800,0.002344401879)};
\addplot[pubRose, opacity=0.18, line width=0.35pt] coordinates {(100,0.01295116916) (200,0.01152314385) (300,0.01108075213) (400,0.01018479152) (500,0.009892436769) (600,0.009777132499) (700,0.009375018733) (800,0.008712905925) (900,0.008490588257) (1000,0.008284120029) (1100,0.007536930498) (1200,0.006888145511) (1300,0.006436493085) (1400,0.006344664446) (1500,0.005982737965) (1600,0.005830487865) (1700,0.005564221018) (1800,0.005644427543) (1900,0.005388212134) (2000,0.005081038736) (2100,0.004957503593) (2200,0.00496180437) (2300,0.004805358103) (2400,0.004645452159) (2500,0.004480584292) (2600,0.004039904755) (2700,0.004069517646) (2800,0.003958204365) (2900,0.003916414524) (3000,0.004069496971) (3100,0.003923091316) (3200,0.003988604969) (3300,0.00395304854) (3400,0.00376526746) (3500,0.003765800991) (3600,0.003680653102) (3700,0.00352256631) (3800,0.003367790813) (3900,0.00324696044) (4000,0.003046207246) (4100,0.003066927753) (4200,0.002925114799) (4300,0.002787376475) (4400,0.002757767402) (4500,0.002645316441) (4600,0.002657930134) (4700,0.002618914517) (4800,0.002585067134) (4900,0.002592425956) (5000,0.002540278691) (5100,0.002541220677) (5200,0.002527443017) (5300,0.002546440181) (5400,0.002545215609) (5500,0.002537249238) (5600,0.002499325131) (5700,0.002462081821) (5800,0.00248738837) (5900,0.0024590282) (6000,0.002482479648) (6100,0.002395721129) (6200,0.002427030634) (6300,0.002484449348) (6400,0.002448061015) (6500,0.002428860613) (6600,0.002413205337) (6700,0.00240363793) (6800,0.002390590007) (6900,0.00243358775) (7000,0.002369235526) (7100,0.002345696627) (7200,0.002248351951) (7300,0.002162415092) (7400,0.002185117058) (7500,0.002133929438) (7600,0.002124421683) (7700,0.00208721878) (7800,0.00216755867) (7900,0.002139752032) (8000,0.002180767898) (8100,0.002213277575) (8200,0.002231653733) (8300,0.002199612185) (8400,0.002180748526) (8500,0.00224153077) (8600,0.002245406795) (8700,0.002321653045) (8800,0.002230629697) (8900,0.002229824592) (9000,0.002192025702) (9100,0.002216772945) (9200,0.002195006958) (9300,0.002236387949) (9400,0.002205698565) (9500,0.002168746269) (9600,0.002192119253) (9700,0.002143239276) (9800,0.002193052275) (9900,0.002204380604) (10000,0.002187291882) (10100,0.002158588055) (10200,0.002129947417) (10300,0.002059079206) (10400,0.002073606977) (10500,0.002080363233) (10600,0.002069894818) (10700,0.002008758916) (10800,0.001976915484) (10900,0.001961344865) (11000,0.001973331335) (11100,0.00194432654) (11200,0.002032392472) (11300,0.002052647737) (11400,0.002030902402) (11500,0.002011986962) (11600,0.002026008419) (11700,0.002093533077) (11800,0.002098865295) (11900,0.002115693106) (12000,0.002111246646) (12100,0.002136322705) (12200,0.002059715556) (12300,0.002058386442) (12400,0.002086588449) (12500,0.002156290237) (12600,0.002196862979) (12700,0.00221525688) (12800,0.002223481413) (12900,0.002209741867) (13000,0.002259213559) (13100,0.002276647638) (13200,0.002343044744) (13300,0.002379415906) (13400,0.002388889738) (13500,0.002380865207) (13600,0.002325102035) (13700,0.002266030246) (13800,0.002247952344) (13900,0.002197819389) (14000,0.002167398878) (14100,0.00216145569) (14200,0.002076888678) (14300,0.002038277686) (14400,0.001957797108) (14500,0.001943913021) (14600,0.001898022194) (14700,0.001972024899) (14800,0.001990249485) (14900,0.001945077942) (15000,0.001956335211) (15100,0.002019398462) (15200,0.002005277085) (15300,0.002016403351) (15400,0.002099015471) (15500,0.002082122234) (15600,0.002107381448) (15700,0.002084755013) (15800,0.002087836852) (15900,0.002145056822) (16000,0.002136795688) (16100,0.002049656166) (16200,0.002063621301) (16300,0.00203402125) (16400,0.002040344069) (16500,0.00203155505) (16600,0.002003150189) (16700,0.00193252078) (16800,0.001976960746) (16900,0.0019726314) (17000,0.001971901616) (17100,0.002030213282) (17200,0.002098853618) (17300,0.002130740648) (17400,0.00208690227) (17500,0.00213248845) (17600,0.002130998974) (17700,0.00216367729) (17800,0.002100091195) (17900,0.002144226804) (18000,0.002105858363) (18100,0.002056878968) (18200,0.001975712762) (18300,0.001988554536) (18400,0.00201272082) (18500,0.001944602444) (18600,0.001966372679) (18700,0.00195865971) (18800,0.001940870541) (18900,0.001874089125) (19000,0.001908443635) (19100,0.001963806944) (19200,0.002021985408) (19300,0.002000140771) (19400,0.001995352469) (19500,0.002042365028) (19600,0.002047479164) (19700,0.002066256176) (19800,0.002049415547) (19900,0.002072959545) (20000,0.002082259336) (20100,0.002072811301) (20200,0.002068634599) (20300,0.002058885584) (20400,0.002022656309) (20500,0.002020509774) (20600,0.001994224521) (20700,0.002050981647) (20800,0.002110036521) (20900,0.002141534782) (21000,0.002178739256) (21100,0.002109884273) (21200,0.002086422592) (21300,0.002149296412) (21400,0.002160444565) (21500,0.002197767526) (21600,0.002242405957) (21700,0.002137621923) (21800,0.002149936429) (21900,0.002140558336) (22000,0.002170479053) (22100,0.002179374325) (22200,0.002177318651) (22300,0.002168094111) (22400,0.002130527608) (22500,0.002066972409) (22600,0.002061796654) (22700,0.002106385375) (22800,0.002057771827) (22900,0.002080201171) (23000,0.001973068691) (23100,0.001979635679) (23200,0.002049862454) (23300,0.002039414318) (23400,0.002089859988) (23500,0.00212393913) (23600,0.002116136975) (23700,0.00220363664) (23800,0.002191149176)};
\addplot[pubRose, opacity=0.18, line width=0.35pt] coordinates {(100,0.005810541101) (200,0.005141387228) (300,0.004413440358) (400,0.004124674131) (500,0.004110829346) (600,0.003890209637) (700,0.003837476451) (800,0.003892918205) (900,0.003928594421) (1000,0.003934420249) (1100,0.003746361355) (1200,0.003697672556) (1300,0.00366081947) (1400,0.003692647233) (1500,0.003579973872) (1600,0.003712058486) (1700,0.003720116406) (1800,0.003683815128) (1900,0.003485965612) (2000,0.00333508586) (2100,0.003185413242) (2200,0.003080118424) (2300,0.003157928842) (2400,0.003099477361) (2500,0.003104147059) (2600,0.002943795547) (2700,0.00291934507) (2800,0.002849368611) (2900,0.00288247047) (3000,0.002876267675) (3100,0.002877596905) (3200,0.00283865321) (3300,0.002708778228) (3400,0.002625113796) (3500,0.002558873012) (3600,0.002538434626) (3700,0.002429799223) (3800,0.002310589538) (3900,0.002271761606) (4000,0.002280281717) (4100,0.002300866763) (4200,0.002279308392) (4300,0.002345061628) (4400,0.002324560075) (4500,0.002309809416) (4600,0.002275145799) (4700,0.002252264391) (4800,0.002213813923) (4900,0.002208343311) (5000,0.002258511609) (5100,0.002221686975) (5200,0.00222552733) (5300,0.002149309195) (5400,0.002161144931) (5500,0.002165186964) (5600,0.002187817427) (5700,0.002195727779) (5800,0.002198499569) (5900,0.00214044305) (6000,0.002031906368) (6100,0.002069938369) (6200,0.001991628716) (6300,0.001982480893) (6400,0.00197587274) (6500,0.001942085684) (6600,0.001882032806) (6700,0.001857088821) (6800,0.001888257219) (6900,0.001916065323) (7000,0.001886873331) (7100,0.001785034069) (7200,0.001915684214) (7300,0.001919855725) (7400,0.001932508999) (7500,0.001998272457) (7600,0.002050876594) (7700,0.002007059194) (7800,0.001992176333) (7900,0.001992084004) (8000,0.001995823113) (8100,0.002063046978) (8200,0.001976890699) (8300,0.001966143213) (8400,0.001996829594) (8500,0.001980453241) (8600,0.002014230355) (8700,0.002095330157) (8800,0.002104960801) (8900,0.002085958852) (9000,0.002159597736) (9100,0.002146372932) (9200,0.002092930197) (9300,0.002183761017) (9400,0.002110033494) (9500,0.002050686174) (9600,0.001996043324) (9700,0.0019486345) (9800,0.001905489655) (9900,0.001948336698) (10000,0.00185436504) (10100,0.00179984346) (10200,0.001853784127) (10300,0.001778399874) (10400,0.00174213961) (10500,0.001735398744) (10600,0.001781566138) (10700,0.001819006575) (10800,0.001883010298) (10900,0.001869530289) (11000,0.001954052818) (11100,0.001984636963) (11200,0.001962465991) (11300,0.001943829574) (11400,0.002019334212) (11500,0.002050595777) (11600,0.00206935734) (11700,0.002068188158) (11800,0.002083749534) (11900,0.002076713182) (12000,0.002082034713) (12100,0.002090771659) (12200,0.002080592292) (12300,0.002096864721) (12400,0.002037755866) (12500,0.002036380209) (12600,0.001995697082) (12700,0.00198792459) (12800,0.001943271817) (12900,0.001961400593) (13000,0.001878757484) (13100,0.001872525399) (13200,0.001977111003) (13300,0.001962106384) (13400,0.001985639485) (13500,0.001976738533) (13600,0.001987483341) (13700,0.001968049374) (13800,0.001977225929) (13900,0.001922214706) (14000,0.001922531449) (14100,0.001920726616) (14200,0.001815884199) (14300,0.001819001115) (14400,0.001841997285) (14500,0.001891954849) (14600,0.001842490514) (14700,0.00186556814) (14800,0.001838517364) (14900,0.001946922706) (15000,0.001983477967) (15100,0.001957967482) (15200,0.002026644524) (15300,0.00203324547) (15400,0.002014271379) (15500,0.001993144071) (15600,0.001995859738) (15700,0.002017220226) (15800,0.00204797464) (15900,0.001986883464) (16000,0.002024281817) (16100,0.002057210007) (16200,0.002029560134) (16300,0.002122430643) (16400,0.002119638096) (16500,0.00211437454) (16600,0.002106411895) (16700,0.002058202459) (16800,0.002078540565) (16900,0.002098882326) (17000,0.002114603051) (17100,0.002077750489) (17200,0.002054585307) (17300,0.002005674411) (17400,0.002062393236) (17500,0.002044258569) (17600,0.002050887991) (17700,0.002059368091) (17800,0.002020459448) (17900,0.002029304334) (18000,0.002019427961) (18100,0.002099317196) (18200,0.002062918269) (18300,0.002022321778) (18400,0.001952506835) (18500,0.001938273269) (18600,0.001943025552) (18700,0.001917436183) (18800,0.001928828482) (18900,0.001888749504) (19000,0.001826886542) (19100,0.001789956552) (19200,0.001833633648) (19300,0.001859332842) (19400,0.001898533397) (19500,0.001961010217) (19600,0.001974691008) (19700,0.001995897049) (19800,0.001968244556) (19900,0.001990846428) (20000,0.00204162742) (20100,0.002033657092) (20200,0.001986386161) (20300,0.001995530771) (20400,0.001980615966) (20500,0.001899941219) (20600,0.001900637965) (20700,0.001850795886) (20800,0.001864729053) (20900,0.001896850811) (21000,0.001874835743) (21100,0.001894366182) (21200,0.001905898121) (21300,0.001828603004) (21400,0.001808813307) (21500,0.00180840909) (21600,0.001800144324) (21700,0.001861043938) (21800,0.001840711711) (21900,0.001801851904) (22000,0.001793441141) (22100,0.001731731836) (22200,0.001826119749) (22300,0.001835561928) (22400,0.001884547144) (22500,0.001901306619) (22600,0.001944784832) (22700,0.001923950913) (22800,0.00200258306) (22900,0.002034415503) (23000,0.002055599866) (23100,0.002123489662) (23200,0.002031303605) (23300,0.002080862317) (23400,0.002034333674) (23500,0.002053737105) (23600,0.001979872386) (23700,0.00197857843) (23800,0.001922212634)};
\addplot[pubRose, opacity=0.18, line width=0.35pt] coordinates {(100,0.00371650001) (200,0.004280363908) (300,0.004073562411) (400,0.003945346281) (500,0.003847093182) (600,0.00396963694) (700,0.003851257564) (800,0.003879533062) (900,0.003823670559) (1000,0.003796084039) (1100,0.00377156334) (1200,0.003594952845) (1300,0.003579968377) (1400,0.003652691934) (1500,0.003652442386) (1600,0.003672292782) (1700,0.003677599691) (1800,0.003749900963) (1900,0.003685218142) (2000,0.003641024628) (2100,0.003607215197) (2200,0.003593930602) (2300,0.003590311669) (2400,0.00343682915) (2500,0.00345504093) (2600,0.003269009432) (2700,0.003273375146) (2800,0.003165662661) (2900,0.003125281841) (3000,0.003074004827) (3100,0.00299514872) (3200,0.002913786215) (3300,0.002763991128) (3400,0.002693402977) (3500,0.002480794676) (3600,0.002421203954) (3700,0.002369649475) (3800,0.002218799666) (3900,0.002244413737) (4000,0.002249648469) (4100,0.002218631236) (4200,0.002236892283) (4300,0.002238550736) (4400,0.002299215528) (4500,0.002319527324) (4600,0.00225319186) (4700,0.00214732684) (4800,0.002149018727) (4900,0.002070924162) (5000,0.00197009499) (5100,0.001926614845) (5200,0.001881454035) (5300,0.001899602229) (5400,0.001835816482) (5500,0.001878565492) (5600,0.001890916878) (5700,0.001934966294) (5800,0.001907701965) (5900,0.001876566105) (6000,0.001872247376) (6100,0.001891706593) (6200,0.001905005041) (6300,0.001889257028) (6400,0.001904662244) (6500,0.001902637503) (6600,0.001920685545) (6700,0.00188699842) (6800,0.001893184194) (6900,0.00193906232) (7000,0.001998373901) (7100,0.001986678783) (7200,0.001956790767) (7300,0.001965493558) (7400,0.001933448005) (7500,0.001878083951) (7600,0.001891030429) (7700,0.001885901077) (7800,0.001920962671) (7900,0.001962254453) (8000,0.001950982842) (8100,0.001949164702) (8200,0.001925632218) (8300,0.001898220694) (8400,0.001873941463) (8500,0.001920476742) (8600,0.001970896451) (8700,0.001980736735) (8800,0.001944703748) (8900,0.001923509082) (9000,0.001918849791) (9100,0.001969414181) (9200,0.001997031551) (9300,0.001990669535) (9400,0.002045600989) (9500,0.002064097475) (9600,0.002008256491) (9700,0.002047787921) (9800,0.002062367077) (9900,0.002108170267) (10000,0.002095739439) (10100,0.002102532296) (10200,0.002114001266) (10300,0.002136816282) (10400,0.002043833409) (10500,0.001967560663) (10600,0.001992809633) (10700,0.001999064675) (10800,0.001977785537) (10900,0.001935363002) (11000,0.002005294664) (11100,0.001960140863) (11200,0.001994235243) (11300,0.001965047547) (11400,0.002109241055) (11500,0.002196200611) (11600,0.002217616374) (11700,0.002204313688) (11800,0.002245131391) (11900,0.002243729471) (12000,0.002172141848) (12100,0.002180746384) (12200,0.002174803242) (12300,0.002211039932) (12400,0.002192887478) (12500,0.002194932452) (12600,0.002180145355) (12700,0.002248388599) (12800,0.00221338456) (12900,0.002249899204) (13000,0.002287270455) (13100,0.00233826153) (13200,0.002419484453) (13300,0.002454679459) (13400,0.002442482021) (13500,0.002472906583) (13600,0.002439522813) (13700,0.002336683124) (13800,0.002313227346) (13900,0.002312226454) (14000,0.00232665278) (14100,0.002285155561) (14200,0.002198984521) (14300,0.0021923183) (14400,0.002203424997) (14500,0.002141478402) (14600,0.002144261426) (14700,0.002196160774) (14800,0.002216290915) (14900,0.002156661719) (15000,0.002092762408) (15100,0.002120467811) (15200,0.002126022498) (15300,0.002174369991) (15400,0.002134394646) (15500,0.002181388857) (15600,0.00219810002) (15700,0.002213894553) (15800,0.002252020733) (15900,0.002258647128) (16000,0.002306623198) (16100,0.00228940011) (16200,0.002292756829) (16300,0.002247974346) (16400,0.002229371131) (16500,0.002187779313) (16600,0.002156800905) (16700,0.002143996803) (16800,0.002107537736) (16900,0.002165803302) (17000,0.002180331096) (17100,0.002155961015) (17200,0.00216162802) (17300,0.002178880095) (17400,0.002266699343) (17500,0.002322505333) (17600,0.002293840074) (17700,0.002290359279) (17800,0.002266390203) (17900,0.002254859032) (18000,0.002187457983) (18100,0.00222723179) (18200,0.002192832017) (18300,0.002148627071) (18400,0.002073243982) (18500,0.002035882394) (18600,0.002133142902) (18700,0.002053444949) (18800,0.0020475809) (18900,0.001979184803) (19000,0.002002380509) (19100,0.002042197622) (19200,0.002053046902) (19300,0.002014631464) (19400,0.002055688517) (19500,0.002041721193) (19600,0.002048066107) (19700,0.002181556204) (19800,0.002172655019) (19900,0.002169907303) (20000,0.002209209278) (20100,0.00212234736) (20200,0.002172785858) (20300,0.002175641619) (20400,0.002099554357) (20500,0.002087471029) (20600,0.001999279612) (20700,0.001947142032) (20800,0.001942122309) (20900,0.001982660161) (21000,0.001954584464) (21100,0.001945498248) (21200,0.001880888338) (21300,0.001899592462) (21400,0.001896607049) (21500,0.001892338495) (21600,0.00189673393) (21700,0.001812231611) (21800,0.001877062873) (21900,0.001817482547) (22000,0.001839650632) (22100,0.001934969379) (22200,0.001944107318) (22300,0.001923343644) (22400,0.00200121454) (22500,0.001987689896) (22600,0.002056018054) (22700,0.002104300819) (22800,0.002092749137) (22900,0.002186566184) (23000,0.002203132177) (23100,0.002176014881) (23200,0.002157004597) (23300,0.002167671185) (23400,0.002113942115) (23500,0.002151549456) (23600,0.002101533848) (23700,0.002110403671) (23800,0.002105466102)};
\addplot[pubRose, opacity=0.18, line width=0.35pt] coordinates {(100,0.007488208357) (200,0.007265596418) (300,0.006973829586) (400,0.006917110179) (500,0.006524895318) (600,0.006057541274) (700,0.005855063575) (800,0.005705661606) (900,0.005478977122) (1000,0.005277049541) (1100,0.004925493058) (1200,0.004520785157) (1300,0.004180586385) (1400,0.003817964578) (1500,0.003637393727) (1600,0.003606361686) (1700,0.003539839177) (1800,0.003447481641) (1900,0.003490903066) (2000,0.003471226594) (2100,0.003337488533) (2200,0.003392349347) (2300,0.003379880264) (2400,0.003361132438) (2500,0.00326786472) (2600,0.003220884292) (2700,0.003138190578) (2800,0.003028370766) (2900,0.002979506343) (3000,0.002995384089) (3100,0.003026718204) (3200,0.002875537821) (3300,0.002846710617) (3400,0.002861812036) (3500,0.002867235476) (3600,0.002905444265) (3700,0.002877471759) (3800,0.002907927427) (3900,0.002751285327) (4000,0.002743169153) (4100,0.002739376086) (4200,0.002744146134) (4300,0.002753042104) (4400,0.002668925794) (4500,0.002730471781) (4600,0.002635050076) (4700,0.002545444365) (4800,0.002459740196) (4900,0.00255308121) (5000,0.002473804238) (5100,0.002436561882) (5200,0.002566921525) (5300,0.002578992327) (5400,0.002582724229) (5500,0.00250900744) (5600,0.00252740446) (5700,0.002583217039) (5800,0.002579853148) (5900,0.002540068608) (6000,0.00260358206) (6100,0.002598805935) (6200,0.002485072799) (6300,0.00256128998) (6400,0.002571744751) (6500,0.00266865585) (6600,0.002733026934) (6700,0.002759399847) (6800,0.002795653464) (6900,0.002738501038) (7000,0.002655403153) (7100,0.00268977324) (7200,0.002757188585) (7300,0.002728073392) (7400,0.002762433398) (7500,0.002800742933) (7600,0.002733131009) (7700,0.002740575117) (7800,0.002764425427) (7900,0.002919281204) (8000,0.00303648964) (8100,0.002966693696) (8200,0.002910505049) (8300,0.002867579204) (8400,0.002933468949) (8500,0.002797860629) (8600,0.002705634956) (8700,0.002675600804) (8800,0.002744675695) (8900,0.002677645173) (9000,0.002622476418) (9100,0.002630608517) (9200,0.002706254774) (9300,0.00269996057) (9400,0.00260906663) (9500,0.002649641305) (9600,0.002718494227) (9700,0.0027174399) (9800,0.002800513688) (9900,0.002789350855) (10000,0.002777263639) (10100,0.002903863485) (10200,0.0028789951) (10300,0.002879992221) (10400,0.002906139172) (10500,0.002898852387) (10600,0.003065444878) (10700,0.003077020147) (10800,0.002917865384) (10900,0.00297485604) (11000,0.002993818605) (11100,0.002897097799) (11200,0.002941081882) (11300,0.002994882641) (11400,0.002998218383) (11500,0.002970419545) (11600,0.002814829559) (11700,0.002922888193) (11800,0.002950842911) (11900,0.002897260501) (12000,0.00287276844) (12100,0.002915279474) (12200,0.002899703034) (12300,0.002847108105) (12400,0.002799748373) (12500,0.002985624992) (12600,0.002925993945) (12700,0.002854180429) (12800,0.002785227983) (12900,0.002712900355) (13000,0.00270321893) (13100,0.002699626889) (13200,0.0026366723) (13300,0.002693034033) (13400,0.002673921874) (13500,0.002532324474) (13600,0.002633892535) (13700,0.002568608779) (13800,0.002610963699) (13900,0.002706264448) (14000,0.00273433479) (14100,0.002777454583) (14200,0.002848158102) (14300,0.002793223667) (14400,0.002900413633) (14500,0.002942509274) (14600,0.002882554824) (14700,0.002905476163) (14800,0.00293631996) (14900,0.002800849767) (15000,0.002713087911) (15100,0.00265445431) (15200,0.002630780835) (15300,0.002550722484) (15400,0.002540233324) (15500,0.002520896995) (15600,0.002617245086) (15700,0.002632287133) (15800,0.002569296083) (15900,0.002621467621) (16000,0.002734314697) (16100,0.002725925506) (16200,0.002733237785) (16300,0.002892671595) (16400,0.002941261115) (16500,0.002894432913) (16600,0.002873157198) (16700,0.002925983304) (16800,0.002918472094) (16900,0.002938863169) (17000,0.002855320042) (17100,0.002854593145) (17200,0.002742669568) (17300,0.002643377637) (17400,0.002562028798) (17500,0.002553216927) (17600,0.002566648275) (17700,0.002464065677) (17800,0.00246808033) (17900,0.002565143001) (18000,0.002520065429) (18100,0.002488515014) (18200,0.002568715182) (18300,0.002683576918) (18400,0.002658466948) (18500,0.002697028453) (18600,0.002609591931) (18700,0.002655888093) (18800,0.002635727194) (18900,0.002585163317) (19000,0.002631442714) (19100,0.002606528788) (19200,0.002615202428) (19300,0.002467847173) (19400,0.002504107472) (19500,0.00245630641) (19600,0.002555124438) (19700,0.002478493657) (19800,0.002548022801) (19900,0.002495964291) (20000,0.002503772778) (20100,0.002479621058) (20200,0.00246577271) (20300,0.002474078035) (20400,0.002430086106) (20500,0.002459215268) (20600,0.002455284062) (20700,0.002446772438) (20800,0.002379374532) (20900,0.002329658042) (21000,0.002416376583) (21100,0.002489538526) (21200,0.002569034125) (21300,0.002721863368) (21400,0.002710473014) (21500,0.002690961992) (21600,0.002547777712) (21700,0.002653791709) (21800,0.002711168118) (21900,0.00276651138) (22000,0.00264851097) (22100,0.002579585789) (22200,0.002448913851) (22300,0.002325313794) (22400,0.002328193747) (22500,0.00227147413) (22600,0.002309321729) (22700,0.002197750122) (22800,0.002211884712) (22900,0.002169796778) (23000,0.002190491976) (23100,0.002257955563) (23200,0.002284390666) (23300,0.002285956126) (23400,0.002221118624) (23500,0.00231992983) (23600,0.002343259088) (23700,0.00242902009) (23800,0.002327593334)};
\addplot[pubRose, ultra thick] coordinates {(100,0.006447343156) (200,0.006212343345) (300,0.005941413265) (400,0.005816810532) (500,0.005595357437) (600,0.005682628679) (700,0.005874917709) (800,0.005786519934) (900,0.005639712559) (1000,0.005578001624) (1100,0.005362844456) (1200,0.005172914453) (1300,0.005028667545) (1400,0.004944446159) (1500,0.004897844989) (1600,0.004746778647) (1700,0.004493617476) (1800,0.004362605896) (1900,0.004242701188) (2000,0.004107952316) (2100,0.003929080477) (2200,0.003853939672) (2300,0.003745832713) (2400,0.003567060083) (2500,0.003417098557) (2600,0.003307779378) (2700,0.00333239187) (2800,0.003171333624) (2900,0.003120948607) (3000,0.00302800223) (3100,0.002983877482) (3200,0.002877082431) (3300,0.002825037809) (3400,0.002858339157) (3500,0.00280383529) (3600,0.002699134941) (3700,0.002700748644) (3800,0.002656368527) (3900,0.002639417606) (4000,0.002590204973) (4100,0.002508933353) (4200,0.002482608007) (4300,0.002420587087) (4400,0.002412305854) (4500,0.002357972786) (4600,0.002327745495) (4700,0.002320894721) (4800,0.002298623242) (4900,0.002347579622) (5000,0.00232952676) (5100,0.002316866606) (5200,0.002345507644) (5300,0.002298383677) (5400,0.002273593732) (5500,0.002257310646) (5600,0.002216597815) (5700,0.002183586528) (5800,0.002166136145) (5900,0.002140366734) (6000,0.002105192433) (6100,0.002130896604) (6200,0.002149969846) (6300,0.002162889269) (6400,0.002205628791) (6500,0.002170472051) (6600,0.0021645375) (6700,0.002083838615) (6800,0.002125712787) (6900,0.002126944828) (7000,0.002170134382) (7100,0.002179732092) (7200,0.002199092373) (7300,0.002153983619) (7400,0.002191984281) (7500,0.002172725572) (7600,0.002170165553) (7700,0.002157919476) (7800,0.002226136968) (7900,0.002153936226) (8000,0.002144452068) (8100,0.002098238119) (8200,0.002042499941) (8300,0.00207546585) (8400,0.002086851187) (8500,0.002132965135) (8600,0.00211686492) (8700,0.002117526782) (8800,0.002164244669) (8900,0.002150781266) (9000,0.002175811719) (9100,0.002181572939) (9200,0.002160331002) (9300,0.002169515437) (9400,0.002123016014) (9500,0.002082875988) (9600,0.002171023289) (9700,0.002145076223) (9800,0.002131497796) (9900,0.002156275435) (10000,0.00214151566) (10100,0.002157789672) (10200,0.002154553658) (10300,0.002136216959) (10400,0.002074146993) (10500,0.002089675999) (10600,0.002072583331) (10700,0.002032005094) (10800,0.002012076706) (10900,0.00201301105) (11000,0.002010645531) (11100,0.001999521937) (11200,0.002046588331) (11300,0.002051421464) (11400,0.002102912008) (11500,0.002160687017) (11600,0.002186166542) (11700,0.002190591249) (11800,0.00217611721) (11900,0.002182667411) (12000,0.002141694247) (12100,0.002113547182) (12200,0.002100475703) (12300,0.002080837131) (12400,0.002100357134) (12500,0.002148686233) (12600,0.002188504167) (12700,0.002231822739) (12800,0.002218432986) (12900,0.002229820535) (13000,0.002273242007) (13100,0.002307454584) (13200,0.002354594972) (13300,0.002414487844) (13400,0.002388650249) (13500,0.002398762933) (13600,0.002353841381) (13700,0.002251006116) (13800,0.00223093665) (13900,0.002239112032) (14000,0.00223116522) (14100,0.002247624699) (14200,0.002226755192) (14300,0.002212221042) (14400,0.002181934257) (14500,0.002142205165) (14600,0.002106319997) (14700,0.002122328465) (14800,0.002132366516) (14900,0.002135249262) (15000,0.002084185008) (15100,0.002110264823) (15200,0.002115731803) (15300,0.002127408516) (15400,0.002169392316) (15500,0.002185168199) (15600,0.002191675524) (15700,0.002178592421) (15800,0.002189459361) (15900,0.002230191982) (16000,0.002237292309) (16100,0.002202198631) (16200,0.002243180922) (16300,0.00223609634) (16400,0.002255328931) (16500,0.002218880551) (16600,0.002204278571) (16700,0.002194148063) (16800,0.002156831307) (16900,0.002191571536) (17000,0.002246092394) (17100,0.002221404231) (17200,0.002192114625) (17300,0.002154810372) (17400,0.002196584112) (17500,0.002243759274) (17600,0.002266964549) (17700,0.002268360567) (17800,0.002238357166) (17900,0.002225117019) (18000,0.002198056568) (18100,0.002227124124) (18200,0.002256323007) (18300,0.002202735731) (18400,0.002151009109) (18500,0.002117136918) (18600,0.002157758357) (18700,0.002098301466) (18800,0.002125942206) (18900,0.002108333958) (19000,0.002178049152) (19100,0.002195269713) (19200,0.002131650626) (19300,0.002161239431) (19400,0.002142912732) (19500,0.002122742165) (19600,0.002122981672) (19700,0.002134055051) (19800,0.002147267899) (19900,0.002144095354) (20000,0.002154862619) (20100,0.002151537978) (20200,0.002194260433) (20300,0.002171428548) (20400,0.002133133757) (20500,0.002148068842) (20600,0.002114516351) (20700,0.002132309211) (20800,0.002162077022) (20900,0.002194080979) (21000,0.002228600404) (21100,0.002177134738) (21200,0.002161328337) (21300,0.0021680801) (21400,0.002186887705) (21500,0.00222054419) (21600,0.002248996677) (21700,0.002228590025) (21800,0.002196510491) (21900,0.002172453463) (22000,0.002162417304) (22100,0.002145949274) (22200,0.002165878477) (22300,0.00219025432) (22400,0.002171510569) (22500,0.002144546882) (22600,0.002147979086) (22700,0.002149844996) (22800,0.002185929328) (22900,0.002178181481) (23000,0.002196812077) (23100,0.002183260018) (23200,0.002163343341) (23300,0.002164856758) (23400,0.002144243632) (23500,0.002190842258) (23600,0.002165219054) (23700,0.002210223541) (23800,0.002249305736)};
\nextgroupplot[title={Actor loss}, xlabel={Gradient update}, scaled y ticks=true]
\addplot[pubRose, opacity=0.18, line width=0.35pt] coordinates {(100,0.227704674) (200,0.2236740068) (300,0.2228807112) (400,0.2207037136) (500,0.2145386487) (600,0.2131878113) (700,0.2108115192) (800,0.2095240243) (900,0.2066535966) (1000,0.2049626902) (1100,0.2006382212) (1200,0.1954978824) (1300,0.1893014789) (1400,0.1840271309) (1500,0.182219258) (1600,0.1794308171) (1700,0.1759318531) (1800,0.1725370541) (1900,0.1705290601) (2000,0.1683709517) (2100,0.1642793164) (2200,0.1632880509) (2300,0.1627684295) (2400,0.1633872673) (2500,0.1616073668) (2600,0.1590655416) (2700,0.1577972531) (2800,0.1564256892) (2900,0.1547400534) (3000,0.152807723) (3100,0.1538961917) (3200,0.1522758618) (3300,0.1511261553) (3400,0.1488473386) (3500,0.1470750466) (3600,0.1460305884) (3700,0.1448199153) (3800,0.143683058) (3900,0.1419013187) (4000,0.1394574195) (4100,0.1371868581) (4200,0.1350542791) (4300,0.1325111307) (4400,0.1305754505) (4500,0.1285836197) (4600,0.1266334735) (4700,0.1239366747) (4800,0.1217649542) (4900,0.1213627197) (5000,0.1195577741) (5100,0.1153215013) (5200,0.1144184612) (5300,0.1141279668) (5400,0.1116496749) (5500,0.1114988886) (5600,0.1100880884) (5700,0.1092824422) (5800,0.107033363) (5900,0.1038952537) (6000,0.1035874471) (6100,0.1044506595) (6200,0.1017278709) (6300,0.09960242584) (6400,0.100145936) (6500,0.0980101116) (6600,0.09595473111) (6700,0.09319885001) (6800,0.09031810015) (6900,0.08819566965) (7000,0.08629255742) (7100,0.0834683232) (7200,0.08122814931) (7300,0.07841106467) (7400,0.07474919148) (7500,0.07060743272) (7600,0.06850207224) (7700,0.0671247758) (7800,0.06558122821) (7900,0.06461346187) (8000,0.06128445156) (8100,0.0596196074) (8200,0.06002737954) (8300,0.0579253085) (8400,0.05535928607) (8500,0.05560414456) (8600,0.05358582288) (8700,0.05218711384) (8800,0.05205055028) (8900,0.04939434081) (9000,0.04969217442) (9100,0.04951112084) (9200,0.04680345356) (9300,0.04606079571) (9400,0.0457154993) (9500,0.04376730807) (9600,0.04277848192) (9700,0.04140839316) (9800,0.03913102262) (9900,0.03901360445) (10000,0.03618693762) (10100,0.03384966217) (10200,0.03209925536) (10300,0.03167352546) (10400,0.03002543468) (10500,0.02962263413) (10600,0.02815672681) (10700,0.02552243494) (10800,0.02406995385) (10900,0.02175991912) (11000,0.02089274474) (11100,0.0188029414) (11200,0.01787181851) (11300,0.01559539922) (11400,0.01431610389) (11500,0.01223697998) (11600,0.0110413691) (11700,0.01078471418) (11800,0.01030105751) (11900,0.009810985927) (12000,0.008473585639) (12100,0.008032973157) (12200,0.007702917932) (12300,0.006566327438) (12400,0.00500985086) (12500,0.002981232014) (12600,0.0008672506548) (12700,-0.0005582165206) (12800,-0.002664642711) (12900,-0.004556892882) (13000,-0.005720788194) (13100,-0.006564126862) (13200,-0.007957705436) (13300,-0.007572887384) (13400,-0.00814323622) (13500,-0.008100662206) (13600,-0.008197986719) (13700,-0.008971659106) (13800,-0.009197743156) (13900,-0.008727277687) (14000,-0.008658917632) (14100,-0.009274180431) (14200,-0.01040365874) (14300,-0.01256097555) (14400,-0.0126459673) (14500,-0.01309539266) (14600,-0.01371456645) (14700,-0.01373861004) (14800,-0.0139114731) (14900,-0.01521097994) (15000,-0.01569661498) (15100,-0.01645431304) (15200,-0.01611801367) (15300,-0.01601460259) (15400,-0.01710183909) (15500,-0.01710038865) (15600,-0.01759686479) (15700,-0.01755874818) (15800,-0.01817878662) (15900,-0.01876977971) (16000,-0.01950799283) (16100,-0.01928395629) (16200,-0.02021883009) (16300,-0.02114719758) (16400,-0.02126928093) (16500,-0.02257158402) (16600,-0.02313837465) (16700,-0.02399291396) (16800,-0.0250209244) (16900,-0.02543744873) (17000,-0.02543333713) (17100,-0.02614356522) (17200,-0.02626247797) (17300,-0.0264921641) (17400,-0.02709751986) (17500,-0.02639300991) (17600,-0.02616781108) (17700,-0.02606838085) (17800,-0.0253835747) (17900,-0.02568228208) (18000,-0.02575297989) (18100,-0.02596505675) (18200,-0.02642400712) (18300,-0.0260070052) (18400,-0.02617797516) (18500,-0.02804924585) (18600,-0.02844741028) (18700,-0.0277841026) (18800,-0.02792412043) (18900,-0.02730493061) (19000,-0.02819490917) (19100,-0.0282848984) (19200,-0.02880411167) (19300,-0.02859964855) (19400,-0.02805518527) (19500,-0.02736322861) (19600,-0.02713609003) (19700,-0.02808105052) (19800,-0.02773409449) (19900,-0.02790598106) (20000,-0.02768895738) (20100,-0.02769008726) (20200,-0.0273935955) (20300,-0.02824452091) (20400,-0.02902303543) (20500,-0.02873747535) (20600,-0.0292180296) (20700,-0.02903732974) (20800,-0.02948449645) (20900,-0.02894681487) (21000,-0.0299596766) (21100,-0.03002510648) (21200,-0.03043713588) (21300,-0.03011711724) (21400,-0.02939651534) (21500,-0.02992751393) (21600,-0.0296761103) (21700,-0.02947579641) (21800,-0.02922333684) (21900,-0.02976884041) (22000,-0.02978718523) (22100,-0.02907420695) (22200,-0.02892209589) (22300,-0.03002685774) (22400,-0.03015478738) (22500,-0.02953066193) (22600,-0.02948094588) (22700,-0.02980325706) (22800,-0.03044079207) (22900,-0.03108704332) (23000,-0.03005356323) (23100,-0.03073414918) (23200,-0.03042131606) (23300,-0.02867422737) (23400,-0.02906845938) (23500,-0.0289580917) (23600,-0.02919532862) (23700,-0.02934907284) (23800,-0.02882746439)};
\addplot[pubRose, opacity=0.18, line width=0.35pt] coordinates {(100,0.1673000157) (200,0.1653966159) (300,0.1694593529) (400,0.1718969978) (500,0.1730568886) (600,0.1752892435) (700,0.1739938791) (800,0.1743069496) (900,0.1737961223) (1000,0.1742417097) (1100,0.1731893122) (1200,0.1735162914) (1300,0.1709372059) (1400,0.1686113819) (1500,0.1674720779) (1600,0.1648889557) (1700,0.163459608) (1800,0.1633223549) (1900,0.1633690923) (2000,0.1611434892) (2100,0.1593573451) (2200,0.1585101098) (2300,0.1576065212) (2400,0.1571451947) (2500,0.1565868005) (2600,0.1574065477) (2700,0.1569101125) (2800,0.1546427235) (2900,0.1534208983) (3000,0.1526665539) (3100,0.1526089489) (3200,0.1508553237) (3300,0.1499964118) (3400,0.1499672711) (3500,0.148192358) (3600,0.1465012625) (3700,0.1463523328) (3800,0.1450307563) (3900,0.1444483578) (4000,0.1438964918) (4100,0.1442980662) (4200,0.1432692304) (4300,0.1426638365) (4400,0.1414963916) (4500,0.1410518542) (4600,0.1382101461) (4700,0.1368756875) (4800,0.1362124786) (4900,0.1325781569) (5000,0.1321716473) (5100,0.1306606874) (5200,0.1297598913) (5300,0.1299825326) (5400,0.1287993416) (5500,0.1277059719) (5600,0.1276902638) (5700,0.1260096736) (5800,0.1243208893) (5900,0.1239344686) (6000,0.1231234699) (6100,0.1206643425) (6200,0.1206046395) (6300,0.1179607287) (6400,0.1145129904) (6500,0.1123442605) (6600,0.1110221379) (6700,0.1098526664) (6800,0.1086575665) (6900,0.1074890807) (7000,0.1044310413) (7100,0.1063092589) (7200,0.1051023252) (7300,0.1040815637) (7400,0.1064681433) (7500,0.1045474328) (7600,0.1025622226) (7700,0.1017784335) (7800,0.1007526092) (7900,0.09954362437) (8000,0.09770035669) (8100,0.09347339794) (8200,0.09033419192) (8300,0.08936580643) (8400,0.08565242887) (8500,0.08360874876) (8600,0.08295354769) (8700,0.08129782155) (8800,0.07972771451) (8900,0.07959618196) (9000,0.07783137709) (9100,0.07806391641) (9200,0.07721810639) (9300,0.07550319359) (9400,0.07268036529) (9500,0.07102257088) (9600,0.07134879977) (9700,0.0710782215) (9800,0.0704586193) (9900,0.06797579899) (10000,0.066955119) (10100,0.06436729468) (10200,0.06390109695) (10300,0.06426232196) (10400,0.06474271454) (10500,0.06541281678) (10600,0.06297649704) (10700,0.06113821715) (10800,0.06025861725) (10900,0.06001006439) (11000,0.05920099504) (11100,0.05879281536) (11200,0.05698326007) (11300,0.05404578224) (11400,0.05310373157) (11500,0.05287633725) (11600,0.05224769488) (11700,0.05235972144) (11800,0.05028896704) (11900,0.04829340242) (12000,0.04852452129) (12100,0.0505029466) (12200,0.05081045181) (12300,0.05082344115) (12400,0.04943171442) (12500,0.04851556346) (12600,0.04690805711) (12700,0.04567133375) (12800,0.04638171569) (12900,0.04586552754) (13000,0.04498963021) (13100,0.04201364592) (13200,0.04194764905) (13300,0.04270430878) (13400,0.04420134984) (13500,0.04366104789) (13600,0.04357981384) (13700,0.04338821769) (13800,0.04269112386) (13900,0.04337122552) (14000,0.04295585975) (14100,0.04176021069) (14200,0.04009640925) (14300,0.03852661848) (14400,0.03786420338) (14500,0.03712679595) (14600,0.03617501575) (14700,0.03489202391) (14800,0.03418124449) (14900,0.0330828039) (15000,0.03320245799) (15100,0.03424322512) (15200,0.03362131026) (15300,0.03420715574) (15400,0.0333233336) (15500,0.03316018712) (15600,0.03260582052) (15700,0.03170129154) (15800,0.03093264662) (15900,0.03132590875) (16000,0.03009190932) (16100,0.02710253736) (16200,0.02772179684) (16300,0.02704235232) (16400,0.02564578531) (16500,0.0275018326) (16600,0.02816412961) (16700,0.02839429872) (16800,0.02787811896) (16900,0.02848733468) (17000,0.02861868227) (17100,0.02981412821) (17200,0.02996537797) (17300,0.02913576607) (17400,0.03013637047) (17500,0.02840222511) (17600,0.02809872981) (17700,0.02762117274) (17800,0.0279544875) (17900,0.02679311782) (18000,0.02701582499) (18100,0.02709570769) (18200,0.02657559998) (18300,0.02645629048) (18400,0.02678043209) (18500,0.02552427966) (18600,0.02514566574) (18700,0.02586196307) (18800,0.02692089174) (18900,0.02627575714) (19000,0.02618408147) (19100,0.02703769449) (19200,0.02681142557) (19300,0.02713172436) (19400,0.02731495872) (19500,0.02813394573) (19600,0.02900456656) (19700,0.03008543067) (19800,0.0298587393) (19900,0.03110769615) (20000,0.03033150099) (20100,0.03018935639) (20200,0.03018844742) (20300,0.02967552785) (20400,0.02807584964) (20500,0.02772662975) (20600,0.02692918126) (20700,0.02503182311) (20800,0.02432006244) (20900,0.024012761) (21000,0.02458927259) (21100,0.02360847834) (21200,0.02261574343) (21300,0.02192388205) (21400,0.02269623233) (21500,0.02218510369) (21600,0.02190343002) (21700,0.02346448908) (21800,0.02347417371) (21900,0.02340640733) (22000,0.02229807694) (22100,0.02184502799) (22200,0.02240844592) (22300,0.02202936383) (22400,0.02210248867) (22500,0.02185783396) (22600,0.02160511659) (22700,0.01951351101) (22800,0.01933125192) (22900,0.01840839582) (23000,0.019235464) (23100,0.01920998478) (23200,0.01881160894) (23300,0.01964634182) (23400,0.01963652642) (23500,0.01960175568) (23600,0.01971727936) (23700,0.02138255769) (23800,0.01965001971)};
\addplot[pubRose, opacity=0.18, line width=0.35pt] coordinates {(100,0.4187889099) (200,0.4154426157) (300,0.421310703) (400,0.4247684181) (500,0.4256152093) (600,0.4263786574) (700,0.4238420512) (800,0.4200610109) (900,0.4151003824) (1000,0.4137229711) (1100,0.4091679186) (1200,0.4044880599) (1300,0.3992731035) (1400,0.3922104537) (1500,0.3862086535) (1600,0.3796085954) (1700,0.3749225497) (1800,0.3717579186) (1900,0.3703344375) (2000,0.3666954547) (2100,0.365011543) (2200,0.3622447491) (2300,0.3591279328) (2400,0.3565660983) (2500,0.3550324768) (2600,0.3547684431) (2700,0.352653712) (2800,0.3482556909) (2900,0.3486900002) (3000,0.3444100499) (3100,0.342010057) (3200,0.3436551154) (3300,0.3441766053) (3400,0.3428395987) (3500,0.3398250699) (3600,0.337252301) (3700,0.3351256311) (3800,0.3376740366) (3900,0.3316914827) (4000,0.3309938014) (4100,0.3273319781) (4200,0.3202670157) (4300,0.3173152566) (4400,0.3097055167) (4500,0.307425347) (4600,0.3017102599) (4700,0.2976942092) (4800,0.2906187922) (4900,0.2861673713) (5000,0.2815961242) (5100,0.2776623398) (5200,0.274797073) (5300,0.2680195361) (5400,0.2686319143) (5500,0.2655592829) (5600,0.2626097813) (5700,0.2581299886) (5800,0.2545400694) (5900,0.2531663552) (6000,0.2473713204) (6100,0.2441182405) (6200,0.2401635483) (6300,0.2347385213) (6400,0.2300070554) (6500,0.2242808044) (6600,0.2202034846) (6700,0.2162781328) (6800,0.2115834534) (6900,0.2055241272) (7000,0.2026569948) (7100,0.196919626) (7200,0.1957378805) (7300,0.1924003139) (7400,0.1856709614) (7500,0.1811211094) (7600,0.1760617182) (7700,0.170253706) (7800,0.1691711649) (7900,0.1654986367) (8000,0.160927929) (8100,0.1555277973) (8200,0.1479358882) (8300,0.142357789) (8400,0.1383317135) (8500,0.1340770245) (8600,0.1281224519) (8700,0.1222556233) (8800,0.1135705702) (8900,0.1069832988) (9000,0.1004957207) (9100,0.09461349659) (9200,0.0888347026) (9300,0.08468685858) (9400,0.08200710975) (9500,0.07349832058) (9600,0.06829044335) (9700,0.06760673188) (9800,0.06330379434) (9900,0.05896006562) (10000,0.05526050068) (10100,0.05262653008) (10200,0.04809877742) (10300,0.04174631983) (10400,0.03443596858) (10500,0.03242521074) (10600,0.02888924966) (10700,0.02202122157) (10800,0.01773933141) (10900,0.01410645448) (11000,0.0102719794) (11100,0.007107431995) (11200,0.003948762262) (11300,0.001447951415) (11400,0.0006329209777) (11500,-0.00255147235) (11600,-0.005162157747) (11700,-0.006093637296) (11800,-0.008143921918) (11900,-0.01025579094) (12000,-0.01204523426) (12100,-0.01418688064) (12200,-0.01501902444) (12300,-0.01517330064) (12400,-0.0170079242) (12500,-0.0178324251) (12600,-0.01835301802) (12700,-0.01894322066) (12800,-0.01947479602) (12900,-0.0194418909) (13000,-0.01928069778) (13100,-0.01867653597) (13200,-0.01963076908) (13300,-0.01980537456) (13400,-0.02038400844) (13500,-0.02129175086) (13600,-0.02195601184) (13700,-0.02314615864) (13800,-0.02411015127) (13900,-0.02519001551) (14000,-0.0251879679) (14100,-0.02637733649) (14200,-0.02588001247) (14300,-0.02621842604) (14400,-0.02635436039) (14500,-0.02572473064) (14600,-0.02550999597) (14700,-0.02537072804) (14800,-0.02545999531) (14900,-0.02557613458) (15000,-0.02598571964) (15100,-0.02477587312) (15200,-0.02538751382) (15300,-0.02662586123) (15400,-0.02596455589) (15500,-0.02590436433) (15600,-0.0259795595) (15700,-0.02599771526) (15800,-0.02616812941) (15900,-0.0272543475) (16000,-0.02801754046) (16100,-0.02866418529) (16200,-0.0293108508) (16300,-0.02853501514) (16400,-0.02921092473) (16500,-0.03002172764) (16600,-0.03009608835) (16700,-0.03070571218) (16800,-0.0300523743) (16900,-0.02908994239) (17000,-0.02929875944) (17100,-0.02967951521) (17200,-0.03005621098) (17300,-0.03032174278) (17400,-0.02961748075) (17500,-0.02942040842) (17600,-0.02933989931) (17700,-0.02928793952) (17800,-0.02956977226) (17900,-0.02930962276) (18000,-0.02807101291) (18100,-0.02743288539) (18200,-0.02631201744) (18300,-0.02642024439) (18400,-0.02675556708) (18500,-0.02670094464) (18600,-0.02664552331) (18700,-0.026318538) (18800,-0.02690149061) (18900,-0.02712387573) (19000,-0.0281627262) (19100,-0.02850026842) (19200,-0.02891151551) (19300,-0.0280553991) (19400,-0.02798539549) (19500,-0.02794721685) (19600,-0.02841402441) (19700,-0.02848106753) (19800,-0.02756571285) (19900,-0.0271426931) (20000,-0.02719767801) (20100,-0.02739798352) (20200,-0.02673300654) (20300,-0.02697359063) (20400,-0.02680099159) (20500,-0.02674905695) (20600,-0.02615641132) (20700,-0.02562785726) (20800,-0.02615235317) (20900,-0.02620393559) (21000,-0.0265357703) (21100,-0.02663391605) (21200,-0.02647878341) (21300,-0.02676184941) (21400,-0.02577589005) (21500,-0.0249399025) (21600,-0.02519647051) (21700,-0.02627948262) (21800,-0.02549109962) (21900,-0.02508776207) (22000,-0.02399206124) (22100,-0.02378194965) (22200,-0.02337918282) (22300,-0.02300750501) (22400,-0.02405860703) (22500,-0.02495039515) (22600,-0.02533413861) (22700,-0.02447005566) (22800,-0.02473194841) (22900,-0.02487575337) (23000,-0.02489760574) (23100,-0.02421425153) (23200,-0.0245764548) (23300,-0.02452032566) (23400,-0.02500131782) (23500,-0.02478645165) (23600,-0.02381535675) (23700,-0.02397013679) (23800,-0.02401576079)};
\addplot[pubRose, opacity=0.18, line width=0.35pt] coordinates {(100,0.04758103564) (200,0.04062645882) (300,0.04132807751) (400,0.04564781673) (500,0.04378821999) (600,0.04320324026) (700,0.04233950802) (800,0.04098733282) (900,0.04027277066) (1000,0.04013669379) (1100,0.04003353454) (1200,0.04061701559) (1300,0.03909032755) (1400,0.0367231328) (1500,0.0365677733) (1600,0.03611386046) (1700,0.0359109398) (1800,0.03567638397) (1900,0.03620361276) (2000,0.03596830443) (2100,0.03474713564) (2200,0.03455149606) (2300,0.03525944091) (2400,0.03552218899) (2500,0.03493136801) (2600,0.03516045995) (2700,0.03500433601) (2800,0.03491793051) (2900,0.03387110848) (3000,0.03309125677) (3100,0.03282322176) (3200,0.03157519717) (3300,0.02981954161) (3400,0.02805850077) (3500,0.02795073036) (3600,0.02746417914) (3700,0.02713571247) (3800,0.02749801148) (3900,0.02715074401) (4000,0.0274112422) (4100,0.02687992286) (4200,0.02703170665) (4300,0.02698279265) (4400,0.02715166658) (4500,0.02554314174) (4600,0.025071574) (4700,0.02452500723) (4800,0.02270493871) (4900,0.0220641925) (5000,0.02124413783) (5100,0.02082745684) (5200,0.02035865569) (5300,0.02172317645) (5400,0.02108161421) (5500,0.02212264454) (5600,0.02151764678) (5700,0.02102310164) (5800,0.02218018323) (5900,0.02179385135) (6000,0.02147873091) (6100,0.02196267871) (6200,0.02117544608) (6300,0.02128103385) (6400,0.02187839774) (6500,0.02205180349) (6600,0.02265202207) (6700,0.0222857791) (6800,0.02242915509) (6900,0.02205425659) (7000,0.02226417204) (7100,0.02136269184) (7200,0.02320772232) (7300,0.02129400354) (7400,0.02076312602) (7500,0.02007050142) (7600,0.01923337355) (7700,0.01970846094) (7800,0.01916458663) (7900,0.01961755278) (8000,0.0196017622) (8100,0.01970031532) (8200,0.01805463107) (8300,0.01906862892) (8400,0.02033921778) (8500,0.02132333592) (8600,0.0202466324) (8700,0.01977998093) (8800,0.01933634039) (8900,0.02079962622) (9000,0.01952593224) (9100,0.0193310936) (9200,0.0190422399) (9300,0.01954314318) (9400,0.01836010404) (9500,0.01815703921) (9600,0.01878678259) (9700,0.0189739028) (9800,0.01930312328) (9900,0.0189449368) (10000,0.01965376688) (10100,0.01925097061) (10200,0.02003605589) (10300,0.02022882998) (10400,0.02144638468) (10500,0.02131506279) (10600,0.02115217652) (10700,0.02157615423) (10800,0.0217359364) (10900,0.02162259202) (11000,0.0214991414) (11100,0.023141417) (11200,0.02313944539) (11300,0.02197381994) (11400,0.02139652418) (11500,0.02108241478) (11600,0.0222361139) (11700,0.02217490477) (11800,0.02146271933) (11900,0.02030175375) (12000,0.02123846635) (12100,0.01935781958) (12200,0.01927630501) (12300,0.01957869502) (12400,0.01964716827) (12500,0.0191786048) (12600,0.01890933802) (12700,0.01868488016) (12800,0.01919531599) (12900,0.01862739958) (13000,0.01824501045) (13100,0.01949250093) (13200,0.01854935186) (13300,0.01808621539) (13400,0.01815228174) (13500,0.01807103809) (13600,0.01699205702) (13700,0.01610499742) (13800,0.01556329299) (13900,0.01547383899) (14000,0.01503916448) (14100,0.01439709356) (14200,0.01493700547) (14300,0.01617321363) (14400,0.01461784393) (14500,0.01484965226) (14600,0.01507706568) (14700,0.01599201132) (14800,0.01711703027) (14900,0.01800956046) (15000,0.01731269192) (15100,0.01769246478) (15200,0.01725913212) (15300,0.01512824926) (15400,0.01572637162) (15500,0.01549911317) (15600,0.01551747085) (15700,0.01507854522) (15800,0.01443084949) (15900,0.01406309414) (16000,0.01477987585) (16100,0.0139177931) (16200,0.01431525811) (16300,0.01648840401) (16400,0.01666564206) (16500,0.01652887566) (16600,0.01659608651) (16700,0.01754708327) (16800,0.01760821175) (16900,0.01752177002) (17000,0.01773299454) (17100,0.01827452946) (17200,0.0180470678) (17300,0.01699547535) (17400,0.01613208652) (17500,0.01631994676) (17600,0.01633628402) (17700,0.01591558103) (17800,0.01540943142) (17900,0.01539157135) (18000,0.01471587839) (18100,0.01423717523) (18200,0.01345583685) (18300,0.01281713312) (18400,0.01365674678) (18500,0.01388464952) (18600,0.01364609702) (18700,0.01336517902) (18800,0.01354166241) (18900,0.01350429794) (19000,0.01439687246) (19100,0.01451665014) (19200,0.01554255337) (19300,0.01582460329) (19400,0.01602702271) (19500,0.01513931174) (19600,0.01565665668) (19700,0.01479836777) (19800,0.01382260253) (19900,0.01376130669) (20000,0.01256642533) (20100,0.0124901488) (20200,0.0126043716) (20300,0.01156394929) (20400,0.01066709459) (20500,0.01105289403) (20600,0.01062731445) (20700,0.01142284637) (20800,0.01289406163) (20900,0.01359570962) (21000,0.01297692202) (21100,0.01283176199) (21200,0.01251739853) (21300,0.01335409812) (21400,0.01337220713) (21500,0.01355060933) (21600,0.01202655321) (21700,0.01162797965) (21800,0.01015734475) (21900,0.01008492026) (22000,0.01099433396) (22100,0.01117561935) (22200,0.01024474553) (22300,0.01036909674) (22400,0.009931482142) (22500,0.009826201526) (22600,0.01090088277) (22700,0.01018932923) (22800,0.01107870662) (22900,0.01027939045) (23000,0.01091002054) (23100,0.0106525322) (23200,0.01209755922) (23300,0.01191179212) (23400,0.01220611222) (23500,0.01324763866) (23600,0.01358916797) (23700,0.01359546054) (23800,0.01376557471)};
\addplot[pubRose, opacity=0.18, line width=0.35pt] coordinates {(100,-0.1962490827) (200,-0.1914049983) (300,-0.1886890729) (400,-0.1891749576) (500,-0.1874208421) (600,-0.1866016388) (700,-0.1849067573) (800,-0.1823775619) (900,-0.1806579183) (1000,-0.1783957511) (1100,-0.1747552142) (1200,-0.1721931472) (1300,-0.170117487) (1400,-0.1652884498) (1500,-0.1622274473) (1600,-0.1590863332) (1700,-0.155343163) (1800,-0.1529428691) (1900,-0.1510924697) (2000,-0.1480440661) (2100,-0.1445830703) (2200,-0.1414498597) (2300,-0.1362129055) (2400,-0.1343423754) (2500,-0.1313909583) (2600,-0.1271183856) (2700,-0.1252028063) (2800,-0.121285826) (2900,-0.1170113049) (3000,-0.1149144255) (3100,-0.1115322851) (3200,-0.1081339419) (3300,-0.1073538363) (3400,-0.104873284) (3500,-0.1027199656) (3600,-0.1017255984) (3700,-0.09713901728) (3800,-0.09631182849) (3900,-0.09397066012) (4000,-0.09160585105) (4100,-0.0898593016) (4200,-0.08983075544) (4300,-0.08726355806) (4400,-0.08560151979) (4500,-0.08333234116) (4600,-0.08134884313) (4700,-0.08154512867) (4800,-0.07846641429) (4900,-0.07727732472) (5000,-0.0756334845) (5100,-0.07450711839) (5200,-0.07083025947) (5300,-0.06904181726) (5400,-0.06531071179) (5500,-0.06262244135) (5600,-0.0596890308) (5700,-0.05821543075) (5800,-0.05618342087) (5900,-0.05356792584) (6000,-0.05206253193) (6100,-0.0492656149) (6200,-0.04825702645) (6300,-0.04667758532) (6400,-0.04569137841) (6500,-0.04430973306) (6600,-0.04283004925) (6700,-0.04094434604) (6800,-0.03960631229) (6900,-0.03781650085) (7000,-0.03683139142) (7100,-0.03723264094) (7200,-0.03519785143) (7300,-0.03453111351) (7400,-0.03363925032) (7500,-0.03265270162) (7600,-0.0311916694) (7700,-0.02953975573) (7800,-0.02914324533) (7900,-0.0289031541) (8000,-0.02715995349) (8100,-0.02565008141) (8200,-0.02397061661) (8300,-0.02209339291) (8400,-0.02112258673) (8500,-0.02038697675) (8600,-0.02026988026) (8700,-0.01934940778) (8800,-0.01743957298) (8900,-0.01613325756) (9000,-0.01448644339) (9100,-0.0125710682) (9200,-0.01252706526) (9300,-0.01153444806) (9400,-0.01048242706) (9500,-0.008296815259) (9600,-0.006432267171) (9700,-0.005293122618) (9800,-0.005720609645) (9900,-0.004654489548) (10000,-0.003564679457) (10100,-0.002934005155) (10200,-0.0007532367923) (10300,-0.0003760625816) (10400,0.001021455534) (10500,0.0009076142011) (10600,0.002009437571) (10700,0.003235278663) (10800,0.004928473488) (10900,0.00539919713) (11000,0.00610911469) (11100,0.006316750008) (11200,0.005520306621) (11300,0.006328667072) (11400,0.006211335375) (11500,0.006202817708) (11600,0.00659865085) (11700,0.007129569352) (11800,0.007493318804) (11900,0.008770391438) (12000,0.008864283655) (12100,0.009314974165) (12200,0.01053669865) (12300,0.01108902057) (12400,0.01319127448) (12500,0.0144535277) (12600,0.01499615749) (12700,0.01551379533) (12800,0.01550376015) (12900,0.0156425911) (13000,0.0162695419) (13100,0.01725669485) (13200,0.01716438476) (13300,0.01801777864) (13400,0.0175684643) (13500,0.01781153437) (13600,0.01769218203) (13700,0.01757680923) (13800,0.01805946948) (13900,0.01787659796) (14000,0.01805374296) (14100,0.01884949626) (14200,0.0196455108) (14300,0.02016343661) (14400,0.02029046565) (14500,0.02119545136) (14600,0.02204814609) (14700,0.02333326694) (14800,0.02415978704) (14900,0.0247379072) (15000,0.02589091267) (15100,0.02566679344) (15200,0.02665374577) (15300,0.02697741408) (15400,0.02682998385) (15500,0.02734243721) (15600,0.02622523885) (15700,0.02548452225) (15800,0.02572193332) (15900,0.02633966561) (16000,0.02619144283) (16100,0.02631780803) (16200,0.02614330556) (16300,0.02637090851) (16400,0.02581749819) (16500,0.02489715312) (16600,0.02641327046) (16700,0.02664513048) (16800,0.02670523617) (16900,0.02682979554) (17000,0.02707464173) (17100,0.02740351409) (17200,0.0274764603) (17300,0.02725700252) (17400,0.02790038437) (17500,0.02834500726) (17600,0.02826540805) (17700,0.02913737651) (17800,0.02966187764) (17900,0.0293266343) (18000,0.02786789276) (18100,0.02792656757) (18200,0.02785304748) (18300,0.02837099023) (18400,0.02950131483) (18500,0.02970044632) (18600,0.02967902366) (18700,0.02875127178) (18800,0.02951113191) (18900,0.0293254707) (19000,0.03116771188) (19100,0.03127839528) (19200,0.03116274439) (19300,0.03119902909) (19400,0.03136232048) (19500,0.03142391667) (19600,0.03114922531) (19700,0.03271630406) (19800,0.03233497664) (19900,0.03260668386) (20000,0.03167247344) (20100,0.03130195532) (20200,0.03152857516) (20300,0.03115982376) (20400,0.02995361704) (20500,0.02977346983) (20600,0.02970895935) (20700,0.02801672481) (20800,0.02703305222) (20900,0.02834871057) (21000,0.02841324694) (21100,0.02909557447) (21200,0.02899128404) (21300,0.02975476496) (21400,0.02963916492) (21500,0.03044263795) (21600,0.03068897519) (21700,0.03245583922) (21800,0.033928027) (21900,0.03313417621) (22000,0.03268554136) (22100,0.0316045247) (22200,0.03192004357) (22300,0.03261229228) (22400,0.0335823549) (22500,0.03204213008) (22600,0.03215502147) (22700,0.0313122062) (22800,0.03022684436) (22900,0.02922930811) (23000,0.03007035162) (23100,0.03127525095) (23200,0.03092553485) (23300,0.02844673321) (23400,0.02827161886) (23500,0.02869626544) (23600,0.02811844461) (23700,0.02768799514) (23800,0.02840571553)};
\addplot[pubRose, opacity=0.18, line width=0.35pt] coordinates {(100,0.1713679731) (200,0.1722728834) (300,0.1728087167) (400,0.1741381958) (500,0.1734680951) (600,0.1725341827) (700,0.1709621911) (800,0.1692236084) (900,0.1685924861) (1000,0.1673579335) (1100,0.1659334779) (1200,0.1649595872) (1300,0.163387306) (1400,0.1601996943) (1500,0.1584188506) (1600,0.1572491094) (1700,0.1555602029) (1800,0.1553656533) (1900,0.1536150396) (2000,0.1523830354) (2100,0.1510914758) (2200,0.1492299393) (2300,0.1478035972) (2400,0.1466643617) (2500,0.1466128573) (2600,0.145065026) (2700,0.1457837492) (2800,0.1448895916) (2900,0.145646961) (3000,0.1466239288) (3100,0.145978789) (3200,0.1454353735) (3300,0.1454627097) (3400,0.1456679538) (3500,0.1449041501) (3600,0.1449939132) (3700,0.1440840706) (3800,0.1444656149) (3900,0.1435523674) (4000,0.1397495553) (4100,0.1398599163) (4200,0.1386944398) (4300,0.1375114396) (4400,0.1378573626) (4500,0.1365434423) (4600,0.1338201888) (4700,0.1326779969) (4800,0.1299997255) (4900,0.1263315111) (5000,0.1275288716) (5100,0.1270255044) (5200,0.1270112514) (5300,0.1254818894) (5400,0.1238968514) (5500,0.1214810312) (5600,0.1227680065) (5700,0.1221655123) (5800,0.1208326265) (5900,0.1216983326) (6000,0.1204055183) (6100,0.1182470545) (6200,0.1172361381) (6300,0.1166845098) (6400,0.1151880346) (6500,0.1151582435) (6600,0.1142238773) (6700,0.1122956976) (6800,0.1109849893) (6900,0.1107587948) (7000,0.1100560829) (7100,0.1069792248) (7200,0.1054277986) (7300,0.1032437608) (7400,0.1039901942) (7500,0.1039663248) (7600,0.1018662252) (7700,0.1030391455) (7800,0.1035610043) (7900,0.1023113891) (8000,0.1007976457) (8100,0.1029214777) (8200,0.1022315614) (8300,0.1023265168) (8400,0.1005158536) (8500,0.09777283818) (8600,0.09875453785) (8700,0.09532766491) (8800,0.0929743059) (8900,0.09135745764) (9000,0.09103350267) (9100,0.08749420941) (9200,0.08595116585) (9300,0.08406428844) (9400,0.08284590468) (9500,0.08571407497) (9600,0.08356666788) (9700,0.08361735195) (9800,0.08304921612) (9900,0.08210066035) (10000,0.08071240634) (10100,0.08273268119) (10200,0.08290991858) (10300,0.0820616588) (10400,0.0798525162) (10500,0.07547739446) (10600,0.07547663525) (10700,0.07387934476) (10800,0.07416115552) (10900,0.0740051724) (11000,0.07293963134) (11100,0.06996238939) (11200,0.06832213365) (11300,0.06829464845) (11400,0.0687099684) (11500,0.06840620972) (11600,0.06590024419) (11700,0.06485963836) (11800,0.0630907245) (11900,0.06294677109) (12000,0.06230323985) (12100,0.0623807855) (12200,0.06191109307) (12300,0.06200751849) (12400,0.0622258123) (12500,0.06179223135) (12600,0.06202273145) (12700,0.0611514207) (12800,0.06213747151) (12900,0.06079337783) (13000,0.05993780233) (13100,0.05889911763) (13200,0.05770226456) (13300,0.0562888708) (13400,0.05300340056) (13500,0.0523723051) (13600,0.05027394928) (13700,0.04946895391) (13800,0.04637034722) (13900,0.04561315328) (14000,0.04629482068) (14100,0.04720902964) (14200,0.04687422216) (14300,0.04629968815) (14400,0.04729210697) (14500,0.04522752054) (14600,0.04560996704) (14700,0.04575868659) (14800,0.04547119066) (14900,0.04369328283) (15000,0.04290195778) (15100,0.04164821357) (15200,0.04062794037) (15300,0.04014221653) (15400,0.03850375824) (15500,0.03882427998) (15600,0.03759940639) (15700,0.03736994602) (15800,0.03748237714) (15900,0.03726247661) (16000,0.03553565852) (16100,0.03425096385) (16200,0.03443615548) (16300,0.03445247896) (16400,0.03493698351) (16500,0.03484649956) (16600,0.03442021962) (16700,0.03398401421) (16800,0.03398775253) (16900,0.03360217717) (17000,0.03320003916) (17100,0.03384117354) (17200,0.0326960396) (17300,0.0313765686) (17400,0.03086090796) (17500,0.02922806088) (17600,0.02870820295) (17700,0.02870701756) (17800,0.02755004223) (17900,0.02771688886) (18000,0.02775218077) (18100,0.02561821807) (18200,0.024449444) (18300,0.02409340143) (18400,0.02273984216) (18500,0.02394170575) (18600,0.02373897545) (18700,0.02229464073) (18800,0.02242066823) (18900,0.02164552175) (19000,0.02031717068) (19100,0.02151321964) (19200,0.02268661791) (19300,0.02228541235) (19400,0.02216327814) (19500,0.02108942429) (19600,0.02133152587) (19700,0.02170088133) (19800,0.02128745755) (19900,0.02157724025) (20000,0.02129101595) (20100,0.02011632295) (20200,0.01942354022) (20300,0.02001137054) (20400,0.02008586666) (20500,0.02106833281) (20600,0.01952194255) (20700,0.01879876256) (20800,0.0182351226) (20900,0.01674110563) (21000,0.01695006276) (21100,0.01575160269) (21200,0.0158641275) (21300,0.01509716362) (21400,0.015489012) (21500,0.01399417929) (21600,0.01356685315) (21700,0.01359561929) (21800,0.01363622849) (21900,0.01337797195) (22000,0.01322632572) (22100,0.01408752557) (22200,0.01349802716) (22300,0.01420797231) (22400,0.01377979079) (22500,0.015385604) (22600,0.01614593253) (22700,0.0155497801) (22800,0.01546693402) (22900,0.01717031794) (23000,0.01756521622) (23100,0.01720651267) (23200,0.01758846566) (23300,0.01666977359) (23400,0.01730189687) (23500,0.01635573683) (23600,0.01666521607) (23700,0.01646376066) (23800,0.01653789636)};
\addplot[pubRose, opacity=0.18, line width=0.35pt] coordinates {(100,0.1922287494) (200,0.1893916726) (300,0.1878062238) (400,0.1873190403) (500,0.1852892399) (600,0.1837077414) (700,0.1811610375) (800,0.1786677334) (900,0.1766707285) (1000,0.1753139257) (1100,0.173895435) (1200,0.1707003847) (1300,0.1698982611) (1400,0.1692705542) (1500,0.1694383875) (1600,0.1699121743) (1700,0.171127592) (1800,0.1723191381) (1900,0.1742843628) (2000,0.1765946627) (2100,0.1754199892) (2200,0.1775321871) (2300,0.1764502272) (2400,0.1769373968) (2500,0.1771448702) (2600,0.176092273) (2700,0.1740880519) (2800,0.1726735845) (2900,0.1716309816) (3000,0.17042) (3100,0.1703469038) (3200,0.169559668) (3300,0.1700534746) (3400,0.1662157014) (3500,0.1650577947) (3600,0.1639185458) (3700,0.163746421) (3800,0.1632775292) (3900,0.162636295) (4000,0.1611130551) (4100,0.1614148498) (4200,0.1600106031) (4300,0.1578373685) (4400,0.1575808689) (4500,0.1543874919) (4600,0.1529597923) (4700,0.1500047401) (4800,0.1482266411) (4900,0.1459055319) (5000,0.1436513126) (5100,0.1400671527) (5200,0.1390579313) (5300,0.1376319125) (5400,0.1370266795) (5500,0.1368729666) (5600,0.136027573) (5700,0.1353189655) (5800,0.1347222723) (5900,0.1335180201) (6000,0.132061293) (6100,0.129542277) (6200,0.1249388091) (6300,0.1221202001) (6400,0.1204528905) (6500,0.1187266834) (6600,0.1163032211) (6700,0.114844849) (6800,0.112360175) (6900,0.1099892862) (7000,0.1065160356) (7100,0.1064044856) (7200,0.1065394253) (7300,0.1052849546) (7400,0.1030150592) (7500,0.1029392898) (7600,0.101304163) (7700,0.1002968192) (7800,0.09866443798) (7900,0.09774464965) (8000,0.09693656936) (8100,0.09490656778) (8200,0.09303045943) (8300,0.09047416374) (8400,0.08846010193) (8500,0.08370560557) (8600,0.08098319024) (8700,0.07879662886) (8800,0.07603529431) (8900,0.07271131203) (9000,0.07171817273) (9100,0.06923775077) (9200,0.06668511443) (9300,0.06510494687) (9400,0.06383802108) (9500,0.06254097298) (9600,0.06090991944) (9700,0.057620373) (9800,0.05611752383) (9900,0.05469129719) (10000,0.05191990435) (10100,0.0504952196) (10200,0.04913602658) (10300,0.04758217596) (10400,0.04483134151) (10500,0.04386952072) (10600,0.04192455746) (10700,0.04092885405) (10800,0.03867777232) (10900,0.03720441666) (11000,0.03599267323) (11100,0.03443514537) (11200,0.03276183028) (11300,0.03165844902) (11400,0.03121986538) (11500,0.02778065009) (11600,0.02772475732) (11700,0.02706229212) (11800,0.02520745648) (11900,0.02519793538) (12000,0.02471609404) (12100,0.0232099642) (12200,0.02206500713) (12300,0.02098489841) (12400,0.01900955262) (12500,0.01788384542) (12600,0.01597933024) (12700,0.01440600865) (12800,0.01421754023) (12900,0.01230827381) (13000,0.01015129879) (13100,0.007755013718) (13200,0.00749466361) (13300,0.006801756728) (13400,0.00576454259) (13500,0.006327113486) (13600,0.005484696035) (13700,0.00498051513) (13800,0.00433490579) (13900,0.003304120433) (14000,0.003212585475) (14100,0.002794677939) (14200,0.001282870758) (14300,-0.0003812958603) (14400,-0.0005061065429) (14500,-0.002336003503) (14600,-0.002993961202) (14700,-0.004320866277) (14800,-0.004109324032) (14900,-0.005361523136) (15000,-0.006508602685) (15100,-0.005782928009) (15200,-0.006360990618) (15300,-0.007062982459) (15400,-0.008546274732) (15500,-0.009306866134) (15600,-0.009806920093) (15700,-0.009828704101) (15800,-0.01008457839) (15900,-0.01014434993) (16000,-0.01057016365) (16100,-0.01167626404) (16200,-0.01163100298) (16300,-0.01149922502) (16400,-0.01203277153) (16500,-0.01254001693) (16600,-0.01383077339) (16700,-0.01432194686) (16800,-0.01535397666) (16900,-0.01656838767) (17000,-0.01642936263) (17100,-0.01688452093) (17200,-0.01702178409) (17300,-0.01803594502) (17400,-0.01761145489) (17500,-0.01820391184) (17600,-0.01713408669) (17700,-0.01809686143) (17800,-0.01851151222) (17900,-0.01851775469) (18000,-0.01952369157) (18100,-0.01994130928) (18200,-0.02025636518) (18300,-0.02062904043) (18400,-0.02072702637) (18500,-0.0199597328) (18600,-0.02180705173) (18700,-0.02217857605) (18800,-0.02281423034) (18900,-0.02341073183) (19000,-0.02315927604) (19100,-0.02344735386) (19200,-0.02499336433) (19300,-0.02472130377) (19400,-0.02539001517) (19500,-0.02514367234) (19600,-0.02486021314) (19700,-0.02438323069) (19800,-0.02374472376) (19900,-0.02383617628) (20000,-0.02444639225) (20100,-0.02400063798) (20200,-0.02386836018) (20300,-0.02439127751) (20400,-0.02439815365) (20500,-0.02443917282) (20600,-0.02480172496) (20700,-0.02510423325) (20800,-0.02619228624) (20900,-0.02622452825) (21000,-0.02542403843) (21100,-0.02514774241) (21200,-0.02468183674) (21300,-0.02449719813) (21400,-0.02495642714) (21500,-0.0252113685) (21600,-0.02498022094) (21700,-0.02462799214) (21800,-0.02471593097) (21900,-0.0250165116) (22000,-0.02576671802) (22100,-0.02592234444) (22200,-0.02676976006) (22300,-0.02682210449) (22400,-0.02699370645) (22500,-0.0273537755) (22600,-0.02786843348) (22700,-0.02786782719) (22800,-0.02734227069) (22900,-0.02708104234) (23000,-0.0276179783) (23100,-0.02841675375) (23200,-0.02818356492) (23300,-0.02820995934) (23400,-0.02897342071) (23500,-0.0297468232) (23600,-0.02910502236) (23700,-0.02965416424) (23800,-0.02982000802)};
\addplot[pubRose, opacity=0.18, line width=0.35pt] coordinates {(100,0.4900595844) (200,0.4904662669) (300,0.4827229977) (400,0.4830254689) (500,0.476803571) (600,0.4701071084) (700,0.4654791185) (800,0.4615025483) (900,0.4573277997) (1000,0.4508665383) (1100,0.442710191) (1200,0.4338734597) (1300,0.4258722603) (1400,0.4185847372) (1500,0.4123116076) (1600,0.4070812076) (1700,0.3994690657) (1800,0.3943040818) (1900,0.3913939625) (2000,0.390647912) (2100,0.3882522613) (2200,0.3830246717) (2300,0.3801846385) (2400,0.3753444552) (2500,0.3732741445) (2600,0.3691548258) (2700,0.3705307633) (2800,0.3679948896) (2900,0.3642116785) (3000,0.3618213773) (3100,0.3569877684) (3200,0.3558035105) (3300,0.3530224562) (3400,0.3492308617) (3500,0.3439811438) (3600,0.3442333162) (3700,0.3377341479) (3800,0.3363304824) (3900,0.3316399693) (4000,0.3238173097) (4100,0.3203259528) (4200,0.3161920816) (4300,0.3116656512) (4400,0.30954386) (4500,0.3075721741) (4600,0.3018734783) (4700,0.2988776803) (4800,0.2910319775) (4900,0.2862313271) (5000,0.2849899471) (5100,0.2809000164) (5200,0.2786564052) (5300,0.2774626106) (5400,0.2708064914) (5500,0.2670381367) (5600,0.2649268121) (5700,0.2638370842) (5800,0.2614646196) (5900,0.2601678729) (6000,0.2578242406) (6100,0.2547036827) (6200,0.2485930011) (6300,0.2439857066) (6400,0.2432303458) (6500,0.2369937479) (6600,0.2344448894) (6700,0.2271667361) (6800,0.2245801374) (6900,0.2188685015) (7000,0.2112337604) (7100,0.2045670614) (7200,0.2012522399) (7300,0.1960530192) (7400,0.1886475369) (7500,0.1891122237) (7600,0.1814882025) (7700,0.1787828311) (7800,0.1732482851) (7900,0.1704879612) (8000,0.1695070922) (8100,0.169607763) (8200,0.1624691948) (8300,0.1594242766) (8400,0.1599174887) (8500,0.1566097453) (8600,0.1533313513) (8700,0.1477664106) (8800,0.1420664176) (8900,0.1361600734) (9000,0.1311219051) (9100,0.1269146122) (9200,0.1276120462) (9300,0.1229680851) (9400,0.1166469298) (9500,0.1103308931) (9600,0.1081487574) (9700,0.1075055033) (9800,0.1065859653) (9900,0.1044015139) (10000,0.1011954948) (10100,0.0979093641) (10200,0.09330702499) (10300,0.09112637714) (10400,0.09001767263) (10500,0.08510770984) (10600,0.08154809438) (10700,0.07795078047) (10800,0.07277431525) (10900,0.06983719692) (11000,0.06798732318) (11100,0.06453033835) (11200,0.06178816296) (11300,0.05921935216) (11400,0.05435871445) (11500,0.05278914422) (11600,0.04955458269) (11700,0.04535918962) (11800,0.04517039042) (11900,0.04394821208) (12000,0.04251120631) (12100,0.04074305985) (12200,0.0398822831) (12300,0.03829714414) (12400,0.0380600499) (12500,0.03643145505) (12600,0.0351493651) (12700,0.0350899417) (12800,0.03318673261) (12900,0.02976759695) (13000,0.02764361226) (13100,0.02571966876) (13200,0.0230627324) (13300,0.02117493651) (13400,0.01924160491) (13500,0.01833690866) (13600,0.0166672776) (13700,0.01513250661) (13800,0.0146480422) (13900,0.01542535769) (14000,0.01336129415) (14100,0.01267025336) (14200,0.01314273668) (14300,0.01280324222) (14400,0.01261140513) (14500,0.01193810813) (14600,0.01104618949) (14700,0.01046110061) (14800,0.01048530503) (14900,0.00968324081) (15000,0.01011001123) (15100,0.009787211102) (15200,0.009102521418) (15300,0.008501206944) (15400,0.006938736187) (15500,0.005546546169) (15600,0.00618773005) (15700,0.005971068796) (15800,0.004591450829) (15900,0.004300947429) (16000,0.004735802056) (16100,0.003940801893) (16200,0.003740132216) (16300,0.003309133311) (16400,0.004350336397) (16500,0.005552760779) (16600,0.004705768113) (16700,0.004135446745) (16800,0.003661943011) (16900,0.002888765785) (17000,0.001567626098) (17100,0.001484501775) (17200,0.0003638856284) (17300,0.001142779522) (17400,5.434690356e-05) (17500,-0.001275954262) (17600,-0.001508320807) (17700,-0.0008704594598) (17800,-0.0003450318734) (17900,-0.0002481008676) (18000,-0.0005810606672) (18100,-0.0005975864158) (18200,5.823747197e-05) (18300,-0.0003897957416) (18400,0.0001577301504) (18500,0.0007798487757) (18600,0.001102067227) (18700,4.123242106e-05) (18800,-0.0008255505236) (18900,-0.0006456536677) (19000,-0.0007122639538) (19100,-0.0001263390033) (19200,-0.0001514897071) (19300,-0.0003189505398) (19400,-0.001145073641) (19500,-0.001313174792) (19600,-0.002139486044) (19700,-0.001908713588) (19800,-0.001674434086) (19900,-0.0009650071501) (20000,-0.0005065078221) (20100,-0.001961696598) (20200,-0.001607561084) (20300,-0.001887520305) (20400,-0.0007835817945) (20500,-0.001032918222) (20600,-0.0006685855114) (20700,-0.0009128340651) (20800,-0.001458368184) (20900,-0.003148605269) (21000,-0.003565433872) (21100,-0.001985635207) (21200,-0.002057837212) (21300,-0.002921591472) (21400,-0.004284541641) (21500,-0.004192613909) (21600,-0.003251928673) (21700,-0.003107195371) (21800,-0.001631054957) (21900,-0.000980097102) (22000,0.0003416045569) (22100,-0.0001848201937) (22200,3.852453374e-05) (22300,0.0005167961615) (22400,0.0004244417913) (22500,0.0007456203137) (22600,-0.0007871024165) (22700,-0.0008193624759) (22800,-0.001503047542) (22900,-0.0008818542206) (23000,-0.002184400344) (23100,-0.001926048275) (23200,-0.00253303264) (23300,-0.002228064113) (23400,-0.002640914661) (23500,-0.002480830191) (23600,-0.00134038947) (23700,7.830573304e-05) (23800,0.000774620325)};
\addplot[pubRose, opacity=0.18, line width=0.35pt] coordinates {(100,-0.1712110788) (200,-0.1712435186) (300,-0.1720325698) (400,-0.1733249016) (500,-0.1706788719) (600,-0.1691518525) (700,-0.1674386795) (800,-0.1651591547) (900,-0.1627222051) (1000,-0.1609442189) (1100,-0.1585007831) (1200,-0.1559180498) (1300,-0.1517521411) (1400,-0.1475895211) (1500,-0.1447649017) (1600,-0.1418033436) (1700,-0.1383261301) (1800,-0.1365552925) (1900,-0.1357937641) (2000,-0.133499749) (2100,-0.1312838733) (2200,-0.1290218003) (2300,-0.1279378876) (2400,-0.1261352979) (2500,-0.1249633238) (2600,-0.1227920517) (2700,-0.1211964905) (2800,-0.1190674514) (2900,-0.1164387092) (3000,-0.1149011359) (3100,-0.1121876501) (3200,-0.1103813149) (3300,-0.1079452991) (3400,-0.1052979365) (3500,-0.1031490535) (3600,-0.1017737232) (3700,-0.1000898294) (3800,-0.098303435) (3900,-0.09587500021) (4000,-0.09301533848) (4100,-0.09074126929) (4200,-0.0884153299) (4300,-0.08602867573) (4400,-0.08505315036) (4500,-0.08205241635) (4600,-0.08070579767) (4700,-0.07808210477) (4800,-0.07550854683) (4900,-0.07452278212) (5000,-0.07365119979) (5100,-0.07286270857) (5200,-0.0722551018) (5300,-0.07137570754) (5400,-0.06944326386) (5500,-0.06919021308) (5600,-0.06737883985) (5700,-0.06613531895) (5800,-0.0647803124) (5900,-0.06176273823) (6000,-0.05886927955) (6100,-0.05700811706) (6200,-0.05405277796) (6300,-0.05247508883) (6400,-0.05136547238) (6500,-0.04966378622) (6600,-0.04764719419) (6700,-0.04682218432) (6800,-0.04688337371) (6900,-0.04704657756) (7000,-0.04753072672) (7100,-0.04673983753) (7200,-0.04620395862) (7300,-0.04547880366) (7400,-0.04360388555) (7500,-0.0424593959) (7600,-0.04166963361) (7700,-0.04088832214) (7800,-0.03933802247) (7900,-0.03760000616) (8000,-0.03717309386) (8100,-0.03704553284) (8200,-0.03646476865) (8300,-0.03603214808) (8400,-0.03585073594) (8500,-0.03489131089) (8600,-0.03425832298) (8700,-0.03347473592) (8800,-0.03325717263) (8900,-0.03302654084) (9000,-0.03139710594) (9100,-0.03091567252) (9200,-0.02995324209) (9300,-0.02832222842) (9400,-0.02770345435) (9500,-0.02702307161) (9600,-0.02662713844) (9700,-0.02593883518) (9800,-0.0238094206) (9900,-0.02347927047) (10000,-0.0222715375) (10100,-0.02031493103) (10200,-0.01873035356) (10300,-0.01873165537) (10400,-0.01789563186) (10500,-0.01748219747) (10600,-0.01708037015) (10700,-0.01723346505) (10800,-0.01678984463) (10900,-0.0167009186) (11000,-0.0171344731) (11100,-0.01768893339) (11200,-0.01862080814) (11300,-0.0189957724) (11400,-0.01899422416) (11500,-0.01908018114) (11600,-0.01835770616) (11700,-0.01785055054) (11800,-0.01865883488) (11900,-0.01760945823) (12000,-0.01738041434) (12100,-0.01673910683) (12200,-0.01651774133) (12300,-0.01488517383) (12400,-0.01491663856) (12500,-0.01373346057) (12600,-0.01269604042) (12700,-0.01237513172) (12800,-0.01201236127) (12900,-0.01157745887) (13000,-0.01068770094) (13100,-0.01057492383) (13200,-0.009419586137) (13300,-0.009622299019) (13400,-0.008894796856) (13500,-0.009137834981) (13600,-0.009541061427) (13700,-0.009366671834) (13800,-0.008626731113) (13900,-0.008858382422) (14000,-0.00958783729) (14100,-0.009284918196) (14200,-0.01003219737) (14300,-0.009813418472) (14400,-0.01012282926) (14500,-0.01036123619) (14600,-0.01054170127) (14700,-0.00967755036) (14800,-0.009766092687) (14900,-0.009618873126) (15000,-0.008805839508) (15100,-0.008945339057) (15200,-0.008602255187) (15300,-0.008959322539) (15400,-0.009284626902) (15500,-0.009308886644) (15600,-0.009157915809) (15700,-0.01028438359) (15800,-0.01085354406) (15900,-0.01022822312) (16000,-0.01068403893) (16100,-0.0104340659) (16200,-0.009977160813) (16300,-0.008751105075) (16400,-0.008698098199) (16500,-0.008494935301) (16600,-0.007642861223) (16700,-0.006362693955) (16800,-0.005884210684) (16900,-0.005896016757) (17000,-0.005753897119) (17100,-0.005716374598) (17200,-0.005445238075) (17300,-0.006719917513) (17400,-0.0061613595) (17500,-0.005654919834) (17600,-0.006922765833) (17700,-0.006913580466) (17800,-0.007667722832) (17900,-0.008347623772) (18000,-0.008070410904) (18100,-0.008107986511) (18200,-0.008334073192) (18300,-0.00773022077) (18400,-0.007710929448) (18500,-0.00830156412) (18600,-0.008186192065) (18700,-0.008901265031) (18800,-0.007667992869) (18900,-0.00628409693) (19000,-0.006184265064) (19100,-0.00627962579) (19200,-0.005849813536) (19300,-0.005572019127) (19400,-0.005161523103) (19500,-0.005116623949) (19600,-0.004510748509) (19700,-0.003653655606) (19800,-0.003132623222) (19900,-0.004560387621) (20000,-0.00482071561) (20100,-0.004608721306) (20200,-0.005974538639) (20300,-0.0062130412) (20400,-0.00628421492) (20500,-0.006072741971) (20600,-0.006459540539) (20700,-0.006654158758) (20800,-0.007487540762) (20900,-0.007385412115) (21000,-0.007087115175) (21100,-0.006899152533) (21200,-0.006777774752) (21300,-0.007611660403) (21400,-0.007962301071) (21500,-0.008135723905) (21600,-0.007385016535) (21700,-0.007951155584) (21800,-0.00708792709) (21900,-0.007669063902) (22000,-0.00729543713) (22100,-0.007577070473) (22200,-0.00653517325) (22300,-0.006133178943) (22400,-0.00595946549) (22500,-0.006294034277) (22600,-0.006949439408) (22700,-0.007457742603) (22800,-0.008196774808) (22900,-0.007126666869) (23000,-0.007918623113) (23100,-0.007341188518) (23200,-0.007721099048) (23300,-0.007070223219) (23400,-0.006577142468) (23500,-0.006129013142) (23600,-0.005109015456) (23700,-0.00477875385) (23800,-0.00477287292)};
\addplot[pubRose, opacity=0.18, line width=0.35pt] coordinates {(100,0.2402989268) (200,0.2444701716) (300,0.2421536644) (400,0.2428329736) (500,0.2433582634) (600,0.2415722633) (700,0.2390530684) (800,0.2373224236) (900,0.2352741741) (1000,0.2354653299) (1100,0.2329679057) (1200,0.2299772784) (1300,0.2297445238) (1400,0.2268784359) (1500,0.2233047783) (1600,0.2207436174) (1700,0.2193581536) (1800,0.2190347403) (1900,0.2190070838) (2000,0.2161388308) (2100,0.2146742091) (2200,0.2116422758) (2300,0.2082344458) (2400,0.2062791094) (2500,0.2055984259) (2600,0.204045099) (2700,0.2032697126) (2800,0.2024608001) (2900,0.1997311279) (3000,0.1977476403) (3100,0.1954075068) (3200,0.193949236) (3300,0.1910999194) (3400,0.1899454117) (3500,0.1887179792) (3600,0.1868007272) (3700,0.1850187376) (3800,0.1819953948) (3900,0.1825219244) (4000,0.179935728) (4100,0.1794103295) (4200,0.181253843) (4300,0.1810838073) (4400,0.1801149637) (4500,0.178281267) (4600,0.1784700558) (4700,0.1788649902) (4800,0.1751554325) (4900,0.1728396043) (5000,0.1733366638) (5100,0.173000662) (5200,0.1694008917) (5300,0.1677796856) (5400,0.1654461056) (5500,0.1635211065) (5600,0.1601245791) (5700,0.1562975481) (5800,0.156486626) (5900,0.1530264601) (6000,0.1496905819) (6100,0.1461885169) (6200,0.1434158519) (6300,0.1393648148) (6400,0.134681996) (6500,0.1320835449) (6600,0.1313648604) (6700,0.1290739544) (6800,0.1249194041) (6900,0.1237513885) (7000,0.1221447915) (7100,0.1195692003) (7200,0.1189549357) (7300,0.1165881008) (7400,0.1185445003) (7500,0.1168282509) (7600,0.1130350806) (7700,0.1103513606) (7800,0.1094769113) (7900,0.1069834098) (8000,0.1038688347) (8100,0.1027436018) (8200,0.0986622721) (8300,0.09914736822) (8400,0.09720000923) (8500,0.09489886835) (8600,0.09479306638) (8700,0.09296799302) (8800,0.09179774672) (8900,0.08836932927) (9000,0.08675037026) (9100,0.0836764127) (9200,0.08356244862) (9300,0.0807385698) (9400,0.07717306092) (9500,0.07583742812) (9600,0.07329960465) (9700,0.07167536318) (9800,0.06744551547) (9900,0.06625774056) (10000,0.06462597884) (10100,0.06469833739) (10200,0.0622520268) (10300,0.06213851124) (10400,0.06106877327) (10500,0.0589012634) (10600,0.05763925686) (10700,0.05488121659) (10800,0.05509592853) (10900,0.05482418798) (11000,0.05379138887) (11100,0.0522675585) (11200,0.05268162787) (11300,0.0503353022) (11400,0.04741530381) (11500,0.0474715393) (11600,0.04504511878) (11700,0.0443915464) (11800,0.04385956451) (11900,0.04182731025) (12000,0.04117166586) (12100,0.03996950127) (12200,0.03574256655) (12300,0.0345517002) (12400,0.03374773134) (12500,0.03122216165) (12600,0.02964366786) (12700,0.02886246182) (12800,0.02904638797) (12900,0.02822419554) (13000,0.02686927095) (13100,0.02550161108) (13200,0.02569273058) (13300,0.02476528622) (13400,0.0243504677) (13500,0.02285192255) (13600,0.02340547275) (13700,0.02339982111) (13800,0.01981939594) (13900,0.01874127188) (14000,0.01875736215) (14100,0.016551953) (14200,0.01565302061) (14300,0.01514319582) (14400,0.01393352565) (14500,0.01344165113) (14600,0.01276604049) (14700,0.01147782924) (14800,0.01159079517) (14900,0.01237729085) (15000,0.01131373518) (15100,0.01196234925) (15200,0.0123690417) (15300,0.0128935162) (15400,0.01358198575) (15500,0.0128683932) (15600,0.01188635676) (15700,0.01253822212) (15800,0.01293203494) (15900,0.0116049907) (16000,0.01142104583) (16100,0.01168316788) (16200,0.01074717993) (16300,0.009869108241) (16400,0.01029109246) (16500,0.0109875387) (16600,0.01081314816) (16700,0.009576224932) (16800,0.009718652559) (16900,0.01036393477) (17000,0.009846594301) (17100,0.009662882541) (17200,0.009728809097) (17300,0.008714389266) (17400,0.008520457917) (17500,0.00870137203) (17600,0.009632114321) (17700,0.009561263584) (17800,0.009250315651) (17900,0.00820006174) (18000,0.008040492842) (18100,0.008072871016) (18200,0.008788605733) (18300,0.009840086615) (18400,0.009428325994) (18500,0.009827972995) (18600,0.01029247898) (18700,0.01089061094) (18800,0.01108047436) (18900,0.01129885237) (19000,0.01136299204) (19100,0.0119965326) (19200,0.01124113523) (19300,0.01119632614) (19400,0.01011454422) (19500,0.01037825212) (19600,0.01039698958) (19700,0.0105198708) (19800,0.01107224669) (19900,0.01073697834) (20000,0.01064146832) (20100,0.01000017599) (20200,0.009683874092) (20300,0.009579300071) (20400,0.009851373918) (20500,0.008715379471) (20600,0.007088748086) (20700,0.006073349318) (20800,0.00572413865) (20900,0.006448072358) (21000,0.006932561495) (21100,0.007595705078) (21200,0.00809950307) (21300,0.008125230111) (21400,0.009706611582) (21500,0.009508131933) (21600,0.01101726524) (21700,0.01156603615) (21800,0.01176357069) (21900,0.01210226365) (22000,0.01239047865) (22100,0.01107780878) (22200,0.01044468975) (22300,0.01019531863) (22400,0.01009186639) (22500,0.009793045645) (22600,0.00847211433) (22700,0.008741918992) (22800,0.008186636347) (22900,0.007436044276) (23000,0.007468706224) (23100,0.008188040456) (23200,0.01011044497) (23300,0.01037709556) (23400,0.01084058016) (23500,0.01069913004) (23600,0.01135610798) (23700,0.01239147894) (23800,0.01261359677)};
\addplot[pubRose, opacity=0.18, line width=0.35pt] coordinates {(100,0.1773792803) (200,0.1777831987) (300,0.1780963937) (400,0.1800256185) (500,0.1800922424) (600,0.1760670965) (700,0.1750635483) (800,0.1746157017) (900,0.1754570239) (1000,0.1741352752) (1100,0.172187376) (1200,0.1711159527) (1300,0.1695218772) (1400,0.1667235762) (1500,0.1666446954) (1600,0.1692646682) (1700,0.1689148784) (1800,0.1691995949) (1900,0.1672823399) (2000,0.1672755986) (2100,0.1679586723) (2200,0.1668667972) (2300,0.1652635664) (2400,0.1660923794) (2500,0.1639609501) (2600,0.1612288624) (2700,0.1598942071) (2800,0.1563980848) (2900,0.1539649189) (3000,0.1508822411) (3100,0.1480548143) (3200,0.146239239) (3300,0.1452097878) (3400,0.1443439886) (3500,0.1413942114) (3600,0.140205133) (3700,0.1376539558) (3800,0.1352711007) (3900,0.1336585477) (4000,0.1322480805) (4100,0.1306470618) (4200,0.1292328432) (4300,0.1276781522) (4400,0.1247035958) (4500,0.124421566) (4600,0.1220887676) (4700,0.120402889) (4800,0.1200387619) (4900,0.1184773877) (5000,0.1169273578) (5100,0.1167725205) (5200,0.1153954387) (5300,0.1137213714) (5400,0.1119407356) (5500,0.1094605811) (5600,0.1077643178) (5700,0.1064517908) (5800,0.1027055703) (5900,0.1001741029) (6000,0.09810775816) (6100,0.0939762488) (6200,0.09089170024) (6300,0.08787722811) (6400,0.08535640761) (6500,0.08302626535) (6600,0.08078310713) (6700,0.07904026955) (6800,0.07811943814) (6900,0.07698529512) (7000,0.07530343458) (7100,0.07420690283) (7200,0.07449384853) (7300,0.07327361256) (7400,0.07061282545) (7500,0.06838710941) (7600,0.06498173662) (7700,0.06184020787) (7800,0.06019401439) (7900,0.05814553723) (8000,0.05692787208) (8100,0.05562731288) (8200,0.05131804235) (8300,0.04904326163) (8400,0.04666438252) (8500,0.04465950951) (8600,0.04362612441) (8700,0.04320238568) (8800,0.04094582126) (8900,0.03851991612) (9000,0.03654634114) (9100,0.03410280831) (9200,0.03363343216) (9300,0.03233252969) (9400,0.03126575835) (9500,0.02908801045) (9600,0.027009496) (9700,0.02473648256) (9800,0.02454924127) (9900,0.02370393192) (10000,0.02148620123) (10100,0.019893663) (10200,0.01763196532) (10300,0.01627200199) (10400,0.01548749441) (10500,0.01494717086) (10600,0.01427022498) (10700,0.0128118807) (10800,0.01054609139) (10900,0.009524987498) (11000,0.008278642478) (11100,0.007871131576) (11200,0.007581238891) (11300,0.006187076867) (11400,0.004663053062) (11500,0.004333295487) (11600,0.003285581688) (11700,0.001462245197) (11800,0.001884675887) (11900,0.001436711056) (12000,0.0008923373418) (12100,-0.0003946663113) (12200,-0.002946886071) (12300,-0.003689993033) (12400,-0.005081406562) (12500,-0.007114701765) (12600,-0.008298939024) (12700,-0.007612020825) (12800,-0.01003071803) (12900,-0.01033053851) (13000,-0.01097062953) (13100,-0.01121166117) (13200,-0.01049382939) (13300,-0.01042409876) (13400,-0.01067043488) (13500,-0.009941684559) (13600,-0.009427936759) (13700,-0.01015627348) (13800,-0.008572818048) (13900,-0.009255290939) (14000,-0.008114357502) (14100,-0.008151160437) (14200,-0.008608793165) (14300,-0.00925522584) (14400,-0.008848425071) (14500,-0.008649028442) (14600,-0.009024144965) (14700,-0.008530455525) (14800,-0.009590884857) (14900,-0.01084252205) (15000,-0.01132964306) (15100,-0.01263211085) (15200,-0.01144912875) (15300,-0.01080981948) (15400,-0.01126308375) (15500,-0.01148670631) (15600,-0.01124308965) (15700,-0.01109253721) (15800,-0.01074404776) (15900,-0.01038509294) (16000,-0.01057991533) (16100,-0.00961410132) (16200,-0.009350721346) (16300,-0.008639461169) (16400,-0.007826956542) (16500,-0.008440241398) (16600,-0.01000312529) (16700,-0.0105647989) (16800,-0.01003156299) (16900,-0.009197835572) (17000,-0.008963017951) (17100,-0.009160664718) (17200,-0.0101562459) (17300,-0.0105252428) (17400,-0.01073173163) (17500,-0.01030273391) (17600,-0.007971382001) (17700,-0.008329306403) (17800,-0.009127446683) (17900,-0.009136212384) (18000,-0.009247495001) (18100,-0.009703394817) (18200,-0.009792358289) (18300,-0.01065598284) (18400,-0.01116851508) (18500,-0.0114365852) (18600,-0.01169337058) (18700,-0.01143642995) (18800,-0.01081547202) (18900,-0.01100527612) (19000,-0.01136640157) (19100,-0.01026440896) (19200,-0.009222298392) (19300,-0.008830210191) (19400,-0.00808558026) (19500,-0.007341327376) (19600,-0.007633295) (19700,-0.007691507132) (19800,-0.007483098234) (19900,-0.005778381505) (20000,-0.005501747315) (20100,-0.006334520548) (20200,-0.007757227728) (20300,-0.007708805474) (20400,-0.00742476345) (20500,-0.00849048486) (20600,-0.00937575018) (20700,-0.009698427375) (20800,-0.009752777405) (20900,-0.0111668007) (21000,-0.01126986765) (21100,-0.01200988195) (21200,-0.01161863902) (21300,-0.01209525722) (21400,-0.01267494555) (21500,-0.01194909071) (21600,-0.01162144332) (21700,-0.01102533392) (21800,-0.01160368333) (21900,-0.01229542354) (22000,-0.01249318356) (22100,-0.01204503989) (22200,-0.01159565938) (22300,-0.01096411357) (22400,-0.01029785885) (22500,-0.00992653463) (22600,-0.008953882754) (22700,-0.009443172254) (22800,-0.01068299599) (22900,-0.01060735211) (23000,-0.01052812543) (23100,-0.009327614447) (23200,-0.01033057449) (23300,-0.01091983994) (23400,-0.01239769952) (23500,-0.01313839783) (23600,-0.01447064108) (23700,-0.01398186288) (23800,-0.0133022306)};
\addplot[pubRose, opacity=0.18, line width=0.35pt] coordinates {(100,0.01675347611) (200,0.01421806589) (300,0.01457024676) (400,0.01498809597) (500,0.01368669383) (600,0.01294984436) (700,0.01175195604) (800,0.01065350225) (900,0.01040737719) (1000,0.01045333366) (1100,0.00876010051) (1200,0.007795123663) (1300,0.006297669141) (1400,0.00521379225) (1500,0.005097287102) (1600,0.00406858637) (1700,0.003567959741) (1800,0.002884573326) (1900,0.003744226345) (2000,0.002944975696) (2100,0.003544075205) (2200,0.00360619307) (2300,0.004357778165) (2400,0.003474381589) (2500,0.002120743203) (2600,0.003328149184) (2700,0.004560765182) (2800,0.004987386568) (2900,0.004017078876) (3000,0.004119354766) (3100,0.002995521855) (3200,0.003147295024) (3300,0.001779589616) (3400,0.001667405572) (3500,0.002774248784) (3600,0.002028326644) (3700,0.001902485499) (3800,0.001882277307) (3900,0.001053136081) (4000,0.001762349909) (4100,0.002667155169) (4200,0.002249880392) (4300,0.003650555458) (4400,0.003773414508) (4500,0.004321063148) (4600,0.004694942875) (4700,0.004600024713) (4800,0.004486245925) (4900,0.005162651506) (5000,0.004679893232) (5100,0.004490482628) (5200,0.004939683469) (5300,0.005015809869) (5400,0.005996203737) (5500,0.005249386898) (5600,0.004862599133) (5700,0.004059494007) (5800,0.004175360892) (5900,0.004240994966) (6000,0.003932404234) (6100,0.004529195123) (6200,0.003687231087) (6300,0.003276579997) (6400,0.003025960813) (6500,0.002276485188) (6600,0.001594543953) (6700,0.001492273093) (6800,0.001352104766) (6900,0.0009539708612) (7000,0.0008059287793) (7100,0.001109310787) (7200,0.002470704901) (7300,0.00189473647) (7400,0.002010253206) (7500,0.002860357607) (7600,0.00273959508) (7700,0.001702723169) (7800,0.001792775141) (7900,0.002372432989) (8000,0.002226214323) (8100,0.001015648097) (8200,-0.0007698635192) (8300,-0.0005507104128) (8400,-0.001017724353) (8500,-0.001351098527) (8600,-0.001132944712) (8700,4.602086556e-05) (8800,-0.0001595746173) (8900,0.0006510257343) (9000,0.0007940373092) (9100,0.001039515523) (9200,0.002505312854) (9300,0.002267100284) (9400,0.003033665608) (9500,0.002952205471) (9600,0.004290947982) (9700,0.003770878806) (9800,0.004743870615) (9900,0.003167777765) (10000,0.003417202563) (10100,0.003719818103) (10200,0.002862042096) (10300,0.0029912047) (10400,0.002692438374) (10500,0.002145545362) (10600,0.001559905327) (10700,0.0009582928993) (10800,-0.0004048903793) (10900,-0.0005988784222) (11000,-0.0006277173845) (11100,-0.001180655634) (11200,-0.0004964349413) (11300,-0.0005329358988) (11400,-0.0005583015867) (11500,-5.519533443e-05) (11600,0.0001320924712) (11700,0.0007135034175) (11800,0.002175335922) (11900,0.002306456109) (12000,0.002279383493) (12100,0.003183353053) (12200,0.003022337345) (12300,0.002737556971) (12400,0.003136158842) (12500,0.002384062196) (12600,0.001764478083) (12700,0.001367080226) (12800,0.0005105677046) (12900,0.0003123301154) (13000,3.725146089e-05) (13100,-0.0003920286181) (13200,-4.212229833e-05) (13300,-0.0001375268432) (13400,-0.0001793137038) (13500,0.0005325418562) (13600,0.0003553244765) (13700,0.001254676633) (13800,0.001743196644) (13900,0.001414618682) (14000,0.001468616514) (14100,0.0008072810713) (14200,-0.0008307982236) (14300,-0.0001625963952) (14400,-0.00202424298) (14500,-0.002203731833) (14600,-0.002414039802) (14700,-0.003000455559) (14800,-0.00334280231) (14900,-0.002794745899) (15000,-0.003757636048) (15100,-0.003233282824) (15200,-0.002373515951) (15300,-0.003296068261) (15400,-0.002031382566) (15500,-0.002040321773) (15600,-0.001740825053) (15700,-0.001167254574) (15800,-0.002313749101) (15900,-0.003331228552) (16000,-0.002640080728) (16100,-0.002629490158) (16200,-0.003434735917) (16300,-0.002772290539) (16400,-0.003510282259) (16500,-0.003963759146) (16600,-0.003927136458) (16700,-0.004306536704) (16800,-0.003286773637) (16900,-0.003482007905) (17000,-0.003025176092) (17100,-0.003064861772) (17200,-0.003133127624) (17300,-0.00283108787) (17400,-0.002627624504) (17500,-0.001652475403) (17600,-0.002243832283) (17700,-0.001398110407) (17800,-0.002112323162) (17900,-0.001347517851) (18000,-0.001746047894) (18100,-0.001713206619) (18200,-0.0002610634547) (18300,-0.0007204144844) (18400,-0.0004690148286) (18500,-0.000938393909) (18600,-0.001186671795) (18700,-0.002015252423) (18800,-0.001943161816) (18900,-0.001918078971) (19000,-0.002924978163) (19100,-0.003678802948) (19200,-0.003982569661) (19300,-0.003537406633) (19400,-0.003835217329) (19500,-0.004490883532) (19600,-0.004010251095) (19700,-0.004038866656) (19800,-0.003683637874) (19900,-0.003474626329) (20000,-0.003089550708) (20100,-0.003397487674) (20200,-0.004651441833) (20300,-0.005191198096) (20400,-0.005074046634) (20500,-0.004059970437) (20600,-0.00485332181) (20700,-0.0042033604) (20800,-0.00487219264) (20900,-0.004788406659) (21000,-0.004878136329) (21100,-0.00411671293) (21200,-0.004038185021) (21300,-0.004881643457) (21400,-0.004851690866) (21500,-0.004916079054) (21600,-0.004596277318) (21700,-0.005952508413) (21800,-0.005283767232) (21900,-0.00517223322) (22000,-0.00568245345) (22100,-0.0065825112) (22200,-0.006500272809) (22300,-0.006561918726) (22400,-0.006643222745) (22500,-0.007071920406) (22600,-0.00756516663) (22700,-0.006924118932) (22800,-0.007369261623) (22900,-0.008608827135) (23000,-0.00815323065) (23100,-0.008411504654) (23200,-0.007573538314) (23300,-0.007269739493) (23400,-0.009044344252) (23500,-0.01004861912) (23600,-0.01042154405) (23700,-0.01074182363) (23800,-0.01185050741)};
\addplot[pubRose, opacity=0.18, line width=0.35pt] coordinates {(100,0.1283351183) (200,0.1265379451) (300,0.1241343046) (400,0.1239180118) (500,0.1213648587) (600,0.1187833324) (700,0.1184742855) (800,0.1179349963) (900,0.1169270029) (1000,0.1156912051) (1100,0.1137323983) (1200,0.1118530504) (1300,0.111667797) (1400,0.1105853587) (1500,0.1098096639) (1600,0.1104256906) (1700,0.1100780286) (1800,0.1098942891) (1900,0.1087244034) (2000,0.1081269078) (2100,0.1072101444) (2200,0.1056947745) (2300,0.1047614127) (2400,0.1043972492) (2500,0.1032205664) (2600,0.102739431) (2700,0.1019078583) (2800,0.09999584481) (2900,0.1004422843) (3000,0.1007976532) (3100,0.09997547194) (3200,0.100084804) (3300,0.09824707955) (3400,0.09606279507) (3500,0.09533281773) (3600,0.09289249927) (3700,0.09150299355) (3800,0.09140394703) (3900,0.09010356665) (4000,0.08925267756) (4100,0.08954501748) (4200,0.08872799277) (4300,0.08883985803) (4400,0.08972937167) (4500,0.09113521501) (4600,0.09240636006) (4700,0.09190173075) (4800,0.09100085497) (4900,0.09022297859) (5000,0.08917125836) (5100,0.08689523265) (5200,0.08732559681) (5300,0.0858101666) (5400,0.08328682855) (5500,0.08165250197) (5600,0.08034928963) (5700,0.08003057763) (5800,0.08162985221) (5900,0.08236908838) (6000,0.08364121541) (6100,0.08551618606) (6200,0.08436951786) (6300,0.08527163938) (6400,0.0865377225) (6500,0.08817160055) (6600,0.08710538149) (6700,0.0866888009) (6800,0.08622021079) (6900,0.08554131463) (7000,0.08428209871) (7100,0.08376776278) (7200,0.08398917019) (7300,0.08317374587) (7400,0.08157219887) (7500,0.07848131284) (7600,0.07934771031) (7700,0.07701067552) (7800,0.0759438768) (7900,0.07602094337) (8000,0.07438674793) (8100,0.07508080155) (8200,0.0738452211) (8300,0.07414104119) (8400,0.07550864443) (8500,0.0762129426) (8600,0.07515811473) (8700,0.07650267482) (8800,0.07627431378) (8900,0.0742304042) (9000,0.07379585356) (9100,0.07001229338) (9200,0.07003352307) (9300,0.06945740096) (9400,0.06936432458) (9500,0.06921630166) (9600,0.06925459541) (9700,0.06885077097) (9800,0.06653047428) (9900,0.06621517874) (10000,0.0666057419) (10100,0.0677087158) (10200,0.0663321238) (10300,0.06536593847) (10400,0.06240554415) (10500,0.05985534415) (10600,0.05855925493) (10700,0.05663036816) (10800,0.05792397559) (10900,0.05812153257) (11000,0.05675439537) (11100,0.05640883446) (11200,0.05637956187) (11300,0.05456252769) (11400,0.05400050655) (11500,0.05365355797) (11600,0.05341252945) (11700,0.05375220254) (11800,0.05082882643) (11900,0.05035939217) (12000,0.05105236694) (12100,0.04921733104) (12200,0.0484671209) (12300,0.04833376557) (12400,0.04769968763) (12500,0.0472187452) (12600,0.04534998573) (12700,0.04533848986) (12800,0.04604697376) (12900,0.04414385073) (13000,0.04198747948) (13100,0.04033250287) (13200,0.03868230525) (13300,0.03784690145) (13400,0.03723124769) (13500,0.03775554802) (13600,0.0377995627) (13700,0.03590366654) (13800,0.03482180536) (13900,0.03516163453) (14000,0.03450100981) (14100,0.03545269854) (14200,0.03591510337) (14300,0.03451669738) (14400,0.03494520523) (14500,0.03291521817) (14600,0.03258141167) (14700,0.03231318705) (14800,0.03080984652) (14900,0.02949236818) (15000,0.0294771716) (15100,0.029033022) (15200,0.02825603317) (15300,0.02874130197) (15400,0.02735094372) (15500,0.02800872866) (15600,0.02724689338) (15700,0.02767830472) (15800,0.02825640105) (15900,0.02841650359) (16000,0.0277340984) (16100,0.02686356474) (16200,0.02600890528) (16300,0.02512849588) (16400,0.02457787655) (16500,0.02313565388) (16600,0.0218641758) (16700,0.02036112584) (16800,0.01955779213) (16900,0.01861727517) (17000,0.01791483629) (17100,0.01720967423) (17200,0.01577029035) (17300,0.01497010041) (17400,0.01511010483) (17500,0.01490002836) (17600,0.01602345149) (17700,0.0159290611) (17800,0.01494650797) (17900,0.01491872922) (18000,0.01572811541) (18100,0.01620939348) (18200,0.01651519393) (18300,0.01697566318) (18400,0.01591315609) (18500,0.01601904188) (18600,0.01416278505) (18700,0.01419061958) (18800,0.0148699706) (18900,0.01465752996) (19000,0.01288274066) (19100,0.01165694385) (19200,0.01191025218) (19300,0.01117820512) (19400,0.01197984763) (19500,0.01129344893) (19600,0.01212227424) (19700,0.01177815809) (19800,0.01114696604) (19900,0.01078518149) (20000,0.01142486162) (20100,0.01137602138) (20200,0.01212971178) (20300,0.01301779291) (20400,0.01154090454) (20500,0.01168664236) (20600,0.01179117354) (20700,0.01143913873) (20800,0.01246877237) (20900,0.01219637356) (21000,0.01151310442) (21100,0.01143328057) (21200,0.01065119993) (21300,0.009527812328) (21400,0.009762345394) (21500,0.009875367722) (21600,0.00863191695) (21700,0.007602989185) (21800,0.007214273768) (21900,0.007227580505) (22000,0.007712179632) (22100,0.008453551237) (22200,0.009039003332) (22300,0.009262728482) (22400,0.009569749446) (22500,0.01070629314) (22600,0.01069096837) (22700,0.01207141695) (22800,0.01132166744) (22900,0.01253316761) (23000,0.01119023144) (23100,0.01042239057) (23200,0.009739402414) (23300,0.01021975885) (23400,0.01120720125) (23500,0.01001478789) (23600,0.00978807284) (23700,0.009024766725) (23800,0.009095042091)};
\addplot[pubRose, opacity=0.18, line width=0.35pt] coordinates {(100,0.2917630672) (200,0.2935343534) (300,0.2888441682) (400,0.2943510339) (500,0.289987582) (600,0.2851322989) (700,0.2831748128) (800,0.2832288481) (900,0.2828929722) (1000,0.28084023) (1100,0.2777728945) (1200,0.2721418872) (1300,0.2715779588) (1400,0.2671135738) (1500,0.2627762541) (1600,0.2620278701) (1700,0.2575509161) (1800,0.2518612951) (1900,0.2467585295) (2000,0.2453515008) (2100,0.2429120988) (2200,0.2427283302) (2300,0.2375569895) (2400,0.2336547375) (2500,0.2318312228) (2600,0.2269723549) (2700,0.2262483567) (2800,0.2250620976) (2900,0.2224850267) (3000,0.2172690615) (3100,0.214002046) (3200,0.2085540205) (3300,0.2058612898) (3400,0.2025125369) (3500,0.2018032834) (3600,0.2024589196) (3700,0.2012322828) (3800,0.1993420914) (3900,0.1969564945) (4000,0.1978500143) (4100,0.1976590276) (4200,0.1961584508) (4300,0.1967589289) (4400,0.1955894321) (4500,0.1913734704) (4600,0.1865584612) (4700,0.1838132858) (4800,0.1812288746) (4900,0.1795304298) (5000,0.1755920127) (5100,0.1724054024) (5200,0.1707513541) (5300,0.1681520388) (5400,0.1676552311) (5500,0.1679624051) (5600,0.1652405441) (5700,0.1633846179) (5800,0.1613407046) (5900,0.1579697132) (6000,0.1536282063) (6100,0.1490796871) (6200,0.1476236992) (6300,0.1422923408) (6400,0.1389362209) (6500,0.1324289553) (6600,0.1306360319) (6700,0.126823698) (6800,0.1219794407) (6900,0.1223677665) (7000,0.1211659871) (7100,0.1207333706) (7200,0.1186220609) (7300,0.1180330388) (7400,0.1154513426) (7500,0.1180170424) (7600,0.1151492737) (7700,0.1138186194) (7800,0.1120939508) (7900,0.1088747293) (8000,0.1077529386) (8100,0.1043036349) (8200,0.1007203989) (8300,0.09579464942) (8400,0.0924796924) (8500,0.08852477148) (8600,0.08575291261) (8700,0.08293634057) (8800,0.08135216832) (8900,0.07838997245) (9000,0.07464327328) (9100,0.07154631093) (9200,0.07012845948) (9300,0.06954080574) (9400,0.06634608284) (9500,0.06226969324) (9600,0.05976592749) (9700,0.05655659847) (9800,0.05573595762) (9900,0.05234788284) (10000,0.05040489584) (10100,0.04987632036) (10200,0.04783613458) (10300,0.04563578218) (10400,0.04371935055) (10500,0.04088737145) (10600,0.03958739731) (10700,0.03759877086) (10800,0.03477084078) (10900,0.03307659905) (11000,0.03121756092) (11100,0.02834817152) (11200,0.02692937739) (11300,0.02408361873) (11400,0.02308704546) (11500,0.02318317769) (11600,0.02193272356) (11700,0.02035899563) (11800,0.01935015908) (11900,0.01904425202) (12000,0.01792512052) (12100,0.01692021275) (12200,0.0145916393) (12300,0.01476631658) (12400,0.01216016365) (12500,0.01110224989) (12600,0.01021919223) (12700,0.01022446603) (12800,0.007967733894) (12900,0.006758776796) (13000,0.005567253774) (13100,0.004641549158) (13200,0.004041085028) (13300,0.002276425072) (13400,0.003627229211) (13500,0.003559887465) (13600,0.004043107721) (13700,0.002419868307) (13800,0.001346434266) (13900,0.00120766052) (14000,0.001551164559) (14100,0.001171128405) (14200,0.001188922813) (14300,0.001923336997) (14400,-0.0002109133871) (14500,-0.001322645438) (14600,-0.003006332624) (14700,-0.00291156345) (14800,-0.002258554799) (14900,-0.003429381526) (15000,-0.004260733747) (15100,-0.005117663229) (15200,-0.006423953269) (15300,-0.008053577808) (15400,-0.00668356705) (15500,-0.006414004113) (15600,-0.006191772246) (15700,-0.006731998222) (15800,-0.007514932891) (15900,-0.008323215088) (16000,-0.008475530846) (16100,-0.008525287593) (16200,-0.008682795102) (16300,-0.007761205966) (16400,-0.00905760061) (16500,-0.009141896362) (16600,-0.009237553854) (16700,-0.01007689952) (16800,-0.01002651907) (16900,-0.00969640871) (17000,-0.01014869006) (17100,-0.01019713113) (17200,-0.01047163473) (17300,-0.010708022) (17400,-0.01077910771) (17500,-0.0120239757) (17600,-0.01335865743) (17700,-0.01352971978) (17800,-0.01362996046) (17900,-0.01359990444) (18000,-0.01391604207) (18100,-0.01355525767) (18200,-0.01294598235) (18300,-0.01296998262) (18400,-0.01452233139) (18500,-0.01413806425) (18600,-0.01310442174) (18700,-0.01219761102) (18800,-0.01232148847) (18900,-0.01179939918) (19000,-0.0119985681) (19100,-0.01217613025) (19200,-0.01272706417) (19300,-0.01305905739) (19400,-0.01250527194) (19500,-0.01267915303) (19600,-0.01251938632) (19700,-0.01298472062) (19800,-0.01348602781) (19900,-0.0143232665) (20000,-0.01465918319) (20100,-0.01473490354) (20200,-0.01407631729) (20300,-0.01429830361) (20400,-0.01415700857) (20500,-0.01403588122) (20600,-0.01512749651) (20700,-0.01540483106) (20800,-0.01456797346) (20900,-0.01440732004) (21000,-0.01420939369) (21100,-0.0140538902) (21200,-0.01418276839) (21300,-0.01404235857) (21400,-0.01305876672) (21500,-0.01401272062) (21600,-0.01283570873) (21700,-0.01229034476) (21800,-0.01254451396) (21900,-0.0125671939) (22000,-0.01286565149) (22100,-0.01322313864) (22200,-0.01376054045) (22300,-0.01374377357) (22400,-0.01445984021) (22500,-0.0131241668) (22600,-0.01370574851) (22700,-0.01438341318) (22800,-0.01527618635) (22900,-0.01518239351) (23000,-0.0147389587) (23100,-0.01442587757) (23200,-0.01448515374) (23300,-0.01483464474) (23400,-0.01457632678) (23500,-0.01516815163) (23600,-0.0152559231) (23700,-0.01485789157) (23800,-0.01373928394)};
\addplot[pubRose, opacity=0.18, line width=0.35pt] coordinates {(100,0.2037498206) (200,0.1943643987) (300,0.1889936129) (400,0.1885908507) (500,0.1852852404) (600,0.1835707501) (700,0.1833973484) (800,0.1819780339) (900,0.1804640873) (1000,0.1777970701) (1100,0.1729470342) (1200,0.1711557969) (1300,0.1695524901) (1400,0.1665685862) (1500,0.1648608521) (1600,0.1633679792) (1700,0.1609345868) (1800,0.1594041169) (1900,0.1571281567) (2000,0.1559090599) (2100,0.1561102688) (2200,0.1532439202) (2300,0.151904799) (2400,0.1509953141) (2500,0.1500295192) (2600,0.1482273147) (2700,0.1470463812) (2800,0.1450073048) (2900,0.1444051012) (3000,0.1436627477) (3100,0.1413723752) (3200,0.1418952182) (3300,0.1409199849) (3400,0.1392142922) (3500,0.1373055324) (3600,0.1358750463) (3700,0.1334526397) (3800,0.1323868074) (3900,0.1302595019) (4000,0.1278238364) (4100,0.1257309712) (4200,0.123897437) (4300,0.1217922166) (4400,0.1204378702) (4500,0.1183833897) (4600,0.1172158003) (4700,0.1151586473) (4800,0.1129961818) (4900,0.1110132337) (5000,0.1109172024) (5100,0.1090047598) (5200,0.1063987568) (5300,0.1042079993) (5400,0.102084735) (5500,0.1012023114) (5600,0.09909203202) (5700,0.09901983738) (5800,0.09807110876) (5900,0.09721024558) (6000,0.0968338564) (6100,0.09574999958) (6200,0.09454959631) (6300,0.09248911887) (6400,0.09071480781) (6500,0.08919393346) (6600,0.08706128448) (6700,0.08478277922) (6800,0.0828599304) (6900,0.08031300828) (7000,0.07663120478) (7100,0.07503764182) (7200,0.07318287715) (7300,0.07240537181) (7400,0.07056814656) (7500,0.06760978252) (7600,0.06701137722) (7700,0.06488990746) (7800,0.06304036416) (7900,0.062123565) (8000,0.06074428409) (8100,0.05920161456) (8200,0.05703940578) (8300,0.05487028845) (8400,0.05401859917) (8500,0.05334843583) (8600,0.0504867088) (8700,0.0496330142) (8800,0.04718924612) (8900,0.04552055784) (9000,0.04401904233) (9100,0.04271245301) (9200,0.04171856716) (9300,0.039704437) (9400,0.03754772898) (9500,0.03687611874) (9600,0.03580071963) (9700,0.03430273607) (9800,0.03465158567) (9900,0.03299212623) (10000,0.03147569168) (10100,0.02999769486) (10200,0.02811184879) (10300,0.02855401579) (10400,0.02773539387) (10500,0.02541880272) (10600,0.02434941083) (10700,0.02136942705) (10800,0.0204358218) (10900,0.01925476769) (11000,0.01847096251) (11100,0.01795907614) (11200,0.01760163594) (11300,0.01557018037) (11400,0.01442415712) (11500,0.01338445821) (11600,0.01258518535) (11700,0.01379075386) (11800,0.0130728621) (11900,0.01258473163) (12000,0.01212597108) (12100,0.01079268097) (12200,0.01043984001) (12300,0.009129747818) (12400,0.008764383686) (12500,0.00916505421) (12600,0.007892271038) (12700,0.006082682125) (12800,0.003949200676) (12900,0.002919136302) (13000,0.002277060819) (13100,0.001857695146) (13200,5.519657861e-05) (13300,-0.0001203577965) (13400,-0.001101043331) (13500,-0.002653284767) (13600,-0.002763515012) (13700,-0.002728027361) (13800,-0.003202666633) (13900,-0.003029593651) (14000,-0.003904127248) (14100,-0.004706186196) (14200,-0.004472134449) (14300,-0.004126463621) (14400,-0.00420026877) (14500,-0.004642986879) (14600,-0.00459425) (14700,-0.004752439074) (14800,-0.005102720205) (14900,-0.005014299159) (15000,-0.006016701669) (15100,-0.006522579747) (15200,-0.007352873893) (15300,-0.008222231083) (15400,-0.00892189946) (15500,-0.008557990612) (15600,-0.00985812603) (15700,-0.0110380763) (15800,-0.01135684929) (15900,-0.01186312265) (16000,-0.01118505225) (16100,-0.01181098642) (16200,-0.01113269408) (16300,-0.01167508229) (16400,-0.0113050139) (16500,-0.01152413953) (16600,-0.01088548992) (16700,-0.01068328042) (16800,-0.01054099062) (16900,-0.01101326542) (17000,-0.01160599277) (17100,-0.01166739971) (17200,-0.01289630234) (17300,-0.01328081628) (17400,-0.01431226917) (17500,-0.01478194054) (17600,-0.01498089479) (17700,-0.01548294509) (17800,-0.01587941376) (17900,-0.01625424847) (18000,-0.01672513047) (18100,-0.01601472925) (18200,-0.01560954135) (18300,-0.0146970137) (18400,-0.01539148306) (18500,-0.01534624705) (18600,-0.01613601781) (18700,-0.01528506698) (18800,-0.01555393348) (18900,-0.01553994445) (19000,-0.01570155686) (19100,-0.01590905832) (19200,-0.01599432528) (19300,-0.01694725892) (19400,-0.0163204872) (19500,-0.01660165982) (19600,-0.01556910696) (19700,-0.01668732595) (19800,-0.01666729841) (19900,-0.01707271319) (20000,-0.01709282398) (20100,-0.01746829329) (20200,-0.01790320445) (20300,-0.01783131808) (20400,-0.01795093063) (20500,-0.01837987825) (20600,-0.01906342972) (20700,-0.01798306108) (20800,-0.01791223343) (20900,-0.01738616927) (21000,-0.01737419581) (21100,-0.01724984748) (21200,-0.01702367188) (21300,-0.01629215218) (21400,-0.01655758955) (21500,-0.01643920038) (21600,-0.0159802461) (21700,-0.01666306313) (21800,-0.01546580233) (21900,-0.01581887333) (22000,-0.01573428409) (22100,-0.01515884036) (22200,-0.01497855708) (22300,-0.01496585654) (22400,-0.01493496103) (22500,-0.01412199307) (22600,-0.01434683902) (22700,-0.01418648688) (22800,-0.01538721854) (22900,-0.016120099) (23000,-0.01483865026) (23100,-0.01525041731) (23200,-0.01533074044) (23300,-0.01594742928) (23400,-0.01513064615) (23500,-0.01539646778) (23600,-0.01636619521) (23700,-0.0159247566) (23800,-0.01552414419)};
\addplot[pubRose, opacity=0.18, line width=0.35pt] coordinates {(100,0.07764587551) (200,0.07181384787) (300,0.07088593394) (400,0.06769961957) (500,0.06417978331) (600,0.06305374391) (700,0.06074587788) (800,0.05867204117) (900,0.05741203328) (1000,0.05643786639) (1100,0.05392428227) (1200,0.05297196768) (1300,0.05140202157) (1400,0.05109192319) (1500,0.05135314018) (1600,0.05123113282) (1700,0.05156637654) (1800,0.05194764584) (1900,0.05201497078) (2000,0.05237290189) (2100,0.05230832808) (2200,0.05158377327) (2300,0.05052038655) (2400,0.04917881228) (2500,0.04857722148) (2600,0.04762788899) (2700,0.04655224979) (2800,0.04611846618) (2900,0.0453333579) (3000,0.04436429217) (3100,0.04406992607) (3200,0.04387056641) (3300,0.04406231865) (3400,0.0454000283) (3500,0.04536570571) (3600,0.04560288265) (3700,0.04608965889) (3800,0.04528391436) (3900,0.04570357352) (4000,0.04680208303) (4100,0.04554088488) (4200,0.04538870603) (4300,0.04440832846) (4400,0.04335293919) (4500,0.04323040247) (4600,0.04148957878) (4700,0.04163427614) (4800,0.04154557735) (4900,0.04108291008) (5000,0.03889484275) (5100,0.03917736989) (5200,0.03856996763) (5300,0.0383988576) (5400,0.03753097076) (5500,0.03663982656) (5600,0.03609030098) (5700,0.03417715784) (5800,0.03413610253) (5900,0.03433869611) (6000,0.03430266324) (6100,0.03451411892) (6200,0.03336211517) (6300,0.03211163729) (6400,0.03188857324) (6500,0.03113129474) (6600,0.03184844311) (6700,0.03219302576) (6800,0.03188926242) (6900,0.02988294605) (7000,0.0294290835) (7100,0.02779073436) (7200,0.02894120999) (7300,0.03027351946) (7400,0.02882872596) (7500,0.02911073491) (7600,0.0277992757) (7700,0.02738971878) (7800,0.02527797446) (7900,0.02564517595) (8000,0.02484513558) (8100,0.02485039178) (8200,0.02297489606) (8300,0.02102188943) (8400,0.02114528352) (8500,0.02031398313) (8600,0.02033814797) (8700,0.02104443861) (8800,0.02280502813) (8900,0.02178503685) (9000,0.02166933659) (9100,0.02090743408) (9200,0.02046667896) (9300,0.02108687302) (9400,0.02182630124) (9500,0.02130029192) (9600,0.02080126358) (9700,0.01911570383) (9800,0.01757508665) (9900,0.01811377918) (10000,0.01673679275) (10100,0.01530948433) (10200,0.01430066142) (10300,0.01314080765) (10400,0.01067632399) (10500,0.0102064061) (10600,0.009265478095) (10700,0.008788086055) (10800,0.007544331696) (10900,0.007001165558) (11000,0.007360075462) (11100,0.008134268971) (11200,0.008420337226) (11300,0.007897741861) (11400,0.008275137447) (11500,0.007062071563) (11600,0.007272029064) (11700,0.006103057449) (11800,0.00659384505) (11900,0.005654166696) (12000,0.005648084321) (12100,0.004300969686) (12200,0.003107958693) (12300,0.002661524776) (12400,0.002255862781) (12500,0.002312434097) (12600,0.002003210368) (12700,0.001958831522) (12800,0.001394900191) (12900,0.001006454043) (13000,0.0001076874556) (13100,0.0002550957375) (13200,0.0008236995549) (13300,0.0006213854824) (13400,0.0003055012436) (13500,0.0004295327468) (13600,-0.0009383194847) (13700,-0.0003099895897) (13800,0.0002022235363) (13900,-0.0002668651578) (14000,-0.0004806073906) (14100,-0.0006144879328) (14200,-0.001784733398) (14300,-0.001881187543) (14400,-0.002347629727) (14500,-0.002645680876) (14600,-0.002518440559) (14700,-0.003418773375) (14800,-0.004199546864) (14900,-0.004543647481) (15000,-0.00524589823) (15100,-0.006110240621) (15200,-0.005993415471) (15300,-0.006056936141) (15400,-0.005410758776) (15500,-0.005957157107) (15600,-0.006074373599) (15700,-0.006783148309) (15800,-0.007822879672) (15900,-0.00880474326) (16000,-0.008874194778) (16100,-0.008879062242) (16200,-0.009755096619) (16300,-0.009747473884) (16400,-0.01063680835) (16500,-0.01105245668) (16600,-0.01061304109) (16700,-0.01103742588) (16800,-0.01078664539) (16900,-0.01028031074) (17000,-0.01035234248) (17100,-0.01088066017) (17200,-0.009846191108) (17300,-0.009599035466) (17400,-0.008833769891) (17500,-0.009070035718) (17600,-0.009465970001) (17700,-0.009444318894) (17800,-0.008849010193) (17900,-0.008511016514) (18000,-0.008873657692) (18100,-0.008659783094) (18200,-0.00905644856) (18300,-0.00932462527) (18400,-0.009226594353) (18500,-0.008719224297) (18600,-0.008597020293) (18700,-0.008025082899) (18800,-0.008568891231) (18900,-0.008441523276) (19000,-0.007889213599) (19100,-0.008114780765) (19200,-0.008482066914) (19300,-0.009001438087) (19400,-0.010032622) (19500,-0.009211233689) (19600,-0.009794994083) (19700,-0.01012177813) (19800,-0.01049216875) (19900,-0.01124227584) (20000,-0.01154094596) (20100,-0.01083373359) (20200,-0.009929984354) (20300,-0.009646681801) (20400,-0.009986984183) (20500,-0.011494593) (20600,-0.01118016061) (20700,-0.0106799406) (20800,-0.009735250892) (20900,-0.009157792805) (21000,-0.008948387997) (21100,-0.00902589974) (21200,-0.009686402325) (21300,-0.008953668014) (21400,-0.008095056913) (21500,-0.007643435872) (21600,-0.006866013259) (21700,-0.007088326616) (21800,-0.007977169938) (21900,-0.007871267572) (22000,-0.007954481617) (22100,-0.008409227058) (22200,-0.008541471511) (22300,-0.009324684297) (22400,-0.009820239502) (22500,-0.01109987197) (22600,-0.01164167095) (22700,-0.01189033678) (22800,-0.01151880194) (22900,-0.01208191374) (23000,-0.01160873659) (23100,-0.01117899111) (23200,-0.01116661131) (23300,-0.01139291297) (23400,-0.01138825719) (23500,-0.01043399521) (23600,-0.01133974148) (23700,-0.01124817859) (23800,-0.01208194196)};
\addplot[pubRose, opacity=0.18, line width=0.35pt] coordinates {(100,0.2015779167) (200,0.20529636) (300,0.2053866486) (400,0.2070587203) (500,0.2028960735) (600,0.1993713851) (700,0.2006993464) (800,0.2016134895) (900,0.2012901488) (1000,0.2010851145) (1100,0.200733085) (1200,0.1976485491) (1300,0.1959517539) (1400,0.1939958036) (1500,0.1934974954) (1600,0.1936528698) (1700,0.1910595149) (1800,0.1883707121) (1900,0.1861807719) (2000,0.1843705058) (2100,0.1840920478) (2200,0.1840691537) (2300,0.1826678738) (2400,0.1802206799) (2500,0.178496103) (2600,0.1757582575) (2700,0.1730193838) (2800,0.1720954373) (2900,0.1707640618) (3000,0.1670503274) (3100,0.1640924469) (3200,0.1624450222) (3300,0.1587404236) (3400,0.1574376166) (3500,0.1565041885) (3600,0.154909867) (3700,0.1522517309) (3800,0.1491568372) (3900,0.1470972434) (4000,0.1468951121) (4100,0.1439376071) (4200,0.1419209331) (4300,0.1417105749) (4400,0.1390336677) (4500,0.1378536999) (4600,0.1388329595) (4700,0.1382160209) (4800,0.1375098176) (4900,0.1367170565) (5000,0.137088453) (5100,0.1373261534) (5200,0.1369260974) (5300,0.1361934371) (5400,0.1365508236) (5500,0.1336450838) (5600,0.1309783697) (5700,0.1291889064) (5800,0.1266420215) (5900,0.125807935) (6000,0.1237126887) (6100,0.1225279868) (6200,0.1213625714) (6300,0.1200700372) (6400,0.1194886997) (6500,0.1195785709) (6600,0.117013026) (6700,0.1184571221) (6800,0.1183253609) (6900,0.1155777961) (7000,0.112573918) (7100,0.1087429479) (7200,0.1049193665) (7300,0.1026099794) (7400,0.09958845899) (7500,0.09659230709) (7600,0.09456344247) (7700,0.09194052592) (7800,0.09146359488) (7900,0.09021094814) (8000,0.09045087993) (8100,0.09133205339) (8200,0.09000535533) (8300,0.08908672258) (8400,0.08665229082) (8500,0.08493698016) (8600,0.08472744673) (8700,0.08310591727) (8800,0.0812077716) (8900,0.07890535519) (9000,0.07585465536) (9100,0.07447064668) (9200,0.07406354174) (9300,0.0744780086) (9400,0.07431740314) (9500,0.07461329028) (9600,0.07262240984) (9700,0.07223081253) (9800,0.0702340547) (9900,0.06985997967) (10000,0.06882365271) (10100,0.06524993852) (10200,0.06427976117) (10300,0.06136075445) (10400,0.05973188765) (10500,0.05855401829) (10600,0.05930374004) (10700,0.05692611635) (10800,0.05476217978) (10900,0.05329077058) (11000,0.05221120603) (11100,0.05180531852) (11200,0.04997545183) (11300,0.04816506132) (11400,0.04559453465) (11500,0.04361856245) (11600,0.0403157942) (11700,0.03872748911) (11800,0.03890862353) (11900,0.03739664685) (12000,0.03536634836) (12100,0.03329720087) (12200,0.03243584298) (12300,0.02986174319) (12400,0.02968225908) (12500,0.02808302809) (12600,0.02804188561) (12700,0.02662128564) (12800,0.02411648352) (12900,0.02323292177) (13000,0.02287705764) (13100,0.02248119991) (13200,0.01992023336) (13300,0.02052957164) (13400,0.01894636815) (13500,0.01865334203) (13600,0.01662458889) (13700,0.01587555716) (13800,0.01594533725) (13900,0.01429649591) (14000,0.01372651923) (14100,0.01299437694) (14200,0.01346588638) (14300,0.01253185915) (14400,0.01286398871) (14500,0.01135009578) (14600,0.01100147436) (14700,0.01101322523) (14800,0.009521644702) (14900,0.009608816635) (15000,0.008498946391) (15100,0.008558070567) (15200,0.008007296361) (15300,0.007155724824) (15400,0.006561392942) (15500,0.005819743662) (15600,0.005000606063) (15700,0.004004078184) (15800,0.004081174277) (15900,0.003275455593) (16000,0.003089931014) (16100,0.002213042101) (16200,0.001940617722) (16300,0.002345284109) (16400,0.001142609573) (16500,0.001614425902) (16600,0.002047606243) (16700,0.002461835765) (16800,0.00205970651) (16900,0.00225135911) (17000,0.002748124429) (17100,0.00303034893) (17200,0.002563866353) (17300,0.002359526124) (17400,0.002868210254) (17500,0.002824150244) (17600,0.001930048256) (17700,0.001329974926) (17800,0.001407342558) (17900,0.001395594471) (18000,0.001116169619) (18100,0.0005671883147) (18200,0.0007764473499) (18300,0.0006310189085) (18400,-1.75033987e-05) (18500,-0.0005052391731) (18600,0.0003754023754) (18700,0.0005224072258) (18800,0.0005609391956) (18900,0.0006165390208) (19000,0.0004929499075) (19100,0.0007159349931) (19200,-0.0002136268929) (19300,-0.0007900683253) (19400,0.0002286773131) (19500,0.0003927264668) (19600,-3.44550208e-06) (19700,-0.0004366447261) (19800,-0.0007394508561) (19900,-0.001121153234) (20000,-0.001669672999) (20100,-0.002204346965) (20200,-0.001962554298) (20300,-0.002130389359) (20400,-0.003210595675) (20500,-0.003218708682) (20600,-0.003636701667) (20700,-0.00366286175) (20800,-0.00329241415) (20900,-0.002676674025) (21000,-0.0027969238) (21100,-0.00307302716) (21200,-0.002816931612) (21300,-0.003553773451) (21400,-0.003221329569) (21500,-0.002904855134) (21600,-0.001966027217) (21700,-0.002391628758) (21800,-0.003112183418) (21900,-0.003210706939) (22000,-0.003011646838) (22100,-0.003143594664) (22200,-0.002873263806) (22300,-0.00156978983) (22400,-0.001593189799) (22500,-0.00229542872) (22600,-0.002211635113) (22700,-0.002042293195) (22800,-0.00206282923) (22900,-0.002408474726) (23000,-0.00299631667) (23100,-0.002965116665) (23200,-0.003019779106) (23300,-0.003855472198) (23400,-0.004588799062) (23500,-0.004012366952) (23600,-0.005052867215) (23700,-0.004905429072) (23800,-0.00452925275)};
\addplot[pubRose, opacity=0.18, line width=0.35pt] coordinates {(100,-0.09799157828) (200,-0.1017213538) (300,-0.1034892748) (400,-0.1065390445) (500,-0.107235384) (600,-0.1096863076) (700,-0.1107180725) (800,-0.1125714974) (900,-0.1124981683) (1000,-0.1139022894) (1100,-0.1164816089) (1200,-0.1167337701) (1300,-0.1166401453) (1400,-0.1154206246) (1500,-0.1144644693) (1600,-0.1127289787) (1700,-0.1114848353) (1800,-0.1099770524) (1900,-0.1085665293) (2000,-0.106968011) (2100,-0.1059254833) (2200,-0.104475683) (2300,-0.1034136891) (2400,-0.1019884966) (2500,-0.1010212079) (2600,-0.100284788) (2700,-0.09826060534) (2800,-0.0960004963) (2900,-0.0952122502) (3000,-0.09369418994) (3100,-0.09087263793) (3200,-0.08990560248) (3300,-0.08979008421) (3400,-0.0890424259) (3500,-0.08724364564) (3600,-0.08447581902) (3700,-0.08277194872) (3800,-0.08018139787) (3900,-0.077436959) (4000,-0.07410782166) (4100,-0.07176496387) (4200,-0.07005913481) (4300,-0.0666796755) (4400,-0.06476383172) (4500,-0.06323322579) (4600,-0.06169647947) (4700,-0.06033728011) (4800,-0.05998637863) (4900,-0.0598385144) (5000,-0.05871693268) (5100,-0.05779684074) (5200,-0.05676317289) (5300,-0.05500244014) (5400,-0.05386513621) (5500,-0.05280084275) (5600,-0.05166564658) (5700,-0.05189858824) (5800,-0.05103610903) (5900,-0.04925784543) (6000,-0.0487446595) (6100,-0.04785487205) (6200,-0.04779400565) (6300,-0.04803303778) (6400,-0.04675439075) (6500,-0.0459465161) (6600,-0.04499603212) (6700,-0.04323240556) (6800,-0.04231153354) (6900,-0.0418089658) (7000,-0.04053279646) (7100,-0.03926439229) (7200,-0.03717812616) (7300,-0.03506780267) (7400,-0.03377156798) (7500,-0.03262658343) (7600,-0.03179886639) (7700,-0.03108869884) (7800,-0.02860269975) (7900,-0.02787840664) (8000,-0.02669216488) (8100,-0.02555696983) (8200,-0.02554586884) (8300,-0.02522390988) (8400,-0.02462122329) (8500,-0.02422458064) (8600,-0.02328044139) (8700,-0.02199210655) (8800,-0.0222144464) (8900,-0.02129128017) (9000,-0.02054083962) (9100,-0.0205410121) (9200,-0.0177752946) (9300,-0.01674911934) (9400,-0.0165001228) (9500,-0.01617337051) (9600,-0.01511646272) (9700,-0.01430295114) (9800,-0.01423089732) (9900,-0.01301606861) (10000,-0.01276169592) (10100,-0.01188314245) (10200,-0.0121844477) (10300,-0.012067765) (10400,-0.01189758647) (10500,-0.01105810944) (10600,-0.01146654282) (10700,-0.01124815848) (10800,-0.009614709119) (10900,-0.008930105582) (11000,-0.007606655528) (11100,-0.007162410003) (11200,-0.00737750156) (11300,-0.006248915259) (11400,-0.005490936729) (11500,-0.004081196984) (11600,-0.004104997491) (11700,-0.002851039026) (11800,-0.003184679605) (11900,-0.002957885203) (12000,-0.003768639153) (12100,-0.004016850682) (12200,-0.003384819327) (12300,-0.004036118175) (12400,-0.003767157151) (12500,-0.003857431529) (12600,-0.00283583907) (12700,-0.003169541134) (12800,-0.002854928284) (12900,-0.002773771703) (13000,-0.002932856313) (13100,-0.001630976691) (13200,-9.507228388e-05) (13300,0.0003804846987) (13400,0.0002811875049) (13500,3.615554306e-06) (13600,0.0002124628722) (13700,0.001142197088) (13800,0.001093690118) (13900,0.0004617631173) (14000,0.001216804725) (14100,0.001190501288) (14200,0.0007589187704) (14300,0.0008376188487) (14400,0.001270834033) (14500,0.001775911185) (14600,0.001945770453) (14700,0.001663180688) (14800,0.002141417356) (14900,0.002610256409) (15000,0.002240275871) (15100,0.00152722795) (15200,0.001213631569) (15300,0.001288935798) (15400,0.001917420246) (15500,0.002288274479) (15600,0.002259791421) (15700,0.002682402928) (15800,0.002894431481) (15900,0.003274807369) (16000,0.004149039101) (16100,0.006001121073) (16200,0.006489555666) (16300,0.007142076315) (16400,0.007800634159) (16500,0.007743904088) (16600,0.007447624591) (16700,0.006798132008) (16800,0.006950528373) (16900,0.006458942743) (17000,0.006198388932) (17100,0.005323465948) (17200,0.005558115465) (17300,0.005203799566) (17400,0.004023807601) (17500,0.003538394056) (17600,0.004040458403) (17700,0.003065011255) (17800,0.002341391522) (17900,0.002219793692) (18000,0.003170104933) (18100,0.002989241949) (18200,0.002081363726) (18300,0.001746256541) (18400,0.001878418991) (18500,0.002304028495) (18600,0.002557409796) (18700,0.003058274477) (18800,0.00413418099) (18900,0.005119371658) (19000,0.003414656978) (19100,0.002680025214) (19200,0.002788525209) (19300,0.002960291965) (19400,0.003159008123) (19500,0.002773666492) (19600,0.001630305377) (19700,0.002070197475) (19800,0.001198358031) (19900,0.0009087288869) (20000,0.001738785033) (20100,0.002563399787) (20200,0.002111631748) (20300,0.002158056013) (20400,0.00178460828) (20500,0.001581497633) (20600,0.00178351613) (20700,0.001231835544) (20800,0.001699658454) (20900,0.001397808091) (21000,0.001340689426) (21100,0.001534641403) (21200,0.001976800369) (21300,0.001595589929) (21400,0.002273128243) (21500,0.003619001841) (21600,0.003877761521) (21700,0.003617055874) (21800,0.002850950274) (21900,0.00280087131) (22000,0.002740795206) (22100,0.002818616427) (22200,0.003355602216) (22300,0.004015269608) (22400,0.003823390638) (22500,0.003031452361) (22600,0.003113274777) (22700,0.003027812403) (22800,0.003879600938) (22900,0.003994200693) (23000,0.004109430744) (23100,0.00366864294) (23200,0.002961418033) (23300,0.002021077345) (23400,0.001352627156) (23500,0.0008763232036) (23600,0.001124943723) (23700,0.002014935133) (23800,0.001157687989)};
\addplot[pubRose, opacity=0.18, line width=0.35pt] coordinates {(100,0.1477684379) (200,0.1512702554) (300,0.1512620548) (400,0.1506913304) (500,0.1497367799) (600,0.1483359163) (700,0.1470076527) (800,0.1448890921) (900,0.1454436729) (1000,0.1449387535) (1100,0.1448051959) (1200,0.1416462317) (1300,0.1399106994) (1400,0.1384395257) (1500,0.1372756302) (1600,0.1354045436) (1700,0.1349824965) (1800,0.1355966851) (1900,0.1325849175) (2000,0.1315022245) (2100,0.128778182) (2200,0.1299884424) (2300,0.1279322639) (2400,0.1275197819) (2500,0.1255014569) (2600,0.1243769646) (2700,0.1222027056) (2800,0.1186344236) (2900,0.1180691488) (3000,0.1161952823) (3100,0.1144562185) (3200,0.1111028709) (3300,0.109696155) (3400,0.1069489837) (3500,0.1065156408) (3600,0.1055930719) (3700,0.1044807717) (3800,0.1042704232) (3900,0.1024372615) (4000,0.1010512196) (4100,0.1003354676) (4200,0.1004224256) (4300,0.0994408831) (4400,0.09642287716) (4500,0.09360916689) (4600,0.09178875089) (4700,0.08903404102) (4800,0.08755246773) (4900,0.08484782726) (5000,0.0818716079) (5100,0.08115061596) (5200,0.07797442675) (5300,0.07549375445) (5400,0.07529476732) (5500,0.0745998472) (5600,0.07370762005) (5700,0.07248952426) (5800,0.07047995515) (5900,0.06957494169) (6000,0.06778096519) (6100,0.06481217444) (6200,0.06410007998) (6300,0.06415890977) (6400,0.06257343143) (6500,0.05916566998) (6600,0.05651610345) (6700,0.05445623733) (6800,0.05243617408) (6900,0.05151864663) (7000,0.0511211928) (7100,0.04830400273) (7200,0.04515031427) (7300,0.04251874872) (7400,0.04113370739) (7500,0.03942088) (7600,0.03703110777) (7700,0.03784517609) (7800,0.03624880947) (7900,0.03375390321) (8000,0.03179000095) (8100,0.03217679374) (8200,0.03113338947) (8300,0.03033826929) (8400,0.03010276426) (8500,0.02917326055) (8600,0.02918708548) (8700,0.02589831054) (8800,0.02503379472) (8900,0.02394455234) (9000,0.02193630943) (9100,0.01952657774) (9200,0.01692009117) (9300,0.01525570988) (9400,0.01190040102) (9500,0.01272330107) (9600,0.01091741985) (9700,0.009766642097) (9800,0.008440533187) (9900,0.006956610944) (10000,0.006656496994) (10100,0.005934516036) (10200,0.006976252444) (10300,0.004304745629) (10400,0.003923743322) (10500,0.0004417109405) (10600,0.0002142363795) (10700,-0.001358421742) (10800,-0.003088288977) (10900,-0.003909678443) (11000,-0.003921413666) (11100,-0.004833450832) (11200,-0.006273142167) (11300,-0.006238204718) (11400,-0.00757441693) (11500,-0.006511952565) (11600,-0.007472290122) (11700,-0.007642783155) (11800,-0.007348021702) (11900,-0.007503893715) (12000,-0.008841461618) (12100,-0.009898596467) (12200,-0.009355262574) (12300,-0.009249806218) (12400,-0.01007013815) (12500,-0.01188637179) (12600,-0.01297987772) (12700,-0.01245921638) (12800,-0.01410390998) (12900,-0.01510931293) (13000,-0.01612570316) (13100,-0.01590453994) (13200,-0.01795574329) (13300,-0.01891719424) (13400,-0.01812046436) (13500,-0.01852122867) (13600,-0.01966442917) (13700,-0.02113294117) (13800,-0.02021162957) (13900,-0.02028752081) (14000,-0.02084520031) (14100,-0.02229102924) (14200,-0.0221492555) (14300,-0.02222530693) (14400,-0.02304323912) (14500,-0.02334207259) (14600,-0.02376555484) (14700,-0.02380692437) (14800,-0.02456616983) (14900,-0.02577577401) (15000,-0.02520363703) (15100,-0.02519593015) (15200,-0.02555565368) (15300,-0.02603315394) (15400,-0.02641168535) (15500,-0.02661068365) (15600,-0.02640794255) (15700,-0.02720966134) (15800,-0.02728292029) (15900,-0.02743519936) (16000,-0.02870798204) (16100,-0.02820219062) (16200,-0.02886083201) (16300,-0.02865611427) (16400,-0.02945530452) (16500,-0.02921743114) (16600,-0.0288375821) (16700,-0.02833432667) (16800,-0.02829262242) (16900,-0.02849639542) (17000,-0.02800059561) (17100,-0.02887678891) (17200,-0.02877960596) (17300,-0.02976511326) (17400,-0.0289281141) (17500,-0.02887131236) (17600,-0.02826247197) (17700,-0.02924581245) (17800,-0.03009632677) (17900,-0.02897937521) (18000,-0.02932098005) (18100,-0.02989999857) (18200,-0.02989495788) (18300,-0.02881441247) (18400,-0.02926584762) (18500,-0.02933286093) (18600,-0.03015586827) (18700,-0.03012118284) (18800,-0.03004077356) (18900,-0.03014086038) (19000,-0.03010697849) (19100,-0.02979663201) (19200,-0.02998505179) (19300,-0.03087197337) (19400,-0.03088975586) (19500,-0.03044786397) (19600,-0.03060466051) (19700,-0.03105305471) (19800,-0.03074577358) (19900,-0.03077020794) (20000,-0.03069906831) (20100,-0.03011199068) (20200,-0.02951934841) (20300,-0.03009154256) (20400,-0.0295125939) (20500,-0.03101488892) (20600,-0.03012698404) (20700,-0.02933918368) (20800,-0.02851903383) (20900,-0.02885786742) (21000,-0.028638082) (21100,-0.0292964993) (21200,-0.02996386979) (21300,-0.02915222887) (21400,-0.02974691615) (21500,-0.02931063212) (21600,-0.03046504613) (21700,-0.03035286274) (21800,-0.03125629593) (21900,-0.03097740021) (22000,-0.03109128941) (22100,-0.02998781931) (22200,-0.0297899371) (22300,-0.02992957998) (22400,-0.03011617735) (22500,-0.03050153218) (22600,-0.03043602183) (22700,-0.03059939388) (22800,-0.03098147362) (22900,-0.03125894442) (23000,-0.03155326694) (23100,-0.03133238424) (23200,-0.03158026617) (23300,-0.03182604071) (23400,-0.03246841393) (23500,-0.03058382589) (23600,-0.03106821235) (23700,-0.03103306741) (23800,-0.03067180216)};
\addplot[pubRose, opacity=0.18, line width=0.35pt] coordinates {(100,0.03588190675) (200,0.03334009834) (300,0.03067751601) (400,0.02938580699) (500,0.02717638388) (600,0.02661322244) (700,0.02573568826) (800,0.02637435682) (900,0.02713559899) (1000,0.02707478814) (1100,0.02666268162) (1200,0.02575389761) (1300,0.0257178789) (1400,0.02689357903) (1500,0.0266561199) (1600,0.02710609008) (1700,0.02805065569) (1800,0.02758795731) (1900,0.02730974741) (2000,0.02882567979) (2100,0.02909908332) (2200,0.03017583545) (2300,0.03256773595) (2400,0.03260425534) (2500,0.03365832064) (2600,0.03396114409) (2700,0.0349197166) (2800,0.03604214694) (2900,0.03559769504) (3000,0.03496961482) (3100,0.03464598767) (3200,0.03474581949) (3300,0.03333865479) (3400,0.03433222473) (3500,0.0357417142) (3600,0.03567673173) (3700,0.03507526759) (3800,0.03453996386) (3900,0.03607039787) (4000,0.03606970645) (4100,0.03653925285) (4200,0.03684833124) (4300,0.03708615117) (4400,0.03552570827) (4500,0.03472367749) (4600,0.03488737606) (4700,0.03507455625) (4800,0.03551744744) (4900,0.03450156413) (5000,0.03402496874) (5100,0.03452383503) (5200,0.03462915085) (5300,0.03313089348) (5400,0.03333650753) (5500,0.03278698009) (5600,0.03301849794) (5700,0.03252296802) (5800,0.03175956775) (5900,0.03304738756) (6000,0.03369555008) (6100,0.03340662103) (6200,0.03429869767) (6300,0.03523428161) (6400,0.03584893551) (6500,0.03705893122) (6600,0.03675205745) (6700,0.0379952576) (6800,0.03811453432) (6900,0.03639893569) (7000,0.03605258726) (7100,0.03497511055) (7200,0.03329060189) (7300,0.03239272665) (7400,0.03193889726) (7500,0.03197799195) (7600,0.03205269109) (7700,0.03194871861) (7800,0.03227252979) (7900,0.03339323383) (8000,0.03434967156) (8100,0.03559740111) (8200,0.03578883689) (8300,0.03767314851) (8400,0.03767958879) (8500,0.03720373698) (8600,0.03737846538) (8700,0.03668627925) (8800,0.03682901636) (8900,0.03549889959) (9000,0.03524773493) (9100,0.03389352355) (9200,0.03323764224) (9300,0.0314862499) (9400,0.03163951952) (9500,0.03137609381) (9600,0.03148399796) (9700,0.03240030091) (9800,0.03273428585) (9900,0.03290157393) (10000,0.03257395551) (10100,0.03330176938) (10200,0.0343546262) (10300,0.03514703382) (10400,0.03448811211) (10500,0.03489721939) (10600,0.03510021605) (10700,0.03380611073) (10800,0.03363533858) (10900,0.03399659768) (11000,0.03428027034) (11100,0.03381005246) (11200,0.03290118501) (11300,0.03298069984) (11400,0.03259011842) (11500,0.03176201973) (11600,0.03158570249) (11700,0.03155637514) (11800,0.0299814444) (11900,0.03021370173) (12000,0.02935071066) (12100,0.02946180664) (12200,0.02968300413) (12300,0.02895189729) (12400,0.02996863164) (12500,0.03091213275) (12600,0.02977476809) (12700,0.02906018309) (12800,0.03050262649) (12900,0.02965218946) (13000,0.02900925558) (13100,0.029103739) (13200,0.02966114525) (13300,0.0288201984) (13400,0.02765679378) (13500,0.02670496516) (13600,0.02698524464) (13700,0.0279346684) (13800,0.02636957932) (13900,0.02617034242) (14000,0.02639346682) (14100,0.02587930448) (14200,0.02551792078) (14300,0.02526227338) (14400,0.02501429869) (14500,0.0246096869) (14600,0.02427792391) (14700,0.02411541259) (14800,0.0246845209) (14900,0.0247745079) (15000,0.02474714117) (15100,0.02464032052) (15200,0.02454979168) (15300,0.02547461633) (15400,0.02607812695) (15500,0.02631076761) (15600,0.02690610755) (15700,0.0255737938) (15800,0.02530632783) (15900,0.02516675126) (16000,0.02519514747) (16100,0.02571311202) (16200,0.02608249895) (16300,0.02558675781) (16400,0.02460732944) (16500,0.02436070722) (16600,0.02343565002) (16700,0.02320911083) (16800,0.02277575899) (16900,0.02204510504) (17000,0.02121823775) (17100,0.02017736295) (17200,0.01845404105) (17300,0.01816415107) (17400,0.01872351384) (17500,0.01773159374) (17600,0.01781512797) (17700,0.01786688417) (17800,0.01701275064) (17900,0.01771983709) (18000,0.01753627975) (18100,0.01672620606) (18200,0.01759929694) (18300,0.01760226563) (18400,0.01686693039) (18500,0.01757749813) (18600,0.01698295539) (18700,0.01679598363) (18800,0.01699154703) (18900,0.01595565816) (19000,0.01578369411) (19100,0.01653155042) (19200,0.01547228806) (19300,0.01564676203) (19400,0.01524275988) (19500,0.01469341135) (19600,0.01403597491) (19700,0.01383583434) (19800,0.0139501499) (19900,0.01476225965) (20000,0.01480228752) (20100,0.01432831697) (20200,0.01390291564) (20300,0.01289865011) (20400,0.0121566494) (20500,0.01178169455) (20600,0.01141061117) (20700,0.01099307151) (20800,0.01116102277) (20900,0.01037356318) (21000,0.009306969726) (21100,0.008914193837) (21200,0.009015077678) (21300,0.008975281985) (21400,0.008474360616) (21500,0.008429206559) (21600,0.008402669267) (21700,0.008045797865) (21800,0.008143456629) (21900,0.007903706538) (22000,0.00834003638) (22100,0.007920246036) (22200,0.007875369885) (22300,0.007989088842) (22400,0.008544290625) (22500,0.008404326439) (22600,0.008276547864) (22700,0.008467477886) (22800,0.007998957066) (22900,0.007330915704) (23000,0.006203189434) (23100,0.006065375998) (23200,0.005524787621) (23300,0.005003804772) (23400,0.004604324687) (23500,0.004423083446) (23600,0.0044017946) (23700,0.003987881227) (23800,0.003382074705)};
\addplot[pubRose, ultra thick] coordinates {(100,0.1693339944) (200,0.1688347496) (300,0.1711340348) (400,0.1730175968) (500,0.1732624918) (600,0.1739117131) (700,0.1724780351) (800,0.171765279) (900,0.1711943042) (1000,0.1707466044) (1100,0.169060427) (1200,0.1678299859) (1300,0.1664545916) (1400,0.1633841403) (1500,0.1616398513) (1600,0.1603085443) (1700,0.1582473949) (1800,0.1573848851) (1900,0.1553715982) (2000,0.1541460477) (2100,0.1536008723) (2200,0.1512369297) (2300,0.1498541981) (2400,0.1488298379) (2500,0.1483211882) (2600,0.1466461703) (2700,0.1464150652) (2800,0.1449484482) (2900,0.1450260311) (3000,0.1451433383) (3100,0.1436755821) (3200,0.1436652958) (3300,0.1430648863) (3400,0.1417791404) (3500,0.1393498719) (3600,0.1380400896) (3700,0.1355532978) (3800,0.133828954) (3900,0.1319590248) (4000,0.1300359584) (4100,0.1281890165) (4200,0.1265651401) (4300,0.1247351844) (4400,0.122570733) (4500,0.1214024778) (4600,0.1196522839) (4700,0.1177807681) (4800,0.1165174719) (4900,0.1147453107) (5000,0.1139222801) (5100,0.1121631306) (5200,0.110408609) (5300,0.1089646854) (5400,0.106867205) (5500,0.1053314462) (5600,0.1034281749) (5700,0.1027358141) (5800,0.1003883395) (5900,0.09869217426) (6000,0.09747080728) (6100,0.09486312419) (6200,0.09272064827) (6300,0.09018317349) (6400,0.08862626515) (6500,0.088682767) (6600,0.08708333299) (6700,0.08573579006) (6800,0.08454007059) (6900,0.08292716146) (7000,0.08045665175) (7100,0.07925298251) (7200,0.07786099892) (7300,0.07584233861) (7400,0.07268100847) (7500,0.06949727107) (7600,0.06775672473) (7700,0.06600734163) (7800,0.06431079619) (7900,0.06336851344) (8000,0.06101436783) (8100,0.05941061098) (8200,0.05853339266) (8300,0.05639779847) (8400,0.05468894262) (8500,0.0544762902) (8600,0.05203626584) (8700,0.05091006402) (8800,0.0496198982) (8900,0.04745744932) (9000,0.04685560837) (9100,0.04611178692) (9200,0.04426101036) (9300,0.04288261635) (9400,0.04163161414) (9500,0.04032171341) (9600,0.03928960077) (9700,0.03785556462) (9800,0.03689130414) (9900,0.03600286534) (10000,0.03438044656) (10100,0.03357571578) (10200,0.03322694078) (10300,0.03341027964) (10400,0.03223070163) (10500,0.03102392243) (10600,0.02852298823) (10700,0.02377182826) (10800,0.02290294513) (10900,0.02169125557) (11000,0.02119594307) (11100,0.0209721792) (11200,0.02050563195) (11300,0.01878460958) (11400,0.01791034065) (11500,0.0172334365) (11600,0.01725895446) (11700,0.01707487474) (11800,0.01621151059) (11900,0.01581449183) (12000,0.0150255458) (12100,0.01385644686) (12200,0.01256416898) (12300,0.01292766858) (12400,0.01267571907) (12500,0.0127778888) (12600,0.01260767486) (12700,0.01231523734) (12800,0.01109263706) (12900,0.009533525305) (13000,0.00785927628) (13100,0.006198281438) (13200,0.005767874319) (13300,0.0045390909) (13400,0.004695885901) (13500,0.004943500475) (13600,0.004763901878) (13700,0.003700191718) (13800,0.003039051217) (13900,0.002359369557) (14000,0.002381875017) (14100,0.001992589614) (14200,0.001235896786) (14300,0.001380477923) (14400,0.0005299603231) (14500,0.0002266328737) (14600,-0.0002341346743) (14700,-0.0006241913813) (14800,-5.856872167e-05) (14900,-9.224474497e-05) (15000,-0.0007586800886) (15100,-0.0008530274368) (15200,-0.000579942191) (15300,-0.001003566232) (15400,-5.698116001e-05) (15500,0.0001239763529) (15600,0.000259483184) (15700,0.0007575741769) (15800,0.0002903411896) (15900,-2.821059134e-05) (16000,0.0002249251429) (16100,-0.0002082240284) (16200,-0.0007470590976) (16300,-0.0002135032148) (16400,-0.001183836343) (16500,-0.001174666622) (16600,-0.0009397651073) (16700,-0.0009223504694) (16800,-0.0006135335632) (16900,-0.0006153243976) (17000,-0.0007287749968) (17100,-0.0007901799987) (17200,-0.001384620998) (17300,-0.0008441541739) (17400,-0.0012866388) (17500,-0.001464214832) (17600,-0.001876076545) (17700,-0.001134284934) (17800,-0.001228677518) (17900,-0.0007978093592) (18000,-0.00116355428) (18100,-0.001155396517) (18200,-0.0001014129914) (18300,-0.000555105113) (18400,-0.0002432591136) (18500,-0.000721816541) (18600,-0.0004056347098) (18700,-0.000987010001) (18800,-0.00138435617) (18900,-0.001281866319) (19000,-0.001818621058) (19100,-0.001902570976) (19200,-0.002098098277) (19300,-0.002163737479) (19400,-0.002490145485) (19500,-0.002902029162) (19600,-0.003074868569) (19700,-0.002781184597) (19800,-0.002403528654) (19900,-0.002297889782) (20000,-0.002379611853) (20100,-0.00280091732) (20200,-0.003306998065) (20300,-0.003660793728) (20400,-0.004142321154) (20500,-0.00363933956) (20600,-0.004245011738) (20700,-0.003933111075) (20800,-0.004082303395) (20900,-0.003968505964) (21000,-0.004221785101) (21100,-0.003594870045) (21200,-0.003427558317) (21300,-0.004217708454) (21400,-0.004568116253) (21500,-0.004554346482) (21600,-0.003924102995) (21700,-0.004529851892) (21800,-0.004197975325) (21900,-0.00419147008) (22000,-0.004347050144) (22100,-0.004863052932) (22200,-0.004686768308) (22300,-0.003851484387) (22400,-0.003776327644) (22500,-0.004294731499) (22600,-0.00458053726) (22700,-0.004483206063) (22800,-0.004716045427) (22900,-0.004767570797) (23000,-0.005457469892) (23100,-0.005153152592) (23200,-0.00529665871) (23300,-0.005462847708) (23400,-0.005582970765) (23500,-0.005070690047) (23600,-0.005080941335) (23700,-0.004842091461) (23800,-0.004651062835)};
\nextgroupplot[title={Twin-Q gap}, xlabel={Gradient update}, scaled y ticks=true]
\addplot[pubRose, opacity=0.18, line width=0.35pt] coordinates {(100,0.0206920132) (200,0.01932712179) (300,0.01954536264) (400,0.0191234299) (500,0.01924743354) (600,0.01901910827) (700,0.01914951764) (800,0.01922441088) (900,0.01895915903) (1000,0.0188837586) (1100,0.01903570313) (1200,0.01886257194) (1300,0.01884665731) (1400,0.01891938876) (1500,0.01853073165) (1600,0.01847061664) (1700,0.01804390233) (1800,0.01797728091) (1900,0.01823761705) (2000,0.01821180023) (2100,0.01797584184) (2200,0.01808175985) (2300,0.01788638514) (2400,0.01793525331) (2500,0.01814007983) (2600,0.01830595247) (2700,0.01862333752) (2800,0.01853367276) (2900,0.0183292333) (3000,0.01838123519) (3100,0.01815153714) (3200,0.01815880854) (3300,0.01807602774) (3400,0.01770163346) (3500,0.017631598) (3600,0.01760519054) (3700,0.01738527641) (3800,0.01704881676) (3900,0.01695324276) (4000,0.01676073745) (4100,0.01657379828) (4200,0.01637851987) (4300,0.01640584432) (4400,0.01667035036) (4500,0.01666476689) (4600,0.01643188726) (4700,0.01632048823) (4800,0.01634394284) (4900,0.01605701577) (5000,0.01598573923) (5100,0.01603723988) (5200,0.0161077708) (5300,0.01574432421) (5400,0.01523663932) (5500,0.01487248307) (5600,0.01465590326) (5700,0.01444519619) (5800,0.01432195483) (5900,0.01441736147) (6000,0.01443530601) (6100,0.01431818036) (6200,0.01434643259) (6300,0.01432171352) (6400,0.01441798639) (6500,0.01441999003) (6600,0.01445303578) (6700,0.01442002412) (6800,0.01427270547) (6900,0.01416487638) (7000,0.0139017785) (7100,0.01375739733) (7200,0.01356590362) (7300,0.01369176991) (7400,0.01378446193) (7500,0.01371461274) (7600,0.01358286925) (7700,0.01366498284) (7800,0.01366304196) (7900,0.01376867583) (8000,0.01371164201) (8100,0.01374139031) (8200,0.01358599206) (8300,0.01327181812) (8400,0.01322526578) (8500,0.01327749845) (8600,0.01321859388) (8700,0.01314839711) (8800,0.01332582291) (8900,0.01319311243) (9000,0.01318589961) (9100,0.01315551354) (9200,0.01313783238) (9300,0.01353564011) (9400,0.01339159664) (9500,0.01330249794) (9600,0.01344076693) (9700,0.01347344108) (9800,0.01336069265) (9900,0.01324440418) (10000,0.01322391341) (10100,0.01321312794) (10200,0.01310481234) (10300,0.01283572009) (10400,0.01281266985) (10500,0.01275450801) (10600,0.01264994284) (10700,0.01259948574) (10800,0.0125228419) (10900,0.01254866989) (11000,0.01230613301) (11100,0.01226924714) (11200,0.0122849605) (11300,0.01251529194) (11400,0.01239609635) (11500,0.01237700358) (11600,0.0123736728) (11700,0.0122159454) (11800,0.01240170468) (11900,0.0123185846) (12000,0.01245625783) (12100,0.01235371111) (12200,0.0124378344) (12300,0.01217234042) (12400,0.01208935203) (12500,0.0119474641) (12600,0.0118190689) (12700,0.01189787844) (12800,0.01170668276) (12900,0.01177997282) (13000,0.01180306775) (13100,0.0117382084) (13200,0.01166952662) (13300,0.01173677035) (13400,0.01181232277) (13500,0.01188613409) (13600,0.01189508624) (13700,0.01193458568) (13800,0.01185322944) (13900,0.01197170923) (14000,0.01215733616) (14100,0.01230452638) (14200,0.01229193183) (14300,0.01212773109) (14400,0.01220049635) (14500,0.01239369679) (14600,0.01231113449) (14700,0.01233574664) (14800,0.01258399878) (14900,0.01259567095) (15000,0.01227622712) (15100,0.01225780696) (15200,0.01229842771) (15300,0.01237306818) (15400,0.01232093545) (15500,0.01219880907) (15600,0.01226586355) (15700,0.01221311204) (15800,0.01209978843) (15900,0.01189541779) (16000,0.01197459344) (16100,0.01190715451) (16200,0.01194542898) (16300,0.01185112288) (16400,0.01175518744) (16500,0.01172649264) (16600,0.0116538926) (16700,0.0116956084) (16800,0.01169294706) (16900,0.01179903606) (17000,0.01187861003) (17100,0.01194645967) (17200,0.01196448607) (17300,0.01230142759) (17400,0.01231055176) (17500,0.01235419763) (17600,0.01261562286) (17700,0.01260961145) (17800,0.01274899859) (17900,0.01286207195) (18000,0.01278005941) (18100,0.01278879857) (18200,0.01274876287) (18300,0.01257563876) (18400,0.01288626073) (18500,0.01271656286) (18600,0.01266265912) (18700,0.01257196059) (18800,0.01238165079) (18900,0.01241830168) (19000,0.01238151304) (19100,0.01215724144) (19200,0.01220648512) (19300,0.0121821573) (19400,0.01195140462) (19500,0.01208590539) (19600,0.01205614572) (19700,0.01216398366) (19800,0.01233886098) (19900,0.01213013446) (20000,0.01237032022) (20100,0.01244307207) (20200,0.01235575667) (20300,0.01233470282) (20400,0.01229073564) (20500,0.01243811725) (20600,0.01233341889) (20700,0.01216071043) (20800,0.01199121932) (20900,0.01204339797) (21000,0.01212906437) (21100,0.01218793336) (21200,0.01224914985) (21300,0.01221281728) (21400,0.0124213268) (21500,0.01231823405) (21600,0.01233822843) (21700,0.01242068196) (21800,0.01241170326) (21900,0.01233993145) (22000,0.01195459422) (22100,0.01200911151) (22200,0.01190539822) (22300,0.01206321642) (22400,0.01194084045) (22500,0.01189057743) (22600,0.01216176264) (22700,0.01218657754) (22800,0.01232712381) (22900,0.01241924772) (23000,0.01267780028) (23100,0.0128079473) (23200,0.01283162823) (23300,0.01271591075) (23400,0.012696735) (23500,0.01270857928) (23600,0.01243686751) (23700,0.01243931577) (23800,0.01216137633)};
\addplot[pubRose, opacity=0.18, line width=0.35pt] coordinates {(100,0.0172390826) (200,0.01586418832) (300,0.01638036501) (400,0.01650785073) (500,0.01643353049) (600,0.01628518679) (700,0.01631602445) (800,0.01627728867) (900,0.01632510053) (1000,0.0165104867) (1100,0.01652836027) (1200,0.01700035483) (1300,0.01700207032) (1400,0.01691238005) (1500,0.01724002995) (1600,0.01784580592) (1700,0.01804875005) (1800,0.01842285525) (1900,0.0186722463) (2000,0.01875930671) (2100,0.01920041908) (2200,0.01911619734) (2300,0.01946135145) (2400,0.01987924986) (2500,0.01969074216) (2600,0.01974420268) (2700,0.01980348025) (2800,0.01947323959) (2900,0.01933639664) (3000,0.01919225641) (3100,0.01905568819) (3200,0.01902512312) (3300,0.01905215904) (3400,0.01904397663) (3500,0.01905444451) (3600,0.01881707963) (3700,0.01896506231) (3800,0.01908648703) (3900,0.01897794195) (4000,0.01909800004) (4100,0.01873240788) (4200,0.01872414201) (4300,0.01848228686) (4400,0.01814484) (4500,0.01835849211) (4600,0.01815950871) (4700,0.01784091741) (4800,0.01769525483) (4900,0.01770069804) (5000,0.01768004131) (5100,0.01811180487) (5200,0.01809336357) (5300,0.01780554093) (5400,0.01801160239) (5500,0.01805795655) (5600,0.01820293013) (5700,0.01806385759) (5800,0.01825215798) (5900,0.01830547098) (6000,0.01838405523) (6100,0.01801676694) (6200,0.01789130718) (6300,0.01801626738) (6400,0.01782495491) (6500,0.0177650502) (6600,0.01754882969) (6700,0.01767751537) (6800,0.01753841713) (6900,0.01755177043) (7000,0.01753077284) (7100,0.0176003186) (7200,0.01749770939) (7300,0.0175341757) (7400,0.01769275852) (7500,0.01743229311) (7600,0.01743774675) (7700,0.01755343415) (7800,0.01778584607) (7900,0.01782163493) (8000,0.01774472892) (8100,0.01763339154) (8200,0.01782636028) (8300,0.01750376187) (8400,0.01747146919) (8500,0.01740689427) (8600,0.0174034005) (8700,0.01730635632) (8800,0.01730527617) (8900,0.01747439597) (9000,0.01734907832) (9100,0.01773157194) (9200,0.01776030567) (9300,0.01805694159) (9400,0.01797576882) (9500,0.01796038318) (9600,0.01788760833) (9700,0.01787502449) (9800,0.01773711443) (9900,0.01738839615) (10000,0.01740394719) (10100,0.01697222255) (10200,0.01683534533) (10300,0.01682528425) (10400,0.01696485765) (10500,0.01711485144) (10600,0.01701946985) (10700,0.01702163834) (10800,0.01673633307) (10900,0.01679691076) (11000,0.01681453642) (11100,0.01671800707) (11200,0.0166366132) (11300,0.01660413332) (11400,0.0162523767) (11500,0.01617648806) (11600,0.01648950987) (11700,0.01626304742) (11800,0.01637420934) (11900,0.01643221751) (12000,0.0164325336) (12100,0.01660874244) (12200,0.01664749961) (12300,0.0167580599) (12400,0.01683739349) (12500,0.01708852556) (12600,0.01672318764) (12700,0.01680273116) (12800,0.01702901628) (12900,0.01712571308) (13000,0.01720047798) (13100,0.01756562535) (13200,0.01767520551) (13300,0.01763734519) (13400,0.01771423854) (13500,0.01756918412) (13600,0.01764006745) (13700,0.01772327982) (13800,0.0173962133) (13900,0.01731764469) (14000,0.01701843347) (14100,0.01648421902) (14200,0.01645361912) (14300,0.01617424758) (14400,0.0164078827) (14500,0.0164315532) (14600,0.0164525439) (14700,0.01638907613) (14800,0.0164733178) (14900,0.01618320486) (15000,0.01640902692) (15100,0.01658392558) (15200,0.01636703284) (15300,0.01673087459) (15400,0.01649789549) (15500,0.0163307827) (15600,0.01615696345) (15700,0.01608432271) (15800,0.01608881168) (15900,0.01629836131) (16000,0.01608953355) (16100,0.01585760666) (16200,0.01611271231) (16300,0.01580032622) (16400,0.01557741035) (16500,0.01545819845) (16600,0.01566459835) (16700,0.01560309175) (16800,0.01532986304) (16900,0.01518728454) (17000,0.01544361422) (17100,0.01564210961) (17200,0.0152647675) (17300,0.01526134638) (17400,0.0155321788) (17500,0.01547401361) (17600,0.01561820097) (17700,0.01567612216) (17800,0.01595909139) (17900,0.01585146068) (18000,0.01584285507) (18100,0.01584689254) (18200,0.01597486129) (18300,0.01594932759) (18400,0.01585809039) (18500,0.01599308262) (18600,0.01585827721) (18700,0.01615237435) (18800,0.01612121314) (18900,0.01630269019) (19000,0.01622076826) (19100,0.01588026285) (19200,0.01584538706) (19300,0.015957186) (19400,0.01584336758) (19500,0.0156549491) (19600,0.01557056485) (19700,0.01528977379) (19800,0.01556446655) (19900,0.01541431872) (20000,0.01514607221) (20100,0.0152559909) (20200,0.01545142485) (20300,0.01532032816) (20400,0.01548580183) (20500,0.01560519366) (20600,0.01560134813) (20700,0.01542317783) (20800,0.01502103144) (20900,0.01504298309) (21000,0.01489062347) (21100,0.01483292487) (21200,0.01465992574) (21300,0.01476231851) (21400,0.01462350031) (21500,0.01442924794) (21600,0.01438202225) (21700,0.01458882671) (21800,0.01463907678) (21900,0.01472453326) (22000,0.01490215193) (22100,0.01474120747) (22200,0.01475301888) (22300,0.01454957388) (22400,0.01443630522) (22500,0.01463098116) (22600,0.01439355593) (22700,0.01421442945) (22800,0.014179081) (22900,0.01393939899) (23000,0.01396558778) (23100,0.01428786805) (23200,0.01436662572) (23300,0.01432893863) (23400,0.01450209031) (23500,0.01459079785) (23600,0.01480088606) (23700,0.01501078662) (23800,0.01470597815)};
\addplot[pubRose, opacity=0.18, line width=0.35pt] coordinates {(100,0.02949843556) (200,0.02747700922) (300,0.02625106027) (400,0.02724464424) (500,0.02748629227) (600,0.02890347255) (700,0.02969095537) (800,0.03113577235) (900,0.031822645) (1000,0.03172075208) (1100,0.03248592) (1200,0.03326546084) (1300,0.03362543862) (1400,0.03324567247) (1500,0.03281729091) (1600,0.03191318307) (1700,0.0309573194) (1800,0.02950828187) (1900,0.02823155783) (2000,0.02762280144) (2100,0.0259499386) (2200,0.02465172112) (2300,0.02414876428) (2400,0.02380979843) (2500,0.02367071714) (2600,0.02284474485) (2700,0.02249096502) (2800,0.02178374175) (2900,0.02095474843) (3000,0.02053497508) (3100,0.02015466727) (3200,0.02015240565) (3300,0.01939335689) (3400,0.01879474483) (3500,0.01814851314) (3600,0.01777916383) (3700,0.01733507812) (3800,0.01711582802) (3900,0.01721263342) (4000,0.01669199169) (4100,0.01688854881) (4200,0.01646109913) (4300,0.016613492) (4400,0.01656890884) (4500,0.01629580297) (4600,0.01615845971) (4700,0.01603658218) (4800,0.01591600608) (4900,0.01560686445) (5000,0.01543698078) (5100,0.01501224143) (5200,0.0148855147) (5300,0.0146830718) (5400,0.01457581064) (5500,0.01502901064) (5600,0.01507946048) (5700,0.01492651347) (5800,0.01460584998) (5900,0.01461050287) (6000,0.0144543929) (6100,0.01445942381) (6200,0.01443505827) (6300,0.01423691437) (6400,0.0139660012) (6500,0.01355825597) (6600,0.0134556449) (6700,0.01319784084) (6800,0.01328692939) (6900,0.01319806846) (7000,0.01323636668) (7100,0.01298155654) (7200,0.01284229523) (7300,0.01294118799) (7400,0.01287621781) (7500,0.01272487417) (7600,0.01259966129) (7700,0.01273367852) (7800,0.0127178031) (7900,0.01267855903) (8000,0.01273453785) (8100,0.01274800459) (8200,0.01268951306) (8300,0.01259047668) (8400,0.01252089459) (8500,0.0125643489) (8600,0.01270385254) (8700,0.01246259799) (8800,0.01230157968) (8900,0.01210280722) (9000,0.01214139601) (9100,0.01217426574) (9200,0.0121687728) (9300,0.01211985704) (9400,0.01231484786) (9500,0.01208115891) (9600,0.01187487906) (9700,0.01196649382) (9800,0.01189646702) (9900,0.01182063734) (10000,0.01137564462) (10100,0.01121043153) (10200,0.01104775965) (10300,0.01089635417) (10400,0.01052266927) (10500,0.01038199319) (10600,0.0103645537) (10700,0.0104494934) (10800,0.01036967849) (10900,0.01033156347) (11000,0.01046891958) (11100,0.01050579399) (11200,0.01060747188) (11300,0.01051472211) (11400,0.01060636118) (11500,0.01070772707) (11600,0.0106889558) (11700,0.01037548063) (11800,0.01022232706) (11900,0.01012710258) (12000,0.009983071871) (12100,0.009830527194) (12200,0.009585522115) (12300,0.00976760583) (12400,0.009723210428) (12500,0.009613408707) (12600,0.009406642988) (12700,0.009359349217) (12800,0.009290669672) (12900,0.009273731522) (13000,0.00927737141) (13100,0.009371688124) (13200,0.00946760457) (13300,0.00924576791) (13400,0.009182257578) (13500,0.009205449745) (13600,0.009214412794) (13700,0.009347511362) (13800,0.009315185621) (13900,0.009390215017) (14000,0.009375380352) (14100,0.009009044943) (14200,0.0087674879) (14300,0.008742859447) (14400,0.008758491604) (14500,0.008767103078) (14600,0.008774464065) (14700,0.008621839294) (14800,0.008795516426) (14900,0.00880255145) (15000,0.008832977945) (15100,0.009022365883) (15200,0.009100799263) (15300,0.009118931834) (15400,0.009073450882) (15500,0.008980429266) (15600,0.008880309295) (15700,0.008979800716) (15800,0.008943632804) (15900,0.008968512528) (16000,0.008957821317) (16100,0.008957381453) (16200,0.009101324622) (16300,0.009248094354) (16400,0.009203637019) (16500,0.009252334572) (16600,0.009330942202) (16700,0.009299017675) (16800,0.009332827758) (16900,0.00933025647) (17000,0.009467701241) (17100,0.00946426373) (17200,0.009292024281) (17300,0.009197404981) (17400,0.009214854334) (17500,0.009082694259) (17600,0.009141905233) (17700,0.009095719177) (17800,0.009000262246) (17900,0.008886892069) (18000,0.008754723985) (18100,0.009030369855) (18200,0.009095196147) (18300,0.009113984648) (18400,0.009248027392) (18500,0.00943424562) (18600,0.009537138604) (18700,0.009563152026) (18800,0.009537120722) (18900,0.009599102847) (19000,0.00962255355) (19100,0.009367947001) (19200,0.009442482609) (19300,0.009542368911) (19400,0.009513888601) (19500,0.009477892704) (19600,0.009378235135) (19700,0.009414839931) (19800,0.00946953306) (19900,0.009559448157) (20000,0.009487195034) (20100,0.00958870789) (20200,0.009452362079) (20300,0.009325293545) (20400,0.009319272358) (20500,0.009286565147) (20600,0.009414026886) (20700,0.009340678714) (20800,0.009630869143) (20900,0.00964869) (21000,0.009779961873) (21100,0.009799270611) (21200,0.01002573362) (21300,0.0100451557) (21400,0.01016081665) (21500,0.01013342636) (21600,0.009978221636) (21700,0.01006641826) (21800,0.009948504437) (21900,0.009921387117) (22000,0.009990899824) (22100,0.009967910685) (22200,0.009909251146) (22300,0.009998341743) (22400,0.00978194885) (22500,0.00994082503) (22600,0.01001965683) (22700,0.01003485695) (22800,0.009786637221) (22900,0.009656375833) (23000,0.009543754905) (23100,0.00945922425) (23200,0.009485640097) (23300,0.009528269339) (23400,0.00958501827) (23500,0.009547284804) (23600,0.009438764025) (23700,0.009444104228) (23800,0.009570384771)};
\addplot[pubRose, opacity=0.18, line width=0.35pt] coordinates {(100,0.02838714421) (200,0.02782628872) (300,0.02561270135) (400,0.02389849117) (500,0.02234762572) (600,0.02101753559) (700,0.02017629759) (800,0.01950265921) (900,0.01906689639) (1000,0.01863045311) (1100,0.01717573581) (1200,0.01596594835) (1300,0.01514356276) (1400,0.01457553571) (1500,0.01460018111) (1600,0.0147363632) (1700,0.01461469717) (1800,0.01469687251) (1900,0.01466046125) (2000,0.01485513467) (2100,0.01508610845) (2200,0.01480787331) (2300,0.01487184195) (2400,0.01508153146) (2500,0.01490281923) (2600,0.01482994603) (2700,0.01526339501) (2800,0.01499439403) (2900,0.01494837133) (3000,0.01485280683) (3100,0.01491473736) (3200,0.01501918416) (3300,0.01544782538) (3400,0.01538191782) (3500,0.0156671063) (3600,0.0157352726) (3700,0.01552000334) (3800,0.01552018421) (3900,0.01552705653) (4000,0.01541170375) (4100,0.01535628038) (4200,0.01547809681) (4300,0.01499091303) (4400,0.01520635411) (4500,0.01507295519) (4600,0.01517549381) (4700,0.01513431035) (4800,0.01575166974) (4900,0.01610318674) (5000,0.01626497824) (5100,0.01618234375) (5200,0.01622067355) (5300,0.01642425414) (5400,0.01627758043) (5500,0.01633372055) (5600,0.01603637151) (5700,0.01596180359) (5800,0.01574689159) (5900,0.01555130137) (6000,0.01562215229) (6100,0.0156019968) (6200,0.01577446461) (6300,0.01578292549) (6400,0.01600098731) (6500,0.01596959913) (6600,0.01608882099) (6700,0.01605711281) (6800,0.01594026983) (6900,0.01572722103) (7000,0.0156374773) (7100,0.01573066721) (7200,0.01593343029) (7300,0.01587959835) (7400,0.01577927368) (7500,0.01578207063) (7600,0.0156951651) (7700,0.01567051355) (7800,0.01563403346) (7900,0.0158357678) (8000,0.01583199892) (8100,0.01574380845) (8200,0.01560184248) (8300,0.0156996415) (8400,0.01564575257) (8500,0.01571310153) (8600,0.01599233728) (8700,0.01621720484) (8800,0.01637439402) (8900,0.01625494128) (9000,0.01616005925) (9100,0.01620325753) (9200,0.01597357094) (9300,0.01588324057) (9400,0.0159494211) (9500,0.01588498401) (9600,0.0156608385) (9700,0.01547710644) (9800,0.01512582889) (9900,0.01524821399) (10000,0.01525775064) (10100,0.01512404205) (10200,0.01492560785) (10300,0.01488906275) (10400,0.01500052493) (10500,0.01460257052) (10600,0.01461200733) (10700,0.01456067767) (10800,0.01479365304) (10900,0.01470761728) (11000,0.01466492033) (11100,0.01481647193) (11200,0.01503161481) (11300,0.01527977036) (11400,0.0150472465) (11500,0.01543461457) (11600,0.01547412248) (11700,0.01561699705) (11800,0.01554059433) (11900,0.0154446899) (12000,0.01547588902) (12100,0.01526467735) (12200,0.01518500913) (12300,0.01513231713) (12400,0.0152063217) (12500,0.01493562534) (12600,0.01475014184) (12700,0.01469003409) (12800,0.01481909119) (12900,0.01509550959) (13000,0.01500698924) (13100,0.01516467771) (13200,0.01527951388) (13300,0.01525882846) (13400,0.01521893172) (13500,0.01555655124) (13600,0.01578691034) (13700,0.01558160065) (13800,0.01539537134) (13900,0.01496968875) (14000,0.01487267241) (14100,0.01462162845) (14200,0.01435379833) (14300,0.01457919106) (14400,0.01454445478) (14500,0.01411730759) (14600,0.01406826265) (14700,0.01397068445) (14800,0.01413495857) (14900,0.01427827161) (15000,0.01426030463) (15100,0.01453385986) (15200,0.01463580597) (15300,0.01417455599) (15400,0.01415831335) (15500,0.01441779705) (15600,0.0144628332) (15700,0.01485174289) (15800,0.01464688489) (15900,0.01457080748) (16000,0.01468940983) (16100,0.01491639419) (16200,0.01518815244) (16300,0.0154218669) (16400,0.01530197496) (16500,0.01512552761) (16600,0.01501377374) (16700,0.01488425732) (16800,0.01491088578) (16900,0.01512161512) (17000,0.01541808331) (17100,0.01517069126) (17200,0.01491114246) (17300,0.01463322369) (17400,0.01493943343) (17500,0.0150459948) (17600,0.01505725281) (17700,0.01516211769) (17800,0.01522969231) (17900,0.01541083064) (18000,0.01527794041) (18100,0.01527926782) (18200,0.01562496917) (18300,0.01569614913) (18400,0.01556187347) (18500,0.01576202642) (18600,0.01577623058) (18700,0.01568989856) (18800,0.01578135053) (18900,0.01547558466) (19000,0.01535796346) (19100,0.0153361395) (19200,0.01515992098) (19300,0.01521943631) (19400,0.01515635261) (19500,0.01494682031) (19600,0.01496302672) (19700,0.01496846462) (19800,0.01475172145) (19900,0.01477407999) (20000,0.01472500842) (20100,0.0148199847) (20200,0.01477706395) (20300,0.01477781683) (20400,0.01461691596) (20500,0.01453311127) (20600,0.01481197709) (20700,0.01482573757) (20800,0.01495070094) (20900,0.0150888633) (21000,0.01505735517) (21100,0.01483587492) (21200,0.01477964185) (21300,0.01482003406) (21400,0.01511216238) (21500,0.0150776268) (21600,0.01465380294) (21700,0.01463522138) (21800,0.01444704821) (21900,0.01421243912) (22000,0.0143000151) (22100,0.01426827908) (22200,0.01405483987) (22300,0.01388382465) (22400,0.01369038615) (22500,0.0136834681) (22600,0.01380762048) (22700,0.01384621188) (22800,0.01390466802) (22900,0.01384529946) (23000,0.01378110247) (23100,0.0138970688) (23200,0.0140748105) (23300,0.01407175343) (23400,0.01402420914) (23500,0.01400466245) (23600,0.01401003525) (23700,0.01413653977) (23800,0.0142781496)};
\addplot[pubRose, opacity=0.18, line width=0.35pt] coordinates {(100,0.02802253701) (200,0.02490384784) (300,0.02343767261) (400,0.02321192343) (500,0.02321163937) (600,0.0224298158) (700,0.02187171126) (800,0.02149569406) (900,0.02103927111) (1000,0.02090404965) (1100,0.01992153134) (1200,0.0195252711) (1300,0.01948452219) (1400,0.0189991653) (1500,0.01873562001) (1600,0.01863138862) (1700,0.01871607453) (1800,0.01860847194) (1900,0.01883327775) (2000,0.01859019585) (2100,0.0185391143) (2200,0.01852360256) (2300,0.01823435519) (2400,0.01828837991) (2500,0.01796793267) (2600,0.01802604236) (2700,0.01785205863) (2800,0.01773074958) (2900,0.01737271659) (3000,0.01728270538) (3100,0.01733122095) (3200,0.01722927243) (3300,0.01704380568) (3400,0.0169214081) (3500,0.01682496965) (3600,0.01686729193) (3700,0.01643590629) (3800,0.01628474221) (3900,0.01645949967) (4000,0.01645025592) (4100,0.01609689621) (4200,0.01605389146) (4300,0.01626837393) (4400,0.01593895769) (4500,0.01576104024) (4600,0.01560340775) (4700,0.01599703031) (4800,0.0160207564) (4900,0.01572266193) (5000,0.01571815247) (5100,0.01584629919) (5200,0.01569498368) (5300,0.01541196322) (5400,0.01576147722) (5500,0.01569451671) (5600,0.01542376168) (5700,0.01515099853) (5800,0.01501766788) (5900,0.01499191085) (6000,0.01493760822) (6100,0.01460835272) (6200,0.01446812088) (6300,0.01442334568) (6400,0.01393224681) (6500,0.01412042677) (6600,0.01418610737) (6700,0.01397659462) (6800,0.0138798533) (6900,0.01401887788) (7000,0.01388202729) (7100,0.01407691762) (7200,0.01430115197) (7300,0.01430751448) (7400,0.01437191386) (7500,0.01449077055) (7600,0.01454524603) (7700,0.01459345408) (7800,0.01474363804) (7900,0.01453165188) (8000,0.01471974794) (8100,0.01460701348) (8200,0.01480886834) (8300,0.01484280564) (8400,0.01518205414) (8500,0.01477070134) (8600,0.01462585302) (8700,0.01482342407) (8800,0.0148333489) (8900,0.0148841083) (9000,0.01484505534) (9100,0.01498006526) (9200,0.01449386962) (9300,0.01456330614) (9400,0.01430397248) (9500,0.01430655178) (9600,0.01429700209) (9700,0.01425643675) (9800,0.01423947448) (9900,0.01421504039) (10000,0.0140978232) (10100,0.01418254636) (10200,0.01432114383) (10300,0.01437175842) (10400,0.01453188332) (10500,0.01464730361) (10600,0.01468958529) (10700,0.01468317155) (10800,0.01469470896) (10900,0.01461170483) (11000,0.01454112828) (11100,0.01444002567) (11200,0.01440716907) (11300,0.01419056542) (11400,0.01446394753) (11500,0.01439013379) (11600,0.01426182603) (11700,0.01437346041) (11800,0.0142645726) (11900,0.01436199415) (12000,0.01439422704) (12100,0.01433181046) (12200,0.01414665645) (12300,0.01431691153) (12400,0.01397597631) (12500,0.01413829196) (12600,0.01436911188) (12700,0.01423424045) (12800,0.01433647592) (12900,0.01418238934) (13000,0.01416184111) (13100,0.0142000881) (13200,0.01446747882) (13300,0.01452346025) (13400,0.0146544571) (13500,0.01447809087) (13600,0.0142418962) (13700,0.01417578356) (13800,0.01425776519) (13900,0.01444512103) (14000,0.01436699349) (14100,0.01425625933) (14200,0.01433221856) (14300,0.0140686295) (14400,0.01405196255) (14500,0.01414076304) (14600,0.01441410128) (14700,0.014617287) (14800,0.01448507793) (14900,0.01432053093) (15000,0.01455711806) (15100,0.014742014) (15200,0.01463973308) (15300,0.01471696794) (15400,0.01462503579) (15500,0.01447235812) (15600,0.01414422104) (15700,0.01393448869) (15800,0.01405011164) (15900,0.01402208935) (16000,0.01373865362) (16100,0.0138054572) (16200,0.01383473584) (16300,0.01402744874) (16400,0.01408434035) (16500,0.01424398841) (16600,0.01436416209) (16700,0.01453183275) (16800,0.01464506425) (16900,0.01480811834) (17000,0.01516663656) (17100,0.01518653743) (17200,0.01499859821) (17300,0.01495438125) (17400,0.01491677752) (17500,0.01491803462) (17600,0.01498536477) (17700,0.0150721618) (17800,0.01478780704) (17900,0.01477013575) (18000,0.0145562592) (18100,0.01453241492) (18200,0.01490473123) (18300,0.01495761732) (18400,0.01494429195) (18500,0.01501323348) (18600,0.01487444593) (18700,0.01463687485) (18800,0.01460923133) (18900,0.01467237556) (19000,0.01472289693) (19100,0.01470344346) (19200,0.01459539682) (19300,0.01446011933) (19400,0.01447508326) (19500,0.01428449312) (19600,0.01434368799) (19700,0.01454764167) (19800,0.01447320078) (19900,0.0143397172) (20000,0.01414838452) (20100,0.01403779509) (20200,0.01388934487) (20300,0.01366844699) (20400,0.01356364992) (20500,0.01353644887) (20600,0.01352333771) (20700,0.01346584586) (20800,0.01375550618) (20900,0.0138453858) (21000,0.01406549141) (21100,0.01408124901) (21200,0.01415819991) (21300,0.01439187936) (21400,0.01452431753) (21500,0.01444503656) (21600,0.01439261781) (21700,0.01468965793) (21800,0.01442693463) (21900,0.01419271957) (22000,0.01412663963) (22100,0.01407658067) (22200,0.01411011834) (22300,0.01408286802) (22400,0.01394953383) (22500,0.01416296726) (22600,0.01431444744) (22700,0.01399874007) (22800,0.01403606348) (22900,0.01402871208) (23000,0.01403817097) (23100,0.01406761631) (23200,0.01404923387) (23300,0.01411961196) (23400,0.0141751945) (23500,0.01423956249) (23600,0.01403151443) (23700,0.01408402557) (23800,0.01425004704)};
\addplot[pubRose, opacity=0.18, line width=0.35pt] coordinates {(100,0.02489212155) (200,0.02316091023) (300,0.02215083502) (400,0.0213231747) (500,0.0211626444) (600,0.02088924932) (700,0.02090930992) (800,0.02037544735) (900,0.02049249659) (1000,0.02056479696) (1100,0.01988293361) (1200,0.01937173214) (1300,0.01911032088) (1400,0.01919076182) (1500,0.01911667511) (1600,0.01874963362) (1700,0.01827809699) (1800,0.01846522018) (1900,0.01813194081) (2000,0.01762750279) (2100,0.01779986694) (2200,0.01781319864) (2300,0.01795588695) (2400,0.01797589324) (2500,0.01777970381) (2600,0.01797491498) (2700,0.01825068165) (2800,0.01821880322) (2900,0.0185155483) (3000,0.01865340434) (3100,0.01839022655) (3200,0.01893778872) (3300,0.01911147255) (3400,0.01897789035) (3500,0.01896192227) (3600,0.01906322483) (3700,0.01917856168) (3800,0.01915525682) (3900,0.0191476373) (4000,0.01925762817) (4100,0.01946208738) (4200,0.01895303987) (4300,0.01879144646) (4400,0.01856697723) (4500,0.01854752582) (4600,0.01834969949) (4700,0.01818242837) (4800,0.01806191951) (4900,0.01793399062) (5000,0.01781946588) (5100,0.01756817345) (5200,0.01797433663) (5300,0.01776317637) (5400,0.01799398661) (5500,0.01791068595) (5600,0.01782901306) (5700,0.01779515315) (5800,0.01769796256) (5900,0.017575421) (6000,0.01770318542) (6100,0.01780627239) (6200,0.01745791137) (6300,0.01736932416) (6400,0.01716729626) (6500,0.01694460716) (6600,0.01690916549) (6700,0.01661149301) (6800,0.01669170447) (6900,0.01675772704) (7000,0.01653621458) (7100,0.01636509988) (7200,0.01616082415) (7300,0.01623377409) (7400,0.01612739228) (7500,0.01641013976) (7600,0.0166024629) (7700,0.01667918265) (7800,0.01655371282) (7900,0.01633080151) (8000,0.01637059432) (8100,0.01640508324) (8200,0.01632969845) (8300,0.01630228925) (8400,0.01639542952) (8500,0.01627123915) (8600,0.01615255196) (8700,0.01641137768) (8800,0.01631838912) (8900,0.01623523096) (9000,0.01590141989) (9100,0.01576923486) (9200,0.01577825639) (9300,0.01561017549) (9400,0.01550421277) (9500,0.01572453501) (9600,0.0155536727) (9700,0.01530176615) (9800,0.01529398765) (9900,0.01525049657) (10000,0.01533062672) (10100,0.0156498705) (10200,0.015643779) (10300,0.0159308874) (10400,0.01584197087) (10500,0.01565647954) (10600,0.01580272922) (10700,0.01567617022) (10800,0.01567858309) (10900,0.01563069411) (11000,0.01576021332) (11100,0.01533497572) (11200,0.0155207851) (11300,0.01524020964) (11400,0.01540857693) (11500,0.01520820279) (11600,0.0150804013) (11700,0.01524425643) (11800,0.01513201827) (11900,0.01500952076) (12000,0.01491945852) (12100,0.01522221956) (12200,0.01511360304) (12300,0.01524855504) (12400,0.01513824333) (12500,0.01504806997) (12600,0.01517537022) (12700,0.01494092988) (12800,0.0152523932) (12900,0.01541207591) (13000,0.01536831046) (13100,0.015112809) (13200,0.01503465418) (13300,0.0147878876) (13400,0.01479107859) (13500,0.01480358113) (13600,0.01453296822) (13700,0.0147034131) (13800,0.01433881801) (13900,0.01409600405) (14000,0.01402918724) (14100,0.0140106353) (14200,0.01394914789) (14300,0.01378621357) (14400,0.01362430202) (14500,0.01358843669) (14600,0.01354153967) (14700,0.01342129884) (14800,0.0135069889) (14900,0.01359851798) (15000,0.0138032685) (15100,0.01391016785) (15200,0.01396616716) (15300,0.01410792088) (15400,0.0140232739) (15500,0.01389672253) (15600,0.01390618701) (15700,0.01388980169) (15800,0.01401341129) (15900,0.01398656396) (16000,0.01378267864) (16100,0.01359496042) (16200,0.01373736691) (16300,0.01381599167) (16400,0.01403437229) (16500,0.01425176188) (16600,0.01452693874) (16700,0.01444372255) (16800,0.01444004858) (16900,0.01459887214) (17000,0.01460143831) (17100,0.01486356743) (17200,0.01476016585) (17300,0.0147858995) (17400,0.01463188445) (17500,0.01459211782) (17600,0.01468175873) (17700,0.01478119334) (17800,0.01461154763) (17900,0.01426377818) (18000,0.01403776854) (18100,0.01358974725) (18200,0.01343231108) (18300,0.0133707785) (18400,0.01354148956) (18500,0.01335063446) (18600,0.0128461251) (18700,0.01276415065) (18800,0.01286275899) (18900,0.01283350606) (19000,0.01297712224) (19100,0.0130577771) (19200,0.01297002155) (19300,0.01293609049) (19400,0.01262832182) (19500,0.01268461393) (19600,0.01300442852) (19700,0.01316071711) (19800,0.01304358933) (19900,0.01321097221) (20000,0.01313711638) (20100,0.01345726307) (20200,0.01356567983) (20300,0.0134139725) (20400,0.01366476454) (20500,0.01385471737) (20600,0.01376074534) (20700,0.01343272664) (20800,0.01329658041) (20900,0.01335317204) (21000,0.01337860627) (21100,0.01307141948) (21200,0.01288745208) (21300,0.01300319806) (21400,0.01275944458) (21500,0.01269713063) (21600,0.01269358713) (21700,0.0128203446) (21800,0.01289518652) (21900,0.01282856278) (22000,0.01284531597) (22100,0.01280427361) (22200,0.01313098036) (22300,0.0132961777) (22400,0.01316593923) (22500,0.01312644612) (22600,0.0129265504) (22700,0.01277402937) (22800,0.01285194103) (22900,0.01284584776) (23000,0.01310833488) (23100,0.01306802286) (23200,0.01288878219) (23300,0.01278551603) (23400,0.01305630887) (23500,0.01317010578) (23600,0.01339509925) (23700,0.0133718933) (23800,0.01333231218)};
\addplot[pubRose, opacity=0.18, line width=0.35pt] coordinates {(100,0.02144422382) (200,0.02029437013) (300,0.02002667263) (400,0.02086854819) (500,0.02045090385) (600,0.02011971238) (700,0.02027889687) (800,0.02030692925) (900,0.02057920541) (1000,0.02053607423) (1100,0.02059734575) (1200,0.02051382698) (1300,0.02064296901) (1400,0.02020191513) (1500,0.02044233512) (1600,0.02080155648) (1700,0.02105988916) (1800,0.02111724932) (1900,0.0208504336) (2000,0.02069435604) (2100,0.02070802245) (2200,0.02100563403) (2300,0.02099979948) (2400,0.02140645739) (2500,0.02124485634) (2600,0.02096643094) (2700,0.02065859437) (2800,0.02054875307) (2900,0.02037690654) (3000,0.0206207877) (3100,0.0201895956) (3200,0.01988477819) (3300,0.01962452941) (3400,0.01908631716) (3500,0.01896118782) (3600,0.01902923249) (3700,0.01880258154) (3800,0.01861529667) (3900,0.01887676213) (4000,0.01879153866) (4100,0.01880736798) (4200,0.01868851222) (4300,0.0186811056) (4400,0.01887105498) (4500,0.01876095571) (4600,0.01859634295) (4700,0.01854671184) (4800,0.01852815263) (4900,0.01828844696) (5000,0.01817117967) (5100,0.01806496456) (5200,0.01814617068) (5300,0.01805407964) (5400,0.01774216574) (5500,0.01801993176) (5600,0.01771300016) (5700,0.01744638523) (5800,0.0171299289) (5900,0.01676463122) (6000,0.01653241599) (6100,0.01621596953) (6200,0.01594518116) (6300,0.01561769759) (6400,0.01567268474) (6500,0.0151801317) (6600,0.01501323543) (6700,0.01510121953) (6800,0.015084254) (6900,0.01480121473) (7000,0.01455322383) (7100,0.01471716007) (7200,0.01449438399) (7300,0.01435540039) (7400,0.01402159929) (7500,0.01381609403) (7600,0.01389848264) (7700,0.01373122912) (7800,0.01351849642) (7900,0.01356812334) (8000,0.01347501269) (8100,0.01328071225) (8200,0.01338583622) (8300,0.01343329903) (8400,0.01322623026) (8500,0.01318830475) (8600,0.01300095776) (8700,0.01282050274) (8800,0.01280192444) (8900,0.01277504424) (9000,0.01257147454) (9100,0.01240932969) (9200,0.01213552011) (9300,0.01212683413) (9400,0.01212543733) (9500,0.01209418736) (9600,0.01213810463) (9700,0.01223628577) (9800,0.01214721156) (9900,0.01218897244) (10000,0.01222874895) (10100,0.01212110538) (10200,0.01213568402) (10300,0.01200262858) (10400,0.01205404075) (10500,0.01205080403) (10600,0.01192126079) (10700,0.01166889416) (10800,0.01156238746) (10900,0.01128613651) (11000,0.01134860767) (11100,0.01132912831) (11200,0.01126687974) (11300,0.01107126307) (11400,0.01094914079) (11500,0.01075345967) (11600,0.01068330389) (11700,0.01073241038) (11800,0.01064018402) (11900,0.01071338505) (12000,0.01044712253) (12100,0.01051259013) (12200,0.01058817683) (12300,0.01076515848) (12400,0.01083744066) (12500,0.01089505246) (12600,0.01093014358) (12700,0.01104983287) (12800,0.01123144245) (12900,0.01118493443) (13000,0.01137789544) (13100,0.01127673024) (13200,0.01120951949) (13300,0.0111577251) (13400,0.01094567347) (13500,0.01089491602) (13600,0.01086737663) (13700,0.01057158513) (13800,0.01055092737) (13900,0.01046513719) (14000,0.01035989868) (14100,0.01021244815) (14200,0.01016822178) (14300,0.01016790392) (14400,0.01022830214) (14500,0.01023026649) (14600,0.01005387604) (14700,0.0100795798) (14800,0.009965702798) (14900,0.01003842466) (15000,0.01002566349) (15100,0.01012525745) (15200,0.01010605907) (15300,0.01000122987) (15400,0.009986345284) (15500,0.009922301117) (15600,0.009956114739) (15700,0.01002767105) (15800,0.01006536847) (15900,0.01007202715) (16000,0.01001773383) (16100,0.0100610543) (16200,0.01000048881) (16300,0.01012517055) (16400,0.01011764975) (16500,0.009984679893) (16600,0.01020190259) (16700,0.01002641683) (16800,0.01004752647) (16900,0.01014905954) (17000,0.01032616105) (17100,0.01043435428) (17200,0.01050439021) (17300,0.01043868605) (17400,0.0103986837) (17500,0.01060989983) (17600,0.01030141292) (17700,0.01047037384) (17800,0.01050494751) (17900,0.01040490428) (18000,0.01012807842) (18100,0.01012288667) (18200,0.01019268762) (18300,0.01014393987) (18400,0.01020604866) (18500,0.01020847401) (18600,0.01021410506) (18700,0.01007096283) (18800,0.009960695356) (18900,0.009803247266) (19000,0.009923032485) (19100,0.00975725688) (19200,0.009784800466) (19300,0.009709510766) (19400,0.009770044312) (19500,0.009434186528) (19600,0.009516341565) (19700,0.009680379415) (19800,0.009756851429) (19900,0.009821097227) (20000,0.009996871511) (20100,0.009950744221) (20200,0.009890155168) (20300,0.009987786179) (20400,0.009869894804) (20500,0.01014236771) (20600,0.01017849818) (20700,0.009978630207) (20800,0.009883173183) (20900,0.01001670361) (21000,0.009823035076) (21100,0.009958119784) (21200,0.009885864332) (21300,0.00973681584) (21400,0.009704214055) (21500,0.009817834478) (21600,0.00966685228) (21700,0.009765696619) (21800,0.009735720605) (21900,0.009743107669) (22000,0.00964158345) (22100,0.009435022343) (22200,0.009397017956) (22300,0.009551656619) (22400,0.009591662139) (22500,0.009513351507) (22600,0.009626752138) (22700,0.009484059922) (22800,0.009688581061) (22900,0.009481568262) (23000,0.009627848957) (23100,0.009810659476) (23200,0.009734473843) (23300,0.009579527099) (23400,0.009638193529) (23500,0.009685395844) (23600,0.009688470699) (23700,0.009772216994) (23800,0.009522177186)};
\addplot[pubRose, opacity=0.18, line width=0.35pt] coordinates {(100,0.03052751347) (200,0.02947921678) (300,0.02871293947) (400,0.03032699134) (500,0.02957106195) (600,0.03000956743) (700,0.02993138161) (800,0.03088359279) (900,0.03116037386) (1000,0.03111223187) (1100,0.03086814079) (1200,0.03130962159) (1300,0.03190811668) (1400,0.03158943001) (1500,0.03189316224) (1600,0.03206932601) (1700,0.03229655344) (1800,0.03114986196) (1900,0.03088221401) (2000,0.03051848374) (2100,0.03044335451) (2200,0.0299495196) (2300,0.02964933421) (2400,0.02881026696) (2500,0.0284295423) (2600,0.02779325582) (2700,0.02732835095) (2800,0.02717834357) (2900,0.02652072515) (3000,0.0261884911) (3100,0.02581507899) (3200,0.02548066899) (3300,0.02469647508) (3400,0.02475323956) (3500,0.02437195797) (3600,0.0233563168) (3700,0.02279920988) (3800,0.02277315259) (3900,0.02256729677) (4000,0.02227264848) (4100,0.0219725024) (4200,0.0216205243) (4300,0.02130860705) (4400,0.02099709418) (4500,0.02083477862) (4600,0.02111486755) (4700,0.02116893642) (4800,0.02067475282) (4900,0.02048221789) (5000,0.02020366024) (5100,0.02005301733) (5200,0.01986712497) (5300,0.01995701212) (5400,0.02003982626) (5500,0.02004139312) (5600,0.01994611602) (5700,0.01978773978) (5800,0.01951797679) (5900,0.01953176968) (6000,0.0197208412) (6100,0.01966139022) (6200,0.01963786278) (6300,0.01963087413) (6400,0.01926013529) (6500,0.01924517639) (6600,0.01941149477) (6700,0.01931110416) (6800,0.01962355897) (6900,0.01933138389) (7000,0.01917823981) (7100,0.01897367164) (7200,0.01905690543) (7300,0.01886981167) (7400,0.01882689744) (7500,0.01882406492) (7600,0.01841011234) (7700,0.0180743495) (7800,0.0179954784) (7900,0.01811692044) (8000,0.0180270426) (8100,0.01791357584) (8200,0.01766531151) (8300,0.0174576873) (8400,0.01762712244) (8500,0.01738434173) (8600,0.01746543907) (8700,0.01763077639) (8800,0.01734183263) (8900,0.01731752455) (9000,0.01732978225) (9100,0.01745689996) (9200,0.01742412988) (9300,0.01739172917) (9400,0.01723910887) (9500,0.01725699417) (9600,0.01725235451) (9700,0.01725081261) (9800,0.01745388769) (9900,0.01713293921) (10000,0.01696373392) (10100,0.01704318151) (10200,0.01695425268) (10300,0.01684441306) (10400,0.01646808125) (10500,0.01608305098) (10600,0.01582927508) (10700,0.01575003127) (10800,0.01554955216) (10900,0.01570057804) (11000,0.01577681908) (11100,0.01585262278) (11200,0.01576986136) (11300,0.01600001408) (11400,0.01590302978) (11500,0.0160538285) (11600,0.01614327813) (11700,0.01590190316) (11800,0.01594929015) (11900,0.01578026619) (12000,0.01587006468) (12100,0.01547410022) (12200,0.01547852466) (12300,0.01540529253) (12400,0.0155597588) (12500,0.0153035624) (12600,0.01514132507) (12700,0.01523058945) (12800,0.01509104008) (12900,0.01508094948) (13000,0.01499909144) (13100,0.01490294319) (13200,0.01493547754) (13300,0.01492823148) (13400,0.01483428376) (13500,0.01515701385) (13600,0.01506710164) (13700,0.01494476022) (13800,0.01478247391) (13900,0.01482833484) (14000,0.01463249503) (14100,0.01472775964) (14200,0.01474403152) (14300,0.01464809636) (14400,0.01465669302) (14500,0.01432035062) (14600,0.01430258052) (14700,0.0143295154) (14800,0.01434223978) (14900,0.01407646639) (15000,0.01430385159) (15100,0.01440330027) (15200,0.01427768683) (15300,0.01407012222) (15400,0.0141491753) (15500,0.01427747598) (15600,0.014273958) (15700,0.01400041273) (15800,0.01398097808) (15900,0.01388396835) (16000,0.01360200839) (16100,0.01327330889) (16200,0.01334312316) (16300,0.01326470487) (16400,0.01316062175) (16500,0.01309841322) (16600,0.01291749915) (16700,0.01318777055) (16800,0.01335900491) (16900,0.01349451663) (17000,0.01354018236) (17100,0.01370298881) (17200,0.01349280495) (17300,0.01356049152) (17400,0.01360988095) (17500,0.0136478398) (17600,0.01376889655) (17700,0.0137370673) (17800,0.01368273189) (17900,0.01367398556) (18000,0.01353761349) (18100,0.01364124985) (18200,0.01370359445) (18300,0.0137706358) (18400,0.01360678328) (18500,0.0136917891) (18600,0.01353334971) (18700,0.01347136069) (18800,0.01353345122) (18900,0.01356011005) (19000,0.01350953542) (19100,0.01347536407) (19200,0.01361495946) (19300,0.01360960538) (19400,0.01373321703) (19500,0.01333131008) (19600,0.01368345534) (19700,0.01363221109) (19800,0.01341960356) (19900,0.01344610341) (20000,0.01358267469) (20100,0.01351901721) (20200,0.01340269363) (20300,0.01338844709) (20400,0.01340204561) (20500,0.01362630622) (20600,0.01358651323) (20700,0.01347557073) (20800,0.01341449236) (20900,0.01338197654) (21000,0.01310813027) (21100,0.01293311529) (21200,0.01284546526) (21300,0.0128641746) (21400,0.01294417772) (21500,0.01306379065) (21600,0.01324359048) (21700,0.01329238256) (21800,0.0132737726) (21900,0.01334091555) (22000,0.01344205271) (22100,0.01357375709) (22200,0.01371396296) (22300,0.01369754886) (22400,0.01351609137) (22500,0.01339320103) (22600,0.01300097974) (22700,0.01318982961) (22800,0.01322156535) (22900,0.01326043094) (23000,0.01348240664) (23100,0.01354095088) (23200,0.01333400058) (23300,0.01340493318) (23400,0.01333831875) (23500,0.01337153893) (23600,0.01345458534) (23700,0.0134055132) (23800,0.01337124836)};
\addplot[pubRose, opacity=0.18, line width=0.35pt] coordinates {(100,0.01686333492) (200,0.01728607528) (300,0.01655177089) (400,0.01598596992) (500,0.01542854607) (600,0.01528442139) (700,0.01522369111) (800,0.0151742145) (900,0.01529806335) (1000,0.01502045356) (1100,0.01500885878) (1200,0.01470338572) (1300,0.01463200226) (1400,0.01453119256) (1500,0.01458817702) (1600,0.01468303772) (1700,0.0148271489) (1800,0.0148246632) (1900,0.01517554941) (2000,0.01568036294) (2100,0.01551099634) (2200,0.0156488427) (2300,0.01564490506) (2400,0.01582510127) (2500,0.01596707944) (2600,0.01593222115) (2700,0.01591609363) (2800,0.01588275013) (2900,0.01542795086) (3000,0.01538832383) (3100,0.01526417388) (3200,0.01529288637) (3300,0.01537225982) (3400,0.01529346798) (3500,0.01540378397) (3600,0.01537642339) (3700,0.01532140411) (3800,0.01547283195) (3900,0.01551598217) (4000,0.01516500255) (4100,0.01538097421) (4200,0.01517562475) (4300,0.01503862385) (4400,0.01515033012) (4500,0.01488521537) (4600,0.01480180398) (4700,0.01479110429) (4800,0.01459593466) (4900,0.01438570758) (5000,0.01433366388) (5100,0.01401421856) (5200,0.01386697721) (5300,0.01381331645) (5400,0.01344331605) (5500,0.01369338529) (5600,0.01371717975) (5700,0.01345611829) (5800,0.01323526008) (5900,0.01345795868) (6000,0.01353008542) (6100,0.01381258778) (6200,0.01379642421) (6300,0.01405982086) (6400,0.01422296222) (6500,0.01396032833) (6600,0.01375988675) (6700,0.01361829387) (6800,0.01376411896) (6900,0.01332650408) (7000,0.01335008908) (7100,0.01305419132) (7200,0.01324931029) (7300,0.01310404073) (7400,0.01318107517) (7500,0.01306002205) (7600,0.01300164936) (7700,0.01320746038) (7800,0.01325801481) (7900,0.01340444311) (8000,0.01341502052) (8100,0.01339690331) (8200,0.01328053838) (8300,0.0131789309) (8400,0.01298781876) (8500,0.01287011979) (8600,0.01298850756) (8700,0.01277860533) (8800,0.01261745524) (8900,0.01256055478) (9000,0.01242511934) (9100,0.01239940925) (9200,0.01228989996) (9300,0.0122988455) (9400,0.01236294527) (9500,0.01262881039) (9600,0.01267449381) (9700,0.01267651841) (9800,0.01270898897) (9900,0.01279279422) (10000,0.0129690621) (10100,0.01294207061) (10200,0.01288656052) (10300,0.01274128556) (10400,0.01277792118) (10500,0.01258001272) (10600,0.01246508891) (10700,0.0131662895) (10800,0.01318776207) (10900,0.0131133683) (11000,0.0129401274) (11100,0.01289042411) (11200,0.01325827595) (11300,0.01325508393) (11400,0.01306008995) (11500,0.01329225637) (11600,0.01323374612) (11700,0.01265184218) (11800,0.01275761714) (11900,0.01279219929) (12000,0.01272789203) (12100,0.01294048242) (12200,0.01255431082) (12300,0.01272735558) (12400,0.01277079675) (12500,0.01273971014) (12600,0.01288314667) (12700,0.01296315026) (12800,0.01279512253) (12900,0.01278886702) (13000,0.01287729954) (13100,0.01279175142) (13200,0.01284849783) (13300,0.01292445725) (13400,0.01312157223) (13500,0.0132478103) (13600,0.01293876767) (13700,0.01296139788) (13800,0.01285890499) (13900,0.01284210598) (14000,0.01278935317) (14100,0.01290198937) (14200,0.01292062663) (14300,0.01274418384) (14400,0.01271776594) (14500,0.01246303339) (14600,0.01268535834) (14700,0.01255884646) (14800,0.01272063218) (14900,0.01274813451) (15000,0.0127536593) (15100,0.01258642003) (15200,0.01245992854) (15300,0.01235761549) (15400,0.01211553412) (15500,0.01201110659) (15600,0.01191368299) (15700,0.01214035153) (15800,0.01217482965) (15900,0.01214124151) (16000,0.01209092643) (16100,0.01214584168) (16200,0.01236647703) (16300,0.01242471756) (16400,0.01269968823) (16500,0.0126189582) (16600,0.01255099252) (16700,0.01226422824) (16800,0.01213752097) (16900,0.01218092861) (17000,0.01221050629) (17100,0.01227398356) (17200,0.01209487161) (17300,0.01224353565) (17400,0.01207594378) (17500,0.01227517314) (17600,0.01262149904) (17700,0.01263038982) (17800,0.01259105084) (17900,0.01256229496) (18000,0.01252463236) (18100,0.01250790227) (18200,0.01277470449) (18300,0.01276210649) (18400,0.01279606586) (18500,0.01270402828) (18600,0.01260957336) (18700,0.0124977109) (18800,0.012302587) (18900,0.01225845683) (19000,0.01211526347) (19100,0.01187090669) (19200,0.01165038515) (19300,0.01134219328) (19400,0.01152127776) (19500,0.01139151817) (19600,0.01125555756) (19700,0.01151258182) (19800,0.01165774325) (19900,0.01195557443) (20000,0.01192497006) (20100,0.0121553164) (20200,0.01207620194) (20300,0.01215044884) (20400,0.01187771326) (20500,0.01202178756) (20600,0.01200109059) (20700,0.01183596905) (20800,0.01193213779) (20900,0.01163492575) (21000,0.01176421959) (21100,0.01158714406) (21200,0.01155292895) (21300,0.01151464554) (21400,0.0114174325) (21500,0.01137962937) (21600,0.01153750392) (21700,0.01220440147) (21800,0.01232799692) (21900,0.0123612727) (22000,0.01245034663) (22100,0.01247686828) (22200,0.01262442088) (22300,0.01270256648) (22400,0.01276252726) (22500,0.01269560382) (22600,0.01242593443) (22700,0.01166355638) (22800,0.01172085321) (22900,0.01166788181) (23000,0.01162186144) (23100,0.01177763073) (23200,0.01160978451) (23300,0.01173294988) (23400,0.01171975443) (23500,0.01188568929) (23600,0.01209867755) (23700,0.01218222333) (23800,0.01193055985)};
\addplot[pubRose, opacity=0.18, line width=0.35pt] coordinates {(100,0.03116114438) (200,0.02795513347) (300,0.02548247327) (400,0.0242131562) (500,0.02311790846) (600,0.02222470598) (700,0.02176716046) (800,0.0214341078) (900,0.02128255367) (1000,0.02083552144) (1100,0.0196313018) (1200,0.01896895375) (1300,0.01864972021) (1400,0.01862158626) (1500,0.01902128533) (1600,0.0191284975) (1700,0.01900003254) (1800,0.01922541745) (1900,0.0193346519) (2000,0.01977333035) (2100,0.01991745867) (2200,0.02048231214) (2300,0.0209663678) (2400,0.02078672647) (2500,0.02080533449) (2600,0.021135002) (2700,0.02149057612) (2800,0.0216080578) (2900,0.02149909697) (3000,0.02159324363) (3100,0.02161182649) (3200,0.02128255311) (3300,0.02107839752) (3400,0.02115616743) (3500,0.02098480463) (3600,0.02084721774) (3700,0.02067279201) (3800,0.02061251327) (3900,0.02069759537) (4000,0.02040480394) (4100,0.02023759913) (4200,0.0202223463) (4300,0.02029830739) (4400,0.02041566595) (4500,0.02048014961) (4600,0.02028383203) (4700,0.02008477356) (4800,0.01997567751) (4900,0.01956567988) (5000,0.01943089869) (5100,0.01955114324) (5200,0.01936271451) (5300,0.01917487588) (5400,0.0190885108) (5500,0.01867281366) (5600,0.0186713228) (5700,0.01866490785) (5800,0.01844464447) (5900,0.01856488399) (6000,0.01857352629) (6100,0.01856815834) (6200,0.01841224898) (6300,0.01802097568) (6400,0.01789596444) (6500,0.01795572611) (6600,0.01760463165) (6700,0.01755539803) (6800,0.01731875511) (6900,0.01705398001) (7000,0.01702837572) (7100,0.01679099444) (7200,0.01679677535) (7300,0.01695356788) (7400,0.01703601303) (7500,0.01721351845) (7600,0.01724993465) (7700,0.01716356752) (7800,0.01721893223) (7900,0.0173677437) (8000,0.0171370348) (8100,0.01706771534) (8200,0.01688289968) (8300,0.01665753666) (8400,0.01646056864) (8500,0.01612693369) (8600,0.0161075104) (8700,0.01597064789) (8800,0.01585055646) (8900,0.01567339115) (9000,0.01579070054) (9100,0.01560840923) (9200,0.01570803905) (9300,0.01598333828) (9400,0.01571184788) (9500,0.01582202278) (9600,0.01561829718) (9700,0.01562183192) (9800,0.01551053366) (9900,0.01553132571) (10000,0.01540429657) (10100,0.01540086865) (10200,0.01548566818) (10300,0.01519452957) (10400,0.0151566986) (10500,0.01477985559) (10600,0.01522605736) (10700,0.01520943744) (10800,0.01533294851) (10900,0.0153119294) (11000,0.01548129106) (11100,0.01543511003) (11200,0.01544946991) (11300,0.01547519043) (11400,0.01529540243) (11500,0.0157049343) (11600,0.01542193331) (11700,0.01534429519) (11800,0.01558928723) (11900,0.01555717569) (12000,0.01539875343) (12100,0.01535328263) (12200,0.01524173627) (12300,0.01511912541) (12400,0.01533432202) (12500,0.01497065835) (12600,0.01488230238) (12700,0.01479441067) (12800,0.01458263109) (12900,0.01448087124) (13000,0.01448419532) (13100,0.01431199275) (13200,0.01398907397) (13300,0.01418404374) (13400,0.01397541352) (13500,0.01375297038) (13600,0.01386168487) (13700,0.01423145849) (13800,0.01413043691) (13900,0.01418179721) (14000,0.0139553057) (14100,0.01411723346) (14200,0.01430925634) (14300,0.01424923399) (14400,0.01442872882) (14500,0.01460402273) (14600,0.01436536238) (14700,0.01421391191) (14800,0.01434404757) (14900,0.014315998) (15000,0.01431847401) (15100,0.01441620635) (15200,0.01442487696) (15300,0.01432697074) (15400,0.01422533076) (15500,0.01413399521) (15600,0.01431831503) (15700,0.01421381179) (15800,0.01381223481) (15900,0.01381470859) (16000,0.01397161735) (16100,0.01403212994) (16200,0.01399851907) (16300,0.01426464394) (16400,0.01424068371) (16500,0.01442236537) (16600,0.01445250688) (16700,0.01445625098) (16800,0.01481359759) (16900,0.01491939276) (17000,0.01481706789) (17100,0.01460816925) (17200,0.01460482897) (17300,0.01433277288) (17400,0.01448806487) (17500,0.01461002808) (17600,0.01457589194) (17700,0.0146962055) (17800,0.01465582829) (17900,0.01444292953) (18000,0.01459899824) (18100,0.01472060597) (18200,0.01451432062) (18300,0.01470117169) (18400,0.0148614238) (18500,0.01464053402) (18600,0.0147031174) (18700,0.01466391366) (18800,0.01461305572) (18900,0.01468851315) (19000,0.01443942254) (19100,0.01440883977) (19200,0.01480072355) (19300,0.0145977756) (19400,0.01434491379) (19500,0.01463773074) (19600,0.01461047595) (19700,0.01445911331) (19800,0.01460260069) (19900,0.0146782374) (20000,0.01467982335) (20100,0.01455923868) (20200,0.01438870989) (20300,0.01444444815) (20400,0.01446342226) (20500,0.01415220816) (20600,0.01408294393) (20700,0.01412521061) (20800,0.01393077988) (20900,0.01388633205) (21000,0.01406174349) (21100,0.01418125313) (21200,0.01433865344) (21300,0.0144475597) (21400,0.01469398579) (21500,0.01473960551) (21600,0.01478457209) (21700,0.01497004135) (21800,0.01497551622) (21900,0.01507656025) (22000,0.01504595494) (22100,0.01491257632) (22200,0.0146073292) (22300,0.0144524063) (22400,0.01441576565) (22500,0.0144080882) (22600,0.01444241041) (22700,0.01416616449) (22800,0.01425021896) (22900,0.01405647928) (23000,0.01416096361) (23100,0.01423960915) (23200,0.01435345272) (23300,0.01416289415) (23400,0.01410786733) (23500,0.01398487492) (23600,0.01403505942) (23700,0.01432944415) (23800,0.01415927643)};
\addplot[pubRose, opacity=0.18, line width=0.35pt] coordinates {(100,0.0275069233) (200,0.02262121718) (300,0.02029133278) (400,0.01979279192) (500,0.01925409697) (600,0.01928938335) (700,0.01885733035) (800,0.01853230922) (900,0.01834670599) (1000,0.01876326539) (1100,0.01757589523) (1200,0.01748590898) (1300,0.01763251256) (1400,0.01742124856) (1500,0.01767376084) (1600,0.01749310959) (1700,0.01772511583) (1800,0.01797473468) (1900,0.01812077034) (2000,0.01780134868) (2100,0.01807520408) (2200,0.01813375149) (2300,0.01805584189) (2400,0.01810206939) (2500,0.01798301339) (2600,0.01810287628) (2700,0.01807334553) (2800,0.01796454433) (2900,0.01765400711) (3000,0.01731785778) (3100,0.01709493902) (3200,0.01700901967) (3300,0.01707577668) (3400,0.01736038309) (3500,0.01699114628) (3600,0.01668928508) (3700,0.0164423693) (3800,0.01601660214) (3900,0.015862629) (4000,0.01589016151) (4100,0.01562311798) (4200,0.01557722911) (4300,0.01527525438) (4400,0.01479187887) (4500,0.01491710637) (4600,0.01491674893) (4700,0.0147501165) (4800,0.01491777385) (4900,0.01505407235) (5000,0.01475302717) (5100,0.01490815962) (5200,0.01489386922) (5300,0.01474140529) (5400,0.01452923045) (5500,0.01432266338) (5600,0.01388692679) (5700,0.0138323904) (5800,0.01371296346) (5900,0.01356045799) (6000,0.01359333834) (6100,0.01336284457) (6200,0.01320396094) (6300,0.01325874142) (6400,0.01336620823) (6500,0.01324533308) (6600,0.01345343655) (6700,0.01345897214) (6800,0.01359584564) (6900,0.01351643438) (7000,0.01328589953) (7100,0.01341761043) (7200,0.0133101983) (7300,0.01333770175) (7400,0.01356927054) (7500,0.01346285259) (7600,0.0133242867) (7700,0.01324172067) (7800,0.01293685706) (7900,0.01298722858) (8000,0.01305376375) (8100,0.01278513009) (8200,0.01288644327) (8300,0.01281124027) (8400,0.01243969919) (8500,0.01256465381) (8600,0.01254822025) (8700,0.01254079873) (8800,0.01253823498) (8900,0.01240844717) (9000,0.01248291247) (9100,0.01248610597) (9200,0.01238578651) (9300,0.0123831138) (9400,0.01230827505) (9500,0.01242715642) (9600,0.01240790486) (9700,0.0123706758) (9800,0.01227489263) (9900,0.01250618463) (10000,0.0123978884) (10100,0.01245985888) (10200,0.01243808093) (10300,0.01244439781) (10400,0.01233756859) (10500,0.01211885875) (10600,0.01222527688) (10700,0.01203914909) (10800,0.01193307089) (10900,0.01176187554) (11000,0.01174051473) (11100,0.0116491532) (11200,0.01158198938) (11300,0.01155374739) (11400,0.0114657809) (11500,0.01151299821) (11600,0.01133062188) (11700,0.01139283776) (11800,0.01165287495) (11900,0.01161675546) (12000,0.01166096218) (12100,0.01166176312) (12200,0.01140485005) (12300,0.01121871294) (12400,0.01135409176) (12500,0.0111196178) (12600,0.01130265798) (12700,0.0113474505) (12800,0.01118018078) (12900,0.01104295915) (13000,0.01092906604) (13100,0.01102523282) (13200,0.0111532351) (13300,0.01122650616) (13400,0.01122885821) (13500,0.0113071911) (13600,0.01106890012) (13700,0.01107973903) (13800,0.01112459311) (13900,0.01129188808) (14000,0.01128346743) (14100,0.01132967779) (14200,0.01134779677) (14300,0.01134075541) (14400,0.01138385367) (14500,0.01141349841) (14600,0.01151966499) (14700,0.01135832211) (14800,0.01122406423) (14900,0.01131149102) (15000,0.0115267437) (15100,0.01128491834) (15200,0.01126281545) (15300,0.01113072876) (15400,0.01138377041) (15500,0.01138488594) (15600,0.01133408071) (15700,0.01157145714) (15800,0.01157298135) (15900,0.01145426808) (16000,0.01114963247) (16100,0.01109863082) (16200,0.01102640433) (16300,0.01114556603) (16400,0.01088860473) (16500,0.01085653752) (16600,0.01083167996) (16700,0.01063143881) (16800,0.01082089143) (16900,0.01072515007) (17000,0.01082758037) (17100,0.01100494638) (17200,0.01119217221) (17300,0.01117970748) (17400,0.01119049704) (17500,0.01119441502) (17600,0.01133244429) (17700,0.01132009812) (17800,0.01127254274) (17900,0.01131931162) (18000,0.01121223345) (18100,0.01110808076) (18200,0.0110230281) (18300,0.01107439455) (18400,0.01101437109) (18500,0.01107644923) (18600,0.01083900612) (18700,0.01072480399) (18800,0.01080176663) (18900,0.01070502307) (19000,0.01089939941) (19100,0.01090870611) (19200,0.01082752869) (19300,0.01080952426) (19400,0.01075474992) (19500,0.01080330862) (19600,0.01105703069) (19700,0.01128383391) (19800,0.01128321504) (19900,0.01137621114) (20000,0.01122707874) (20100,0.01153219119) (20200,0.01152433529) (20300,0.01150540691) (20400,0.01145743644) (20500,0.01129948674) (20600,0.01120354868) (20700,0.01102519222) (20800,0.01083979169) (20900,0.01093199765) (21000,0.01105552763) (21100,0.01078457776) (21200,0.01079567606) (21300,0.01105492543) (21400,0.0112033654) (21500,0.01119109057) (21600,0.01115936898) (21700,0.01125735343) (21800,0.01146872668) (21900,0.0114230915) (22000,0.01142095188) (22100,0.01168031944) (22200,0.01193592278) (22300,0.01180524789) (22400,0.01194245107) (22500,0.01223175535) (22600,0.01229671203) (22700,0.01234070752) (22800,0.01213548398) (22900,0.01203707801) (23000,0.01227227887) (23100,0.01219400605) (23200,0.01223317571) (23300,0.01230541319) (23400,0.01223037336) (23500,0.01197173009) (23600,0.01197267445) (23700,0.01201175852) (23800,0.01194525687)};
\addplot[pubRose, opacity=0.18, line width=0.35pt] coordinates {(100,0.01432010718) (200,0.0134911132) (300,0.01267833076) (400,0.01235429058) (500,0.01185942739) (600,0.01187112307) (700,0.01174626393) (800,0.01162186463) (900,0.01157568881) (1000,0.01141049964) (1100,0.01120254081) (1200,0.01096831989) (1300,0.01091882596) (1400,0.01090970971) (1500,0.01103787106) (1600,0.01088613141) (1700,0.01103670904) (1800,0.01116434615) (1900,0.01131170075) (2000,0.01159830159) (2100,0.01170317288) (2200,0.01182958139) (2300,0.01222638879) (2400,0.01235369183) (2500,0.01227582199) (2600,0.01253295923) (2700,0.01243214523) (2800,0.012401665) (2900,0.01225913279) (3000,0.01207295442) (3100,0.01200364036) (3200,0.01207335135) (3300,0.01183228176) (3400,0.01183658475) (3500,0.01229474414) (3600,0.012257009) (3700,0.01242457544) (3800,0.0125304372) (3900,0.01253042994) (4000,0.01276380327) (4100,0.01290052403) (4200,0.01295669517) (4300,0.01295286873) (4400,0.01313105701) (4500,0.01301041991) (4600,0.01307640607) (4700,0.01291535264) (4800,0.01294448879) (4900,0.01298282277) (5000,0.01298543261) (5100,0.01271417774) (5200,0.01288135555) (5300,0.01298474949) (5400,0.01288519157) (5500,0.01292430256) (5600,0.01278210273) (5700,0.01292225132) (5800,0.0129221064) (5900,0.01310715973) (6000,0.01303475518) (6100,0.01322503109) (6200,0.01324763587) (6300,0.01323122429) (6400,0.01318286061) (6500,0.01320339777) (6600,0.01337841405) (6700,0.01330250315) (6800,0.0131862307) (6900,0.01309996406) (7000,0.01316561913) (7100,0.01327243093) (7200,0.01331077181) (7300,0.0133102702) (7400,0.01343711661) (7500,0.01330615235) (7600,0.01323164692) (7700,0.01323714759) (7800,0.01332935002) (7900,0.01333095692) (8000,0.01323766774) (8100,0.01311745923) (8200,0.01288938317) (8300,0.01289708652) (8400,0.01263348656) (8500,0.01273226039) (8600,0.01265987419) (8700,0.01250723731) (8800,0.01238796441) (8900,0.01268007057) (9000,0.01253122222) (9100,0.01239791801) (9200,0.01232496286) (9300,0.01233713534) (9400,0.01223668512) (9500,0.0121927334) (9600,0.01220322866) (9700,0.01225403622) (9800,0.01229182454) (9900,0.01193811679) (10000,0.01201365674) (10100,0.01218074868) (10200,0.01213523531) (10300,0.01204054486) (10400,0.01213237736) (10500,0.01194578214) (10600,0.01199500412) (10700,0.01231236663) (10800,0.01232987065) (10900,0.01241044989) (11000,0.01242704568) (11100,0.01229522871) (11200,0.01245170543) (11300,0.01236928022) (11400,0.01232518358) (11500,0.01231900239) (11600,0.01231324188) (11700,0.01210269909) (11800,0.01202132981) (11900,0.01214821115) (12000,0.01228011502) (12100,0.01209091777) (12200,0.01192609742) (12300,0.01193592101) (12400,0.01190428808) (12500,0.01188988183) (12600,0.01164061259) (12700,0.01162548214) (12800,0.01154907383) (12900,0.01140831178) (13000,0.0112386331) (13100,0.01129249297) (13200,0.01126254918) (13300,0.01133429129) (13400,0.01138846548) (13500,0.01137660854) (13600,0.01143515464) (13700,0.0113587467) (13800,0.01160533531) (13900,0.01145720929) (14000,0.01138190217) (14100,0.01153073553) (14200,0.01146657728) (14300,0.01147365216) (14400,0.01165599367) (14500,0.01166017195) (14600,0.01167092016) (14700,0.01196481539) (14800,0.01192466319) (14900,0.01176726399) (15000,0.01189662041) (15100,0.0118458367) (15200,0.01179471398) (15300,0.0117580465) (15400,0.01158652427) (15500,0.01139347767) (15600,0.01146884989) (15700,0.01121669142) (15800,0.01114634024) (15900,0.01138261706) (16000,0.01158297723) (16100,0.0114658094) (16200,0.01159196729) (16300,0.0114852475) (16400,0.01138339359) (16500,0.01172558684) (16600,0.0117019888) (16700,0.01165890926) (16800,0.01152680842) (16900,0.01147375703) (17000,0.01124873608) (17100,0.01149013676) (17200,0.01168596316) (17300,0.01171211284) (17400,0.01161357164) (17500,0.0113457608) (17600,0.01133959852) (17700,0.01132145748) (17800,0.01131431907) (17900,0.01143470658) (18000,0.01126396665) (18100,0.01104874546) (18200,0.01103757462) (18300,0.01104749301) (18400,0.01120048566) (18500,0.01140507367) (18600,0.01145323068) (18700,0.01159521732) (18800,0.0115429854) (18900,0.01144335689) (19000,0.01196280364) (19100,0.01189779481) (19200,0.01162483534) (19300,0.01162674483) (19400,0.01175426766) (19500,0.01164134163) (19600,0.01154777156) (19700,0.01137461727) (19800,0.01139075877) (19900,0.01123677399) (20000,0.01086358046) (20100,0.01089022998) (20200,0.01078017857) (20300,0.01080071712) (20400,0.0106298686) (20500,0.01061683837) (20600,0.01059751641) (20700,0.01067676209) (20800,0.01058354136) (20900,0.01063204203) (21000,0.01047969926) (21100,0.01050634505) (21200,0.01057317797) (21300,0.01050249459) (21400,0.01046542404) (21500,0.01061176639) (21600,0.01066043181) (21700,0.01062446767) (21800,0.01083985409) (21900,0.01078887992) (22000,0.0108246197) (22100,0.01076358333) (22200,0.01080137966) (22300,0.01071558092) (22400,0.010870976) (22500,0.01079835352) (22600,0.01058659833) (22700,0.01057870165) (22800,0.01055807918) (22900,0.01072425339) (23000,0.01078657992) (23100,0.01083992617) (23200,0.01094971877) (23300,0.01106652096) (23400,0.01143064788) (23500,0.01137729259) (23600,0.01141353911) (23700,0.01134888632) (23800,0.01120281881)};
\addplot[pubRose, opacity=0.18, line width=0.35pt] coordinates {(100,0.0236206837) (200,0.02023031749) (300,0.01817245005) (400,0.01699605747) (500,0.01616778858) (600,0.01579545314) (700,0.0152863932) (800,0.01495230687) (900,0.01481593328) (1000,0.01480558505) (1100,0.01393369809) (1200,0.01359827984) (1300,0.01366571616) (1400,0.01364113679) (1500,0.01381715462) (1600,0.01392328162) (1700,0.01400580611) (1800,0.01440900136) (1900,0.01459572669) (2000,0.0148590479) (2100,0.01490199436) (2200,0.01514278576) (2300,0.01528810291) (2400,0.01561035253) (2500,0.01577164903) (2600,0.01594459452) (2700,0.01631886754) (2800,0.01623258665) (2900,0.01619583666) (3000,0.01589165134) (3100,0.0157813536) (3200,0.01556869606) (3300,0.01518599931) (3400,0.01476753363) (3500,0.01476145172) (3600,0.01463319035) (3700,0.01435851743) (3800,0.01421493171) (3900,0.01423234846) (4000,0.01438185191) (4100,0.01454765471) (4200,0.01461692639) (4300,0.01489379928) (4400,0.0154720529) (4500,0.01539709959) (4600,0.01525648898) (4700,0.01525845574) (4800,0.01531198453) (4900,0.01523427665) (5000,0.01541565079) (5100,0.01557206567) (5200,0.01545942165) (5300,0.0153408234) (5400,0.01508902805) (5500,0.01526824143) (5600,0.01524374494) (5700,0.01530929934) (5800,0.01561965188) (5900,0.01544904457) (6000,0.01508951606) (6100,0.01500658235) (6200,0.01515931943) (6300,0.01519913794) (6400,0.01498516658) (6500,0.01495454265) (6600,0.01497840472) (6700,0.01522275982) (6800,0.01503817523) (6900,0.01526301643) (7000,0.0151447271) (7100,0.01465663398) (7200,0.01455695909) (7300,0.01440817071) (7400,0.01472335858) (7500,0.01434636628) (7600,0.01453689877) (7700,0.01421652399) (7800,0.0142193282) (7900,0.01400033385) (8000,0.01453543156) (8100,0.0149762081) (8200,0.01512101749) (8300,0.01539746393) (8400,0.01532861963) (8500,0.01545662312) (8600,0.01542668948) (8700,0.01557580819) (8800,0.01531993477) (8900,0.01559800776) (9000,0.01528098183) (9100,0.01481500883) (9200,0.01488247048) (9300,0.01490831301) (9400,0.01474494198) (9500,0.015226404) (9600,0.01497335136) (9700,0.01478364747) (9800,0.01495828778) (9900,0.01484189723) (10000,0.01476502139) (10100,0.01509179948) (10200,0.01498673446) (10300,0.01488361843) (10400,0.01485679811) (10500,0.01422651904) (10600,0.01436281689) (10700,0.01442327369) (10800,0.0143095376) (10900,0.01422408298) (11000,0.01404022872) (11100,0.01393947974) (11200,0.01373036569) (11300,0.01340315351) (11400,0.01348608015) (11500,0.01363583822) (11600,0.01366134984) (11700,0.01362202875) (11800,0.01367026679) (11900,0.01365608033) (12000,0.01377542606) (12100,0.01382896202) (12200,0.01417525532) (12300,0.01444431888) (12400,0.0144909896) (12500,0.01445736131) (12600,0.01425591931) (12700,0.01417562719) (12800,0.0140293302) (12900,0.01399936117) (13000,0.01416791938) (13100,0.0140340846) (13200,0.01370601654) (13300,0.01371386219) (13400,0.0134068626) (13500,0.0134588887) (13600,0.01357422052) (13700,0.0136431098) (13800,0.01354424972) (13900,0.0136033291) (14000,0.01341013378) (14100,0.01355370739) (14200,0.01356532173) (14300,0.01362482002) (14400,0.0139225916) (14500,0.01355071822) (14600,0.01359555926) (14700,0.01364673544) (14800,0.01367361369) (14900,0.01356146736) (15000,0.01341725765) (15100,0.01342163496) (15200,0.01349116638) (15300,0.01336427964) (15400,0.01341590676) (15500,0.01377114626) (15600,0.01374514848) (15700,0.01363148596) (15800,0.01372708557) (15900,0.01370835928) (16000,0.01387971872) (16100,0.01360066347) (16200,0.01358333686) (16300,0.01341556376) (16400,0.01303487457) (16500,0.01294461759) (16600,0.0127055211) (16700,0.01260529561) (16800,0.0125169755) (16900,0.01261111349) (17000,0.01269192277) (17100,0.01300388491) (17200,0.01303152293) (17300,0.01317265658) (17400,0.01333795078) (17500,0.01328095309) (17600,0.0134161517) (17700,0.01370941801) (17800,0.01392777869) (17900,0.0139570469) (18000,0.01390752029) (18100,0.0136828254) (18200,0.01380714569) (18300,0.0139010299) (18400,0.01361386264) (18500,0.01374044502) (18600,0.01392263621) (18700,0.01361992303) (18800,0.01333518364) (18900,0.01325319521) (19000,0.01337466165) (19100,0.01354432926) (19200,0.01324237371) (19300,0.01293719783) (19400,0.01329217954) (19500,0.01321492968) (19600,0.01304956079) (19700,0.01322077215) (19800,0.01316519529) (19900,0.01326048933) (20000,0.01294920975) (20100,0.01284142192) (20200,0.01317316452) (20300,0.01336056711) (20400,0.01313410988) (20500,0.01328771198) (20600,0.01347903404) (20700,0.01333942562) (20800,0.01361718532) (20900,0.01349415695) (21000,0.01375188911) (21100,0.01398679232) (21200,0.0136790378) (21300,0.01360980626) (21400,0.01374317007) (21500,0.01355182705) (21600,0.01348570045) (21700,0.01349871606) (21800,0.0134391062) (21900,0.01353401775) (22000,0.01327573117) (22100,0.01307160668) (22200,0.01308167987) (22300,0.01312044822) (22400,0.01307989946) (22500,0.01294442164) (22600,0.01284772782) (22700,0.01285514385) (22800,0.01275243163) (22900,0.01281856354) (23000,0.01289855903) (23100,0.0129627334) (23200,0.01301666796) (23300,0.01332500782) (23400,0.0132354998) (23500,0.01339025097) (23600,0.01339048846) (23700,0.01346096806) (23800,0.01365483142)};
\addplot[pubRose, opacity=0.18, line width=0.35pt] coordinates {(100,0.03922845423) (200,0.03334250487) (300,0.03228376309) (400,0.03042530874) (500,0.02908570878) (600,0.02822237431) (700,0.02707553469) (800,0.02660544473) (900,0.02630257834) (1000,0.02574861143) (1100,0.02385594621) (1200,0.02302028239) (1300,0.02218809538) (1400,0.02195453923) (1500,0.02169674095) (1600,0.02148733307) (1700,0.02174480986) (1800,0.02170811761) (1900,0.02156534437) (2000,0.02192078512) (2100,0.02204917185) (2200,0.02222760189) (2300,0.02230434138) (2400,0.02228388954) (2500,0.02194753792) (2600,0.02174937408) (2700,0.02155237962) (2800,0.02085359246) (2900,0.02060321122) (3000,0.01985785309) (3100,0.01945459563) (3200,0.01928986777) (3300,0.01850897744) (3400,0.01785754208) (3500,0.01775232572) (3600,0.01739961114) (3700,0.01695032604) (3800,0.01704553682) (3900,0.01665026154) (4000,0.01668698601) (4100,0.01669491567) (4200,0.01639279947) (4300,0.01637301371) (4400,0.01638489384) (4500,0.01613388909) (4600,0.01597808152) (4700,0.01591664515) (4800,0.01568180257) (4900,0.01563537112) (5000,0.01532494836) (5100,0.01492417548) (5200,0.01498906277) (5300,0.01501709558) (5400,0.01473278031) (5500,0.01493530991) (5600,0.01497752666) (5700,0.0149600802) (5800,0.01491726739) (5900,0.01472537816) (6000,0.0146443788) (6100,0.01474352498) (6200,0.01455439478) (6300,0.01468141787) (6400,0.01527433842) (6500,0.01493936591) (6600,0.01487461114) (6700,0.01476855548) (6800,0.01488938853) (6900,0.01505042473) (7000,0.01508946456) (7100,0.01502464991) (7200,0.01498073023) (7300,0.0150314251) (7400,0.01454872275) (7500,0.01477299733) (7600,0.01460655006) (7700,0.01457190821) (7800,0.01454726625) (7900,0.01438597068) (8000,0.01446980098) (8100,0.01444645794) (8200,0.01433876241) (8300,0.01402113065) (8400,0.01404973436) (8500,0.01401866712) (8600,0.01401513582) (8700,0.0137987501) (8800,0.01364943609) (8900,0.01370561169) (9000,0.01364844553) (9100,0.01349771693) (9200,0.01340097366) (9300,0.01323825866) (9400,0.01317880768) (9500,0.01305913981) (9600,0.01311865645) (9700,0.01323903315) (9800,0.01333709331) (9900,0.01313306019) (10000,0.0129566459) (10100,0.01308621503) (10200,0.01312126685) (10300,0.01321176169) (10400,0.01292306287) (10500,0.01284980923) (10600,0.01262633307) (10700,0.0126483391) (10800,0.01229209369) (10900,0.01225287309) (11000,0.01229735566) (11100,0.0122986313) (11200,0.01218836587) (11300,0.01202416671) (11400,0.01216101684) (11500,0.01208199514) (11600,0.0119615579) (11700,0.01181486873) (11800,0.01206926452) (11900,0.01221272601) (12000,0.01211789409) (12100,0.01192518808) (12200,0.01184076667) (12300,0.0118506724) (12400,0.0117426266) (12500,0.01172466306) (12600,0.0117270112) (12700,0.01161618382) (12800,0.01133813038) (12900,0.01119264467) (13000,0.01115795122) (13100,0.01113924524) (13200,0.01114155576) (13300,0.01118736565) (13400,0.01136837015) (13500,0.01138728531) (13600,0.01187272901) (13700,0.01202571737) (13800,0.01225466104) (13900,0.01231267322) (14000,0.01221840018) (14100,0.01211936427) (14200,0.01206431398) (14300,0.0121217926) (14400,0.01179148043) (14500,0.01173198633) (14600,0.01122680753) (14700,0.01117339442) (14800,0.01099691074) (14900,0.01084604524) (15000,0.01104984479) (15100,0.01119682342) (15200,0.01124707749) (15300,0.01098872172) (15400,0.01124454429) (15500,0.01123342877) (15600,0.01134064132) (15700,0.01126473397) (15800,0.01133913752) (15900,0.01141626723) (16000,0.01134006046) (16100,0.01124340463) (16200,0.01126335403) (16300,0.01138502657) (16400,0.01113346098) (16500,0.01119652633) (16600,0.01140105464) (16700,0.01142464904) (16800,0.01131618097) (16900,0.0112601636) (17000,0.01102571674) (17100,0.01105152247) (17200,0.01112275021) (17300,0.01088605756) (17400,0.01078289766) (17500,0.01064489847) (17600,0.01037470885) (17700,0.01033943361) (17800,0.0103191114) (17900,0.01038002009) (18000,0.0105013499) (18100,0.01055195816) (18200,0.01047294568) (18300,0.01080429601) (18400,0.01107700802) (18500,0.01121486248) (18600,0.01132460944) (18700,0.01140641375) (18800,0.01139680054) (18900,0.0113344755) (19000,0.01142476434) (19100,0.01122244587) (19200,0.01119446931) (19300,0.01117366077) (19400,0.01099634692) (19500,0.01077124579) (19600,0.01058909893) (19700,0.01046331795) (19800,0.0103306327) (19900,0.01032278854) (20000,0.01028412944) (20100,0.01049330672) (20200,0.01043438921) (20300,0.01015537595) (20400,0.01018613875) (20500,0.01022493988) (20600,0.01019768184) (20700,0.01022839313) (20800,0.0102998225) (20900,0.01004203493) (21000,0.01000719555) (21100,0.009914204106) (21200,0.009925281256) (21300,0.01016931813) (21400,0.01016958989) (21500,0.01012203842) (21600,0.01030932032) (21700,0.0102861356) (21800,0.01017838558) (21900,0.0104604804) (22000,0.010367392) (22100,0.01029139617) (22200,0.01003608378) (22300,0.009919993114) (22400,0.009988959506) (22500,0.009875871148) (22600,0.009648831934) (22700,0.009513601847) (22800,0.009699818306) (22900,0.009623670392) (23000,0.009473060537) (23100,0.009260467533) (23200,0.009506559186) (23300,0.009518710524) (23400,0.009420638345) (23500,0.009573032521) (23600,0.009809613507) (23700,0.0099465671) (23800,0.009688033629)};
\addplot[pubRose, opacity=0.18, line width=0.35pt] coordinates {(100,0.02615479007) (200,0.02672945894) (300,0.02573071793) (400,0.0245386106) (500,0.02290231735) (600,0.02171361881) (700,0.0215245637) (800,0.02126315795) (900,0.0210477983) (1000,0.02098757736) (1100,0.02000713516) (1200,0.01922711153) (1300,0.01875636037) (1400,0.01848396659) (1500,0.01858488824) (1600,0.0187951196) (1700,0.01855315808) (1800,0.01841057613) (1900,0.01809014026) (2000,0.01756125623) (2100,0.01773347585) (2200,0.01751877749) (2300,0.01746799601) (2400,0.01725402819) (2500,0.01737111183) (2600,0.01734721931) (2700,0.0172669542) (2800,0.01705606738) (2900,0.01720355945) (3000,0.0173409408) (3100,0.01700413395) (3200,0.01687391419) (3300,0.01661182977) (3400,0.01672110278) (3500,0.01669360008) (3600,0.01656030491) (3700,0.01625924818) (3800,0.01622461732) (3900,0.01613871679) (4000,0.01602095962) (4100,0.01632294282) (4200,0.01631623562) (4300,0.01606227085) (4400,0.01585696843) (4500,0.0156839272) (4600,0.01555654248) (4700,0.01556912968) (4800,0.01567962384) (4900,0.01550906654) (5000,0.01557800733) (5100,0.01525539933) (5200,0.01495500477) (5300,0.01502012378) (5400,0.01500364672) (5500,0.01463037208) (5600,0.01450450253) (5700,0.0144287223) (5800,0.01417003814) (5900,0.01409677295) (6000,0.01369851073) (6100,0.01353547815) (6200,0.01369306911) (6300,0.01368415169) (6400,0.0135140839) (6500,0.01367356721) (6600,0.01365628242) (6700,0.01352655124) (6800,0.01354443654) (6900,0.01345810518) (7000,0.01356140506) (7100,0.01359515721) (7200,0.01340664988) (7300,0.01345765768) (7400,0.01332251197) (7500,0.01323593585) (7600,0.01303227264) (7700,0.01312016584) (7800,0.01286568493) (7900,0.01277259747) (8000,0.01270084027) (8100,0.01265526647) (8200,0.0125084877) (8300,0.01232050993) (8400,0.01238305019) (8500,0.01230317671) (8600,0.01218149215) (8700,0.01211379524) (8800,0.0122004251) (8900,0.01208163667) (9000,0.01203598082) (9100,0.01200501528) (9200,0.01234249277) (9300,0.01219914202) (9400,0.01205578241) (9500,0.01198535739) (9600,0.01217744732) (9700,0.01201870954) (9800,0.01203533793) (9900,0.01212646952) (10000,0.01208871026) (10100,0.01203532191) (10200,0.01176985493) (10300,0.01185819823) (10400,0.01187851587) (10500,0.01188591765) (10600,0.01177993426) (10700,0.01193424249) (10800,0.01174436957) (10900,0.01167291012) (11000,0.01167994989) (11100,0.0115537216) (11200,0.01144859223) (11300,0.01138673378) (11400,0.01134277573) (11500,0.01140715051) (11600,0.01154317688) (11700,0.01133750994) (11800,0.01154676955) (11900,0.01149245566) (12000,0.01150258761) (12100,0.01138309697) (12200,0.01134981373) (12300,0.01120251603) (12400,0.01124662142) (12500,0.01101916768) (12600,0.01084614517) (12700,0.01090426911) (12800,0.01061640494) (12900,0.01059044134) (13000,0.01051625712) (13100,0.01066154353) (13200,0.01073998101) (13300,0.0108655923) (13400,0.01068244465) (13500,0.01078074053) (13600,0.01084647132) (13700,0.01090786643) (13800,0.01109162532) (13900,0.0109924336) (14000,0.01104039233) (14100,0.01096697524) (14200,0.01097894181) (14300,0.01092489082) (14400,0.0111114393) (14500,0.01096728332) (14600,0.01089101341) (14700,0.0108625426) (14800,0.01073338827) (14900,0.0110338808) (15000,0.01118563516) (15100,0.01121721566) (15200,0.01115875393) (15300,0.01106001437) (15400,0.01102290098) (15500,0.01107040495) (15600,0.01104545332) (15700,0.01100165341) (15800,0.01105572507) (15900,0.01094696587) (16000,0.01076905131) (16100,0.01073189788) (16200,0.01067887209) (16300,0.01065702569) (16400,0.01079985602) (16500,0.01086712535) (16600,0.01079493351) (16700,0.01084844889) (16800,0.01076071886) (16900,0.01070899311) (17000,0.01059948616) (17100,0.01071635988) (17200,0.0106714502) (17300,0.01086727092) (17400,0.010712195) (17500,0.0105600113) (17600,0.01074826838) (17700,0.01067814948) (17800,0.01081457259) (17900,0.01067572357) (18000,0.01083212923) (18100,0.01079771332) (18200,0.01080598319) (18300,0.01088716863) (18400,0.01081555001) (18500,0.01091866773) (18600,0.0106905289) (18700,0.01064539896) (18800,0.01064628167) (18900,0.01073172791) (19000,0.01077648876) (19100,0.01063130153) (19200,0.01070821676) (19300,0.01063011857) (19400,0.01065474059) (19500,0.01070985012) (19600,0.01082337666) (19700,0.01085377242) (19800,0.0108756043) (19900,0.01096913274) (20000,0.01099391729) (20100,0.01120743593) (20200,0.0113440874) (20300,0.01127090389) (20400,0.01120851804) (20500,0.01114392038) (20600,0.01111621624) (20700,0.01118265511) (20800,0.01112542106) (20900,0.01119820103) (21000,0.01102462616) (21100,0.01088423282) (21200,0.01075885296) (21300,0.010687037) (21400,0.01083421754) (21500,0.01082158554) (21600,0.01088541672) (21700,0.01077056536) (21800,0.01059585148) (21900,0.01051165275) (22000,0.0104965318) (22100,0.01068525966) (22200,0.01092320532) (22300,0.0109762311) (22400,0.01098080091) (22500,0.0110263302) (22600,0.01089264574) (22700,0.01096286057) (22800,0.01110436479) (22900,0.01095332094) (23000,0.01104747392) (23100,0.01115905624) (23200,0.01097183749) (23300,0.01101048682) (23400,0.01092415052) (23500,0.01085334411) (23600,0.01077329786) (23700,0.0107103534) (23800,0.01071729958)};
\addplot[pubRose, opacity=0.18, line width=0.35pt] coordinates {(100,0.02065422758) (200,0.0183063969) (300,0.01717310989) (400,0.01642560749) (500,0.01575208753) (600,0.01500151539) (700,0.01436948949) (800,0.01410057524) (900,0.01401348867) (1000,0.01401421055) (1100,0.01325960625) (1200,0.01298168693) (1300,0.01277074795) (1400,0.01275055911) (1500,0.01267764857) (1600,0.01277781138) (1700,0.0130634021) (1800,0.01325097661) (1900,0.01325310254) (2000,0.01317072744) (2100,0.01309090285) (2200,0.01301916568) (2300,0.01293065371) (2400,0.0128656921) (2500,0.01292443974) (2600,0.01313215578) (2700,0.01299285172) (2800,0.01280042157) (2900,0.01281508319) (3000,0.01277308483) (3100,0.01279599853) (3200,0.01272225687) (3300,0.01278605498) (3400,0.01273874464) (3500,0.01259913454) (3600,0.01245039701) (3700,0.01238239119) (3800,0.01231132103) (3900,0.01219349699) (4000,0.01198156169) (4100,0.01192591926) (4200,0.0120382865) (4300,0.012214451) (4400,0.0122538913) (4500,0.01234224066) (4600,0.01251774114) (4700,0.01269705854) (4800,0.012915642) (4900,0.01299782917) (5000,0.0132794884) (5100,0.01334589357) (5200,0.01334020831) (5300,0.01305343276) (5400,0.01296239449) (5500,0.01301561166) (5600,0.01268937122) (5700,0.01273220396) (5800,0.01262825998) (5900,0.01252536057) (6000,0.01241915813) (6100,0.01235565729) (6200,0.01220605224) (6300,0.01230304213) (6400,0.0122936165) (6500,0.01215109667) (6600,0.01233822582) (6700,0.01198992673) (6800,0.01191983679) (6900,0.01197732259) (7000,0.01175425891) (7100,0.01165487925) (7200,0.01183780143) (7300,0.01197440522) (7400,0.01211302569) (7500,0.01214220738) (7600,0.01190403104) (7700,0.01197594358) (7800,0.01196365785) (7900,0.01180996299) (8000,0.01204522317) (8100,0.01203289703) (8200,0.01184463277) (8300,0.01165571418) (8400,0.01157646524) (8500,0.01156677995) (8600,0.01173796924) (8700,0.01189824473) (8800,0.01187431021) (8900,0.01201803964) (9000,0.01195467953) (9100,0.01211442621) (9200,0.01202905858) (9300,0.0119818463) (9400,0.01194471968) (9500,0.01200506818) (9600,0.01203912413) (9700,0.01197611047) (9800,0.01213896358) (9900,0.01210578009) (10000,0.01209062692) (10100,0.01199347563) (10200,0.01210167129) (10300,0.01222801087) (10400,0.01211007508) (10500,0.01203836156) (10600,0.01189315282) (10700,0.01190042198) (10800,0.01158737037) (10900,0.01158811012) (11000,0.01150896959) (11100,0.01144913351) (11200,0.01152287675) (11300,0.01138482802) (11400,0.01149562038) (11500,0.01144341221) (11600,0.01143534565) (11700,0.01138707548) (11800,0.01175523959) (11900,0.01169166472) (12000,0.01179202031) (12100,0.01177165732) (12200,0.01184346629) (12300,0.01193445595) (12400,0.01174995722) (12500,0.01178763565) (12600,0.01172292642) (12700,0.01164513752) (12800,0.01184953805) (12900,0.01193012111) (13000,0.01187336938) (13100,0.01211925969) (13200,0.01189491274) (13300,0.01186718028) (13400,0.01205452932) (13500,0.01191369891) (13600,0.01194695011) (13700,0.01207275456) (13800,0.01155968281) (13900,0.01144345151) (14000,0.01158685349) (14100,0.01141374307) (14200,0.01141213039) (14300,0.01118827518) (14400,0.01119094919) (14500,0.01130564185) (14600,0.01128864167) (14700,0.01121505918) (14800,0.01123213852) (14900,0.01122222859) (15000,0.01112219961) (15100,0.01101802159) (15200,0.01099068802) (15300,0.01120960377) (15400,0.01116185244) (15500,0.01134176403) (15600,0.0116151874) (15700,0.01158534335) (15800,0.0115789108) (15900,0.01172024673) (16000,0.01163243931) (16100,0.01153430426) (16200,0.01157232169) (16300,0.01146001825) (16400,0.01134328004) (16500,0.01116368286) (16600,0.0110381119) (16700,0.01105149342) (16800,0.01113339029) (16900,0.01105384296) (17000,0.01098901853) (17100,0.01127593052) (17200,0.01137407674) (17300,0.01134826206) (17400,0.01151319165) (17500,0.01147478651) (17600,0.01136526177) (17700,0.01126529248) (17800,0.01117613064) (17900,0.01120956773) (18000,0.01139058182) (18100,0.01133390181) (18200,0.01126350183) (18300,0.0115901595) (18400,0.01157075893) (18500,0.0117946188) (18600,0.01211139392) (18700,0.01230032416) (18800,0.01238827771) (18900,0.01218998069) (19000,0.01228020871) (19100,0.01220679646) (19200,0.01234029504) (19300,0.01218999997) (19400,0.01206946308) (19500,0.01201197878) (19600,0.01177020343) (19700,0.01172020994) (19800,0.01178277293) (19900,0.01207887065) (20000,0.0119417429) (20100,0.01207100758) (20200,0.01202583117) (20300,0.01195316231) (20400,0.01210984811) (20500,0.01216686768) (20600,0.01216019457) (20700,0.01226149341) (20800,0.01222255165) (20900,0.01219696486) (21000,0.01221193569) (21100,0.0123219246) (21200,0.01224830234) (21300,0.01237340868) (21400,0.01240320029) (21500,0.01243878398) (21600,0.01240939759) (21700,0.01237955214) (21800,0.01229252703) (21900,0.01203243937) (22000,0.0119426026) (22100,0.01178735476) (22200,0.0117196396) (22300,0.01143540749) (22400,0.01149362586) (22500,0.0112498302) (22600,0.01109902877) (22700,0.01096112374) (22800,0.01091064895) (22900,0.01118553365) (23000,0.01134256348) (23100,0.01141604614) (23200,0.01152758887) (23300,0.01160272192) (23400,0.01145589836) (23500,0.01134327352) (23600,0.01148577631) (23700,0.01169995721) (23800,0.01185991718)};
\addplot[pubRose, opacity=0.18, line width=0.35pt] coordinates {(100,0.03022029623) (200,0.02921857312) (300,0.02777812754) (400,0.02661279822) (500,0.02614533864) (600,0.02548211844) (700,0.02488792794) (800,0.02403876488) (900,0.02347011512) (1000,0.02322972734) (1100,0.02221774161) (1200,0.02190561742) (1300,0.02139748391) (1400,0.02131130733) (1500,0.02084371503) (1600,0.02072044704) (1700,0.02084361035) (1800,0.02109124251) (1900,0.02124962695) (2000,0.021126673) (2100,0.02124382854) (2200,0.02103613317) (2300,0.02090816814) (2400,0.02072155625) (2500,0.02067977954) (2600,0.02042154483) (2700,0.0203649452) (2800,0.02035631593) (2900,0.02030526586) (3000,0.02061134055) (3100,0.02028048467) (3200,0.02001338471) (3300,0.02030117381) (3400,0.02011795305) (3500,0.02014679015) (3600,0.02046751324) (3700,0.02022484262) (3800,0.0198657671) (3900,0.01959745102) (4000,0.01909441855) (4100,0.01914229635) (4200,0.01887050904) (4300,0.01866655834) (4400,0.01857436504) (4500,0.01835790183) (4600,0.01785693187) (4700,0.01745183691) (4800,0.01726551466) (4900,0.0172366567) (5000,0.01714686062) (5100,0.01708784811) (5200,0.01685716119) (5300,0.01648265403) (5400,0.01641333066) (5500,0.01603855435) (5600,0.01587085668) (5700,0.01586549841) (5800,0.01592819635) (5900,0.01582984421) (6000,0.01564441137) (6100,0.01540855039) (6200,0.01546672117) (6300,0.01530352468) (6400,0.01508958004) (6500,0.0153247552) (6600,0.01530407816) (6700,0.01522395071) (6800,0.01513928529) (6900,0.01508452995) (7000,0.01525144847) (7100,0.0150590416) (7200,0.01497075837) (7300,0.01491020974) (7400,0.01495435322) (7500,0.01484032338) (7600,0.01473274454) (7700,0.01472336845) (7800,0.01480237506) (7900,0.01460443428) (8000,0.01429018108) (8100,0.01436144523) (8200,0.01405903511) (8300,0.01414198177) (8400,0.01393112754) (8500,0.01390321432) (8600,0.01407985538) (8700,0.01400401099) (8800,0.01381559521) (8900,0.01376753608) (9000,0.01377611971) (9100,0.0136788208) (9200,0.01385480799) (9300,0.01370892301) (9400,0.01357049346) (9500,0.01341768168) (9600,0.01318252943) (9700,0.01339062396) (9800,0.01329277391) (9900,0.01336767878) (10000,0.01330628349) (10100,0.0131931142) (10200,0.01312325569) (10300,0.01313092802) (10400,0.01327097248) (10500,0.0132885579) (10600,0.01362029724) (10700,0.01329969857) (10800,0.01318500191) (10900,0.01320387926) (11000,0.01312714526) (11100,0.01335652731) (11200,0.01342078978) (11300,0.01329366015) (11400,0.01331465011) (11500,0.01324214386) (11600,0.01283499608) (11700,0.01283802064) (11800,0.01290128771) (11900,0.0127480939) (12000,0.01277015395) (12100,0.01263874676) (12200,0.01244821148) (12300,0.01255683135) (12400,0.01233555963) (12500,0.01224935604) (12600,0.01246585418) (12700,0.01247371351) (12800,0.01246066894) (12900,0.01237296369) (13000,0.01224701125) (13100,0.01207559491) (13200,0.01207208997) (13300,0.01217194982) (13400,0.01243512873) (13500,0.01255085785) (13600,0.01225843532) (13700,0.01226580599) (13800,0.01223188695) (13900,0.01217831839) (14000,0.01222136533) (14100,0.0122908745) (14200,0.01235853415) (14300,0.01208850397) (14400,0.01189186256) (14500,0.01171710771) (14600,0.0118000662) (14700,0.01159551041) (14800,0.01168616638) (14900,0.01173347747) (15000,0.01176875699) (15100,0.0120031666) (15200,0.01191351116) (15300,0.01198667651) (15400,0.01210256079) (15500,0.0121144128) (15600,0.01226137532) (15700,0.01258138502) (15800,0.01251191292) (15900,0.01250924328) (16000,0.01245869799) (16100,0.01212765388) (16200,0.01209936542) (16300,0.01223681392) (16400,0.01219632793) (16500,0.01235993551) (16600,0.01203277856) (16700,0.01170334509) (16800,0.01169229997) (16900,0.01177214039) (17000,0.01193014411) (17100,0.01197794173) (17200,0.01204071138) (17300,0.01196989017) (17400,0.01192874499) (17500,0.01207821947) (17600,0.0122049069) (17700,0.01219846876) (17800,0.01229721224) (17900,0.01223322591) (18000,0.0122567025) (18100,0.01242561648) (18200,0.0123716699) (18300,0.01247492535) (18400,0.0124501477) (18500,0.01224804763) (18600,0.0122535795) (18700,0.01245142696) (18800,0.01226856522) (18900,0.01242658813) (19000,0.0122735342) (19100,0.01211060043) (19200,0.01214897707) (19300,0.01211500494) (19400,0.01205490213) (19500,0.01187685709) (19600,0.01179724149) (19700,0.01177117908) (19800,0.01207131855) (19900,0.01204899997) (20000,0.01204508431) (20100,0.01215868555) (20200,0.01214623619) (20300,0.0122318944) (20400,0.01228553923) (20500,0.01246843422) (20600,0.01231897259) (20700,0.01237359913) (20800,0.01227262635) (20900,0.01202838663) (21000,0.01195316818) (21100,0.01188557753) (21200,0.01192357922) (21300,0.01159028476) (21400,0.01156468289) (21500,0.0116171605) (21600,0.01192543413) (21700,0.01189131234) (21800,0.01184498351) (21900,0.01210313737) (22000,0.01225037891) (22100,0.01234588698) (22200,0.01241643336) (22300,0.01264403192) (22400,0.0127497538) (22500,0.01247960879) (22600,0.01253214562) (22700,0.01243551709) (22800,0.01244309442) (22900,0.01240402469) (23000,0.01241965611) (23100,0.01231919331) (23200,0.01219807556) (23300,0.01216566237) (23400,0.01199134523) (23500,0.01234305492) (23600,0.01216804283) (23700,0.01227127295) (23800,0.01236445233)};
\addplot[pubRose, opacity=0.18, line width=0.35pt] coordinates {(100,0.02349870652) (200,0.02028304152) (300,0.01814469757) (400,0.017596127) (500,0.0171731228) (600,0.01675780661) (700,0.01639749948) (800,0.01629597985) (900,0.01632127424) (1000,0.01613393044) (1100,0.01513775596) (1200,0.01546465093) (1300,0.01592921419) (1400,0.01579877241) (1500,0.01561761051) (1600,0.01565904561) (1700,0.01580230352) (1800,0.01571936319) (1900,0.0155568718) (2000,0.01574493246) (2100,0.01587965796) (2200,0.01535821268) (2300,0.01503826538) (2400,0.01522544296) (2500,0.01539484011) (2600,0.01544051049) (2700,0.01542130867) (2800,0.01541321203) (2900,0.01532360427) (3000,0.01498421654) (3100,0.01488712532) (3200,0.01499834778) (3300,0.01507737106) (3400,0.01487792972) (3500,0.01468235515) (3600,0.01462238133) (3700,0.01453882456) (3800,0.01448937533) (3900,0.01454572324) (4000,0.01460072696) (4100,0.01445144052) (4200,0.01416929616) (4300,0.0139382123) (4400,0.01374792662) (4500,0.01364080515) (4600,0.01345577436) (4700,0.01328857085) (4800,0.01333987992) (4900,0.01326647457) (5000,0.0134638763) (5100,0.0133962092) (5200,0.01328555765) (5300,0.01345070666) (5400,0.01347446879) (5500,0.01344878981) (5600,0.01345139565) (5700,0.01341180438) (5800,0.01311679753) (5900,0.01292415634) (6000,0.01262327228) (6100,0.01276872363) (6200,0.01281302394) (6300,0.01258620983) (6400,0.01267085364) (6500,0.01262376234) (6600,0.01247508563) (6700,0.01249183696) (6800,0.01265621986) (6900,0.0125340716) (7000,0.01257139202) (7100,0.01254403191) (7200,0.01235754751) (7300,0.01242417553) (7400,0.01226532822) (7500,0.01227963902) (7600,0.0122479517) (7700,0.01219935454) (7800,0.01213317337) (7900,0.01232526843) (8000,0.01236107498) (8100,0.0123130654) (8200,0.0124048803) (8300,0.01235629525) (8400,0.01247841781) (8500,0.01254419424) (8600,0.01269137384) (8700,0.01256488794) (8800,0.01241082232) (8900,0.0122497486) (9000,0.01199608035) (9100,0.01177090276) (9200,0.01169957975) (9300,0.01140546799) (9400,0.01140426993) (9500,0.01140779918) (9600,0.0113286457) (9700,0.01157792835) (9800,0.01154322447) (9900,0.01174125327) (10000,0.01190050095) (10100,0.01208663341) (10200,0.01219179612) (10300,0.01231790129) (10400,0.01217983756) (10500,0.01235387921) (10600,0.01246189727) (10700,0.01231393972) (10800,0.01232187012) (10900,0.01212041155) (11000,0.01210351726) (11100,0.01207548836) (11200,0.01196031692) (11300,0.0120883855) (11400,0.01212243196) (11500,0.01193152275) (11600,0.01189812394) (11700,0.01182433041) (11800,0.01180030815) (11900,0.01181702586) (12000,0.01171744419) (12100,0.01173267877) (12200,0.01193567896) (12300,0.01171279633) (12400,0.01157392114) (12500,0.01163175888) (12600,0.01151520433) (12700,0.01156371925) (12800,0.01178019159) (12900,0.01180607453) (13000,0.01177082369) (13100,0.01173384888) (13200,0.01165652638) (13300,0.01179322572) (13400,0.01198274214) (13500,0.01178338602) (13600,0.01180076916) (13700,0.01173412902) (13800,0.01150378007) (13900,0.01124888705) (14000,0.01120731253) (14100,0.01131496225) (14200,0.01135969544) (14300,0.0111358705) (14400,0.01103143394) (14500,0.01112116082) (14600,0.01104504867) (14700,0.01107418528) (14800,0.01107077738) (14900,0.01121213129) (15000,0.0114716812) (15100,0.01139143398) (15200,0.01124630095) (15300,0.01156020286) (15400,0.01167549901) (15500,0.01172670154) (15600,0.01182791013) (15700,0.01185964216) (15800,0.01201156871) (15900,0.01213093726) (16000,0.01199346585) (16100,0.01210327474) (16200,0.01213151123) (16300,0.0120019475) (16400,0.01175158024) (16500,0.01168312561) (16600,0.01153401975) (16700,0.0114051722) (16800,0.01132243192) (16900,0.01130110286) (17000,0.01127415802) (17100,0.01107872799) (17200,0.01102360142) (17300,0.01102523962) (17400,0.01100909794) (17500,0.01101557901) (17600,0.01106192814) (17700,0.01110468535) (17800,0.01110240789) (17900,0.01118731322) (18000,0.01115837423) (18100,0.01129256459) (18200,0.01129153501) (18300,0.01140367799) (18400,0.01156369867) (18500,0.01149783395) (18600,0.01146096764) (18700,0.01145751895) (18800,0.01144681433) (18900,0.01134819305) (19000,0.01141201789) (19100,0.01146853529) (19200,0.01133782305) (19300,0.01100343941) (19400,0.01088984599) (19500,0.01101401234) (19600,0.01094022263) (19700,0.01085564252) (19800,0.01099584186) (19900,0.01111997915) (20000,0.01109167654) (20100,0.01103552254) (20200,0.01112380484) (20300,0.0112192817) (20400,0.0112191421) (20500,0.01106847469) (20600,0.01139991358) (20700,0.01141432971) (20800,0.01129332269) (20900,0.01109075295) (21000,0.01106487494) (21100,0.01097171316) (21200,0.01109718969) (21300,0.01121629421) (21400,0.01122140661) (21500,0.01130960267) (21600,0.011175506) (21700,0.01144460356) (21800,0.01119816722) (21900,0.01122725355) (22000,0.01125813434) (22100,0.01127007138) (22200,0.01118989903) (22300,0.01116642905) (22400,0.01117674755) (22500,0.01110676648) (22600,0.01101643173) (22700,0.01071389671) (22800,0.01115127318) (22900,0.01126796212) (23000,0.01109443521) (23100,0.01108701658) (23200,0.01114775948) (23300,0.01108102035) (23400,0.01106595434) (23500,0.01099003572) (23600,0.011121016) (23700,0.01135463826) (23800,0.01106614061)};
\addplot[pubRose, opacity=0.18, line width=0.35pt] coordinates {(100,0.02408476919) (200,0.02206420898) (300,0.02026365201) (400,0.01893610135) (500,0.01832246259) (600,0.01795135128) (700,0.01724472163) (800,0.01705457503) (900,0.01657626095) (1000,0.01621031379) (1100,0.01528583076) (1200,0.01453766432) (1300,0.0142615566) (1400,0.01411196832) (1500,0.01390304063) (1600,0.01402237546) (1700,0.0141214421) (1800,0.01410561185) (1900,0.01450155824) (2000,0.01489491826) (2100,0.01517305449) (2200,0.01555335075) (2300,0.01540404139) (2400,0.01559651261) (2500,0.01571425889) (2600,0.01538352259) (2700,0.01556000449) (2800,0.01556770578) (2900,0.01538764164) (3000,0.01513240729) (3100,0.01495595109) (3200,0.01487476844) (3300,0.01498339316) (3400,0.01497656545) (3500,0.01476347009) (3600,0.01492730705) (3700,0.01479354165) (3800,0.01459735371) (3900,0.01445605895) (4000,0.01455581179) (4100,0.01429725504) (4200,0.01415272905) (4300,0.01443969412) (4400,0.01452459609) (4500,0.0146987481) (4600,0.01471237559) (4700,0.01454059957) (4800,0.01462408686) (4900,0.01469425838) (5000,0.01456164103) (5100,0.01460558092) (5200,0.01445927098) (5300,0.01416754201) (5400,0.01378022861) (5500,0.01390161105) (5600,0.01366294529) (5700,0.013718965) (5800,0.01377097024) (5900,0.01378515577) (6000,0.013780188) (6100,0.01382414512) (6200,0.01389453653) (6300,0.01378233247) (6400,0.01374359131) (6500,0.01336856075) (6600,0.01329982029) (6700,0.0131655274) (6800,0.01289283223) (6900,0.01292788889) (7000,0.01291730627) (7100,0.01266796943) (7200,0.0123999781) (7300,0.01246610619) (7400,0.01247442523) (7500,0.01251943856) (7600,0.01234786948) (7700,0.01218544673) (7800,0.0120138837) (7900,0.01167706316) (8000,0.0113581406) (8100,0.01145500792) (8200,0.01153511032) (8300,0.0113605449) (8400,0.01132477494) (8500,0.01114212265) (8600,0.01110516405) (8700,0.01126530599) (8800,0.01143678194) (8900,0.0113667204) (9000,0.0111620998) (9100,0.01099903993) (9200,0.01087590298) (9300,0.010752628) (9400,0.01076775845) (9500,0.01086946344) (9600,0.01108298106) (9700,0.0110344219) (9800,0.01105504297) (9900,0.01099864496) (10000,0.0113063802) (10100,0.01130639883) (10200,0.01136141783) (10300,0.01144335466) (10400,0.01134365303) (10500,0.01133376127) (10600,0.01123947417) (10700,0.01126085548) (10800,0.0111860822) (10900,0.01135199694) (11000,0.0112787433) (11100,0.01124063293) (11200,0.01117197582) (11300,0.0111231518) (11400,0.01102397079) (11500,0.01099191969) (11600,0.0109724435) (11700,0.01103324322) (11800,0.01122005088) (11900,0.01097790031) (12000,0.01133062011) (12100,0.01149221007) (12200,0.01165716238) (12300,0.0118858411) (12400,0.01215424193) (12500,0.01212707562) (12600,0.01210771436) (12700,0.01190719549) (12800,0.01161613343) (12900,0.0118475371) (13000,0.01153920302) (13100,0.01142915599) (13200,0.0114029929) (13300,0.01120169982) (13400,0.01139764888) (13500,0.01148077827) (13600,0.01141080474) (13700,0.01139892014) (13800,0.01142125223) (13900,0.01124219764) (14000,0.01115539148) (14100,0.01107380716) (14200,0.01089901226) (14300,0.01099248705) (14400,0.01097572576) (14500,0.01100855106) (14600,0.01096162032) (14700,0.011276739) (14800,0.0113206788) (14900,0.01160267359) (15000,0.01178358225) (15100,0.01188805001) (15200,0.01218826631) (15300,0.01210555146) (15400,0.01170400744) (15500,0.01170244971) (15600,0.01187267592) (15700,0.01180023877) (15800,0.01189281195) (15900,0.01162252799) (16000,0.01174054844) (16100,0.01174436603) (16200,0.01165926047) (16300,0.01173617495) (16400,0.01196125643) (16500,0.0120014186) (16600,0.012090596) (16700,0.01190595226) (16800,0.01191525068) (16900,0.01217287537) (17000,0.01193788247) (17100,0.01208624821) (17200,0.01198009532) (17300,0.01224376019) (17400,0.01205153121) (17500,0.01193218501) (17600,0.01172079295) (17700,0.01194411432) (17800,0.01194921853) (17900,0.01192687191) (18000,0.01221119389) (18100,0.01215666179) (18200,0.01217664909) (18300,0.01197244693) (18400,0.01210447932) (18500,0.01227529244) (18600,0.01247794405) (18700,0.01209855024) (18800,0.01211153781) (18900,0.01225981601) (19000,0.01209086161) (19100,0.01211084398) (19200,0.01216482483) (19300,0.01224012682) (19400,0.01254731286) (19500,0.01231934149) (19600,0.01217384255) (19700,0.0124370764) (19800,0.01225467166) (19900,0.01215343159) (20000,0.01205633562) (20100,0.01214000946) (20200,0.01222203225) (20300,0.01210463243) (20400,0.01188679691) (20500,0.0120585274) (20600,0.01206540205) (20700,0.01216930766) (20800,0.012191422) (20900,0.01210096963) (21000,0.01200001352) (21100,0.01173582915) (21200,0.01135093793) (21300,0.01131332107) (21400,0.01112062223) (21500,0.01114228666) (21600,0.01120283296) (21700,0.01088391831) (21800,0.01107871579) (21900,0.01096417671) (22000,0.01101914821) (22100,0.01109680841) (22200,0.01124488777) (22300,0.01128822574) (22400,0.01135225939) (22500,0.01120230295) (22600,0.01108895512) (22700,0.01118067121) (22800,0.01096195756) (22900,0.01106311958) (23000,0.01107661035) (23100,0.01114473678) (23200,0.01118944958) (23300,0.01132056173) (23400,0.01135056596) (23500,0.01139241671) (23600,0.01133243535) (23700,0.01146468185) (23800,0.01162557993)};
\addplot[pubRose, opacity=0.18, line width=0.35pt] coordinates {(100,0.02315595187) (200,0.02122516464) (300,0.01972457953) (400,0.0194468121) (500,0.01910330132) (600,0.01822371657) (700,0.01758556321) (800,0.01720658096) (900,0.01679870641) (1000,0.01653223457) (1100,0.01545568658) (1200,0.01498056902) (1300,0.01462362157) (1400,0.01393396873) (1500,0.01344981398) (1600,0.01324510109) (1700,0.01309549315) (1800,0.01313511375) (1900,0.01311443541) (2000,0.01302654818) (2100,0.01311221709) (2200,0.01268295255) (2300,0.01246950915) (2400,0.01260016123) (2500,0.01257436285) (2600,0.01273663985) (2700,0.0127347487) (2800,0.01259739287) (2900,0.01248228941) (3000,0.01245961273) (3100,0.01244501024) (3200,0.01268852139) (3300,0.01292847283) (3400,0.01284225406) (3500,0.01290903213) (3600,0.01293697534) (3700,0.01307168351) (3800,0.01287829159) (3900,0.01300668987) (4000,0.01315419208) (4100,0.01310357396) (4200,0.01300391806) (4300,0.01306693545) (4400,0.01312719854) (4500,0.01310215043) (4600,0.01297622519) (4700,0.012876096) (4800,0.0129100075) (4900,0.01294501266) (5000,0.01267127851) (5100,0.012827418) (5200,0.01296930192) (5300,0.01280255131) (5400,0.01283686943) (5500,0.01280756565) (5600,0.01295240643) (5700,0.0130595454) (5800,0.01317968434) (5900,0.01313544298) (6000,0.01330134599) (6100,0.01328469878) (6200,0.01319320323) (6300,0.01322752349) (6400,0.01322507299) (6500,0.01332207788) (6600,0.01325827818) (6700,0.01325944951) (6800,0.0130973652) (6900,0.01305198902) (7000,0.0130791612) (7100,0.01290489985) (7200,0.01312144762) (7300,0.01335617034) (7400,0.01349503035) (7500,0.01361930966) (7600,0.0135241827) (7700,0.01349869817) (7800,0.01379326722) (7900,0.01402789885) (8000,0.01417120304) (8100,0.01424140055) (8200,0.01422955533) (8300,0.01413558722) (8400,0.01420650166) (8500,0.01421908494) (8600,0.01413592715) (8700,0.01425575614) (8800,0.01393772811) (8900,0.01387694543) (9000,0.01375913061) (9100,0.01400878271) (9200,0.01401059693) (9300,0.01392375324) (9400,0.01383032165) (9500,0.01372075798) (9600,0.01374529069) (9700,0.01368071577) (9800,0.014133514) (9900,0.01394955423) (10000,0.01404930661) (10100,0.01384164691) (10200,0.01367064761) (10300,0.013758964) (10400,0.0137294692) (10500,0.01366728358) (10600,0.01383284573) (10700,0.01374162659) (10800,0.01350131137) (10900,0.01343220389) (11000,0.01331095463) (11100,0.01322772438) (11200,0.01321643172) (11300,0.01306761922) (11400,0.01301868381) (11500,0.01303515313) (11600,0.01317656152) (11700,0.01323127206) (11800,0.01304708337) (11900,0.01320360107) (12000,0.0132944379) (12100,0.01337892842) (12200,0.01363865789) (12300,0.01370568546) (12400,0.01354514332) (12500,0.01351293996) (12600,0.0133186928) (12700,0.01362696504) (12800,0.01388807651) (12900,0.01382042645) (13000,0.01370296422) (13100,0.01372200847) (13200,0.01336480025) (13300,0.01340733962) (13400,0.01346874628) (13500,0.01342327893) (13600,0.01347753247) (13700,0.01307032453) (13800,0.01293741949) (13900,0.01307191141) (14000,0.01310720537) (14100,0.01312891478) (14200,0.01355749303) (14300,0.01352559021) (14400,0.01356804315) (14500,0.01346158162) (14600,0.01332585709) (14700,0.01344249761) (14800,0.0134397543) (14900,0.01315704696) (15000,0.01301892782) (15100,0.01281970423) (15200,0.01253724499) (15300,0.01233059391) (15400,0.01234162608) (15500,0.01254378166) (15600,0.01275932509) (15700,0.01270108204) (15800,0.01260973755) (15900,0.01287197685) (16000,0.0127701316) (16100,0.01281786971) (16200,0.01257520318) (16300,0.01255274527) (16400,0.01256076926) (16500,0.01251201397) (16600,0.01230560848) (16700,0.01225680616) (16800,0.0123776109) (16900,0.01240175236) (17000,0.01251085717) (17100,0.01247564079) (17200,0.01270421762) (17300,0.01278007142) (17400,0.01269467007) (17500,0.01264874619) (17600,0.01279916111) (17700,0.01270721639) (17800,0.01257565608) (17900,0.01249456648) (18000,0.01233446822) (18100,0.01237173397) (18200,0.01250492968) (18300,0.01269070851) (18400,0.01264960524) (18500,0.01258694232) (18600,0.01256996542) (18700,0.01267707963) (18800,0.01268584831) (18900,0.01254409887) (19000,0.01286971699) (19100,0.01281296322) (19200,0.0126936431) (19300,0.01235356107) (19400,0.01233405983) (19500,0.0123354801) (19600,0.01203272706) (19700,0.01185211176) (19800,0.01181965778) (19900,0.01175360726) (20000,0.01146498201) (20100,0.01153315334) (20200,0.01147336438) (20300,0.01162096709) (20400,0.01160434261) (20500,0.01165251313) (20600,0.01168423006) (20700,0.01187140169) (20800,0.01204449283) (20900,0.01200878052) (21000,0.01215537209) (21100,0.0121115217) (21200,0.01206602305) (21300,0.01210049177) (21400,0.01214887528) (21500,0.01200385876) (21600,0.01220541876) (21700,0.01208815388) (21800,0.0120727662) (21900,0.01218512096) (22000,0.01211008951) (22100,0.01188056115) (22200,0.01183506278) (22300,0.01193927936) (22400,0.01204299144) (22500,0.01199467331) (22600,0.01187340543) (22700,0.01173875146) (22800,0.01150543587) (22900,0.01138772545) (23000,0.01158043798) (23100,0.01176083833) (23200,0.01185442619) (23300,0.01167041538) (23400,0.01172067812) (23500,0.01193317249) (23600,0.01207799688) (23700,0.0124616459) (23800,0.01267289864)};
\addplot[pubRose, ultra thick] coordinates {(100,0.02448844537) (200,0.02234271308) (300,0.02027749239) (400,0.02033067006) (500,0.01985250041) (600,0.01970454787) (700,0.01966290761) (800,0.01936353504) (900,0.01901302771) (1000,0.01882351199) (1100,0.01830579918) (1200,0.01817424046) (1300,0.01814111639) (1400,0.01795260757) (1500,0.01810224624) (1600,0.01815821128) (1700,0.01804632619) (1800,0.01819392852) (1900,0.01812635558) (2000,0.01771442574) (2100,0.01788785439) (2200,0.01794747924) (2300,0.01792113604) (2400,0.01795557328) (2500,0.01787381824) (2600,0.01800047867) (2700,0.01796270208) (2800,0.01784764696) (2900,0.01751336185) (3000,0.01732939929) (3100,0.01721307999) (3200,0.01711914605) (3300,0.01705979118) (3400,0.01714089559) (3500,0.01690805797) (3600,0.0167782885) (3700,0.01643913779) (3800,0.01625467977) (3900,0.01629910823) (4000,0.01623560777) (4100,0.01620991952) (4200,0.01618506354) (4300,0.01616532239) (4400,0.01589796306) (4500,0.01572248372) (4600,0.01557997512) (4700,0.01574288742) (4800,0.01571673616) (4900,0.01562111778) (5000,0.01550749405) (5100,0.0154137325) (5200,0.01522424221) (5300,0.01518047359) (5400,0.01504633739) (5500,0.01498216027) (5600,0.01502849357) (5700,0.01494329683) (5800,0.01476155869) (5900,0.01466794051) (6000,0.01454938585) (6100,0.01453388827) (6200,0.01445158958) (6300,0.0143725296) (6400,0.0143204743) (6500,0.0142702084) (6600,0.01431957157) (6700,0.01419830937) (6800,0.01407627938) (6900,0.01409187713) (7000,0.01389190289) (7100,0.01391715747) (7200,0.0139335278) (7300,0.0139996422) (7400,0.01390303061) (7500,0.01376535338) (7600,0.01374067594) (7700,0.01369810598) (7800,0.01372815459) (7900,0.01388450484) (8000,0.01394142252) (8100,0.01399139543) (8200,0.01382251359) (8300,0.01372721484) (8400,0.0135786789) (8500,0.01359035638) (8600,0.01361686485) (8700,0.01347357361) (8800,0.0134876295) (8900,0.01344936206) (9000,0.01341717257) (9100,0.01332661523) (9200,0.01326940302) (9300,0.01338694938) (9400,0.01328520216) (9500,0.01318081887) (9600,0.01315059294) (9700,0.01331482856) (9800,0.01331493361) (9900,0.01318873218) (10000,0.01309648775) (10100,0.01313966461) (10200,0.01311303959) (10300,0.01298332405) (10400,0.01286786636) (10500,0.01280215862) (10600,0.01263813796) (10700,0.0129073143) (10800,0.0128539219) (10900,0.0128310191) (11000,0.01268358654) (11100,0.0125945277) (11200,0.01283406857) (11300,0.01279145558) (11400,0.01270739008) (11500,0.01270607836) (11600,0.01260433444) (11700,0.01243389379) (11800,0.01257966091) (11900,0.01253333925) (12000,0.01259207493) (12100,0.01249622894) (12200,0.01244302294) (12300,0.01236458588) (12400,0.01224490078) (12500,0.01218821583) (12600,0.01228678427) (12700,0.0121904545) (12800,0.0121551035) (12900,0.0121515424) (13000,0.01206019032) (13100,0.0120974273) (13200,0.01198350135) (13300,0.01201956505) (13400,0.01224482902) (13500,0.01223227838) (13600,0.01210269271) (13700,0.01216928028) (13800,0.01224327399) (13900,0.0122454958) (14000,0.01221988276) (14100,0.01229770044) (14200,0.01232523299) (14300,0.01212476185) (14400,0.01204617945) (14500,0.01206284156) (14600,0.01205560034) (14700,0.01215028102) (14800,0.01225433098) (14900,0.01218146747) (15000,0.01208642377) (15100,0.01213048678) (15200,0.01224334701) (15300,0.01221807268) (15400,0.01210904745) (15500,0.0120627597) (15600,0.01208752915) (15700,0.01217673179) (15800,0.01213730904) (15900,0.01213608938) (16000,0.01204219614) (16100,0.01211546431) (16200,0.01211543833) (16300,0.01211938071) (16400,0.01207879218) (16500,0.01218067706) (16600,0.01206168728) (16700,0.01180464868) (16800,0.01180409887) (16900,0.01198595571) (17000,0.01193401329) (17100,0.01203209497) (17200,0.01201040335) (17300,0.01224364792) (17400,0.01206373749) (17500,0.0121766963) (17600,0.01241026488) (17700,0.01240404011) (17800,0.01243643416) (17900,0.01236389619) (18000,0.01229558536) (18100,0.01239867522) (18200,0.01243829979) (18300,0.01252528206) (18400,0.01254987647) (18500,0.01243111738) (18600,0.01252395473) (18700,0.01247456893) (18800,0.01234211889) (18900,0.01233905884) (19000,0.01227687146) (19100,0.01213404271) (19200,0.01218565498) (19300,0.01218607863) (19400,0.0120621826) (19500,0.01204894208) (19600,0.01191498428) (19700,0.01181164542) (19800,0.01194548816) (19900,0.01206393531) (20000,0.0119934136) (20100,0.01214766293) (20200,0.01211121907) (20300,0.01212754063) (20400,0.01199832251) (20500,0.01211269754) (20600,0.01211279831) (20700,0.01216500904) (20800,0.01211795742) (20900,0.0120358923) (21000,0.01206453894) (21100,0.01199854962) (21200,0.01199480114) (21300,0.01184538826) (21400,0.01185677908) (21500,0.01181050963) (21600,0.01206542645) (21700,0.01214627768) (21800,0.01218264662) (21900,0.01214412916) (22000,0.01203234186) (22100,0.01194483633) (22200,0.0119206605) (22300,0.01200124789) (22400,0.01199272126) (22500,0.01211321433) (22600,0.01222923733) (22700,0.0119626645) (22800,0.01192816859) (22900,0.01185247991) (23000,0.01194707016) (23100,0.01198581839) (23200,0.01202625087) (23300,0.01194930612) (23400,0.01185601167) (23500,0.01195245129) (23600,0.01208833721) (23700,0.01222674814) (23800,0.0120533166)};
\nextgroupplot[title={Actor gradient norm}, xlabel={Gradient update}, scaled y ticks=true]
\addplot[pubRose, opacity=0.18, line width=0.35pt] coordinates {(100,0.01819197647) (200,0.01780218538) (300,0.01765456796) (400,0.01755484007) (500,0.01733709425) (600,0.01685848677) (700,0.0161413887) (800,0.01558269525) (900,0.01517973695) (1000,0.01484835409) (1100,0.01413797103) (1200,0.01331247501) (1300,0.01268164255) (1400,0.01174031561) (1500,0.01106224023) (1600,0.01072729779) (1700,0.01052679606) (1800,0.01033675438) (1900,0.01014356734) (2000,0.00962553178) (2100,0.009282469889) (2200,0.009177455818) (2300,0.008961563976) (2400,0.00902459654) (2500,0.008769137179) (2600,0.007983317156) (2700,0.007643873454) (2800,0.007452754606) (2900,0.007098629349) (3000,0.007111728727) (3100,0.007245456963) (3200,0.008191432967) (3300,0.007803566218) (3400,0.008180869767) (3500,0.008464106382) (3600,0.008923265291) (3700,0.009262149455) (3800,0.009689723561) (3900,0.009697412234) (4000,0.00959723657) (4100,0.009740758222) (4200,0.009246361628) (4300,0.009568562964) (4400,0.008926839242) (4500,0.009172894014) (4600,0.009307966335) (4700,0.008798398403) (4800,0.008534465311) (4900,0.008525953302) (5000,0.009157394292) (5100,0.009049299685) (5200,0.008343950892) (5300,0.008706138795) (5400,0.008699823776) (5500,0.008546247473) (5600,0.00836032615) (5700,0.009007127816) (5800,0.008759846305) (5900,0.009077817854) (6000,0.00869910894) (6100,0.009400929976) (6200,0.00978702493) (6300,0.01009173281) (6400,0.0108600148) (6500,0.01032902193) (6600,0.01094669676) (6700,0.01023482783) (6800,0.01040719496) (6900,0.01007460412) (7000,0.009813675936) (7100,0.009258240554) (7200,0.009498205595) (7300,0.008789851516) (7400,0.008493650798) (7500,0.008451711433) (7600,0.008163869986) (7700,0.009196740203) (7800,0.008898713719) (7900,0.009549479559) (8000,0.009867506288) (8100,0.009763659257) (8200,0.009296522941) (8300,0.009163230518) (8400,0.008567928756) (8500,0.009015090903) (8600,0.009099909617) (8700,0.008861675067) (8800,0.00984331728) (8900,0.009948858572) (9000,0.009968900541) (9100,0.01060188315) (9200,0.01108632046) (9300,0.0111641617) (9400,0.01213858244) (9500,0.0118315333) (9600,0.01125587849) (9700,0.01105169132) (9800,0.01018963614) (9900,0.0103321637) (10000,0.009908195212) (10100,0.00967347566) (10200,0.009185984125) (10300,0.00944420509) (10400,0.008732876787) (10500,0.008764896635) (10600,0.00951791727) (10700,0.009946918814) (10800,0.01072893939) (10900,0.01083496283) (11000,0.01123166145) (11100,0.01093833749) (11200,0.01142893247) (11300,0.011434469) (11400,0.01219386584) (11500,0.01307151848) (11600,0.01348599754) (11700,0.01330548879) (11800,0.01298144739) (11900,0.01213443391) (12000,0.01369872363) (12100,0.01375676505) (12200,0.01361621413) (12300,0.01410568394) (12400,0.01418610979) (12500,0.01459002569) (12600,0.01498194579) (12700,0.01533999396) (12800,0.01606185874) (12900,0.0162380727) (13000,0.01520778164) (13100,0.01585983941) (13200,0.01579216421) (13300,0.01527060224) (13400,0.01498661404) (13500,0.01645891061) (13600,0.01619626721) (13700,0.01627099356) (13800,0.01614167923) (13900,0.01639325805) (14000,0.01570543218) (14100,0.0154516574) (14200,0.01581151793) (14300,0.01584361158) (14400,0.01592670986) (14500,0.0137176468) (14600,0.01483260011) (14700,0.01398590598) (14800,0.01383384559) (14900,0.0141124214) (15000,0.01487906352) (15100,0.01487375032) (15200,0.01547865076) (15300,0.01700767614) (15400,0.01785847927) (15500,0.01820842931) (15600,0.01690608701) (15700,0.01859972998) (15800,0.0178147193) (15900,0.01890482614) (16000,0.01921221074) (16100,0.01891207676) (16200,0.01852128245) (16300,0.01733277775) (16400,0.01703511775) (16500,0.01686010174) (16600,0.01696054116) (16700,0.01662567817) (16800,0.01727014892) (16900,0.01648632381) (17000,0.01674448904) (17100,0.01775224674) (17200,0.01902951561) (17300,0.01967717279) (17400,0.01899539325) (17500,0.01953335498) (17600,0.02036619578) (17700,0.02026630063) (17800,0.02138303742) (17900,0.02232286725) (18000,0.02183555402) (18100,0.02107904414) (18200,0.01981920032) (18300,0.01986169415) (18400,0.02130639562) (18500,0.02247152058) (18600,0.02207007864) (18700,0.02252180697) (18800,0.02234549234) (18900,0.02211242141) (19000,0.02408709386) (19100,0.02416263176) (19200,0.02485243166) (19300,0.02465141742) (19400,0.02511064121) (19500,0.02489579665) (19600,0.0236862028) (19700,0.02258293321) (19800,0.02328530839) (19900,0.02368269963) (20000,0.02240023883) (20100,0.02214223482) (20200,0.02136601508) (20300,0.02139291819) (20400,0.02094728798) (20500,0.01946587386) (20600,0.02005949691) (20700,0.02107665455) (20800,0.01974167014) (20900,0.01980379811) (21000,0.01928877076) (21100,0.02065303782) (21200,0.02015783489) (21300,0.02068665903) (21400,0.02013955321) (21500,0.02030876512) (21600,0.0209618154) (21700,0.01996499859) (21800,0.02060114946) (21900,0.02006808035) (22000,0.01987769399) (22100,0.01898337416) (22200,0.0202843179) (22300,0.01976683727) (22400,0.02022182038) (22500,0.02104637856) (22600,0.02196144657) (22700,0.02394472267) (22800,0.02349278741) (22900,0.02319447827) (23000,0.02406063396) (23100,0.02465209942) (23200,0.02387682619) (23300,0.02349494053) (23400,0.02290708115) (23500,0.02217765404) (23600,0.0216225368) (23700,0.0221005355) (23800,0.02123423424)};
\addplot[pubRose, opacity=0.18, line width=0.35pt] coordinates {(100,0.02152289264) (200,0.01962547284) (300,0.01853129702) (400,0.01738338196) (500,0.0161221277) (600,0.01518330195) (700,0.01477405881) (800,0.01442843361) (900,0.01426666064) (1000,0.01420830935) (1100,0.01334828464) (1200,0.01300433427) (1300,0.01241541347) (1400,0.01225320743) (1500,0.01209023809) (1600,0.01217856454) (1700,0.012268674) (1800,0.01212199312) (1900,0.01178246187) (2000,0.01139561059) (2100,0.01088621151) (2200,0.01105578421) (2300,0.01086987527) (2400,0.01054113237) (2500,0.01040345998) (2600,0.01112513626) (2700,0.01081011994) (2800,0.01072770869) (2900,0.01053453442) (3000,0.01152763143) (3100,0.01133950297) (3200,0.01135330014) (3300,0.01119684107) (3400,0.01177484137) (3500,0.0119943636) (3600,0.01083204178) (3700,0.01120649166) (3800,0.01157363961) (3900,0.01187350852) (4000,0.01099334247) (4100,0.0114860272) (4200,0.01081523513) (4300,0.01074387738) (4400,0.01021590494) (4500,0.01020739144) (4600,0.01027730722) (4700,0.009812793508) (4800,0.009072695114) (4900,0.00880367728) (5000,0.00858867364) (5100,0.007990723895) (5200,0.007862320589) (5300,0.008078656439) (5400,0.008413767163) (5500,0.008132779412) (5600,0.008143154578) (5700,0.008047257131) (5800,0.01051768227) (5900,0.01094551957) (6000,0.0111320457) (6100,0.01164099006) (6200,0.01165610421) (6300,0.01285727425) (6400,0.01371161081) (6500,0.01495879227) (6600,0.01560645895) (6700,0.01647247076) (6800,0.01466182359) (6900,0.01547887139) (7000,0.01656907964) (7100,0.01646875832) (7200,0.01776530016) (7300,0.0178223053) (7400,0.01722859778) (7500,0.0165356176) (7600,0.01659182701) (7700,0.01586160352) (7800,0.01536544096) (7900,0.01407349408) (8000,0.01308816904) (8100,0.01408742042) (8200,0.01429964146) (8300,0.0145888607) (8400,0.01442170516) (8500,0.01536614932) (8600,0.01558115408) (8700,0.01613470707) (8800,0.0185905478) (8900,0.01876180954) (9000,0.02037734604) (9100,0.01992719127) (9200,0.01849744786) (9300,0.01775822509) (9400,0.0178945235) (9500,0.01822855109) (9600,0.01877935855) (9700,0.02041672394) (9800,0.01863680566) (9900,0.02098606825) (10000,0.01988601126) (10100,0.02043380477) (10200,0.02051552013) (10300,0.02058974672) (10400,0.02278627902) (10500,0.02162537118) (10600,0.02188945366) (10700,0.01981958868) (10800,0.02094820365) (10900,0.01937478054) (11000,0.01909935074) (11100,0.01829239223) (11200,0.02053805748) (11300,0.02057734663) (11400,0.01886761049) (11500,0.01906450745) (11600,0.01780526172) (11700,0.01813637139) (11800,0.01701880088) (11900,0.01715175072) (12000,0.01899092533) (12100,0.02115145298) (12200,0.02088077003) (12300,0.0212818577) (12400,0.02225010702) (12500,0.02380963573) (12600,0.02404045491) (12700,0.02453226112) (12800,0.02517747283) (12900,0.02504530512) (13000,0.02465979699) (13100,0.02271406008) (13200,0.02176065398) (13300,0.02313483777) (13400,0.02307279846) (13500,0.02218749868) (13600,0.02342627225) (13700,0.02525077062) (13800,0.02725142306) (13900,0.02647846155) (14000,0.02482478246) (14100,0.0265854503) (14200,0.02659777282) (14300,0.02699580556) (14400,0.02700726194) (14500,0.02575015109) (14600,0.02512203269) (14700,0.02265865412) (14800,0.02065627705) (14900,0.02080552308) (15000,0.02249300359) (15100,0.02240871144) (15200,0.02492750799) (15300,0.02426562002) (15400,0.02293227399) (15500,0.02351142783) (15600,0.02578172963) (15700,0.02809891477) (15800,0.02866199929) (15900,0.02963674106) (16000,0.02984582298) (16100,0.02932716906) (16200,0.02661840059) (16300,0.02492355686) (16400,0.0252075525) (16500,0.02671343181) (16600,0.02646721955) (16700,0.02594859377) (16800,0.02625960708) (16900,0.02656875029) (17000,0.02687092181) (17100,0.02920123469) (17200,0.03076988366) (17300,0.03078878988) (17400,0.0316693116) (17500,0.03070822954) (17600,0.02891702596) (17700,0.02925286815) (17800,0.02875540927) (17900,0.02992994972) (18000,0.02845598403) (18100,0.02583656386) (18200,0.02509945799) (18300,0.02587009016) (18400,0.02507903334) (18500,0.02560020536) (18600,0.02518495582) (18700,0.02433483154) (18800,0.02359381597) (18900,0.02485520337) (19000,0.02454352938) (19100,0.02412272003) (19200,0.02307460513) (19300,0.02308907825) (19400,0.02504818421) (19500,0.02457549106) (19600,0.02534780297) (19700,0.0241781841) (19800,0.02572640097) (19900,0.02286265949) (20000,0.02298742821) (20100,0.02411164166) (20200,0.02399511794) (20300,0.02455826001) (20400,0.02272666069) (20500,0.02302932022) (20600,0.0224925383) (20700,0.02415606398) (20800,0.02328207288) (20900,0.02434243634) (21000,0.02461697683) (21100,0.02624590453) (21200,0.02764816526) (21300,0.02652903032) (21400,0.02620918564) (21500,0.02533835061) (21600,0.0257571701) (21700,0.0271802213) (21800,0.02888337262) (21900,0.02792700455) (22000,0.02840945274) (22100,0.02866681777) (22200,0.02825727891) (22300,0.02802983858) (22400,0.02883142885) (22500,0.03145845793) (22600,0.03106544465) (22700,0.03158262633) (22800,0.03080538213) (22900,0.03115629125) (23000,0.0311069563) (23100,0.02850760017) (23200,0.0279550083) (23300,0.02995232865) (23400,0.0297409324) (23500,0.02740178052) (23600,0.02666284991) (23700,0.02424996654) (23800,0.02371086767)};
\addplot[pubRose, opacity=0.18, line width=0.35pt] coordinates {(100,0.007709387224) (200,0.009934552712) (300,0.01036464656) (400,0.009946956648) (500,0.009208057541) (600,0.009214694689) (700,0.009433559076) (800,0.009820098348) (900,0.01018582103) (1000,0.01061423984) (1100,0.01164997043) (1200,0.01256118463) (1300,0.01383158499) (1400,0.01530667255) (1500,0.01674238201) (1600,0.01768874303) (1700,0.01724265558) (1800,0.01621273528) (1900,0.01506588543) (2000,0.01382642714) (2100,0.01229006464) (2200,0.01049954602) (2300,0.008466681512) (2400,0.006643393543) (2500,0.005093654711) (2600,0.00373528786) (2700,0.003590481635) (2800,0.003727440513) (2900,0.003922238946) (3000,0.004279522644) (3100,0.004592382186) (3200,0.004664030438) (3300,0.004801797261) (3400,0.00474254638) (3500,0.004893973609) (3600,0.007185674133) (3700,0.007584219985) (3800,0.01007106027) (3900,0.01152435751) (4000,0.01223156517) (4100,0.01354250088) (4200,0.01355482242) (4300,0.01454989966) (4400,0.01496886648) (4500,0.01509121354) (4600,0.01359460712) (4700,0.01322591645) (4800,0.01151673659) (4900,0.01014434351) (5000,0.01047213692) (5100,0.009457670897) (5200,0.009807093814) (5300,0.00931932414) (5400,0.008868313581) (5500,0.009086683951) (5600,0.008269735612) (5700,0.008623903431) (5800,0.008211120078) (5900,0.008422331465) (6000,0.00738300588) (6100,0.007006708439) (6200,0.007729710406) (6300,0.008064552909) (6400,0.008157615364) (6500,0.00795395961) (6600,0.008270190377) (6700,0.008494754415) (6800,0.009564483492) (6900,0.009611713747) (7000,0.009535454726) (7100,0.009374860837) (7200,0.00888037018) (7300,0.00827039571) (7400,0.008711537835) (7500,0.008313548542) (7600,0.008103738283) (7700,0.008226248273) (7800,0.007891231426) (7900,0.008209096338) (8000,0.008439636906) (8100,0.00976736024) (8200,0.009543421306) (8300,0.009479825571) (8400,0.00954667693) (8500,0.01086216611) (8600,0.01208906639) (8700,0.01169328894) (8800,0.01279841252) (8900,0.01243443443) (9000,0.01500804825) (9100,0.01671905974) (9200,0.01732676625) (9300,0.01754899807) (9400,0.01975532016) (9500,0.01895226194) (9600,0.01863162527) (9700,0.01890873462) (9800,0.0217082141) (9900,0.02335946439) (10000,0.02082859185) (10100,0.01901351009) (10200,0.02238386702) (10300,0.02271362152) (10400,0.02221098598) (10500,0.02226367276) (10600,0.02196934465) (10700,0.02265943279) (10800,0.01784221004) (10900,0.01691507208) (11000,0.0172879918) (11100,0.01725449101) (11200,0.01322410414) (11300,0.01370604979) (11400,0.01227420373) (11500,0.01379516022) (11600,0.01606969656) (11700,0.01733209277) (11800,0.01894625244) (11900,0.01889297464) (12000,0.01912001087) (12100,0.01896550967) (12200,0.01927303001) (12300,0.01879058136) (12400,0.01846311223) (12500,0.01843341682) (12600,0.01847737152) (12700,0.0172405649) (12800,0.0179977132) (12900,0.01833410133) (13000,0.01730653632) (13100,0.01804354794) (13200,0.01782809009) (13300,0.01749870002) (13400,0.01794631879) (13500,0.01707225821) (13600,0.01547869917) (13700,0.01418532354) (13800,0.01390744017) (13900,0.01284562491) (14000,0.01313771131) (14100,0.01260727802) (14200,0.01264791088) (14300,0.01318966672) (14400,0.01517088981) (14500,0.01448462568) (14600,0.01376236686) (14700,0.0144445857) (14800,0.01283633038) (14900,0.01694245604) (15000,0.01916478192) (15100,0.01893476956) (15200,0.01908120411) (15300,0.02113380395) (15400,0.01901465002) (15500,0.01893039607) (15600,0.0183436628) (15700,0.01824028543) (15800,0.02074736501) (15900,0.01695913938) (16000,0.01611524713) (16100,0.0157987569) (16200,0.01588307568) (16300,0.01479038787) (16400,0.0163606741) (16500,0.01767544313) (16600,0.01833892055) (16700,0.01912414189) (16800,0.01681726296) (16900,0.01797865937) (17000,0.0183746404) (17100,0.0202274479) (17200,0.02047204673) (17300,0.01977045983) (17400,0.01796994479) (17500,0.01769650457) (17600,0.01741453735) (17700,0.01721941708) (17800,0.01658345358) (17900,0.01568520144) (18000,0.0146527743) (18100,0.01253516017) (18200,0.01363254827) (18300,0.01325923023) (18400,0.01391779236) (18500,0.01381448987) (18600,0.01562063103) (18700,0.01554734693) (18800,0.01599828047) (18900,0.01770295254) (19000,0.02161794486) (19100,0.02302114945) (19200,0.02310599182) (19300,0.02417371757) (19400,0.02368745916) (19500,0.02307327827) (19600,0.02157563018) (19700,0.02074664757) (19800,0.02290261313) (19900,0.02098805113) (20000,0.01860549552) (20100,0.01872667419) (20200,0.01945419973) (20300,0.01832257994) (20400,0.01865930948) (20500,0.01856466755) (20600,0.01931237765) (20700,0.02101400523) (20800,0.01990931677) (20900,0.02159720771) (21000,0.02055775039) (21100,0.02004128415) (21200,0.01795038413) (21300,0.01784992209) (21400,0.01815489205) (21500,0.01824660357) (21600,0.01747169793) (21700,0.01648724535) (21800,0.01548694419) (21900,0.01481238781) (22000,0.0149496601) (22100,0.01519598896) (22200,0.01542012673) (22300,0.01551114442) (22400,0.01575539587) (22500,0.01572719859) (22600,0.01660325872) (22700,0.01803053711) (22800,0.01917880643) (22900,0.01863576025) (23000,0.01957190111) (23100,0.01917884499) (23200,0.0194599201) (23300,0.01967914039) (23400,0.01846446181) (23500,0.01971716369) (23600,0.01827617711) (23700,0.01674146722) (23800,0.0154013)};
\addplot[pubRose, opacity=0.18, line width=0.35pt] coordinates {(100,0.006018080283) (200,0.007504918845) (300,0.007678725912) (400,0.007945799385) (500,0.007707326021) (600,0.008204533951) (700,0.008977147418) (800,0.009357468516) (900,0.00975958646) (1000,0.01009223307) (1100,0.01064498732) (1200,0.01065111756) (1300,0.01056527114) (1400,0.01037100204) (1500,0.01030779821) (1600,0.00997346458) (1700,0.009260373283) (1800,0.008862811234) (1900,0.008383326791) (2000,0.007768205274) (2100,0.007258386025) (2200,0.007136324933) (2300,0.007010010956) (2400,0.007155412063) (2500,0.007272976357) (2600,0.007231442956) (2700,0.007719199173) (2800,0.007446136139) (2900,0.00797212543) (3000,0.008141332027) (3100,0.008096355991) (3200,0.008108799858) (3300,0.009020499373) (3400,0.009458717192) (3500,0.009519917378) (3600,0.009801605716) (3700,0.009241296584) (3800,0.00993407988) (3900,0.01018051053) (4000,0.0119604791) (4100,0.0118668416) (4200,0.01237390707) (4300,0.01196924294) (4400,0.01208240283) (4500,0.01322710966) (4600,0.0138278658) (4700,0.01426270651) (4800,0.01489561619) (4900,0.01457019504) (5000,0.01287969118) (5100,0.01661176858) (5200,0.01631900743) (5300,0.01710513653) (5400,0.01810479099) (5500,0.01696299277) (5600,0.01655225661) (5700,0.01921667596) (5800,0.01931616617) (5900,0.01979213543) (6000,0.01996599212) (6100,0.01706746137) (6200,0.01848193156) (6300,0.01812079521) (6400,0.01757837655) (6500,0.01908032428) (6600,0.02139871288) (6700,0.01898815362) (6800,0.01906290809) (6900,0.01991286865) (7000,0.02148813196) (7100,0.0220810608) (7200,0.0211216935) (7300,0.02116586911) (7400,0.02095699674) (7500,0.01955440752) (7600,0.01710204193) (7700,0.01869480051) (7800,0.01762550334) (7900,0.01815619087) (8000,0.01671950696) (8100,0.01639746623) (8200,0.01627935125) (8300,0.01615559896) (8400,0.01662432412) (8500,0.01762922518) (8600,0.01791863283) (8700,0.01742591383) (8800,0.01775815375) (8900,0.01809763983) (9000,0.01846186239) (9100,0.01847889759) (9200,0.01927275117) (9300,0.01931004608) (9400,0.01978433421) (9500,0.02006420465) (9600,0.02094770866) (9700,0.02095367229) (9800,0.02332919603) (9900,0.02331599323) (10000,0.02519691149) (10100,0.02494370276) (10200,0.02497119708) (10300,0.02503452916) (10400,0.02459294312) (10500,0.02505965792) (10600,0.02508431431) (10700,0.02592581566) (10800,0.02579290513) (10900,0.02422042284) (11000,0.02291088849) (11100,0.02359216977) (11200,0.02557066083) (11300,0.02598911989) (11400,0.02613437101) (11500,0.0248552477) (11600,0.02448316775) (11700,0.02341365907) (11800,0.02317250092) (11900,0.02434774991) (12000,0.02328892723) (12100,0.02438054513) (12200,0.02338923346) (12300,0.02300751284) (12400,0.02443641461) (12500,0.02537421118) (12600,0.02584735136) (12700,0.02506197002) (12800,0.02462347876) (12900,0.02303317962) (13000,0.02322226483) (13100,0.02212862466) (13200,0.0201590755) (13300,0.02047695639) (13400,0.01889313245) (13500,0.021062556) (13600,0.02024800107) (13700,0.0205864517) (13800,0.01928112982) (13900,0.02043493912) (14000,0.02011093451) (14100,0.02106324183) (14200,0.02384437788) (14300,0.02560960706) (14400,0.02620359268) (14500,0.02351081334) (14600,0.02556925202) (14700,0.02521549221) (14800,0.02511584144) (14900,0.02407663465) (15000,0.02585552894) (15100,0.0241017119) (15200,0.02224779176) (15300,0.02056002365) (15400,0.01904924922) (15500,0.01937397085) (15600,0.01813746691) (15700,0.01986815268) (15800,0.02050915565) (15900,0.02277903426) (16000,0.02357594632) (16100,0.02508528261) (16200,0.02546807686) (16300,0.0256393888) (16400,0.02668931112) (16500,0.02692711111) (16600,0.02676571943) (16700,0.02610536348) (16800,0.02669248432) (16900,0.02495199461) (17000,0.02411498483) (17100,0.02322083898) (17200,0.0219105673) (17300,0.02279058145) (17400,0.02269837251) (17500,0.02218331108) (17600,0.02220015926) (17700,0.02242382346) (17800,0.02194128605) (17900,0.02098118858) (18000,0.02135700947) (18100,0.0210800265) (18200,0.02302060891) (18300,0.02365451101) (18400,0.02444579527) (18500,0.02399918865) (18600,0.02452456038) (18700,0.02420931943) (18800,0.02488998789) (18900,0.02516623996) (19000,0.02286562175) (19100,0.02483258229) (19200,0.02406794596) (19300,0.02176013235) (19400,0.0209180953) (19500,0.02136516832) (19600,0.01959435642) (19700,0.01899457816) (19800,0.01743049286) (19900,0.0175407256) (20000,0.01831984343) (20100,0.01686821086) (20200,0.01601175433) (20300,0.01700975848) (20400,0.01679330869) (20500,0.01735282904) (20600,0.01838183878) (20700,0.02068676604) (20800,0.02054659911) (20900,0.02091546049) (21000,0.02156664999) (21100,0.02180830212) (21200,0.02165800678) (21300,0.02169116447) (21400,0.02262471011) (21500,0.02263007415) (21600,0.02282823743) (21700,0.02066614823) (21800,0.0228217138) (21900,0.02308821939) (22000,0.02373900637) (22100,0.02361315899) (22200,0.02410782669) (22300,0.02439162545) (22400,0.02319387756) (22500,0.02357288562) (22600,0.02295601042) (22700,0.02447458068) (22800,0.0229690914) (22900,0.02235166142) (23000,0.02093040328) (23100,0.0206364397) (23200,0.02080386188) (23300,0.02022670824) (23400,0.02057516947) (23500,0.02081529126) (23600,0.02169081206) (23700,0.02138905199) (23800,0.02263509901)};
\addplot[pubRose, opacity=0.18, line width=0.35pt] coordinates {(100,0.01760025136) (200,0.01521840272) (300,0.01387741324) (400,0.01263901941) (500,0.01203181259) (600,0.01158185676) (700,0.01128724577) (800,0.01128562482) (900,0.01121785812) (1000,0.01132391393) (1100,0.01064655157) (1200,0.01011561672) (1300,0.00967325978) (1400,0.009313334292) (1500,0.008956756303) (1600,0.008604958095) (1700,0.008339580381) (1800,0.007825548528) (1900,0.007293934654) (2000,0.006596926041) (2100,0.005981725175) (2200,0.006014552107) (2300,0.005856689811) (2400,0.005777454423) (2500,0.005763980048) (2600,0.005438152491) (2700,0.005234105396) (2800,0.004913117713) (2900,0.004818465584) (3000,0.004632886313) (3100,0.004467683402) (3200,0.004483332229) (3300,0.004702244629) (3400,0.004623105424) (3500,0.004439938162) (3600,0.004609338357) (3700,0.004441625392) (3800,0.00530643533) (3900,0.005903785815) (4000,0.007469947427) (4100,0.007679166179) (4200,0.00762495352) (4300,0.007765374519) (4400,0.008443999826) (4500,0.0087076389) (4600,0.009616582119) (4700,0.01048071552) (4800,0.01044424362) (4900,0.0109204473) (5000,0.01056032656) (5100,0.0117817285) (5200,0.01219255603) (5300,0.01240396225) (5400,0.01262188009) (5500,0.01331980806) (5600,0.01306263851) (5700,0.01283962401) (5800,0.01292527765) (5900,0.0125236216) (6000,0.01195389424) (6100,0.01099150283) (6200,0.01067585442) (6300,0.01121359235) (6400,0.01191699365) (6500,0.01291089533) (6600,0.01303876201) (6700,0.01337475497) (6800,0.01387361111) (6900,0.01389857372) (7000,0.01507459702) (7100,0.01559342695) (7200,0.0163862112) (7300,0.01569310566) (7400,0.01493194439) (7500,0.01337368125) (7600,0.01399832861) (7700,0.01530477647) (7800,0.01443259669) (7900,0.01469335379) (8000,0.0139239328) (8100,0.0135495442) (8200,0.01465026615) (8300,0.01553662857) (8400,0.01727231182) (8500,0.01955009354) (8600,0.01899309149) (8700,0.0172187346) (8800,0.01842691842) (8900,0.0186734031) (9000,0.02106195968) (9100,0.02134092534) (9200,0.02226094818) (9300,0.02121756901) (9400,0.01986193717) (9500,0.01873873365) (9600,0.01864370904) (9700,0.01929446361) (9800,0.01810992393) (9900,0.01713587847) (10000,0.01472479478) (10100,0.01500727078) (10200,0.01378799221) (10300,0.01704017459) (10400,0.01871414143) (10500,0.01914332169) (10600,0.01923029064) (10700,0.02074523042) (10800,0.02209973275) (10900,0.02456826149) (11000,0.02463582167) (11100,0.02507423488) (11200,0.02544245543) (11300,0.02329250267) (11400,0.02195696933) (11500,0.02057239637) (11600,0.02171381628) (11700,0.02101053922) (11800,0.02072041361) (11900,0.02108432716) (12000,0.02183557833) (12100,0.02180726966) (12200,0.0205448281) (12300,0.02005595323) (12400,0.01951178284) (12500,0.02205602154) (12600,0.02109100558) (12700,0.02130226921) (12800,0.02109352443) (12900,0.02185975369) (13000,0.02231284231) (13100,0.0230763644) (13200,0.02360129664) (13300,0.02351049576) (13400,0.02794339126) (13500,0.02573083648) (13600,0.02749826713) (13700,0.02837165399) (13800,0.03031674875) (13900,0.02831610749) (14000,0.02737756362) (14100,0.02666921159) (14200,0.02661721604) (14300,0.0268998934) (14400,0.02349369861) (14500,0.02391372491) (14600,0.02287419625) (14700,0.02467947789) (14800,0.02669717595) (14900,0.02702437975) (15000,0.02658768622) (15100,0.02633980354) (15200,0.02815061184) (15300,0.03091370901) (15400,0.03124536136) (15500,0.03157949122) (15600,0.03102648733) (15700,0.0271867644) (15800,0.02342452481) (15900,0.02396894414) (16000,0.02471092092) (16100,0.02790250601) (16200,0.02657612981) (16300,0.02372542461) (16400,0.02341102595) (16500,0.02520860406) (16600,0.02542520193) (16700,0.02570161242) (16800,0.02490744675) (16900,0.02282960163) (17000,0.02230967246) (17100,0.01890841229) (17200,0.01887005186) (17300,0.01886025593) (17400,0.01897466332) (17500,0.01933358237) (17600,0.01898991484) (17700,0.01985333785) (17800,0.02099324344) (17900,0.02138453862) (18000,0.0213912569) (18100,0.02556650573) (18200,0.024531883) (18300,0.02670290517) (18400,0.02744249469) (18500,0.02523297751) (18600,0.02498803474) (18700,0.02467998993) (18800,0.02483079415) (18900,0.0249230586) (19000,0.02639173418) (19100,0.02479656003) (19200,0.02638076311) (19300,0.02612837991) (19400,0.02600996895) (19500,0.02518333262) (19600,0.02788597196) (19700,0.02788842935) (19800,0.02775253654) (19900,0.02808166035) (20000,0.02789584156) (20100,0.02572607733) (20200,0.02612694558) (20300,0.02474092674) (20400,0.0240505809) (20500,0.02492189985) (20600,0.02497659754) (20700,0.02723302729) (20800,0.02711587772) (20900,0.02947030757) (21000,0.02949193344) (21100,0.02881604256) (21200,0.02883268977) (21300,0.02920178762) (21400,0.02948039519) (21500,0.02987820683) (21600,0.02933292864) (21700,0.02748816097) (21800,0.02674653474) (21900,0.02464336734) (22000,0.02418434713) (22100,0.02568991138) (22200,0.02536344035) (22300,0.02624805765) (22400,0.02547686407) (22500,0.0246601426) (22600,0.02412038976) (22700,0.02531540273) (22800,0.02517716363) (22900,0.02533241306) (23000,0.02546073962) (23100,0.02580946051) (23200,0.02485007197) (23300,0.02488522604) (23400,0.02514071353) (23500,0.02545773853) (23600,0.02612697389) (23700,0.02600944843) (23800,0.02707361411)};
\addplot[pubRose, opacity=0.18, line width=0.35pt] coordinates {(100,0.02004911192) (200,0.0178157473) (300,0.0159836635) (400,0.01463018055) (500,0.01322510401) (600,0.01199514139) (700,0.01127262593) (800,0.01084001805) (900,0.01051202054) (1000,0.01043782439) (1100,0.009288880974) (1200,0.008794298023) (1300,0.008461065032) (1400,0.008146755863) (1500,0.008134840848) (1600,0.008047785517) (1700,0.007776540099) (1800,0.007436318323) (1900,0.007117809448) (2000,0.006631806772) (2100,0.006565167103) (2200,0.00615021619) (2300,0.005728477472) (2400,0.005841997126) (2500,0.006003890233) (2600,0.006208544597) (2700,0.006481942022) (2800,0.006808858551) (2900,0.007401895523) (3000,0.007824220974) (3100,0.007946991175) (3200,0.007879175758) (3300,0.007915877691) (3400,0.008802467911) (3500,0.008856513212) (3600,0.009156983532) (3700,0.009935682593) (3800,0.00957382624) (3900,0.01072313474) (4000,0.01040848438) (4100,0.01175769893) (4200,0.01247595279) (4300,0.0131672645) (4400,0.0123123466) (4500,0.01277453993) (4600,0.01267104056) (4700,0.0133094931) (4800,0.01498645139) (4900,0.01532305782) (5000,0.01588861011) (5100,0.01452796124) (5200,0.01507026702) (5300,0.01498866659) (5400,0.01561286058) (5500,0.01514357897) (5600,0.01615725178) (5700,0.01490772059) (5800,0.01642366042) (5900,0.01525826501) (6000,0.01535718068) (6100,0.01637096945) (6200,0.01734744683) (6300,0.01990180016) (6400,0.01942182668) (6500,0.02059571985) (6600,0.02084201239) (6700,0.02339133089) (6800,0.02407125449) (6900,0.02433431717) (7000,0.02578863883) (7100,0.02528222119) (7200,0.0272985653) (7300,0.02564632325) (7400,0.0260563245) (7500,0.02652185038) (7600,0.03004870582) (7700,0.02836525943) (7800,0.02906603683) (7900,0.02947063744) (8000,0.02932307404) (8100,0.03108986951) (8200,0.02792071328) (8300,0.02737664757) (8400,0.0302695795) (8500,0.0292456883) (8600,0.02749116262) (8700,0.0283607333) (8800,0.02608524105) (8900,0.02603468997) (9000,0.02574999677) (9100,0.02417364819) (9200,0.02588875992) (9300,0.02610515598) (9400,0.02332351767) (9500,0.02414991539) (9600,0.02300852481) (9700,0.02310683671) (9800,0.02236462887) (9900,0.021297104) (10000,0.02299603028) (10100,0.02356052725) (10200,0.02194844512) (10300,0.0214045262) (10400,0.02335925065) (10500,0.02447545771) (10600,0.02501190547) (10700,0.02547433432) (10800,0.02606777791) (10900,0.02977311024) (11000,0.0278632096) (11100,0.02725816695) (11200,0.02719704015) (11300,0.02776240222) (11400,0.02550017918) (11500,0.02316798037) (11600,0.0224208151) (11700,0.0213141623) (11800,0.0211730347) (11900,0.01909074681) (12000,0.01953560105) (12100,0.02037756881) (12200,0.02025642348) (12300,0.0214554348) (12400,0.02381820567) (12500,0.02411313271) (12600,0.02291464731) (12700,0.02311129924) (12800,0.02229082268) (12900,0.02228615023) (13000,0.02132035494) (13100,0.02048963588) (13200,0.02169270571) (13300,0.02250467502) (13400,0.02059472948) (13500,0.02085401071) (13600,0.02212320752) (13700,0.02235162724) (13800,0.02445873618) (13900,0.02403961308) (14000,0.02406092659) (14100,0.02439168338) (14200,0.02482600864) (14300,0.02215135433) (14400,0.02374183685) (14500,0.02435536273) (14600,0.02542363405) (14700,0.02802051138) (14800,0.02626095731) (14900,0.02698043659) (15000,0.02665995806) (15100,0.027061476) (15200,0.02735824343) (15300,0.03006904013) (15400,0.02957635429) (15500,0.03039172832) (15600,0.02804099526) (15700,0.02549227364) (15800,0.02552392855) (15900,0.02609516568) (16000,0.02723478265) (16100,0.03446264155) (16200,0.0344040513) (16300,0.03225464076) (16400,0.03312862013) (16500,0.03246467207) (16600,0.03476855326) (16700,0.033588346) (16800,0.03472773219) (16900,0.03333712937) (17000,0.03325405875) (17100,0.02749178195) (17200,0.02685020929) (17300,0.02808828438) (17400,0.02647021702) (17500,0.02768026842) (17600,0.02742443373) (17700,0.02777611017) (17800,0.02660855539) (17900,0.02715731096) (18000,0.0272400856) (18100,0.02550356071) (18200,0.02467535548) (18300,0.02731477227) (18400,0.02748559322) (18500,0.02520704418) (18600,0.02609692179) (18700,0.02688831389) (18800,0.02834175844) (18900,0.03364345208) (19000,0.03362965435) (19100,0.03420947865) (19200,0.03554234989) (19300,0.03296858594) (19400,0.03402688261) (19500,0.03557669166) (19600,0.0336062327) (19700,0.03438862897) (19800,0.03758109398) (19900,0.03236940224) (20000,0.03245939929) (20100,0.03363838345) (20200,0.03173516672) (20300,0.03082879912) (20400,0.03111364674) (20500,0.03098145891) (20600,0.0316653654) (20700,0.03218462169) (20800,0.03038644642) (20900,0.0307887977) (21000,0.03336316757) (21100,0.03264129981) (21200,0.03377345055) (21300,0.03367567919) (21400,0.03485804051) (21500,0.03784432393) (21600,0.03735387549) (21700,0.03676040284) (21800,0.03381021079) (21900,0.03434958179) (22000,0.03139548916) (22100,0.03238010127) (22200,0.03301115707) (22300,0.03500462659) (22400,0.0332488507) (22500,0.03041051961) (22600,0.03371771481) (22700,0.03373101968) (22800,0.03344510915) (22900,0.03331616642) (23000,0.03620296763) (23100,0.03598533357) (23200,0.03684438402) (23300,0.03749917978) (23400,0.0361833903) (23500,0.03499918329) (23600,0.03122467315) (23700,0.03016535332) (23800,0.03169054221)};
\addplot[pubRose, opacity=0.18, line width=0.35pt] coordinates {(100,0.009353153408) (200,0.01079320302) (300,0.01300185701) (400,0.01384710823) (500,0.01471693832) (600,0.01522058109) (700,0.0151947307) (800,0.01499959) (900,0.01465112354) (1000,0.01396428742) (1100,0.01387557476) (1200,0.01345884907) (1300,0.01250990178) (1400,0.01163304895) (1500,0.01067077192) (1600,0.00982871158) (1700,0.009461604338) (1800,0.008973382413) (1900,0.008785775304) (2000,0.008911175141) (2100,0.008868198516) (2200,0.008886542777) (2300,0.008316142275) (2400,0.007862130506) (2500,0.007459274447) (2600,0.006850318634) (2700,0.006024499051) (2800,0.005402918602) (2900,0.004601748777) (3000,0.003950895183) (3100,0.003526175558) (3200,0.002932938188) (3300,0.002919263043) (3400,0.003223421983) (3500,0.003307930799) (3600,0.005276158801) (3700,0.005905928928) (3800,0.00641518964) (3900,0.0075186298) (4000,0.007654548157) (4100,0.007701940485) (4200,0.007941105077) (4300,0.008130093804) (4400,0.008580797771) (4500,0.008495521406) (4600,0.007362519996) (4700,0.007000937685) (4800,0.008902756777) (4900,0.009586627223) (5000,0.009752896987) (5100,0.00965293271) (5200,0.01034177684) (5300,0.01090485805) (5400,0.01057533335) (5500,0.01049567836) (5600,0.01114715135) (5700,0.01169156346) (5800,0.009667420271) (5900,0.008020856255) (6000,0.007833735039) (6100,0.008258037432) (6200,0.007798888185) (6300,0.008980722097) (6400,0.009490661114) (6500,0.01062703973) (6600,0.01063327419) (6700,0.01142712731) (6800,0.01216187088) (6900,0.01367656186) (7000,0.01579969306) (7100,0.01669844133) (7200,0.01728029437) (7300,0.01629910851) (7400,0.01623536982) (7500,0.01574248374) (7600,0.01542990562) (7700,0.01560076401) (7800,0.01522053089) (7900,0.0149212339) (8000,0.01518021626) (8100,0.01482042) (8200,0.01578337355) (8300,0.01672440451) (8400,0.0171107091) (8500,0.01696114372) (8600,0.01618198827) (8700,0.01548365047) (8800,0.01525792573) (8900,0.01652385248) (9000,0.01562062008) (9100,0.01611525603) (9200,0.01511841547) (9300,0.01527007595) (9400,0.01549774613) (9500,0.01774064107) (9600,0.02073906912) (9700,0.02282521278) (9800,0.02392690424) (9900,0.02335959412) (10000,0.0232105704) (10100,0.02281412119) (10200,0.02510038419) (10300,0.02414349467) (10400,0.02331168177) (10500,0.02208982436) (10600,0.02011672342) (10700,0.01898299055) (10800,0.02110745767) (10900,0.0211490362) (11000,0.02101702495) (11100,0.02220974825) (11200,0.02037952226) (11300,0.02345198942) (11400,0.02635721248) (11500,0.02529675122) (11600,0.02670900561) (11700,0.02702984698) (11800,0.02443630677) (11900,0.02622335777) (12000,0.02630961966) (12100,0.02476067496) (12200,0.02587708188) (12300,0.02569825789) (12400,0.02268784251) (12500,0.02455844041) (12600,0.02341374457) (12700,0.02389220037) (12800,0.02324150437) (12900,0.02282328168) (13000,0.02293131975) (13100,0.02554547992) (13200,0.0237610966) (13300,0.02301251069) (13400,0.02507874491) (13500,0.02393930703) (13600,0.02402331159) (13700,0.02384011112) (13800,0.02481988212) (13900,0.02243630076) (14000,0.02156917574) (14100,0.01993778059) (14200,0.02040942349) (14300,0.01973710284) (14400,0.01773282299) (14500,0.01862888467) (14600,0.01896605771) (14700,0.01663245643) (14800,0.01547088684) (14900,0.01636007992) (15000,0.0168061431) (15100,0.01660045083) (15200,0.01639764044) (15300,0.01478273985) (15400,0.01666057385) (15500,0.01600369415) (15600,0.01504644067) (15700,0.01859669681) (15800,0.01905506914) (15900,0.01891004564) (16000,0.01984063927) (16100,0.01992656961) (16200,0.02055525323) (16300,0.02144153379) (16400,0.01974676931) (16500,0.02035891032) (16600,0.02005135249) (16700,0.01809263844) (16800,0.01862673005) (16900,0.01763855452) (17000,0.01722215852) (17100,0.01837262427) (17200,0.01751552224) (17300,0.01631189035) (17400,0.01917110737) (17500,0.0181148245) (17600,0.01880100882) (17700,0.01835325519) (17800,0.01784598809) (17900,0.01773417145) (18000,0.01765658464) (18100,0.015865173) (18200,0.01611999655) (18300,0.0169706516) (18400,0.01473759562) (18500,0.01453777552) (18600,0.01373443902) (18700,0.01315236632) (18800,0.0142998212) (18900,0.01529216925) (19000,0.01461828183) (19100,0.01448753728) (19200,0.01487830831) (19300,0.01500806129) (19400,0.01521983305) (19500,0.01585653853) (19600,0.0161890558) (19700,0.01623586621) (19800,0.01538894707) (19900,0.01419544592) (20000,0.01605170025) (20100,0.01700756587) (20200,0.01632881677) (20300,0.01559139742) (20400,0.01541705718) (20500,0.01457476942) (20600,0.0139786128) (20700,0.01372180441) (20800,0.0136604514) (20900,0.01412964668) (21000,0.01433789171) (21100,0.01425579544) (21200,0.01528572133) (21300,0.01652304074) (21400,0.01742148427) (21500,0.01731523713) (21600,0.01789495526) (21700,0.02057031309) (21800,0.02077564057) (21900,0.02171625113) (22000,0.02001590068) (22100,0.01984149711) (22200,0.01865219343) (22300,0.01719168415) (22400,0.01548279487) (22500,0.01759878453) (22600,0.01853453293) (22700,0.01655851677) (22800,0.01737281745) (22900,0.01839573784) (23000,0.01801745072) (23100,0.01803238485) (23200,0.01908617178) (23300,0.01989377551) (23400,0.01995347133) (23500,0.01791004976) (23600,0.01641435204) (23700,0.0170584877) (23800,0.01637355024)};
\addplot[pubRose, opacity=0.18, line width=0.35pt] coordinates {(100,0.01791929826) (200,0.01224304992) (300,0.01203562583) (400,0.01310894138) (500,0.01444133157) (600,0.01518212763) (700,0.0160560401) (800,0.01668501861) (900,0.01664483035) (1000,0.01617692015) (1100,0.01542573343) (1200,0.01556701995) (1300,0.01498816754) (1400,0.01391054811) (1500,0.01274854345) (1600,0.0116533285) (1700,0.01032985132) (1800,0.009123304952) (1900,0.008731274493) (2000,0.008471291978) (2100,0.008509893343) (2200,0.008504614793) (2300,0.008629223378) (2400,0.008984052204) (2500,0.009341088869) (2600,0.009460973833) (2700,0.009696064889) (2800,0.009506021021) (2900,0.00923662479) (3000,0.008982655965) (3100,0.008507169271) (3200,0.008464422077) (3300,0.008384990785) (3400,0.008290154394) (3500,0.007907185145) (3600,0.008377571311) (3700,0.007913990412) (3800,0.008344055479) (3900,0.008902245155) (4000,0.009424945526) (4100,0.01127686133) (4200,0.01124191745) (4300,0.01180426548) (4400,0.01162890396) (4500,0.01193893556) (4600,0.01132773808) (4700,0.01145252585) (4800,0.01213735798) (4900,0.01170395305) (5000,0.01125951852) (5100,0.01110903393) (5200,0.012065366) (5300,0.01463520015) (5400,0.01498039924) (5500,0.01468299003) (5600,0.01499861488) (5700,0.01642879373) (5800,0.01526698237) (5900,0.01498831855) (6000,0.0149942799) (6100,0.01481509754) (6200,0.01431553173) (6300,0.01350497431) (6400,0.01456698175) (6500,0.01554989279) (6600,0.01520680576) (6700,0.01435760432) (6800,0.01713006077) (6900,0.01714855107) (7000,0.0190616345) (7100,0.01894539129) (7200,0.01915446194) (7300,0.01755782776) (7400,0.01847041566) (7500,0.01888570897) (7600,0.02353664124) (7700,0.02608271614) (7800,0.02531083953) (7900,0.02666179212) (8000,0.02867229404) (8100,0.02889229199) (8200,0.02849978581) (8300,0.02793268068) (8400,0.02805187395) (8500,0.02856032616) (8600,0.02405200489) (8700,0.02233876456) (8800,0.0223940298) (8900,0.02389529739) (9000,0.02008779971) (9100,0.01892695026) (9200,0.01860055886) (9300,0.02184597403) (9400,0.02051677778) (9500,0.01861585658) (9600,0.01855380191) (9700,0.0192159866) (9800,0.02046026895) (9900,0.01869946895) (10000,0.01877308786) (10100,0.02005190384) (10200,0.0231371643) (10300,0.02102684602) (10400,0.02143739536) (10500,0.02577853072) (10600,0.02735367604) (10700,0.02603680827) (10800,0.02471959777) (10900,0.02471922208) (11000,0.0249902104) (11100,0.02859765319) (11200,0.02635109639) (11300,0.0254670267) (11400,0.02460179953) (11500,0.02288821889) (11600,0.0226303895) (11700,0.02247797111) (11800,0.02425422529) (11900,0.02361939037) (12000,0.02431437075) (12100,0.01942669814) (12200,0.02135079233) (12300,0.02141577322) (12400,0.02128667179) (12500,0.0211494172) (12600,0.02022430645) (12700,0.02285473691) (12800,0.02218154231) (12900,0.02452580081) (13000,0.02393501932) (13100,0.02571834475) (13200,0.02588843666) (13300,0.02637148248) (13400,0.02760232026) (13500,0.0283621476) (13600,0.0288434159) (13700,0.0277200738) (13800,0.0257674573) (13900,0.02402560581) (14000,0.02519422937) (14100,0.02419823986) (14200,0.02189961728) (14300,0.02219212074) (14400,0.0219721647) (14500,0.02169134002) (14600,0.02148466697) (14700,0.02085447321) (14800,0.02121023862) (14900,0.02064812323) (15000,0.01954444619) (15100,0.02126639476) (15200,0.02169432612) (15300,0.02230761638) (15400,0.02154012425) (15500,0.02076573754) (15600,0.02057035333) (15700,0.01987157892) (15800,0.0211684458) (15900,0.02166307718) (16000,0.02147813104) (16100,0.02061366495) (16200,0.02041396732) (16300,0.02015067851) (16400,0.02085423088) (16500,0.02100667851) (16600,0.02116768043) (16700,0.02201438835) (16800,0.02122822041) (16900,0.02082819091) (17000,0.02325372873) (17100,0.02161149085) (17200,0.02371783806) (17300,0.0236773557) (17400,0.02355943145) (17500,0.02355863536) (17600,0.02414653813) (17700,0.02432671906) (17800,0.02349877777) (17900,0.02476949776) (18000,0.023830646) (18100,0.02530175112) (18200,0.02449181154) (18300,0.02453374118) (18400,0.02422317062) (18500,0.02369970139) (18600,0.02327422947) (18700,0.02244741675) (18800,0.02280242797) (18900,0.02244609892) (19000,0.02269756086) (19100,0.02237277403) (19200,0.02155321222) (19300,0.02092378642) (19400,0.02269693352) (19500,0.02199883666) (19600,0.02284171004) (19700,0.02290291898) (19800,0.02178637758) (19900,0.02121735234) (20000,0.02004208546) (20100,0.02005391996) (20200,0.02297725789) (20300,0.02226367742) (20400,0.02065127697) (20500,0.02084940095) (20600,0.02187985908) (20700,0.02190034539) (20800,0.02208981486) (20900,0.02158618737) (21000,0.02187419832) (21100,0.02227654327) (21200,0.01898407768) (21300,0.01994769182) (21400,0.02000875007) (21500,0.02070326265) (21600,0.01998247355) (21700,0.02029108237) (21800,0.02253111741) (21900,0.02463841643) (22000,0.0245594833) (22100,0.0236106256) (22200,0.02408381216) (22300,0.02449395694) (22400,0.02416369934) (22500,0.02254047925) (22600,0.02289861059) (22700,0.02397418926) (22800,0.02190494137) (22900,0.02012507776) (23000,0.01991528561) (23100,0.02311387481) (23200,0.02440241827) (23300,0.02547428375) (23400,0.02570279529) (23500,0.02742198296) (23600,0.0269822605) (23700,0.0256419871) (23800,0.02610470466)};
\addplot[pubRose, opacity=0.18, line width=0.35pt] coordinates {(100,0.00778008718) (200,0.008432340808) (300,0.008538659662) (400,0.008087373106) (500,0.007800318021) (600,0.00754341778) (700,0.007312778051) (800,0.007166189724) (900,0.00679713767) (1000,0.00659885346) (1100,0.006556671066) (1200,0.00650907266) (1300,0.006574551435) (1400,0.006710297568) (1500,0.007018055581) (1600,0.007441445766) (1700,0.007711529732) (1800,0.00807365845) (1900,0.008576800069) (2000,0.008884862252) (2100,0.008971489314) (2200,0.008968260698) (2300,0.008626804547) (2400,0.008220982226) (2500,0.007690378232) (2600,0.007374702208) (2700,0.007717359159) (2800,0.008649498224) (2900,0.008634417225) (3000,0.009218019806) (3100,0.009473207407) (3200,0.01041411385) (3300,0.01050275369) (3400,0.01062755981) (3500,0.01166069345) (3600,0.01157472352) (3700,0.01147200544) (3800,0.01023116787) (3900,0.01064948463) (4000,0.01036423654) (4100,0.01040896648) (4200,0.009993528249) (4300,0.01065844586) (4400,0.01160256942) (4500,0.01139653833) (4600,0.01230031075) (4700,0.01324845143) (4800,0.01553555336) (4900,0.0157546632) (5000,0.01620247802) (5100,0.01841843473) (5200,0.01888397327) (5300,0.01937623993) (5400,0.01943972679) (5500,0.02053679172) (5600,0.02122188956) (5700,0.02053677859) (5800,0.01929312171) (5900,0.01980211046) (6000,0.02081076363) (6100,0.02154550003) (6200,0.02433726573) (6300,0.02442110954) (6400,0.0243341729) (6500,0.02321485355) (6600,0.02220345754) (6700,0.02198278317) (6800,0.024654523) (6900,0.02332667788) (7000,0.02213536059) (7100,0.01965372809) (7200,0.01782026445) (7300,0.01788305291) (7400,0.01778791207) (7500,0.0172933185) (7600,0.01826136461) (7700,0.01822482026) (7800,0.01524089254) (7900,0.01618135255) (8000,0.01669442244) (8100,0.01614309726) (8200,0.0145419823) (8300,0.01414180677) (8400,0.01383984601) (8500,0.01501352722) (8600,0.01572450222) (8700,0.01645626826) (8800,0.01584776128) (8900,0.01553244917) (9000,0.01519611878) (9100,0.01563698119) (9200,0.01780267712) (9300,0.01718444508) (9400,0.01687942189) (9500,0.01692986195) (9600,0.01601638147) (9700,0.01515085171) (9800,0.02144195228) (9900,0.02526324256) (10000,0.02469162182) (10100,0.02443699888) (10200,0.0217090196) (10300,0.02434309414) (10400,0.02473568483) (10500,0.0234255054) (10600,0.02210032912) (10700,0.02272616103) (10800,0.01645028591) (10900,0.01256015478) (11000,0.01281864308) (11100,0.01465060003) (11200,0.01544270511) (11300,0.01349727525) (11400,0.01851233817) (11500,0.01987491958) (11600,0.022321398) (11700,0.02499678079) (11800,0.02561891535) (11900,0.02629777603) (12000,0.02717250083) (12100,0.02610765826) (12200,0.02579601724) (12300,0.02593951337) (12400,0.02065097345) (12500,0.024982055) (12600,0.02553630555) (12700,0.02345405603) (12800,0.02361921538) (12900,0.02454398535) (13000,0.02782924585) (13100,0.02829817012) (13200,0.02941701068) (13300,0.029058332) (13400,0.03041276252) (13500,0.02570844097) (13600,0.02478508735) (13700,0.02389410846) (13800,0.0237667202) (13900,0.02209542803) (14000,0.0187547246) (14100,0.01876611281) (14200,0.01770820757) (14300,0.01841538772) (14400,0.01831321474) (14500,0.01840441767) (14600,0.01756518781) (14700,0.01879971242) (14800,0.01891261451) (14900,0.0191589796) (15000,0.01816978138) (15100,0.01727039572) (15200,0.01788943009) (15300,0.01710171597) (15400,0.0160481873) (15500,0.01646072147) (15600,0.01674181828) (15700,0.01622860665) (15800,0.01808679225) (15900,0.01803769255) (16000,0.01870412342) (16100,0.01846059114) (16200,0.01837191209) (16300,0.01857932676) (16400,0.01945295818) (16500,0.01875377772) (16600,0.01804436445) (16700,0.01756260293) (16800,0.01548948959) (16900,0.01573034748) (17000,0.01737882774) (17100,0.01822857112) (17200,0.02007655054) (17300,0.01971007688) (17400,0.01944681732) (17500,0.02050239965) (17600,0.02131637083) (17700,0.02273729201) (17800,0.02258324176) (17900,0.02789876331) (18000,0.02799545173) (18100,0.02748540733) (18200,0.02606627159) (18300,0.02741715005) (18400,0.02725213645) (18500,0.02758096764) (18600,0.02847358966) (18700,0.02746007172) (18800,0.02808251455) (18900,0.02367950976) (19000,0.02232512757) (19100,0.0221720472) (19200,0.0241816273) (19300,0.02516146041) (19400,0.02621145677) (19500,0.02554787751) (19600,0.02440843061) (19700,0.0250194652) (19800,0.02550941501) (19900,0.02353333095) (20000,0.02366584232) (20100,0.02527372064) (20200,0.02256385116) (20300,0.02120129867) (20400,0.02200385937) (20500,0.02393411798) (20600,0.02368512945) (20700,0.02345758313) (20800,0.02232495137) (20900,0.02463413374) (21000,0.02372649238) (21100,0.02233788539) (21200,0.02175615933) (21300,0.02226839513) (21400,0.02062714696) (21500,0.01805525357) (21600,0.01817378094) (21700,0.01902221506) (21800,0.01919058757) (21900,0.01762409844) (22000,0.01948731085) (22100,0.01926804325) (22200,0.01975661404) (22300,0.01852321355) (22400,0.01880342783) (22500,0.02196914516) (22600,0.02502944935) (22700,0.0259219693) (22800,0.02646806454) (22900,0.02644818723) (23000,0.02745946124) (23100,0.02818339765) (23200,0.03139221268) (23300,0.03121719994) (23400,0.03065053942) (23500,0.02851768536) (23600,0.025682477) (23700,0.02447639732) (23800,0.02462518392)};
\addplot[pubRose, opacity=0.18, line width=0.35pt] coordinates {(100,0.01744457148) (200,0.0141248377) (300,0.01258421162) (400,0.01088556973) (500,0.009531205427) (600,0.008871359595) (700,0.008479325633) (800,0.008384111512) (900,0.00836636245) (1000,0.008589344332) (1100,0.007732963888) (1200,0.00772155798) (1300,0.00766090597) (1400,0.007874651114) (1500,0.008201201167) (1600,0.008367265854) (1700,0.008525259607) (1800,0.008469661931) (1900,0.008465400292) (2000,0.008493755152) (2100,0.008384596463) (2200,0.008148002159) (2300,0.008247976378) (2400,0.008489729557) (2500,0.008369879518) (2600,0.008245303109) (2700,0.008129225625) (2800,0.008593398565) (2900,0.008459829213) (3000,0.008003355749) (3100,0.007867541397) (3200,0.007822761824) (3300,0.008104987675) (3400,0.007664763462) (3500,0.008013168722) (3600,0.009689999698) (3700,0.01026464887) (3800,0.01074714819) (3900,0.01110668574) (4000,0.01182964654) (4100,0.01233361368) (4200,0.01306779827) (4300,0.01237061918) (4400,0.01370217637) (4500,0.01360188588) (4600,0.01313276119) (4700,0.01353546544) (4800,0.01356727309) (4900,0.01335606263) (5000,0.01283064443) (5100,0.01330112075) (5200,0.01258589583) (5300,0.01271529556) (5400,0.01201815852) (5500,0.01219992614) (5600,0.01168959052) (5700,0.01165978522) (5800,0.01196058081) (5900,0.01204213868) (6000,0.01205407153) (6100,0.01193671082) (6200,0.01288908417) (6300,0.01470111469) (6400,0.01528009297) (6500,0.01585782785) (6600,0.01700698128) (6700,0.01645156052) (6800,0.01581246788) (6900,0.01683565611) (7000,0.01684993627) (7100,0.01691781618) (7200,0.0161334347) (7300,0.01467307881) (7400,0.01414352125) (7500,0.01420389405) (7600,0.01368108178) (7700,0.01374290586) (7800,0.01374302274) (7900,0.01260577003) (8000,0.01280289833) (8100,0.01349716149) (8200,0.01380041614) (8300,0.01442455444) (8400,0.01549409153) (8500,0.01468032394) (8600,0.0141473826) (8700,0.01367390342) (8800,0.01332712593) (8900,0.01393760126) (9000,0.01517021963) (9100,0.01587213716) (9200,0.01778325262) (9300,0.01716827648) (9400,0.01695613293) (9500,0.01852244064) (9600,0.0184008372) (9700,0.02045568069) (9800,0.02057444574) (9900,0.02174535915) (10000,0.02061128942) (10100,0.01951970058) (10200,0.018066487) (10300,0.01943420069) (10400,0.01895055147) (10500,0.01877972381) (10600,0.01934084389) (10700,0.01960975341) (10800,0.01985164694) (10900,0.01872208901) (11000,0.01890145522) (11100,0.01816810146) (11200,0.01893333513) (11300,0.01900000498) (11400,0.01925729848) (11500,0.01846842011) (11600,0.01765312105) (11700,0.01585627347) (11800,0.01680231476) (11900,0.0170392978) (12000,0.01724802461) (12100,0.01804461498) (12200,0.01766682751) (12300,0.0163000823) (12400,0.01676541353) (12500,0.0169574894) (12600,0.01900765756) (12700,0.01919905711) (12800,0.01835103408) (12900,0.01850725729) (13000,0.01898742057) (13100,0.01967251115) (13200,0.01952525415) (13300,0.02292582756) (13400,0.02239181502) (13500,0.02258027745) (13600,0.02270096494) (13700,0.02398718204) (13800,0.02652994078) (13900,0.02877286412) (14000,0.02858628556) (14100,0.02812832072) (14200,0.02709149234) (14300,0.02440797035) (14400,0.02578513026) (14500,0.0255207818) (14600,0.02448947821) (14700,0.02535583489) (14800,0.02335063759) (14900,0.02100429814) (15000,0.02294190824) (15100,0.02446695678) (15200,0.02707521226) (15300,0.02905162908) (15400,0.02720422326) (15500,0.02922875909) (15600,0.02995171389) (15700,0.02768507218) (15800,0.02882841909) (15900,0.02815618617) (16000,0.02796943998) (16100,0.02576048179) (16200,0.02667055754) (16300,0.02446170347) (16400,0.02429547254) (16500,0.02343331464) (16600,0.02296034116) (16700,0.02319001434) (16800,0.02264518505) (16900,0.02383465627) (17000,0.02324888287) (17100,0.02336914176) (17200,0.02166574588) (17300,0.02236420764) (17400,0.02413707851) (17500,0.02436732566) (17600,0.02431331882) (17700,0.02484742329) (17800,0.02365935873) (17900,0.0256654812) (18000,0.02459985483) (18100,0.02539652921) (18200,0.02416518601) (18300,0.02698590374) (18400,0.0259691854) (18500,0.0278557173) (18600,0.02846924709) (18700,0.02975650998) (18800,0.03029704662) (18900,0.02873372184) (19000,0.02942042286) (19100,0.02927674716) (19200,0.02991620619) (19300,0.02654859722) (19400,0.02847907748) (19500,0.02562470548) (19600,0.02520565167) (19700,0.02487050816) (19800,0.02578648794) (19900,0.02570115793) (20000,0.02625942864) (20100,0.02591934949) (20200,0.02630646657) (20300,0.02714918628) (20400,0.02646826096) (20500,0.0267266864) (20600,0.02593328338) (20700,0.02688689176) (20800,0.02678121459) (20900,0.0266317632) (21000,0.02671079524) (21100,0.02734754365) (21200,0.02735354342) (21300,0.02745745219) (21400,0.02743665054) (21500,0.0276334146) (21600,0.02841008119) (21700,0.02605593754) (21800,0.02478965651) (21900,0.02486615945) (22000,0.0246631952) (22100,0.02580145579) (22200,0.02488102671) (22300,0.02501530498) (22400,0.02446402404) (22500,0.02600596622) (22600,0.02633835264) (22700,0.02820601957) (22800,0.02931288183) (22900,0.03053821567) (23000,0.03031356093) (23100,0.02876600754) (23200,0.03058020025) (23300,0.03020004071) (23400,0.02884773072) (23500,0.02720144838) (23600,0.02734921388) (23700,0.02575501241) (23800,0.0251000274)};
\addplot[pubRose, opacity=0.18, line width=0.35pt] coordinates {(100,0.02356719039) (200,0.01743590832) (300,0.01504046004) (400,0.01382694859) (500,0.01336060986) (600,0.01282179713) (700,0.01268900252) (800,0.01281518734) (900,0.01296396057) (1000,0.01287097707) (1100,0.01135497456) (1200,0.01101766536) (1300,0.01085098423) (1400,0.01086552935) (1500,0.01076740529) (1600,0.0108463305) (1700,0.01084035775) (1800,0.01044515437) (1900,0.01017034622) (2000,0.009826302156) (2100,0.009690557513) (2200,0.009627477685) (2300,0.009590142081) (2400,0.009462740971) (2500,0.009282097546) (2600,0.009270451916) (2700,0.008668859163) (2800,0.008875252074) (2900,0.008634639392) (3000,0.009297046112) (3100,0.009189788392) (3200,0.008955710474) (3300,0.008569570864) (3400,0.009014978679) (3500,0.009113378963) (3600,0.008569952892) (3700,0.009917273512) (3800,0.009397479519) (3900,0.01208497342) (4000,0.01291615656) (4100,0.01463266332) (4200,0.01601565806) (4300,0.01701095956) (4400,0.01927472656) (4500,0.01891683456) (4600,0.01982558961) (4700,0.01984893787) (4800,0.02063941425) (4900,0.01772929421) (5000,0.01663999949) (5100,0.01892326958) (5200,0.01870977785) (5300,0.01945316419) (5400,0.01663964363) (5500,0.01769306222) (5600,0.01763686696) (5700,0.01825556001) (5800,0.01864028107) (5900,0.01907107793) (6000,0.02144810623) (6100,0.0187120405) (6200,0.01978567531) (6300,0.02269509127) (6400,0.02549534719) (6500,0.02486443603) (6600,0.0240647113) (6700,0.02366358582) (6800,0.0251175167) (6900,0.0281527169) (7000,0.02808836065) (7100,0.02952561248) (7200,0.03122178465) (7300,0.03229827881) (7400,0.029549699) (7500,0.03412104286) (7600,0.03524438338) (7700,0.03424321003) (7800,0.03208131716) (7900,0.02925184714) (8000,0.03132662373) (8100,0.02897237884) (8200,0.02710638223) (8300,0.02215652065) (8400,0.02452687025) (8500,0.02202937175) (8600,0.02345197182) (8700,0.02492154436) (8800,0.02516577886) (8900,0.02577988952) (9000,0.02169622816) (9100,0.02469692715) (9200,0.02460479327) (9300,0.03061326761) (9400,0.03034810703) (9500,0.02889673486) (9600,0.02654477004) (9700,0.02529982282) (9800,0.02684343094) (9900,0.02608891018) (10000,0.02583572306) (10100,0.02479814161) (10200,0.02602650523) (10300,0.02198175639) (10400,0.02294107638) (10500,0.02462034151) (10600,0.02604983579) (10700,0.02675635228) (10800,0.02540136138) (10900,0.02762933243) (11000,0.02887562308) (11100,0.02718275683) (11200,0.02506495835) (11300,0.02458527433) (11400,0.02307677222) (11500,0.02209813027) (11600,0.02352531916) (11700,0.02530987831) (11800,0.02613761751) (11900,0.0239117736) (12000,0.02457081322) (12100,0.02635294972) (12200,0.02631469956) (12300,0.02696233271) (12400,0.02603240982) (12500,0.02488932908) (12600,0.02327872254) (12700,0.0223189991) (12800,0.02228322104) (12900,0.02532936716) (13000,0.0232672831) (13100,0.02405381091) (13200,0.02447163351) (13300,0.02420131732) (13400,0.02426871061) (13500,0.02486510407) (13600,0.02455318533) (13700,0.02316270843) (13800,0.0225247642) (13900,0.02120956648) (14000,0.02132400423) (14100,0.01883924948) (14200,0.01798130637) (14300,0.01736284746) (14400,0.01731190523) (14500,0.01857775776) (14600,0.01907110466) (14700,0.02044654032) (14800,0.02038880652) (14900,0.01925275894) (15000,0.02011116836) (15100,0.02107272921) (15200,0.02151904991) (15300,0.02111947788) (15400,0.0224119829) (15500,0.02176418621) (15600,0.02173608933) (15700,0.020072868) (15800,0.02166240076) (15900,0.02182846433) (16000,0.02291719159) (16100,0.02237816667) (16200,0.02312491229) (16300,0.02255095523) (16400,0.02181483191) (16500,0.02109241989) (16600,0.02104949746) (16700,0.02154510682) (16800,0.0221754971) (16900,0.0237652055) (17000,0.02308787825) (17100,0.02343861172) (17200,0.02496107472) (17300,0.02586853616) (17400,0.02489486802) (17500,0.02475767117) (17600,0.02527556047) (17700,0.02562353462) (17800,0.02306625107) (17900,0.02115075896) (18000,0.0209287622) (18100,0.02239497853) (18200,0.02232387206) (18300,0.02148084128) (18400,0.02412459096) (18500,0.02679571668) (18600,0.02670697244) (18700,0.02683476107) (18800,0.02728107758) (18900,0.02868875097) (19000,0.02914741617) (19100,0.02732799184) (19200,0.0316581985) (19300,0.03146832641) (19400,0.03178102169) (19500,0.02985468227) (19600,0.02875310685) (19700,0.02812171951) (19800,0.0290110223) (19900,0.02888202909) (20000,0.02940762155) (20100,0.03158512227) (20200,0.02598209493) (20300,0.02611739319) (20400,0.02418514937) (20500,0.02341086287) (20600,0.02391496971) (20700,0.02469654139) (20800,0.02410843) (20900,0.02340855561) (21000,0.02273137476) (21100,0.02083037999) (21200,0.02061699927) (21300,0.02066364456) (21400,0.02021281347) (21500,0.01991626434) (21600,0.01970569305) (21700,0.01981032062) (21800,0.02181979045) (21900,0.02081121458) (22000,0.02096139761) (22100,0.02055049473) (22200,0.01928946646) (22300,0.01982162511) (22400,0.01932019033) (22500,0.02007838786) (22600,0.01956891902) (22700,0.01894927621) (22800,0.01614216939) (22900,0.01725257626) (23000,0.01643158933) (23100,0.01798203802) (23200,0.0192364811) (23300,0.01910429941) (23400,0.02026320957) (23500,0.02074667476) (23600,0.02336405013) (23700,0.02389824223) (23800,0.02373624658)};
\addplot[pubRose, opacity=0.18, line width=0.35pt] coordinates {(100,0.01255756523) (200,0.009655521018) (300,0.00854251394) (400,0.008444617153) (500,0.007777305972) (600,0.007506286105) (700,0.00750423755) (800,0.007386020035) (900,0.007495765264) (1000,0.007659587264) (1100,0.0075074655) (1200,0.007893806463) (1300,0.008372752229) (1400,0.008713138243) (1500,0.009436975187) (1600,0.01011077054) (1700,0.01064211465) (1800,0.01133205323) (1900,0.01153690834) (2000,0.01170353862) (2100,0.01154220495) (2200,0.0113579399) (2300,0.01073892633) (2400,0.009946536599) (2500,0.009281715518) (2600,0.008474764135) (2700,0.007637819741) (2800,0.006556866085) (2900,0.005901352013) (3000,0.005286243116) (3100,0.004877634184) (3200,0.004722193838) (3300,0.006603005971) (3400,0.007620705455) (3500,0.007779101771) (3600,0.008413373749) (3700,0.008932618587) (3800,0.009338297904) (3900,0.009669242008) (4000,0.009632496675) (4100,0.009839466773) (4200,0.009669143613) (4300,0.008022469096) (4400,0.007709961385) (4500,0.007417733734) (4600,0.007649994595) (4700,0.007360892138) (4800,0.007713961555) (4900,0.007874405477) (5000,0.009634338273) (5100,0.009916761238) (5200,0.01039223261) (5300,0.01054789741) (5400,0.01019734861) (5500,0.01075074486) (5600,0.01139507741) (5700,0.01162585272) (5800,0.01175476331) (5900,0.01158857262) (6000,0.01095884601) (6100,0.01064537927) (6200,0.01004579854) (6300,0.009761918802) (6400,0.009804841364) (6500,0.01120146322) (6600,0.01092865593) (6700,0.01077008387) (6800,0.01029466852) (6900,0.01157857277) (7000,0.01127033574) (7100,0.01150951022) (7200,0.01277799588) (7300,0.01371969231) (7400,0.01497813752) (7500,0.01385935806) (7600,0.01293813959) (7700,0.01289640046) (7800,0.01302043553) (7900,0.01261617653) (8000,0.01290554376) (8100,0.01291210381) (8200,0.01469283877) (8300,0.0148601451) (8400,0.013701565) (8500,0.01399720488) (8600,0.01414627405) (8700,0.01407886245) (8800,0.01474941778) (8900,0.01548200562) (9000,0.01441755262) (9100,0.0142536418) (9200,0.01192723601) (9300,0.01113127484) (9400,0.0112309054) (9500,0.01021162611) (9600,0.0116092809) (9700,0.01250492679) (9800,0.01244973959) (9900,0.01123823905) (10000,0.01290427567) (10100,0.01330542695) (10200,0.01307095289) (10300,0.01358945863) (10400,0.01372903185) (10500,0.01626664801) (10600,0.01461325181) (10700,0.01445377064) (10800,0.01403957834) (10900,0.01383674718) (11000,0.01294964696) (11100,0.01263299971) (11200,0.01234286437) (11300,0.01214917805) (11400,0.01273763273) (11500,0.01059027044) (11600,0.01128119542) (11700,0.01099794698) (11800,0.01175512434) (11900,0.0129616396) (12000,0.01255262936) (12100,0.01330258292) (12200,0.01430791067) (12300,0.01375131244) (12400,0.01345407292) (12500,0.01566967601) (12600,0.01516297143) (12700,0.0159518376) (12800,0.01525287544) (12900,0.01416716324) (13000,0.01469611442) (13100,0.01574789369) (13200,0.01480326639) (13300,0.01540313954) (13400,0.01501462017) (13500,0.01271564304) (13600,0.01243848298) (13700,0.01208012123) (13800,0.01208374267) (13900,0.01365415128) (14000,0.01382370312) (14100,0.01231520707) (14200,0.01389746587) (14300,0.01488514715) (14400,0.01501035825) (14500,0.01510465643) (14600,0.01538623804) (14700,0.01499746898) (14800,0.0157014939) (14900,0.01512278924) (15000,0.01553885471) (15100,0.0161637486) (15200,0.01486142045) (15300,0.01346497405) (15400,0.01449514646) (15500,0.0146069455) (15600,0.01490676245) (15700,0.0160310979) (15800,0.01600036984) (15900,0.01511589) (16000,0.01392487776) (16100,0.01335915257) (16200,0.01347879325) (16300,0.01348733823) (16400,0.01365481555) (16500,0.01572462884) (16600,0.01593617206) (16700,0.01581967999) (16800,0.01592509192) (16900,0.01553063462) (17000,0.01601327406) (17100,0.01583727896) (17200,0.01564364927) (17300,0.01625205223) (17400,0.01530282414) (17500,0.0139341577) (17600,0.01330447183) (17700,0.01321799518) (17800,0.01440377967) (17900,0.01496566595) (18000,0.0153732508) (18100,0.01665765052) (18200,0.01744929138) (18300,0.01845020899) (18400,0.01797756874) (18500,0.01798080527) (18600,0.01795818405) (18700,0.01859654812) (18800,0.01785467183) (18900,0.01815594379) (19000,0.01800642097) (19100,0.0166545813) (19200,0.01599768642) (19300,0.01454415694) (19400,0.01596334139) (19500,0.0175124635) (19600,0.01755994326) (19700,0.01577779725) (19800,0.01487979097) (19900,0.01595883491) (20000,0.01599371433) (20100,0.01927370746) (20200,0.01926272959) (20300,0.02232257016) (20400,0.02249845471) (20500,0.02001523413) (20600,0.02166452808) (20700,0.02262265719) (20800,0.02370692166) (20900,0.02208331982) (21000,0.02217163211) (21100,0.01873501278) (21200,0.02065593312) (21300,0.01899815397) (21400,0.01801650189) (21500,0.01834096219) (21600,0.0170652559) (21700,0.01668307669) (21800,0.01661190428) (21900,0.0167341426) (22000,0.01602785699) (22100,0.01766262976) (22200,0.01562683824) (22300,0.01370266536) (22400,0.01388436221) (22500,0.01504647546) (22600,0.01561764535) (22700,0.01517563686) (22800,0.01421262901) (22900,0.01633452792) (23000,0.01798385158) (23100,0.01741133835) (23200,0.01812977539) (23300,0.01945161987) (23400,0.01871601325) (23500,0.01824930366) (23600,0.01742111081) (23700,0.01865523299) (23800,0.01910212031)};
\addplot[pubRose, opacity=0.18, line width=0.35pt] coordinates {(100,0.02940732241) (200,0.02226070594) (300,0.0188272701) (400,0.01716098539) (500,0.01637107953) (600,0.01612190654) (700,0.0161480414) (800,0.01620494924) (900,0.0161625567) (1000,0.01592591358) (1100,0.01408700775) (1200,0.01361846905) (1300,0.01322126649) (1400,0.01263536345) (1500,0.01191113442) (1600,0.01107079801) (1700,0.01023631901) (1800,0.009270773875) (1900,0.008649120806) (2000,0.008232780872) (2100,0.007956633484) (2200,0.007629449572) (2300,0.007747427467) (2400,0.00788984173) (2500,0.008042784501) (2600,0.008023225237) (2700,0.007946799975) (2800,0.008496542927) (2900,0.008111992385) (3000,0.00771013815) (3100,0.007421364496) (3200,0.007342795748) (3300,0.007162419427) (3400,0.007596786087) (3500,0.007342780801) (3600,0.00739935711) (3700,0.007455718331) (3800,0.00682937433) (3900,0.006781685306) (4000,0.007379163196) (4100,0.00729167047) (4200,0.007196514728) (4300,0.007508639665) (4400,0.007503691735) (4500,0.007653531991) (4600,0.007836753409) (4700,0.007889939193) (4800,0.00796939861) (4900,0.008140073996) (5000,0.007518171379) (5100,0.00824895408) (5200,0.008546323515) (5300,0.008784434106) (5400,0.008472172543) (5500,0.008555320418) (5600,0.008733674744) (5700,0.008969217027) (5800,0.00899202791) (5900,0.008810854796) (6000,0.008766992949) (6100,0.00851404015) (6200,0.008736933582) (6300,0.008281331789) (6400,0.009596978407) (6500,0.01020206097) (6600,0.01040599537) (6700,0.01079171621) (6800,0.01067019487) (6900,0.01103262124) (7000,0.01125073442) (7100,0.01145328255) (7200,0.01122355117) (7300,0.01117702229) (7400,0.009848910384) (7500,0.009217585064) (7600,0.008749779034) (7700,0.009164684359) (7800,0.009873036481) (7900,0.009637154592) (8000,0.01070358218) (8100,0.01030227197) (8200,0.01035091393) (8300,0.01130829789) (8400,0.01299690558) (8500,0.01424093014) (8600,0.01420406299) (8700,0.01413519355) (8800,0.01400121227) (8900,0.01380127668) (9000,0.01379833482) (9100,0.0144816793) (9200,0.01465321989) (9300,0.01396110584) (9400,0.01238256632) (9500,0.01116368771) (9600,0.01277106181) (9700,0.01203467483) (9800,0.01436919319) (9900,0.01548169623) (10000,0.01563024255) (10100,0.01595967733) (10200,0.01673163357) (10300,0.01676508798) (10400,0.0175177325) (10500,0.01892447975) (10600,0.01739260461) (10700,0.01813296592) (10800,0.01623709006) (10900,0.01588551831) (11000,0.01642430853) (11100,0.01596646681) (11200,0.01600990873) (11300,0.01667013001) (11400,0.01711173002) (11500,0.01649494227) (11600,0.01839666441) (11700,0.01788036861) (11800,0.0174198282) (11900,0.01759150587) (12000,0.01668010913) (12100,0.01755878404) (12200,0.01730060596) (12300,0.01895847954) (12400,0.01791867549) (12500,0.02006928856) (12600,0.01843587356) (12700,0.01935586156) (12800,0.01908877669) (12900,0.01939307433) (13000,0.01849170481) (13100,0.01895577502) (13200,0.01938337446) (13300,0.01769311358) (13400,0.01750237052) (13500,0.01553952512) (13600,0.01648843051) (13700,0.01518096463) (13800,0.01528151552) (13900,0.01589107825) (14000,0.01847157301) (14100,0.01800698368) (14200,0.01661720104) (14300,0.01686120401) (14400,0.01779944715) (14500,0.01718337671) (14600,0.01661939924) (14700,0.01767419376) (14800,0.01821001903) (14900,0.01671646847) (15000,0.01431988128) (15100,0.01401559976) (15200,0.01506681554) (15300,0.01462179935) (15400,0.01394264214) (15500,0.01455924064) (15600,0.01461763717) (15700,0.01364377262) (15800,0.01282787067) (15900,0.01342088603) (16000,0.01432081102) (16100,0.0148205494) (16200,0.01679514493) (16300,0.01746069272) (16400,0.01830058857) (16500,0.01908035292) (16600,0.0190863742) (16700,0.01904423381) (16800,0.01942340797) (16900,0.02068784963) (17000,0.01995019317) (17100,0.01970231831) (17200,0.018053554) (17300,0.01736950278) (17400,0.01656186394) (17500,0.01607767437) (17600,0.01651156759) (17700,0.01641986966) (17800,0.01642399728) (17900,0.01482826145) (18000,0.01523301806) (18100,0.01424285537) (18200,0.01429856876) (18300,0.01481797425) (18400,0.01567381984) (18500,0.01616250025) (18600,0.01600002619) (18700,0.01712959316) (18800,0.01876416458) (18900,0.0186620296) (19000,0.0207205249) (19100,0.02117406214) (19200,0.0206161893) (19300,0.01960538942) (19400,0.01927569862) (19500,0.01861756295) (19600,0.01752303713) (19700,0.0185399205) (19800,0.01757989158) (19900,0.01830751817) (20000,0.01580899912) (20100,0.0165780691) (20200,0.01683263159) (20300,0.01881237333) (20400,0.01883371067) (20500,0.01962142126) (20600,0.02113916185) (20700,0.02067525964) (20800,0.02109556552) (20900,0.02038919441) (21000,0.0252990555) (21100,0.02413687417) (21200,0.02344266381) (21300,0.02163259024) (21400,0.02051859102) (21500,0.01910178387) (21600,0.01877625664) (21700,0.01807279163) (21800,0.01733216667) (21900,0.01804859103) (22000,0.01438537245) (22100,0.01415731735) (22200,0.01488890373) (22300,0.01549872789) (22400,0.01692448026) (22500,0.01959066205) (22600,0.02136327587) (22700,0.02143232543) (22800,0.02103636079) (22900,0.02130561657) (23000,0.02116205841) (23100,0.02147744615) (23200,0.02157339491) (23300,0.02174870837) (23400,0.0216605017) (23500,0.01921052421) (23600,0.01719773347) (23700,0.01726217968) (23800,0.01905119838)};
\addplot[pubRose, opacity=0.18, line width=0.35pt] coordinates {(100,0.0291455742) (200,0.02788423002) (300,0.02760507353) (400,0.02760548284) (500,0.02800499201) (600,0.02765160147) (700,0.02728807261) (800,0.02673189808) (900,0.02622708264) (1000,0.02567866407) (1100,0.02458595354) (1200,0.02357294187) (1300,0.02223701365) (1400,0.02067696787) (1500,0.01889546467) (1600,0.0174327042) (1700,0.01593884919) (1800,0.01465115929) (1900,0.01331442669) (2000,0.01228530826) (2100,0.01121896082) (2200,0.01061614598) (2300,0.009916563844) (2400,0.009595158463) (2500,0.009117651219) (2600,0.008299687784) (2700,0.007557271561) (2800,0.006984340074) (2900,0.006786847487) (3000,0.006232348783) (3100,0.00612650062) (3200,0.005587574467) (3300,0.005441901181) (3400,0.004770037625) (3500,0.004280626704) (3600,0.00424874546) (3700,0.004238365544) (3800,0.004299205914) (3900,0.004401851213) (4000,0.004489849694) (4100,0.004432305554) (4200,0.004798432812) (4300,0.005037603341) (4400,0.005561629403) (4500,0.006891350984) (4600,0.007076886995) (4700,0.007423223881) (4800,0.008553468157) (4900,0.008734986372) (5000,0.009628347401) (5100,0.01098534754) (5200,0.01143087894) (5300,0.0113292092) (5400,0.01155877211) (5500,0.01100626062) (5600,0.0124362363) (5700,0.01247951179) (5800,0.0124119963) (5900,0.01209063148) (6000,0.01230153767) (6100,0.01109981574) (6200,0.01191586144) (6300,0.01306548878) (6400,0.01353352745) (6500,0.01405227217) (6600,0.01357950424) (6700,0.01422425476) (6800,0.0133406152) (6900,0.01366905286) (7000,0.01342410599) (7100,0.0139396315) (7200,0.01297407788) (7300,0.01197565426) (7400,0.01137647848) (7500,0.01068566008) (7600,0.01107018017) (7700,0.01242515878) (7800,0.01383352336) (7900,0.01507443851) (8000,0.01458010729) (8100,0.01428526174) (8200,0.01435832186) (8300,0.01484260084) (8400,0.01657111058) (8500,0.01775035746) (8600,0.01723128306) (8700,0.01554352464) (8800,0.015000151) (8900,0.01374733085) (9000,0.0144200909) (9100,0.01467900053) (9200,0.01453709835) (9300,0.01470738919) (9400,0.01329464866) (9500,0.01246152436) (9600,0.0121611258) (9700,0.01222823523) (9800,0.0111758695) (9900,0.01138063809) (10000,0.01152726943) (10100,0.01132361917) (10200,0.0117204682) (10300,0.01136080986) (10400,0.01472491315) (10500,0.01580492584) (10600,0.01600921908) (10700,0.01735732825) (10800,0.01790262219) (10900,0.01987538151) (11000,0.0195511682) (11100,0.02045290573) (11200,0.02113267761) (11300,0.02086264621) (11400,0.01732129874) (11500,0.01695153425) (11600,0.01708080089) (11700,0.01572409766) (11800,0.01704987222) (11900,0.01456497051) (12000,0.01691228971) (12100,0.01576462481) (12200,0.01489483528) (12300,0.0152385097) (12400,0.01511223651) (12500,0.01457030615) (12600,0.01479767533) (12700,0.01559837302) (12800,0.01586735742) (12900,0.01712393202) (13000,0.01412313525) (13100,0.01525641624) (13200,0.0164861368) (13300,0.01878468357) (13400,0.01970941043) (13500,0.01994084734) (13600,0.02004663125) (13700,0.02135997526) (13800,0.01994882543) (13900,0.01979674716) (14000,0.02096484415) (14100,0.02099400945) (14200,0.02073596753) (14300,0.0198864447) (14400,0.02051849291) (14500,0.02304429263) (14600,0.0242183907) (14700,0.02411877774) (14800,0.02480732258) (14900,0.02420310229) (15000,0.02457929291) (15100,0.02403441658) (15200,0.02360880813) (15300,0.02241076725) (15400,0.02183701554) (15500,0.01901268074) (15600,0.01736308923) (15700,0.01515102256) (15800,0.01404423593) (15900,0.01509665782) (16000,0.01500501884) (16100,0.01481118621) (16200,0.01429712083) (16300,0.0153069431) (16400,0.01432967) (16500,0.01461878251) (16600,0.01447935179) (16700,0.01467197947) (16800,0.01623177016) (16900,0.0150692421) (17000,0.0141765276) (17100,0.01493052999) (17200,0.01453299532) (17300,0.01306288233) (17400,0.01396270441) (17500,0.01326280693) (17600,0.01370513961) (17700,0.01385480491) (17800,0.01344382679) (17900,0.01341628972) (18000,0.01315162079) (18100,0.01339332135) (18200,0.0133634571) (18300,0.01320771808) (18400,0.01297513954) (18500,0.01310022688) (18600,0.01287284717) (18700,0.01340326332) (18800,0.01308460198) (18900,0.01419227393) (19000,0.01411191514) (19100,0.01280101556) (19200,0.01385052912) (19300,0.0145195676) (19400,0.0161811701) (19500,0.01626969911) (19600,0.0168653762) (19700,0.01606251923) (19800,0.01576149641) (19900,0.01486985702) (20000,0.01449795137) (20100,0.01428609062) (20200,0.01298187771) (20300,0.01205479228) (20400,0.0106556511) (20500,0.01059016814) (20600,0.009498300171) (20700,0.009403596865) (20800,0.009179182956) (20900,0.009063461842) (21000,0.009625527542) (21100,0.01038034591) (21200,0.01086573461) (21300,0.01183203496) (21400,0.01154825641) (21500,0.01271591904) (21600,0.01376982611) (21700,0.01381203709) (21800,0.01375669856) (21900,0.01463947603) (22000,0.01436618902) (22100,0.01457294468) (22200,0.01425399045) (22300,0.01362861954) (22400,0.01302008582) (22500,0.01211930597) (22600,0.01153423591) (22700,0.01265474148) (22800,0.01405944219) (22900,0.01316505829) (23000,0.01405719128) (23100,0.0138450685) (23200,0.01370604127) (23300,0.01487905127) (23400,0.01590154287) (23500,0.01684184694) (23600,0.01649196274) (23700,0.01631107344) (23800,0.01541044493)};
\addplot[pubRose, opacity=0.18, line width=0.35pt] coordinates {(100,0.01237902883) (200,0.0111930836) (300,0.0119875862) (400,0.01263326756) (500,0.01286656596) (600,0.0132666508) (700,0.01344428038) (800,0.01372415491) (900,0.01391644238) (1000,0.01395679638) (1100,0.01398132853) (1200,0.01409698706) (1300,0.01395486286) (1400,0.0137133766) (1500,0.0134236319) (1600,0.01298160627) (1700,0.01239675796) (1800,0.01179863028) (1900,0.01114612082) (2000,0.01058599148) (2100,0.01010322613) (2200,0.009670035541) (2300,0.008980986569) (2400,0.008250085451) (2500,0.00741857898) (2600,0.006743000145) (2700,0.006345381844) (2800,0.005825059465) (2900,0.005577869923) (3000,0.005160908517) (3100,0.005276299617) (3200,0.005005922564) (3300,0.005039291945) (3400,0.005260343826) (3500,0.005952277267) (3600,0.006322921719) (3700,0.007080669887) (3800,0.007891500369) (3900,0.008243698301) (4000,0.008907093806) (4100,0.009033249179) (4200,0.01044548168) (4300,0.01128269662) (4400,0.01122310129) (4500,0.01270039878) (4600,0.01294759004) (4700,0.01368592731) (4800,0.01380728553) (4900,0.01451818948) (5000,0.01422699909) (5100,0.0147695432) (5200,0.01451935214) (5300,0.01386094373) (5400,0.01515727607) (5500,0.01449703784) (5600,0.01399629177) (5700,0.01264242479) (5800,0.01192254289) (5900,0.01166766002) (6000,0.01322554438) (6100,0.01346922382) (6200,0.01274167863) (6300,0.01345343832) (6400,0.01427023429) (6500,0.01399532109) (6600,0.01610703869) (6700,0.0183036116) (6800,0.01917402884) (6900,0.01903724447) (7000,0.01738670114) (7100,0.01709622936) (7200,0.01876704078) (7300,0.01903299773) (7400,0.0179777571) (7500,0.017538168) (7600,0.01609227881) (7700,0.01406451268) (7800,0.01404379113) (7900,0.01407215688) (8000,0.0174795093) (8100,0.01666588872) (8200,0.01547775962) (8300,0.01562409592) (8400,0.01522722417) (8500,0.0174286189) (8600,0.0185888344) (8700,0.01886121458) (8800,0.02040389986) (8900,0.02135349032) (9000,0.01880720467) (9100,0.01909276657) (9200,0.02101885388) (9300,0.02020722879) (9400,0.02009210866) (9500,0.01989574246) (9600,0.01943235584) (9700,0.01987393824) (9800,0.01917739036) (9900,0.01786761312) (10000,0.01811652733) (10100,0.01824457627) (10200,0.01623703837) (10300,0.01749067269) (10400,0.01779860333) (10500,0.01729285829) (10600,0.01721357666) (10700,0.01669521211) (10800,0.01662419299) (10900,0.01753806053) (11000,0.01674110601) (11100,0.01697805012) (11200,0.01705824668) (11300,0.01643360993) (11400,0.01652573347) (11500,0.01773847323) (11600,0.01805768162) (11700,0.01895235209) (11800,0.01881172219) (11900,0.01830319809) (12000,0.02021545647) (12100,0.01972934455) (12200,0.02094871085) (12300,0.02004724289) (12400,0.01971166837) (12500,0.01891820705) (12600,0.01859390377) (12700,0.01824022154) (12800,0.01819910719) (12900,0.01738910754) (13000,0.01564098629) (13100,0.01694433703) (13200,0.01619890737) (13300,0.01647066409) (13400,0.01737353941) (13500,0.01546702813) (13600,0.01532914173) (13700,0.01489441516) (13800,0.01699722474) (13900,0.01773191281) (14000,0.01859884253) (14100,0.01915858714) (14200,0.01927608224) (14300,0.02032356393) (14400,0.02009611782) (14500,0.02163423421) (14600,0.02224717988) (14700,0.02259167274) (14800,0.02049406366) (14900,0.02228356758) (15000,0.02211632105) (15100,0.02132723285) (15200,0.02120407028) (15300,0.02085986966) (15400,0.02163877757) (15500,0.02070057606) (15600,0.02038259367) (15700,0.02149033658) (15800,0.02147624176) (15900,0.02040887997) (16000,0.02031604368) (16100,0.02159079667) (16200,0.02086105086) (16300,0.02205273136) (16400,0.02117090467) (16500,0.02148171179) (16600,0.02128294725) (16700,0.01953868046) (16800,0.01878253417) (16900,0.01968957791) (17000,0.01917809714) (17100,0.01892004907) (17200,0.02044935357) (17300,0.0193481667) (17400,0.01910597328) (17500,0.01888128314) (17600,0.02043959107) (17700,0.02265285309) (17800,0.02308137072) (17900,0.0225697876) (18000,0.02316187751) (18100,0.02350615151) (18200,0.02313739229) (18300,0.02368046232) (18400,0.02413405143) (18500,0.02411838975) (18600,0.02232631659) (18700,0.02149148518) (18800,0.02156919828) (18900,0.02245478677) (19000,0.02174279289) (19100,0.02026778506) (19200,0.01984851481) (19300,0.02046466609) (19400,0.02035120381) (19500,0.02011777889) (19600,0.02148786159) (19700,0.02160351863) (19800,0.02172177499) (19900,0.01957884068) (20000,0.02063602386) (20100,0.02127192328) (20200,0.02033780515) (20300,0.01958010513) (20400,0.01942714658) (20500,0.02036717338) (20600,0.01930115735) (20700,0.02118846988) (20800,0.0203483806) (20900,0.02010101853) (21000,0.01938912384) (21100,0.01861417722) (21200,0.02108700247) (21300,0.02175950324) (21400,0.02130682617) (21500,0.02052879483) (21600,0.02064712793) (21700,0.01861009393) (21800,0.01837759353) (21900,0.01901910771) (22000,0.02148281587) (22100,0.02179341419) (22200,0.02165969526) (22300,0.0216202681) (22400,0.0219607899) (22500,0.02377597019) (22600,0.02332649389) (22700,0.02319562184) (22800,0.02424908513) (22900,0.02434485471) (23000,0.02191264508) (23100,0.02200849848) (23200,0.02334870445) (23300,0.02484977366) (23400,0.02524598753) (23500,0.02382968618) (23600,0.02521630935) (23700,0.02504449785) (23800,0.02870816067)};
\addplot[pubRose, opacity=0.18, line width=0.35pt] coordinates {(100,0.01430358924) (200,0.01392003894) (300,0.01341001193) (400,0.01291020936) (500,0.01289278194) (600,0.01348578651) (700,0.01413668293) (800,0.01461304072) (900,0.01488347269) (1000,0.01451685801) (1100,0.0139334254) (1200,0.01349976119) (1300,0.01319819) (1400,0.01310493592) (1500,0.01274976078) (1600,0.01202156283) (1700,0.01102128429) (1800,0.01000941163) (1900,0.008953370107) (2000,0.008684920939) (2100,0.008596298564) (2200,0.008405959606) (2300,0.008235313185) (2400,0.008072645217) (2500,0.007927475497) (2600,0.00766623863) (2700,0.007534701983) (2800,0.007614438748) (2900,0.007897160947) (3000,0.008255274408) (3100,0.008714115294) (3200,0.009936163155) (3300,0.01241904288) (3400,0.01307368153) (3500,0.01348677124) (3600,0.0135586143) (3700,0.01469464148) (3800,0.01498414711) (3900,0.01498228773) (4000,0.01555990172) (4100,0.01496329498) (4200,0.01453833752) (4300,0.01301002712) (4400,0.01247501303) (4500,0.01177083361) (4600,0.01267712917) (4700,0.01258446854) (4800,0.01212904076) (4900,0.01207787632) (5000,0.01149459626) (5100,0.01241820599) (5200,0.01206113356) (5300,0.01099583791) (5400,0.0109391307) (5500,0.01189121292) (5600,0.01276186653) (5700,0.0119973754) (5800,0.01277537886) (5900,0.01332496209) (6000,0.01300503528) (6100,0.01337208338) (6200,0.0130272422) (6300,0.01312210411) (6400,0.0134961172) (6500,0.01330915736) (6600,0.01232395908) (6700,0.01219220366) (6800,0.01264234129) (6900,0.01258026408) (7000,0.0124171963) (7100,0.01225435622) (7200,0.01206454136) (7300,0.01234138221) (7400,0.0121359278) (7500,0.01347561856) (7600,0.01435488225) (7700,0.01464788145) (7800,0.01380034457) (7900,0.01339551373) (8000,0.01527556372) (8100,0.01412282088) (8200,0.01541962773) (8300,0.01621833881) (8400,0.01638260279) (8500,0.01471595326) (8600,0.01352535374) (8700,0.01323851421) (8800,0.01359541174) (8900,0.01331756716) (9000,0.01172469943) (9100,0.0120493243) (9200,0.01082146917) (9300,0.009956545336) (9400,0.009666265408) (9500,0.01010038634) (9600,0.01007560189) (9700,0.01015364439) (9800,0.009944686247) (9900,0.009783172328) (10000,0.01034478322) (10100,0.01077228319) (10200,0.01283376124) (10300,0.01302858526) (10400,0.0135685984) (10500,0.01326138517) (10600,0.01351795951) (10700,0.01419101907) (10800,0.01402380234) (10900,0.01605208507) (11000,0.01573097184) (11100,0.01555005023) (11200,0.01526127523) (11300,0.01594431745) (11400,0.01587234121) (11500,0.01675152536) (11600,0.01661428288) (11700,0.01747670528) (11800,0.01803944437) (11900,0.01626690151) (12000,0.01688808762) (12100,0.01717172507) (12200,0.01610715147) (12300,0.01676388681) (12400,0.01614976013) (12500,0.0151454119) (12600,0.01593059208) (12700,0.01533880979) (12800,0.01459923619) (12900,0.01552246632) (13000,0.01625219723) (13100,0.01604784625) (13200,0.01527908845) (13300,0.01430988442) (13400,0.01520423731) (13500,0.0162312882) (13600,0.01724867932) (13700,0.01670276029) (13800,0.0174405532) (13900,0.01676223939) (14000,0.01559240371) (14100,0.01552427327) (14200,0.01691544438) (14300,0.01735839061) (14400,0.017731026) (14500,0.01689717313) (14600,0.01641147742) (14700,0.01664495943) (14800,0.01601385837) (14900,0.01761729959) (15000,0.01751076579) (15100,0.01793193771) (15200,0.01829248639) (15300,0.01901603108) (15400,0.01836349815) (15500,0.01936382381) (15600,0.01914243) (15700,0.01920764931) (15800,0.0195230905) (15900,0.01811409881) (16000,0.01838191729) (16100,0.01812485028) (16200,0.01685167002) (16300,0.01649933951) (16400,0.01659830678) (16500,0.01573847011) (16600,0.01556748971) (16700,0.01585598756) (16800,0.01697134301) (16900,0.01725109769) (17000,0.01709427768) (17100,0.01708315695) (17200,0.01778111588) (17300,0.01809448786) (17400,0.01851888318) (17500,0.01835838445) (17600,0.01760742124) (17700,0.01883017085) (17800,0.01834166609) (17900,0.01802438665) (18000,0.01759001575) (18100,0.01807245314) (18200,0.01790397745) (18300,0.01620344473) (18400,0.01553272484) (18500,0.01687807078) (18600,0.01796096163) (18700,0.01835822845) (18800,0.01856234493) (18900,0.02020789618) (19000,0.02269607307) (19100,0.02168119187) (19200,0.02170291794) (19300,0.02202696409) (19400,0.02284098845) (19500,0.02199128848) (19600,0.02127040103) (19700,0.01954444069) (19800,0.01918473728) (19900,0.02033712268) (20000,0.01850419026) (20100,0.01945860442) (20200,0.01917994637) (20300,0.01894655554) (20400,0.01857338166) (20500,0.01878687637) (20600,0.0190574999) (20700,0.01916427733) (20800,0.02020534417) (20900,0.01862331079) (21000,0.02005777573) (21100,0.02047925489) (21200,0.02202761555) (21300,0.02481382098) (21400,0.02541895844) (21500,0.0280875802) (21600,0.03097590357) (21700,0.03130459692) (21800,0.03108050637) (21900,0.03077363595) (22000,0.02997626755) (22100,0.02937572449) (22200,0.02784489356) (22300,0.02692559082) (22400,0.02725842372) (22500,0.02420270815) (22600,0.0227635812) (22700,0.02262626551) (22800,0.02164990939) (22900,0.02134347185) (23000,0.0210047299) (23100,0.02319427077) (23200,0.02393262219) (23300,0.02312928624) (23400,0.02210326474) (23500,0.0240774354) (23600,0.02211220935) (23700,0.02165551968) (23800,0.0225068301)};
\addplot[pubRose, opacity=0.18, line width=0.35pt] coordinates {(100,0.02612727135) (200,0.02594851982) (300,0.0256556732) (400,0.02456427878) (500,0.02326177955) (600,0.02125197618) (700,0.01959425544) (800,0.0184973143) (900,0.01800307559) (1000,0.01772585306) (1100,0.01672794949) (1200,0.01577354819) (1300,0.01441079564) (1400,0.01308260858) (1500,0.01224324834) (1600,0.01231821943) (1700,0.01266700272) (1800,0.0127284917) (1900,0.01261785207) (2000,0.01228435254) (2100,0.01169114318) (2200,0.01121110953) (2300,0.01118649123) (2400,0.01134567047) (2500,0.01108998568) (2600,0.01061384883) (2700,0.01007027086) (2800,0.00975055648) (2900,0.009402859118) (3000,0.008999559935) (3100,0.008559212042) (3200,0.008475747192) (3300,0.0084522611) (3400,0.008543634275) (3500,0.00938209733) (3600,0.00934803593) (3700,0.009180696262) (3800,0.009142629756) (3900,0.009424563358) (4000,0.009375744546) (4100,0.009783389699) (4200,0.01022230107) (4300,0.009876451455) (4400,0.009671609662) (4500,0.01033553667) (4600,0.01152713597) (4700,0.01215569666) (4800,0.01252091881) (4900,0.01283589751) (5000,0.01308195544) (5100,0.01453794325) (5200,0.01381984465) (5300,0.0142972922) (5400,0.01487330059) (5500,0.01324368124) (5600,0.01358934906) (5700,0.01335645397) (5800,0.01387926438) (5900,0.01290469547) (6000,0.01444797446) (6100,0.0132025966) (6200,0.01431473349) (6300,0.01431848682) (6400,0.01418413338) (6500,0.01496015172) (6600,0.01383778797) (6700,0.01418745834) (6800,0.01471634135) (6900,0.01537201805) (7000,0.01370089855) (7100,0.01400376614) (7200,0.01471554413) (7300,0.01406015577) (7400,0.01472186716) (7500,0.01487541562) (7600,0.01513442155) (7700,0.01514469329) (7800,0.01445857277) (7900,0.01594061684) (8000,0.01635057917) (8100,0.01625061147) (8200,0.01462777108) (8300,0.01490392778) (8400,0.01462311689) (8500,0.01542851804) (8600,0.01553528728) (8700,0.0157883076) (8800,0.01631322736) (8900,0.01514172954) (9000,0.01566169681) (9100,0.01525270892) (9200,0.0156700084) (9300,0.01550612273) (9400,0.01531324759) (9500,0.01532798931) (9600,0.01553328019) (9700,0.01526940987) (9800,0.01433644602) (9900,0.01472881185) (10000,0.01442505885) (10100,0.01600916293) (10200,0.01650447892) (10300,0.01669479758) (10400,0.01707721846) (10500,0.01658494482) (10600,0.01664869608) (10700,0.01680454984) (10800,0.01782298321) (10900,0.01746750483) (11000,0.01653624554) (11100,0.01580633209) (11200,0.01620735978) (11300,0.0162160384) (11400,0.01689310395) (11500,0.01707312451) (11600,0.01639311719) (11700,0.01684193523) (11800,0.01593524949) (11900,0.0154082458) (12000,0.01669304306) (12100,0.01636827067) (12200,0.01530138217) (12300,0.0156445181) (12400,0.01396209309) (12500,0.01352024944) (12600,0.01373799182) (12700,0.01282820432) (12800,0.01412117211) (12900,0.01521553863) (13000,0.01494219666) (13100,0.0158096686) (13200,0.01650790498) (13300,0.01679012058) (13400,0.0171117248) (13500,0.0184444081) (13600,0.01863635695) (13700,0.01855490217) (13800,0.01783740623) (13900,0.01712656328) (14000,0.01754466314) (14100,0.01893801894) (14200,0.01855929671) (14300,0.01814805577) (14400,0.01822551396) (14500,0.01703500003) (14600,0.01732549323) (14700,0.01758856997) (14800,0.01811257154) (14900,0.01780151417) (15000,0.01690157428) (15100,0.01475592777) (15200,0.01432986539) (15300,0.01489525437) (15400,0.01553466329) (15500,0.01513589127) (15600,0.01449185191) (15700,0.01672002459) (15800,0.01750981184) (15900,0.01791565996) (16000,0.01852410706) (16100,0.01806778302) (16200,0.01836693892) (16300,0.0196228073) (16400,0.01901271977) (16500,0.0200147206) (16600,0.01990806768) (16700,0.01872801529) (16800,0.0170822015) (16900,0.01864785329) (17000,0.01812008331) (17100,0.01780205891) (17200,0.01840288294) (17300,0.01797395879) (17400,0.01897752509) (17500,0.01840142496) (17600,0.01993287196) (17700,0.01948659187) (17800,0.01889133267) (17900,0.01775311455) (18000,0.01859659208) (18100,0.01959930472) (18200,0.0194637021) (18300,0.01861859336) (18400,0.01832303386) (18500,0.0179593225) (18600,0.01764571154) (18700,0.01823645746) (18800,0.01996945469) (18900,0.02039313177) (19000,0.01984371105) (19100,0.0193378509) (19200,0.02003662931) (19300,0.0198887445) (19400,0.01971029434) (19500,0.02062094668) (19600,0.02006831681) (19700,0.01987703731) (19800,0.01915164618) (19900,0.01794665745) (20000,0.01976625705) (20100,0.01997675234) (20200,0.01870448114) (20300,0.01826030267) (20400,0.01736618746) (20500,0.01703447551) (20600,0.0166945708) (20700,0.01546469303) (20800,0.01493720859) (20900,0.01549480744) (21000,0.0137647327) (21100,0.01350110546) (21200,0.01365698576) (21300,0.01537178801) (21400,0.01568753421) (21500,0.01598602626) (21600,0.01893175906) (21700,0.02018879307) (21800,0.02002975028) (21900,0.02021922097) (22000,0.01988145113) (22100,0.02069597691) (22200,0.02198269833) (22300,0.02118834872) (22400,0.02166183107) (22500,0.02120863125) (22600,0.01864886358) (22700,0.01839646325) (22800,0.0199194273) (22900,0.01966868024) (23000,0.02161841113) (23100,0.02051500902) (23200,0.01924395962) (23300,0.01814127909) (23400,0.0174795337) (23500,0.01726730783) (23600,0.01746445168) (23700,0.01745881271) (23800,0.01621598313)};
\addplot[pubRose, opacity=0.18, line width=0.35pt] coordinates {(100,0.02255468071) (200,0.02111336403) (300,0.01878693451) (400,0.01614492969) (500,0.0140385163) (600,0.01257801506) (700,0.01180170057) (800,0.01126282685) (900,0.01078910458) (1000,0.01044608885) (1100,0.008942545112) (1200,0.007824559882) (1300,0.007329372037) (1400,0.00737770861) (1500,0.007846806245) (1600,0.0083202851) (1700,0.008571381355) (1800,0.008824214153) (1900,0.009038686799) (2000,0.009340962674) (2100,0.009648298752) (2200,0.009688173048) (2300,0.009675305523) (2400,0.009558255924) (2500,0.009022794897) (2600,0.008700537682) (2700,0.008286700724) (2800,0.007916239416) (2900,0.008291440038) (3000,0.008019752428) (3100,0.007751273271) (3200,0.007214984763) (3300,0.006778016966) (3400,0.006423933804) (3500,0.007478099503) (3600,0.01028095982) (3700,0.01154258167) (3800,0.01151107987) (3900,0.01119506266) (4000,0.01146392706) (4100,0.01116253976) (4200,0.01189384023) (4300,0.01276698895) (4400,0.01272512916) (4500,0.01163310786) (4600,0.008491609385) (4700,0.007424601261) (4800,0.007541469857) (4900,0.008141302038) (5000,0.008286592457) (5100,0.008264809428) (5200,0.008107873937) (5300,0.00860020495) (5400,0.009345680848) (5500,0.009628811479) (5600,0.01003879034) (5700,0.009792963928) (5800,0.009554007463) (5900,0.008496581903) (6000,0.008631223952) (6100,0.008914137445) (6200,0.008536438923) (6300,0.007188709918) (6400,0.006856304267) (6500,0.006770351622) (6600,0.006621755473) (6700,0.00682458356) (6800,0.007159635099) (6900,0.007364143757) (7000,0.006939215539) (7100,0.006765971659) (7200,0.007424435066) (7300,0.007244652393) (7400,0.007816720498) (7500,0.008519792953) (7600,0.009103765409) (7700,0.008856197377) (7800,0.009207729367) (7900,0.009901427175) (8000,0.009665206238) (8100,0.009758294211) (8200,0.009003462153) (8300,0.009489482408) (8400,0.009867579443) (8500,0.009985576244) (8600,0.01014472391) (8700,0.01028922978) (8800,0.01081197043) (8900,0.01035673912) (9000,0.01123369909) (9100,0.0110206319) (9200,0.0123434505) (9300,0.01206537005) (9400,0.01193993227) (9500,0.01145935892) (9600,0.01064535268) (9700,0.01204573652) (9800,0.01093991078) (9900,0.01054227785) (10000,0.01019527228) (10100,0.01041688756) (10200,0.009262124822) (10300,0.00997798387) (10400,0.009592928924) (10500,0.009608720895) (10600,0.009591091191) (10700,0.01081193923) (10800,0.01110832831) (10900,0.01145588681) (11000,0.011317572) (11100,0.0124707554) (11200,0.01296930695) (11300,0.01291526547) (11400,0.01284268163) (11500,0.01292338409) (11600,0.01354800621) (11700,0.01164202951) (11800,0.01150637753) (11900,0.0121478267) (12000,0.01206692569) (12100,0.01076599453) (12200,0.01183484867) (12300,0.01127110966) (12400,0.01171859535) (12500,0.01201632447) (12600,0.01148978746) (12700,0.01139904964) (12800,0.01197320749) (12900,0.01312921732) (13000,0.01318251356) (13100,0.01335104164) (13200,0.01191517818) (13300,0.01298084445) (13400,0.01445449051) (13500,0.0140860959) (13600,0.01429848364) (13700,0.01391901439) (13800,0.01472468153) (13900,0.01338497056) (14000,0.01292436416) (14100,0.012958579) (14200,0.01436023908) (14300,0.01372316927) (14400,0.01135182693) (14500,0.01092818431) (14600,0.01127766497) (14700,0.01099353447) (14800,0.009779199539) (14900,0.009218295244) (15000,0.009963104036) (15100,0.009861268103) (15200,0.008641695697) (15300,0.00902235359) (15400,0.009941255907) (15500,0.01025746441) (15600,0.0103929807) (15700,0.01078164992) (15800,0.01056654677) (15900,0.01058069682) (16000,0.01058011753) (16100,0.01092189052) (16200,0.01068354319) (16300,0.009892722266) (16400,0.008931409474) (16500,0.009292729385) (16600,0.009198144171) (16700,0.01048220564) (16800,0.01075895247) (16900,0.01090966244) (17000,0.01164545151) (17100,0.01191497059) (17200,0.01298975335) (17300,0.01358938455) (17400,0.01359662488) (17500,0.01384543246) (17600,0.01360870088) (17700,0.01298512341) (17800,0.01328900065) (17900,0.013361724) (18000,0.01242899327) (18100,0.01200801246) (18200,0.01101806574) (18300,0.01023393222) (18400,0.01062156088) (18500,0.009638603451) (18600,0.009888132615) (18700,0.009646611242) (18800,0.009281004919) (18900,0.009215563396) (19000,0.009094134951) (19100,0.009269671841) (19200,0.009521536576) (19300,0.009934643283) (19400,0.009802567121) (19500,0.01060836697) (19600,0.0103995272) (19700,0.009886707179) (19800,0.01009971993) (19900,0.01027455963) (20000,0.0109264411) (20100,0.01110708062) (20200,0.01078456519) (20300,0.0116658519) (20400,0.01184070893) (20500,0.01253111996) (20600,0.01282928493) (20700,0.01353954272) (20800,0.01392694339) (20900,0.01421532929) (21000,0.01356264828) (21100,0.01331797484) (21200,0.01363321561) (21300,0.01326446449) (21400,0.01322942497) (21500,0.01278844932) (21600,0.01226583151) (21700,0.0123945049) (21800,0.01179679791) (21900,0.01127989898) (22000,0.01160700787) (22100,0.01154407179) (22200,0.01187420432) (22300,0.01193921054) (22400,0.01262861928) (22500,0.01300328122) (22600,0.01363976486) (22700,0.01314725643) (22800,0.01336892275) (22900,0.01380589269) (23000,0.01358758686) (23100,0.0141296153) (23200,0.01377212564) (23300,0.0130323139) (23400,0.01278553698) (23500,0.01189609403) (23600,0.01116164871) (23700,0.01131556625) (23800,0.0119091297)};
\addplot[pubRose, opacity=0.18, line width=0.35pt] coordinates {(100,0.01158707589) (200,0.01072907494) (300,0.01028953089) (400,0.009838163154) (500,0.009603624605) (600,0.009546967534) (700,0.009719193248) (800,0.009955816669) (900,0.01011783609) (1000,0.01023341417) (1100,0.01026296029) (1200,0.01038847733) (1300,0.0105417978) (1400,0.01100710537) (1500,0.01118280394) (1600,0.01122679897) (1700,0.01107532941) (1800,0.01072858982) (1900,0.01035196679) (2000,0.009920806112) (2100,0.009419898689) (2200,0.008987506246) (2300,0.008671367215) (2400,0.008126116637) (2500,0.007822004706) (2600,0.007632350316) (2700,0.007596552698) (2800,0.007482544892) (2900,0.007272142172) (3000,0.007247531554) (3100,0.007035036106) (3200,0.006682841736) (3300,0.00622066725) (3400,0.005916114897) (3500,0.006056800392) (3600,0.005736937514) (3700,0.005641974555) (3800,0.005674153753) (3900,0.006215279549) (4000,0.006150659313) (4100,0.006625674991) (4200,0.006759681809) (4300,0.007051572204) (4400,0.007412126614) (4500,0.007054858888) (4600,0.0073829656) (4700,0.007584099332) (4800,0.007267561695) (4900,0.006840573624) (5000,0.006666567083) (5100,0.006155204028) (5200,0.006230638595) (5300,0.00617671092) (5400,0.005708654108) (5500,0.00623551521) (5600,0.00579314928) (5700,0.005142848007) (5800,0.005546524469) (5900,0.005276973778) (6000,0.005645445827) (6100,0.005645881454) (6200,0.005840798002) (6300,0.006031175004) (6400,0.006132101151) (6500,0.005616790685) (6600,0.005758977984) (6700,0.0058416076) (6800,0.005478120176) (6900,0.005849148566) (7000,0.006239767699) (7100,0.006463280041) (7200,0.00617400012) (7300,0.006046990817) (7400,0.006202194095) (7500,0.006571461167) (7600,0.006566789467) (7700,0.00709401099) (7800,0.007161926571) (7900,0.00712782247) (8000,0.006572027598) (8100,0.006396372896) (8200,0.006425617076) (8300,0.006786993146) (8400,0.00718137268) (8500,0.006795520848) (8600,0.006909425464) (8700,0.006639027409) (8800,0.007012491347) (8900,0.006768561853) (9000,0.006821975717) (9100,0.006714875204) (9200,0.007653120765) (9300,0.007241874794) (9400,0.006856805878) (9500,0.007441694383) (9600,0.007760831155) (9700,0.007864286564) (9800,0.008277516998) (9900,0.009211064596) (10000,0.009271488432) (10100,0.01044937102) (10200,0.009959376231) (10300,0.01031552954) (10400,0.01019971226) (10500,0.00973013835) (10600,0.009605370741) (10700,0.009586733999) (10800,0.009104260197) (10900,0.008945075003) (11000,0.008842481999) (11100,0.00810288121) (11200,0.008610462537) (11300,0.008766321139) (11400,0.009069164563) (11500,0.009674113197) (11600,0.0102184881) (11700,0.01065674489) (11800,0.01073631756) (11900,0.01129454095) (12000,0.01493520718) (12100,0.01546228323) (12200,0.01562046222) (12300,0.01632095454) (12400,0.01655793348) (12500,0.01653288137) (12600,0.01683490863) (12700,0.01624455992) (12800,0.01673148889) (12900,0.01648741439) (13000,0.01327573247) (13100,0.01290837917) (13200,0.01252232343) (13300,0.01207048213) (13400,0.01453438569) (13500,0.01548199477) (13600,0.01497119488) (13700,0.01508487104) (13800,0.01437810892) (13900,0.01360976491) (14000,0.01509205997) (14100,0.01502044927) (14200,0.01506374711) (14300,0.01577035319) (14400,0.01517032217) (14500,0.0143197774) (14600,0.0146393653) (14700,0.01707143686) (14800,0.01929164641) (14900,0.0205967118) (15000,0.02046785494) (15100,0.02113314318) (15200,0.02146412209) (15300,0.02068617484) (15400,0.01942559564) (15500,0.01948868101) (15600,0.01903531644) (15700,0.01674615396) (15800,0.01496826503) (15900,0.01383346152) (16000,0.01228612829) (16100,0.01168069476) (16200,0.01143000294) (16300,0.01125741843) (16400,0.01056866795) (16500,0.01031636102) (16600,0.01202137833) (16700,0.01198350331) (16800,0.01403983478) (16900,0.01494403817) (17000,0.01601255666) (17100,0.01723752776) (17200,0.01805042922) (17300,0.01920910506) (17400,0.0204508272) (17500,0.02229402671) (17600,0.02188307596) (17700,0.02275992502) (17800,0.02230234556) (17900,0.02339121811) (18000,0.02532711979) (18100,0.02468299689) (18200,0.02480411222) (18300,0.02436368773) (18400,0.02332368558) (18500,0.02177078733) (18600,0.02405887423) (18700,0.02667909032) (18800,0.02633312577) (18900,0.02688490292) (19000,0.02471409021) (19100,0.02403965024) (19200,0.0249744569) (19300,0.02398985382) (19400,0.02368351845) (19500,0.02358421488) (19600,0.02350323452) (19700,0.02125062505) (19800,0.02086510109) (19900,0.01846151184) (20000,0.01875040811) (20100,0.0194553067) (20200,0.0186695748) (20300,0.01872290829) (20400,0.01881215321) (20500,0.01859293291) (20600,0.01560259713) (20700,0.01479608836) (20800,0.01487715403) (20900,0.01563294157) (21000,0.0149387477) (21100,0.01528835986) (21200,0.01411291519) (21300,0.01536933165) (21400,0.0174800179) (21500,0.01931739142) (21600,0.02203777293) (21700,0.02248031953) (21800,0.02203289578) (21900,0.02247059168) (22000,0.02411005432) (22100,0.02354468955) (22200,0.0239513034) (22300,0.02349042259) (22400,0.02217503767) (22500,0.02005252829) (22600,0.01891586529) (22700,0.01936130458) (22800,0.02064514859) (22900,0.01948254444) (23000,0.01752759535) (23100,0.01876874808) (23200,0.01845003534) (23300,0.01871382017) (23400,0.01807545098) (23500,0.01870192066) (23600,0.01798430979) (23700,0.01750725992) (23800,0.01747346278)};
\addplot[pubRose, opacity=0.18, line width=0.35pt] coordinates {(100,0.01939801872) (200,0.01779694203) (300,0.0157846141) (400,0.01386172953) (500,0.01263778368) (600,0.01172318038) (700,0.01083746579) (800,0.01008544036) (900,0.009661951361) (1000,0.009383246955) (1100,0.008298844285) (1200,0.007442748174) (1300,0.006863728724) (1400,0.006688974937) (1500,0.006500200322) (1600,0.006295424234) (1700,0.006126189465) (1800,0.00607821811) (1900,0.006058441848) (2000,0.006047440274) (2100,0.005793136498) (2200,0.005673972424) (2300,0.00579266306) (2400,0.005797674414) (2500,0.005922622001) (2600,0.005758023355) (2700,0.00579002439) (2800,0.006052121846) (2900,0.005907735182) (3000,0.005560026644) (3100,0.005430885963) (3200,0.005553254345) (3300,0.005227239523) (3400,0.00532678524) (3500,0.005369828502) (3600,0.005427437741) (3700,0.005484395754) (3800,0.005254685739) (3900,0.005187126808) (4000,0.005327221751) (4100,0.005521395011) (4200,0.005502954125) (4300,0.005505607324) (4400,0.005517345201) (4500,0.005784900254) (4600,0.005863821134) (4700,0.005883060489) (4800,0.005995989777) (4900,0.006870635133) (5000,0.007000672678) (5100,0.006841520546) (5200,0.006993400399) (5300,0.007762157079) (5400,0.007726160251) (5500,0.007742583193) (5600,0.008068570774) (5700,0.008135045366) (5800,0.008197439974) (5900,0.007468396099) (6000,0.007392616384) (6100,0.007944797538) (6200,0.008184171934) (6300,0.007676679268) (6400,0.008154110704) (6500,0.007874069689) (6600,0.007603449235) (6700,0.008411974134) (6800,0.008168656705) (6900,0.009154368239) (7000,0.009321700642) (7100,0.009198196279) (7200,0.009450630611) (7300,0.009589687595) (7400,0.01007332089) (7500,0.01050381465) (7600,0.01152600478) (7700,0.01135321092) (7800,0.01318568895) (7900,0.01226588571) (8000,0.01250113519) (8100,0.0128112182) (8200,0.01316845287) (8300,0.01341088414) (8400,0.0124260447) (8500,0.01291423822) (8600,0.01254800903) (8700,0.01279807715) (8800,0.01121592922) (8900,0.01136003211) (9000,0.01101409323) (9100,0.01045213696) (9200,0.009968200885) (9300,0.00972424224) (9400,0.01071568681) (9500,0.00992898643) (9600,0.009506055852) (9700,0.008885143278) (9800,0.009078950016) (9900,0.009318846883) (10000,0.009438828146) (10100,0.009536722302) (10200,0.008892327175) (10300,0.009367931075) (10400,0.008359957766) (10500,0.008180422708) (10600,0.008566615405) (10700,0.009741798276) (10800,0.009752938198) (10900,0.009838773077) (11000,0.009973221179) (11100,0.01152926413) (11200,0.01216181582) (11300,0.01275860742) (11400,0.01326519949) (11500,0.01394782802) (11600,0.01507712882) (11700,0.01465090681) (11800,0.01574591994) (11900,0.01512467493) (12000,0.0160741406) (12100,0.01478739162) (12200,0.01525534862) (12300,0.0145182163) (12400,0.0156333473) (12500,0.01532086744) (12600,0.01381227444) (12700,0.01313386583) (12800,0.01288404171) (12900,0.01296685082) (13000,0.01191166095) (13100,0.01167779141) (13200,0.01080480819) (13300,0.0106218155) (13400,0.009086875292) (13500,0.008657381497) (13600,0.008777183481) (13700,0.008753767423) (13800,0.009676871262) (13900,0.01067359592) (14000,0.01090055811) (14100,0.01227644184) (14200,0.0123475947) (14300,0.01225155485) (14400,0.01231580614) (14500,0.01266471781) (14600,0.01332639726) (14700,0.01466848236) (14800,0.01276797862) (14900,0.01396178403) (15000,0.01406699396) (15100,0.0137502668) (15200,0.01360037541) (15300,0.01386410985) (15400,0.01441058572) (15500,0.01462321589) (15600,0.01424118504) (15700,0.01413660571) (15800,0.01482319082) (15900,0.01405445519) (16000,0.01437104712) (16100,0.01397652975) (16200,0.01431681402) (16300,0.0145794238) (16400,0.01515519237) (16500,0.01473548207) (16600,0.01483760467) (16700,0.01427814392) (16800,0.01398442732) (16900,0.01269513615) (17000,0.01264501303) (17100,0.0129618512) (17200,0.01318377848) (17300,0.01276904806) (17400,0.01160211978) (17500,0.01285024355) (17600,0.01241894946) (17700,0.01169993975) (17800,0.01195597253) (17900,0.01355668101) (18000,0.01382907527) (18100,0.01337707741) (18200,0.01273396332) (18300,0.01314640399) (18400,0.01402370958) (18500,0.01288579954) (18600,0.01446174895) (18700,0.01446778979) (18800,0.0147827548) (18900,0.01374230944) (19000,0.01390091088) (19100,0.01395089282) (19200,0.0146881748) (19300,0.0153558501) (19400,0.01528555816) (19500,0.01614554012) (19600,0.01541404603) (19700,0.0155978146) (19800,0.01512116641) (19900,0.01464550379) (20000,0.0150972025) (20100,0.01444013426) (20200,0.01545329718) (20300,0.01493202206) (20400,0.01494781543) (20500,0.01390011176) (20600,0.01369697466) (20700,0.01437316895) (20800,0.01527547869) (20900,0.01639062846) (21000,0.01627439195) (21100,0.01679849024) (21200,0.01641645036) (21300,0.01601997414) (21400,0.01582246101) (21500,0.01726561636) (21600,0.01725614136) (21700,0.01673232159) (21800,0.01548768682) (21900,0.01462726789) (22000,0.01430511149) (22100,0.01532168426) (22200,0.01411354924) (22300,0.01523892405) (22400,0.01565484107) (22500,0.01486143516) (22600,0.01413579499) (22700,0.01363841747) (22800,0.01365946787) (22900,0.01491250112) (23000,0.0141097114) (23100,0.01310455897) (23200,0.01338043571) (23300,0.01243992406) (23400,0.01277520577) (23500,0.01318081529) (23600,0.0142214892) (23700,0.01598361265) (23800,0.01722411783)};
\addplot[pubRose, ultra thick] coordinates {(100,0.01775977481) (200,0.01467162021) (300,0.01364371258) (400,0.01346794498) (500,0.01305894298) (600,0.01269990609) (700,0.01224535155) (800,0.01205040608) (900,0.01209090935) (1000,0.0120974455) (1100,0.0115024725) (1200,0.011789425) (1300,0.01163319885) (1400,0.01132007716) (1500,0.01091482276) (1600,0.01078681415) (1700,0.01042832369) (1800,0.009640092752) (1900,0.008996028453) (2000,0.009126068908) (2100,0.009126979602) (2200,0.008977883472) (2300,0.008628013963) (2400,0.008235533838) (2500,0.007985129999) (2600,0.007824777893) (2700,0.007640846597) (2800,0.007467649749) (2900,0.007649528235) (3000,0.007767179562) (3100,0.007586318883) (3200,0.007582778786) (3300,0.007482992823) (3400,0.007642734458) (3500,0.007843143458) (3600,0.00849166332) (3700,0.009056657425) (3800,0.009367888712) (3900,0.009683327121) (4000,0.009614866623) (4100,0.01012421662) (4200,0.01033389138) (4300,0.01070116162) (4400,0.01071950311) (4500,0.0108660375) (4600,0.01080252265) (4700,0.01096662069) (4800,0.01098049011) (4900,0.0105323954) (5000,0.01051623174) (5100,0.01104719073) (5200,0.01174600625) (5300,0.01116252355) (5400,0.01124895141) (5500,0.01144873677) (5600,0.01206291341) (5700,0.01184446943) (5800,0.01194156185) (5900,0.01185489935) (6000,0.01200398288) (6100,0.0113704029) (6200,0.01178598283) (6300,0.01296138151) (6400,0.01351482233) (6500,0.01365223923) (6600,0.01330913312) (6700,0.01378110666) (6800,0.01360711316) (6900,0.01378756779) (7000,0.01438774779) (7100,0.01479859655) (7200,0.01542448942) (7300,0.01436661729) (7400,0.01482690577) (7500,0.01403162605) (7600,0.01417660543) (7700,0.01435619707) (7800,0.01393865724) (7900,0.01407282548) (8000,0.01425202005) (8100,0.01410512065) (8200,0.01458487669) (8300,0.01485137297) (8400,0.01492517053) (8500,0.01518983827) (8600,0.01555822068) (8700,0.01551358756) (8800,0.01512903837) (8900,0.01531186758) (9000,0.01518316921) (9100,0.01575455917) (9200,0.01649838733) (9300,0.0163371996) (9400,0.01618858401) (9500,0.01733525151) (9600,0.01720860933) (9700,0.01708907224) (9800,0.01837336479) (9900,0.01750174579) (10000,0.01687338494) (10100,0.0171268696) (10200,0.01661805625) (10300,0.01726542364) (10400,0.01825637238) (10500,0.01885210178) (10600,0.01831144763) (10700,0.01855797824) (10800,0.01783259662) (10900,0.01750278268) (11000,0.0170145489) (11100,0.01711627056) (11200,0.01663280323) (11300,0.01655186997) (11400,0.01721651438) (11500,0.01740579887) (11600,0.01772919139) (11700,0.01767853694) (11800,0.01772963628) (11900,0.01737162829) (12000,0.01811947497) (12100,0.01850506233) (12200,0.01846992876) (12300,0.01887453045) (12400,0.01819089386) (12500,0.01867581194) (12600,0.01853563765) (12700,0.01871963933) (12800,0.01827507063) (12900,0.01842067931) (13000,0.01789912056) (13100,0.01849966148) (13200,0.01860573228) (13300,0.01823889858) (13400,0.01841972562) (13500,0.01919262772) (13600,0.0193414941) (13700,0.01957067694) (13800,0.01855926802) (13900,0.01876432998) (14000,0.01867678356) (14100,0.01888863421) (14200,0.01827030154) (14300,0.01828172174) (14400,0.01801248055) (14500,0.01849108771) (14600,0.01826562276) (14700,0.01823695309) (14800,0.01910213046) (14900,0.01920586927) (15000,0.01935461406) (15100,0.02000374938) (15200,0.02014263719) (15300,0.02062309925) (15400,0.01903194962) (15500,0.01918825228) (15600,0.01824056485) (15700,0.0185982134) (15800,0.01928907982) (15900,0.01850946248) (16000,0.01895816708) (16100,0.01868633395) (16200,0.01844659727) (16300,0.01910106703) (16400,0.01923283897) (16500,0.01954753676) (16600,0.01949722094) (16700,0.01888612455) (16800,0.01794843948) (16900,0.01831325633) (17000,0.01824736185) (17100,0.01864051828) (17200,0.01894978373) (17300,0.01927863588) (17400,0.01905068327) (17500,0.01910743276) (17600,0.02014953387) (17700,0.02005981924) (17800,0.02118814043) (17900,0.02106597377) (18000,0.02114288583) (18100,0.02107953532) (18200,0.02107153619) (18300,0.02067126771) (18400,0.0223150406) (18500,0.02212115396) (18600,0.02219819762) (18700,0.02196945096) (18800,0.02195734531) (18900,0.02227926017) (19000,0.02251060032) (19100,0.02227241062) (19200,0.02238876154) (19300,0.02189354822) (19400,0.02276896099) (19500,0.02199506257) (19600,0.02153174588) (19700,0.02099863631) (19800,0.02129343804) (19900,0.01995798168) (20000,0.01925833258) (20100,0.01971767838) (20200,0.01935846466) (20300,0.01926333033) (20400,0.01913042862) (20500,0.01954364756) (20600,0.01968593728) (20700,0.02104532989) (20800,0.02044748985) (20900,0.02065232745) (21000,0.02106220019) (21100,0.02056614636) (21200,0.02063646619) (21300,0.0206751518) (21400,0.02017618334) (21500,0.01961682788) (21600,0.0198440833) (21700,0.02007689583) (21800,0.02068839502) (21900,0.02051521777) (22000,0.02048864914) (22100,0.02062323582) (22200,0.02097200658) (22300,0.02050498691) (22400,0.02094182572) (22500,0.0211275049) (22600,0.02166236122) (22700,0.02202929547) (22800,0.02134313509) (22900,0.02071534717) (23000,0.02096756659) (23100,0.02105694292) (23200,0.02118862839) (23300,0.02098770831) (23400,0.02111783559) (23500,0.02078098301) (23600,0.02165667443) (23700,0.02152228584) (23800,0.02187053217)};
\end{groupplot}
\end{tikzpicture}

\caption{Pre-training optimizer diagnostics for all 20 frozen seeds. Thin
transparent lines identify individual seeds and the dark line is the
cross-seed median; mutually exclusive fine-tuning selection branches are
excluded. Evidence class: training-prefix diagnostic.}
\label{fig:appendix_optimizer}
\end{figure}

\clearpage
\subsection*{D. Supplementary Causal and Post-Holdout Evidence}

Figure \ref{fig:appendix_causal_economics} places the principal representation
and training alternatives in turnover--CAGR space. The historical-only control
achieves the greatest realised growth but also considerably higher turnover,
whereas the compressed scenario-CVaR representation improves on the full raw
state with a much smaller implementation increase. These are explanatory
comparisons made after the primary holdout was consumed.

Figure \ref{fig:appendix_causal_seed_effects} prevents an aggregate contrast
from masking optimisation heterogeneity. The matched design compares the same
seed across the reference and alternative intervention, but the resulting
points still describe optimisation variability rather than new market
realisations.

\begin{figure}[!tbp]
\centering
\begin{minipage}[t]{0.56\textwidth}
\centering
% AUTO-GENERATED. DO NOT EDIT BY HAND.
% Figure: Causal CAGR-turnover comparison
% Evidence class: post_holdout_explanatory
% Regenerate with publication_pipeline_draft/generate_publication_tikz.py.
\begin{tikzpicture}
\begin{axis}[
  publication axis,
  width=0.62\linewidth,
  height=0.36\linewidth,
  xlabel={Mean monthly turnover}, ylabel={CAGR (\%)},
  title={Economic location of the principal causal representations},
  xmin=0.28566034,
  xmax=1.0173409,
  ymin=20.032756,
  ymax=49.268929,
  legend style={at={(1.04,0.50)}, anchor=west, legend columns=1,
    column sep=7pt, row sep=1pt},
]
\addplot[only marks, mark=square*, mark size=3pt, pubSlate] coordinates {(0.3466337236,25.7453442)};
\addlegendentry{1 Full state}
\node[font=\scriptsize\bfseries, text=pubSlate, anchor=south west, inner sep=3.2pt] at (axis cs:0.3466337236,25.7453442) {1};
\addplot[only marks, mark=diamond*, mark size=3pt, pubGold] coordinates {(0.6138921794,46.07122213)};
\addlegendentry{2 Historical only}
\node[font=\scriptsize\bfseries, text=pubGold, anchor=south west, inner sep=3.2pt] at (axis cs:0.6138921794,46.07122213) {2};
\addplot[only marks, mark=pentagon*, mark size=3pt, pubTeal] coordinates {(0.9563675568,26.53712051)};
\addlegendentry{3 Moving-block bootstrap}
\node[font=\scriptsize\bfseries, text=pubTeal, anchor=south west, inner sep=3.2pt] at (axis cs:0.9563675568,26.53712051) {3};
\addplot[only marks, mark=triangle*, mark size=3pt, pubRose] coordinates {(0.3752570548,23.23046204)};
\addlegendentry{4 No visible dependence}
\node[font=\scriptsize\bfseries, text=pubRose, anchor=south west, inner sep=3.2pt] at (axis cs:0.3752570548,23.23046204) {4};
\addplot[only marks, mark=*, mark size=3pt, pubNavy] coordinates {(0.3926317597,33.20850815)};
\addlegendentry{5 Compressed CVaR}
\node[font=\scriptsize\bfseries, text=pubNavy, anchor=south west, inner sep=3.2pt] at (axis cs:0.3926317597,33.20850815) {5};
\end{axis}
\end{tikzpicture}

\captionof{figure}{CAGR and mean monthly turnover of selected ten-seed causal
ensembles. Evidence class: post-holdout explanatory.}
\label{fig:appendix_causal_economics}
\end{minipage}\hfill
\begin{minipage}[t]{0.41\textwidth}
\centering
% AUTO-GENERATED. DO NOT EDIT BY HAND.
% Figure: Causal matched-seed heterogeneity
% Evidence class: post_holdout_explanatory
% Regenerate with publication_pipeline_draft/generate_publication_tikz.py.
\begin{tikzpicture}
\begin{axis}[
  publication axis,
  width=0.84\linewidth,
  height=0.27\linewidth,
  xlabel={Paired annual CRRA CE effect (percentage points)}, ylabel={},
  ytick={1,2,3}, yticklabels={Raw vine state,Joint visible dependence,Synthetic pre-training},
  y dir=reverse,
  ymin=0.85, ymax=3.15,
  enlarge y limits=false,
  title={Matched-seed heterogeneity in key mechanism effects},
]
\addplot[pubSlate, thin] coordinates {(0,0.5) (0,3.5)};
\addplot[only marks, mark=*, mark size=1.5pt, pubNavy!58] coordinates {(-25.28263736,1)};
\addplot[only marks, mark=*, mark size=1.5pt, pubNavy!58] coordinates {(-13.98786945,1)};
\addplot[only marks, mark=*, mark size=1.5pt, pubNavy!58] coordinates {(-13.72154551,1)};
\addplot[only marks, mark=*, mark size=1.5pt, pubNavy!58] coordinates {(-13.71597026,1)};
\addplot[only marks, mark=*, mark size=1.5pt, pubNavy!58] coordinates {(-9.971709805,1)};
\addplot[only marks, mark=*, mark size=1.5pt, pubNavy!58] coordinates {(-6.474043193,1)};
\addplot[only marks, mark=*, mark size=1.5pt, pubNavy!58] coordinates {(-4.502080986,1)};
\addplot[only marks, mark=*, mark size=1.5pt, pubNavy!58] coordinates {(1.844916942,1)};
\addplot[only marks, mark=*, mark size=1.5pt, pubNavy!58] coordinates {(4.512094507,1)};
\addplot[only marks, mark=*, mark size=1.5pt, pubNavy!58] coordinates {(8.058394782,1)};
\addplot[only marks, mark=diamond*, mark size=3pt, pubRose] coordinates {(-8.222876499,1)};
\addplot[only marks, mark=*, mark size=1.5pt, pubNavy!58] coordinates {(-21.63778853,2)};
\addplot[only marks, mark=*, mark size=1.5pt, pubNavy!58] coordinates {(-9.473124774,2)};
\addplot[only marks, mark=*, mark size=1.5pt, pubNavy!58] coordinates {(-7.312523557,2)};
\addplot[only marks, mark=*, mark size=1.5pt, pubNavy!58] coordinates {(-7.200112262,2)};
\addplot[only marks, mark=*, mark size=1.5pt, pubNavy!58] coordinates {(1.153245086,2)};
\addplot[only marks, mark=*, mark size=1.5pt, pubNavy!58] coordinates {(5.335583942,2)};
\addplot[only marks, mark=*, mark size=1.5pt, pubNavy!58] coordinates {(10.69941866,2)};
\addplot[only marks, mark=*, mark size=1.5pt, pubNavy!58] coordinates {(12.04604507,2)};
\addplot[only marks, mark=*, mark size=1.5pt, pubNavy!58] coordinates {(15.3786761,2)};
\addplot[only marks, mark=*, mark size=1.5pt, pubNavy!58] coordinates {(22.10299076,2)};
\addplot[only marks, mark=diamond*, mark size=3pt, pubRose] coordinates {(3.244414514,2)};
\addplot[only marks, mark=*, mark size=1.5pt, pubNavy!58] coordinates {(-62.00833121,3)};
\addplot[only marks, mark=*, mark size=1.5pt, pubNavy!58] coordinates {(-43.13682314,3)};
\addplot[only marks, mark=*, mark size=1.5pt, pubNavy!58] coordinates {(-38.2723466,3)};
\addplot[only marks, mark=*, mark size=1.5pt, pubNavy!58] coordinates {(-32.85823607,3)};
\addplot[only marks, mark=*, mark size=1.5pt, pubNavy!58] coordinates {(-17.88618445,3)};
\addplot[only marks, mark=*, mark size=1.5pt, pubNavy!58] coordinates {(-14.53751475,3)};
\addplot[only marks, mark=*, mark size=1.5pt, pubNavy!58] coordinates {(-6.670537349,3)};
\addplot[only marks, mark=*, mark size=1.5pt, pubNavy!58] coordinates {(3.103687324,3)};
\addplot[only marks, mark=*, mark size=1.5pt, pubNavy!58] coordinates {(4.993725472,3)};
\addplot[only marks, mark=*, mark size=1.5pt, pubNavy!58] coordinates {(13.7873646,3)};
\addplot[only marks, mark=diamond*, mark size=3pt, pubRose] coordinates {(-16.2118496,3)};
\end{axis}
\end{tikzpicture}

\captionof{figure}{Matched-seed annual CRRA certainty-equivalent effects.
Blue points are individual seeds and red diamonds are medians. Evidence class:
post-holdout explanatory.}
\label{fig:appendix_causal_seed_effects}
\end{minipage}
\end{figure}

\begin{figure}[!tbp]
\centering
% AUTO-GENERATED. DO NOT EDIT BY HAND.
% Figure: Causal ensemble wealth
% Evidence class: post_holdout_explanatory
% Regenerate with publication_pipeline_draft/generate_publication_tikz.py.
\begin{tikzpicture}
\begin{axis}[
  publication axis,
  width=0.88\linewidth,
  height=0.34\linewidth,
  date coordinates in=x,
  xticklabel={\year--\month},
  xticklabel style={rotate=30, anchor=north east},
  xlabel={Holding-period end}, ylabel={Net wealth index},
  title={Causal representation ensembles on the common realized path},
  legend pos=north west,
]
\addplot[pubSlate, dashed, thick] coordinates {(2024-08-30,1.014023535) (2024-09-30,1.040576042) (2024-10-31,1.05987806) (2024-11-29,1.084681031) (2024-12-31,1.07545241) (2025-01-27,1.11163551) (2025-02-28,1.114612456) (2025-03-31,1.105792829) (2025-04-30,1.110626938) (2025-05-30,1.160201257) (2025-06-30,1.198027448) (2025-07-31,1.218014544) (2025-08-29,1.268052645) (2025-09-30,1.396662225) (2025-10-31,1.438349182) (2025-11-28,1.442124126) (2025-12-31,1.45384133) (2026-01-30,1.53240183) (2026-02-27,1.520893441) (2026-03-31,1.430508006) (2026-04-30,1.551238249) (2026-05-29,1.582375951) (2026-06-30,1.512264295) (2026-07-06,1.526150995)};
\addlegendentry{Full raw state}
\addplot[pubNavy, ultra thick] coordinates {(2024-08-30,1.014527007) (2024-09-30,1.046080476) (2024-10-31,1.066911835) (2024-11-29,1.093585952) (2024-12-31,1.083273728) (2025-01-27,1.122974343) (2025-02-28,1.131142688) (2025-03-31,1.116388673) (2025-04-30,1.135363127) (2025-05-30,1.193000131) (2025-06-30,1.224959109) (2025-07-31,1.255715656) (2025-08-29,1.337422789) (2025-09-30,1.48777045) (2025-10-31,1.533848045) (2025-11-28,1.533463708) (2025-12-31,1.541783334) (2026-01-30,1.642361112) (2026-02-27,1.618376314) (2026-03-31,1.509699531) (2026-04-30,1.654305458) (2026-05-29,1.722506402) (2026-06-30,1.669340635) (2026-07-06,1.659007044)};
\addlegendentry{Compressed scenario-CVaR}
\addplot[pubRose, dashdotted, thick] coordinates {(2024-08-30,1.012215631) (2024-09-30,1.038564124) (2024-10-31,1.028429312) (2024-11-29,1.063215289) (2024-12-31,1.057543115) (2025-01-27,1.080666892) (2025-02-28,1.077192644) (2025-03-31,1.052030919) (2025-04-30,1.06296054) (2025-05-30,1.114790855) (2025-06-30,1.14614399) (2025-07-31,1.176936733) (2025-08-29,1.227108451) (2025-09-30,1.327626874) (2025-10-31,1.368097189) (2025-11-28,1.366376452) (2025-12-31,1.370924896) (2026-01-30,1.440841827) (2026-02-27,1.419959834) (2026-03-31,1.333531538) (2026-04-30,1.456513328) (2026-05-29,1.501662739) (2026-06-30,1.446877918) (2026-07-06,1.458743252)};
\addlegendentry{No visible dependence}
\addplot[pubGold, densely dashed, thick] coordinates {(2024-08-30,0.983078679) (2024-09-30,1.146384863) (2024-10-31,1.158517309) (2024-11-29,1.160965476) (2024-12-31,1.175079889) (2025-01-27,1.195975117) (2025-02-28,1.221680278) (2025-03-31,1.240474026) (2025-04-30,1.243748472) (2025-05-30,1.303367027) (2025-06-30,1.346089661) (2025-07-31,1.397936194) (2025-08-29,1.537655777) (2025-09-30,1.672990287) (2025-10-31,1.741963644) (2025-11-28,1.755236317) (2025-12-31,1.773784186) (2026-01-30,1.86553589) (2026-02-27,1.849250189) (2026-03-31,1.751936532) (2026-04-30,1.938231479) (2026-05-29,2.020560512) (2026-06-30,1.989392725) (2026-07-06,1.973498237)};
\addlegendentry{Historical-only diagnostic}
\end{axis}
\end{tikzpicture}

\caption{Net wealth paths of selected causal ensembles on the common realised
asset-return sequence under common drift-aware costs. The historical-only path
is a diagnostic alternative selected within a post-holdout experiment and is
not a new confirmatory strategy. Evidence class: post-holdout explanatory.}
\label{fig:appendix_causal_wealth}
\end{figure}

Finally, Figure \ref{fig:appendix_compressed_reconciliation} compares the
compressed scenario-CVaR candidate with each frozen financial benchmark. The
positive point estimates describe economic magnitude, but the simultaneous
inference does not establish a fresh benchmark-superiority claim because the
candidate was emphasised after inspection of the consumed holdout.

\begin{figure}[!tbp]
\centering
% AUTO-GENERATED. DO NOT EDIT BY HAND.
% Figure: Compressed-vine benchmark reconciliation
% Evidence class: post_holdout_exploratory
% Regenerate with publication_pipeline_draft/generate_publication_tikz.py.
\begin{tikzpicture}
\begin{axis}[
  publication axis,
  width=0.84\linewidth,
  height=0.39\linewidth,
  xlabel={Annual CRRA CE difference (percentage points)}, ylabel={},
  ytick={1,2,3,4,5,6},
  yticklabels={Equal weight,Shrinkage MV,DCC--GARCH,Static vine,Rolling vine,Dynamic NN-vine},
  y dir=reverse,
  title={Compressed scenario-CVaR policy minus frozen benchmarks},
  xmin=-17.967347,
  xmax=33.612028,
]
\addplot[pubSlate, thin] coordinates {(0,0.5) (0,6.5)};
\addplot[pubNavy, thick, mark=|, mark options={solid,scale=1.2}] coordinates {(-5.034492627,1) (20.90161079,1)};
\addplot[publication point, pubNavy, fill=pubNavy] coordinates {(8.768120217,1)};
\addplot[pubNavy, thick, mark=|, mark options={solid,scale=1.2}] coordinates {(0.1350626719,2) (19.43811232,2)};
\addplot[publication point, pubNavy, fill=pubNavy] coordinates {(9.171694683,2)};
\addplot[pubNavy, thick, mark=|, mark options={solid,scale=1.2}] coordinates {(-2.072900159,3) (13.12460148,3)};
\addplot[publication point, pubNavy, fill=pubNavy] coordinates {(4.545264459,3)};
\addplot[pubNavy, thick, mark=|, mark options={solid,scale=1.2}] coordinates {(-14.44150795,4) (16.5669172,4)};
\addplot[publication point, pubNavy, fill=pubNavy] coordinates {(2.622782245,4)};
\addplot[pubNavy, thick, mark=|, mark options={solid,scale=1.2}] coordinates {(-7.79189509,5) (30.87329967,5)};
\addplot[publication point, pubNavy, fill=pubNavy] coordinates {(13.60229458,5)};
\addplot[pubNavy, thick, mark=|, mark options={solid,scale=1.2}] coordinates {(-14.77216461,6) (18.92737313,6)};
\addplot[publication point, pubNavy, fill=pubNavy] coordinates {(3.871338016,6)};
\end{axis}
\end{tikzpicture}

\caption{Annual CRRA certainty-equivalent differences between the post-hoc
compressed scenario-CVaR candidate and each frozen financial benchmark. Points
are paired estimates and horizontal segments are 95\% circular moving-block
bootstrap intervals; none establishes superiority after Holm correction.
Evidence class: post-holdout exploratory.}
\label{fig:appendix_compressed_reconciliation}
\end{figure}
